% 本文件由 LaTeX 排版 Agent 根据有效 translation.json 自动生成。
% author=John Chrysostom; work_id=2201; title=格林多前书讲道集

\chapter{格林多前书讲道集}

% structure: p0001 source=h1 target=omitted reason=page-title-already-rendered-by-parent

\Emph{绪论}\newline{}\Emph{讲道一} \BibleRef{格前 1:1-3}\newline{}\Emph{讲道二} \BibleRef{格前 1:4-9}\newline{}\Emph{讲道三} \BibleRef{格前 1:10-17}\newline{}\Emph{讲道四} \BibleRef{格前 1:18-25}\newline{}\Emph{讲道五} \BibleRef{格前 1:26-31}\newline{}\Emph{讲道六} \BibleRef{格前 2:1-5}\newline{}\Emph{讲道七} \BibleRef{格前 2:6-16}\newline{}\Emph{讲道八} \BibleRef{格前 3:1-11}\newline{}\Emph{讲道九} \BibleRef{格前 3:12-18}\newline{}\Emph{讲道十} \BibleRef{格前 3:18-4:2}\newline{}\Emph{讲道十一} \BibleRef{格前 4:3-5}\newline{}\Emph{讲道十二} \BibleRef{格前 4:6-10}\newline{}\Emph{讲道十三} \BibleRef{格前 4:10-16}\newline{}\Emph{讲道十四} \BibleRef{格前 4:17-21}\newline{}\Emph{讲道十五} \BibleRef{格前 5:1-8}\newline{}\Emph{讲道十六} \BibleRef{格前 5:9-6:11}\newline{}\Emph{讲道十七} \BibleRef{格前 6:12-14}\newline{}\Emph{讲道十八} \BibleRef{格前 6:15-20}\newline{}\Emph{讲道十九} \BibleRef{格前 7:1-40}\newline{}\Emph{讲道二十} \BibleRef{格前 8:1-13}\newline{}\Emph{讲道二十一} \BibleRef{格前 9:1-12}\newline{}\Emph{讲道二十二} \BibleRef{格前 9:13-23}\newline{}\Emph{讲道二十三} \BibleRef{格前 9:24-10:12}\newline{}\Emph{讲道二十四} \BibleRef{格前 10:13-24}\newline{}\Emph{讲道二十五} \BibleRef{格前 10:25-11:1}\newline{}\Emph{讲道二十六} \BibleRef{格前 11:2-16}\newline{}\Emph{讲道二十七} \BibleRef{格前 11:17-27}\newline{}\Emph{讲道二十八} \BibleRef{格前 11:28-34}\newline{}\Emph{讲道二十九} \BibleRef{格前 12:1-11}\newline{}\Emph{讲道三十} \BibleRef{格前 12:12-20}\newline{}\Emph{讲道三十一} \BibleRef{格前 12:21-26}\newline{}\Emph{讲道三十二} \BibleRef{格前 12:27-13:3}\newline{}\Emph{讲道三十三} \BibleRef{格前 13:4-8}\newline{}\Emph{讲道三十四} \BibleRef{格前 13:8-13}\newline{}\Emph{讲道三十五} \BibleRef{格前 14:1-19}\newline{}\Emph{讲道三十六} \BibleRef{格前 14:20-33}\newline{}\Emph{讲道三十七} \BibleRef{格前 14:34-40}\newline{}\Emph{讲道三十八} \BibleRef{格前 15:1-11}\newline{}\Emph{讲道三十九} \BibleRef{格前 15:11-28}\newline{}\Emph{讲道四十} \BibleRef{格前 15:29-34}\newline{}\Emph{讲道四十一} \BibleRef{格前 15:35-46}\newline{}\Emph{讲道四十二} \BibleRef{格前 15:47-58}\newline{}\Emph{讲道四十三} \BibleRef{格前 16:1-9}\newline{}\Emph{讲道四十四} \BibleRef{格前 16:10-24}\par

\SourceRecord{Translated by Talbot W. Chambers.；Nicene and Post-Nicene Fathers, First Series；Vol. 12.；Edited by Philip Schaff.；Buffalo, NY: Christian Literature Publishing Co.,；1889.；<http://www.newadvent.org/fathers/2201.htm>.}

% structure: p0001 source=h1 target=section reason=page-title

\section{格林多前书讲道集 （序言）}

\SourceRecord{Translated by Talbot W. Chambers.；Nicene and Post-Nicene Fathers, First Series；Vol. 12.；Edited by Philip Schaff.；Buffalo, NY: Christian Literature Publishing Co.,；1889.；<http://www.newadvent.org/fathers/220100.htm>.}

% structure: p0001 source=h1 target=section reason=page-title

\section{格林多前书第一讲道词}

% structure: p0002 source=h2 target=subsection reason=relative-heading-rank; base=h2

\subsection{格林多前书 1:1-3}

\BibleQuote{保禄，蒙召为耶稣基督的宗徒，因天主的旨意，与索斯特乃弟兄，致书给在格林多的天主的教会，就是在基督耶稣内受圣化、蒙召为圣徒的人，以及所有在各处呼求我们主耶稣基督之名的人——基督是他们的主，也是我们的主：愿恩宠与平安，由我们的父天主和主耶稣基督赐与你们。}\par

1. 请看他如何从一开始就立刻摧毁他们的骄傲，打破他们一切虚妄的想像：他称自己是\Quote{蒙召的}。他说：我所学到的，并非我自行发现，也非凭自己的智慧获得，而是当我迫害并摧毁教会时，我被召叫了。如今，召叫者的一切都是召叫者的，被召叫者（可以说）什么也没有，除了服从。\par

\Quote{耶稣基督。} 你们的老师是基督，而你们却将人名，作为你们教义的守护者？\par

\Quote{因天主的旨意。} 因为天主愿意你们这样得救。我们自己没有行任何善事，而是因天主的旨意，我们获得了这救恩；因为他认为好，我们便被召叫，并非因为我们配得上。\par

\Quote{及我们的弟兄索斯特乃。} 这是他谦逊的另一例证：他将一个次于阿颇罗的人置于与自己同等的地位；因为保禄与索斯特乃之间有很大的差距。既然在差距如此悬殊之处，他让远不如他的人与自己并列，那么那些鄙视同等的人又能说什么呢？\par

\Quote{致天主的教会。} 不是\Quote{这个或那个人的}，而是\Quote{天主的}。\par

\Quote{在格林多的。} 你看见他如何在每一句话中压下他们膨胀的骄傲，从各方面训练他们的思想趋向天上吗？他也称这教会为\Quote{天主的教会}，表明它应当是合一的。因为如果它是\Quote{天主的}，它就是合一的，不仅是在格林多，而且在整个世界；因为教会的名称（\OriginalTerm{教会}{\GreekText{ἐκκλησία}}：本意为\Quote{集会}）不是分离的名称，而是合一与和谐的名称。\par

\Quote{在基督耶稣内受圣化的。} 再次出现耶稣的名字；他没有给人名留下位置。但圣化是什么？是洗礼，是净化。因为他提醒他们自己从前的不洁，他解放了他们；从而劝他们谦卑自下；因为不是由于他们自己的善行，而是因天主的仁慈，他们才得以圣化。\par

\Quote{蒙召为圣徒。} 因为即使是因信德得救，他说，也不是出于你们自己；因为你们没有首先亲近，而是被召叫的；因此，连这件小事也不完全是你们的。然而，即使你们亲近了，你们犯有无数罪恶，即便如此，这恩宠也不是你们的，而是天主的。因此，他在给厄弗所人的信中写道：\BibleRef{弗 2:8} \BibleQuote{你们得救是由于恩宠，藉着信德，这也不是出于自己。} 连信德也不完全是你们的；因为你们并非先相信，而是听从了召唤。\par

\Quote{与所有呼求我们主耶稣基督之名的人。} 不是\Quote{这个或那个人的名字}，而是\Quote{主的名}。\par

2. \Quote{在各处，是他们也是我们的。} 虽然这封信只写给格林多人，他却提到了遍及全地的所有信徒；表明普世的教会必须是合一的，尽管在各地分离；而在格林多的教会更应如此。虽然地方分离，但主将他们结合在一起，因为主是众人共有的。因此，为了将他们联合起来，他加上：\Quote{是他们也是我们的。} 这比分离更有力量［使合一］。因为正如在一处的人，有许多彼此对立的主人，就会分心，他们所在的一处不能帮助他们同心一志，因为他们的主人发出彼此矛盾的命令，各人拉向自己的方向，正如基督所说：\BibleRef{玛 6:24} \BibleQuote{你不能服侍天主而又服侍钱财。} 同样，在不同地方的人，如果他们没有不同的主人，而只有一位主人，那么地方并不能损害他们的同心，因为唯一的主将他们联合在一起。\Quote{我并不是说，（他这样说，）你们既然是格林多人，就只应与格林多人同心，而是应与普世所有的人同心，因为你们有共同的主人。} 这也是他第二次加上\Quote{我们的}的原因；因为他已经说过\Quote{耶稣基督我们的主的名}，以免粗心的人以为他作了区分，他再次加上：\Quote{我们的主和他们的主。}\par

3. 为使我的意思更清楚，我按其意义这样读：\Quote{保禄与索斯特乃，致书给在格林多的天主的教会，以及所有在各处呼求那既是我们的主也是他们的主的名的人，无论在罗马或在任何其他地方：愿恩宠与平安，由我们的父天主和主耶稣基督赐与你们。}\par

或者又这样读；我相信这更正确一些：\Quote{保禄与索斯特乃，致书给在格林多的、已受圣化、蒙召为圣徒的人，与所有在各处呼求我们的主耶稣基督之名的人，是他们也是我们的；} 也就是说，\InnerQuote{愿恩宠与平安赐与你们，你们在格林多的，已受圣化并蒙召的人；} 不是给你们单独，而是\InnerQuote{与所有在各处呼求耶稣基督之名的人，他是我们的主也是他们的。}\par

如今，若我们的平安出自恩宠，你为何自高自大？你既因恩宠得救，为何如此骄傲？你若与天主和好，为何又想归附于他人？因为这就是分离的后果。假如你与这个人\Quote{平安}相待，与另一个人甚至得到了\Quote{恩宠？}我的祈祷是，愿这两样都来自天主，我再说，来自祂并归于祂。因为若不享有从上而来的影响，它们便不能稳固地持久（\OriginalTerm{持久}{\GreekText{μένει}}，萨维勒边注）；除非以天主为对象，否则它们对你毫无益处：因为我们若与天主为敌，即使与众人和睦，也毫无益处；同样，若我们与天主和好，即使被所有人视为仇敌，也毫无损害。再者，若得众人称许而得罪了主，也毫无益处；若与天主蒙受接纳与爱，即使众人回避憎恨我们，也毫无危险。因为真正的恩宠、真正的平安，来自天主。凡在天主眼中蒙恩的人，即使遭受千万恐怖，也无所畏惧；我不仅说，不怕任何人，甚至不怕魔鬼本身；但那得罪天主的人，虽然看似安稳，却怀疑所有人。因为人性无常，不仅朋友和兄弟，甚至父亲，也曾在过去完全改变；常常为了一点小事，他们所生的、亲手栽培的枝条，比一切仇敌更成为他们迫害的对象。同样，子女也曾背弃父亲。因此，你若注意，达味在天主前蒙恩，阿贝沙隆在众人前蒙恩。他们各自的结局如何，谁获得更大的光荣，你是知道的。亚巴郎在天主前蒙恩，法郎在众人前蒙恩；因为为取悦法郎，他们交出了义人的妻子。\EditorialNote{圣金口若望\BibleRef{创 12:17}注释}那么，这两人中谁更显赫、更有福？人人都知道。我何必再提义人呢？以色列子民在天主前蒙恩，却被人（埃及人）憎恨；然而他们胜过并征服了恨他们的人，胜利何等辉煌，你们全知道。\par

因此，我们所有人都当为此殷勤努力；若有人是奴隶，让他为此祈祷，求在天主前蒙恩，而非在主子前；若是妻子，让她从天主、她的救主那里寻求恩宠，而非从丈夫；若是士兵，让他追求从上而来的眷顾，胜过追求他的君王和统帅。因为这样，你在人中间也将成为被爱之人。\par

4. 然而，人如何能在天主前蒙恩？此外，除了谦逊之心，还能怎样？因为经上说\BibleRef{雅 4:6}\Quote{天主拒绝骄傲人，却赐恩宠给谦逊人；以及\BibleRef{咏 50:17}天主的祭献是痛悔的精神，天主决不轻看谦逊痛悔的心。}因为若在人间谦逊如此可爱，何况在天主面前呢？同样，外邦人因谦逊蒙恩，犹太人却无端失去恩宠；因为\BibleRef{罗 10:13}\Quote{他们不服从天主的正义。}我所说的谦逊之人，令众人喜悦而愉快，常享平安，心中没有争执的缘由。因为即使你侮辱他、辱骂他，无论你说什么，他都会沉默，温柔地忍受，并对众人怀有极大的平安，无人能描述。是的，与天主也是如此。因为天主的诫命是要与人和睦相处：这样，我们的一生因彼此和平而兴盛。因为无人能伤害天主：祂的本性不朽，超乎一切苦难。没有什么比谦逊之心更使基督徒显为可敬。请听，例如亚巴郎说\BibleRef{创 18:27}\Quote{我不过是尘土和炉灰}；又，论及梅瑟，天主说\BibleRef{户 12:3}\Quote{他是世上最谦和的人。}因为从未有人比他更谦逊；他身为如此伟大民族的领袖，曾在海中淹没埃及王和全军，如同淹死苍蝇；在埃及、红海和旷野行了如此多的奇迹，获得了如此崇高的见证，却觉得自己如同一个普通人；作为女婿，他比岳父更谦逊\BibleRef{出 18:24}，听取他的劝告，并不愤怒，也没有说\Quote{这是什么？成就了如此宏业之后，你竟来给我们进言？}这正是大多数人的感受；即使某人提出最好的建议，他们也因提出者的卑微而藐视。但他并非如此：相反，他因谦逊之心，妥善处理了一切。因此，他也藐视君王的宫廷\BibleRef{希 11:24}，因为他确实谦逊：因为健全的心思和高尚的精神是谦逊的果实。你看，藐视君王的宫殿与餐桌，是何等高贵与宽宏的标记？因为在埃及，君王被视为神，享有无尽的财富与宝藏。但是，他放弃这一切，抛弃埃及的权杖，急忙归附于被囚的、劳苦憔悴的人，他们的力量耗尽在泥土和制砖中，连他自己的奴隶都憎恨他们（原文：\OriginalTerm{憎恨}{\GreekText{ἐβδελύσσοντο}}，七十贤士译本，\BibleRef{出 1:2}\Quote{埃及人憎恨他们}），他奔向这些人，将他们置于其主人之上。由此显然，谦逊之人正是灵魂高尚伟大的人。因为骄傲出自平庸的心思和卑下的精神，而谦逊出自伟大的胸怀和高尚的灵魂。\par

5. 若你们愿意，让我们以实例来逐一考察。请告诉我，有什么人曾比\Person{亚巴郎}更为崇高？然而正是他说了：\Quote{我只是灰土和尘埃；}也正是他，在\BibleRef{创 13:8}中说：\Quote{你我之间不要有纷争。}但这位如此谦卑的人，在\BibleRef{创 14:21}中却鄙视（\Quote{波斯，}即可能\Quote{以拦。}）波斯的战利品，不看重野蛮人的战利品；他这样做是出于极大的高尚精神和高贵培养的性情。因为真正谦卑的人才是真正崇高的（不是谄媚者，也不是伪善者）；因为真正的伟大是一回事，而傲慢是另一回事。这从以下可以清楚看出：如果一个人视泥土为泥土而鄙视它，另一个人则视泥土如黄金而珍视它；请问，哪一个是心灵崇高的人？岂不是那个拒绝赞赏泥土的人吗？哪一个是卑贱渺小的人？岂不是那个赞赏它并非常看重它的人吗？同样，你也应当这样判断：那个称自己只是灰土和尘埃的人是崇高的，虽然他出于谦卑而这样说；但那个不认为自己只是灰土和尘埃，反而疼爱自己、自高自大的人，反而应当被视为渺小，因为他把小事当作大事。由此可知，这位圣祖是出于崇高的思想说出了那句\Quote{我只是灰土和尘埃；}是出于崇高的思想，而非傲慢。\par

因为在身体上，健康丰满是一回事（\OriginalTerm{强健而有弹性的}{\GreekText{σφριγῶντα}}），而肿胀是另一回事，尽管两者都表示肉体的充实（但后者是不健康的，前者是健康的）；同样在此：傲慢是一回事，可以说如同肿胀；而心灵高尚是另一回事，是健康的状态。再者，有一个人因自身身高而高大；另一个人矮小，穿上厚底靴后变高；现在告诉我，这两人中我们应称哪一位为高大？岂不是那高度来自自身的人吗？因为另一个人是借用别人的高度；踏在低矮的东西上，结果成了高个子。许多人就是如此，将自己架在财富和荣耀之上；这不是崇高，因为崇高的人不需要这些东西，反而鄙视它们，他的伟大来自自身。因此，让我们谦卑，好使我们被高举；如\BibleRef{路 14:11}所说：\Quote{因为自谦的必被高举。}任性的人并非如此；相反，他是所有性格中最平凡的。因为泡沫也是膨胀的，但膨胀并不健康；因此我们称这些人为\Quote{骄傲自大。}而清醒的人没有高傲的思想，即使在好运中也不高傲，因为他知道自己的卑微；但庸俗的人即使在小事上也纵容骄傲的幻想。\par

6. 那么，让我们获得那由谦卑而来的崇高。让我们审视人性之事，以点燃我们对未来之事的渴望；因为没有其他方法能使人谦卑，除非热爱神圣之事而鄙视现世之事。因为正如一个即将获得王国的人，如果有人不献上紫袍，却献上一些微不足道的恭维，他会视之为无物；同样，如果我们渴望另一种荣誉，我们也会嘲笑现世的一切。你们没有看见孩子们在玩耍时，组成一队士兵，有先驱和扈从，一个男孩在中间扮演将军，这一切是多么幼稚吗？所有人类之事都是如此；是的，甚至比这些更无价值：今天存在，明天就不复存在。因此，让我们超越这些事物；不但不要渴望它们，而且如果有人向我们提供它们，我们甚至要感到羞愧。因为这样，通过驱除对这些事物的爱，我们将拥有那神圣的另一种爱，并享受不朽的荣耀。愿天主恩赐我们众人获得这一切，因我主耶稣基督的恩宠和慈爱；愿光荣与权能归于圣父，偕同圣善的圣神，至于永世。阿们。\par

\SourceRecord{Translated by Talbot W. Chambers.；Nicene and Post-Nicene Fathers, First Series；Vol. 12.；Edited by Philip Schaff.；Buffalo, NY: Christian Literature Publishing Co.,；1889.；<http://www.newadvent.org/fathers/220101.htm>.}

% structure: p0001 source=h1 target=section reason=page-title

\section{格林多前书讲道第二篇}

% structure: p0002 source=h2 target=subsection reason=relative-heading-rank; base=h2

\subsection{格前 1:4-5}

我时常为你们感谢我的天主，因天主在基督耶稣内赐给你们的恩宠，使你们在一切事上因他富足。\par

他所劝勉别人去做的事，自己也照样实行，说：\Quote{\BibleRef{斐 4:6} 应带着感恩将你们的请求告知天主}，又以此教导我们，务要常从这些话开始，并在一切事上首先感谢天主。因为再没有什么比人存感恩之心，为自己也为他人感谢天主，更蒙天主悦纳的了。故此，他几乎每封书信开头都以此开篇。但在此处，他这样做的缘由比在其他书信中更为迫切。因为感恩的人，既是因为自己处境因泰，也是承认所受的恩惠；而恩惠并非债款，不是偿还，也不是报酬。这一点在任何场合都十分重要，但在对那些渴望教会分裂的格林多人讲话时尤其如此。\par

他带着深厚的情感，抓住那个共有的名号，使之成为自己的，就如先知们也时常所说\BibleRef{咏 43:4}：\Quote{天主，我的天主}；并以此鼓励他们也使用同样的言辞。因为这样的说法属于那些远离一切世俗之事、专心向所热切呼求的那一位迈进的人：唯独那常从今世之事向上主攀登、总是以祂为优先、不断感恩的人——不仅为已经赐下的恩宠，也为后来任何时候所施予的祝福——因而献上同样的赞美。故他不仅仅说：\Quote{我感谢}，而是说：\Quote{时常，为你们}；以此教导他们应当常怀感恩，且只向唯一的天主感恩。\par

\Quote{因为天主的恩宠。}你看出他是从各方面引出话题来纠正他们？因为哪里有\Quote{恩宠}，就没有\Quote{行为}；哪里有\Quote{行为}，就不再是\Quote{恩宠}。如果因此是\Quote{恩宠}，为何你们心高气傲？又因何自高自大？\par

\Quote{是赐给你们的。}是谁赐下的？是我，还是另一位宗徒？绝对不是，而是\Quote{藉着耶稣基督}。因为\Quote{在基督耶稣内}这句话正表明此意。请注意他在多处使用\OriginalTerm{在}{\GreekText{ἐν}}一词，即\Quote{在内}，代替\OriginalTerm{藉着}{\GreekText{δἰ} \GreekText{οὗ}}，即\Quote{藉着祂}，故其含义并无不同。\par

\Quote{使你们在一切事上富足。}又是谁使你们如此？答复是：由祂。这不只是\Quote{你们富足}，而是\Quote{在一切事上}。既然首先是\Quote{财富}，然后是\Quote{天主的财富}，接着是\Quote{在一切事上}，最后是\Quote{藉着独生子}，请思想这不可言喻的宝藏！\par

% structure: p0009 source=h3 target=subsubsection reason=relative-heading-rank; base=h2

\subsubsection{格前 1:5}

\Quote{在一切言论和一切知识上。}\Quote{言论}（或作\Quote{言语}）不是外邦人的那种，而是天主的。因为有没有\Quote{言论}的知识，也有有\Quote{言论}的知识。因为有许多人拥有知识，却没有演讲的能力，如同那些未受教育、不能清楚表达心中所想的人。他说，你们却不是这样，而是既能理解又能言谈。\par

% structure: p0011 source=h2 target=subsection reason=relative-heading-rank; base=h2

\subsection{格前 1:6}

\Quote{正如基督的见证在你们中间得以证实。}在赞扬和感谢的掩饰下，他严厉地触痛他们。他说：\Quote{并非藉着外邦人的哲学，也不是藉着外邦人的学问，而是藉着天主的恩宠，}以及藉着由祂所赐的\Quote{财富}、\Quote{知识}和\Quote{言论}，你们才得以学习真理的教义，并得以坚定于主的见证，即福音。\par

% structure: p0013 source=h2 target=subsection reason=relative-heading-rank; base=h2

\subsection{格前 1:7}

4. \Quote{使你们在恩赐上不落后于人。}这里出现一个大问题。那些曾\Quote{在一切言论上富足}，以致在任何恩赐上都\Quote{不落后于人}的人，难道还是属血肉的吗？因为如果他们起初便是如此，现在更是如此。那么，他怎称他们为\Quote{属血肉的}？因为他说：\BibleRef{格前 3:1} \Quote{我不能把你们当作属神的，只能当作属血肉的。}那么，我们该说什么？或许他们在起初相信时，获得了所有恩赐（因为他们确实热切寻求），但后来懈怠了。或者并非如此，而是这些恩赐和那些批评并非针对所有人；而是前者针对那些接受他赞许的人，后者针对那些受他责备的人。因为事实上他们仍然有恩赐：\BibleRef{格前 14:26}他说：\Quote{每人有圣咏、有启示、有语言、有解释；一切都当为建树而行。}以及\Quote{先知可以两三个人说话。}此外，我想他暗示了自己的作为，因为他也曾显过征兆，正如他在第二封书信中对他们说：\BibleRef{格后 12:12} \Quote{宗徒的征兆在你们中间以各种坚忍得以实现}；又说：\Quote{你们在哪方面不如别的教会？}\par

或者，如我所说，他既提醒他们自己的奇迹，也说话针对那些仍然被认可的人。因为那里有许多圣人，他们\Quote{献身服务圣徒}，并成了\Quote{阿哈雅初果}；正如他在书信末所宣布的\BibleRef{格前 16:15}。\par

5. 无论如何，虽然这些赞美并不十分接近真理，但仍是作为一种预防措施而插入的（\OriginalTerm{以安排的方式}{\GreekText{οἰκονομικῶς}}），为他的言论预先铺路。因为谁若在开初就讲令人不悦的事，就会使自己的话在软弱者中失去听众：因为如果听者与他地位相等，他们会感到恼怒；若远低于他，他们则会烦恼。为避免此，他以看似赞美的话开始。我说\Quote{看似}，因为连这赞美也不属于他们，而是属于天主的恩宠。因为他们得蒙罪赦并获得成义，这来自上天的恩赐。因此他也着重于这些显示天主慈爱的点，为能更彻底地清除他们的病症。\par

6. \BibleQuote{等待我们主耶稣基督的\OriginalTerm{启示}{\GreekText{ἀποκάλυψιν}}。} \Quote{你们为何大惊小怪，}他说，\Quote{为何因基督没有来临而烦恼？不，祂已来临；那日子如今已在门前。}请看他的智慧：他如何将人从世俗的思虑中拉开，以可畏的审判宝座来威慑他们，从而暗示不仅开端必须良好，结局也应如此。因为有了这一切恩赐以及一切美善，我们还必须顾及那日子：并且需要许多劳苦才能到达终点。他用了\Quote{启示}一词，暗示虽然祂未被看见，但祂存在，且现在就在，那时将显现。因此需要忍耐：因为你们领受奇事正是为此，使你们能坚定不移。\par

% structure: p0018 source=h2 target=subsection reason=relative-heading-rank; base=h2

\subsection{格前 1:8}

7. \BibleQuote{祂也必坚固你们到底，使你们在我主耶稣基督的日子无可指责。}此处他似乎讨好他们，但这话毫无谄媚；因为他也知道如何严厉责备他们；如他所说：\BibleRef{格前 4:18}\BibleQuote{有些人自高自大，以为我不会到你们那里去；}又说：\BibleQuote{你们愿意怎样呢？是愿意我带着刑杖到你们那里去，还是以爱德和温柔的精神？}以及\BibleRef{格后 13:3}\BibleQuote{既然你们寻求基督在我内说话的证据。}但他也暗中指责他们：因为说\Quote{祂必坚固}以及\Quote{无可指责}一词标明了他们仍在摇摆，并当受责备。\par

但请思考他如何总是将人如钉子般牢牢拴在基督的名字上。他记起的不是任何个人或导师，而不断是那所渴望的本身。他仿佛要唤醒那些醉酒后昏昏沉沉的人。因为在其他任何书信中，基督的名字出现都不如此频繁。但在这里，短短几节中多次出现；并且可以说，他藉此编织了整个序言。请看从头开始：\BibleQuote{保罗蒙召为耶稣基督的宗徒，致书给在耶稣基督内成圣、蒙召为圣徒、以及所有各处呼求我们主耶稣基督之名的人。愿恩宠与平安从天主我们的父和主耶稣基督赐与你们。我时时为你们感谢我的天主，为了天主在基督耶稣内所赐与你们的恩宠，因为你们藉着祂在一切事上，在一切言论和知识上，都成了富有的，正如基督的见证在你们中得以坚定，使你们在我主耶稣基督的日子上无可指责。天主是忠信的，因为你们是由祂所召，与祂的圣子我们的主耶稣基督合而为一。弟兄们，我因我们主耶稣基督之名，劝告你们……}你们看见基督之名反复出现吗？由此即使最不经心的人也明白，他这样做并非偶然或无意，而是为了藉不断应用这荣耀之名，平息他们的炎症，清除疾病的腐败。\par

% structure: p0021 source=h2 target=subsection reason=relative-heading-rank; base=h2

\subsection{格前 1:9}

8. \BibleQuote{天主是忠信的，你们原是由祂所召，与祂的圣子合而为一。}何等奇妙！他此处所说何其伟大！他所宣告的恩赐何其宏大！你们蒙召与独生子共融，却竟投靠世人？还有什么比这更可怜呢？你们是如何蒙召的？由圣父。因为他常提及圣子时用\Quote{藉着祂}、\Quote{在祂内}等说法，免得人以为他这样说是将圣子视为较低者，便将同样的说法也归于圣父。他说，不是藉着这个或那个，而是\Quote{由圣父}你们蒙召；也是由祂\Quote{得以富足}。再说，\Quote{你们蒙召}，并非你们自己前来。但\Quote{与祂的圣子合而为一}是什么意思？请听他别处更清楚地说明此事：\BibleRef{弟后 2:12}\BibleQuote{如果我们受苦，也必与祂一同为王；如果我们与祂同死，也必与祂同生。}然后，因他所说的是大事，便加了一个无可辩驳的论证；因为，他说，\Quote{天主是忠信的}，即\Quote{真实的}。如果\Quote{真实}，祂所应许的也必实现。祂应许使我们成为祂独生子的分享者；为此祂也召叫了我们。因为\BibleRef{罗 11:29}\BibleQuote{天主的恩赐和召叫是决不会撤回的。}\par

他以一种神圣的技巧如此早地插入这些话，免得在严厉责备之后他们陷入绝望。因为若我们不完全不受祂的驾驭（\OriginalTerm{不受驾驭}{\GreekText{ἀφηνιάσωμεν}}），天主的部分必然会实现。如犹太人，虽蒙召却不愿接受祝福；但这不再是召唤者的问题，而是他们缺乏理智。因为祂确愿赐予，但他们因拒绝接受而自弃。因为，若祂召唤他们从事痛苦艰辛的事业，即便那样他们也不能原谅自己的推辞；不过，他们或许可以说事出有因：但若召唤是为了洁净\BibleRef{格前 1:4}、成义、圣化、救赎、恩宠、白白的恩赐以及那眼所未见、耳所未闻的将来美善；并且是天主召唤，亲自召唤；那么，那些不跑向祂的人还能得到什么宽恕呢？因此，没有人可责怪天主；因为不信并非来自召唤者，而是来自那些\OriginalTerm{背离}{\GreekText{ἀποπηδῶντας}}祂的人。\par

9. 但有人会说：\Quote{祂应当把人带进来，即使违背他们的意愿。}切莫如此。祂不用暴力，也不强迫；因为谁会邀请人去享荣誉、冠冕、盛宴和庆典，却强行拖拽不愿、被捆绑的人呢？没有人。因为这是侮辱人的行为。祂违背人的意愿送他们下地狱，却召唤甘心情愿的人进入天国。祂捆绑着、哀哭着的人带到火里；对永福的境界却不是这样。否则，对福乐本身也是一种羞辱，如果它们的本性不是让人自愿地、满怀感激地奔向它们的话。\par

\Quote{那么，为什么所有人都不选择它们呢？}你说，\Quote{出于他们自己的软弱。}\Quote{为什么祂不除去他们的软弱呢？}请告诉我，祂应当以何种方式除去它？难道祂没有创造一个世界，教导祂的慈爱和权能吗？因为\BibleRef{咏 18:1}\BibleQuote{诸天，}有人说，\BibleQuote{宣扬天主的光荣。}难道祂没有派遣先知吗？难道祂没有召叫并尊重我们吗？难道祂没有行奇事吗？难道祂没有赐下成文的和自然的法律吗？难道祂没有派遣祂的圣子吗？难道祂没有委任宗徒吗？难道祂没有完成罪过？难道祂没有以地狱相威胁？难道祂没有应许天国？难道祂不是每日使太阳升起吗？祂所命令的事不是如此简单容易，以致许多人在克己的伟大中超越了祂的诫命吗？\BibleQuote{我向我的葡萄园还有什么可做的，而我没有做呢？}\BibleRef{依 5:4}\par

10. \Quote{那么，}你说，\Quote{为什么祂不让知识和德行成为我们自然的本性呢？}谁这样说？是希腊人还是基督徒？两者都有，但并非针对同一件事：因为前者以知识为出发点提出异议，后者则以行为为出发点。首先，我们要答复我们这边的人；因为我不太在意那些外人，而是更关心我们自己的成员。\par

那么基督徒说什么呢？\Quote{应当把德行本身的知识植入我们内。}祂已经植入了；因为如果祂没有这样做，我们怎会知道哪些事该做，哪些不该做？法律和法庭从何而来？但\Quote{天主不仅应当赐予知识，还应当赐予行动本身（德行）。}那么，如果一切都出于天主，你还有什么可被奖赏的呢？请告诉我，天主惩罚你和惩罚犯罪的希腊人，方式是一样的吗？当然不是。因为直到某一点，你拥有自信，即来自真知识的自信。那么，如果有人现在说，在知识方面，你和希腊人将被视为同等的功德，这难道不会令你反感吗？我想确实如此。因为你可能会说，希腊人本来有能力获得知识，却不愿意。如果后者也声称天主应当把知识自然植入我们内，难道你不会嘲笑他，并对他说：\Quote{但你为什么不去寻求它？你为什么不像我一样认真呢？}你会满怀自信地站稳脚跟，说责备天主没有自然植入知识是极其愚蠢的。你会这样说，因为你已经获得了知识方面的东西。同样地，如果你也完成了实践方面的事，你就不会提出这些问题了；但你厌倦了德行的实践，因此你用这些轻率的话来庇护自己。然而，以必然的方式使人成为善，这怎么可能合理呢？接下来，我们岂非要和野兽争论德行，因为其中一些比我们更加节制？\par

但你说：\Quote{我宁愿被必然所迫成为善，从而丧失一切赏报，也不愿因自由选择而作恶，以致受罚遭受报应。}但一个人永远不可能因必然成为善。因此，如果你不知道应当做什么，就请指出来，然后我们会告诉你应当说什么。但如果你知道不洁是邪恶的，为什么不逃离这邪恶的事呢？\par

\Quote{我不能，}你说。但那些做过比这更大事的人会反驳你，并且足以堵住你的口。因为你也许和妻子同住，却并不贞洁；而另一个人甚至没有妻子，却保持贞洁不受侵犯。既然别人甚至跃过了界限所划定的标记，你又有什么借口不遵守规则呢？\par

但你却说：\Quote{我的体质或心态并非如此。}这不是缺乏能力，而是缺乏意愿。因为我如此证明：所有人都对德行有某种倾向：一个人做不到的事，即使迫于必要，他也做不到；但如果迫于必要，他就能做到，那么他未做便是出于选择。我所说的意思是：带着沉重的身体飞翔并升向天空，这简直是不可能的事。那么，如果一位君王命令某人做这事，并以死亡相威胁，说：\InnerQuote{不飞的人，我下令斩首、焚烧或施以其他惩罚：}有人会服从他吗？当然不会。因为本性无法做到。但如果对贞洁也做同样的事，他制定法律，污秽者应受惩罚、被焚烧、被鞭打、受尽酷刑，难道许多人不会服从这法律吗？\Quote{不会，}你会说：\Quote{因为现在就有禁止通奸的法律，但并非所有人都遵守。}不是因为恐惧失去效力，而是因为大部分人期望不被发现。所以如果当他们即将行污秽之事时，立法者和审判者出现在他们面前，恐惧就足以驱散情欲。不，即使我施加另一种比这更轻的强制力；如果我把这个人带走，使他远离所爱的人，并把他用锁链严密关起来，他也能忍受，而不受太大伤害。那么，我们不要说某人生来邪恶：因为如果一个人生来善良，他绝不会在任何时候变得邪恶；如果他生来邪恶，他也绝不会变得善良。但我们看到变化迅速发生，人们很快从这边转向那边，又从那边跌回这边。这些事我们不仅在圣经中看到，例如，税吏成为宗徒；门徒成为叛徒；娼妓变得贞洁；强盗成为有声誉的人；术士敬拜；不敬神的人转向虔诚，在新约和旧约中都有；而且每天都能看到许多这样的事。如果事物是本性的，它们就不会改变。因为我们这些有感觉的本性，绝不能通过任何努力变得无感觉。因为本性如何，就永远不会脱离其本然状态。例如，没有人从睡眠变为不睡眠；没有从腐朽变为不朽；没有从饥饿变为永无饥饿感。因此，这些事也不是可责备的，我们也不为此自责；也从没有人在责备别人时说：\InnerQuote{你啊，可朽坏的、易受情欲的；}但我们总是把通奸、淫乱或诸如此类的事归咎于那些负责之人；我们把他们带到法官面前，法官责备并惩罚他们，而在相反的情况下则给予荣誉。\par

11. 既然从我们彼此的行为，从别人审判我们时的行为，从我们制定法律的事，从我们虽无人控告却自己定罪的事，从我因懒惰而变坏、因恐惧而变好的例子，从我们看到别人行善并达到自制\OriginalTerm{哲学}{\GreekText{φιλοσοφίας}}顶峰的事例，都清楚地表明我们也有行善的能力：那么，为什么我们大多数人用无心的借口和托词徒然欺骗自己，不仅得不到宽恕，反而招致无法忍受的惩罚？我们本应将那可怕的日子摆在眼前，专心于德行；经过少许劳苦，获得不朽的冠冕。因为这些话不能为我们辩护；相反，我们的同仆，那些实践相反德行的人，将定所有继续犯罪者的罪：残忍的人被仁慈者定罪；邪恶的被善良者定罪；凶暴的被温柔者定罪；吝啬的被慷慨者定罪；虚荣的被克己者定罪；懒惰的被勤勉者定罪；放纵的被节制者定罪。天主将这样审判我们，将两群人各归其位：对一群施予赞扬，对另一群施予惩罚。但愿天主不让任何在场者成为受罚和受辱的人，而是成为得冠冕和获得天国的人。愿我们众人借着我们的主耶稣基督的恩宠和慈爱，都能获得这一切；愿光荣、权能、尊崇归于圣父、圣子和圣神，从现在直到永远，并至无穷之世。阿们。\par

\SourceRecord{Translated by Talbot W. Chambers.；Nicene and Post-Nicene Fathers, First Series；Vol. 12.；Edited by Philip Schaff.；Buffalo, NY: Christian Literature Publishing Co.,；1889.；<http://www.newadvent.org/fathers/220102.htm>.}

% structure: p0001 source=h1 target=section reason=page-title

\section{格前书讲道集第三篇}

% structure: p0002 source=h2 target=subsection reason=relative-heading-rank; base=h2

\subsection{格前 1:10}

\BibleQuote{弟兄们，我因我们主耶稣基督之名，求你们众人言谈一致，不要有分裂，但要同心合意，全然相合。}\par

我不断说的话，就是我们当以温柔渐进的方式提出责备，保禄在此也是这样做的；因为他正要进入一个充满许多危险、足以将教会连根拔起的主题，却使用非常温和的语言。他的话是他\Quote{恳求}他们，并且\Quote{因基督}恳求他们，仿佛他自己单独不足以作此恳求而获胜。\par

但这是什么意思？\Quote{我因基督恳求你们？}\Quote{我以基督为我作战，帮助我，即祂那受伤害和受侮辱的名号。}这真是一种可怕的说话方式！免得他们变得刚硬和不知羞耻：因为罪使人不安。因此，如果你立刻（\OriginalTerm{如果立刻责斥}{\GreekText{ἄν} \GreekText{μὲν} \GreekText{εὐθέως} \GreekText{ἐπιπλήξης}} Savil. \OriginalTerm{若不}{\GreekText{ἄν} \GreekText{μὴ}} Ben.）严厉责斥，你就会使人变得粗暴无礼；但如果你使他羞愧，就会折服他的颈项，抑制他的自信，使他低头。保禄也怀着这个目的，暂时满足于以基督的名恳求他们。他要恳求的，究竟是什么呢？\par

\BibleQuote{叫你们众人言谈一致，不要有\OriginalTerm{分裂}{schisms}。}\Quote{分裂}一词的强调力——我指的是这个词本身——已足够作为控诉。因为他们并非变成了许多各自完整的部分，而是那唯一的（\OriginalTerm{原初存在的身体}{Body which originally existed}）已经消亡。因为如果他们是完整的教会，可能有许多个；但如果是分裂，那么那最初的唯一就消失了。因为那本身完整的，不仅不会因分裂成许多部分而变成许多，甚至连原来的唯一也丢失了。分裂就是这样的性质。\par

2. 接下来，因为他用了\Quote{分裂}一词严厉对待他们，便又缓和安抚他们说：\BibleQuote{但应同心合意，全然相合。}也就是说：既然他说了\BibleQuote{叫你们众人言谈一致}，他接着说：\Quote{不要以为我说和谐只在言语上；我寻求的是心灵上的和谐。}但是，由于存在言语上的一致，并且是真诚的，然而并非在所有事上一致，因此他加上了这句：\BibleQuote{全然相合。}因为在一件事上合一，在另一件事上分歧的人，就不再是\Quote{完全}的，也没有切合到完全一致。还有一种意见的和谐，却没有情感的和谐；例如，当我们有同样的信仰，却没有在爱中结合：因为这样，我们在意见上是一（因为我们想同样的事），但在情感上并非如此。当时的情况就是这样：这个人选择一个领袖，那个人选择另一个。为此，他说必须在\Quote{心思}和\Quote{判断}上都一致。因为这些分裂并非来自信仰的差异，而是由于人的好争论导致判断分裂。\par

3. 但是，既然被责备的人若没有证人便恬不知耻，请看他是如何不允许他们否认事实，而是引出一些证人来作证。\par

% structure: p0009 source=h2 target=subsection reason=relative-heading-rank; base=h2

\subsection{格前 1:11}

\BibleQuote{因为关于你们，我的弟兄们，我由克罗厄的家人得知了。}他并没有在开头就这样说，而是先提出了指控，像一个信任告密者的人。因为若不是这样，他就不会责备了：因为保禄不是轻易相信的人。因此他既没有立刻说\Quote{我得知了}，以免显得是凭他们的权威而责备；也没有完全不提他们，以免显得是凭自己说话。而且，他称他们为\Quote{弟兄们}；因为虽然过错是明显的，但称呼人弟兄并无不妥。再看他谨慎之处：不说某个具体的人，而是整个家庭，以免使他们对告密者怀有敌意。因为这样，他既保护了告密者，又无所畏惧地提出了指控。因为他不只关心一方的益处，也关心另一方的益处。因此他没有说\Quote{有人告诉我}，而是指出了家庭，免得他们认为他在捏造。\par

4. 他得知了什么呢？\Quote{你们中间有纷争。}这样，当他责备他们时，他说：\BibleQuote{不要有分裂}；但当报告别人的话时，他更温和地说：\Quote{因为有人告诉我……你们中间有纷争；}以免给告密者带来麻烦。\par

% structure: p0012 source=h2 target=subsection reason=relative-heading-rank; base=h2

\subsection{格前 1:12}

接着他也说明了纷争的种类。\par

\BibleQuote{你们各人说：我是属保禄的，我是属阿颇罗的，我是属刻法的。} \Quote{我说：有纷争，}他说，\Quote{我是指，不是关于私事，而是属于更严重的那一类。} \BibleQuote{你们各人说；} 因为这败坏不是涉及一部分，而是涉及整个教会。然而他们并不是在谈论他自己，也不是伯多禄，也不是阿颇罗；但他的意思是，如果连这些人都不可依靠，那别人就更不用说了。因为他们并没有谈论他们，他接着说：\BibleQuote{我把我自己和阿颇罗当作一个比喻，使你们从我们身上学习，不要超过圣经所写的。} 因为假如连称呼自己为保禄、阿颇罗、刻法的名字都不应该，那么其他人的名字就更不应该了。如果在这位教师、宗徒之首、且教导过这么多人的门下，都不该登记自己的名字，那么在那虚无之人的门下就更不应该了。于是，他用夸张的方式，想使他们脱离这种病态，就列出这些名字。此外，他使他的论点不那么尖锐，不点名指出这些粗野的教会分裂者，而是用宗徒的名字像面具一样把他们掩盖起来。\par

\BibleQuote{我是属保禄的，我是属阿颇罗的，我是属刻法的。} 他并非自视高于伯多禄而将自己的名字放在最后，而是大大地偏爱伯多禄。他以递进的方式（\OriginalTerm{按递进法}{\GreekText{κατὰ} \GreekText{αῦξησιν}}）安排他的陈述，免得人们以为他这样做是出于嫉妒；或者出于忌妒而贬低他人的荣誉。因此他也把自己的名字放在最前面。因为一个把自己放在最前面的人被拒绝，不是出于爱荣誉，而是出于对这种名声的极端蔑视。你看，他把自己置于整个攻击的路上，然后提到阿颇罗，然后提到刻法。他这样做不是为了抬高自己，而是在谈论错误的事情时，首先在自己身上施行必要的纠正。\par

5. 但是，那些沉迷于这个或那个人的人确实错了。他正确地责备他们说：\BibleQuote{你们做得不对，因为你们说：\InnerQuote{我是属保禄的，我是属阿颇罗的，我是属刻法的。}} 但他为什么加上 \BibleQuote{我是属基督的？} 因为虽然这些沉迷于人的是错了，但那些献身于基督的人肯定不是确实 \OriginalTerm{并不确实}{\GreekText{οὔδε} \GreekText{που}}（Bened.) \OriginalTerm{并不确实}{\GreekText{οὐ} \GreekText{δήπου}}（Savil.）错了。但这不是他的指控，即他们以基督的名自称，而是他们并未全都唯独以那名为名。我认为他是自己加上这句，想要使指控更加严重，并指出按照这个规则，基督被认为只属于一个党派：尽管他们自己并没有这样使用那名。因为他所暗示的，他在后续中宣告说，\par

% structure: p0017 source=h2 target=subsection reason=relative-heading-rank; base=h2

\subsection{\BibleRef{格前 1:13}}

\BibleQuote{基督被分开了吗？} 他所说的是：\Quote{你们把基督切碎了，分开了祂的身体。} 这是愤怒！这是斥责！这是充满愤慨的话语！因为每当他不争论而只是质问时，他的做法就暗示了一种公认的荒谬。\par

但有人说他说 \BibleQuote{基督被分开了} 是指别的事：好像他说 \Quote{祂把教会分配给人并分开了，自己取了一部分，把另一部分给他们。} 接着，他努力推翻这种谬论，说：\BibleQuote{保禄曾为你们被钉十字架吗？或者你们受洗是归入保禄的名下吗？} 请注意他热爱基督的心；他从此把整个事情集中到自己的名字上，表明并且超越了表明，这荣誉不属于任何人。并且为确保没有人认为他是出于嫉妒而说这些话，他不断地把自己放在前面。也请注意他体贴的方式，他没有说：\Quote{保禄创造了世界吗？保禄从无中创造了你们吗？} 而只提及那属于信友的珍宝，且被人极度关心的事物——他特别指出十字架、圣洗圣事以及随之而来的福祉。因为天主对人的慈爱也通过世界的创造显示出来，但没有什么比通过十字架的 \OriginalTerm{屈尊俯就}{\GreekText{τῆς} \GreekText{συγκαταβάσεως}} 更显示这一点。他没有说 \Quote{保禄为你们死了吗？}而是 \BibleQuote{保禄被钉十字架了吗？} 也指出了死亡的方式。\par

\BibleQuote{或者你们受洗是归入保禄的名下吗？} 他又没有说：\Quote{保禄给你们付洗了吗？} 因为他确实付洗了许多人：但这问题不在于谁付了洗，而在于他们受洗归入了谁的名下！因为这也是导致分裂的一个原因，即他们以付洗者的名字自称，他同样纠正了这个错误，说：\BibleQuote{你们受洗是归入保禄的名下吗？} \Quote{不要告诉我，}他说，\Quote{谁付了洗，而要问归入谁的名下。因为受调查的不是付洗的人，而是在圣洗圣事中被呼求的那一位。因为那一位是赦免我们罪过的。}\par

至此他停住了议论，不再继续这个话题。因为他没有说：\Quote{保禄曾向你们宣告未来的美事吗？保禄曾应许你们天国吗？} 那么，我问，他为什么不也加上这些问题呢？因为应许天国和被钉十字架并不相同。前者既没有危险也没有带来羞辱；而后者两者都有。此外，他以后者证明前者：因为说了 \BibleRef{罗 8:32} \BibleQuote{祂没有怜惜自己的独生子，} 他接着又说：\BibleQuote{怎能不也把一切同祂一起白白赐给我们呢？} 又，\BibleRef{罗 5:10} \BibleQuote{如果我们还作仇敌的时候，藉着祂圣子的死，与天主和好了，那么，和好之后，更要因祂的生命得救了。} 这是他没有加上我刚刚提到的问题的原因之一：也是因为前者他们尚未获得，而后者他们已经历过。前者尚在应许中，后者已经成就。\par

% structure: p0022 source=h2 target=subsection reason=relative-heading-rank; base=h2

\subsection{\BibleRef{格前 1:14}}

6. \Quote{我感谢天主，除了克黎斯颇和加约外，我没有给你们中任何人付洗。} \Quote{你们因付洗而自高自大，而我却因没有这样做而感谢！}这样，他藉着一种\OriginalTerm{神圣的安排}{\GreekText{οἰκονομικῶς}}消除了他们在这点上的傲慢；不是消除圣洗的效力（天主不许），而是消除那些因付洗而自夸者的愚蠢：首先，他指出这恩赐不是他们的；其次，他为此感谢天主。因为圣洗确实是一件大事：但它的伟大不在于付洗者的工作，而在于在圣洗中被呼求的那一位；因为付洗本身对人的劳苦而言算不得什么，远不如宣讲福音。是的，我再说，圣洗确实伟大，没有圣洗不可能获得天国。但一个并无卓越才能的人也能付洗，而宣讲福音则需要巨大的劳苦。\par

% structure: p0024 source=h2 target=subsection reason=relative-heading-rank; base=h2

\subsection{格前 1:15}

他也说明为什么他感谢自己没有给任何人付洗。这原因是什么？\Quote{免得有人说：你们是在我名下受洗的。}怎么，他是指在其他情况下他们这样说吗？完全不是；而是说：\Quote{我怕，}他说，\Quote{这病甚至会蔓延到那种地步。因为如果卑微无足轻重的人付洗都会引起异端，那么，我作为圣洗的首位宣讲者，若给许多人付洗，他们很可能会结党，不仅以我的名字自称，而且还会把圣洗归功于我。}因为如果从地位较低的人中产生了如此大的恶，那么从地位较高的人中，或许会发展成更严重的事。\par

% structure: p0026 source=h2 target=subsection reason=relative-heading-rank; base=h2

\subsection{格前 1:16}

接着，为使在这方面有缺陷的人感到羞愧，他又补充说：\Quote{我也给斯特法纳一家付过洗，}然后进一步压下他们的骄傲，说：\Quote{我不知道是否还付洗了别人。}由此他表明，他并不极力追求由此从众多的人获得的尊荣，也不是为了荣耀而做这事。\par

% structure: p0028 source=h2 target=subsection reason=relative-heading-rank; base=h2

\subsection{格前 1:17}

不仅通过这些，而且通过接下来的话，他极大地抑制了他们的骄傲，说：\Quote{基督派遣我，不是为付洗，而是为宣讲福音：}因为更劳苦的部分，需要大量辛劳和钢铁般灵魂的部分，以及一切所依赖的部分，正是这个。因此，保禄才被委以此任。\par

那么，他既然不是被派去付洗的，为何又付洗呢？不是与派遣他的那一位相争，而是在此超越了他的任务而劳苦。因为他说：\Quote{我不是被禁止，}而是说：\Quote{我不是为此被派遣，而是为那最必要的事。}因为宣讲福音可能是一两个人的工作；而付洗，却是每个拥有司祭职的人都能做的。因为一个人受了教导并被说服，带他来付洗，任何一个人都能做：其余的，一切都由前来者的意愿和天主的恩宠完成。但当要教导不信者时，就需要巨大的劳苦和智慧。而且在当时，还伴随有危险。在前一种情况下，整个事情已完成，即将领受入门圣事的人已被说服：当一个人被说服时，给他付洗并非大事。但在后一种情况下，劳苦是巨大的：要改变审慎的意愿，改变思想的趋向，连根拔起错误，并在其位置上栽种真理。\par

他并没有完全说出这一切，也没有明确论证付洗没有劳苦，而是宣讲有劳苦。因为他知道总是如何缓和语气，而在与世俗智慧的比较中，他却非常热切，因为主题使他能够使用更激烈的言辞。\par

因此，他付洗并非与派遣他的那一位相悖；而是如同在寡妇的事上，尽管宗徒们说过，\Quote{我们不应放弃天主的圣言而去管理饭食，}\BibleRef{宗 6:2}他却履行了执事的职务\BibleRef{宗 12:25}，不是与他们相悖，而是作为超乎他任务的事；此处也是如此。因为即使现在，我们也把这件工作托付给较单纯的司铎，而将教义的宣讲托付给较有智慧的人：因为那里才有劳苦和汗水。因此他本人说：\Quote{那些善于管理的长老，应受加倍的尊敬，尤其是那些在圣言和教导上劳苦的人。}\BibleRef{弟前 5:17}\par

7. \Quote{并且不是凭智慧的言辞，免得基督的十字架失去效力。}\par

他打下那些因付洗而骄傲之人的气焰后，转而迎击那些以世俗智慧夸口的人，并以更激烈的态度披挂上阵。对那些因付洗而自高自大的人，他说：\Quote{我感谢我没有给任何人付洗；}又说：\Quote{因为基督派遣我，不是为付洗。}他说话既不激烈也不辩论，只是淡淡地暗示自己的意思，就迅速跳过去了。但在这里，他在开头就给以重击，说：\Quote{免得基督的十字架成为空的。}那么，你为什么在这件本该使你羞愧的事上夸口呢？因为如果这智慧与十字架为敌并与福音相争，就不应该夸耀它，而应羞愧退避。因为宗徒们之所以不智慧，并非由于恩赐的软弱，而是怕所宣讲的福音受到损害。因此，上述那类人并非被选用来辩护圣言的，反而是在诽谤它。无学问的人才是建立圣言的人。这能抑制虚荣，这能压制骄傲，这能催逼谦逊。\par

\Quote{但如果是\InnerQuote{不是凭智慧的言语}，他们为什么派有口才的\Person{阿颇罗}去呢？}他回答说，并非因为信赖他的口才，而是因为他是\BibleRef{宗 18:24}\Quote{在圣经上很有能力}的，并且\Quote{驳倒了犹太人}。此外，所讨论的要点是，福音的领袖和最初传播者并非有口才的人；因为这些人正是需要在最初驱逐错误时拥有某种大能的人；而且在最初阶段，正是需要丰沛的力量。现在，祂起初可以不用受过教育的人，如果后来允许一些有口才的人加入，祂这样做并非因为需要他们，而是因为祂不分别待人。因为正如祂不需要智者来成就祂所愿意的任何事，同样，如果后来发现有这样的人，祂也不会因此而拒绝他们。\par

8. 但是，请你们证明给我看，\Person{伯多禄}和\Person{保禄}是有口才的。你们不能，因为他们\Quote{是没有学问且愚昧的人！}因此，正如基督在派遣祂的门徒进入世界之前，先已在巴勒斯坦向他们显示了自己的能力，并且说（见：\BibleRef{路 22:35}，按标准文本有\OriginalTerm{鞋}{\GreekText{ὑποδηματων}}）：\Quote{当我没有钱囊、口袋和鞋，打发你们出去的时候，你们缺少什么吗？}然后从那时起，允许他们既有钱囊又有口袋；同样，祂在这里也这样做了：因为重点是彰显基督的大能，而不是因他们外邦人的智慧而拒绝那些亲近信德的人。当希腊人指责门徒没有受过教育时，让我们比他们更积极地提出这个指责。也不要让任何人说：\Quote{\Person{保禄}是有智慧的；}相反，当我们在他们中间称赞那些智慧高超、口才出众的人时，让我们承认我们这边所有的人都是没有受过教育的；因为在那方面，他们将因我们而遭受不小的挫败：而胜利确实是辉煌的。\par

我曾说过这些事，因为我曾听见一个基督徒与一个希腊人辩论，方式荒谬，双方在互相攻击中都伤害了自己。因为基督徒本该说的话，希腊人反倒说了；而人们自然期望希腊人会说的话，基督徒却替自己争辩。例如，关于\Person{保禄}和\Person{柏拉图}的争论，希腊人竭力证明\Person{保禄}没有学问、愚昧无知；而基督徒出于单纯，急于证明\Person{保禄}比\Person{柏拉图}更有口才。结果胜利属于希腊人，因为这种论点占了上风。因为如果\Person{保禄}比\Person{柏拉图}更重要，许多人可能会反对说，他不是靠恩宠，而是靠口才取胜；所以基督徒的断言反而有利于希腊人。而希腊人所说的话却有利于基督徒；因为如果\Person{保禄}没有受过教育却战胜了\Person{柏拉图}，那么正如我所说，胜利是辉煌的；后者的门徒全体被前者吸引，尽管\Person{保禄}没有学问，却说服了他们，并将他们带到了自己一边。由此显而易见，福音不是人类智慧的结果，而是天主的恩宠。\par

因此，为避免我们陷入同样的错误，并在与希腊人辩论时被他们嘲笑；让我们指责宗徒们缺少学问；因为这样的指责本身就是赞美。当他们说宗徒们是粗鲁的，让我们更进一言，说他们也是未受过教导的、不识字的、贫穷的、卑贱的、愚钝的、无名的。这样说并非诽谤宗徒，反倒是光荣，因为这样的人竟能照耀整个世界。因为这些未经训练、粗鲁、目不识丁的人，完全战胜了有智慧、有权势、有暴君、以及那些在财富、荣耀和所有外在事物上兴盛的人，好像他们根本就不是人一样：由此显明，十字架的大能是何等伟大；这些事不是靠着人的力量完成的。因为结果并不遵循自然的进程，相反，所发生的事超越了一切自然。当任何事超越自然，且极其超越，并趋向正确和有益时，很明显，这些事是靠某种神圣的力量和合作完成的。请注意：渔夫、帐棚匠、税吏、无知者、目不识丁者，来自遥远的巴勒斯坦，在自己的土地上击退了哲学家、演说大师、善辩者，独自在短时间内战胜了他们；身处重重危险中：人民和君王的反对，自然本身的挣扎，时间的长度，根深蒂固的习俗的强烈抵抗，武装的魔鬼，列阵的恶魔煽动一切：君王、统治者、人民、民族、城市、蛮族、希腊人、哲学家、演说家、诡辩家、历史学家、法律、法庭、各种刑罚、无数各类的死亡。然而，当那渔夫发言时，所有这些都被击溃、退缩了；就像轻尘无法抵挡狂风的冲击。现在我要说的是，让我们学会这样与希腊人辩论；使我们不致像野兽和牲畜，而是准备回答\Quote{我们心中所怀的希望}\BibleRef{伯前 3:15}。且让我们停下来仔细探讨这个并非不重要的话题；让我们对他们说：弱者如何战胜强者？十二宗徒如何战胜世界？不是使用同样的盔甲，而是赤身裸体与武装的人交战。\par

试问，假如有十二个不谙战事的人，跳进一支装备精良、人数众多的军队中，他们自己不仅手无寸铁，而且体格孱弱；却未受任何伤害，即使被千万件武器攻击，也未被击伤；当箭矢击中他们时，他们赤身裸体，徒手击倒所有敌人，不用任何武器，仅以手打击，最终杀死一些人，又俘虏并带走另一些人，自己却连一个伤口也没有；有谁会认为这是出于人的能力呢？然而，宗徒们所立的战功比这更令人惊叹得多。因为一个赤身露体的人逃过伤害，远不如一个平凡而目不识丁的人——一个渔夫——能战胜如此高超的\OriginalTerm{才干}{\GreekText{δεινότητος}}那样神奇；并且，无论是因为人数稀少，还是因为贫穷，或是因为危险，或是因为习惯的成见，或是因为所命诫命的严峻，或是因为每日的死亡，或是因为受骗者众多，或是因为欺骗者的盛名，他都没有动摇其目标。\par

9. 我说，这就是我们胜过他们并与之作战的方式；我们要以生活方式，而不是以言语使他们惊异。因为这是主要的战斗，这是无可辩驳的论据，即从行为而来的论据。即使我们以言语教导万般哲理，但若我们没有展现出比他们更好的生活，则一无所获。因为吸引他们注意的不是所说的话，他们所探究的是我们所行的；他们说：\Quote{先遵守你自己的话，然后再劝诫别人。但如果你说，来世有无限的福乐，而你自己却似乎被钉在这个世界上，好像没有这些事一样，你的行为对我来说比你的言语更可信。因为当我看见你夺取别人的财物，为死者过度哀哭，在许多其他事上作恶，我怎能相信你说有复活呢？}即使人们没有用言语说出这些，他们心里也是这么想，并常常反复思量。这正是阻止不信者成为基督徒的原因。\par

因此，让我们以生活赢得他们。许多未受过教育的人，曾以这种方式使哲学家的心灵惊异，因为他们在自己身上也显出了那寓于行为中的哲学，并以他们的生活方式和克己发出了比号角更响亮的声音。因为这是比口舌更强有力的。但当我说\Quote{人不该怀恨}，却对那希腊人行各种恶事时，我怎能以言语赢得他，而我的行为却在吓跑他呢？那么，让我们以生活方式捕捉他们；并以这些灵魂建立教会，以此来积聚我们的财富。没有一样东西能与灵魂相比，即使全世界也不例外。所以，即使你给穷人无数的财宝，也比不上那使一个灵魂归化的人所做的事。\BibleRef{耶 15:19}\Quote{你若将宝贵的从卑贱中分别出来，你必像我的口。}祂如此说。怜悯穷人是大善，我承认，但这与将他们从错误中领回相比，则微不足道。因为这样做的人相似\Person{保禄}和\Person{伯多禄}；我们被允许拿起他们的福音，不是像他们那样冒着危险——忍受饥荒、瘟疫和一切其他灾祸（因为现在是和平时期）——而是要显出一种出自热忱的勤勉。因为即使我们坐在家里，也可以从事这种捕鱼。谁有朋友、亲属或同住的人，就向他说这些事，行这些事；他就要像\Person{伯多禄}和\Person{保禄}一样。我为什么说\Person{伯多禄}和\Person{保禄}呢？他将是基督的口。因为祂说：\Quote{你若将宝贵的从卑贱中分别出来，你必像我的口。}即使你今天没有说服，明天你将会说服。即使你从未说服，你也将得到你完全的赏报。即使你没有说服所有人，但许多人中的少数人说服了所有人；但他们仍然与所有人交谈，并为所有人得到赏报。因为天主通常不是按善事的结果，而是按行善者的意向赐予荣冠；即使你只付出两文钱，祂也收纳；祂在寡妇身上所行的，也要在教导者身上施行。因此，不要因为你不能拯救世界，就轻视那少数人；也不要因渴求大事，而离开小事。你若不能救一百人，就照管十人；你若不能救十人，连五人也不要轻视；你若不能救五人，不可忽略一人；你若连一人也不能救，也不可因此灰心，更不可保留你所能做的。难道你们看不到，在生意场上，那些经商的人不仅从金子上获利，也从银子上获利吗？因为如果我们不小看小事，我们也就能把握大事。但如果我们轻视小事，我们也不容易抓住大事。个人就是这样富起来的，积少成多。让我们也这样做，好使我们在一切事上富足，获得天国；借着我们的主耶稣基督的恩宠和慈爱，愿光荣、权能、尊崇借着祂，偕同祂，及圣神，归于圣父，从今世直到永远。阿们。\par

\SourceRecord{Translated by Talbot W. Chambers.；Nicene and Post-Nicene Fathers, First Series；Vol. 12.；Edited by Philip Schaff.；Buffalo, NY: Christian Literature Publishing Co.,；1889.；<http://www.newadvent.org/fathers/220103.htm>.}

% structure: p0001 source=h1 target=section reason=page-title

\section{格林多前书讲道词 第四篇}

% structure: p0002 source=h2 target=subsection reason=relative-heading-rank; base=h2

\subsection{格林多前书 1:18-20}

原来十字架的道理，为丧亡的人是愚妄，为我们得救的人，却是天主的德能。因为经上记载：\BibleQuote{我要摧毁智者的智慧，废除贤者的聪明。}\BibleQuote{智者在哪里？经师在哪里？今世的辩士在哪里？}\par

患病和奄奄一息的人，连有益健康的食物也觉得无味，亲友成了负担；往往甚至认不出他们，反把他们视为闯入者。那些在灵魂上丧亡的人，情形往往也是如此。因为他们不认识那有助于得救的事物，又认为关心他们的人是烦扰。这并非事物的本性使然，而是他们的病态所致。正如精神错乱者憎恨照顾他们的人，甚至辱骂他们，同样，不信者也如此。然而，正如对待前者，受辱者那时更同情他们，为他们哭泣，把这当作病情严重的最大症状，即他们不认识最好的朋友；同样，我们也当如此对待外邦人；是的，我们更当为他们哀哭，甚于为自己的妻子，因为他们不认识那共同的救恩。因为人对妻子的爱，应当不如我们对众人的爱，并引领他们得救；无论他是外邦人，还是别的什么人。因此，我们该为他们哭泣，因为\BibleQuote{十字架的道理，对他们确是愚妄}，而它本身是智慧与德能。因为他说：\BibleQuote{十字架的道理，为丧亡的人是愚妄。}\par

因为很可能，当十字架被希腊人嘲弄时，他们会凭藉那自以为是的智慧来抗辩和争闹，因受希腊人的言论所困扰；保禄安慰他们说：不要以为所发生的事是奇怪和不可理解的。这事物的本性就是这样，它的能力不被丧亡的人所认识。因为他们心神错乱，行为如同疯子；所以他们辱骂和厌恶那带来健康的药物。\par

2. 但是，人啊，你怎么说？基督为你成了奴隶，\BibleQuote{取了奴仆的形体}\BibleRef{斐 2:7}，并被钉死在十字架上，又复活了。当你因此该钦崇复活的祂，并赞美祂的慈爱，因为无论父亲、朋友或儿子都没有为你做的事，主——你那原是仇敌和冒犯者——却为你成就了这一切；我说，当你该因这些事赞美祂时，你却称那充满如此大智慧的事为愚妄？这并不奇怪，因为丧亡的人的特征就是不认识那导向救恩的事物。所以不要烦恼，因为真正伟大的事物被心神错乱的人嘲弄，这并非奇怪或不可理解的事。如今，怀着这种心态的人，你无法用人的智慧说服他们。不，你若要这样说服他们，反而适得其反。因为超越理性的事物只需信德。例如，我们若想用理性来证明天主如何成为人、如何进入童贞玛利亚的胎中，而不将这事交托给信德，他们只会更加嘲笑。因此，那些凭理性追问的人，正是丧亡的人。\par

我说的是天主吗？即使对于受造之物，我们这样做也会引来极大的嘲笑。假如有人想用理性弄清一切，并让他藉你的言谈说服自己我们如何看见光；你试着用理性去说服他。不，你不能；因为如果你说睁眼看见就够了，你表达的只是事实，而非方式。因为有人会说：\Quote{为什么我们不用耳朵看，用眼睛听？为什么我们不用鼻孔听，用耳朵闻？}那么，如果他对这些事疑惑，而我们又不能解释其原因，他开始发笑，难道我们不该嘲笑他吗？\Quote{因为两者都出于同一个大脑，两个器官彼此靠近，为什么它们不能做同样的工作？}我们将无法说明这不可言传和奇妙运作的原因与方法；若我们尝试，只会被嘲笑。因此，让我们把这些事交给天主的德能和无限的智慧，保持静默。\par

同样，对于天主的事，我们若想用外在的智慧来解释，就会引来极大的嘲笑，这并非由于这些事的软弱，而是由于人的愚昧。因为一切伟大的事物，没有言语可以解释。\par

3. 现在，注意：当我说\Quote{祂被钉死了}，希腊人说：\Quote{这怎能合理？祂在十字架上受折磨和痛苦时，没有帮助自己；这些事后，祂怎能复活并帮助别人？因为如果祂有能力，死前才是适当的时机。}（这其实是犹太人说过的话。）\BibleRef{玛 27:41}\Quote{祂连自己都帮助不了，怎能帮助别人？这毫无道理，}他说。不错，人啊，因为这事确实超越理性；十字架的德能是不可言传的。因为祂在惊恐之中，却显出超越一切惊恐；祂在敌人的掌握中，却能得胜；这是源于无限的德能。正如三个青年的情形：他们不进火窑也许不那么令人惊奇，但进入火窑后践踏火焰——更为惊奇；又如同约纳的情形：他若不被鱼吞下，远不如被吞下后不受那怪物伤害更为伟大——同样，论到基督：祂若不死，也许不那么不可理解；但祂死后解开死亡的束缚，才是不可理解的。所以不要说：\Quote{祂为什么在十字架上不帮助自己？}因为祂正急速要去与死亡本身作最后的搏斗。（参看 \Person{Hooker}, \WorkTitle{E. P.} v. 48. 9.）祂没有从十字架上下来，不是因为不能，而是因为不愿。因为那不被死亡的暴政所束缚的，十字架的钉子又怎能束缚祂？\par

4. 然而，这些事我们虽然知道，对不信者却尚未如此。因此他说：\BibleQuote{十字架的话，为那些丧亡的人是愚妄，但为我们得救的人，却是天主的德能。因为经上记载：我要摧毁智者的智慧，废除明者的明智。} 到此为止，他并没有提出任何可能引起反感的话，而是先引用圣经的见证，然后从那里获得勇气，采用更激烈的话语说：\par

% structure: p0011 source=h2 target=subsection reason=relative-heading-rank; base=h2

\subsection{\BibleRef{格前 1:20-21}}

\BibleQuote{天主岂不是使这世界的智慧变成了愚妄吗？智者在哪里？经师在哪里？这世代的辩士在哪里？天主岂不是使这世界的智慧变成了愚妄吗？因为世上凭自己的智慧，既然没有认识天主，天主就乐意借着宣讲的愚妄，拯救那些相信的人。} 说了\BibleQuote{经上记载：我要摧毁智者的智慧，} 接着他从事实作出论证，说道：\BibleQuote{智者在哪里？经师在哪里？} 同时指向外邦人和犹太人。因为哪一位哲学家，哪一位研究逻辑的人，哪一位熟悉犹太事物的人，曾拯救了我们并使我们认识真理？没有一个。全是渔夫的工作。\par

这样，他得出了他心中所想的结论，挫败了他们的骄傲，并说：\BibleQuote{天主岂不是使这世界的智慧变成了愚妄吗？} 他也陈述了这些事为何这样发生的理由。\BibleQuote{因为世人凭自己的智慧，既然没有认识天主，} 他说，\BibleQuote{世界凭着它的智慧没有认识天主，} 十字架便出现了。那么，\BibleQuote{在天主的智慧中}是什么意思？就是在那些祂愿意藉以显示自己的作品中显出的智慧。因为祂造它们，并使它们成为这样，是为了让人从所见之物\OriginalTerm{按比例}{\GreekText{ἀναλόγως}}学习对造物主的赞叹。难道天不大，地不广吗？那么惊叹创造了它们的那一位吧。因为这天，虽然浩大，不仅是由祂所造，而且是轻易造成的；那无边的地也是如此，好像无中生有。所以关于前者，祂说：\BibleRef{咏 101:25}（七十贤士译本）\BibleQuote{祢的手指所造的是诸天，} 关于大地，\BibleRef{依 40:23} \BibleQuote{祂使大地好像虚无。} 既然世界凭这种智慧不愿认识天主，祂便用那看似愚妄的，即福音，来说服人；不是靠推理，而是靠信德。所以，哪里有天主的智慧，就不再需要人的智慧了。因为从前，推断那创造如此伟大世界的，必须是拥有某种不可抑制、不可言喻之能力的天主，并以此认知祂——这是人的智慧的本分。但现在，我们不再需要推理，只需要信德。因为相信被钉十字架并被埋葬的那一位，并完全相信这一位自己复活并坐在高天之上；这不需要智慧，也不需要推理，只需要信德。因为宗徒们自己也不是靠智慧，而是靠信德来到的，并且在智慧和崇高上超越了外邦的智者，甚至超越得更多，因为提出争辩不如以信德接受天主的事。因为这事超越所有人类的理解。\par

但是祂如何\BibleQuote{摧毁了智慧？} 祂藉由保禄和类似他的人使我们认识，表明智慧是无用的。因为对于接受福音的宣讲，有智慧的人并不因智慧而得益，无学问的人也不因无知而受损。但若可以说一些奇妙的事，无知比智慧更适合接受那种印象，且更容易处理。因为牧人和乡下人会更快地接受它，一劳永逸地抑止所有怀疑的想法，并将自己交付给主。就这样，祂摧毁了智慧。因为智慧首先自贬，此后永远无用。因此，当她本该展示自己的适当能力并藉着工作看见主时，她却不愿意。所以，即使她现在愿意自我介绍，也不能。因为事情不是那样的；这种认识天主的方式远超过另一种。那么你看，需要信德和纯朴，我们应在各处寻求它，并优先于外来的智慧。因为他说：\BibleQuote{天主使智慧成了愚妄。}\par

但是\BibleQuote{他使成了愚妄}是什么意思？祂表明它在接受信德方面是愚妄的。因为他们以此自夸，祂就立即揭露它。因为在不能发现善事中首要者时，这算什么智慧？因此，祂使它显得愚妄，在它首先自我定罪之后。因为如果当发现可能通过推理做出时，它却证明不了什么，如今当事情在更大尺度上进行时，它又能成就什么呢？如今当只需要信德，而不需要敏锐时？那么你看，天主已表明它是愚妄的。\par

祂也乐意借着福音的愚妄来拯救；愚妄，我说，不是真正的，而是看似如此。因为更奇妙的是，祂不是引入另一种比前者更优越的智慧，而是看似愚妄的东西而获胜。例如，祂驱逐了\Person{柏拉图}，不是通过另一位更有技巧的哲学家，而是通过一个未受教育的渔夫。因为这样，失败变得更大，胜利更辉煌。\par

% structure: p0017 source=h2 target=subsection reason=relative-heading-rank; base=h2

\subsection{\BibleRef{格前 1:22-24}}

5. 接下来，为显示十字架的能力，他说：\BibleQuote{犹太人要求神迹，希腊人寻求智慧，而我们却宣讲被钉十字架的基督：为犹太人就是绊脚石，为希腊人是愚妄，但为那些蒙召的人，无论是犹太人还是希腊人，基督是天主的德能，天主的智慧。}\par

这里所说的事含义深广！因他要说明天主如何以相反的事取胜，以及福音不是出于人。他所说的大意如下：他说，当我们对犹太人说：\Quote{相信吧}；他们回答：\Quote{复活死人，治好附魔的，向我们显示征兆}。但我们却说什么呢？我们说：\Quote{那被传扬的基督被钉十字架而死了}。这不但不足以吸引不愿相信的人，甚至足以把愿意相信的人也赶走。然而，它非但不赶走，反而吸引、紧握并征服。\par

再者，希腊人要求我们提供修辞风格和诡辩的尖锐。但我们也向他们宣讲十字架：那在犹太人看来是软弱的，在希腊人看来却是愚妄。因此，当我们不仅没有满足他们的要求，反而提出了与他们要求截然相反的东西时——（因为十字架非但没有显露出是理性追寻的征兆，反而简直是征兆的毁灭；非但不被视为能力的证明，反而被视为软弱的证据；非但不是智慧的展示，反而是愚妄的暗示）——那么，当寻求征兆和智慧的人不仅没有得到他们所求的，反而听到了与愿望相反的事，并借着相反的事而被说服——那被传扬者的能力又怎能不是不可言喻的呢？这好比有人被暴风颠簸，渴望一个港口，你却向他展示的不是港口，而是另一片更汹涌的海域，如此却能使他感恩地跟随你？又如一个医生，向那受伤需要医治的人许诺，不用药物，而是再用火灼来治愈他，以此吸引那人归向自己？这实在是巨大能力的结果。宗徒们也是这样得胜的：不单没有征兆，反而借着那与一切已知征兆相反的事。基督在瞎子身上也是这样行的：因为祂要治愈他，便借着加重瞎眼的方式除去瞎眼：即抹上了泥。\BibleRef{若 9:6}正如祂用泥治愈了瞎子，同样，祂也用十字架将世界吸引到自己身边。那确实是添加绊脚石，而非除去绊脚石。在创造中祂也是如此行事，以相反的事物成就事工。例如，祂用沙围住了海，使软弱之物成为强者的缰绳。祂将大地安置在水上，安排重而密的东西承载在软而流的东西上。又借着先知，用一小片木头从水底把铁提上来。\BibleRef{列下 6:5}同样，祂也用十字架将世界吸引到自己身边。因为水承载大地，同样，十字架承载世界。你现在看，借着那些直接反对我们的事来说服人，正是大能和大智的证明。因此，十字架似乎是绊脚石；然而它远非绊脚石，反而吸引人。\par

% structure: p0021 source=h2 target=subsection reason=relative-heading-rank; base=h2

\subsection{\BibleRef{格前 1:25}}

6. 因此，保禄顾念这一切，惊奇不已，说：\Quote{天主的愚妄比人更智，天主的软弱比人更强。}这是就十字架而言，所说的愚妄和软弱并非真的，而是表面的。因为他是在回应对方的观点。因为哲学家们借着推理无法成就的事，那看似愚妄的事却成就得非常好。那么，谁是更智的呢？是说服多数人的，还是说服少数人，甚至无人说服的？是说服有关最大之事的人，还是说服关于微不足道之事的人？（\OriginalTerm{}{\GreekText{μηδὲν} \GreekText{όντων}} Reg. ms. \OriginalTerm{}{\GreekText{μη} \GreekText{δεόντων}} Bened.）柏拉图和其追随者忍受了何等大的辛劳，向我们谈论线、角和点，谈论奇数和偶数，彼此相等和不相等，以及诸如此类的蜘蛛网——（因为这些网对人的生活并不比这些题材更有用）——而他们凭借这些方法并没有给任何人带来任何大小益处，就这样结束了性命。他何等努力地试图证明灵魂是不死不灭的！但他来去空空，没有说出任何确定的话，也没有说服任何听众。但是十字架借着无学问的人实现了说服；是的，它说服了整个世界：不是关于寻常的事，而是关于天主的道理、合乎真理的虔诚、福音的生活方式，以及将来之事的审判。它使所有人成为哲学家：那些乡野之人、全然无学问的人。看哪，\Quote{天主的愚妄比人更智}，\Quote{天主的软弱比人更强}。如何\Quote{更强}？因为十字架遍行普世，以强力夺取一切，而当人们千方百计要消灭被钉者的名字时，相反的事却发生了：那名字越发兴旺昌盛，而他们却衰败消亡；活人对抗死人，却毫无能力。所以，当希腊人称我为愚拙时，他自己显得格外愚拙：因为我这被他视为愚拙的人，显然比智者更智。当他称我为软弱时，他便显出自己更为软弱。因为税吏和渔夫靠着天主的恩宠所能成就的高贵事业，哲学家、修辞学家、暴君，以及整个世界四处奔波，甚至无法构想。十字架带来了什么？灵魂不死不灭的道理，身体复活的道理，轻视现世之物的道理，渴慕未来之物的道理。它使凡人变成天使，使所有人在各处实践克己（\OriginalTerm{}{\GreekText{φιλοσοφοῦσι}}），并表现出各种刚毅。\par

7. 但在他们中间，也有人会说，许多人是轻视死亡的。告诉我谁？是那个饮毒芹的人吗？但如果你愿意，我能从教会内提出成千上万这样的人。因为假如在遭受迫害时，允许他们饮毒芹而离世，那么所有的人都变得比他更出名。此外，他饮毒芹时，并没有自由选择饮或不饮；而是无论愿意与否，他都必须承受：这显然不是刚毅的表现，而是出于必然，仅此而已。因为就连强盗和杀人犯，在受到法官的判决后，也遭受了更痛苦的事。但在我们这里，情况完全相反。因为殉道者并非违背自己的意愿而忍受，而是出于自愿，并且有自由选择不忍受；他们展现了比金刚石还坚硬的刚毅。那么，你就会看到，我提到的这个人饮毒芹并不算什么大事；因为他已无法不饮，而且当他达到很大年龄时。因为当他轻视生命时，他说自己已七十岁；如果这可以称为轻视的话。就我而言，我不能肯定这一点；而且更重要的是，没有其他人能肯定。但请给我指出一个为了虔敬而坚忍痛苦的人，正如我在世界各地指出成千上万的人。谁在指甲被拔时英勇承受？谁在身体关节\OriginalTerm{被撕裂}{\GreekText{ἀνασκαπτομένων}}时承受？谁在身体被一片片毁掉（\OriginalTerm{身体被一片片毁掉}{\GreekText{τοῦ} \GreekText{σώματος} \GreekText{κατὰ} \GreekText{μέρος} \GreekText{πορθουμένου}}；头\OriginalTerm{头}{\GreekText{τῆς} \GreekText{κεφαλῆς}}）肢解？或者头？谁在骨头被杠杆撬出（\OriginalTerm{被撬开}{\GreekText{ἀναμοχλευομένων}}）时承受？谁被不间断地放在煎锅上？谁被扔进大锅？请指出这些例子。因为死于毒芹就像一个人处于睡眠状态。不仅如此，如果传闻是真的，这种死亡甚至比睡眠更甜美。但如果他们中某些人确实忍受了折磨，那么这些人的赞美也归于无有。因为他们是在一些可耻的场合中死去的；有些人因为泄露奥秘；有些人因为图谋统治；另一些人被发现在最肮脏的罪行中；还有些人鲁莽、无益且愚蠢地，没有任何理由地自杀。但我们不是这样。因此，那些人的行为无人提及；但这些行为却兴旺且日益增长。\Person{保禄}心中记着这些，说道：\BibleQuote{天主的懦弱强于人}。\par

8. 因为福音是神圣的，甚至由此也显而易见；即，这怎么会发生在十二个无知的人身上，让他们尝试如此伟大的事？他们曾在沼泽、河流、旷野中居住；或许从未进入过城市或广场；——他们怎么会想到要对抗全世界呢？因为他们是胆小且懦弱的，记述他们的人这样指出，既不为他们辩解，也不愿掩饰他们的缺点：这本身就是真理的一个非常伟大的标志。那么，关于他们，他（指福音作者）说了什么？当基督在行了无数奇迹后被逮捕时，他们逃跑了；而留下的那位，是其余门徒的首领，却否认了。那么，当基督活着时，他们尚且经受不住犹太人的攻击；现在他死了，埋葬了，照你们所说，没有复活，也没有与他们交谈，没有激发他们的勇气——他们又怎么会起来对抗如此广大的世界呢？他们难道不会自问说：\Quote{这是什么意思？他自己不能拯救自己，他还会保护我们吗？他自己活着时没有为自己辩护，现在死了，还会向我们伸手吗？他自己活着时连一个国家都没有征服，我们能靠呼号他的名来说服全世界吗？} 请问，这一切怎么可能合理？我不说是作为可行之事，甚至作为可想象之事？由此显然，如果他们在他复活后没有看见他，没有充分体验到他的能力，他们绝不会冒如此大的风险。\par

9. 假如他们拥有无数的朋友，他们岂不立刻就会使所有朋友变成敌人，扰乱古老的习俗，并移除他们祖先的界标（\OriginalTerm{界标}{\GreekText{ὅρια}} Ms. Reg. \OriginalTerm{习俗}{\GreekText{ἔθη}} Ben.）？但事实上，他们所有人都成了敌人，包括本国人和外国人。因为即使他们因外在一切而受到尊崇，难道所有人不也会因为他们引入新的政体而厌恶他们吗？但现在他们甚至一无所有；很可能正因为如此，所有人会立即憎恨并鄙视他们。你会提名谁？犹太人？不，他们因主人所遭受的事而对犹太人怀有无法形容的仇恨。那么希腊人呢？首先，希腊人曾拒绝了一个并不比他们差的人；没有人比希腊人更清楚这一点。因为\Person{柏拉图}想要创立一种新的政府形式，或者更确切地说，政府的一部分；而且不是通过改变与神明相关的习俗，而仅仅是以一种行为准则代替另一种；却被赶出\Place{西西里}，几乎丧命。但这事没有发生：所以他只是失去了自由。如果没有某个蛮族人比西西里的暴君更温和，这位哲学家就永远无法从异乡的奴役中得救。然而，在王国的政事和宗教敬拜上进行创新并非一回事。因为后者比其他任何事都更会引起骚乱和困扰人心。因为说：\Quote{让某某人娶某女人，让（共和国的）监护人以如此这般的方式行使其监护权}，这不足以引起任何大骚乱：尤其是当这一切都被写入一本书，而立法者并不急于将提案付诸实践时。另一方面，说：\Quote{人们敬拜的不是神，而是魔鬼；被钉十字架的那位是天主}，你们很清楚这激起了多大的愤怒，人们必须为此付出多么沉重的代价，它煽起了多么大的战争火焰。\par

因为\Person{普罗塔戈拉}，他们中的一位，竟敢说：\Quote{我不知道有神}，不是走遍世界宣扬这话，而是在一个城市里，就陷入了生命的极大危险。\Place{米利都}的\Person{狄亚戈拉斯}和被称为无神论者的\Person{特奥多鲁斯}，虽有朋友和来自雄辩的影响力，并因其哲学而受人钦佩，然而这些无一能帮助他们。伟大的\Person{苏格拉底}，他在他们当中超越所有哲学家，也为此饮下毒芹，因为他在关于诸神的言论中被怀疑稍有偏离。如果单单是革新的嫌疑就给哲学家和智者，以及那些获得无限声望的人带来如此大的危险，他们不仅不能实现自己的愿望，反而被驱逐出生命和乡国；那么，当你看到渔夫对世界产生了如此大的影响，并完成了他的目的，战胜了所有化外人和希腊人，你怎能不感到钦佩和惊讶呢？\par

10. 但你可能会说：他们并没有像其他人那样引入异神。不错，你所说的正是最令人惊叹之处：革新是双重的，既要推倒那些现有的神，又要宣布被钉十字架者。因为这种念头从何进入他们心中？从何使他们对自己的事成充满信心？在他们之前的人中，他们能看到谁在类似尝试中成功了呢？难道不是所有人都崇拜邪魔吗？难道不是所有人都以元素为神吗？差异不就在于不敬的方式吗？然而，他们攻击了所有，推翻了所有，并在短时间内像某种有翼之物一样跑遍了整个世界；他们不计算危险、死亡、事情的困难、自身的稀少、对手的众多、那些与他们作战之人的权威、权力和智慧。因为他们有一位比这些都要大的盟友，即被钉十字架并复活者的能力。如果他们以字面意义（\OriginalTerm{感性战争}{\GreekText{πόλεμον} \GreekText{αἰσθητόν}}）与世俗战争，就像实际发生的那样，那倒不那么令人惊奇。因为按照战斗的法则，他们可能站在敌人对面，占据一些有利地势，相应列阵迎敌，并选择攻击和近战的时间。但这里并非如此。因为他们没有自己的营帐，而是与敌人混在一起，从而战胜了他们。即使他们在敌人中间行走，也逃脱了敌人的抓住（\OriginalTerm{抓住}{\GreekText{λαβὰς}}，或称\OriginalTerm{伤害}{\GreekText{βλαβὰς}}），变得优越，并取得了辉煌的胜利；这胜利应验了那预言：\BibleQuote{你在仇敌中必掌握权柄}\BibleRef{咏 109:2}。因为这件事充满了所有的惊奇：他们的敌人抓住他们，把他们投入监狱和锁链，却不仅没有战胜他们，反而最终自己也向他们低头：鞭打着向被鞭打者低头，捆绑者向被捆绑者低头，迫害者向逃亡者低头。所有这些事，我们都能对希腊人讲，甚至更多；因为真理丰富（\OriginalTerm{真理的丰富}{\GreekText{πολλή} \GreekText{τῆς} \GreekText{ἀληθείας} \GreekText{ἡ} \GreekText{περιουσία}}）。如果你愿意跟随论证，我们将教你对抗他们的整个方法。同时，让我们在此持守两个要点：弱者如何战胜强者？以及，他们这些人，从何产生这样的计划，除非他们享有神圣的帮助？\par

11. 到此为止，我们所说的已经足够。但让我们通过行动显明一切品行上的卓越，并炽燃美德的火焰。因为他说：\BibleQuote{你们是光，照耀在世界中}\BibleRef{斐 2:15}。天主交给我们职能，比祂交给太阳的更大：比天、地、海更大，正如属灵的事物比感性的事物更优越。当我们仰望太阳圆盘，欣赏其美丽、本体和光辉时，让我们再思：我们内的光更大、更好，正如黑暗若我们不留意也更可怕。事实上，深沉的夜压迫着整个世界。这是我们必须驱散和消融的。这夜不仅在异端者和希腊人中，也在我们中的许多教友中，在教义和生活方面。因为许多人完全不信复活；许多人用\OriginalTerm{星命}{\GreekText{γὲνεσιν}}为自己\OriginalTerm{筑起壁垒}{\GreekText{ἐπιτειχίζουσι}}；许多人执着于迷信的观察、预兆、占卜和征兆。还有些人使用护身符和咒文。但对这些事，我们以后当完成对希腊人的讲论后再讲。\par

在此期间，你们要持守所说的教导，并要与我一同在争战中合作；藉着你们的生活方式吸引他们归向我们，并改变他们。因为，正如我常说的，教导\OriginalTerm{崇高道德}{\GreekText{περὶ} \GreekText{φιλοσοφίας}}的人，首先应当以身作则，成为听众所不可或缺的。让我们因此成为这样的人，使希腊人对我们怀有好感。如果我们下定决心，不仅不行恶，而且还要忍受恶，这事便会成就。难道我们没有看见，当小孩子被父亲抱在怀中时，他们会在抱他们的人的腮帮上打耳光，而父亲多么温柔地任凭孩子发泄怒气，当他看见孩子怒气已消时，面容便喜笑颜开？同样，让我们也如此行；正如父亲对待孩子一样，让我们与希腊人交谈。因为所有的希腊人都是孩子。这一点，他们自己的作家也曾说过：\Quote{那个民族永远是孩子，没有一个希腊人是老人。}如今，孩子不能忍受为任何有用之事操心；同样，希腊人也永远想要玩耍；他们躺在地上，姿势和情感都匍匐于地。此外，孩子往往在我们谈论重要事情时，根本不理会所说的话，甚至一直发笑：希腊人正是如此。当我们谈论天国时，他们发笑。正如从婴孩口中大量滴下的唾沫，常常弄脏他们的饮食，同样，从希腊人口中流出的言语也是虚妄和不洁的。即使你给孩子必需的食物，他们仍会用恶言喋喋不休地烦扰提供食物的人，我们必须一直容忍他们（\OriginalTerm{容忍}{\GreekText{διαβαστάζεσθαι}}）。再者，孩子看见盗贼进来搬走家具时，非但不抵抗，反而对那心怀不轨的人微笑；但若有人拿走他们的小篮子或\OriginalTerm{响器}{\GreekText{σεῖστρα}}或其他任何玩具，他们便耿耿于怀，烦恼，撕扯自己，跺脚；希腊人同样如此：当他们看见魔鬼窃取他们所有的祖业，甚至维持生命之物时，他们发笑，并跑向魔鬼如同跑向朋友；但若有人拿走任何财产，无论财富或任何此类幼稚之物，他们就哭喊，撕扯自己。正如孩子不知羞耻地暴露肢体，希腊人也是如此，沉溺于淫乱和通奸，暴露自然律，引入非法的结合，而不知羞愧。\par

你们对我报以热烈的掌声和欢呼，但你们在所有的欢呼中要留意，免得你们就是那些被提及的人中间的一个。因此，我恳求你们都成为成人；因为，只要我们仍是孩子，我们怎能把男子气概教导给他们？我们怎能阻止他们脱离孩子的愚昧？让我们因此成为成人；使我们能达到基督为我们所标定的身量的尺度，并获得那未来的恩赐：藉着祂的恩宠和慈爱，等等。\par

\SourceRecord{Translated by Talbot W. Chambers.；Nicene and Post-Nicene Fathers, First Series；Vol. 12.；Edited by Philip Schaff.；Buffalo, NY: Christian Literature Publishing Co.,；1889.；<http://www.newadvent.org/fathers/220104.htm>.}

% structure: p0001 source=h1 target=section reason=page-title

\section{\BibleBook{格林多}{前书}讲道词第五篇}

% structure: p0002 source=h2 target=subsection reason=relative-heading-rank; base=h2

\subsection{格林多前书 1:26-27}

\BibleQuote{\Emph{你们看，} \Emph{弟兄们，你们蒙召，按肉身来说，有智慧的并不多，有权势的并不多，有尊贵的也不多；} 但天主却拣选了世上愚拙的，叫那智慧的羞愧。}\par

他已说过：\BibleQuote{天主的愚拙比人更有智慧}；他已经表明，人的智慧被弃绝，既因圣经的见证，也因事件的结果：因见证，他说：\BibleQuote{我要摧毁智者的智慧}；因事件，他提出一个问题说：\BibleQuote{智者在哪里？经师在哪里？}再者，他又同时证明这件事不是新的，而是古老的，因为从一开始它就被预表和预言了。因为他说：经上记载：\BibleQuote{我要摧毁智者的智慧}。他还表明，事情这样发展并非不适宜或不可解释：（因为他说：\BibleQuote{天主以其智慧，世人没有认识祂，天主就乐意藉着愚拙的宣讲来拯救那些信的人}；）并且十字架显示了不可言喻的能力和智慧，天主的愚拙远比人的智慧强大。他又进一步证明这一点，不是藉着教师，而是藉着门徒本身。因为他说：\BibleQuote{你们看，你们的蒙召}：不仅未受训练的教师，而且同样没受过训练的门徒，都是祂所拣选的；祂拣选了\Quote{按肉身来说，有智慧的并不多}（这是他的话）。所以他所谈论的事证明它既在力量上也在智慧上超越，因为它使许多人和不智慧的人信服：说服一个无知的人极其困难，尤其是在谈论伟大而必要的事情时。然而，他们确实使人信服。他邀请格林多人自己为此作证。因为他说：\BibleQuote{弟兄们，你们看，你们的蒙召}：想一想，考察一下：因为这样智慧、甚至比一切更智慧的教义，被普通人所接受，证明了教师的极大智慧。\par

2. 但\Quote{按肉身来说}是什么意思？指肉眼可见的，指今生，指外邦人的教养。然后，免得他显得自相矛盾（因为他曾使方伯\BibleRef{宗 13:12}和亚略巴古的官员\BibleRef{宗 17:34}以及阿颇罗\BibleRef{宗 18:26}信服，我们见过其他智慧的人也皈依了福音），他没有说\Quote{没有智者}，而是说\Quote{有智慧的并不多}。因为他并非\OriginalTerm{特意地}{\GreekText{ἀποκεκληρωμένως}}召叫愚昧者而忽略智慧者，而是他也接纳了这些人，但前者的人数多得多。为什么？因为按肉身来说的智者充满了极端的愚妄；当他不肯放弃自己败坏的教义时，他正是特别符合\Quote{愚拙}这个词的人。正如一个医生想要教导某些人他的医术秘密，那些懂得一点、但有一种恶劣而倒错的技艺实践方式并坚持不改的人，不会愿意安静学习；而那些什么都不知道的人，会最容易接受所说的话。这里也是如此。没有学问的人更易信服，因为他们没有那种自以为聪明的极端疯狂。因为极度的愚昧在这些人身上比在其他任何人身上更甚，这些人，我说，是将那些只能凭信心确定的事交由推理判断的人。例如，铁匠用钳子拉出烧红的铁；如果有人坚持用手去拿，我们会认为他极度愚蠢。同样，那些坚持靠自己探究这些事的哲学家们，轻视了信仰。正是因此，他们没有找到他们所寻求的任何东西。\par

\BibleQuote{有权势的并不多，有尊贵的也不多}；因为这些人也充满了骄傲。对于确切认识天主，没有什么比傲慢和\OriginalTerm{被钉住}{\GreekText{προσηλῶσθαι}}于财富更无用的了：因为这些使人仰慕眼前的事物，不关心将来；并且因着许多挂虑堵塞了耳朵。但\BibleQuote{天主拣选了世上愚拙的}，这是胜利的最大标志，即祂用来得胜的人是没有受过教育的。因为希腊人不是在被\Quote{智者}击败时感到羞愧，而是在看到工匠和集市上随便遇到的人比他们更有哲学智慧时，他们才感到困惑。因此，他自己也说：\BibleQuote{叫那智慧的羞愧}。不仅如此，他还同样做了其他人生优势的事。因为接着，\BibleQuote{天主拣选了世上软弱的，叫那强壮的羞愧}。因为祂不仅召叫没有学问的人，也召叫贫穷、卑贱和默默无闻的人，为要叫那些在高位上的人谦卑。\par

% structure: p0007 source=h2 target=subsection reason=relative-heading-rank; base=h2

\subsection{格林多前书 1:28}

\BibleQuote{天主又拣选了世上卑贱的，被人轻视的，以及那无有的，为要废掉那有的。}现在，祂称什么为\Quote{无有的}？那些因其极其渺小而被视为无有的人。这样，祂显示了祂的大能，藉着那些似乎无有的人推倒了伟大的人。他在别处也这样表达，\BibleRef{格后 12:9}：\BibleQuote{因为我的能力是在软弱中显得完全。}因为教导被遗弃的人、从未涉猎任何学问的人，忽然间对天上的事谈论智慧，这是一件伟大的能力。设想一个医生、一个演说家或任何其他人：我们最钦佩他时，是他使完全未受教育的人信服并受教。如果向一个未受教育的人灌输技艺规则是一件非常奇妙的事，那么对于如此崇高的哲学的事情，更是如此。\par

3. 但祂行此事，不仅为奇事之故，亦非只显祂自己的权能，也为抑制骄傲的人。因此祂先前既说：\Quote{为要叫那有智慧的羞愧，叫那刚强的羞愧，使那有的归于无有，}这里又说：\par

% structure: p0010 source=h2 target=subsection reason=relative-heading-rank; base=h2

\subsection{\BibleRef{格前 1:29}}

\Quote{为使凡有血肉的，在天主面前一个也不能自夸。}因为天主行这一切，正是为此：抑制虚荣与骄傲，打倒夸耀。祂说：\Quote{你们也当从事此事。}祂行一切，为使我们不归功于自己，而将一切归于天主。而你们竟将自己交付于这个人或那个人？那你们将获得什么宽恕呢？\par

因为天主亲自显明，我们单凭自己不可能得救；祂从起初就如此行。因为那时人也不能凭自己得救，而需要他们观察天空的美丽、大地的广阔及受造界的整体，从而被引导至万有的伟大工程师。祂这样行，是预先抑制将要发生的自满。正如一位师傅命令弟子无论他领到哪里都要跟随，当看见弟子抢先且渴望独自学习一切时，便任他走入歧途；及至证明他不足以获得知识，这才终于将他自己所要教导的传授给他。同样，天主起初命令人借着受造物所给予的观念来追寻祂；但因他们不愿如此，祂便以经验显明他们靠自己并不足，从而另开一条路引导他们到祂跟前。祂将世界当作一面展板；但哲学家们并未在这些事上学习，也不愿听从祂，更不愿经由祂所命令的方式接近祂。于是祂引入另一条比前者更显明的路，一条能令人信服人不能单凭自己而自足的路。因为那时，那些祂借着受造物引导的人，可能会生出理性的疑惑并运用外邦人的智慧；但现在，除非人成为愚拙，即舍弃一切理性与智慧，将自己交付于信德，否则不可能得救。你看见，除了使道路容易之外，祂也借此根除了一个不小的弊病，即禁止夸耀并心存高傲：\Quote{为使凡有血肉的不能自夸。}因为罪由此而来：人坚持要比天主的法律更明智，不愿按祂所规定的方式获得知识；因此他们根本未能获得知识。从起初也是如此。祂对亚当说：\Quote{你当行这，那样的事你不可行。}他以为能发现更多，于是悖逆，反而失去了他所拥有的。祂对后来的人说：\Quote{不要停留在受造物上；而应借着它默观造物主。}他们却好像要造出比所命令的更为明智的东西，发动了无数的曲折。于是他们彼此碰撞，也撞向自己，既未找到天主，对受造物也没有清晰的认识，更没有合宜而真实的想法。因此，祂又极其有力地\OriginalTerm{极其有力地}{\GreekText{ἐκ} \GreekText{πολλοῦ} \GreekText{τοῦ} \GreekText{περίοντος}}压下他们的狂妄，让未受教育的人先行进入，借此显明所有人都需要从上而来的智慧。不仅是在知识方面，而且在所有其他事上，祂也将人及一切受造物构设成如此需要祂，使他们也有这极有力的服从与依附的动机，免得他们离弃而丧亡。为此，祂没有使他们足以自足。因为即使如今，许多人尽管缺乏仍藐视祂，倘若不是这样，他们将在傲慢中漂泊至何处呢？因此，祂制止他们自夸，并非出于嫉妒，而是为引他们脱离由此而来的毁灭。\par

% structure: p0013 source=h2 target=subsection reason=relative-heading-rank; base=h2

\subsection{\BibleRef{格前 1:30-31}}

4. \Quote{但你们得以在基督耶稣内，是由天主；基督耶稣成了我们的智慧、正义、圣化与救赎。}\par

我想祂在此用\Quote{由天主}一词，并非指我们的被造，而是指信德而言：即我们成为天主的子女，\Quote{不是由血，也不是由肉欲，也不是由人欲}\BibleRef{若 1:13}。\Quote{那么，你们不要以为祂夺去了我们的夸耀，便留下我们如此：因为另有更大的夸耀，即祂的恩赐。因为你们已是祂的子女，而在祂面前原不宜夸耀；你们乃借着基督而成为如此。}既然祂说过：\Quote{天主选了世上愚拙的、卑贱的，}祂便表明他们比众人更尊贵，因有天主为父。而这尊贵的原因，不是这人或那人，乃是基督，祂使我们成为有智慧的、正义的、圣洁的。因为这些话的意思便是：\Quote{祂成了我们的智慧。}\par

那么，谁比我们更聪明呢？我们并无\OriginalTerm{柏拉图}{Plato}的智慧，却拥有基督自己——这是天主所愿的。\par

但\Quote{由天主}是什么意思？每当祂论及\OriginalTerm{独生子}{Only-Begotten}的尊大时，祂便加上对圣父的提及，免得有人以为圣子非受生。既然祂确认了圣子的权能如此伟大，并将一切归于圣子，说祂\Quote{成了我们的智慧、正义、圣化与救赎}——祂再次借着圣子将一切归于圣父，所以说\Quote{由天主}。\par

但为什么祂不说\Quote{祂使我们成为智慧}，却说\BibleQuote{祂成了我们的智慧？}以显示恩赐的丰盈。仿佛祂说：祂将自己赐给了我们。且看祂如何依次进行。因为首先，祂使我们因脱离谬误而成为智慧；然后，因赐给我们圣神而使我们成为正义和圣洁；祂如此将我们从一切邪恶中拯救出来，以致我们\Quote{属于祂}，这话并非表达\OriginalTerm{本体的相通}{\GreekText{οὐσιώσεως}}，而是关乎信德而言。在别处我们见祂说：\BibleQuote{我们在祂内成为正义}；并以这些话说：\BibleQuote{祂使那不认识罪的，替我们成了罪，好叫我们在祂内成为天主的正义}；\BibleRef{格后 5:21}但如今祂说：\BibleQuote{祂为我们成了正义，以致凡愿意的都可以丰富地分享。}因为使我们成为智慧的，不是这个人或那个人，而是基督。因此，\BibleQuote{夸耀的，该在主内夸耀}，不在这样或那样的人内。一切都出于基督。因此，说了\BibleQuote{祂成了我们的智慧、正义、圣化、救赎}之后，祂又加上：\Quote{为应经上所记载的：\BibleQuote{夸耀的，该在主内夸耀。}}\par

为此，他也曾猛烈抨击希腊人的智慧，为教训人这一功课（\OriginalTerm{正是为此}{\GreekText{τοῦτο} \GreekText{αὐτὸ}}萨维勒本；\OriginalTerm{借着这个}{\GreekText{τούτῳ} \GreekText{αὐτῶ}}本笃本），别无其他：即（这确实公正）他们应在主内夸耀。因为当我们凭自己寻求超越我们的事物时，没有比我们更愚拙、更软弱的。在这种情况下，我们可能有磨利的舌头，但却无法拥有教义的稳固。相反，单独的推理如同蛛网。因为有些人竟疯狂到如此地步，以至于说整个存在中没有真实的东西：是的，他们坚称一切事物都与表象相反。\par

所以不要说任何事出于你自己，而要在一切事上荣耀天主。任何时候都不要归功于人。因为如果连\Person{保禄}都不应归功，那么别人更不应归功。因为他说：\BibleRef{格前 3:6}\BibleQuote{我栽种了，阿颇罗浇灌了，但使之生长的，却是天主。}那学会在主内夸耀的，决不会高傲，而会时常谦和，并在一切境况下感恩。但希腊人的心思并非如此；他们将一切归于自己，因此甚至将人当作神。如此狂妄的傲慢使他们陷入了极大的羞耻（\OriginalTerm{陷入羞耻}{\GreekText{ἐξετραχήλισεν}}）。\par

5. 现在是时候了，在剩下的部分，去与这些作战。回想我们前一日讲到何处。我们曾说，按人的因果来看，渔夫们胜过哲学家是不可能的。然而它却成了可能：由此清楚可见，是借着恩宠成为可能的。我们曾说，他们甚至不可能想象如此伟大的功业；并且我们表明，他们不仅想象了，而且极其容易地完成了。今天我们来处理同一个论点：即，除非他们看见了复活后的基督，否则他们心中怎么会期望战胜世界？什么？他们疯了吗，轻率随意地指望这样的事？因为如果没有天主的恩宠，而指望如此伟大事业的成就，那甚至超越了一切疯狂。如果他们癫狂失常，又怎能成功呢？但如果他们神志清醒——如事实所显明的——他们若不是从天上得到可靠的许诺，并且享受从上而来的影响，又怎能出发从事如此伟大的战争，冒险对抗天地，赤身露体而如此英勇地坚守阵地，为改变世上长久以来固定的习俗，而他们不过是十二个人呢？\par

而且，更有甚者，什么使他们期望通过邀请人们到天堂和天上的居所而说服听众呢？即使他们是在尊荣、财富、权力和学识中长大的，也不大可能会被激发去承担如此沉重的使命。然而，他们的期望会更有一些理由。但就实际情况而言，他们中有些人曾从事湖泊、皮革、税关等职业；没有比这些追求更不利于哲学，更不利于说服人拥有崇高思想的；尤其当一个人没有榜样可循时。不，他们不仅没有使成功显得可能的榜样，反而有完全否定成功可能性的榜样，而且这些榜样就在他们家门口（\OriginalTerm{近在眼前}{\GreekText{ἔναυλα}}）。因为许多企图革新的都被彻底消灭了——我不是说希腊人中，那算不了什么——而是说就在那时的犹太人中；他们不是用十二个人，而是用大批人从事这工作。例如，\Person{特敖达}和\Person{犹达斯}，带领大批人，结果连他们的门徒一起灭亡了。这些榜样所引起的恐惧足以阻碍他们，除非他们坚信没有天主的能力胜利是不可能的。\par

是的，即使他们果真期望得胜，他们是怀着怎样的希望才冒如此大的危险，除非他们着眼于来世？但让我们假设他们至少期望胜利；那么他们从引领众人归向祂（你们说\Quote{祂没有复活}）中期望得到什么？因为如果现在那些相信天国和无数祝福的人尚且不情愿地面对危险，那么他们又怎能白白地——甚至是为恶——忍受如此多呢？因为如果所发生的事并未发生，如果基督没有升天；那么，他们顽固热心地捏造这些事并说服全世界相信它们，必定是得罪了天主，并且必会期望从天上降下万道雷霆。\par

6. 或者从另一角度看：如果他们在基督活着时怀着如此大的热忱，然而在祂死后，这份热忱就会熄灭。因为若祂没有复活，在他们眼中，祂就不过是一个欺骗者和冒充者。你们不知道吗？当统帅和君王活着时，即使军队软弱，仍能团结一致；但当身居此位者离世后，无论他们多么强大，都会土崩瓦解？\par

那么请问，当他们将要接受福音并出发前往普世时，是哪些诱人的论据促使他们行动？难道缺乏任何阻止他们的障碍吗？如果他们疯了（因为我不愿停止重复），他们根本不可能成功；因为没有人会听从疯子的建议。但如果他们成功了，正如他们确实成功，事件也证明了，那么就没有人比他们更智慧。既然没有人比他们更智慧，那么显然，他们不会轻易开始宣讲。若非他们看到祂复活后，有什么足以吸引他们投入这场战争？有什么不会使他们退避三舍？祂对他们说：\Quote{三天后我要复活}，并许下了天国。祂说，他们领受圣神后，将征服整个世界；此外还有成千上万超越一切本性的事。所以，如果这些事一件也没有发生，尽管他们活着时相信祂，死后他们也不会相信祂，除非他们看到祂复活。因为他们会说：\Quote{祂说\InnerQuote{三天后我要复活}，却没有复活。祂许诺要赐下圣神，却没有派遣。那么祂关于另一世界的言论怎能令我们信服呢？当祂关于今世的言论经检验而落空时？}\par

再者，如果祂没有复活，他们为何宣讲祂复活了？\Quote{因为他们爱祂，}你会说。但显然，他们后来应会恨祂，因为祂欺骗并出卖了他们；祂用无数的希望抬举他们，使他们离开家庭、父母和一切，并使他们与整个犹太民族为敌，最后却背叛了他们。如果这事是出于软弱，他们或许能原谅；但现在它会被视为极端恶意的结果。因为祂本应说实话，而不是许诺天堂，既然祂只是一个凡人，如你们所说。因此，他们最可能采取的做法恰恰相反：揭露骗局，宣告祂是个冒充者和骗子。这样，他们本可摆脱一切危险，结束战争。此外，既然犹太人给士兵钱让他们说偷走了尸体，如果门徒们站出来说：\Quote{我们偷了祂，祂没有复活，}他们会得到多大的荣耀？这样，他们本可以享受荣誉，甚至被加冕。那么，他们为何为了侮辱和危险而放弃这些呢？除非是某种神圣的力量影响了他们，比这一切更强大。\par

7. 如果你们仍不信服，也请考虑这一点：若非如此，即便他们再怎么善意，也不会以祂的名宣讲这福音，反而会憎恶祂。因为你们知道，甚至那些欺骗我们的人的名字，我们也不愿听。但他们为何又宣讲祂的名？指望通过祂获得胜利吗？他们本应期待相反的情况：即使他们快要得胜，也会因提出一个欺骗者的名字而自取灭亡。如果他们想遮掩先前的事件，他们应该保持沉默；极力为之辩护只会激起更严重的敌意和嘲笑。那么，他们是从哪里想到编造这些事的呢？我说\Quote{编造}：因为他们已忘记了所听到的事。但是，如果当没有恐惧时，他们忘记了许多事，有些甚至不明白（如同福音作者自己所说），那么当如此大的危险降临时，他们怎可能不把一切都抛在脑后呢？我何必提及言语？就连他们对师傅的爱，也因对将来的恐惧而逐渐消褪；祂也曾为此责备他们。因为此前，他们依恋祂，不断问：\Quote{祢往哪里去？}但此后，当祂详细讲述，并宣告在十字架时和十字架后将临到他们的恐怖时，他们因恐惧而沉默僵住了——请听祂如何指出这一点：\BibleRef{若 16:5}\BibleQuote{你们中没有人问我：祢往哪里去？只因我给你们说了这些，你们心中就充满了忧愁。} 如果祂即将死亡并复活的期待已使他们如此忧伤，那么若他们未见祂复活，这忧伤岂不成了毁灭？是的，他们宁愿沉入地底深处，既因被欺骗而沮丧，又因恐惧未来而深感窘迫。\par

再者，他们高深的教义从何而来？因为更高深的事，祂说他们以后才会听到。因为祂说：\BibleRef{若 16:12}\BibleQuote{我还有许多事要告诉你们，然而你们现在不能肩负。}因此，未说的话是更高深的。其中一个门徒甚至不愿与祂一同往犹太地去，当听到危险时，却说：\BibleQuote{我们也去，好同祂一起死。}\BibleRef{若 11:16} 他因预料会死而难以接受。如果那门徒在祂身边时，预料会死并因此退缩，那么与他分开后，当其他门徒也离开，那无耻行为的暴露如此彻底时，他还会期待什么呢？\par

8. 此外，他们出去时有什么可说的呢？因为全世界都知道祂的受难：祂是在正午时分，在京城，在最重要的节日——一个无人可以缺席的节日——被悬挂在高处的\OriginalTerm{木架}{\GreekText{ἰκρίου}}上。但那些外人中没有人看见祂的复活，这对他们说服人来说是一个不小的障碍。再者，祂被埋葬是人人皆谈的话题；士兵和所有的犹太人都说祂的门徒偷走了祂的身体；但祂已经复活了，那些外人中没有一个人亲眼看见。那么，他们凭什么指望说服世界呢？因为，当奇迹发生时，某些士兵被说服作了相反的见证，那么这些人凭什么指望在没有奇迹的情况下完成宣讲者的工作，并且身无分文地说服陆地和海洋关于复活的事呢？再者，如果他们是出于对荣誉的渴望而尝试此事，那么他们就会各自把教义归给自己，而不是归给那位已经死去的人。或许有人说，人们不会相信他们？那么，这两者中哪一个更可能赢得人们的相信呢？是那位被捕并被钉十字架的人，还是那些逃脱了犹太人手的人？\par

9. 其次，请告诉我，他们这样做是出于什么目的？他们没有立即离开犹太进入外邦人的城市，而是在其境内来回走动。但是，除非他们行了奇迹，否则他们如何说服人呢？因为如果他们真的行了奇迹（而他们确实行了），那是出于天主的德能。另一方面，如果他们没有行任何奇迹却仍然成功了，那么这件事就更奇妙了。他们难道不知道犹太人——请告诉我——他们的恶行和充满怨恨的灵魂吗？因为在徒步渡过红海之后，在胜利之后，借着梅瑟的手对奴役他们的埃及人树立起那无血的奇妙胜利标记之后，在玛纳之后，在岩石及其涌出的溪流源泉之后，在埃及地、红海和旷野的万千奇迹之后，他们甚至用石头砸梅瑟\BibleRef{户 14:10}。他们将耶肋米亚投入坑中，杀害了许多先知。例如，请听厄里亚在那可怕的饥荒、奇妙的雨、他从天上降下的火炬和奇异的全燔祭之后，被赶到他们国家极远边境时所说的话：\BibleQuote{主，他们杀了你的先知，拆毁了你的祭台，只剩下我一个人，他们还在寻索我的性命。}\BibleRef{列上 19:10}然而那些受迫害的人并没有扰乱任何既定的规矩。那么，请告诉我，人们有什么理由相信我们现在所谈论的这些人呢？因为一方面，他们比任何先知都更卑微；另一方面，他们引入的正是那些导致犹太人将他们的师傅钉在十字架上的新鲜事物。\par

并且从另一方面看，基督说出这些话比他们说出这些话似乎更不难以理解；因为祂，他们可能会想，这样做是为了为自己获取光荣；但这些人，他们会更加憎恨，因为他们在为另一个人与犹太人作战。\par

10. 但是，罗马的法律帮助他们吗？不，这些法律反而使他们陷入更大的困境。因为他们的说法是\BibleRef{若 19:12}：\BibleQuote{凡自命为君王的，就不是凯撒的朋友。}所以仅此一点就足以阻碍他们：他们首先是那位被视为篡位者的人的门徒，然后又想加强祂的事业。那么，到底是什么促使他们冲入如此巨大的危险？而且，关于祂的哪些话会使他们获得信任？是祂被钉十字架？是祂出生于一位与犹太木匠有婚约的贫穷犹太妇女？是祂来自一个被世界仇恨的民族？不，所有这些不仅不能说服和吸引听众，反而足以使每个人反感；尤其是当这些话是由制帐幕者和渔夫说出来的时候。那么，门徒们难道不会记住这一切吗？胆怯的本性会想象得比现实更可怕，而他们的本性正是如此。那么，他们凭什么希望成功呢？不，更确切地说，他们根本没有希望，因为有无数的因素会使他们退缩，如果基督没有复活的话。岂不是连最轻率的人也看得很清楚，除非他们享有丰沛而强大的恩宠，并得到复活的保证，否则他们不仅无法去做和承担这些事，甚至连想都想不到？因为即使有如此巨大的阻碍妨碍他们的计划——我指的不是成功，而是计划——他们却仍然计划并付诸实施，成就了远远超乎一切预期的事，我想每个人都能看出，他们不是靠人的能力，而是靠天主的恩宠成就了这些事。\par

这些论证，我们不仅应当独自操练，也应当彼此操练；这样，其余之事的发现对我们也就会更容易。\par

11. 你不可因为自己是工匠，就以为这类操练超出了你的范围；因为连保禄也是制帐幕者。\par

\Quote{是，}有人说，\Quote{但那时他也充满丰盛的恩宠，并由此说出了一切。}然而，在这恩宠之前，他曾在\Person{加玛里耳}脚前；不仅如此，他之所以领受恩宠，是因为他表现出配得上恩宠的心志；此后他又重新操起手艺。因此，凡有手艺的人都不应羞愧；羞愧的应是那些无所事事、游手好闲、雇佣众多仆役、被大批随从服侍的人。因为靠持续辛劳维持生计是一种苦修。(\OriginalTerm{哲学的一种形式}{\GreekText{φιλοσοφίας}?\GreekText{ἶδος}}，参：Hooker, E. P. V. lxxii. 18.)这样的人灵魂更清澈，心智更敏锐。因为无所事事的人更容易胡言乱语、胡乱行事，终日忙碌于虚无，完全被巨大的麻木占据。而有事可做的人不容易在行为、言语或思想上容纳任何无用之物；因为他整个灵魂都专注于劳苦的谋生之道。因此，我们不可轻视那些靠双手劳动养活自己的人，反而应因此称他们有福。请告诉我，当你从父亲那里得到产业，却不从事任何职业，反而随意挥霍殆尽时，你还有什么可感谢的？难道你不知道我们不会都接受同样的审判吗？那些在此生享有更大自由的人，将受更严格的审判；而那些受劳苦、贫穷或其他类似痛苦的人，审判则不那么严厉。这一点从\Person{拉匝禄}和富人的故事中可以明显看出。正如你因忽视正当使用闲暇而受责难，同样，穷人因勤勉工作并将剩余时间用于正当事务，将获得极大的冠冕。但你是否主张，军人的职责至少可以为你开脱？并以此抱怨你缺少闲暇？这种借口毫无道理。因为\Person{科尔乃略}是百夫长，但军人的腰带丝毫没有妨碍他严格的生活准则。而你，在假期与舞女和戏子厮混，将整个生命浪费在舞台上时，从未想到用服兵役的必要性或对统治者的恐惧来为自己开脱这些享乐；但当我们召唤你到教会时，却出现了无穷无尽的障碍。\par

在那一天，当你看见火焰、火河和永不松脱的锁链，并听到切齿声时，你将说什么？谁能在那一天为你辩护？当你看见亲手劳作、正直生活的人享受一切荣耀，而你自己，如今身着软衣、香气四溢，却陷入无可救药的痛苦时，你的财富和过剩对你有什么益处？而工匠的贫穷又能给他带来什么害处？\par

因此，为使我们将来不受苦，让我们现在就畏惧所警告的一切，并将我们的全部时间用于真正必要的事务上。这样，如果我们为过去的罪过求得了天主的宽恕，并在未来增添善行，我们就能获得天国：藉着他的仁慈与恩爱等。\par

\SourceRecord{Translated by Talbot W. Chambers.；Nicene and Post-Nicene Fathers, First Series；Vol. 12.；Edited by Philip Schaff.；Buffalo, NY: Christian Literature Publishing Co.,；1889.；<http://www.newadvent.org/fathers/220105.htm>.}

% structure: p0001 source=h1 target=section reason=page-title

\section{格林多前书第六篇讲道}

% structure: p0002 source=h2 target=subsection reason=relative-heading-rank; base=h2

\subsection{\BibleBook{格林多}{前书} 2:1-2}

\BibleQuote{弟兄们，当我来到你们那里时，并不是用高超的言论或智慧，给你们宣讲天主的奥迹；因为我曾决定，在你们中除了耶稣基督，就是那被钉在十字架上的基督以外，什么都不知道。}\par

再没有比保禄的心灵更善于战斗的了；或者不如说，我该说并非他的心灵（因为这一切不是他自己发明的），而是那在他内工作、战胜万物的恩宠，无物可与之相比。因为前面所说的已足以打倒那些夸耀智慧之人的骄傲；其实，甚至其中一部分已经足够。但是，为了增强胜利的光彩，他重新论证他所肯定的一切，踩踏倒地之敌。我们且看他是如何做的。他引出了那预言说：\Quote{我要摧毁智慧人的智慧。}他显示了天主的智慧：藉着那似乎是愚妄的事，摧毁了外邦人的哲学；他显示了\Quote{天主的愚妄比人还智慧}；他显示了：祂不仅藉无知的人施教，也拣选无知的人向祂学习。现在他表明：所宣讲的事本身以及宣讲的方式，都足以使人动摇，然而却没有使人动摇。正如他所说的：\Quote{不仅门徒们没有学问，连我这个宣讲者本人也是如此。}\par

因此他说：\Quote{而且，弟兄们，}（他再次使用\InnerQuote{弟兄们}这词，以缓和话语的严厉）\Quote{我来到你们那里时，并不是用高超的言论，给你们宣讲天主的奥迹。}\Quote{那么，请告诉我，如果你选择用高超的言论，你能否做到呢？}\Quote{我确实，如果选择的话，是不能做到的；但基督如果选择，是能做的。然而祂不愿意，为要使祂的胜利更加辉煌。}因此，在先前的一段中，他表明这是祂所做的工程，祂的旨意是要用没有学问的方式宣讲圣言，他说：\Quote{因为基督派遣我，不是为施洗，而是为传福音，且不用智慧的言词。}但保禄愿意这件事，远不如基督愿意这件事——是的，远远不如。\par

\Quote{因此，}他说，\Quote{我不是靠雄辩的展示，也不是用外在的论据来武装自己，而向你们宣讲天主的奥迹。}他没有说\Quote{宣讲}，而是说\Quote{天主的奥迹}；这个词本身足以使他收敛。因为他四处宣讲死亡；为此他加上：\Quote{因为我曾决定，在你们中除了耶稣基督，就是那被钉在十字架上的基督以外，什么都不知道。}这就是他想要表达的意思：他完全没有外在的智慧；正如他上面所说的：\Quote{我不是用高超的言论来的。}因为他显然也可能拥有这种技能；因为那凭他的衣服就能使死人复活、凭他的影子就能驱逐疾病的人，他的灵魂岂不更能接受雄辩？因为这是一种可以学习的东西；而前者超越一切技巧。那么，一个知道超乎技巧之事的人，必然更有能力做更小的事。但基督不允许，因为这不合适。因此他正确地说：\Quote{因为我曾决定什么都不知道：}因为我本人与基督有同样的旨意。\par

在我看来，他向他们说话的语气似乎比对任何人都低，为要抑制他们的骄傲。因此，\Quote{我决定什么都不知道}这句话是与外在的智慧相对而言的。因为我来不是编织三段论或诡辩，也不是对你们讲别的，而是\Quote{基督被钉在十字架上}。他们确实有无穷无尽的话要说，论及无穷无尽的事，他们展开长长的言词之链，构造论据和三段论，编织无止境的诡辩。但我到你们这里来，只说一件事：\Quote{基督被钉在十字架上}，而我胜过了他们所有人；这无疑表明了我所宣讲的那一位的能力，非言语所能形容。\par

% structure: p0008 source=h2 target=subsection reason=relative-heading-rank; base=h2

\subsection{\BibleBook{格林多}{前书} 2:3}

2. \Quote{而且，我在你们那里时，又软弱，又恐惧，又战栗。}\par

这又是另一个主题：因为不仅信徒是没有学问的人；不仅宣讲者是没有学问的；不仅整个教导的方式都带着没有学问的特点；不仅所宣讲的事本身就足以令人动摇（因为所传的讯息是十字架和死亡）；而且还有其他的障碍：危险、阴谋、每日的恐惧、被追捕。因为\Quote{软弱}这个词，在他许多地方是指迫害而言；正如在别处所说：\Quote{你们并没有轻视我肉体上的软弱}\BibleRef{迦 4:13}；又\Quote{如果我必须夸耀，我就夸耀关于我软弱的事}\BibleRef{格后 11:30}。什么样的[软弱]？\Quote{阿勒特王的总督把守了大马士革城，要捉拿我}\BibleRef{格后 5:32}。又\Quote{所以我喜欢软弱}\BibleRef{格后 12:10}；然后，他解释是什么软弱，加上\Quote{在凌辱中，在急需中，在迫害中。}他在这里也作了同样的陈述；因为他说了\Quote{而且我在软弱中}等等之后，并没有停在这里，而是解释\Quote{软弱}这个词，提到了他的危险。他又加上：\Quote{又恐惧，又战栗，我在你们那里时。}\par

\Quote{你怎么说？保禄也害怕危险吗？}他确实害怕，而且极其恐惧；因为尽管他是保禄，但他仍然是人。但这不是对保禄的指责，而是人性的软弱；这反而赞扬了他坚定的心志，即当他甚至惧怕死亡和鞭打时，他并未因这种恐惧而做错事。因此，那些声称他不惧怕鞭打的人，不仅没有尊敬他，反而大大削减了他的赞美。如果他不惧怕，那么在承受危险时，有什么忍耐或自制呢？以我而言，正因如此我钦佩他；因为他虽在恐惧中，且不仅仅是\Quote{恐惧}，甚至在他危险中的\Quote{战栗}中，他仍如此奔跑，以致永远保持他的冠冕；并在任何危险面前都不屈服，在他净化世界的任务中，无论海陆，到处播种福音。\par

% structure: p0012 source=h2 target=subsection reason=relative-heading-rank; base=h2

\subsection{\BibleRef{格前 2:4}}

3. \Quote{我的言论和我的宣讲，不在于智慧动听的言词：}也就是说，没有外来的智慧。如果所宣讲的道理并不深奥，蒙召的人又未受过教育，宣讲者也是如此，再加上迫害、战栗和恐惧；告诉我，他们如何能不靠天主的能力而获胜？正因为如此，在说了\Quote{我的言论和我的宣讲，不在于智慧动听的言词}之后，他又加上\Quote{而是在于圣神和能力的明证。}\par

你看出\InnerQuote{天主的愚妄比人智慧，天主的软弱比人刚强}吗？他们这些未受过教育的人，宣讲这样的福音，在锁链和迫害中战胜了迫害他们的人。靠什么呢？难道不是靠他们提供那来自圣神的证据吗？因为这确实是公认的明证。因为，告诉我，谁看见死人复活、魔鬼被驱逐之后，还能不承认呢？\par

但既然也有骗人的奇事，比如巫士的，他也消除了这个怀疑。因为他不是简单地说\Quote{能力}，而是先说\Quote{圣神}，然后\Quote{能力}：表明所行的事是属神的。\par

因此，福音不是借着智慧来宣讲，这并不是贬低，反而是极大的荣美。因为这将被认可为最清楚的标记，证明它是神圣的，根源于天上。所以他又加上：\par

% structure: p0017 source=h2 target=subsection reason=relative-heading-rank; base=h2

\subsection{\BibleRef{格前 2:5}}

\BibleQuote{为使你们的信德不在于人的智慧，而在于天主的能力。}\par

你看见他如何清楚地从各方面指出这种\Quote{无知}的巨大收获和这种\Quote{智慧}的巨大损失吗？因为后者使十字架落空，而前者宣扬了天主的能力；后者除了使他们未能发现他们最需要的任何事物外，还使他们自夸；前者除了使他们接受真理外，还使他们以天主为荣。再者，智慧会使许多人怀疑这道理是出于人；而这清楚地证明它是神圣的，从天降下。当用言语的智慧来证明时，甚至更差的一方常常战胜更好的一方，因为更有辞令；虚假超越真理。但在此情形下并非如此：因为圣神不进入不洁的灵魂，即使进入，也永远不会被征服；即使所有可能的言辞巧妙来攻击它。因为靠着工作和奇迹的证明远比言语的证明更为明显。\par

4. 但也许有人会说：\Quote{如果福音要得胜，不需要言语，免得十字架落空；那么为什么现在不行奇迹呢？}为什么？你是在不信中说话，不承认即使在宗徒时代也行过奇迹，还是你真想求问？如果是不信，我将首先反驳这一点。我说，如果那时没有行奇迹，他们怎么被追赶、迫害、战栗、锁链加身，成为世界的公敌，成为众人施暴的目标，而且没有任何吸引人的东西——没有言语、外表、财富、城市、民族、家族、职业\OriginalTerm{追求}{\GreekText{ἐπιτήδευμα}}、荣耀等类似之物；反而面临一切相反的东西——无知、卑贱、贫穷、仇恨、敌意，并反对整个国家，且宣布这样的信息；他们如何能说服人呢？因为诫命带来许多辛劳，教义带来许多危险。而听到并要服从的人，是在奢侈、醉酒和极大的邪恶中长大的。告诉我，他们如何说服人？他们从哪里获得可信度？因为我刚说过，如果不用奇迹他们就能说服人，那奇迹就更大了。所以不要用现在不行奇迹的事实，作为那时没有行奇迹的证据。因为正如那时行奇迹是有用的，现在不再这样行了。\par

也不必然因为现在只有言语是说服的工具，就说现在的\Quote{宣讲}在于\Quote{智慧}。因为当初播种圣言的人都是没有学问的外行人\OriginalTerm{外行人}{\GreekText{ὶδιῶται}}，没有受过教育，他们不说自己的话；而是将从天主那里领受的东西分施给世界。我们自己在此时也不引入自己的发明；而是将从他们那里领受的东西告诉众人。甚至现在我们也不是靠着辩论说服人；而是从圣经和那时所行的奇迹中引证我们所说的话。另一方面，他们那时也不仅靠奇迹说服人，也靠宣讲。奇迹和旧约圣经的见证，而不是所说之事的巧妙，使他们的言语显得更有力量。\par

5. 你会说，既然如此，为何那时奇迹有益，而今却不适宜？让我们假设一个情形（因为我还在与希腊人争辩，所以我对必然要发生的事作假设性陈述），我说，让我们假设一个情形；让不信者同意相信我们的肯定，即便是\OriginalTerm{姑且附和}{\GreekText{κἄν} \GreekText{κατὰ} \GreekText{συνδρομήν}}：例如，基督将要来临。那时基督与祂的众天使一同来临，显明为天主，万有都归顺于祂，连希腊人也不会相信吗？很明显，他也会俯伏朝拜，宣认祂是天主，即使他的顽固超乎想象。因为谁看见天开了，祂驾云降临，周围有天上万军集合，火河涌来，众人都站立战栗，会不俯伏在祂面前，相信祂是天主呢？那么，告诉我，那朝拜与认识能算作希腊人的信德吗？绝不，绝不可算。为何不能？因为这不是信德。因为这是必然性造成的，是所见之事的明证，并非出于选择，而是因景象的浩大，心灵的能力被拖曳。由此可见，事件过程越明显、越压倒性，信德的分量就越减少。正因如此，如今不行奇迹。\par

为了证实这个真理，请听祂对多默所说的话：\BibleRef{若 20:29} \BibleQuote{那些没有看见而相信的，是有福的。}因此，奇迹展示的明证越大，信德的赏报就越减少。所以，如果现今也行奇迹，同样的事也会发生。因为那时我们不再凭信德认识祂，保禄已指明说：\BibleQuote{因为我们现今是藉信德行走，并非凭眼见。} \BibleRef{格后 5:7} \EditorialNote{不在公认的经文内}。正如那时，即使你相信，也不会计入你的账，因为事情如此明显；同样现在，假若行出从前那样的奇迹。因为当我们接受那些丝毫不能凭推理理解的事时，那才是信德。正因如此，地狱被警告，却未被展示；因为若展示了，同样的事又会发生。\par

6. 此外，如果你寻求的是奇迹，即使现在，你也能看到奇迹，尽管不是同一种类：无数的预言，涉及无尽多样的主题；世界的归化；野蛮人的\OriginalTerm{哲学}{\GreekText{φιλοσοφίαν}}生活；从野蛮习俗的改变；虔诚的更大强度。\Quote{什么预言？}你会说。\Quote{因为刚才提到的所有事情都是在现今状况开始后才写下来的。}何时？何处？何人？告诉我。多少年前？五十年，还是一百年？那时，一百年前，他们根本没有写任何东西。那么，世界如何保存教义及其余一切，既然记忆不足够？他们如何知道伯多禄\OriginalTerm{被钉十字架}{\GreekText{ἀνεσκολοπίσθη}}？那些在事件发生之后的人，怎么会想到预言，例如，福音要传遍全世界？犹太制度应终止，永不恢复？那些为福音舍弃性命的人，怎能容忍看到福音被篡改？作者们如何获得信任，既然奇迹已经停止？著作如何能传入野蛮人、印度人区域，直至海洋的边界，如果讲述者不值得信任？作者们，他们是谁？他们何时、如何、为何写作？是为了给自己挣得荣耀吗？那为何他们用别人的名字来题写这些书？\Quote{为什么？是出于推荐教义的愿望。}作为真的，还是假的？因为如果你说，他们坚持它，作为假的；他们加入它完全不可能；但如果作为真理，就不需要像你所说的那些编造。此外，预言是那样的一种性质，以至于直到现在，时间也无法强行改变\OriginalTerm{已说过的话}{\GreekText{ὡς} \GreekText{μὴ} \GreekText{δύνασθαι} \GreekText{βιαζὲσθαι} \GreekText{χρόνῳ} \GreekText{τα} \GreekText{εἰρημένα}}：因为耶路撒冷的毁灭确实发生在许多年前；但还有其他预言从那时一直延伸到祂来临的日子；你可以随意检验：例如，这个：\BibleQuote{我天天与你们同在，直到今世的终结。} \BibleRef{玛 28:20} 和：\BibleQuote{在这磐石上，我要建立我的教会，阴间的门决不能战胜她。} \BibleRef{玛 16:18} 和：\BibleQuote{这天国的福音要传遍天下。} \BibleRef{玛 24:14} 以及那个曾是妓女的妇人所行的事；还有许多其他超过这些的。那么，这个预言的真实性从何而来，如果它真是伪造的？阴间的门如何未能战胜教会？基督如何一直与我们同在？因为如果祂没有与我们同在，教会就不会得胜。福音如何散布到全世界各地？那些反对我们的人也足以证明这些书的古老性；我指比如\Person{责尔苏}和\Place{巴塔内亚}的那位，他紧随其后。因为我想，他们不是在反对他们时代之后才写的书。\par

[7] 此外，还有整个世界一致接受了福音。若不是圣神的恩宠，从地极到地极绝不可能有这样大的共识；而伪造之事的作者也会很快被揭露。如此卓越的事也不可能从虚构和谎言中产生。难道你没有看见整个世界正在皈依，谬误被消除，古代隐修士的严峻\OriginalTerm{智慧}{\GreekText{φιλοσοφία}}比太阳更明亮，贞女的歌咏团，野蛮人中的虔诚，众人在同一轭下服役？因为这些事不仅由我们预言，也从起初由先知预言。我想，你们也不会挑剔他们的预言：因为这些书卷在敌人手中，并因某些希腊人的热忱被译成希腊文。所以，这些书卷也预言了许多关于这些事，表明将要来到我们中间的是天主。\par

[8] 那么为什么现在不是所有人都相信呢？因为情况已经退化了：这要归咎于我们。（因为从此处的话语也是对我们说的。）因为即使在当初，他们也不仅凭神迹，也凭着许多归化者的生活模式而被吸引。因为祂说：\BibleQuote{你们的光也当这样照在人前，叫他们看见你们的好行为，便将荣耀归给你们在天上的父。}\BibleRef{玛 5:16} 又说：\BibleQuote{众人都一心一意，没有人说他的财物中有什么是自己的，一切都是公有的；照各人所需用的分给各人。}\BibleRef{宗 4:32} 他们过着天使般的生活。如果现在同样的事得以实行，即使没有神迹，我们也能使整个世界皈依。但在此期间，愿那些愿意得救的人留意圣经；因为他们会在其中找到这些高尚的事迹，以及比这些更伟大的事。因为还可以补充说，教师们本身超过了其他人的事迹：他们生活在饥饿、干渴和赤身露体中。但我们却贪图极大的奢华、安逸和舒适；他们却不是这样：他们大声呼喊：\BibleQuote{直到现在，我们仍是又饥又渴，赤身露体，受人捶打，居无定所。}\BibleRef{格前 4:11} 有的从耶路撒冷跑到伊利里库姆，\BibleRef{罗 15:19} 有的到印度人地区，有的到摩尔人地区，有的到世界的这一部分，有的到另一部分。而我们却连离开自己的家乡都没有勇气；反而追求奢侈的生活、华丽的房屋和所有其他多余之物。因为我们中间有谁曾为天主的话语而忍饥挨饿？有谁曾在旷野居住？有谁曾踏上遥远的朝圣之旅？我们中间有哪位教师曾亲手做工帮助别人？有谁天天面对死亡？因此，那些与我们在一起的人也变得懒散。假如看到士兵和将领与饥饿、干渴、死亡以及一切可怕的事争斗，像狮子一样忍受寒冷、危险和一切，并因此而兴旺；然后后来放松了那种严格，变得软弱，贪爱财富，沉迷于生意和交易，然后被敌人击败，那么去寻找这一切的原因就是极其愚蠢的。现在让我们这样思考我们自己和祖先的事；因为我们也都变得比所有人更软弱，被钉在现世生活中。\par

如果有人被发现有一点古代智慧的痕迹，就离开城市和市场、脱离世俗的交往和他人的管理，投奔到深山里；若有人问起那退隐的原因，他便会提出一个不能容允的借口。因为他说：\Quote{免得我也毁灭，我的良善之刃被磨钝，我就躲开了。}那么，你变得迟钝一点而赢得他人，岂不比停留在高处而忽视你垂死的弟兄要好得多？\par

然而，当有一类人对德行漠不关心，而那些真正关注德行的人却远离我们的人群，我们又怎能制伏我们的敌人？因为即使现在有神迹发生，谁又会相信？或者那些教外的人，眼见我们如此邪僻，谁会留意我们？因为确实，在许多人看来，我们的正直生活是更可信的论据；神迹可能被顽固的恶人曲解，而纯洁的生活却足以封住魔鬼自己的口。\par

9. 这些话我向治理者和被治理者都说，而且首先向我自己说；好叫我们身上显出的生活方式真正令人钦佩，好叫我们各守其位，藐视一切现世之物；轻视财富，而不轻视地狱；忽略荣耀，而不忽略救恩；在此忍受劳苦和艰辛，以免落入那里的惩罚。让我们这样与希腊人作战；这样以胜过自由的俘虏身份俘虏他们。\par

当我们不停地反复说这些话时，它们却很少发生。然而，无论是否实行，不断地提醒你们这些是应当的。因为若有人用花言巧语迷惑人，那么那些将人唤回真理的人，更不应厌倦讲有益的话。再者：若那些迷惑者使用了如此多的计谋——他们花费金钱、运用论证、经历危险、炫耀他们的庇护——那么，我们这些将人从欺骗中赢得的人，岂不更应当忍受危险、死亡和一切，好使我们既赢得自己和他人，又在敌人面前不可抗拒，从而获得所应许的福分，藉着恩宠和慈爱，等等。\par

\SourceRecord{Translated by Talbot W. Chambers.；Nicene and Post-Nicene Fathers, First Series；Vol. 12.；Edited by Philip Schaff.；Buffalo, NY: Christian Literature Publishing Co.,；1889.；<http://www.newadvent.org/fathers/220106.htm>.}

% structure: p0001 source=h1 target=section reason=page-title

\section{格林多前书第七篇讲道}

% structure: p0002 source=h2 target=subsection reason=relative-heading-rank; base=h2

\subsection{格前 2:6-7}

\BibleQuote{然而，在成全的人中，我们讲的却是智慧，但不是这世界的智慧，也不是这世界将要终结的统治者的智慧；我们讲的乃是天主奥秘的、隐藏的智慧，这智慧是天主在万世以前预定的，为使我们获得光荣。}\par

黑暗对于眼疾患者似乎比光明更合适；因此他们宁愿躲进一间完全遮暗的房间。属灵的智慧也是如此。因为天主的智慧在外人看来似乎是愚妄；同样，他们自己的智慧，实际上是愚妄，却被他们看作智慧。结果就像这样：一个人精通航海，却承诺不用船和帆就能横渡无垠的大海，然后试图通过推理证明这是可能的；而另一个对此一无所知的人，却将自己托付给一艘船、一位舵手和水手，因而安全航行。因为后者的看似无知比前者的智慧更为明智。因为航海术是卓越的；但当它过于自夸时，就变成了一种愚妄。每一种不自量力的艺术也是如此。同样，那没有圣神之助的智慧，若是曾领受这帮助，本会是真智慧。但因它全凭自己，自以为不需要那帮助，就成了愚妄，尽管看似智慧。因此，他先用事实揭露了它，然后才称它为愚妄；并先按他们的看法称天主的智慧为愚妄，然后才证明它是智慧（因为只有在论证之后，而非之前，我们才能最好地使反对者羞愧）。\par

于是他说：\Quote{然而，在成全的人中，我们讲的却是智慧}。因为当我这个被视为愚妄、宣讲愚妄的人胜过智慧时，我不是凭愚妄，而是凭更成全的智慧得胜；这智慧如此广大、如此高超，以致另一方的智慧显得是愚妄。因此，他先前按他们当时所称的名号称呼它，并用事实证明了胜利，显明了对方的极端愚妄；此后便赐予它正确的名称，说道：\Quote{然而，在成全的人中，我们讲的却是智慧。}他称福音、救恩的方式、藉十字架得救为\Quote{智慧}。\Quote{成全的人}就是那些相信的人。因为他们的确是\Quote{成全的}，知道一切人间之事全然无助，并因深信这些事对他们毫无益处而藐视它们；这样的人就是真信徒。\par

\Quote{但不是这世界的智慧。}因为那没有圣神的智慧有何用处？它在此终止，不再前进，甚至在此也不能使拥有者受益？\par

此处所称的\Quote{世界的统治者}，按一些人的猜疑，并非指某些魔鬼，而是指那些当权者、执政者、那些视此事为角逐对象的人，即哲学家、修辞学家和\OriginalTerm{演说词作者}{\GreekText{λογογράφους}}。因为这类人是主导者，常成为民众的领袖。\par

他称他们为\Quote{世界的统治者}，因为他们的统治不超越现世。因此，他进一步补充道：\Quote{终归乌有}，既从智慧本身贬低它，也从那些运用它的人贬低它。因为他已证明它是虚假的、愚妄的、一无所能、软弱的，又进一步指出它不过是短暂的。\par

2. \Quote{我们讲的乃是天主奥秘的智慧。} 什么奥秘？因为基督肯定说：\BibleRef{玛 10:27}（公认文本\OriginalTerm{你们听}{\GreekText{ἀκούετε}}）\Quote{你们在暗中听到的，要在屋顶上宣讲出来。} 那么，他为何称它为\Quote{奥秘}？因为这事尚未成就之前，无论是天使、总领天使，还是任何受造的能力都不知道。因此他说：\BibleRef{弗 3:10}\Quote{为使天上的率领者和掌权者，藉着教会，如今得知天主的各色智慧。} 天主这样做是为了尊荣我们，使他们不能没有我们而听到奥秘。我们也一样，当我们结交朋友时，常将这事作为友谊的确实证明：即我们优先向某人透露秘密。让那些使福音的奥秘蒙羞、不加区别地向所有人展示\Quote{珍珠}和教义、并将\Quote{圣物}丢给\Quote{狗}和\Quote{猪}及无用的推理的人听吧。因为奥秘无需论证；只需宣告它本身。如果你自己再加添什么，它就不再是神圣的、各部分完整的奥秘了。\par

再从另一种意义上说，之所以称为奥秘，是因为我们所见的并非所见之物，而是我们看见一些，相信另一些。因为我们的奥迹正是如此。例如，我在这些事上与外教人的感受不同。我听说\Quote{基督被钉死了}，便立刻钦佩祂对人的慈爱；外教人听到，却视为软弱。我听说\Quote{祂成了仆人}，便惊叹祂对我们的眷顾；外教人听到，却视为羞辱。我听说\Quote{祂死了}，便惊异于祂的权能，因为祂虽在死亡中却未被拘禁，反而挣断了死亡的锁链；外教人听到，却猜测那是无能。他听到复活，便说这是传说；我深知证明复活的实事，便俯伏朝拜天主的计划。他听到圣洗，只把它当作水；但我所见的不单是可见之物，还有由圣神而来的灵魂净化。他仅考虑我的身体被洗了；但我相信灵魂也成为了纯洁而神圣的；并视之为坟墓、复活、圣化、成义、救赎、义子名分、继承、天国、圣神的\OriginalTerm{充分倾注}{\GreekText{χορηγίαν}}。因为我不是凭肉眼判断可见之事，而是凭心灵之眼。我听到\Quote{基督的身体}，我理解的是一种意义，外教人理解的却是另一种。\par

正如儿童看着书本，不知道字母的含义，也不知自己所看的是什么；甚至一个成年人若不识字，也会遭遇同样的事；但识字的人会在字母中发现许多意义，甚至是完整的人生和历史：一封书信在不识字者手中，只会被当作纸和墨；但懂得阅读的人却会听到声音，与不在场的人交谈，并通过书写随意回复：奥秘也是如此。外教人虽听，却似乎听不见；但信友拥有来自圣神的技能，能看见其中蕴含的意义。例如，这正是保禄所表示的，他说甚至现在所传讲的圣言也是隐藏的：因为\Quote{为丧亡的人}，他说，\Quote{它是隐藏的。}\BibleRef{格后 4:3}\par

从另一个角度看，这个词也指福音完全超乎一切期待。圣经惯常以此名称呼那些超出人类一切希望和思想的事。因此在另一处也有：\Quote{我的奥秘是为我}，以及为我的。保禄又说：\BibleRef{格后 15:51}\Quote{看，我告诉你们一个奥秘：我们众人不全睡眠，但众人却要改变。}\par

3. 尽管它到处被传讲，却仍是奥秘；因为正如我们受命\Quote{在屋顶上宣讲你们在耳中所听到的}，我们也受命\Quote{不要把圣物给狗，也不要把你们的珍珠丢在猪前。}\BibleRef{玛 7:9}因为有些人是属血肉的，不能明白；另一些人心上蒙着帕子，看不见；因此那奥秘尤胜一切，它到处被传讲，却不被没有正确心智的人所知；它不是由智慧，而是由圣神启示，按照我们所能接受的限度。为此，一个人若在这方面也称它为奥秘，即\OriginalTerm{不可外传的}{\GreekText{ἀπόῤῥητον}}，也并非错误。因为甚至对于我们信友，也未赐予完全的确切和精确。所以保禄也说：\BibleRef{格前 13:9}\Quote{我们如今仿佛对着镜子观看，模糊不清；到那时就要面对面了。}\par

4. 为此他说：\Quote{我们在奥秘中讲智慧，是隐藏的智慧，是天主在万世以前预定使我们得荣耀的。隐藏的：}即任何在上掌权者在我们之前都不曾知道它；如今多数人也不认识它。\par

\Quote{祂预定使我们得荣耀}，然而在别处他说\Quote{为祂自己的荣耀}，因为他视我们的得救为祂自己的荣耀；正如他也称之为祂的财富，\BibleRef{弗 3:8}尽管祂本身就富足于美善，无需任何事物才能富足。\par

\Quote{预定的}，他说，指出祂对我们的眷顾。因为人们认为那些从一开始就尽心为我们行善的人是最尊重和爱我们的：这正如父亲对待子女一样。虽然他们直到以后才赐予财产，但最初从一开始就预定了。这正是保禄现在所强调的；天主从起初就一直爱我们，甚至在我们尚未存在之时。因为除非祂爱我们，祂就不会预定我们的财富。所以不要思虑中间发生的敌意；因为这友谊比敌意更古老。\par

至于\Quote{在万世以前}（\OriginalTerm{在万世以前}{\GreekText{πρὸ} \GreekText{τῶν} \GreekText{αἰώνων}}），意思是永恒的。因为在另一处祂也这样说：\Quote{祂是在万世以前的。}圣子，如果你留意，也会被发现同样是永恒的。因为关于祂，保禄说：\BibleRef{希 1:2}\Quote{藉着祂，祂创造了诸世。}这等于说祂在万世以前已存在；因为显然制造者在被造之物之先。\par

% structure: p0018 source=h2 target=subsection reason=relative-heading-rank; base=h2

\subsection{格林多前书 2:8}

5. \Quote{这智慧，今世有权势的人没有一个认识它；因为如果他们认识了，就不至于把光荣的主钉在十字架上了。}\par

如今，他们若不知道，主又如何对他们说\BibleRef{若 7:28}：\Quote{你们认识我，也知道我从哪里来？}诚然，关于比拉多，圣经说他不知道。\BibleRef{若 19:9}很可能黑落德也不知道。这些人可称为这世界的元首；但若有人说这也是指犹太人和司祭而言，他也不会错。因为主也对他们说：\BibleRef{若 8:19}\Quote{你们不认识我，也不认识我的父。}那么，祂怎么又在稍前说：\Quote{你们认识我，也知道我从哪里来？}然而，这种认识方式和那种认识方式已在福音中论及（\WorkTitle{论圣若望福音}第49篇）；为避免不断重复同一主题，我们便将读者引向那里。\par

那么，他们有关十字架的罪被赦免了吗？诚然，祂确实说过：\BibleQuote{父啊，宽赦他们吧。}\BibleRef{路 23:34}他们若悔改，就得了赦免。因为那曾向斯德望设置无数袭击者并迫害教会的人，甚至保禄，都成了教会的捍卫者。同样，那些愿意悔改的人也得了赦免：保禄自己正表明这点，当他高呼：\BibleRef{罗 11:1}\BibleQuote{那么我说：他们失足是要跌倒吗？绝对不是。}\BibleQuote{那么我说：难道天主摈弃了祂所预知的百姓吗？绝对不。}然后，为显示他们的悔改并非被排除，他提出自己的归化作为决定性的证据，说：\BibleQuote{因为我自己也是一个以色列人。}\par

至于\Quote{他们不知道}这句话，在我看来，此处并非指基督的位格，而是仅指隐藏在那事件中的救恩计划（\OriginalTerm{关于这事本身的计划}{\GreekText{περὶ} \GreekText{αὐτῆς} \GreekText{τοῦ} \GreekText{πράγματος} \GreekText{τῆς} \GreekText{οἰκονομὶας}}）；好像他说：他们所不知道的是\Quote{死亡}和\Quote{十字架}的意义。因为在那段经文中，祂也不是说\Quote{他们不认识我}，而是说\Quote{他们不知道他们做的是什么}；也就是说，他们对其正在完成的计划和奥秘一无所知。因为他们不知道十字架将如此辉煌地闪耀；它成为世界的救恩和天主与人和好的途径；他们的城将陷落，他们将遭受极端的痛苦。\par

他用\Quote{智慧}的名称既指基督，也指十字架和福音。他也适时地称祂为\Quote{光荣的主}。因为十字架被视为耻辱之事，他表明十字架是极大的光荣；但需要极大的智慧，不仅是为了认识天主，也是为了学习天主的这个计划：而那外在的智慧不仅成为认识天主的障碍，也成为理解此计划的障碍。\par

% structure: p0024 source=h2 target=subsection reason=relative-heading-rank; base=h2

\subsection{格前 2:9-13}

6. \Quote{而是正如经上所载：\BibleQuote{天主为爱祂的人所准备的，是眼所未见，耳所未闻，人心所未想到的。}}\par

这些话在哪里记载着呢？为什么说它们\Quote{记载着}呢？因为即使在字句未曾记录而只在实际事件中，如历史书中所载；或者当同一含义被表达，但并非完全相同的字句时，也可以说\Quote{记载着}，就如这里：因为\BibleQuote{那些未曾被告知祂的人将要看见，那些未曾听见的人将要领悟}\BibleRef{依 52:15}与\BibleQuote{眼所未见，耳所未闻}是相同的。要么这就是他的意思，要么很可能这些话确实在某些书中记载着，而抄本失传了。因为许多书被毁，即使在第一次被掳时期，也只有少数被完整保存。这一点从现存于我们手中的书中可以清楚看出。因为宗徒说：\BibleRef{宗 3:24}\BibleQuote{从撒慕尔及以后的众先知，都曾预言过这些日子。}而他们的话并非全部存留。然而，保禄既通晓法律，又因圣神说话，自然会知道所有细节。我为什么提到被掳呢？甚至在被掳之前，许多书已经消失，犹太人堕落到极度的不虔敬；这从\WorkTitle{列王纪下}的结尾可以清楚看出：\BibleRef{列下 22:8}因为\WorkTitle{申命纪}几乎找不到，它被埋在某个粪堆里。\par

此外，在许多地方有双重预言，容易被智慧者理解；从这些预言中，我们可以发现许多隐晦的事情。\par

7. 那么，\BibleQuote{眼未曾看见天主所预备的}吗？没有。因为世人谁看见过那将要完成的事呢？同样，\BibleQuote{耳未曾听见，人心也未曾想到}。这是怎么回事？因为如果先知说了，他怎么说\BibleQuote{耳未曾听见，人心也未曾想到}呢？它并未进入人心；因为他不是在说自己一人，而是说全人类。那么，先知们，难道他们没有听见吗？是的，他们听见了；但先知的耳朵不是\Quote{人的}耳朵；因为他们不是作为人听见，而是作为先知。因此他说：\BibleRef{依 50:4}\BibleQuote{主赐给我一个受教的耳朵，}意指从圣神而来的\Quote{赐予}。由此可知，在听见之前，人心未曾想到。因为获得圣神的恩赐后，先知的心不再是人的心，而是属神的心；正如他自己也说：\BibleQuote{我们拥有基督的心意}\BibleRef{格前 2:16}，如同说：\Quote{在我们获得圣神的祝福并学习了那无人能言说的事之前，我们中无人，甚至先知，也未曾在心中想到它们。我们怎能想到呢？因为连天使也不知晓。\InnerQuote{这世界的元首}又有何必要说呢？}他说，\Quote{因为没有人知道它们，连天上的掌权者也不知道。}\par

这些事究竟是怎样的呢？就是借着那被视为宣讲之愚拙的，祂将战胜世界，万民将被带入，天主与人和好，如此大的福佑将降临在我们身上！我们怎么\Quote{知道？向我们}他说，\Quote{天主借着圣神将它们启示给了我们}，不是借着外来的智慧；因为这智慧像一位失宠的婢女，未被准许进入，弯腰窥探\BibleRef{若 20:5}与主有关的奥秘。你看见这智慧与那智慧之间的差别何其大吗？天使都不知道的事，这些正是她教导我们的；但那外来的，却做了相反的事。她不仅未能教导，反而阻碍和拦阻，并在事后试图遮掩祂的作为，使十字架归于无效。那么，他描述赐予我们的恩荣，并非仅仅因我们接受了知识，也非因我们与天使一同接受，而是——更重要的是——因祂的圣神将之传达给我们。\par

7. 然后，为了显示其伟大，他说：如果知道天主奥秘的圣神没有启示它们，我们就不会学到它们。这个主题对天主而言是如此关注，以至于成了祂的秘密。因此，我们也需要那位完全知道这些事的导师；因为他说：\BibleQuote{圣神洞察一切，甚至天主的深奥。}\BibleRef{格前 2:10}因为\Quote{洞察}一词在这里不是指无知，而是指准确的知识：这正是他论及天主时所用的同一说法，他说：\BibleQuote{洞察人心的那位知道圣神的意愿是什么。}\BibleRef{罗 8:27}然后，在精确讲述了圣神的知识，并指出它完全等同于天主的知识，如同人的知识与其自身相等；并指出我们已从它学到一切，且必须从它学到之后，他补充说：\Quote{这些事我们也讲论，不是用人的智慧所教的言语，而是用圣神所教的言语，将属灵的事与属灵的事相比较。}你看见他因导师的尊贵将我们高举到何等程度吗？因为我们比他们更智慧，其差别有如\Person{柏拉图}和圣神之间的差别；他们以异教修辞学家为师，而我们的导师是圣神。\par

8. 但\Quote{将属灵的事与属灵的事相比较}是什么意思？当一件事是属灵的且含义模糊时，我们从属灵的事中引出证据。例如，我说：基督复活了——由童贞女所生；我引出证据、预像和论证：\Person{约纳}在鲸鱼腹中及其后来的获释；不生育者\Person{撒辣}、\Person{黎贝加}等生子；乐园中树木的生长\BibleRef{创 2:5}，那时没有撒种，没有降雨，没有犁沟。因为将来的事由先前的事塑造和预表出来，如同影子，为的是使现在这些事在它们来临时得以被相信。我们又展示，人如何出自尘土，女人如何单独来自男人，且没有任何交合；大地本身如何出自无有，大工匠的力量处处足以成全一切。因此，我用\Quote{属灵的事}来\Quote{比较属灵的事}，我从不需外来的智慧——既不需它的推理，也不需它的修饰。因为这种人只会搅扰软弱的理解力，使其迷惑；不能清楚证明他们所肯定的任何一件事，反而产生相反的效果。他们更扰乱心思，使其充满黑暗和许多困惑。因此他说：\Quote{将属灵的事与属灵的事相比较。}你看见他显示这是多么多余吗？不仅多余，甚至敌对和有害：因为\Quote{免得基督的十字架归于无效}和\Quote{使我们的信德（你们的信德，通用文本）不在于人的智慧}等说法正是此意。他指出，对于那些将一切自信托付于它的人，要学到任何有益的事是不可能的；因为\par

% structure: p0032 source=h2 target=subsection reason=relative-heading-rank; base=h2

\subsection{\BibleRef{格前 2:14-16}}

9. \Quote{属血气的人不接受圣神的事。}\par

因此必须先把它搁置一旁。有人会说：\Quote{那么，外来的智慧就受贬抑吗？但它却是天主的化工。}这怎能清楚？因为天主并未创造它，它是你的发明。因为在此处他用\Quote{智慧}一词指好奇的探究和多余的辞藻。但若有人说他指的是人的理智，即使在此意义上，错误也在你。因为你滥用它，给它带来恶名；你为了伤害和阻挠天主，向它要求它从未有过的事物。既然你以此自夸并同天主争斗，祂就揭露了它的软弱。因为身体的力气也是好东西，但当\Person{加音}没有善用它时，天主就使他软弱、颤抖\BibleRef{创 4:12}。酒也是好东西，但因犹太人过度沉溺其中，天主便完全禁止司祭使用这果实。既然你也滥用智慧来拒绝天主，并向它要求超过其自身能力的事，为了将你从人的希望中拉出来，祂已向你显示了它的软弱。\par

因为（继续）他是\Quote{一个属血气的人}，凡事归因于心思的推理，不认为自己需要来自上天的帮助；这简直是愚昧的标志。因为天主赋予它（理智）是为了让它学习并接受祂的帮助，而不是让它认为自己自给自足。因为眼睛美丽有用，但若它们选择在没有光的情况下观看，它们的美丽对它们毫无益处，连同它们天生的能力也没有，甚至产生损害。所以，你若留意，任何灵魂，若选择在没有圣神的情况下观看，甚至会成为自己的障碍。\par

\Quote{那么，在此之前，}有人会说，\Quote{她怎会靠自己看见一切？}她从未靠自己这样做，但她有受造物作为一本放在她面前敞开的书。但当人们不再遵行天主吩咐的道路，不再藉可见之物的美丽认识伟大的工匠，而将知识的权杖托付给辩论时，他们便变得软弱，沉没在无神的海洋中；因为他们很快就引出了万恶的深渊，断言无物从无中产生，而是从未受造的质料中产生；由此，他们成为了无数异端的根源。\par

此外，在他们极端荒谬的观点上，他们是一致的；但在那些他们似乎隐约梦见某些有益之事（尽管只是如同在影子里）的事上，他们彼此争执，好让双方都被嘲笑。因为从无中产生无物，几乎所有的人都一致主张并书写；并且是带着极大的热忱。在这些荒谬之事上，他们受到魔鬼的驱使。但在他们有益的言论中，他们似乎（虽然是\OriginalTerm{隐晦地}{\GreekText{ἐν} \GreekText{αἰνίγματι}}）找到了他们所寻求的一部分，在这些事上他们彼此争战：例如，灵魂是不朽的；德行无需外在之物；行善或行恶并非出于必然或命运。\par

你看到魔鬼的诡计了吗？如果他在任何地方看到人们说败坏的话，他就使所有人意见一致；但如果任何地方看到人们说正确的话，他就兴起别人反对他们；这样，荒谬之论因普遍同意而得巩固，有益之事则因被不同理解而消亡。请注意，灵魂在各方面都是\OriginalTerm{松弛的}{\GreekText{ἄτονος}}，不足以自足。这是可以预料的。因为，如果灵魂在这样状态下尚且渴望一无所缺并远离天主；设想她若没有陷入这种状态，她会不知不觉地沉入何等的极端疯狂中？如果她带着必死的身体，却因魔鬼的虚假应许——\BibleRef{创 3:4}——而期待更大的事，那么她若一开始就接受了不朽的身体，又会将自己抛弃到何种地步？因为，甚至在那之后，她还通过摩尼教派的腐败之口宣称自己是未受生的、属于天主的本体，正是这种混乱给了她发明希腊众神的机会。因此，在我看来，天主使德行变得艰辛，目的是压制灵魂并使之归于节制。为了使你信服这是真实的（如同从微小之事可以猜测伟大之事），让我们从以色列人身上学习。众所周知，当他们不过辛劳的生活而耽于安逸时，不能承受繁荣，便陷入了无神。那么天主对此做了什么？他将许多法律加在他们身上，以约束他们的放纵。为了使你信服这些法律无助于任何德行，而是作为某种缰绳赐给他们，为他们提供不断辛劳的机会，请听先知关于他们所说的：\Quote{我赐给他们不良的规律。} \BibleRef{则 20:25}。\Quote{不良的}是什么意思？是那些不大有助于德行的。因此他又加上：\Quote{并使他们不能生活的法令。}\par

10. \Quote{然而属血气的人不接受圣神的事。}\par

因为正如用这双眼睛无人能学习天上的事一样，灵魂若没有帮助也无法学习圣神的事。我为何谈论天上的事呢？它甚至不能领会地上的一切事。例如，从远处看一座方形的塔，我们以为它是圆的；这样的看法不过是眼睛的错觉。同样，我们可以确信，当人单凭自己的理智审视遥远的事物时，会产生许多可笑的结果。因为他不仅不能看见事物的真相，甚至会把它们当作相反的东西。因此他又加上：\Quote{因为在他看来是愚拙的。}但这并非出于事物的本性，而是出于他的软弱，他无法通过灵魂的眼睛达到它们的伟大。\par

11. 接着，他继续对比，说明其原因，说：\Quote{他不能知晓，因为那是凭圣神审断的：}即所主张的事需要信德，而凭推理无法领会，因为它们的伟大远远超过我们理解的渺小。因此他说：\Quote{惟独属神的人审断一切，但他自己却不受任何人审断。}因为那有视力的人，自己看见一切与无视力者有关的事；但无视力的人不能看见另一人在做什么。同样，在我们眼前的情形中，我们自己的事和不信者的事，我们全都知道；但我们的这些事，他们从今以后再也无法知晓。我们知道现世之物的本质，未来之事的尊严；当这世态不再有时，世界会成为什么样子，罪人要遭受什么，义人要享受什么。我们还知道现世之物毫无价值，并揭露它们的卑贱；（因为\Quote{审断}也是揭露；\OriginalTerm{审断、揭露}{\GreekText{ἀνακρίνειν}, \GreekText{ἐλέγχειν}}）并且知道未来之事是不朽且不变的。这一切事，属神的人都知道；属血气的人离世时将要遭受什么，以及信者完成今世旅程后将享受什么：这些事，属血气的人一样也不知道。\par

12. 因此，他在明确证明了所断言之事后，紧接着说：\Quote{谁能知道主的心，谁会教导祂？但我们拥有基督的心意。}也就是说，我们知道基督心意中的事，正是祂所愿意并启示的。因为他曾说：\Quote{圣神启示了这些事}；为避免有人轻视子，他接着指出基督也向我们显示了这些事。这并不是说，祂所知道的一切我们都知道；而是说，我们所知道的一切并非属人的、容易引起怀疑的，而是属于祂的心意和属神的。\par

因为关于这些事，我们所拥有的心意来自基督；也就是说，我们关于信仰之事的认知是属神的；因此我们有理由\Quote{不受任何人判断。}因为属本性的人不可能知道属神的事。所以他也说：\Quote{谁能知道主的心？}暗示我们自己关于这些事的心意就是祂的心意。而他加上\Quote{谁会教导祂}也非没有理由，而是针对他刚说过的\Quote{属神的人不被任何人识透。}因为若无人能知道主的心，就更不能教导和修正它。因为\Quote{谁会教导祂}正是此意。\par

你看，他从各方面驱散外来的智慧，并表明属神的人知道更多、更大的事。因为他看到这些理由：\Quote{使一切血肉之躯无法夸耀}；\Quote{为此祂拣选了愚拙的，为使智慧人蒙羞}；以及\Quote{免得基督的十字架落了空}，在外人看来似乎不太值得信赖，也缺乏吸引力、必要性或益处，于是他最终提出了主要理由：因为如此我们最容易看出，我们可以从谁那里获得学习高超、隐密和超越我们之事的方法。因为由于我们无法通过外邦人的智慧理解超越我们的事，理性完全失效了。\par

你也可以注意到，以这种方式从圣神学习更为有益。因为那是最容易、最清晰的教学。\par

\Quote{但我们拥有基督的心意。}也就是说，属神的、神圣的、毫无属人成分的心意。因为它不是柏拉图或毕达哥拉斯的，而是基督亲自将祂自己的事放在我们心中。\par

那么，亲爱的，即使没有别的，让我们敬畏这一点，并使我们的生命发出最卓越的光辉；因为祂自己也把这作为深厚友谊的确证，即将祂的秘密启示给我们：正如祂所说：\BibleRef{若 15:15} \Quote{以后我不再称你们为仆人，因为你们都是我的朋友；凡我从父所听见的，我都已经告诉了你们。}也就是说，我信任你们。如果信任本身是友谊的证明，那么当祂不仅将\OriginalTerm{用言语表达的奥秘}{\GreekText{τὰ} \GreekText{διὰ} \GreekText{ῥημάτων} \GreekText{μυστηρία}}托付给我们，还将以行动表达的奥秘（\OriginalTerm{以行动表达的奥秘}{\GreekText{διὰ} \GreekText{τῶν} \GreekText{ἔργων}}，即圣事行为）也传授给我们，想想这份爱是多么硕大。即使没有别的，让我们敬畏这一点；即使我们不那么在意地狱，然而对这样一位朋友和恩人忘恩负义，应当比地狱更加可怕。让我们不要像雇工，而是像儿子和自由人，为了对父的爱做一切事；让我们最终停止依恋世界，好使外邦人也感到羞愧。因为现在，即使想对他们使出全力，我也犹豫不决，唯恐尽管我们以论证和我们所教导的真理胜过他们，却因我们生活方式的比较而招致大量嘲笑；因为他们虽然执着于错误且没有这样的确信，却坚守哲学，而我们却恰恰相反。然而，我还是要说。因为在与他们较量时，我们也许、也许会渴望在生活方式上也变得比他们更好。\par

14. 我不久前说过，宗徒们若没有领受天主的恩宠，根本就不会想到宣讲他们所宣讲的；他们不仅不会成功，甚至不会设计出这样的事。那么，让我们今天也在讲论中继续同一主题；让我们表明，若基督没有与他们同在，这事不可能被他们选择或想到：不是因为他们是软弱的对阵强壮的，不是因为少数对阵多数，不是因为贫穷对阵富有，不是因为无学问对阵有智慧，而是因为他们的偏见的力量也是强大的。因为你知道，没有什么比古老习俗的专制对人更有力量。因此，即使他们不是仅有十二人，也不是如此卑微，不是如此本色，而是有另一个像这个世界一样大的世界，并有同等数量的人站在他们一边，甚至更多，这件事也仍然是难以实现的。因为另一方有习俗支持，而对他们来说，他们的新奇性是一个障碍。因为没有比革新和引入新奇之事更能搅动人心的了，即使是为了某种有益的目的；尤其是在涉及敬拜天主和天主的荣耀之事上。这种情况有多大的力量，我现在要说明；首先提出以下声明，即关于犹太人还有另一重困难。因为对于希腊人，他们彻底摧毁了他们的神祇和教义；但对犹太人，他们并非如此辩论，而是废除了许多教义，同时却要求他们敬拜颁布同一法律的天主。并且声称人应尊重立法者，他们说：\Quote{不要在所有方面服从祂的法律}；例如，守安息日、行割礼、献祭或做任何类似的事。因此，不仅习俗是障碍，而且还有这样一个事实：当他们命令人敬拜天主时，他们命令人违背祂的许多法律。\par

15. 然而在希腊人中，习俗的暴政极为强大。假若这习俗仅存续十年，我说的是并非如此漫长的岁月，并且它只占据了少数人，我说的是并非整个世界，当他们这些人前来接触时，即使在这种情况下，要带来变革也是困难的。但如今，诡辩家、演说家、父亲、祖父，以及比所有这些都更古老的许多人，早已被谬误所占据；大地、海洋、山岳、丛林，所有蛮族国家，所有希腊部落，智者和愚者，统治者与臣民，妇女与男子、老少、主人与奴隶、工匠与农夫、城居者与乡居者——全都如此。那些受教导的人自然会问：\Quote{这究竟是怎么回事？难道世上所有的人都受骗了吗？诡辩家、演说家、哲学家、历史学家、当代人以及更早的人、毕达哥拉斯学派、柏拉图学派、将军、执政官、国王，从起初在各城邦作公民和殖民者的人，无论蛮族还是希腊人，全都受骗了吗？而十二个渔夫、帐幕匠和税吏却比所有这些人都更有智慧？谁又能忍受这样的说法呢？}然而，他们并没有这样说，心中也没有这样的念头，反而容忍了他们，承认他们比所有人都更有智慧。正因如此，他们甚至胜过了一切。习俗虽因其历经岁月而全盛时被视为不可战胜，但在这事上并未成为阻碍。\par

为使你明白习俗的力量何等强大，它常常胜过了天主的诫命。何以我说诫命？甚至胜过极大的恩赐。例如，犹太人有玛纳时，却要求葱蒜；他们享有自由，却惦记着奴役；他们不断向往埃及，因为习惯了那里。习俗真是专横之物。\par

如果你也愿意听异教徒的例子：据说柏拉图虽然清楚知道关于众神灵的一切不过是一种欺骗，却顺从了所有节庆及其他事项，因为他无力与习俗抗争；实际上这是从他的老师那里学到的。因为他的老师也涉嫌类似的革新，结果非但未能如愿，反而丧了性命，且是在为自己辩护之后。如今我们看见多少人因偏见而固守偶像崇拜，在被指责为希腊人时，却说不出任何有说服力的话，只举出父亲、祖父和曾祖父。异教徒中有些人称习俗为第二天性，无他原因。但当教义成为习俗的内容时，它便扎根更深。因为人改变其他一切，远比改变与宗教有关的事容易。羞耻感加上习俗，足以造成障碍；而在极度年迈时，似乎要向并不那么聪明的人学习新道理，也是障碍。既然连身体上习俗都有巨大力量，灵魂上发生这种事又何足为奇呢？\par

16. 然而，在宗徒们的情形中，还有一个比这些更强大的障碍；他们不仅是在改变如此古老原始的习俗，而且这改变还伴随着危险。因为他们并非简单地将人从一种习俗引向另一种习俗，而是从一种没有恐惧的习俗，引向一种带有危险威胁的事业。因为信徒必须立即面临没收财产、迫害、被逐出祖国；必须遭受最坏的恶事，被众人憎恨，成为自己人和陌生人的共同敌人。所以，即使他们是邀请人从一种习俗转向新奇之事，这已是困难之事。但当是从习俗转向革新，且还附加所有这些恐怖时，请思考这障碍是何等巨大！\par

再者，另一件事，其重要不亚于前述者，更增加了改变的难度。因为除了习惯和危险之外，这些诫命更为沉重，而他们劝人脱离的那些事却是容易而轻松的。因为他们所召唤的是从淫乱到贞洁；从贪生怕死到各种死亡；从醉酒到斋戒；从欢笑到眼泪和痛悔；从贪婪到彻底贫乏；从安全到危险：并且在一切事上要求极其严格的谨慎。因为，\BibleQuote{淫秽} \BibleRef{弗 5:4} 他说，\BibleQuote{以及愚妄的谈吐和戏笑的话，总不可出口。}这些话是对那些只知醉酒和服侍肚腹的人说的；他们举行的盛宴不过是\Quote{淫秽}、欢笑和种种喧闹嬉戏（\OriginalTerm{各种戏谑}{\GreekText{κωμῳδίας} \GreekText{ἁπάσης}}）。因此，不仅那些关于生活严谨的教义本身是沉重的，而且因为它们是对那些在无忧无虑、\InnerQuote{淫秽}、\InnerQuote{愚妄的谈吐}、欢笑和狂欢中长大的人说的。因为这些沉溺于其中的人，当听到\BibleRef{玛 10:38}\BibleQuote{谁不背起自己的十字架跟随我，不配是我的；}和\BibleRef{玛 10:34}\BibleQuote{我来不是为送和平，而是带刀剑，叫人与父亲生分，女儿与母亲生分，}谁不会感到\OriginalTerm{浑身僵冷}{\GreekText{ἐνάρκησε}}呢？当听到\BibleQuote{谁不告别家园、田产和财物，不配是我的}，谁不会犹豫、拒绝呢？然而，竟有人不但不感到寒冷，也不退缩，反而迎上前去，冲入艰难之中，热切地接受所命令的诫命。再者，当被告知：\BibleQuote{凡所说的闲话，在审判之日，都要交账；}\BibleRef{玛 12:36} 以及 \BibleQuote{凡注视妇女，有意贪恋她的，就已经在心里与她犯了奸淫；}\BibleRef{玛 5:28} 以及 \Quote{凡无缘无故向弟兄发怒的，就要受地狱的罚；}——这些事哪一样不会吓跑当时的人呢？然而众人都涌了进来，许多人甚至越过了赛程的界限。那么，是什么吸引了他们？难道不是那被宣讲者的能力吗？假使情况并非如此，而是恰恰相反——这边是那边，那边是这边——请问，能否轻易抓住并拖住那些抗拒的人？我们不能这样说。因此，无论如何，那如此卓越运作的能力被证明是神圣的。否则，请告诉我，他们如何能说服那些轻浮放荡之人，促使他们走向那严峻而崎岖的生活道路？\par

17. 好了，诫命的本性就是如此。但让我们看看道理是否吸引人。不，在这方面也足以吓跑不信的人。因为传道者说什么？我们必须敬拜那被钉十字架的，并视祂为天主，祂出自一位犹太女子。现在，有谁会受这些话的劝说呢，除非有神圣的能力在引领？祂确实被钉十字架并被埋葬，这是所有人都知道的；但祂复活并升天，除了宗徒们，没有人看见。\par

但你会说，他们用许诺来激动他们，用言语的空响欺骗了他们。不，正是这个主题最特别地表明（即使撇开所有已说的）我们的道理并非欺骗。因为所有的艰苦都发生在今世，而安慰则要在复活后应许。因此，我重复，这一点表明我们的福音是神圣的。因为为什么没有一位信徒这样说：\Quote{我不接受这个，也不忍受它？你以今世的艰苦威胁我，却把好事应许在复活之后。怎么，如何能知道会有复活？已离世的人中谁曾回来？安息的人中谁曾复活？他们中谁曾说过我们离世后的景况？}但这一切都没有进入他们的心中；相反，他们为被钉十字架者献出了自己的生命。因此，这个赤裸裸的事实比什么都更能证明那伟大的能力：第一，他们立时使那些从未听说过此类事情的人，对如此重要的事产生确信；第二，他们说服他们以考验来接受艰难，并将福乐视为希望之事。现在，如果他们是欺骗者，他们会做相反的事：他们会把好事应许为\OriginalTerm{今世}{\GreekText{ἐντεῦθεν}}，正如\BibleRef{若 18:36}；而可怕的事他们则不会提及，无论是关乎今生还是来世。因为欺骗者和谄媚者就这样行事。他们不提出任何严厉、刺激或沉重的事，而完全相反。因为这就是欺骗的本性。\par

18. 但\Quote{愚昧}，有人会说，\Quote{使大部分人相信了他们被告知的事。}你怎么说？当他们处于希腊人中时，他们并不愚昧；但当他们来到我们这里时，他们的愚昧才开始吗？然而，他们并非来自另一种族或另一个世界，宗徒们才得以说服他们：他们也不过是持有希腊人意见的人，但他们在危险中接受了我们的意见。因此，如果他们更有理由坚持前者，就不会背离，何况他们已经受了那么长时间的教育，尤其是背离并非没有危险。但是，当他们从事情的本质认识到那一边全是嘲笑和欺骗时，于是即使在各种死亡的威胁下，他们也从他们的习惯中\OriginalTerm{跳脱}{\GreekText{ἀπεπήδησαν}}，自愿来到这新的道理；因为后者符合本性，而前者违反本性。\par

但\Quote{被说服的人}，据说，\Quote{是奴隶、妇女、乳母、产婆和阉人}。首先，我们的教会并不仅由这些人组成，这是尽人皆知的。但即使如此，这恰恰使福音值得赞叹：因为像柏拉图和其追随者都无法领悟的道理，渔夫们却能突然有能力说服最无知的众人接受。因为如果他们只说服了明智的人，结果就不会如此奇妙；但使奴隶、乳母和阉人达到如此严格的生活方式，以致成为天使的竞争者，这提供了他们受神感最有力的证明。再者，如果他们命令的是一些微不足道的小事，那么提出这些人被说服的情况来显示所说之事的琐碎也许是合理的；但如果他们智慧课程的内容是伟大、崇高、几乎超越人性、需要崇高思想的，那么你显示被说服者的愚昧越多，你就越清楚地表明说服者是有智慧、充满天主恩宠的。\par

但你或许会说，他们借着应许的极其伟大而说服了他们。但请告诉我，难道这不正是令你惊奇的事吗？他们如何说服人们盼望死后的奖赏和报酬？因为这一点，即使没有别的，也足以令我惊叹。但这也，有人会说，是出于愚昧。请告诉我，这些事哪里愚昧了：灵魂是不朽的；一个公正的法庭将在今世之后接纳我们；我们要向那知晓一切隐秘的天主交账，连我们的行为、言语和思想；我们将看见恶人受罚，好人头戴冠冕。不，这些事不是出于愚昧，而是智慧的最高教导。愚昧在于与这些相反的意见。\par

19. 如果只有这一件事，即轻视眼前之物、重视德行、不在此世寻求赏报、而在希望中远远超越、并保持灵魂如此专注和信实，以致不受任何现世的恐惧阻碍对将来之事的盼望；请告诉我，这属于何等崇高的哲学？但你是否也愿意学习应许和预言本身的力量，以及那些在现存状态之前和之后所发出的预言的真理？看，我向你展示一条金链，从起初就巧妙地编织而成！祂告诉他们一些关于祂自己、关于教会、以及关于将来的事；祂说话时，行了大能的事。因此，借着祂所说之事的应验，显然所行的奇事是真实的，将来的和应许的事也是真实的。\par

但为使我的意思更清楚，让我用实际事例来说明。祂只用一句话就复活了拉匝禄，使他活过来。此外，祂说：\BibleQuote{阴间的门不能胜过教会}\BibleRef{玛 16:18}，并且\BibleQuote{凡舍弃父亲或母亲的，必在今世领受百倍，并承受永生。}\BibleRef{玛 19:29}奇迹只有一个，就是复活拉匝禄；但预言有两个：一个在此世实现，另一个在来世实现。现在看它们如何互相证明。因为如果有人不相信拉匝禄复活，从关于教会发出的预言，让他学习相信奇迹。因为那在多年前所说的话，那时应验了，得到了实现：因为\BibleQuote{阴间的门没有胜过教会。}你看见那在预言中说真理的，显然也行过奇迹；而那既行奇迹又使祂所说的话实现的，显然也在关于未来之事的预言中说真理，当祂说：\Quote{轻视眼前之物的人，必在今世领受百倍，并承受永生。}因为已经成就和说过的事，祂已作为最确实的保证，给予将来要发生的事。\par

因此，从福音书中汇集所有这些事以及类似的事，让我们告诉它们，从而堵住他们的口。但如果有人说：为什么错误没有被完全消灭？我们可以这样回答：你们自己应归咎，你们反抗自己的救恩。因为天主如此\OriginalTerm{安排了}{\GreekText{ᾠκονόμησεν}}这事，以致连一点旧日不敬神的残余都不必留下。\par

20. 现在，简略重述一下所说的：事物的自然趋势是什么？是弱者被强者克服，还是相反？是说易事的人，还是说更严厉之事的人？是以危险吸引人的人，还是以安全吸引人的人？是创新者，还是巩固习惯的人？是引人进入崎岖之路，还是平坦之路的人？是使人脱离祖先传统的人，还是不立奇怪新律的人？是应许所有好事在离世之后的人，还是在此生奉承的人？是少数胜过多数，还是多数胜过少数？\par

但是，有人會說：\Quote{可是你也曾對今世之事許下承諾。}那麼，我們在今生許諾了什麼？是罪的赦免與重生的聖洗聖事。首先，聖洗聖事本身主要關乎未來之事；保祿高聲說：\BibleQuote{因為你們已經死了，你們的生命與基督一同藏在天主內；當你們的生命顯現時，你們也要與祂一同在光榮中顯現。}\BibleRef{哥 3:4}然而，即使今世也有其益處——確實如此——這更是萬分令人驚奇的事實：那些做了無數惡事、甚至無人做過的事的人，竟能說服他們徹底洗淨一切，並且不必為任何過犯交賬。因此，最令人驚奇的是，他們說服了蠻族接受這樣的信仰，對未來之事懷有希望；他們拋棄了以往的罪擔，以最大的熱忱投身於未來德行所要求的勞苦，不再貪戀任何感官之物，而是提升到超越一切有形之物的高度，領受純粹屬神的恩賜：是的，波斯人、薩爾馬提亞人、摩爾人、印度人都認識了靈魂的淨化、天主的大能、祂對人類不可言喻的仁慈、信仰的嚴格紀律、聖神的降臨、身體的復活以及永生的道理。因為在所有這些事以及超越這些的事上，漁夫們藉著聖洗感化各類蠻族，說服他們\OriginalTerm{按高尚原則生活}{\GreekText{φιλοσοφεῖν}}。\par

因此，我們要精確觀察這一切，向外邦人講論，並再次向他們展示我們生活的證據：好使我們自己得救，也使他們因我們而被引領到天主的光榮。因為光榮歸於祂，直到永遠。阿們。\par

\SourceRecord{Translated by Talbot W. Chambers.；Nicene and Post-Nicene Fathers, First Series；Vol. 12.；Edited by Philip Schaff.；Buffalo, NY: Christian Literature Publishing Co.,；1889.；<http://www.newadvent.org/fathers/220107.htm>.}

% structure: p0001 source=h1 target=section reason=page-title

\section{格林多前书第八篇讲道}

% structure: p0002 source=h2 target=subsection reason=relative-heading-rank; base=h2

\subsection{格林多前书 3:1-3}

\BibleQuote{弟兄们，我从前对你们说话，不能把你们当作属神的，只能当作属血肉的，当作在基督里的婴孩。我给你们喝奶，没有给干粮；因为那时你们不能吃，就是如今还是不能；因为你们仍是属血肉的。}\par

在推翻了外在的哲学并挫败了其一切傲慢之后，他转向另一个论点。因为他们很可能会说：\Quote{如果我们提出柏拉图、毕达哥拉斯或任何其他哲学家的意见，你倒有理由这样长篇大论地指责我们。但如果我们传讲属神的事，你为何要\OriginalTerm{翻来覆去}{\GreekText{ἄνω} \GreekText{καὶ} \GreekText{κάτω} \GreekText{στρέφεις}}地搬弄外在的智慧呢？}\par

请听他是如何反驳这一点的。\BibleQuote{弟兄们，我从前对你们说话，不能把你们当作属神的。}为什么？首先，他说，即使你们在属神的事上也已完美，也绝不该自高自大；因为你们所宣讲的不是你们自己的，也不是你们凭自己发现的。但现在，连这些事你们也不晓得，正如你们应当晓得的，你们还是学生，是最末后的。因此，若外邦人的智慧是你们高傲的缘由，那智慧已被证明为虚无，而且就属神的事而言，甚至是对我们有害的；或者若是由于属神的事，在这些事上你们也欠缺，居末位。因此他说：\BibleQuote{我不能把你们当作属神的。}他没有说\BibleQuote{我没有说}，以免显得他有所保留；而是通过两种方式压制他们的高傲：第一，因为他们不知道完美的事；第二，因为他们的无知是由于自己；还有，第三种方式，他指出\BibleQuote{就是如今他们还是不能}。因为起初他们的无能或许是由于情况使然。但事实上，他连这个借口也没给他们留下。因为他说他们没能领受，并非由于他们无法接受高深的道理，而是因为他们\BibleQuote{属血肉}。然而，开始时这并不那么可责；但过了这么久，他们仍未达到更完美的知识，这是最迟钝的征兆。\par

可以发现，他对希伯来人也提出了同样的指控，但语气没那么强烈。他说，那些人如此，部分是因为患难；而这些人，则是因为某种对邪恶的贪恋。这两者并不相同。他也暗示，当他讲这些真理之言时，在一个场合他意在责备，在另一个场合则是激励。因为他对这些格林多人说：\BibleQuote{就是如今你们还是不能}；但对另一群人，\BibleRef{希 6:1}\BibleQuote{所以，我们应当离开基督初步的道理，努力进到完全的地步}；又，\BibleRef{希 5:9}\BibleQuote{我们深信你们有更好的事，就是关乎救恩的事，虽然我们这样说。}\par

他如何称那些获得了如此大量圣神的人为\BibleQuote{属血肉的}，并曾在开头大大赞美他们？因为主对他们也说：\BibleRef{玛 7:22}\BibleQuote{离开我吧，你们这些作恶的人，我不认识你们；}然而他们也曾驱魔、复活死人、发预言。所以，即使行奇迹的人也可能是属血肉的。因为天主藉巴郎行事，向法郎启示未来之事，向拿步高启示；盖法预言，却不知自己所说；还有，有些人奉主名驱魔，虽\BibleRef{路 9:49}\BibleQuote{不与主同行}；因为这些事不是为了行者的缘故而做，而是为了他人；并且，那些完全不配的人被用作工具，也并非罕见。那么，如果这些事为了他人的缘故而藉不配之人成就，又何必惊奇呢？因为即使由圣人施行时也是如此。因为保禄说：\BibleRef{格前 3:22}\BibleQuote{一切都是你们的；无论是保禄、阿颇罗、克法，或是生命、死亡；}又，\BibleRef{弗 4:11}\BibleQuote{祂赐予一些作宗徒，一些作先知，一些作牧者及教师，为成全圣徒，从事服事的工作。}因为若不如此，就无法防止普遍的败坏。因为很可能掌权者是邪恶污秽的，而他们的子民是善良正直的；平信徒度敬虔的生活，而司铎却行邪恶；若每次恩宠都要求功德，就不可能通过这样的人施行圣洗、基督的身体或祭献。但事实上，天主甚至藉不配之人行事，圣洗的恩宠丝毫不因司铎的品行受损；否则领受者就会受损。因此，虽然这类事很少发生，但必须承认，确实发生。我说这些话，免得任何旁观者因关注司铎的生活而对\OriginalTerm{所举行的圣事}{\GreekText{τὰ} \GreekText{τελούμενα}}感到反感。\Quote{因为人不能把任何东西引入我们面前所摆列的事中，而是全凭天主大能的作为，祂正是\OriginalTerm{引导你进入奥迹者}{\GreekText{ὁ} \GreekText{μυσταγωγῶν}}。}\par

\BibleQuote{弟兄们，我从前对你们说话，不能把你们当作属神的，只能当作属血肉的，当作在基督里的婴孩。我给你们喝奶，没有给干粮。因为你们那时不能担当，就是如今还是不能。}\par

因为怕他显得是\OriginalTerm{为了夸耀}{\GreekText{φιλοτιμίας} \GreekText{ἕνεκα}}而说了刚才那些话（即：\BibleQuote{属神的人看透万事}和\BibleQuote{他自己不受任何人判断}和\BibleQuote{我们有基督的心}），同时也为了抑制他们的骄傲：请注意他说什么。他说：\Quote{并不是因为这个缘故}他说，\Quote{我保持沉默，不是因为我不能告诉你们更多，而是因为\BibleQuote{你们是属血肉的：就是如今还是不能。}}\par

为什么他不说：\Quote{你们不愿意}，却说：\Quote{你们不能够}？这正是因为他以后者代替前者。因为能力的缺乏源于意愿的缺乏。这虽然对他们是指责，但对他们导师而言却是借口。因为如果他们天生无能为力，或许还可以被原谅；但既然是因为选择，他们就完全失去了借口。接着他谈到使他们成为属血肉的具体原因：\BibleQuote{因为在你们中间有嫉妒、纷争和分裂，你们岂不是属血肉、照着人的样子行事吗？}虽然他也谈到他们中间有淫乱和不洁，但他还是先列出这个他一直努力纠正的过错。如果\Quote{嫉妒}使人成为属血肉的，那么我们就该痛哭流涕，披上麻衣，躺在灰烬中了。因为谁没有沾染这种情欲呢？除非我只是从自己推测别人的情况。如果\Quote{嫉妒}使人成为\Quote{属血肉的}，不容许他们成为\Quote{属神的}，即使他们能说预言并显其他奇事；那么，当我们连这样的恩宠也没有时，我们自己的行为将在哪里立足呢？我们不仅在这一点上，而且在其他更重要的事上也都被定罪了。\par

4. 由此我们得知，基督有理由说：\BibleRef{若 3:20}\BibleQuote{那作恶的，不来就光}；不洁的生活是接受高深道理的一种阻碍，不让悟性的明辨力显现出来。正如一个错误的人，如果他正直地生活，就不可能留在错误中；同样，一个在罪恶中长大的人，很难迅速仰望所传给我们的道理的高深之处；他必须摆脱一切情欲，才能追求真理：因为谁摆脱了这些，也将摆脱他的错误而获得真理。因为，我请求你们，不要以为仅仅戒除贪婪或奸淫就足以达到这个目的。不是这样。凡寻求真理的人，必须在他身上一切美德都协调一致。因此\Person{伯多禄}说：\BibleRef{宗 10:34}\BibleQuote{我真明白了：天主是不看情面的；但在每个民族中，凡敬畏祂、履行正义的，都为祂所悦纳。}这就是说，祂召唤并吸引他归向真理。你们没有看到\Person{保禄}吗？他比任何人在争战和迫害时都更激烈。然而，因为他过着无可指责的生活，而且做这些事并非出于人的激情，所以他被接纳，并且达到了超乎一切的目标。但如果有人说：\Quote{某个希腊人，他友善、善良、有人情味，怎么会继续留在错误中呢？}我的回答是：他有某种别的私欲，如虚荣心、思想怠惰、或对自己的救恩不关心，以为所有关于他的一切都是随意和偶然地漂流。\Person{伯多禄}称一个在一切事情上无可指摘的人为\Quote{履行正义的}，[保禄说：]\Quote{按法律中的正义，是无可指责的。}又说：\BibleQuote{我所事奉的天主，是我祖先所事奉的，并以纯洁的良心事奉的，我感谢祂，}\BibleRef{弟后 1:3}那么，你们会说，不洁的人怎么配得福音呢？因为他们愿意并渴望它。所以一种人虽然处于错误，但因他们摆脱了情欲而被祂吸引；另一种人自愿接近，也不会被推开。许多人甚至从他们的祖先起就接受了真宗教。\par

% structure: p0012 source=h3 target=subsubsection reason=relative-heading-rank; base=h2

\subsubsection{格林多前书 3:3}

5. \BibleQuote{因为在你们中间有嫉妒和纷争。}\par

在这一点上，他准备与那些顺从的人争辩：因为前面他已经推翻了教会的领导者，他说言谈的智慧毫无价值。但在这里，他攻击那些在属下的人，他说：\par

% structure: p0015 source=h2 target=subsection reason=relative-heading-rank; base=h2

\subsection{格林多前书 3:4}

\BibleQuote{因为有人说：我是属保禄的，我是属阿颇罗的，你们岂不是属血肉的吗？}\par

他指出，这一点不但没有帮助他们，也没有使他们获得任何东西，反而成了他们在更大事上获益的阻碍。因为这事产生了嫉妒，嫉妒使他们成了\Quote{属血肉的}；成了\Quote{属血肉的}之后，他们就不能自由地听那些更崇高的事理了。\par

% structure: p0018 source=h2 target=subsection reason=relative-heading-rank; base=h2

\subsection{格林多前书 3:5}

\BibleQuote{那么，保禄是谁？阿颇罗是谁？}\par

这样，在提出并证明事实之后，他从此更加公开地提出控告。而且他用了自己的名字，消除了所有的严酷，不让他们对所讲的话感到愤怒。因为如果\Person{保禄}什么也不是，而且不抱怨，那么他们更不应该觉得自己受了亏待。你看，他有两种安慰他们的方法：第一，提出他自己的人；第二，不剥夺他们的一切，好像他们什么也没有贡献。相反，他允许他们有一点小小的份额：虽然小，他还是允许的。因为他说了：\Quote{保禄是谁？阿颇罗是谁？}之后，他加上：\Quote{不过是服务者，你们藉着他们信了。}这件事本身是伟大的，值得大赏报：虽然相对于原型和一切美善的根源来说，它算不了什么。（因为赐予我们美善的不是\Quote{服侍}的人，而是提供并赐予它们的人，祂才是我们的恩主。）他没有说\Quote{福音传道者}，而是说\Quote{服务者}，这更胜一筹。因为他们不仅传了福音，而且也服侍了我们；前者仅是言语之事，而后者还包含行为。所以，如果连基督也只是美善的服务者，而不是根源和泉源（我是指身为子的意义上），那么看这件事被引向何等深沉崇高的地步：\OriginalTerm{事情被引向何等深度与高度}{\GreekText{ποῦ} \GreekText{τὸ} \GreekText{πρᾶγμα} \GreekText{κατάγεται}}.\Quote{它被引向何等深沉崇高。}那么，你们会说，祂怎么被称为割礼的服务者呢？\BibleRef{罗 15:8}他在那里是在肉身中奥秘的救恩计划，与我们刚才说的意义不同。因为在那里，\Quote{服务者}的意思是\Quote{完成者}（\OriginalTerm{完成者}{\GreekText{πληρωτὴν}}，即预象的完成者），而不是说祂从自己的库存中赐予恩赐。\par

再者，他没有说：\Quote{那些引导你们进入信仰的人，}而是说：\Quote{你们藉以相信的人；}又将更大的份归给他们，并由此表明执事的从属类别（\OriginalTerm{并由此表明执事}{\GreekText{τοὺς} \GreekText{διακόνους} \GreekText{κὰντεῦθεν} \GreekText{δηλῶν}}）。既然他们是从属于另一位的，他们又怎能夺取权柄呢？但是，我要你们注意，他绝不怪罪他们夺取权柄，而是怪罪那些赋予他们权柄的人。因为错误的根基在于群众；既然一方堕落了，另一方也就瓦解了。他巧妙地预备了两点：首先，如同\OriginalTerm{挖掘}{\GreekText{ὑπορύξας}}一样，在必须推翻祸患的地方预备了；其次，在他们方面，不招致憎恨，也不使他们更加好斗。\par

\BibleQuote{正如基督（\OriginalTerm{主}{\GreekText{ὁ} \GreekText{Κύριος}}，公认文本）赐给各人。}\par

因为连这小事本身也不是出于他们，而是出于天主，是祂放在他们手中的。免得他们说：\Quote{那么，我们岂不该爱那些服事我们的人吗？}是的，他说；但你们要知道到什么程度。因为连这事本身也不是出于他们，而是出于赐予的天主。\par

% structure: p0024 source=h2 target=subsection reason=relative-heading-rank; base=h2

\subsection{格林多前书 3:6}

\BibleQuote{我栽种了，阿颇罗浇灌了，但使之生长的是天主。}\par

也就是说，我首先将圣言撒在地上；但为使种子不因诱惑而枯萎，阿颇罗加上了他的部分。但这一切都是出于天主。\par

% structure: p0027 source=h2 target=subsection reason=relative-heading-rank; base=h2

\subsection{格林多前书 3:7}

6. \BibleQuote{所以，栽种的不算什么，浇灌的也不算什么，唯有使之生长的天主才算什么。}\par

你注意到他如何安抚他们，使他们不致因听见\Quote{这个人是谁？}和\Quote{那个人是谁？}而过分恼怒吗？\Quote{不，两者都令人反感，即是说：\InnerQuote{这个人是谁？那个人是谁？}以及说\InnerQuote{栽种的和浇灌的都不算什么。}}那么，他如何缓和这些说法呢？首先，将轻蔑归在自己的身上：\Quote{保禄是谁？阿颇罗是谁？}其次，将一切归于赐予万物的天主。因为他说了\Quote{某人栽种了}，又加上\Quote{栽种的不算什么}，随后补充\Quote{但使之生长的是天主。}他甚至不在此止步，又用另一种治愈性的话语，即接下来的话。\par

% structure: p0030 source=h2 target=subsection reason=relative-heading-rank; base=h2

\subsection{格林多前书 3:8}

\BibleQuote{栽种的和浇灌的，都是一体。}\par

因为藉此他又确立了另一点，即他们不应彼此高举。他断言他们是一体，是指他们若无\Quote{使之生长的天主}便不能做什么。这样说，他既不容许那些劳苦较多的人凌驾于贡献较少的人之上，也不让后者嫉妒前者。其次，既然这有使人更加怠惰的倾向，即无论劳苦多寡，都被视为一体；请注意他如何纠正这一点，说：\Quote{但各人要照自己的劳苦，领取自己的报酬。}他仿佛说：不要因为我曾说你们是一体而惧怕；因为与天主的工程相比，他们是一体；然而，就劳苦而言，并非如此，而是\Quote{各人要领取自己的报酬。}\par

然后，达成了他所愿的，他更缓和了语气，并在允许的范围内慷慨地取悦他们。\par

% structure: p0034 source=h2 target=subsection reason=relative-heading-rank; base=h2

\subsection{格林多前书 3:9}

因为我们是天主的同工；\Quote{你们是天主的田园，天主的建筑物。}\par

你看到吗？他也给他们分派了不小的工作，先前已申明一切都是出于天主。因为他总是劝导他们服从那些统治他们的人，为此他避免过分轻视他们的教师。\par

\Quote{你们是天主的田园。}\par

因为他曾说过\Quote{我栽种了}，所以就保持比喻。如果你们是天主的田园，你们就不应被称为那些耕种你们之人的，而应被称为天主的。因为田园不被称为农夫家的，而是家主家的。\par

\Quote{你们是天主的建筑物。}\par

再者，建筑物不是工匠的，而是主人的。如果你们是建筑物，就不应被拆散；因为那样就不是建筑物了。如果你们是农田，就不应被分割，而应用一道围墙围起来，即同心合意。\par

% structure: p0041 source=h2 target=subsection reason=relative-heading-rank; base=h2

\subsection{格林多前书 3:10}

\BibleQuote{按照天主赐给我的恩宠，我好似一位精明的建筑师，奠立了根基。}\par

在这里，他称自己为精明的，并非抬高自己，而是给他们一个榜样，并指出这是精明人的部分，即奠立根基。你可以留意他谦虚态度的一个实例：在自称精明时，他不让这一点被视为出于自己；而是首先将一切归于天主，然后才用那个名字称呼自己。因为他说：\Quote{按照天主赐给我的恩宠。}这样，他同时表明一切都是出于天主，而恩宠最在于此，即不被分裂，而是立定于一个根基之上。\par

7. \BibleQuote{别人在这根基上建造；但各人要谨慎怎样在上面建造。}\par

我认为，在此处和下文，他使他们接受关于实践的考验，在他已将他们合而为一之后。\par

% structure: p0046 source=h2 target=subsection reason=relative-heading-rank; base=h2

\subsection{格林多前书 3:11}

\BibleQuote{因为除已奠立的根基，即耶稣基督外，谁也不能奠立别的根基。}\par

我说，无人能奠立它，只要他是精明的建筑师；但若他奠立（\OriginalTerm{奠立}{\GreekText{τιθῃ}}校勘读法，反对\OriginalTerm{已奠立}{\GreekText{τεθῃ}}。\Emph{杜纳乌斯，见萨维尔八卷·注页261。}），他就不是精明的建筑师了。\par

请看，他甚至从人的普通观念就证明了他的全部论点。他的意思是：\Quote{我宣讲了基督，我把根基交付了给你们。你们要谨慎怎样在这根基上建造，免得是在虚荣中，免得是将门徒引向人。}那么，我们就不当注意异端。\Quote{因为除了已经奠立的根基外，没有人能奠立别的根基。}那么，我们就当在这根基上建造，并且如同根基一样紧贴它，如同枝条紧贴葡萄树；在我们与基督之间不应有任何间隔。因为如果有任何间隔，我们立刻就丧亡。因为枝条因紧贴而吸收肥汁，建筑物因被粘合而屹立。因为如果它分离，就丧亡了，没有东西可以支持自己。那么，我们不要仅抓住基督，而要与他粘合，因为如果我们分离，就丧亡。\BibleQuote{因为凡远离祢的，必会丧亡；}\BibleRef{咏 72:27}经上这样说了。那么，让我们紧贴祂，并藉我们的行为紧贴祂。\BibleQuote{因为那遵守我诫命的，就住在我内}\BibleRef{若 14:21}因此，有许多比喻，祂藉以使我们合而为一。所以，如果你注意，祂是\Quote{头}，我们是\Quote{身体}：在头与身体之间能有什么空隔吗？祂是\Quote{根基}，我们是\Quote{建筑物}：祂是\Quote{葡萄树}，我们是\Quote{枝条}：祂是\Quote{新郎}，我们是\Quote{新娘}：祂是\Quote{牧人}，我们是\Quote{羊}；祂是\Quote{道路}，我们是\Quote{行路者}。再者，我们是\Quote{一座圣殿}，祂是\Quote{住者}；祂是\Quote{首生者}，我们是\Quote{弟兄}；祂是\Quote{继承者}，我们是\Quote{同祂一起的继承者}；祂是\Quote{生命}，我们是\Quote{活着的人}；祂是\Quote{复活}，我们是\Quote{复活的人}；祂是\Quote{光}，我们是\Quote{蒙光照的人}。这一切都表明合一；不容许任何空隔，连最小的也不。因为那移动一点点距离的，会继续直到变得非常遥远。因为同样，身体虽受刀只小小一割，也会丧亡；建筑物虽只有一条小裂缝，也会朽坏；枝条虽被从根上切下只有一会儿，也会无用。所以这小事不是小事，而是几乎就是全部。因此，每当我们犯了些小过失或甚至疏忽，我们不要忽视那微小的；因为那被忽略的，迅速变成大的。同样，当一件衣服开始破裂而被忽视时，它往往使裂缝扩大至全衣；屋顶，当几片瓦落下而被忽视时，会使整座房屋倒塌。\par

8. 那么，我们当牢记这些事，永远不要轻视小事，免得我们陷入大事。但如果我们轻视了它们，进入了罪恶的深渊，即使到了那里，我们也不要灰心，免得我们陷入\OriginalTerm{麻木}{\GreekText{καρηβαρίαν}}。因为从那以后出来是困难的，对那极不警醒的人；不仅因为距离，也因为我们所处的境况本身。因为罪也是深坑，常常压碎人。正如那些掉进井里的人不容易出来，但需要别人把他们拉上来；同样，那陷入任何罪之深处的人也是如此。对于这样的人，我们必须放下绳子，将他们拉上来。不，我们不仅需要别人，也需要我们自己，为使我们自己能够抓住自己并上升，我说不是像我们下降那么多，而是更远，如果我们愿意：为什么呢？天主也帮助：因为祂不愿罪人死亡，而愿他悔改。那么，不要让任何人绝望；不要让任何人有无神之人的感觉；因为这类罪正属于他们：\Quote{无神之人陷入任何罪恶的深处，就轻视它。}所以，不是人罪的众多造成他们的绝望，而是他们的无神之心。\par

那么，即使你在邪恶中已走尽一切极端，仍要对自己说：天主是爱人并渴望我们得救的；因为祂说：\BibleQuote{纵然你们的罪如朱红，我必使你们白如雪，}\BibleRef{依 1:18}并且我要将你们改变成相反的习惯。因此，我们不要绝望放弃；因为跌倒不像躺卧在跌倒之处那样可悲；受伤不像受伤后拒绝医治那样可怕。\BibleQuote{谁能夸口说他的心是纯洁的？谁能自信地说他无罪？}\BibleRef{箴 20:9}我说这些话，不是为使你们更疏忽，而是为防止你们绝望。\par

你想知道我们的主人多么良善吗？那个税吏满身万恶，上到圣殿，只说了句\Quote{求你开恩可怜我}，就下去成了义人。\EditorialNote{圣 48:13,14}是的，天主藉先知说：\Quote{因罪我暂时使他忧伤，\BibleRef{依 57:17}我看见（\OriginalTerm{我看见}{\GreekText{εἶδον} \GreekText{δτι}}不在七十七译本中）他忧伤哀戚，\OriginalTerm{我医治了他的道路}{\GreekText{ἰασάμην} \GreekText{αὐτὸν}}（七十七译本）。}有什么能与这种慈爱相比呢？条件是他\OriginalTerm{使他忧伤}{\GreekText{ἲνα} \GreekText{στυγνάση}}（参\BibleRef{若 8:56}）\OriginalTerm{使他看见那日子}{\GreekText{ἲνα} \GreekText{ἴδη} \GreekText{τὴν} \GreekText{ἡμέραν}}只是\Quote{忧伤}，祂说：\Quote{我赦免了他的罪。}但我们连这事也不做：因此我们特别惹天主发怒。因为那甚至因小事就被安抚的，当祂连这些也遇不到时，当然愤怒，要求我们最后的刑罚；这出自极端的藐视。例如，有谁曾为自己的罪而忧伤？有谁自怨自艾？有谁捶胸？有谁焦虑？依我看来，没有一人。然而人们为死去的仆人、为金钱的损失，无日不哭；而对于我们日日糟蹋的灵魂，却毫不关心。那么，当你甚至不知道自己犯了罪时，你怎能平息天主呢？\par

\Quote{是啊，}有人说，\Quote{我犯了罪。} \Quote{是啊，}你只是用舌头向我说：要用你的心思向我说，并用言语深切哀恸，好使你常有喜乐。因为，如果我们为自己的罪过忧伤，如果我们为自己的过犯深切哀恸，就没有别的事能使我们悲伤，这一种痛苦就会驱散所有沮丧。那么，从彻底的忏悔中，我们还得到另一件事：就是不致被现世生活的痛苦所淹没（\OriginalTerm{被淹没}{\GreekText{βαπτίζεσθαι}}），也不因它的光辉而骄傲。这样，我们又能更完全地平息天主的怒气；正如我们目前的行为惹动祂的怒气。请告诉我，如果你有一个仆人，他在其他仆人手中受了许多苦，却毫不介意其他任何人，只担心不激怒主人；难道仅凭这一点他不能消除你的怒气吗？但是，如果他不在意得罪你，却对得罪其他仆人满心忧虑；你难道不给他更重的惩罚吗？天主也是这样：当我们忽视祂的忿怒时，祂就重重地加在我们身上；但当我们留意它时，就轻一些。是的，甚至祂根本不再加给我们。祂愿意我们亲自为自己的过犯施行报复，此后祂自己就不再施行。这正是祂之所以威胁要惩罚的原因：藉着恐惧消除我们的藐视；当单单警告就足以使我们恐惧时，祂就不让我们实际经受考验。请看，例如，祂对耶肋米亚所说的话，\BibleRef{耶 7:17} \BibleQuote{难道你们看不见他们所作的事？他们的父亲点火，他们的儿子拾柴，他们的女人揉面。}恐怕同样的事也会论到我们说：\Quote{难道你们看不见他们所作的事？没有人寻求基督的事，只寻求自己的事。他们的儿女奔向不洁，他们的父亲奔向贪婪和强夺，他们的妻子不但不阻止丈夫脱离生活的虚荣和浮华，反而刺激他们的欲望。}只要站在市场上，盘问来来往往的人，你不见一个人忙于属灵的事，却都奔向属肉体的东西。我们何时才从醉梦中醒来？我们还要沉入沉睡多久？我们受的恶事还不够吗？\par

9. 然而有人可能认为，即使没有言语，经验本身也足以教你明白现世之物的虚无和极其卑劣。无论如何，曾有一些人，仅仅运用外邦人的智慧，对未来一无所知，却因已经证实了现世之物极大的无价值，就为此缘故放弃了它们。那么，你匍匐在地，不为那伟大而永恒的事物轻视那微小而短暂的事物，又亲耳听到天主亲自向你宣告并启示这些事，并从祂那里得到如此多的应许，你还能期望得到什么宽恕呢？因为这里的事物没有足够的力量留住一个人，这一点已由那些即使没有更大事物的应许也远离它们的人所表明。他们期望什么财富，以至于变得贫穷？没有。而是因为他们深知这样的贫穷胜过财富。他们期望什么样的生活而放弃奢华，投身严格的纪律？没有。而是因为他们已经意识到事物的本质，并觉察到在这两者中，哪一个更适合灵魂的严格训练和身体的健康。\par

我们既恰当地衡量这些事，并不断在心中思虑将来的福乐，就让我们从现世世界退出，好能获得那将要来临的另一个世界；藉着我们的主耶稣基督的恩宠和慈爱，祂和圣父及圣神，等等，等等。\par

\SourceRecord{Translated by Talbot W. Chambers.；Nicene and Post-Nicene Fathers, First Series；Vol. 12.；Edited by Philip Schaff.；Buffalo, NY: Christian Literature Publishing Co.,；1889.；<http://www.newadvent.org/fathers/220108.htm>.}

% structure: p0001 source=h1 target=section reason=page-title

\section{\WorkTitle{格林多前书第九篇讲道}}

% structure: p0002 source=h2 target=subsection reason=relative-heading-rank; base=h2

\subsection{\BibleBook{格林多}{前书} 3:12-15}

\BibleQuote{若有人在这根基上建造，或金、或银、或宝石，或木、草、禾稭；各人的工程将来必显露出来；因为那日子要指明它，它要在火中启示出来，那火要试验各人的工程如何。若人的工程存留得住，他就要得赏报；若人的工程被烧了，他就要受亏损；他自己固然可得救，但如从火中经过一样。}\par

我们所提出的并非一个小的探究题目，而是关于那些最必要且众人都探究的事；即地狱之火是否有终结。因为基督确实说过它没有终结：\Quote{他们的火永不熄灭，他们的虫子不死。}\BibleRef{谷 8:44,46,48}\par

我知道，你们听到这些事时，会感到\OriginalTerm{战栗}{\GreekText{ναρκᾶτε}}；但我能做什么呢？因为这是天主的命令，要不断在你们耳中重述这些事，祂说：\Quote{你要吩咐这百姓}\BibleRef{出 19:10}（七十士译本，此处希腊文\OriginalTerm{吩咐}{\GreekText{διάστειλαι}}），而我们既被委任为圣言的职务，就必须使听者感到痛苦，不是出于自愿，而是出于不得已。不，如果你们愿意，我们可以避免使你们痛苦。因为祂说：\BibleRef{罗 13:3}\Quote{如果你行善，就不必惧怕。}\par

正如我刚才所说，它没有终结，基督已经宣告了。保禄也指出惩罚的永恒性，说罪人\Quote{要受毁灭的惩罚，且是永远的}\BibleRef{得后 1:9}。又说\BibleRef{格前 6:9}\Quote{你们不要自欺：无论是淫乱的、奸淫的、作娈童的，都不能继承天主的国。}在给希伯来人书中也说\BibleRef{希 12:14}\Quote{你们要努力与众人和平，并追求圣化，没有圣化，谁也不能见主。}基督也对那些说\Quote{我们因祢的名行了许多奇事}的人说：\Quote{离开我吧，我不认识你们，你们这些作恶的人。}\BibleRef{玛 7:22}还有那些被关在门外的童贞女，她们再也没有进去。至于那些没有给祂食物的人，祂说\BibleRef{玛 25:46}\Quote{他们要进入永罚里去。}\par

2. 不要对我说：\Quote{如果惩罚没有终结，正义的原则如何完全保持？}相反，当天主行事时，要服从祂的决定，不要用人的理性来揣度祂说的话。而且，一个人从起初就经历了无数恩惠，却行了应受惩罚的事，既不受威胁也不受恩惠而改过，他受惩罚怎么可能是非正义的呢？因为如果你探究绝对的正义，按照严格正义的定义，我们本该从一开始就立即灭亡。不，即使按照正义的原则，也不仅仅如此；如果我们遭受这个，结果也包含了仁慈。因为当一个人侮辱那个未曾亏待他的人时，按照正义的原则，他受惩罚；但当他的恩人，以前不受任何恩惠约束，却赐予了无数恩惠，他是他存在的唯一作者，他是天主，将灵魂吹入他内，赐予千万恩宠，愿意将他提到天上——我说，当这样一位恩人，在接受了如此大的祝福之后，却受到对方的侮辱，每日的侮辱，那人如何可以被认为配得宽恕？你没有看到天主如何因一个罪惩罚了亚当吗？\par

\Quote{是的，}你会说，\Quote{但祂给了他乐园并使他享受了许多恩宠。}不，肯定这不是一回事：一个人在享受安全和安逸时犯罪，与在极度痛苦中犯罪是不同的。事实上，可怕的是你的罪并非在乐园中犯的，而是在今生无数的恶事中；即使遭受痛苦，你也没有清醒，好像一个在监狱中的人仍然行恶一样。然而，祂向你应许了比乐园更大的事。但祂现在没有赐予，免得你在战斗的时节松懈；也没有保持沉默，免得你在劳苦中完全沮丧。至于亚当，他犯了仅一个罪，就给自己带来了确实的死亡；而我们每天犯成千上万的过犯。那么，如果他因那一次行为给自己带来如此大的灾祸并引入了死亡，我们这些不断生活在罪恶中，且代替乐园期望天堂的人，将遭受什么呢？\par

这论证令人不快，使听者痛苦：仅凭我自己的感受，我知道这一点。因为我的心确实不安和悸动；我越看见地狱的描述被证实，我就越因恐惧而战栗和退缩。但有必要说这些事，免得我们堕入地狱。你所接受的不是乐园，也不是树木和植物，而是天堂和天上的美好事物。现在，如果那接受较少的人被定了罪，没有任何考虑可以豁免他，那么何况我们这些犯罪更多、并被召叫接受更大事物的人，将遭受无法补救的灾祸呢？\par

请想一想，仅因一个罪，我们的种族就在死亡中待了多久。五千年多过去了，死亡仍未消除，只因一个罪。我们甚至不能说亚当曾听过先知，见过别人因罪受罚，理应因之恐惧而改正，哪怕仅仅通过榜样。因为他那时是第一个，也是唯一的；然而他仍受了罚。但你们却没有什么可以推诿，你们在这么多榜样之后变得更糟；你们被赐予了如此卓越的圣神，却给自己招来了不是一个、两个、三个，而是无数的罪！因为，不要因为罪是在一瞬间犯下的，就认为惩罚也必是一瞬间的事。你们难道没有看见那些人，仅因一次偷盗或一次奸淫，在一瞬间犯下，却常常在监狱和矿坑中度过一生，忍受着持续的饥饿和各样的死亡？没有人能释放他们，或说：\Quote{罪行发生在一瞬间；惩罚也应与罪的时间相等。}\par

3. 但有人会说：\Quote{做这些事的是人；至于天主，祂是爱人的。} 首先，就连人这样行也不是出于残忍，而是出于人道。而天主自己，正如\Quote{祂是爱人的}，也以同样的性情惩罚罪恶。\BibleRef{宗 16:12} \BibleQuote{因为祂的仁慈何其大，祂的惩罚也何其大。} 所以，当你对我说\Quote{天主是爱人的}时，你就告诉我惩罚的更大理由：就是我们将这样一位存在得罪了。因此保禄也说：\BibleRef{希 10:31} \BibleQuote{落入永生天主的手里，真是可怕。} 我恳求你，忍受这言语的火热力量，因为或许——或许你会由此得到一些安慰！谁能像天主那样惩罚呢？当祂引起洪水，完全毁灭了一个如此众多的种族；稍后又从天降火，将他们全部消灭？人的惩罚怎能与之相比？难道你没有看见，今世的惩罚几乎是永恒的吗？四千年过去了，索多玛人的惩罚仍然高悬。因为祂的仁慈何其大，祂的惩罚也何其大。\par

再者，如果祂加给我们任何沉重或不可能的事，或许我们可以强调律法的困难；但如果它们极其容易，我们还有什么可说，因为我们连这些也不遵守？假设你不能守斋或守贞；尽管你若愿意就能，而那些能做到的人恰是我们的定罪。然而，天主并未对我们如此严格；祂也没有规定这些或立为法律，而是将选择留给听者的自由。然而，你能在婚姻内守节；你能戒除醉酒。你不能舍弃你所有的财产？不，你当然能；那些这样做的人证明了这一点。但尽管如此，祂没有规定这个，而是命令不可贪婪，并用自己的财物帮助那些有需要的人。但如果有人说：\Quote{我连只满足于一个妻子都不能}，他是在自欺并错误地推理；而那些没有妻子而守独身的人则定了他的罪。可是，告诉我，你怎么能不使用辱骂的言语？你不能不咒骂吗？何必呢，做这些事是令人厌烦的，而不是禁止做它们。那么，我们还有什么借口不遵守如此容易而轻松的诫命？我们完全无法提出任何借口。那么，从以上所说，惩罚是永恒的就显而易见了。\par

4. 但由于保禄的话在某些人看来似乎支持相反的观点，让我们也把它拿出来彻底检查。因为他说过：\Quote{若有人在这根基上建造的工程存留，他将获得赏报；若有人所建造的工程被烧毁，他将遭受损失}，然后又加上：\Quote{但他自己却要得救，却要像从火中经过一样。} 对此我们该说什么呢？首先让我们考虑什么是\Quote{根基}，什么是\Quote{金子}，什么是\Quote{宝石}，什么是\Quote{草}，什么是\Quote{禾秸}。\par

\Quote{根基}他自己明确指明是基督，说：\Quote{因为除已奠立的根基以外，没有人能奠立别的}，他说：\Quote{就是耶稣基督。}\par

接下来，这建筑在我看来是指行为。虽然有些人认为这也是指教师和门徒以及关于腐败的异端；但推理不允许这样。因为如果是这样，那么当\Quote{工程被毁灭}时，为什么那个\Quote{建筑者}要\Quote{得救}，尽管是\Quote{像从火中经过}？按理，作者更应失丧；但现在却会发现，那被建造在工程中的人承受更重的刑罚。因为如果教师是恶行的原因，他就应受更重的惩罚：那么他怎能\Quote{得救}？反之，如果他不是原因，而是门徒因自己的乖戾而变成这样，他根本没有受罚的资格，也不该遭受损失——我指的是建造得如此好的人。那么他说\Quote{他将遭受损失}是什么意思？\par

从这一点可以看出，所谈论的是行为。因为既然他接下来要全力对付那犯奸淫的人，他就从高处和很久以前开始预备前提。因为他知道，在讨论一个主题时，如何在关于那件事的论述中，为他要转而讨论的另一件事打下基础。例如，他在责备他们吃饭时彼此不等候时，就为论及奥迹的论述打下了基础。也正因如此，既然他现在正赶着要对付那奸淫者，他在谈论\Quote{基础}时，就补充说：\BibleQuote{你们不知道你们是天主的殿，天主的圣神住在你们内吗？若有人毁灭（\OriginalTerm{毁灭}{\GreekText{Φθείρη}}，修订版作\InnerQuote{玷污}）天主的殿，天主必要毁灭这人。}他说这些话，是为了开始用恐惧来震撼那失节者的灵魂。\par

% structure: p0017 source=h3 target=subsubsection reason=relative-heading-rank; base=h2

\subsubsection{\BibleRef{格前 3:12}}

5. \BibleQuote{若有人在根基上建造金银、宝石、木、草、禾秸。}因为信德之后需要建树，因此他在别处说：\BibleQuote{你们要用这些话彼此劝勉。}\BibleRef{得前 5:11}因为工匠和学习者都参与建树。因此他说：\BibleQuote{但各人要谨慎怎样在上面建造。}\BibleRef{格前 3:10}但如果这些话是关于信德的，那么所肯定的就不合理。因为信德方面，所有人都应当平等，因为\BibleQuote{只有一个信德；}\BibleRef{弗 4:5}但在生活的良善上，不可能所有人都一样。因为信德不是在此处更小，在彼处更卓越，而是在所有真正相信的人身上相同。但在生活中，有的人更勤奋，有的更懒惰；有的更严格，有的更平常；有的在更大的事上做得好，有的在较小的事上；有的错误更严重，有的不太显著。因此他说：\BibleQuote{金银、宝石、木、草、禾秸——各人的工程必将显露：}——他所谈论的是品行：——\BibleQuote{若有人在所建造的工程上存留，他必得赏报；若人的工程被烧毁，他必受亏损。}然而，如果这句话与门徒和教师有关，他不应该因为门徒不听而\Quote{受亏损}。因此他说：\BibleQuote{每个人要照自己的劳苦得自己的赏报}，不是照结果，而是照\Quote{劳苦}。因为如果听者不留意呢？所以这段经文也证明这话是关于行为的。\par

现在他的意思是：如果一个人有正确的信德但生活邪恶，他的信德不能保护他免受惩罚，他的工作被烧毁。\Quote{被烧毁}的意思是\Quote{不能承受火的威力。}但就像一个人穿着金盔甲穿过火河，他从穿越中变得更为明亮；但如果他带着禾秸穿过，他不仅无益，反而毁灭自己；人的工作也是如此。因为他说这话，不是讨论物质的东西被烧毁，而是为了使他们的恐惧更强烈，并显示那留在邪恶中的人毫无防御。因此他说：\BibleQuote{他必受亏损：}看，这是第一惩罚：\BibleQuote{但他自己却要得救，可是这样得救，就像从火中经过的一样；}看，这又是第二个。他的意思是：他自己不会像他的工作那样灭亡，化为乌有，而是留在火中。\par

6. 然而，他称此为\Quote{救恩}，你会说；这就是为什么他加上\Quote{这样得救，就像从火中经过的一样}的原因；因为当我们谈论那些不会立即燃烧成灰的物质时，我们也常说\Quote{它在火中被保存。}因为听到火这个字时，不要以为那些被焚烧的人就化为无有。尽管他称这种惩罚为救恩，也不要惊讶。因为他的习惯是在听起来不好的事上使用好的词语，而在好的事上则相反。例如，\Quote{俘虏}一词似乎是邪恶的名称，但\Person{保禄}却用在好的意义上，当他说：\BibleQuote{将一切思想掳去，使之顺服基督。}\BibleRef{格后 10:5}同样，他又用好的词语称呼邪恶的事，说：\BibleQuote{罪恶为王，}\BibleRef{罗 5:21}这里\Quote{为王}一词显然是好听的。所以在这里他说\Quote{他必得救}，只是暗指刑罚的严重性；好像说：\Quote{但他自己要永远留在惩罚中。}然后他推论说：\par

% structure: p0021 source=h2 target=subsection reason=relative-heading-rank; base=h2

\subsection{\BibleRef{格前 3:16}}

7. \BibleQuote{你们不知道你们是天主的殿吗？}因为他在前面一段中谈论了那些分裂教会的人，然后也攻击那犯了不洁的人；虽然不是明说，而是概括地；暗示他腐败的生活方式，并强调那已赐给他的恩赐的罪过。然后他使其他所有人羞愧，从他们已经拥有的这些祝福出发来进行论证：因为他一贯这样做，要么从未来，要么从过去，无论是令人悲伤的还是令人鼓舞的。首先，从未来：\BibleQuote{因为主的日子要把它揭露出来，因为主的日子要在火中出现。}再次，从已发生的事：\BibleQuote{你们不知道你们是天主的殿，天主的圣神住在你们内吗？}\par

% structure: p0023 source=h2 target=subsection reason=relative-heading-rank; base=h2

\subsection{\BibleRef{格前 3:17}}

\BibleQuote{若有人毁灭天主的殿，天主必要毁灭这人。}你看到他的话语的横扫之势了吗？然而，只要那人未被指明，所说的话就不那么令人反感，大家都分担责备的恐惧。\par

\BibleQuote{天主必要毁灭这人，}就是使他灭亡。这并非咒诅的话，而是预言。\par

\BibleQuote{因为天主的殿是圣的：}但那犯奸淫的人是亵渎的。\par

然后，为使他不致显得仅仅在这一件事上全力以赴，当他说\Quote{因为天主的殿是圣的}时，他又加上\Quote{这殿就是你们}。\par

% structure: p0028 source=h2 target=subsection reason=relative-heading-rank; base=h2

\subsection{格林多前书 3:18}

8. \Quote{谁也不要自欺。}这同样是指那个人，他自以为有点什么，并以智慧自夸。但为使他不致显得在离题中冗长地责备他，他先让他陷入一种痛苦，将他交付于恐惧，然后将话题转回普遍的过错，说\Quote{若你们中有人在这世上自以为有智慧，就让这人变为愚妄，好使他\OriginalTerm{成为}{\GreekText{γένηται} . rec. vers.~'be'}有智慧的。}后来他以极大的胆量这样做，因为他已经充分击垮了他们，并用那恐惧震撼了不仅是那污秽之人的心，也是所有听众的心：他如此精确地衡量了他所要讲的话的限度。因为若一个人富有，若他高贵，当他被罪恶俘获时，他比所有卑贱者更为卑贱。正如若一个人是国王却被蛮族奴役，他是所有人中最悲惨的，对于罪恶也是如此：因为罪恶是蛮族，一旦俘虏了一个灵魂，她不懂得宽恕，反而暴虐地毁灭所有接纳她的人。\par

9. 因为没有什么是比罪恶更轻率的：没有什么如此无感觉，如此完全愚蠢和狂暴。凡它所落之处，一切都倾覆、混乱、毁灭。观之不雅，可憎又沉重。若有一位画家描绘她的图像，我想他不会错画成这样的：一个具有野兽形态的女人，野蛮，吐着火焰，丑陋，黝黑，如同异教诗人所描绘的斯库拉。因为她用千万只手攫取我们的思想，出其不意地袭来，将一切撕碎，如同那些偷偷咬人的狗。\par

但不如说，何需画家的技艺，我们更应举出那些按照罪恶的模样所造的人？\par

那么你们愿意我们先描绘谁？贪婪和强夺的人？还有比那眼睛更无耻的吗？还有比那更不谦逊、更像贪婪的狗的吗？因为当他攫取所有人的财物时，没有一条狗能像他那样厚颜无耻地坚守阵地。还有比那手更污秽的吗？还有比那口更放肆的吗？吞下一切，永不满足。不，不要看那面容和眼睛像是人的。因为这样的面容不属于人的眼睛。他不看人如人；他不看天如天。他甚至不抬头向主；但一切在他眼中都是金钱。人的眼睛惯于看贫穷的人受苦，并会变得柔和；但贪婪之人的眼睛，一见穷人，就如野兽般怒视。人的眼睛不把别人的财物看作自己的，反而把自己的看作别人的；他们不贪图给予别人的东西，反而用自己的东西周济别人：但这些人的眼睛，除非夺取所有人的财产，绝不满足。因为他们所有的不是人的眼睛，而是野兽的。人的眼睛不能忍受看到自己赤身露体（因为即使那是别人的身体，也是自己的），但这些人的眼睛，除非剥夺每个人的一切，并把所有人的财产都放在自己家里，永远不会餍足；是的，他们永不知足。以至于人们可能会说，他们的手不仅是野兽的，甚至比这些更为野蛮和残忍。因为熊和狼在饱足时会停止它们的吃法：但这些不知任何饱足。为此，天主造了我们手，是为了帮助别人，而不是为了谋害他们。若我们为此目的使用它们，倒不如砍掉它们，让我们没有手更好。但你，若野兽撕碎一只羊，你会忧伤；但当你对自己骨肉同胞做同样的事时，你以为你的行为丝毫不残暴吗？那么你怎么能是人呢？你没有看到吗？我们称一件事为人道的，是指它充满怜悯和仁爱。但当一个人做出任何残忍或野蛮的事时，我们称他为不人道的。你看，描绘人的标记就是显示怜悯；野兽则相反，正如常言所说：\Quote{为何是野兽或狗？}\BibleRef{列下 8:13}。因为人救济贫困，而不是加重贫困。同样，这些人的口是野兽的口；不如说，这些比野兽更凶猛。因为他们说出的话所释放的毒液，比野兽的牙齿更致命，制造杀戮。若有人逐一细究，就会清楚地看到，不人道如何使那些实行它的人从人变为野兽。\par

10. 但若他探查那类人的心灵，他就不再只称他们为野兽，而是邪魔。首先，他们充满极大的残忍和对同仆的仇恨：\BibleRef{玛 18:33}那里既没有天国之爱，也没有地狱之惧；没有对人的敬畏，没有怜悯，没有同情；只有无耻和胆大，以及藐视一切未来的事。天主关于惩罚的话语在他们看来像是寓言，他的威胁如同玩笑。因为贪婪之人的心就是这样。既然内在是邪魔，外在是野兽；是的，比野兽更坏；我们该把这样的人置于何处？因为他们甚至比野兽更坏，从这一点可以清楚看出：野兽因本性如此；而这些人，本性温柔，却强力反抗本性，把自己训练成野蛮的。邪魔也有那些谋害者为人相助，以致若没有这种帮助，他们对我们的阴谋大半就会消失：但这些人，当那些被他们恶意对待的人在同样对待他们时，他们仍试图比他们更恶毒。再者，魔鬼与人争战，不与同类邪魔争战；但我们所说的这个人，千方百计地伤害自己的同族和同家，甚至不尊重本性。\par

我知道许多人因这些话而恨我们；但我并不恨他们，反而怜悯并哀哭那些如此存心的人。即使他们愿意动手，我也乐意忍受，只要他们能摆脱这种野蛮的心境。因为不仅是我，还有先知也与我一起，将一切这样的人逐出人类家庭，说：\BibleRef{咏 48:20} \BibleQuote{人在尊贵中毫无智慧，反倒如同无知的畜生。}\par

那么，让我们终于成为人，让我们仰望上天；并让我们领受并恢复那按照祂肖像的，\BibleRef{哥 3:10} 使我们也能获得将来的福分，借着我们的主耶稣基督的恩宠和慈爱，愿光荣、权能、尊崇归于祂，与圣父及圣神，从今时直到永远。阿们。\par

\SourceRecord{Translated by Talbot W. Chambers.；Nicene and Post-Nicene Fathers, First Series；Vol. 12.；Edited by Philip Schaff.；Buffalo, NY: Christian Literature Publishing Co.,；1889.；<http://www.newadvent.org/fathers/220109.htm>.}

% structure: p0001 source=h1 target=section reason=page-title

\section{第十篇讲道：\WorkTitle{格林多前书}}

% structure: p0002 source=h2 target=subsection reason=relative-heading-rank; base=h2

\subsection{格前 3:18-19}

\BibleQuote{不要自欺。若有人（\OriginalTerm{}{\GreekText{ἐνὗμῖν}}省略）以为自己在这个世界是聪明的，就让他变成愚拙，好使他成为聪明的。因为今世的智慧在天主面前是愚拙的。}\par

正如我先前所说的，他在适当时机之前就发起了对行淫者的指控，并用寥寥数语半遮半掩地透露了这件事，使那人的良心战栗之后，他又急忙转向与外邦智慧的争战，以及那些因此自高自大、分裂教会之人的指控：为要补全剩余的部分，精确地完成整个主题，然后他才能让自己的舌头带着强烈的冲动去攻击那不洁之人，而在此之前他只是与他有过初步的交锋。因为这句话：\BibleQuote{不要自欺}——主要是针对他，并用恐惧预先压制他；而关于\BibleQuote{碎秸}的说法，最适合暗示他。同样，这句话：\BibleQuote{你们不知道你们是天主的圣殿，天主的圣神住在你们内吗？}——也最为贴切。因为这两件事最能使我们远离罪恶：当我们记起为罪所定的惩罚时，以及当我们计算我们真正尊严的分量时。所以，他提出\BibleQuote{干草}和\BibleQuote{碎秸}来恐吓他们；但提及他们高贵出身的尊严时，又使他们羞愧；前者努力改正那些较为麻木的人，后者则改正那些较为周详的人。\par

2. \BibleQuote{不要自欺；若有人以为自己在这个世界是聪明的，就让他变成愚拙。}\par

正如他吩咐人，可以说，向世界死去——这种死去毫无害处，反而有益，成为生命的原因——同样，他也吩咐人向这个世界变成愚拙，借此向我们引进真正的智慧。现在，成为世界上的愚拙的人，是那些轻视外邦智慧、并相信它对理解信仰毫无贡献的人。所以，如同按照天主而言的贫穷是财富的原因，谦卑是高举的原因，轻视荣耀是荣耀的原因；同样，成为愚拙使人比所有人更聪明。因为在我们这里，一切都是相反的。\par

再者：为什么他没有说\Quote{让他放下智慧}，而是说\Quote{让他成为愚拙}？是为了最大限度地贬低外邦学问。因为说\Quote{放下你们的智慧}与说\Quote{成为愚拙}不是同一回事。此外，他也在训练人们不以我们中的粗陋为耻；因为他全然嘲笑一切外邦的事物。出于同样的理由，他毫不避讳这些名称，因为他信赖他所谈论之事的能力。\par

因此，正如十字架虽被视为屈辱，却成为无数祝福的创始者，以及说不出的荣耀的基础和根源；同样，那被视为愚拙的，却成了我们智慧的原因。因为，若有人学错了东西，除非他完全清除，使自己的灵魂平坦洁净，然后献给他要书写的人，否则他绝不能确切地认识任何健全的真理；同样，对于外邦的智慧也是如此。除非你彻底清除，清洁你的心思，如同无知的人一样将自己交托给信仰，你就不会精确地认识任何卓越之事。因为那些视力不佳的人，若不肯闭上眼睛，将自己托付给别人，反而相信自己的有缺陷的视力，他们将会犯比那些看不见的人更多的错误。\par

但你会说，人们如何放下这种智慧？通过不遵行它的教条。\par

3. 然后，既然他如此急切地吩咐人远离它，他便加上原因说：\BibleQuote{因为今世的智慧在天主面前是愚拙的。}因为它不仅无益，甚至有害。我们必须远离它，因为它造成伤害。你注意到他以何等高的手法夺取胜利的成果，证明它不但对我们毫无益处，甚至还是对手吗？\par

他不仅满足于自己的论证，还再次引用证言说：\Quote{因为经上记载，\BibleRef{约 5:13}他在自己的诡诈中捉住智慧人。}通过\Quote{诡诈}，即用他们自己的武器胜过他们。因为他们用智慧来消除对天主的一切需要，天主便用这智慧本身（而非其他方式）驳斥了他们，表明他们特别需要天主。如何及通过什么方法？因为他们因此成了愚拙，所以正当地被捉住。那些自认为不需要天主的人，竟陷入如此大的窘境，以至于显得不如渔夫和没有学问的人；并且从那时起就离不开他们。因此他说：\Quote{在他们的诡诈中}他捉住了他们。因为\Quote{我要毁灭他们的智慧}这句话是针对它没有引入任何有用的事物；而\Quote{他在自己的诡诈中捉住智慧人}则是为了显示天主的能力。\par

接下来，他也宣告天主捉住他们的方式，加上另一个证言：\par

% structure: p0013 source=h2 target=subsection reason=relative-heading-rank; base=h2

\subsection{格前 3:20}

他说：\BibleQuote{因为主知道人的思念\BibleRef{咏 93:11}是虚妄的（七十贤士译本）。}当那无限的智慧关于它们宣告这道谕令，并宣称它们如此时，你还寻求什么其他证据证明它们极端的愚拙？因为人的判断在许多情况下确实会出错；但天主的决断在任何情况下都是无懈可击、不腐败的。\par

4. 在如此光辉地树立了来自上天的审判的纪念之后，他在接下来的话中采用了某种激昂的风格，转而针对那些受他管辖的人（\OriginalTerm{受管辖的人}{\GreekText{ἀρχομένους}}），说：\par

% structure: p0016 source=h2 target=subsection reason=relative-heading-rank; base=h2

\subsection{格前 3:21}

\Quote{所以，谁也不可拿人来夸耀，因为一切都是你们的。}他再次回到先前的主题，指出他们甚至不应该因为属灵的事物而自高自大，因为自己一无所有。\Quote{既然外在的智慧是有害的，而神恩不是你们所赐，你们还有什么可夸耀的呢？}至于外在的智慧，他说：\Quote{谁也不可自欺，}因为他们自负于一种实际上害多于益的事物。但在这里，因为所说的事确实有益，他说：\Quote{谁也不可夸耀。}他更加温和地调整了他的话：\Quote{因为一切都是你们的。}\par

% structure: p0018 source=h2 target=subsection reason=relative-heading-rank; base=h2

\subsection{\BibleRef{格前 3:22-23}}

\Quote{无论是保禄，或是阿颇罗，或是刻法，或是世界，或是生命，或是死亡，或是现在，或是将来，一切都是你们的；你们却是基督的，基督是天主的。}因为他严厉地对待了他们之后，又使他们振作起来。正如他先前所说，\BibleRef{格前 3:9} \Quote{我们与天主同工；}并用许多其他表达安抚了他们：同样地，他在这里也说，\Quote{一切都是你们的；}以此打压教师的骄傲，表明他们非但没有给他们任何恩惠，反而应当感激别人。因为为了他们的缘故，他们才成为那样，而且甚至领受了恩宠。但看到这些人也肯定要夸耀，他就预先根除了这种疾病，说：\Quote{就如天主赐给了每人，}\EditorialNote{参见上文，卷六5.6节}以及\Quote{天主赐予生长：}这样，无论哪一方都不会因自己是善事的施予者而自高自大；另一方也不会在第二次听到\Quote{一切都是你们的}时再次得意洋洋。\Quote{因为，的确，即使是为了你们的缘故，但全是天主的工作。}我希望你观察他是如何一贯地以自己和伯多禄的名义作出假设。\par

\Quote{或是死亡？}就是说，即使他们死去，也是为你们而死，为你们的得救而冒危险。你注意到他如何再次打压门徒的高傲，而提升教师的精神吗？事实上，他对待他们就像对待出身高贵的孩童，他们有导师，并且要继承一切。\par

我们也可以从另一个意义上说，亚当之死是为了我们，好使我们受改正；基督之死也是为了我们，好使我们得救。\par

\Quote{你们却是基督的，基督是天主的。}在一个意义上\Quote{我们是基督的，}在另一个意义上\Quote{基督是天主的，}在第三个意义上\Quote{世界是我们的。}因为我们的确是基督的，作为他的工程：\Quote{基督是天主的，}作为真正的子嗣，而非工程：在这个意义上，世界也不是我们的。所以，虽然说法相同，但含义不同。因为\Quote{世界是我们的，}作为为我们的缘故而被造之物；但\Quote{基督是天主的，}因为天主是他存在的根源，因为他是父。而\Quote{我们是基督的，}因为他塑造了我们。现在他说：\Quote{如果他们是你们的，}他说，\Quote{那么你们做了什么与此相反的事，竟用他们的名字称呼自己，而不是以基督和天主的名呢？}\par

% structure: p0023 source=h2 target=subsection reason=relative-heading-rank; base=h2

\subsection{\BibleRef{格前 4:1}}

5. \Quote{人应当这样看待我们：作为基督的仆役，以及天主奥秘的管理者。}在他打压了他们的精神之后，注意他如何再次使之振作，说：\Quote{作为基督的仆役。}那么，你不可放弃主人，而从仆役和管理者那里接受名称。他说：\Quote{管理者；}表示我们不应将这些事物给予所有人，而只应给予那些应得的人，以及那些我们应当服务的人。\par

% structure: p0025 source=h2 target=subsection reason=relative-heading-rank; base=h2

\subsection{\BibleRef{格前 4:2}}

\Quote{此外，对于管理者，要求的是，他要显出忠信；}就是说，他不将主人的财物据为己有，不象主人一样为自己索取，而是作为管理者来管理。因为管理者的职责是妥善管理托付给他的一切：不说主人的东西是他自己的；反而，他自己的东西是主人的。让每一个人都思考这些事，无论是那有口才的，还是那拥有财富的，就是：他受了主人财物的委托，这些财物不是他自己的；他不可把它们留给自己，也不可记在自己的账上；而应把它们归于那赐予一切的天主。你愿看见忠信的管理者吗？请听伯多禄说：\Quote{你们为什么注视我们，好像是我们凭自己的能力或虔诚使这人行走呢？}\BibleRef{宗 3:12} 他也对科尔乃略说：\Quote{我们也是和你们一样的人；}并对基督自己说：\Quote{看，我们舍弃了一切，跟随了你。}\BibleRef{玛 19:27} 保禄也不例外，当他说：\Quote{我比他们众人更劳苦，}\BibleRef{格前 15:10} 又加上：\Quote{然而不是我，而是与我同的天主的恩宠。}在别处，他强烈反对同一批人，说：\Quote{你有什么不是你领受的呢？}\BibleRef{格前 4:7} \Quote{因为你没有一样是自己的，既没有财富，也没有口才，甚至没有生命本身；因为这也是主的。所以，当必要时，你也当舍弃它。但如果你贪恋生命，被命令舍弃时拒绝，你就再也不是忠信的管理者了。}\par

\Quote{当天主召唤时，人怎能抵抗呢？}这正是我所说的：正因此我尤其钦佩天主的慈爱，因为祂本可以违背你的意愿收取那些东西，却不愿你们不情愿地\OriginalTerm{缴纳}{\GreekText{εἰσενεχθῆναι}}它们，好让你们还可获得赏报。例如，祂可以不经过你的同意就取走生命；但祂愿意你自愿同意，好使你能与保禄一同说：\BibleQuote{我天天冒死}，\BibleRef{格前 15:31}。祂可以不经过你的同意就夺去你的荣耀，使你卑微；但祂愿意你心甘情愿地献上，好为你编成冠冕。祂可以使你贫穷，即使你不愿意，但祂要你自愿成为贫穷，好为你编织冠冕。你看到天主对人的仁慈吗？你看到我们自己的愚钝吗？\par

如果你已获得高位，并曾取得教会管理的某些职务，不要心高气傲。你并未取得荣耀，而是天主将其加于你。因此，要像用别人的东西一样，节俭地使用它；既不要滥用，也不要用于不当之事，更不要自高自大，或将其归为己有；反要自视为贫穷和卑微。因为——假若你受托保管君王的紫色袍，绝不可滥用袍子而糟蹋它，反而要更加谨慎地为赐予者保管它。你得了口才吗？不要自高，不要傲慢；因为恩赐不是你的。不要吝惜主的美物，要分施给你的同仆；不要因这些事而得意，仿佛它们是属你的，也不要吝于分施。再者，如果你有儿女，他们是天主的，你只是受托。你若这样想，就会因养育他们而感恩，若失去他们也不会过于悲伤。约伯就是这样，他说：\BibleRef{约 1:21} \BibleQuote{上主赐的，上主收回。}\par

因为我们所有的一切都来自基督。我们本身的存在就是藉着祂，还有生命、气息、光明、空气和大地。祂若将我们排除在任何一样之外，我们就丧亡了。因为\BibleRef{伯前 2:11} \BibleQuote{我们是旅客和寄居的}。所有关于\Quote{我的}和\Quote{你的}都不过是空洞的言辞，不代表实际。因为若你说房子是你的，那是没有实质的话：连空气、大地、物质都属于造物主；连你自己，建造房子的人，也是如此。即使你认为使用是你的，这也靠不住，不仅因为死亡，而且在死亡之前，因为事物的不稳定。\par

6. 我们应不断将这一切放在眼前，过严谨的生活；这样我们将获得两大益处。首先，我们无论拥有还是失去，都会感恩；我们也不会被转瞬即逝、不属于自己的事物所奴役。因为，无论祂取走财富，祂只是取走祂自己的；或名誉、荣耀、身体，甚至生命本身；若祂取走你的儿子，祂取走的不是你的儿子，而是祂自己的仆役。因为你并没有形成他，而是祂造了他。你只是为他的出现服务；整个作品是天主的工作。因此，我们应感谢蒙恩成为此事中的仆人。但是怎样？你希望他永远活着吗？这又显出你的吝啬，并且你不知道你拥有的是别人的孩子，不是你自己的。因此，那些甘愿放手的人只是意识到他们拥有不属于自己的东西；反之，悲伤的人实际上是将君王的财产视为己有。因为我们并不属于自己，他们怎能属于我们呢？我说，我们：因为我们在两方面属于祂：由于我们的受造，也由于信德。为此达味说：\Quote{我的实质与祢同在。} \EditorialNote{\GreekText{ὑπόστασις}} 七译本作\Quote{希望}，校正版，第6节；\BibleRef{咏 139:14}。保禄也说：\BibleQuote{因为我们生活、行动、存在都在祂内}，\BibleRef{宗 17:28}；他在申论信德时又说：\BibleRef{格前 6:19} \BibleQuote{你们不是自己的}，并且\BibleQuote{你们是用高价买来的}。因为一切都是天主的。因此，当祂召唤并愿意收取时，我们不要像吝啬的仆人一样逃避算账，也不要偷窃主人的财物。你的灵魂不是你的，你的财富又怎能是你的呢？既然如此，你怎么能把这些不属于你的东西浪费在无谓的事上呢？你不知道为这些我们很快就要受审判吗？也就是说，如果我们糟糕地使用了它们。然而，既然它们不是我们的，而是我们主人的，应该把它们用在我们的同仆身上。值得思考的是，这正是那富人受谴责的原因；也是那些没有给主食物的人受谴责的原因。\BibleRef{路 14:21}\par

7. 那么，不要说：\Quote{我不过是在花自己的钱，用自己的东西享乐。}那不是你的，而是别人的。我说是别人的，因为这是你自己的选择：天主的旨意原是让那些托付给你的东西，为你的弟兄们所用，成为你的。现在，不属于你的东西，若用在别人身上，就成为你的；但若你无度地用在自已身上，你自己的东西就不再是你的。因为，既然你残酷地使用它们，还说：\Quote{我自己的东西完全用于自己的享乐是公平的。}所以，我不称它们为你的。因为它们是你和你的同仆所共有的；就像太阳、空气、大地和其他一切都是共有的。关于身体，每一个功能属于整个身体也属于每个肢体；但若只应用于一个肢体，就会毁坏那个肢体本身的功能：财富也是如此。为使我的话更清楚：身体的食粮是众肢体共同分享的，若它只进入一个肢体，最终甚至对那个肢体也成了异己。因为当它不能被消化或提供营养时，我甚至说，对那个肢体，它也成了异己。但若它被共享，那个肢体和所有其他肢体都把它当作自己的。\par

同样地，关于财富。如果你独自享用，你也失去了它：因为你将得不到它的赏报。但如果你与他人共同拥有，那么它会更属于你，你也将从中获益。你难道没有看见手在服务，口在软化，胃在接收吗？难道胃会说：既然我接收了，我就应当全部保留吗？那么，我恳求你，在财富这件事上，不要这样说。因为接收者本应施予。正如胃留住食物不分配是一种恶习（因为这对全身有害），同样，富人将自己所拥有的留给自己也是一种恶习。因为这既毁灭了他们自己，也毁灭了别人。再者，眼睛接收所有的光，但它并不独自保留，而是照亮整个身体。因为只要它是眼睛，它的本性就不是将光只留给自己。再者，鼻孔能闻到香气，但它们并不全部留给自己，而是传递到大脑，以甜美的气味影响胃，并通过它们使整个人得到更新。双脚独自行走，但它们不仅移动自己，也带动整个身体。同样，无论你被托付了什么，不要只留给自己，因为你是在损害全身，而且首先损害你自己。\par

不仅是在肢体上可以看到这种情况；因为铁匠若选择不将自己的技艺传授给任何人，也会毁灭自己和所有其他工匠。同样，鞋匠、农夫、面包师，以及所有从事必要行业的人，若选择不将技艺的成果分享给任何人，不仅会毁灭别人，也会连同他们一起毁灭自己。\par

我为什么说\Quote{富有的}呢？因为穷人若也效法你们这些贪婪而富有的人的邪恶，也会极大地伤害你们，并很快使你们贫穷；甚至，他们会完全毁灭你们，假如他们因你们的匮乏而不愿分享自己所拥有的：例如，农夫分享他双手劳动的果实；水手分享他航海所得的收益；士兵分享他在战争中赢得的荣誉。\par

因此，如果别无他法，至少这件事应使你们羞愧，你们要效法他们的仁爱。你们不把自己的财富分施给任何人吗？那么你们也不应从别人那里得到任何东西；这样一来，世界就会天翻地覆。因为凡事有施予和接收，是众多福佑的原则：在种子、学生、技艺上都是如此。因为若有人想要把自己的技艺留给自己，他既颠覆了自己，也颠覆了整个事物秩序。农夫若把种子埋藏和保存在自己家里，就会导致严重的饥荒。同样，富人若如此对待自己的财富，必在穷人之前毁灭自己，将更严厉的地狱之火堆积在自己头上。\par

8. 因此，正如教师无论有多少学生，都向每人传授一些学识；同样，愿你们的财富成为众多受惠者。愿众人都说：\Quote{某人，他使这人脱离贫穷，使那人脱离危险。此人本会丧亡，若非借天主的恩宠而得到了你的庇护。你治好了这个人的病，使另一人摆脱诬告，又收留了作客旅的，给赤身露体的穿上衣服。}取之不尽的财富和众多宝藏也不及这样的赞誉。它们比你的金衣、马匹和奴隶更能吸引众人的目光。因为后者使人显得令人厌恶（希腊文：\OriginalTerm{}{\GreekText{φορτικόν}}，Saville校勘为\OriginalTerm{}{\GreekText{φορτικά}}），以致被人憎恨如公敌；而前者却宣告他是公共的父亲和恩人。而且，最重要的是，天主的恩宠在你所做的一切事上等候你。我的意思是，一人说：他帮助嫁了我的女儿；另一人说：他使我的儿子得以在众人中显要（\OriginalTerm{成为显要之人}{\GreekText{εἰς} \GreekText{ἄνδρας} \GreekText{ἐμφανῆναι}}）；另一人说：他终止了我的灾祸；另一人说：他救我脱离危险。这些话胜过金冠，一个人在他的城中应有无数人宣扬他的善行。这样的声音远比走在官长前面的传令官的声音更令人愉悦、更甜美；被称为救主、恩人、保护者（这些正是天主的尊称），而不是贪婪、骄傲、贪得无厌、吝啬。我恳求你们，我们不要贪爱这些称号中的任何一个，而应贪爱相反的。因为如果这些在地上说出的话能使人如此光彩照人；当它们被写在天上，天主在将来的日子宣告它们时，想想你将享有的何等的声望、何等的荣光！愿我们众人都有幸获得这一切，因我们的主耶稣基督的恩宠和慈爱；愿光荣、权能、尊崇归于祂及圣父和圣神，从今世直到永远，世世无尽。阿们。\par

\SourceRecord{Translated by Talbot W. Chambers.；Nicene and Post-Nicene Fathers, First Series；Vol. 12.；Edited by Philip Schaff.；Buffalo, NY: Christian Literature Publishing Co.,；1889.；<http://www.newadvent.org/fathers/220110.htm>.}

% structure: p0001 source=h1 target=section reason=page-title

\section{论格林多前书第十一篇讲道}

% structure: p0002 source=h2 target=subsection reason=relative-heading-rank; base=h2

\subsection{格林多前书 4:3-5}

\BibleQuote{但对我来说，我受你们审断，或受人间审断，是一件极小的事；我连我自己也不审断。因为我虽自觉良心无愧，但并非因此成义；那审断我的是主。}\par

与所有其他祸患一起，我不知道怎么竟然，人的本性染上了躁动刺探和不合时宜好奇的疾病，基督自己也责备过，说：\BibleRef{玛 7:1} \BibleQuote{你们不要判断，免得你们受判断。} 这种东西不像所有其他罪那样有快感，却只有惩罚和报复。因为虽然我们自己充满万般邪恶，眼睛里还扛着\Quote{大梁}，却变成邻人过错的严格审讯官，而邻人的过错根本不比\Quote{木屑}更大。因此，格林多教会就发生了这样的事：虔诚并蒙天主所爱的人因缺少学识而被嘲笑、被排斥；而另一些满身无数罪恶的人却因口才流利而被高度评价。于是，他们像坐在公堂上审案的人一样，鲁莽地投票说：\Quote{这个人配得上；那个人比另一个更好；这个人不如那个人；那个人比这个人好。} 他们不再为自己的恶行哀痛，反倒成了别人的法官；这样就再次点燃了惨烈纷争。\par

请注意，保禄如何明智地纠正他们，消除这种弊病。因为他曾说过：\Quote{此外，在管家中要求的是，人必须被证明是忠信的。} 这似乎给了他们评判和刺探每个人生活的空隙，从而加剧了派系情绪；唯恐对他们产生这样的影响，他便将他们从那种琐碎的争辩中拉开，说：\Quote{但对我来说，受你们审断是极小的事。} 他再次以自身为例继续讲论。\par

2. 但\Quote{但对我来说，受你们审断或受\OriginalTerm{人的日子}{\GreekText{ἡμέρας}}审断是极小的事？} 是什么意思？\Quote{我自认不配}他说，\Quote{受你们审断。} 为什么我说\Quote{受你们审断？} 我还要加上\Quote{受（\OriginalTerm{\GreekText{καὶ} \GreekText{τό}}{\GreekText{καὶ} \GreekText{τό}} [\OriginalTerm{\GreekText{τοῦ}}{\GreekText{τοῦ}}]）任何其他人审断。} 然而，不要让人谴责保禄傲慢，尽管他说没有人配对他判案。因为首先，他说这些话不是为自己，而是希望把别人从他们因格林多人所遭到的憎恶中拯救出来。其次，他不只将事情限于格林多人，他自己也放弃了这个审断权，说断定这样的事超出他的决定。至少他加上：\Quote{我不审断我自己。}\par

但除了上述所说，我们必须探究说出这些言论的基础。因为他很知道在许多情况下如何以高昂的精神说话：这不是出于骄傲或傲慢，而是出于某种\OriginalTerm{卓越的管理}{\GreekText{οἰκονομίας} \GreekText{ἀρίστης}}，因为在目前的情况下，他这样说也不是抬高自己，而是压下别人的气焰，热切地寻求赋予圣人们应有的荣耀。为了证明他是极其谦卑的人，请听他说的话，他引用敌人在这方面的见证：\BibleQuote{他的身体软弱，言语毫无价值。} \BibleRef{格后 10:10} 又引：\BibleQuote{最后，也如同是流产儿，祂也显现给我。} \BibleRef{格后 15:8} 然而，看这谦卑的人，当时机召唤他时，他将门徒的精神提升到何等的高度，不是教导骄傲，而是灌输健康的勇气。因为与这些人谈论时他说：\BibleQuote{你们将要审判世界，难道不配审判最小的事吗？} \BibleRef{格前 6:2} 因为基督徒应当远离傲慢，也当远离谄媚和卑劣的精神。因此，若有人说：\Quote{我把金钱视为无物，世上一切都如影、如梦、如儿童游戏。} 我们绝不能指责他傲慢；因为这样我们就得指控撒罗满本人傲慢，他\OriginalTerm{严峻地}{\GreekText{φιλοσοφοῦντα}}谈论这些事，说：\BibleQuote{虚而又虚，万事皆虚。} \BibleRef{训 1:2} 但天主禁止我们把严谨的生活规范称为傲慢。因此轻视这些事不是傲慢，而是灵魂的伟大；尽管我们看到国王、统治者、权贵们重视它们。但许多穷人，过着严谨的生活，轻视它们；我们因此不称他为傲慢，而是高尚：正如另一方面，如果有人极度沉迷于它们，我们不称之为心地谦卑、有节制，而是软弱、胆怯、卑贱。因为同样，如果儿子蔑视父亲应有的追求而喜欢奴才习气，我们不应称赞他谦卑，而应谴责他卑贱和奴性。我们应当欣赏的是，他轻视那些卑下的事物，而重视来自父亲的东西。因为认为自己比同工更好是傲慢：但对事物做出正确判断不是来自自夸，而是来自生活的严谨。\par

为此，保禄也不是为了抬高自己，而是为了谦卑别人，压制那些越位出头的人，并劝他们谦虚，说：\Quote{但对我来说，受你们审断或受人的日子审断是极小的事。} 注意他如何也安抚另一方。因为无论谁听说他轻视一切人，不屑受任何人审断，都不会再感到痛苦，好像自己是唯一被排除的人。如果他只说：\Quote{受你们审断，} 就不再言语，这就足以刺痛他们，仿佛被藐视。但现在，他加了\Quote{也不受人的日子审断，} 减轻了打击；使他们有同伴分担这种轻视。不仅如此，他甚至再次缓和这点，说：\Quote{我连自己也不审断。} 注意这说法多么没有傲慢：因为连他自己，他说，也不足以做如此严格的判断。\par

3. 然后，因为这句话也像是极大赞扬自己的人说的，他对此也加以纠正，说：\BibleQuote{但我并不因此就成义。}那么怎样？我们难道不该判断自己和我们自己的过错吗？当然应该：我们在犯罪时极需要这样做。但保禄说的不是这个，他说：\BibleQuote{我什么也不知道，}他说，\BibleQuote{反对我自己。}那么，当他\BibleQuote{不知道自己有什么过错}时，他要判断什么过错呢？然而，他说：\BibleQuote{他并未成义。} \BibleRef{格前 6:3}我们这些良心布满千万伤痕，自觉毫无善功，反而恰恰相反的人，又能说什么呢？\par

怎么可能会这样：如果他不知道自己有什么过错，他却没有成义？因为他有可能犯了某些罪，却不知道它们是罪。由此你该估量将来的审判将何等严厉。你瞧，他并非认为自己无可指责才说不该由他们审判他，而是为了堵住那些不合理地这样做之人的口。至少在其他地方，即使人的罪昭然若揭，他也不允许别人判断，因为时机需要。他说：\BibleQuote{你为什么判断你的弟兄？} \BibleRef{罗 14:10} 或 \BibleQuote{你，为什么轻看你的弟兄？} 因为你没有被命令，人啊，去判断别人，而是要省察你自己的行为。那么，你为什么抢占主的职分？审判是他的，不是你的。\par

为此，他又说：\BibleQuote{因此，在时候未到以前，你们什么也不要判断，只等主来；他要照亮隐藏的黑暗事物，并显露心中的计谋；那时，各人要从天主那里得到称赞。}那么怎样？我们的教师不该这样做吗？在公开承认的罪上，且在有适当机会时，甚至带着痛苦和内心的烦恼，这样做是对的：不像他们当时那样行事，出于虚荣和傲慢。因为在此处，他并非在谈论那些所有人都承认是罪的事，而是在谈论偏爱某一个人胜过另一个人，以及比较生活方式。因为这些事，只有那位要审判我们隐秘行为的主才能准确判断，哪些应受更大的惩罚或尊荣，哪些应受较小的。而我们都凭眼所见做这一切。他说：\Quote{如果在我自己的错误中，我什么也不清楚，我怎么能配得上去判断别人呢？如果我不准确知道自己的情况，我又怎能判断别人的境况呢？}如果保禄都这样感受，何况我们呢？因为（继续下去）他说这些话，并非要显明自己无过，而是为了说明，即使他们中间有某个人是这种毫无过犯的人，他也不配判断别人的生命；而且，如果他虽然自觉一无所是却宣称自己有罪，那么那些内心有千万罪恶的人更是如此。\par

4. 你看，他这样堵住了那些做出这样判断之人的口之后，接着便带着强烈的感情，几乎要爆发出来，冲向那个污秽的人。正如当风暴来临时，一些乌云先带着黑暗奔跑；之后，当雷声大作，使整个天空变成一片乌云时，大雨突然倾泻到地上：当时发生的事也是如此。因为他虽然可能十分愤慨地处理那行奸淫的人，但他并没有这样做；而是先用可怕的话语压制了那人膨胀的骄傲，因为事实上，所发生的是双重的罪，即奸淫，以及比奸淫更糟的，即对所犯之罪不感到忧伤。因为他哀叹的不是那罪本身，而是犯罪者，那还没有悔改的人。因此，他说：\BibleQuote{我要为许多人哀哭，}他不只是说\BibleQuote{那些从前犯了罪的人，}而是加上\BibleQuote{那些没有悔改他们所行污秽和邪淫的人。} \BibleRef{格后 12:21}因为犯罪后行了悔改的人，不是可悲的对象，而是值得庆贺的，因为他已转入义人的队伍。因为，\BibleRef{依 43:26}\BibleQuote{你首先要陈明你的过犯，以便你能成义；}但如果一个人犯罪后毫无羞耻，那么他跌倒固然可怜，但躺卧在跌倒之处更可怜。\par

如果犯罪后不悔改是一个严重的过失；因罪而骄傲，该受怎样的惩罚呢？因为如果一个人因自己的善行而得意是不洁的，那么因为自己的罪而有这种感受的人，将得到怎样的宽恕呢？\par

既然那行淫者是这样的人，并因他的罪而使自己的心变得如此顽固和倔强，他当然以打击他的骄傲开始。他既没有首先提出指控，以免他因被单独指责于众人之前而变得刚硬；也没有等到最后，免得他认为与他相关的事只是附带提及。而是先通过向其他人直言不讳地说话，在他心中激起极大的恐惧，然后，直到那时，才转而针对他，在斥责他人的过程中预先给了那人的顽固执拗一份。\par

因为同一句话，即\BibleQuote{我什么也不知道反对我自己，但我并不因此成义，}以及这句\BibleQuote{审判我的是主，他要照亮隐藏的黑暗事物，并显露心中的计谋，}不仅轻轻触及那个人，也触及那些与他同谋并轻视圣人的人。他说：\Quote{因为，如果某些人外表看来是品德高尚、令人钦佩的人，那又怎样？那位审判者不只是辨明外表，还要照亮一切秘密。}\par

5. 你看见，为了两个原因，或者更确切地说，为了三个原因，正确的判断不属于我们。其一，因为即使我们自觉毫无所愧，我们仍然需要有人严格地指责我们的罪过。其二，因为所行的大部分事情都逃脱我们的注意而被隐藏了。其三，除此之外，许多别人所行的事情，在我们看来似乎是好的，但它们并非出自正确的心思。那么，你为何说这个人没有犯过罪？那个人比另一个人更好？因为我们不可作此论断，甚至对那自己不觉有愧的人也不可。因为洞悉秘密的那一位，祂才确实地审判。例如，看哪；就我而言，我自觉毫无所愧：然而即便如此，我仍未成义，也就是说，我并未免除应交的账目和应回答的控告。因为祂不是说：\Quote{我不列于义人之中；}而是说\Quote{我并未从罪中得洁净。}因为在别处祂也说，\BibleRef{罗 6:7} \BibleQuote{已死的人是脱离了罪，}也就是说，\Quote{被释放了。}\par

再者，我们行许多善事，但并非出自正确的心思。因为我们也如此称赞许多人，并非出于使他们显扬的愿望，而是为了藉着他们来伤害别人。所行的事本身是对的——那行善者受称赞；但意图是败坏的：因为它是出于一种撒殚般的动机。因为此人时常这样做，并非与弟兄一同喜乐，而是想要伤害另一方。\par

再者，一个人犯了大错；另一个人想要取而代之，便说他什么也没做，并藉着诉诸人性的普通软弱来安慰他，似乎在安慰他的错误。但此人时常这样做，并非出于同情之心，而是为了使他在过错中更加安逸。\par

再者，一个人时常斥责，与其说是为了责备和劝诫，不如说是为了公开地\OriginalTerm{显示和夸大}{\GreekText{ἐκπομπεῦσαι} \GreekText{καὶ} \GreekText{ἐκτραγωδῆσαι}}邻人的罪过。然而，我们的计谋本身是人所不知的，但\BibleRef{罗 8:27} \BibleQuote{那鉴察人心的}完全知道它们；并且祂将在那时使这一切显现。因此祂说：\Quote{祂要显明黑暗中的隐密事，并显明心中的谋略。}\par

6. 那么，既然即使在我们\Quote{自觉毫无所愧}之处，我们也不能免于控告，并且当我们行任何善事，但并非出于正确的心思时，我们仍当受惩罚；那么，请思量人们在判断上是何等受骗。因为这一切事情，人不能触及，唯独那不眠之眼才能触及；尽管我们可以欺骗人，我们的诡辩却永远不能胜过祂。那么，不要说：黑暗围绕着我，墙壁围绕着我；谁看见我呢？因为那位亲自塑造了我们心肠的祂，知道一切。\BibleRef{咏 138:12} \BibleQuote{黑暗对祂不是黑暗的。}然而，正在犯罪的人确实会说：\Quote{黑暗围绕着我，墙壁围绕着我；}因为倘若他的心中没有黑暗，他就不可能驱逐对天主的敬畏而任意妄为。除非那主宰的原则首先被黑暗笼罩，否则罪过之进入而无恐惧是不可能的。那么，不要说：谁看见我呢？因为有一位\BibleRef{希 4:12} \BibleQuote{刺入甚至灵魂和神、关节和骨髓；}但你看不见自己，也不能刺透那云层；反而如同四围都有墙壁环绕你，你无力仰望上天。\par

因为无论你愿意犯什么罪，首先让我们查考，你将会看到罪的产生正是如此。因为如同强盗和挖墙的人在想要偷窃任何贵重物品时，先吹灭蜡烛，然后动手；同样，在犯罪之人身上，人的乖张推论也是如此。因为在我们内，也实在有一种光，即理性的光，永远燃烧。但如果罪恶之灵带着强风急来，扑灭那火焰，它便立刻使灵魂黑暗，并胜过它，随即掠取其中所积蓄的一切。因为当灵魂被不洁的欲望俘虏时，那欲望如同云雾遮蔽肉体的眼睛，同样，那欲望截断了心灵的远见，不容它看见任何远处，无论是悬崖、地狱还是恐惧；反而，从那以后，他有了那欺骗作为暴君在他之上，便容易为罪所胜；在他眼前竖立起一座无窗的墙壁，不容正义的光芒照入心灵，肉欲的荒谬幻想如同壁垒四面围困他。从那时起，不贞洁的女人处处与他相遇：出现在他的眼前、他的心意、他的思想之前。正如盲人，即使站在正午太阳的正下方，也因眼睛紧闭而接不到光；同样，这些人，尽管万般救恩的教义从四面八方传入他们耳中，但灵魂已被此激情占据，便堵塞耳朵不听这些言论。那些经历过的人都深知这一点。但愿天主不让你们从实际经验中知道这事。\par

7. 不仅这罪有如此效果，一切偏离正轨的情欲亦然。你若愿意，请将论题从不贞洁的妇人转到金钱，我们也要在此看到同样浓厚而持续的黑暗。因为在前一种情况，所爱之物为单一且固着于一地，情感不甚猛烈；但金钱处处显现——在银匠铺、酒馆、金匠作坊、富人的家中——情欲便刮起猛烈的风暴。当仆人们在广场上招摇过市，饰有黄金马具的马匹、穿着昂贵衣服的人，被患此病态者看见时，笼罩他的黑暗就变得深沉。何需提及房屋和银匠铺？在我看来，这样的人即使只在图画或影像中看见财富，也会抽搐、狂乱、咆哮。因此，黑暗从四面八方聚集在他们周围。倘若他们偶然看见君王的画像，他们不会赞赏宝石之美，也不欣赏黄金或紫袍，而是憔悴。正如前面提到的可怜恋人，即使只看见所爱妇人的画像，也会依恋那无生命之物；同样，此人看见财富的无生命影像，也会更强烈地受其影响，因为他被更残暴的情欲所束缚。从此他要么待在家中，要么冒险进入广场，然后带着无数伤害回家。因为有许多事物刺痛他的眼睛。正如前者只看见那妇人，后者也匆匆绕过穷人及一切事物，以免得到丝毫缓解。但他定睛观看富人，藉此将火焰猛力引入自己的灵魂。因为这火焰悲惨地吞噬着陷入其中的人；纵然没有地狱的威胁，也没有惩罚，这种状态本身就是惩罚：不断地受折磨，永远无法找到疾病的终结。\par

8. 这些事本身已足以劝我们逃避此病。但没有比轻率更大的恶，它使人固着于带来悲伤而无益的事物。因此，我劝你们在情欲之初将其斩断：因为正如热病初发时并不使患者剧烈口渴，但随其加剧及火焰升高，便从此引发不治之渴；纵使让人喝满饮料，也不能扑灭炉火，反而使其更炽烈。这情欲亦然：除非当它初次侵入我们的灵魂时我们阻止它并关闭门户，一旦进入，便使接纳者的疾病成为不治。因为好事与坏事，在我们内停留越久，就越强大。\par

在所有其他事物中，任何人也能看见这事发生。正如刚种下的植物容易拔起，但长久扎根后则否，那时需用杠杆大力。新砌的建筑物易被推倒，但一旦牢固，便给试图拆毁者带来极大麻烦。长期在某处栖息的野兽也难以被驱走。\par

因此，我劝尚未被此情欲掌控的人不要被俘虏。因为预防陷入比陷入后脱身更容易。\par

9. 但对于那些已被抓住并击垮的人，如果他们愿意把自己交托给医治的圣言，我靠着天主的恩宠，许诺他们获得救恩的极大希望。因为如果他们考虑那些曾受苦、陷入此病而痊愈的人，他们就会对除去疾病抱有希望。那么，有谁曾陷入此病而轻易摆脱呢？众所周知的匝凯。有谁比税吏更贪财？但他一下子成为生活严谨的人（\OriginalTerm{有德之人}{\GreekText{Φιλόσοφος}}），熄灭了那全部烈焰。玛窦同样：他也是税吏，生活在不断的掠夺中。但他也一下子摆脱了祸害，消了渴，并追随属灵的收获。因此，想到这些和类似的人，你也不要绝望。因为若你愿意，你很快就能痊愈。若你愿意，我们将按医生的规矩，精确地开出你当行的药方。\par

因此，在一切事之先，必须正确对待这一点：我们永不灰心，也不对我们得救绝望。其次，我们不仅要看那些行善之人的榜样，也要看那些坚持犯罪之人的苦楚。因为我们既想到匝凯和玛窦，也当计算犹达斯、革哈齐、阿干（见：若苏厄书第七章）、阿哈布、阿纳尼雅和撒斐辣，为藉前者除去一切绝望，藉后者除去一切怠惰；使灵魂不轻视所提出的疗法。我们要教导他们自己说出那天犹太人走近伯多禄时所说的话：\BibleQuote{我们该做什么才能得救？}\BibleRef{宗 2:37}并让他们听听他们该做什么。\par

10. 那么我们必须做什么？我们当知道这些事物何等无价值，财富乃是一个逃跑的奴隶，冷酷无情，给它的拥有者带来无数的灾祸。让我们不断用这样的话像符咒一样在他们耳边回响。如同医生为安抚那些要求冷水的病人，说他们愿意给，但以泉水、容器、合适的时间以及许多诸如此类的借口搪塞（因为若他们立刻拒绝，会使病人狂乱），同样，我们也当如此对待贪爱金钱的人。当他们说\Quote{我们渴望致富}，我们不要立即说财富是邪恶的；而是赞同，并说我们也渴望它，但要适时，是的，真正的财富，是的，那有永不衰残的喜乐的财富，是的，那为你自己积聚、而非为他人（且往往为我们的仇敌）积聚的财富。让我们提出真智慧的教训，说：\Quote{我们不禁止财富，而禁止不义之财。}因为人可以合法地致富，但不可贪婪，不可掠夺、强暴，并遭众人恶评。先用这些论调使他们缓和下来，暂且不谈地狱。因为病人还受不了这样的话。因此，让我们用这个世界的一切论据来讨论这些事，并说：你为什么选择通过贪婪致富呢？难道是为了金银储存给他人，而为你留下咒骂和无数指责？难道是为了让你所欺骗的人因缺乏生活必需品而痛苦，自怨自艾，招致千万人对你的谴责，并在黄昏时分在市场上徘徊，在小巷中逢人便诉，极度困惑，甚至不知道当晚能否安身？因为他如何能入睡？腹中饥饿，饥荒无情地折磨，且常常是寒风刺骨，雨水淋身；而你，沐浴后回家，穿着柔软的衣服，心中快乐，喜气洋洋，急匆匆赶向一个预备好的丰盛宴席；而他，被寒冷和饥饿驱赶，在市场上到处流浪，弯腰驼背，伸出双手；他甚至没有勇气不颤抖地向你这样酒足饭饱、一心安闲的人恳求食物；反而常常带着侮辱离去。因此，当你回到家，躺在床上，四周灯火通明，餐桌已备好并丰盛时，那时要记起那个可怜的穷人，像小巷里的狗一样在黑暗和泥泞中流浪；只不过，他常常不得不离开那里，不是去房屋、妻子或床铺，而是去一个草垫子；正如我们看见狗整夜吠叫一样。而你，若看见屋顶滴下一小滴水，就搅得全家不安，呼唤你的奴隶，扰乱一切；而他，躺在破布、稻草和泥土中，忍受全部寒冷。\par

什么野兽不会被这些事软化？有谁如此野蛮残忍，这些事还不能使他温和呢？然而，有些人竟达到了如此残忍的地步，甚至说他们活该受这些苦。是，正当他们应该同情、哭泣、帮助减轻人们的灾难时，他们反而用野蛮凶残的谴责来对待他们。我很想问这些人：\Quote{请告诉我，他们为什么活该受这些苦？是因为他们想要吃饱而不饿死吗？}\par

不，你会回答；而是因为他们想在懒惰中得到饱食。而你，难道你不沉迷于懒惰吗？我说什么？你难道不是常常从事一项比任何懒惰都更痛苦的职业，即攫取、压迫、贪婪吗？你若是也这样懒惰倒更好；因为这样懒惰总比贪婪好。但现在你甚至践踏他人的灾难，不仅懒惰，不仅从事比懒惰更坏的职业，而且还诽谤那些在悲惨中度日的人。\par

让我们进一步向他们叙述他人的灾祸：意外的丧亲、监狱中的居民、在法庭上被撕裂的人、为生命战栗的人；妇女意想不到的寡居；富人突发的破产：用这些来软化他们的心。因为通过我们对他人遭遇的叙述，我们必会诱使他们也惧怕这些祸患临到自己身上。因为当他们听到某个贪婪攫取之人的儿子，或（\OriginalTerm{或某人之子}{\GreekText{ἠ} \GreekText{τοῦ} \GreekText{δεῖνος}}，原文作\OriginalTerm{是}{\GreekText{ἦν}}；\OriginalTerm{某人}{\GreekText{τοῦ} \GreekText{δεῖνος}}）某个做过许多暴虐行为的某人的妻子，在丈夫死后遭受无尽的痛苦；受害者攻击妻子和儿女，从四面八方对他家发起全面战争；即使一个人是最无感的，但预期自己也会遭受同样的事，并为自己的人担心他们遭遇相同的命运，他会变得更为节制。如今我们发现生活中有许多这样的故事，我们不会缺乏这种纠正的方法。\par

但当我们说这些事时，不要以劝告或建议的方式说，免得我们的谈话变得太令人厌倦；而是要按叙述的顺序并借着与其他事物的关联，在每个场合都转向那种类型的谈话，并不断引导他们进入这类故事，允许他们只谈论以下这些话题：某人的辉煌著名宅邸如何倒塌；它如何如此彻底荒废，以至于里面的一切都落入了他人之手；每日有多少关于同一财产的审判，多么大的骚动；那个人的\OriginalTerm{亲属}{\GreekText{οἴκεται}}（可能为\OriginalTerm{亲属}{\GreekText{οἰκεῖοι}}）有多少要么作为乞丐，要么作为监狱居民死去。\par

这一切，让我们怀着对亡者的怜悯来说，并轻看现世的事物；为的是藉着恐惧与怜悯，我们能软化那残酷的心。当我们看见人们因这些叙述而畏缩时，那时，而非在此之前，我们才将地狱的道理介绍给他们注意，并非为恐吓这些人，而是出于对他人的同情。让我们说：但何必提现世的事呢？因为我们的关切远非止于此；还有更严厉的惩罚等待着所有这样的人：甚至是一条火河，一条有毒的虫，无边的黑暗，以及不止息的折磨。如果藉着这样的言辞我们成功地对他们施加影响，我们将纠正我们自己和他们，并迅速克服我们的软弱。\par

在那一天，我们将有天主来赞美我们：正如\Person{保禄}所说：\Quote{那时，各人要从天主那里获得称赞。}因为来自人的称赞，既是短暂的，有时也并非出于善意。但来自天主的称赞，既长存不变，又明澈发光。因为当那在万物受造之前就已知道一切，且毫无私情偏心的那一位给予赞美时，我们德行的证据甚至是无可置疑的。\par

因此，知道这些事，让我们行事为要得天主的称赞，并获享最大的福分；愿天主恩赐我们众人获得这一切，因我们的主耶稣基督的恩宠与慈爱，偕同圣父及圣神，享受荣耀、权能、尊崇，从今时直到永远，及于无穷之世。阿们。\par

\SourceRecord{Translated by Talbot W. Chambers.；Nicene and Post-Nicene Fathers, First Series；Vol. 12.；Edited by Philip Schaff.；Buffalo, NY: Christian Literature Publishing Co.,；1889.；<http://www.newadvent.org/fathers/220111.htm>.}

% structure: p0001 source=h1 target=section reason=page-title

\section{论《格林多前书》第十二讲}

% structure: p0002 source=h2 target=subsection reason=relative-heading-rank; base=h2

\subsection{\BibleBook{格林多}{前書} 4:6}

\BibleQuote{弟兄們，我為你們的緣故，把這些事轉用在自己和\Person{阿颇罗}身上，叫你們在我們身上學到不可超過聖經所記的。}\par

當需要如此嚴厲的話語時，他沒有拉開帷幔，而是繼續爭論，彷彿他自己就是受責備的人；這樣，被責備者的尊嚴不致於與責備者對抗，也不給他們留下因指控而發怒的餘地。但當需要更溫和的處理時，他就剝掉那層外衣，摘下假面具，顯出那些被\Person{保禄}和\Person{阿颇罗}的名字所隱藏的人。為此緣故，他說：\Quote{弟兄們，我為你們的緣故，把這些事藉著比喻轉用在自己和\Person{阿颇罗}身上。}\par

這正如病人的情況：當孩子身體不適，踢開醫生遞來的食物時，看護者會叫來父親或監護人，請他們從醫生手中接過食物拿給他，這樣孩子因敬畏他們而吃下並安靜下來。同樣，\Person{保禄}打算責備他們對待某些人的態度——他認為有些人受到傷害，有些人則受到過分尊崇——他沒有直接點出那些人，而是以自己和\Person{阿颇罗}的名義進行論證，好使他們因敬畏這兩位而接受他的治療方式。一旦治療被接受，他立刻表明他是為誰而這樣表述。\par

這不是偽善，而是\OriginalTerm{屈就}{\GreekText{συγκατάβασις}}和\OriginalTerm{管理}{\GreekText{οἰκονομία}}。因為如果他公開說：\Quote{至於你們，你們所判斷的那些人是聖徒，值得一切讚賞。}他們可能會不高興，甚至\OriginalTerm{甚至會跳開}{\GreekText{κἂν} \GreekText{ἀπεπήδησαν}}。但如今他說：\Quote{但你們判斷我，我卻認為是極小的事。}以及\Quote{\Person{保禄}是誰？\Person{阿颇罗}是誰？}這就使他的話容易被接受。\par

如果你留意，這正是他這裡說\Quote{我為你們的緣故，把這些事藉著比喻轉用在自己身上，叫你們在我們身上學到不可超過聖經所記}的原因，意思是：如果他把論證直接用在他們身上，他們就不會學到所需要的一切，也不會因所說的話而煩惱接受糾正。但實際上，由於敬畏\Person{保禄}，他們很好地承受了責備。\par

2. 但\Quote{不可超過聖經所記}是什麼意思呢？經上記著：\BibleRef{瑪 7:3}\BibleQuote{為什麼你看見你兄弟眼中的木屑，卻不注意自己眼中的大樑呢？}以及\BibleQuote{你們不要判斷，免得你們受判斷。}因為如果我們是一體、彼此相連，就不該互相攻擊。因為經上說：\BibleQuote{自謙者必被高舉。}並且\BibleRef{瑪 20:26}\BibleQuote{誰若願意在眾人中為首，就當作眾人的僕役。}這些就是\Quote{所記的}事。\par

\BibleQuote{免得你們有人自高自大，偏袒這個反對那個。}再者，他離開了師傅們，轉而責備門徒們。因為是他們使前者驕傲自大。\par

此外，那些領袖們不會平靜地接受那種話語，因為他們貪求表面的榮耀；他們甚至被這種熱情弄瞎了眼。而門徒們，既然自己不享受榮耀的果實，而是為他人謀求榮耀，就會更溫和地忍受責備，而且比領袖們更有能力消除這種病症。\par

看來，因別人的事而自高自大，即使一個人對自己的事沒有這種感覺，這也是\Quote{自高自大}的症狀。因為一個人若是因別人的財富而驕傲，那是由於傲慢；同樣，因別人的榮耀而驕傲也是如此。\par

他稱之為\Quote{自高自大}，很是恰當。因為當某一個肢體高過其他肢體時，這無非是發炎和疾病；因為除非腫脹發生，一個肢體不會比另一個高。（在我們的語言中稱為\Quote{驕肉}。）同樣，在教會的身體中，凡是發炎並自高自大的人，必定是有病的；因為他腫脹得超過了其他肢體的比例。這正是我們所說的\Quote{腫脹}。在身體中也是如此，當某種虛假的、有害的體液積聚，取代了正常的營養時，便會發生這種情況。傲慢也是如此：不屬於我們的觀念侵入我們。請注意，他說\Quote{不要自高自大}是多麼準確：因為自高自大的人因充滿敗壞的體液而有一種精神上的腫瘤。\par

然而，他說這些話，並非禁止一切安撫，而是禁止那種導致危害的安撫。\Quote{你要服事這個人或那個人嗎？我不禁止你；但不要做損害別人的事。}因為師傅們賜給我們，不是為了使我們彼此對立，而是為了使我們都互相合一。正如將軍被派在軍隊之上，是為了把分散的人合成一個身體。但如果他要瓦解軍隊，他就成了敵人而不是將軍。\par

% structure: p0014 source=h2 target=subsection reason=relative-heading-rank; base=h2

\subsection{\BibleBook{格林多}{前書} 4:7}

3. \BibleQuote{誰使你與眾不同呢？你有什麼不是領受的呢？}\par

從這一節開始，他離開了被治理的人，轉向治理者。他的話大意是：從哪裡看出你配得讚賞？難道有過任何審判？任何調查？任何考驗？任何嚴格的測試？不，你不能這樣說：即使人們投票，他們的判斷也不正直。但假設你真的配得讚賞，確實擁有恩賜，而且人的判斷也不腐敗：即便如此，也不該驕傲；因為你所有的一切不是出自你自己，而是從天主領受的。那麼，你為什麼假裝擁有你沒有的東西呢？你會說：\Quote{你擁有它。}別人也與你一同擁有。那麼，你是藉著領受而擁有它：不只是這一件事或那一件事，而是所有你擁有的一切。\par

因为这些卓绝之处并不属于你，而是属于天主的恩宠。无论你提及信德，它来自于祂的召叫；或你谈论罪的赦免、神恩、教导的话语、奇迹，你这一切都从那里领受。那么，告诉我，你有什么不是领受的，而是凭自己成就的？你无话可说。好吧：你领受了，这就使你自高自大吗？不，它应当使你谦卑自守。因为所赐给你的不是你的，而是施予者的。若你领受了，你从他领受。若你从他领受，你所领受的便不是你的；若你只是领受了非你所有的，你为何如同拥有自己之物而自高自大？因此他又加上：\Quote{若是你领受了，为何你夸耀，好像没有领受呢？}\par

4. 你瞧，他通过让步论证（\OriginalTerm{归纳论证}{\GreekText{κατὰ} \GreekText{συνδρομὴν}}）完善了他的论点后，指出他们仍有欠缺，并且不在少数。他说：\Quote{首先，即使你领受了一切，也不该夸耀，因为没有任何事物是你自己的；但实际情况是，你缺少许多东西。}开头他只暗示了这一点，说：\Quote{我不能把你们当作属神的人讲话}；以及\Quote{我曾决定，在你们中间不知道别的，只知道耶稣基督，这被钉十字架的基督}。但在此处他以使他们羞愧的方式说道：\par

% structure: p0019 source=h2 target=subsection reason=relative-heading-rank; base=h2

\subsection{\BibleRef{格前 4:8}}

\Quote{你们已经饱满了，已经富足了：}意思是，你们从此不需要什么了；你们成了成全的；你们达到了顶峰；按你们的想法，你们不需要任何人，无论是宗徒还是教师。\par

\Quote{你们已经饱满了。}他特意说\Quote{已经}，从时间上指出他们论断的不可信和他们对自身的无理看法。因此他是以嘲讽的态度对他们说：\Quote{这么快你们就到了终点}，这在时间上是不可能的：因为所有更成全的事物都在未来长久等待；但一点就\Quote{饱满}，证明灵魂软弱；从一点就自以为\Quote{富足}，证明灵魂病态可怜。因为虔敬是一种永不满足的事；从刚开始就以为自己已获得全部，是幼稚的想法；尚未进入事物序曲的人却自以为已掌握结局，这是高傲。\par

然后通过接下来的话，他更加使他们难为情；在说了\Quote{你们已经饱满了}之后，他又加上：\Quote{你们已经富足了，你们没有我们已为王了；我恨不得你们果真为王，叫我们也能与你们一同为王。}这话极其严厉；因此它被放在最后，在他大量责备之后引入。因为我们的劝告受人尊重并容易被接受，是在我们谴责之后引入\OriginalTerm{使人羞愧的话}{\GreekText{τὰ} \GreekText{ἐυτρεπτικὰ} \GreekText{ῤήματα}}之时。这足以抑制甚至无耻的灵魂，比直接指责更尖锐地打击它，并纠正因指控而产生的苦涩和僵硬感。可以肯定，这些诉诸羞耻感的论点最令人钦佩之处在于，它们具有两种相反的优势。一方面，比公开抨击更深地切入；另一方面，它使受斥责之人更耐心地承受那更重的刺痛。\par

5. \Quote{你们没有我们已为王了。}这里有很大的力量，既涉及教师，也涉及门徒；并且也指出了他们对自己的无知（\OriginalTerm{无自知之明}{\GreekText{τὸ} \GreekText{ἀσυνείδητον}}）以及他们的极大轻率。因为他所说的意思是：\Quote{在劳苦中，确实，我们和你们共同承担一切，但在赏报和冠冕上，你们领先。我这样说并非出于恼怒：}因此他又加上：\Quote{我恨不得你们果真为王：}然后，为避免显得有讽刺，他加上：\Quote{叫我们也能与你们一同为王；}因为，他说，我们也应分享（\OriginalTerm{分享}{\GreekText{ἐπιτύχοιμεν}}，Reg.抄本作\OriginalTerm{分享}{\GreekText{ἐπιτύχωμεν}}）这些福分。你看到他如何同时在自身中展现了他的严厉、他对他们的关怀以及他的克己精神？你看到他如何打击他们的骄傲？\par

% structure: p0024 source=h2 target=subsection reason=relative-heading-rank; base=h2

\subsection{\BibleRef{格前 4:9}}

\Quote{我以为天主把我们宗徒列在最后，好像定了死罪的人。}\par

他说\Quote{我们}时再次有很深的含义和严厉；他甚至不满足于此，还加上了他们的身份，猛烈地打击他们：\Quote{我们宗徒}，就是我们这些忍受无数患难的人；我们这些播种虔敬之言的人；我们这些引导你们走向这严格生活规范的人。这些\Quote{祂列在最后，好像定了死罪的人}，即是说，如同被定罪的人。因为他之前说了\Quote{叫我们也能与你们一同为王}，用那句话缓和了严厉，以免使他们沮丧；现在他更加严肃地重新拾起，说：\Quote{我以为天主把我们宗徒列在最后，好像定了死罪的人。}\Quote{因为据我所见，}他说，\Quote{并且从你们的话可知，所有人类中最卑微、明确被定罪的人，正是我们这些被置于持续受苦中的人。但你们已经在幻想中有了王国、尊荣和巨大的赏报。}为了将他们的推理引向更荒谬的地步，并显示出其极端的不可信，他不仅说\Quote{我们是\InnerQuote{最后}}，还说\Quote{天主使我们成为最后}；他又不止于说\Quote{最后}，还加了\Quote{定了死罪的人}：为的是使即使完全无知的人也感到这话极不可信，他这话是出于恼怒并强烈地使他们羞愧。\par

也要注意保禄的明智。当合适的时候，他用以高抬自己、显示自己的尊荣、使自己显大的话题，现在他正是用这些来使他们羞愧，称自己为\Quote{被判罪的人}。在此处，他所说的\Quote{被判死刑}的意思是\Quote{被判罪}，而且当受万死。\par

\BibleQuote{我们成了一台戏，给世界、天使和众人观看。}\par

\Quote{我们成了一台戏，给世界看}是什么意思？他说：\Quote{我们遭受这些事，不是在一个角落，也不是在世界的一小部分，而是处处、在众人面前。}但\Quote{给天使看}是什么意思？人可能\Quote{成为给众人的一台戏}，但不是给天使，除非所做的事是超凡的。然而我们的搏斗是如此值得天使的观看。请看，他是如何从那些贬低自己的事上再次显出伟大；又从那些他们引以为傲的事上，显示出他们的渺小。因为被认为愚昧比显得智慧更可鄙；软弱比强壮更可鄙；不被尊重比光荣显赫更可鄙；而且他即将把前一种称呼加在他们身上，而他自己却接受了后一种；他表明后者优于前者，至少因为后者的缘故，他使群众——不单是人，而且连天使本身——都凝望着他们。因为我们不仅与人类搏斗，也与无形的力量搏斗。因此，\OriginalTerm{一个巨大的剧场}{\GreekText{μέγα} \GreekText{θέατρον} \GreekText{κάθηται}}也设立了。\par

% structure: p0030 source=h2 target=subsection reason=relative-heading-rank; base=h2

\subsection{\BibleRef{格前 4:10}}

\BibleQuote{我们为基督的缘故成了愚拙的，你们在基督里却是聪明的。}\par

同样，这句话他也说得使他们羞愧；暗示这些对立面不可能一致，相距如此遥远的事不能同时并存。他说：\Quote{怎能这样呢？你们聪明，我们却在有关基督的事上愚拙？}意思是：一类人被打、被藐视、被羞辱，被视为无物；另一类人却享有尊荣，被许多人视为聪明谨慎的人；这给他机会这样说：仿佛在说：\Quote{怎么能这样呢？宣讲这些事的人，却实际上被看作是在行相反的事？}\par

\BibleQuote{我们软弱，你们却强壮。}这就是说，我们被驱赶、被迫害；你们却享有平安，备受服侍；然而福音的本性并不容忍这种情况。\par

\BibleQuote{我们被人藐视，你们却受尊敬。}在这里，他针对那些显贵和倚仗外在优势而自傲的人。\par

\BibleQuote{直到现今，我们还是又饥、又渴、又赤身露体、又挨打、又无定所，并且劳苦，亲手作工。}这就是说：\Quote{我说的不是古老的故事，而是当前的时间本身为我作证：我们不关心人间的事，也不关心任何外在的排场；我们只仰望天主。}这也是我们在各处都需要实行的。因为不仅天使在观看，而且比他们更甚的是，那主持这演出的一位也在观看。\par

所以，我们不要渴望任何其他人为我们鼓掌。因为这就是侮辱祂；绕过祂，仿佛祂不足以欣赏我们，我们就径自奔向同侪。正如那些在小剧场中竞赛的人会寻找更大的剧场，仿佛这个小剧场不足以展示他们；同样，那些在天主眼前竞赛之后又寻求人的喝彩的人，放弃了更大的赞美，急切追求较小的，给自己招来严厉的惩罚。难道这不就是一切颠倒的原因吗？这是使整个世界混乱的原因：我们做任何事都着眼于人，甚至为了我们的好处，我们也不把天主作为欣赏者当一回事，反而寻求同侪的认可；相反地，对于坏事，我们藐视天主，却惧怕人。然而，他们必定会与我们一同站在那个审判台前，却对我们毫无益处。但我们现在藐视的天主，祂自己要向我们宣判。\par

然而，虽然我们明知这些事，我们仍然向人张口，这是最大的罪恶之一。如果没有人看见，没有人会选择犯奸淫；但即使他被那祸害燃烧一万次，对人类的敬畏也能克服这种情欲的暴政。然而在天主眼前，人们不仅犯奸淫和私通，而且还有许多更可怕的事，许多人已经做了，并且仍在做。那么，仅此一点，难道还不足以从天上引来一万个霹雳吗？我说奸淫和私通了吗？不，甚至比这些小得多的事，我们在人面前也不敢做；但在天主眼前，我们不再畏惧。实际上，世上的一切邪恶都由此而来：因为在真正邪恶的事上，我们敬畏的不是天主，而是人。\par

因此你看，那些真正美好的事物，不被一般人视为如此，反而成了我们厌恶的对象，我们不探究事物的本性，而是尊重多数人的意见；同样地，在邪恶的事上，我们也依此原则行事。所以，有些事物并非真正美好，但在众人眼中显得美好，我们却出于同一习性而追求它们，如同追求美好事物一样。因此，我们两面都走向毁灭。\par

或许许多人会觉得这个评论有点晦涩。因此我们必须更清楚地表达它。当我们行不洁之事时（因为必须从所举的例子开始），我们更怕人甚于天主。因此，当我们这样把自己置于他们之下，使他们成为我们的主人时，还有许多其他事情，在这些主人看来是邪恶的，而其实并非如此；我们也同样逃避它们。例如：贫穷，许多人视为可耻；我们逃避贫穷，不是因为它是可耻的，也不是因为我们如此认为，而是因为我们的主人认为它可耻，我们害怕他们。同样，不被尊重、被人轻视、没有权柄，在大多数人看来也是极大的羞耻和卑贱。我们也逃避这个；不是谴责事情本身，而是因为主人的判决。\par

另一方面，我们也遭受同样的祸害。正如财富被视为好东西，骄傲、炫耀和引人注目也被视为好东西。同样地，我们追求这些，并非因为考虑到事物本身的善，而是受我们主人意见的左右。因为人民是我们的主人，还有大群\OriginalTerm{暴民}{\GreekText{ὅ} \GreekText{πολὺς} \GreekText{όχλος}}；一个野蛮的主人和严厉的暴君：我们根本不需要命令来听从他们；只要我们知道他想要什么，就不需要命令而服从：我们对他怀有如此大的善意。再说，天主天天威胁和劝诫，却未被听从；但平民百姓，充满混乱，由各种渣滓组成，却不需要一句命令；只要他们示意什么是他们所喜悦的，我们就立即在一切事上服从。\par

\Quote{但怎样，}有人说，\Quote{才能摆脱这些主人呢？} 通过拥有一颗比他们更大的心灵；通过观察事物的本性；通过谴责大众的呼声；最重要的是，通过训练自己在真正可耻的事上，不怕人，而怕那不眠之目；同样，在一切善事上，寻求来自祂的冠冕。因为这样，我们也不会在其他类型的事上容忍他们。因为凡是行善却认为他们不配知道自己的善行，并满足于天主的称许的人，也不会在意相反之事上的评价。\par

\Quote{这怎么可能？}你会说。想想人是什么，天主是什么；你离弃谁，又投奔谁；你很快就会完全正确。人同你一样处在同样的罪中，同样的定罪，同样的惩罚。\BibleQuote{人似虚幻，}\BibleRef{咏 143:4}，他没有正确的判断，需要从上而来的纠正。\Quote{人是尘土和灰烬，}如果他给予赞美，常常是随意地，或出于偏爱或恶意。如果他毁谤和控告，同样也是出于同样的目的。但天主不是这样：祂的判决无可指摘，祂的审判纯洁。因此我们必须常常投奔祂；不仅因为这些原因，还因为祂造了你，并且超过一切地宽恕你，爱你胜过你爱自己。\par

那么，为什么我们忽略如此可钦佩（\OriginalTerm{可钦佩的}{\GreekText{θαυμαστόν}}）的认可者，而投奔于人，那虚无、全凭冲动、全然随意的人呢？他若在你并非如此时称你为邪恶和污秽，你更当可怜他，为他败坏而哭泣；并鄙视他的意见，因为他理智的眼睛是昏暗的。连宗徒们也这样被恶评，他们却嘲笑诽谤者。但他若称你为良善和仁慈，你若果真如此，也不因这意见而自傲；你若并非如此，就更当鄙视它，视之为嘲弄。\par

你想知道大多数人的判断是如何败坏、无用、可笑吗？有些来自癫狂错乱的人，有些来自吃奶的婴孩？听听从一开始的情况。我将告诉你判断，不仅是民众的，还有那些被认为最智慧的人，最早期的立法者。因为谁会认为比那些被认为配为城邦和人民立法的人更智慧呢？然而，这些智慧的人却认为邪淫并非邪恶，也不值得惩罚。至少，没有一条外教法律惩处它，或将人因此事受审。如果有人因这类事将他人告上法庭，众人会嘲笑，法官也不会允许。掷骰子同样不受任何惩罚：无人因此受罚。酗酒与暴食，不仅不是罪行，反被许多人视为美事。在军中狂欢中，这是一大竞赛；那些最需要清醒头脑和强壮身体的人，却最屈服于酗酒的暴政，既削弱身体，又使灵魂昏暗。然而，立法者中没有一个惩罚这一过错。还有什么比这疯狂更糟糕的呢？\par

那么，如此堕落之人的好评，对你来说是值得渴望的吗？你难道不躲到地底下吗？即使所有这样的人称赞你，难道你不该羞愧掩面，为被如此败坏判断的人鼓掌吗？\par

再说，一般立法者并不认为亵渎是可怕的事。至少，从未有人因亵渎天主而受审和受罚。但若有人偷别人的衣服，或割钱包，他的两侧被剥皮，常被处死；而亵渎天主的人，外教立法者却不加追究。若有人有妻室又诱奸女仆，外教法律和一般人看来似乎算不了什么。\par

10. 你还想听听另一类显示他们愚昧的事吗？因为他们不惩罚这些事，所以另有一些事他们用法律强制执行。什么事呢？他们聚集人群填满剧场，并在那里引入妓女和卖淫孩童的合唱队，是的，那些践踏自然本身的人；他们让众人坐在高处，这样俘获他们的城市；他们如此加冕那些因战功和胜利而永远受崇拜的强大的王。然而，还有什么比这荣誉更乏味？什么比这快乐更不快乐？那么，你从这些人中寻找评判者来赞扬你的行为吗？你愿意与舞者、柔弱者、小丑和娼妓一起享受称赞的声音吗？回答我。\par

这些事岂不都是极端愚昧的明证吗？我愿意问他们：颠覆自然法则、引入非法交合是不是一件可怕的事？他们一定会说，是可怕的：至少他们佯装对那种罪行施加惩罚。那么，你为什么把那些受虐待的不幸之人带上舞台？不仅带他们上场，还用数不尽的荣誉和说不尽的礼物来尊崇他们？在其他地方，你惩罚那些胆敢如此行事的人；但在这里，你却像对待城市的公共恩人一样，在他们身上花钱，以公费供养他们。\par

\Quote{然而，}你会说，\Quote{他们是\OriginalTerm{卑贱的}{\GreekText{ἄτιμοι}}。}那么为什么还要训练他们（\OriginalTerm{教练}{\GreekText{παιδοτριβεῖς}}）？为什么选择卑贱的人来向国王致敬？为什么毁坏我们的城市（\OriginalTerm{你使……颠覆}{\GreekText{ἐκτραχηλίζεις}}，普鲁塔克《\OriginalTerm{论儿童教育}{\GreekText{περὶ} \GreekText{παίδων} \GreekText{ἀγωγῆς}}》，第17章）？或者为什么在这些人物身上花费如此之多？既然他们是卑贱的，驱逐对卑贱者来说才是最合适的。因为你为什么使他们成为卑贱的？是出于赞扬还是谴责？当然是谴责。接下来难道要这样：虽然你在定罪之后使他们成为卑贱的，却好像他们是可敬的，跑去看他们，赞美、钦佩、鼓掌？何须我提赛马中的那种魅力？或者野兽搏斗的竞赛？因为那些地方也充满了各种无意义的刺激，训练民众养成一种残忍、野蛮和无人性的性情，使他们习惯于看人肢体撕裂、血流成河、野兽的凶猛混乱一切。所有这些都是我们聪明的立法者从起初就引入的，成为这么多的灾祸！而我们的城市却鼓掌赞赏。\par

11. 但是，如果你愿意，暂且搁置这些明显公认可憎的事——尽管在异教立法者看来并非如此（\OriginalTerm{不被认为}{\GreekText{οὐκ} \GreekText{ἐδοξεν}}，或许是\Quote{未被规定}）——让我们来看看他们严肃的诫命；你会看到这些诫命也因大众的意见而被败坏。这样，婚姻在我们和外人看来都是可尊敬的\BibleRef{希 13:4}，而且它确实是可尊敬的。但是，当婚礼举行时，就会发生像你马上会听到的那些可笑事情：因为大多数人被习俗所迷住和欺骗，甚至没有意识到它们的荒谬，需要别人来教导他们。因为舞蹈、钹、笛子、可耻的话语、歌曲、醉酒、狂欢，以及所有魔鬼的一大堆（\OriginalTerm{魔鬼的大量垃圾}{\GreekText{πολὺς} \GreekText{ὁ} \GreekText{τοῦ} \GreekText{διαβόλου} \GreekText{φορυτός}}）垃圾都被引入了。\par

我确实知道，我指责这些事情会显得可笑，并且会因扰乱古法而被大众指责为极大的愚昧：因为正如我以前所说，习俗的欺骗力量是巨大的。但是，我仍不会停止重复这些话：因为确实有、确实有希望，虽然不是所有人，但至少有一些人会接受我们的说法，并选择与我们一同被嘲笑，而不是我们与他们一同发出那种应当流泪、招致无尽惩罚和报复的嘲笑。\par

因为一个完全在家中度日、自幼受训于谦逊的少女，竟被迫突然抛掉一切羞耻，从结婚之初就被教导鲁莽，置身于放荡、粗野、不贞和柔弱的男人中间，这岂非最该受谴责的吗？从那天起，新娘心中会种下什么邪恶？不自重、轻浮、傲慢、虚荣：因为这些自然会发展下去，她们会希望所有的日子都像这样。因此，我们的妇女变得奢侈浪费，因此她们缺乏谦逊，因此她们生出无数的邪恶。\par

不要跟我提习俗：因为如果一件事是邪恶的，那就连一次也不要做；但如果它是好的，那就经常做。告诉我，行淫难道不是邪恶的吗？那么我们可以允许这件事发生一次吗？绝不。为什么？因为即使只做一次，它仍然是邪恶的。同样，如果这样款待新娘是邪恶的，那就连一次也不要做；如果它不是邪恶的，那就始终做。\par

\Quote{那么，}有人说，\Quote{你是在指责婚姻吗？告诉我。}但愿不会。我没有那么愚蠢：而是指责那些如此不配地附加于婚姻的事物——涂脂抹粉、画眉毛，以及所有诸如此类的讲究。因为事实上，从那天起，她在命定的配偶之前就会吸引许多情人。\par

\Quote{但许多人会因为她的美貌而赞赏那女人。}那又怎样？即使她谨慎，也很难避免恶意的猜疑；但如果她粗心，她将很快被征服，因为她在当天就获得了放荡行为的起点。\par

尽管邪恶如此之大，省略这些程序却被某些无异于禽兽的人称为侮辱，他们愤慨于女人没有被展示给众人，没有被摆出来作为舞台上的奇观，为所有观看者所共有；而最确实地，他们反而应该把这些事情发生视为侮辱，视为笑柄和闹剧。因为即使现在，我知道人们会指责我极其愚昧，并把我当作笑柄：但只要能从中获得任何益处，我就能忍受嘲弄。因为我一边劝人藐视多数人的意见，一边自己比任何人更被那种感觉所征服，那我确实该被嘲笑。\par

请看，这所带来的是何等后果。不仅是在白天，甚至在傍晚，他们故意找一些已经喝醉、神志不清、因奢华饮食而兴奋的男人，去观赏少女的容貌；不止是在屋内，甚至穿过市场，她们被隆重地引领着游行，以便让所有人都能看到；在深夜用火把引导，以便所有人都能看到她。他们的行为无非是在推荐一件事：让她从此抛弃一切羞耻。他们甚至不止于此，还用可耻的话语引导她。这在群众中成了一条法律。甚至逃亡的奴隶和罪犯，成千上万、穷凶极恶之徒，可以肆无忌惮地说出任何他们想说的话，既针对她也针对将要娶她回家的人。没有任何庄严可言，一切都卑贱且充满不雅。新娘看到并听到这些事，岂不是一堂很好的贞洁课吗？\EditorialNote{Savile 将此句读作疑问句。}在这些放荡之徒中，有一种魔鬼般的竞争，看谁更热心地滥用辱骂和污言秽语，从而让整个聚会陷入窘迫；那些找到了最多的辱骂和最无耻的话来攻击邻舍的人，就得意洋洋地离去。\par

现在我知道自己是一个令人讨厌、令人不快、脾气乖戾的人，仿佛我正在剥夺生活中的某些乐趣。然而，这正是我悲伤的原因：如此不愉快的事竟被视为一种乐趣。请问，被侮辱和谩骂怎能不令人不愉快？连同你的新娘一起被所有人责备？若有人在市场上说你妻子的坏话，你会没完没了地纠缠，觉得活着没有价值；那么，在全体市民面前，与你未来的配偶一起蒙羞，你反而感到高兴、面露喜色吗？这岂不是奇怪的疯狂！\par

有人说：\Quote{但这是习惯。}但正因为如此，我们才最应该哀叹，因为魔鬼已经用习惯把这事围了起来。事实上，婚姻是庄严的事，它延续我们的种族，是无数祝福的缘由；那恶者内心嫉妒，知道婚姻被设立为对抗不洁的屏障，于是通过一种新的手段，把各种不洁引入其中。无论如何，在这种聚会上，许多贞女甚至已被败坏。即使并非每次都如此，那也是因为魔鬼暂时满足于那些可耻的话语和歌曲，满足于公开展示新娘，并带着新郎胜利地穿过市场。\par

此外，因为这一切发生在傍晚，甚至连黑暗都不能遮掩这些恶事，所以许多火把被拿来，不让这丢脸的场面被隐藏。那大群人是为什么？那宴饮是为什么？那笛子又是为什么？最明显是为了防止甚至那些在家中、\OriginalTerm{浸没}{\GreekText{βαπτιζόμενοι}}在沉睡中的人，对这些事一无所知；他们被笛声唤醒，倚着窗户向外看，成为这场闹剧的见证人。\par

对那些歌曲本身，又能说什么呢？它们充满了各种不洁，引入了畸形的恋情、非法的结合、家庭的颠覆和无休止的悲剧场景，不断提及\Quote{朋友和情人}、\Quote{情妇和爱人}这些名号。更严重的是，年轻女子也出席这些场合，她们抛弃了所有的羞耻；名义上是尊敬新娘，不如说是侮辱她，甚至拿自己的得救当儿戏，在放荡的青年中间，用她们混乱的歌曲、污秽的话语、魔鬼般的和声扮演无耻的角色。那么告诉我：你还在问\Quote{通奸从何而来？淫乱从何而来？婚姻的侵犯从何而来？}吗？\par

12. 你会说：\Quote{但做这些事的不是高贵端庄的妇女。}那么，你既然在我说话之前就自己知道了这条法律，为什么还要嘲笑我的规劝呢？我说，如果这些做法是对的，那就允许那些出身高贵的妇女也这样做。因为假如这些其他人生活在贫困中，难道她们就不是贞女吗？难道她们就不应该注意贞洁吗？但现在，这里有一个贞女在浪荡青年的公共剧场里跳舞；我问你，她不是比一个妓女更受辱吗？\par

但如果你说：\Quote{女仆才做这些事；}即便如此，我也不会免除你的责任；因为连这些事也不应允许她们做。所有这些邪恶的根源在于，我们不关心自己家中的人。但轻蔑地说\Quote{他是奴隶}和\Quote{她们是使女}就足够了。然而，我们天天听见：\BibleRef{迦 3:28}\BibleQuote{在基督耶稣内，既没有奴隶，也没有自由人。}再者，若是马或驴，你绝不疏忽，而是竭尽全力不让它成为劣等；而你的奴隶，同你一样有灵魂，你却不关心吗？我说奴隶，为什么不说儿子和女儿呢？然后必然的结果是什么？当这一切走向毁灭时，悲伤（\OriginalTerm{悲伤}{\GreekText{λύπην}}，意指\OriginalTerm{伤害}{\GreekText{λύμην}}）必然立即进入。而且常常也造成非常大的损失，因为贵重的金饰在人群和混乱中丢失了。\par

13. 然后，结婚之后，若偶然生下一个孩子，我们将会再次看到同样的愚行和许多充满荒谬的习俗（\OriginalTerm{符号}{\GreekText{σύμβολα}}）。因为当为婴儿取名的时候，他们不关心像古人起初那样以圣人之名命名，而是点灯并给灯命名，然后以燃烧最久的那个灯的名字给孩子命名，由此推测他会活得很久。毕竟，如果孩子多次夭折（也确实常有），魔鬼便会大大嘲笑，因为它戏弄了他们，仿佛他们是愚蠢的孩子。关于挂在手上的护符和铃铛、红线以及其他充满极端愚行的事物，我们还能说什么呢？那时他们本应以十字架的保护来装备孩子，别无他物。但现在，那曾转化整个世界、重创魔鬼并推翻其一切权能的十字架却被轻视，而线、线和那类护符却被托付来保护孩子的安全。\par

我能再提一件比这更可笑的事吗？但愿无人指责我们说话不合时宜，如果我们的论证也要举这个例子。因为要清洗溃疡的人，必不先怕弄脏自己的手。那么，这件极其可笑的习俗是什么呢？它确实被看作无关紧要（这正是我悲哀的原因），但它是愚行和极度疯狂的开始。女人们在浴池中，乳母和侍女们拿起泥，用指头涂抹，在孩子的前额上做一个记号。若有人问：泥和粘土是什么意思？答案是：\Quote{它能驱除邪眼、巫术和嫉妒。}令人惊讶！泥竟有这等效力！粘土竟有这等威力！它有多大的力量呢？它竟能避开魔鬼的一切大军。告诉我，你们能不羞愧得无地自容吗？你们永远不能明白魔鬼的诡计，它如何从早年开始逐渐引入它所设计的种种邪恶吗？因为如果泥有这种效果，为什么当你成年、品格已形成时，你自己不为自己的额头也这样做呢？你比孩子更可能招人嫉妒。为什么你不全身涂满泥呢？我说，如果泥在额头上的功效如此之大，为什么你不在全身涂上泥呢？这一切对撒旦来说都是娱乐和做戏，不只是嘲弄，而是地狱之火，这些受骗者正趋向那个结局。\par

14. 在希腊人中发生这样的事不足为奇；但在十字架的崇拜者（\OriginalTerm{崇拜十字架}{\GreekText{τὸν} \GreekText{σταυρὸν} \GreekText{προσκυνοῦσι}}）和不可言喻的奥秘的分享者，以及声称持守\OriginalTerm{如此高尚的道德}{\GreekText{τοσαῦτα} \GreekText{φιλοσοφοῦσιν}}的人中间，竟流行这等不雅之事，这尤其令人一再痛惜。天主以属神的傅油尊荣了你，而你却用泥污秽你的孩子？天主尊荣了你，而你自己却羞辱自己？当你应该在他额上刻上那提供无敌保障的十字架时，你竟放弃这个，而投身于撒旦的疯狂之中？\par

若有人视这些事为小事，就当知道它们是重大邪恶的根源；而且就连保禄也不认为忽略小事是合适的。因为，告诉我，有什么比一个人蒙头更小的事呢？然而请看他是何等认真地对待这事，并何等恳切地禁止它；在诸多话中说：\Quote{他羞辱自己的头。}\BibleRef{格前 11:4}。如果蒙头的人\Quote{羞辱自己的头}，那么用泥涂污自己孩子的人，又怎能不使其成为可憎的呢？因为，我想知道，他怎能把孩子带到司铎的手中呢？你怎能要求在那你涂了泥的额头上，由司铎的手放上印记呢？不，我的弟兄们，不要这样做，而要从早年开始用属神的武器环绕他们，教导他们用手（\OriginalTerm{用手教导他们封印额头}{\GreekText{τῇ} \GreekText{χειρὶ} \GreekText{παιδεύτε} \GreekText{σφραγίζειν} \GreekText{τὸ} \GreekText{μέτωπον}}）封印额头；在他们自己尚不能这样做之前，你当在他们身上印上十字架。\par

何必再提在分娩和生育中的其他撒旦式习俗呢？这些习俗是产婆们自招祸害引入的。何必提每个人临终时和在埋葬时发出的哭喊、非理性的哀号、葬礼上的愚行、对墓碑的热心、哭丧妇女的纠缠不休和可笑的聚集，以及对日子的遵守——我是指出生和离世的日子？\par

15. 我恳求你们，难道你们追求的就是这些人的好评吗？热切寻求那些思想如此败坏、行为全无定准之人的称赞，若不是极端的愚行又是什么呢？我们岂不应常常投靠那不眠之眼，并在我们一切言行中注视祂的审判吗？因为这些人的赞许，即使他们认可，也不能使我们获益。但祂若接纳我们的行为，就必在此世使我们光荣，并在将来之日将那不可言喻的美善赐予我们：愿我们众人都有分获得，赖我们的主耶稣基督的恩宠和慈爱；愿光荣、权能、尊崇归于祂及圣父和圣神，从今世直到永远。阿们。\par

\SourceRecord{Translated by Talbot W. Chambers.；Nicene and Post-Nicene Fathers, First Series；Vol. 12.；Edited by Philip Schaff.；Buffalo, NY: Christian Literature Publishing Co.,；1889.；<http://www.newadvent.org/fathers/220112.htm>.}

% structure: p0001 source=h1 target=section reason=page-title

\section{格林多前书第十三篇讲道}

% structure: p0002 source=h2 target=subsection reason=relative-heading-rank; base=h2

\subsection{\BibleRef{格前 4:10}}

\BibleQuote{我们为基督成了愚拙的人}（因为从此刻起必须恢复我们的讲论）\BibleQuote{但你们在基督内却是聪明的；我们软弱，你们却强壮；你们光荣，我们却受羞辱。}\par

他先用极其严厉的话语充实了他的讲论，这比任何直接的指责都更具打击力，并且说了：\BibleQuote{你们没有我们，已经为王了}和\BibleQuote{天主把我们宗徒列在最后，好像定了死刑的人}，然后通过下文展示他们如何是\BibleQuote{定了死刑的人}，说道：\Quote{我们成了愚拙的、软弱的、受轻视的，又饥又渴，赤身露体，被人拳打脚踢，居无定所，并且劳苦，亲手作工}：这些都是真正导师和宗徒的标志。而其他人却以相反的事——智慧、光荣、财富、体面——自夸。\par

因此，为了压低他们的自负，并指出在这些事上，他们非但不可自夸，反而应当羞愧，他首先讥讽他们说：\BibleQuote{你们没有我们，已经为王了}。仿佛在说：\Quote{我的判断是：现在不是荣耀的时候，也不是光荣的时候——你们所享受的正是这些——而是迫害和凌辱的时候，正如我们所受的。如果不是这样，如果此时反倒是报偿的时候，那么据我看，}（但他说这话是带着讽刺的）\Quote{你们这些门徒，已经变得不比君王差；而我们这些导师和宗徒，并且最有资格获得报酬的，不仅大大落后于你们，甚至像定了死罪的人，即定了罪的囚犯，完全生活在耻辱、危险和饥饿中：还被侮辱为愚拙的，被驱赶，忍受一切难以忍受的事。}\par

他说这些话，是为了借此也使他们思考：他们应当热切追求宗徒们的处境——他们的危险和凌辱，而不是他们的荣誉和光荣。因为福音所要求的是这些，而不是别的。但他不是直接这样说，以免显得令人反感；而是以他自己特有的方式来进行这项责备。如果他以直接的方式表达，他会这样说：\Quote{你们错了，被迷惑了，远远偏离了教导的方式。因为每一位宗徒和基督的仆人，都应当被视为愚拙的，应当生活在痛苦和耻辱中：这正是我们的情形；而你们却处于相反的情况。}\par

但这样表达可能会更加得罪他们，因为其中包含着对宗徒的赞美；并且可能使他们因为被指责为懒惰、虚荣和奢侈而变得更加激烈。因此，他不以这种方式陈述，而是以另一种更突出但较少冒犯的方式；这就是为什么他继续以如下的方式讲论，带着讽刺说：\Quote{但你们是强壮的、光荣的}；因为，如果他没有用讽刺，他会这样说：\Quote{不可能一个人被认为是愚拙的，而另一个是聪明的；一个强壮，另一个软弱；因为福音同时要求这两者。如果这事物的本性是有人是这样，有人是那样，或许你们所说的有些道理。但现在，无论是被视为聪明、光荣，还是免于危险，都是不被允许的。否则，必然得出你们在天主眼中比我们更受优待的结论：你们这些门徒在我们这些导师之前，而且是在我们经历了无限的艰辛之后。}如果这对任何人来说都太过分，那么你们就应当以我们的处境为目标。\par

并且，他说：\Quote{不要有人以为我只是说过去的事：}\par

% structure: p0009 source=h2 target=subsection reason=relative-heading-rank; base=h2

\subsection{\BibleRef{格前 4:11}}

\BibleQuote{直到此时此刻，我们还是又饥又渴，赤身露体。}你看见了吗？基督徒的全部生活都应当是这样，而不仅仅是一两天。因为那在单次比赛中获胜的角力者，若只取得一次胜利，就被加冕；但若他摔倒，就不再被加冕。\par

\BibleQuote{又饥又渴}——针对奢侈的人。\BibleQuote{被人拳打脚踢}——针对自高自大的人。\BibleQuote{居无定所}——因为我们被驱赶。\BibleQuote{赤身露体}——针对富有的人。\par

% structure: p0012 source=h2 target=subsection reason=relative-heading-rank; base=h2

\subsection{\BibleRef{格前 4:12-13}}

\BibleQuote{并且劳苦}——现在又针对那些既不肯辛劳也不肯冒险，却自己享受果实的假宗徒。\Quote{但我们却不是这样}，他说：\Quote{除了外来的危险，我们还用不断的劳动竭尽全力。而且更甚的是，没有人能说我们为此烦恼，因为我们对如此对待我们的人反而以相反的态度回报：我说，这才是要点：不是我们受苦（因为这是所有人的常事），而是我们受苦却不沮丧或烦恼。我们非但不沮丧，反而充满喜乐。一个确切的证据就是，我们对那些错待我们的人以相反的态度回报。}\par

至于他们确实如此行的事实，请听下文。\par

\BibleQuote{被人咒骂，我们就祝福；被人迫害，我们就忍受；被人毁谤，我们就劝勉；我们成了世界的渣滓。}这就是\OriginalTerm{为基督成了愚拙的}{fools for Christ's sake}的含义。因为那受委屈而不报复、也不烦恼的人，被外邦人视为愚拙、受屈辱和软弱。为了使他的言语不至于因提到他们所遭受的痛苦而太令人不快，他说了什么？\BibleQuote{我们成了渣滓}，不是\Quote{你们城市的}，而是\Quote{世界的}。又说，\BibleQuote{万物的渣滓}，不仅是你们的，也是所有人的。正如当他谈论基督的眷顾时，他略过大地、天空、整个受造界，只提出十字架；同样，当他想要吸引他们归向自己时，他略过自己所有的奇迹，只提自己为他们所受的苦难。所以，当我们被任何人伤害和轻视时，我们也应当把我们为他们所受的一切摆出来。\par

\Quote{万人的渣滓，直到如今。}这是他最后给予的有力一击，\Quote{万人}；\Quote{不限于迫害我们的人，}他说，\Quote{也包括那些我们为他们受苦的人：我何等亏欠他们。}这是一个深切关怀者的表达；他自己并不痛苦，而是渴望使他们感到（\OriginalTerm{打击}{\GreekText{πλῆξαι}}）那有无数怨言可发的人竟向他们致意。为此，基督命令我们温和地忍受侮辱，好使我们既在崇高的德行中操练自己，又使对方更加羞愧。因为责备的效果不如沉默。\par

% structure: p0017 source=h2 target=subsection reason=relative-heading-rank; base=h2

\subsection{\BibleRef{格前 4:14}}

3. 接着，他见那打击难以承受，便迅速医治，说道：\Quote{我写这些话，不是要羞辱你们，而是劝诫你们如同我亲爱的儿女。} \Quote{因为我不是为使你们羞愧，}他说，\Quote{才说这些话。}他自己在言语中所做的，他却说没有做；反而承认自己做了，但并非出于恶意和怨恨。其实，这种抚慰的方式是最好的：我们说出当说的话，再从动机上加以辩解。因为若不说，便不可能，否则他们将得不到纠正；另一方面，说了之后，不加医治，又太残忍。因此，他在严厉中加以辩解：这非但不会削弱刀锋的效力，反而使它更深入，同时减轻伤口的全部疼痛。因为当一个人被告知，这话不是出于责备，而是出于爱时，他就更乐意接受纠正。\par

然而，即使在这里也有严厉之处，并且强烈地唤起他们的羞耻感（\OriginalTerm{羞耻}{\GreekText{ἐντροπή}}）：因为他没有说\Quote{作为主人}，也没有说\Quote{作为宗徒}，更没有说\Quote{作为你们是我的门徒}（这本符合他对他们的权利要求），而是说：\Quote{如同我亲爱的儿女，我劝诫你们。}并且不单是\Quote{儿女}，而是\Quote{深爱的}。他说：\Quote{请原谅我。} \Quote{如果说了什么不愉快的话，那全是出于爱。}现在，谁会忍心不听一位悲伤的父亲在给予善意劝告呢？因此，他之前没有说这话，而是在打击之后才说。\par

\Quote{那么怎么办？}有人可能会说：\Quote{别的教师难道不宽容我们吗？} \Quote{我不是这样说，而是他们没有如此容忍。}然而，他没有立即说出这一点，而是通过他们的称呼和头衔暗示出来：他使用了\Quote{导师}和\Quote{父亲}这些词。\par

% structure: p0021 source=h2 target=subsection reason=relative-heading-rank; base=h2

\subsection{\BibleRef{格前 4:15}}

4. \Quote{因为}他说，\Quote{你们在基督里虽有上万导师，但为父的不多。}他在这里不是炫耀自己的尊荣，而是凸显爱的无比伟大。这样，他也没有伤害其他教师：因为他加了\Quote{在基督里}这一限制，反而安慰了他们，称他们当中那些热心且耐心劳苦的人为\Quote{导师}，而不是\Quote{奉承者}，同时也表明了他自己对他们的殷切关怀。为此，他没有说\Quote{主人不多}，而是说\Quote{为父的不多}。他的目的并非列出任何尊荣的名称，或争论他们从他那里获得了更大的益处；而是承认其他人对格林多人所付出的巨大辛劳（因为\Quote{导师}一词的意义就在于此），而爱的优越则保留为自己的部分（因为\Quote{父亲}一词的意义正在于此）。\par

他不单说\Quote{没有人像你这样爱我们}，这是一个不容置疑的陈述，而且他还提出一个事实。这事实是什么呢？\Quote{因为我在基督耶稣里，借着福音生了你们。在基督耶稣里。}这不是归功于我自己。他再次打击那些将自己的名字赋予自己教导的人。因为\Quote{你们}他说\Quote{就是我宗徒职分的印证。}又说：\Quote{我栽种了。}在此处则说：\Quote{我生了。}他用的不是\Quote{我传了道}，而是\Quote{我生了}，使用了表示\OriginalTerm{自然关系的词}{\GreekText{τοῖς} \GreekText{τῆς} \GreekText{φύσεως} \GreekText{ὀνόμασι}}。因为他此刻唯一的关切是显明他对他们的爱。\Quote{因为他们确实从我这里接受了你们，并带领了你们；但你们成为信徒，完全是因着我。}如此，因为他曾说\Quote{如同儿女}，免得你以为那是奉承之词，他又摆出了事实。\par

% structure: p0024 source=h2 target=subsection reason=relative-heading-rank; base=h2

\subsection{\BibleRef{格前 4:16}}

5. \Quote{我劝你们效法我，如同我也效法基督。} \OriginalTerm{如同我也效法基督}{\GreekText{κάθως} \GreekText{κὰγὼ} \GreekText{Χριστοῦ}}（我们的译本遗漏了这一句；武加大译本有此句，参阅\BibleRef{格前 11:1}）令人惊叹！教师的讲论何等大胆！形象何等完美，当他甚至能以此劝勉他人！他并非出于自高自傲，而是暗示德性是容易的事。就好像他说：\Quote{\InnerQuote{我不能效法你。你是导师，而且是伟大的导师。}因为你我之间的差距不如基督与我之间的差距大，然而我却效法了他。}\par

另一方面，他在写信给厄弗所人时，没有提到自己，而是直接引导他们到一个点上：他的话语是\Quote{你们要效法天主。}\BibleRef{弗 5:1}但在此时，由于他的讲话是针对软弱的人，他便顺便把自己放进去。\par

此外，他也表明，即使如此也能效法基督。因为那复制完美印章印记的人，也就复制了原版模型。\par

那么我们来看他是如何跟随基督的：因为这种效法不需要时间和技巧，只需要坚定的心志。好比我们走进画室看画家作画，即使看上千次也无法复制肖像。但藉着聆听，我们就能效法他。你们愿意我将画板摆在你们面前，为你们勾勒出保禄的生活方式吗？好，就让我们把它呈现出来——那幅比所有皇帝画像更光辉的图画：因为它的材料不是胶合的木板，也不是张开的画布；而是天主的化工——它既是灵魂又是肉身：灵魂是天主的化工，不是人的；肉身亦然。\par

你们在这里鼓掌了吗？不，现在不是鼓掌的时候；而是在接下来的部分：我说的是鼓掌，也是效法。因为到目前为止，我们只看到那毫无例外的共同材料：就灵魂之为灵魂而言，灵魂与灵魂并无差别；但内心的意向显出了差异。正如一个身体之为身体并无差别，但保禄的身体与常人无异，唯有危险使一个身体比另一个更光彩照人；灵魂也是如此。\par

6. 假设我们的画板是保禄的灵魂：这块画板不久前还覆盖着烟灰，布满蜘蛛网；（因为没有什么比亵渎更糟糕的了；）但当那转化万物的主来到，看见他的线条之所以如此，并非由于懈怠或迟钝，而是由于缺乏经验和没有真正虔敬的\OriginalTerm{色彩}{\GreekText{τὰ} \GreekText{ἄνθη}}（因为他确有热诚，却没有\Quote{按知识的热诚}）；主便将真理的色彩，即恩宠，赐予了他；顷刻间，他便展现了君王的形象。因为他获得了色彩，学会了他所不知道的，毫不耽搁，立刻显为一位最卓越的艺术家。他首先展示了君王的头——宣讲基督；然后也展示了身体的其他部分——完美基督徒生命的身体。我们知道画家们通常闭门不出，极其精细而安静地完成他们的作品，不向任何人开门；但这个人，在众目睽睽之下，在普遍的反对、喧嚷、纷扰中，竟在这样的环境下完成了这幅君王肖像，并且没有受到阻碍。因此他说：\Quote{我们成了世界的一台戏，}在地、海、天以及整个可居住的地球，在物质与理智的世界中，他正在绘制那幅自画像。\par

你愿意看看这肖像从头部向下的其他部分吗？或者你希望我们从下往上描述？那么请看一座金像，或者比金更贵重的东西，一座可以立在天空的像；不是用铅固定在一个地方，而是从耶路撒冷直到依里黎苛\BibleRef{罗 15:19}，再进入西班牙，仿佛插上翅膀飞遍世界各地。因为还有什么比这些\Quote{美丽}的\Quote{脚}更\Quote{美丽}的呢？它们走遍了天下的每个角落。先知古时也宣告这\Quote{美丽}，说：\BibleQuote{传报和平福音者的脚步是多么美丽啊！}\BibleRef{依 52:7}你看见那双脚多么美丽了吗？你也要看看胸怀吗？来，让我也给你看这个，你会发现它比那些美丽的脚，甚至比古立法者的胸怀更加辉煌。因为梅瑟揣着石版；但此人心中却有基督自己：正是他佩戴的君王肖像。\par

因此，他比赎罪盖和革鲁宾更令人敬畏。因为从它们那里没有发出像从他那里发出的声音；但藉着它们，主要与人们谈论感官之物；而藉着保禄的口舌，则谈论天上之上的事。再者，赎罪盖仅向犹太人传谕神谕；但从此处却向全世界传谕：那里是无生命之物；这里却是富有德行的灵魂。\par

这赎罪盖甚至比天空更明亮，不是因繁星或太阳的光芒照耀，而是公义的太阳本身就在那里，从那里发出祂的光线。同样，在我们的天空中，不时有云彩飘过，使其昏暗；但那个胸怀从未有任何风暴横扫而过。或者说，多次有风暴横扫而过，但光明并未被遮蔽；反而在试探和危险中，光明闪耀。因此，当他被锁链捆绑时，仍不断呼喊：\BibleQuote{天主的话决不会被束缚。}\BibleRef{弟后 2:9}就这样，通过那口舌，祂不断发出光芒。没有恐惧，没有危险能使那胸怀昏暗。也许这胸怀似乎超过了脚；然而，脚是美丽的，胸怀也是美丽的。\par

你也要看看那拥有其独特美丽的肚子吗？请听他关于它所说的：\BibleQuote{若食物使我的弟兄跌倒，我就永远不吃肉。}\BibleRef{格前 8:13}\BibleQuote{不吃肉、不喝酒，不做任何使你弟兄跌倒、或冒犯他、或使他软弱的事，是好的。}\BibleRef{罗 14:21}\BibleQuote{食物是为肚腹，肚腹是为食物。}\BibleRef{格前 6:13}还有什么比这肚子更美丽的呢？它被如此训练得安静，学会了一切节制，懂得如何忍受饥饿、饥渴和干渴。因为它像一匹戴着金缰绳的骏马，以匀称的步伐行走，战胜了本性的需要。因为基督在其中行走。既然它是如此节制，显然其他的恶习都被清除了。\par

你要看看那双手吗？那些他现在拥有的手？或者你宁愿先看看他们从前的邪恶？\BibleRef{宗 8:3} \BibleQuote{进入各家（这同一个人），他拉走}，不久之前，\BibleQuote{连男带女}，用的不是人的手，而是某种凶猛的野兽的手。但他一接受了真理的色彩和灵性的经验，这些就不再是人的手，而是属神的手；日复一日地被锁链捆绑。它们从未击打过任何人，却无数次被击打。有一次甚至一条毒蛇 \BibleRef{宗 28:3} 也尊敬了那双手：因为它们不再是人的手了；因此毒蛇甚至没有咬住它们。\par

你也要看看那背部吗？它与其他肢体相似。听听他对此怎么说。\BibleRef{格后 11:24} \BibleQuote{我五次从犹太人那里受了四十下鞭打，减去一下；三次被棍打，一次被石头砸，三次遭遇船难，一昼夜在深海里。}\par

但恐怕我们也跌入无底深渊，被远远带走，逐一细数他的每个肢体；让我们离开身体，看看另一种美，即来自他衣裳的美；就连魔鬼也对它表示尊敬；因此它们都逃走了，疾病也消散了。无论保禄出现在哪里，它们都退避让路，仿佛全世界的勇士出现了。正如那些在战争中多次受伤的人，若只是看到那伤者的一部分盔甲，也会战栗；同样，魔鬼也只看到\Quote{手帕}就惊惶失措。如今富人在哪里？那些对财富有高傲想法的人在哪里？那些数算自己的头衔和昂贵袍服的人在哪里？若他们与这些事相比，他们自己的一切在他们眼中将如泥土和灰尘。我为什么提到衣裳和金饰？如果一人把全世界都赐给我，我却认为保禄的一片指甲比那全权柄更强：他的贫穷比一切奢华更强：他的羞辱比一切荣耀更强：他的赤身比一切财富更强：没有安全感能与那神圣头颅的击打相比：没有冠冕能与作为他目标的石头相比。亲爱的，让我们渴望这冠冕，若现在没有迫害，让我们在此期间预备自己。因为我们所谈论的那位并不仅因迫害而光荣：他也说过，\BibleRef{格前 9:27} 修订文本 \OriginalTerm{}{\GreekText{ὑπωπίαζω}} \BibleQuote{我克制我的身体}；现在没有迫害，人也能在这方面达到卓越。他还劝戒不要\BibleRef{罗 3:14} \BibleQuote{为肉体安排，去满足私欲}。又说，\BibleRef{弟前 6:8} \BibleQuote{只要有食物和衣服，我们就当知足}。因为这些目的，我们不需要迫害。他也试图节制富人，说：\BibleRef{弟前 6:9} \BibleQuote{那些想要致富的人，就陷于诱惑}。\par

因此，如果我们这样操练自己，当我们进入竞赛时，我们将得冠冕；即使没有迫害摆在面前，我们也将因这些事获得许多赏报。但如果我们纵容身体，过着猪一样的生活，即使在平安中，我们也常常犯罪并蒙羞。\par

难道你们看不见我们与谁角力吗？是与无形的权势。既然我们自己是有肉身的，又怎能胜过它们呢？因为与人格斗尚需节制饮食，何况与恶神角力呢？但当我们既肉体丰满又被财富束缚时，我们怎能战胜对手？因为财富是一条锁链，是一条沉重的锁链，对于那些不知道如何使用它的人；是一个野蛮而不人道的暴君，以凌辱的方式对服侍它的人发号施令。然而，如果我们愿意，我们就能废除这苦毒的暴政，使它服从我们，而不是发号施令。这如何做到呢？通过将我们的财富分施给众人。因为只要它单枪匹马地反对我们，它就像旷野中的强盗一样实现其所有恶果；但当我们把它带到众人面前，它就不再统治我们，因为它将被所有人从四面八方用锁链捆住。\par

我说这些，并非因为财富是罪：罪在于不分施给穷人，在于滥用财富。因为天主所造的一切都是好的，且都非常好；所以财富也是好的；即如果它不主宰其主人，如果它能消除邻人的匮乏。因为那光并非好，如果它不仅不驱散黑暗，反而使黑暗更浓；我也不称之为财富，如果它非但不能消除贫穷，反而增加贫穷。因为富人不是要从别人那里取，而是要帮助别人；但那些寻求从别人那里获取的，就不再是富人，而是极度贫穷。所以，财富不是恶，而是那将财富变成贫穷的匮乏心态。这些人比那些在窄巷中乞讨、背着行囊、身体残缺的人更可怜。我说，他们虽衣衫褴褛，却不像那些穿着绸缎发光衣服的人那样悲惨。那些在市场昂首阔步的人，比那些在路口徘徊、进入庭院、在地窖里呼喊、求施舍的人更可怜。因为这些人确实向天主发出赞美，说出慈悲和严格道德的话语。因此我们怜悯他们，伸手相助，从不指责他们。但那些恶意致富的人，残暴、不人道、贪婪和撒旦般的欲望，都在他们吐出的言语中。因此所有人都厌恶嘲笑他们。你只须考虑：在所有人中，这两者谁更可耻，是向富人乞讨还是向穷人乞讨？我想每个人都立刻明白：是向穷人乞讨。如今，如果你留意，富人正是这样做的；因为他们不敢向比自己更富有的人求助；而那些乞讨的人却向富人乞讨：因为一个乞丐不会向另一个乞丐求施舍，而是向富人求；但富人却撕碎穷人。\par

再者，请告诉我，哪一样更体面：从那些心甘情愿且欠你人情的人那里接受，还是当人们不情愿时，强迫和纠缠他们？显然是不去麻烦那些不情愿的人。但富人也这样做：因为穷人从心甘情愿且欠他们人情的人手中接受；而富人却从不情愿且抗拒的人手中接受，这表明他们更加贫穷。因为如果没有人喜欢去赴宴，除非邀请者觉得欠了客人的人情，那么通过强迫来分享财产，又如何能算得上光荣呢？我们难道不正是因为狗纠缠不休，几乎逼迫我们，才躲避它们并逃离它们的吠叫吗？我们的富人也是如此。\par

\Quote{但是，让恐惧伴随礼物，这才是更体面的。}不，这恰恰是最可耻的。因为那个为获得利益而翻天覆地的人，还有谁能像他那样受人嘲笑呢？因为即便是对狗，我们也常常出于恐惧而扔掉手中所有之物。我再次请问，哪一样更可耻？是一个衣衫褴褛的人乞讨，还是一个身穿丝绸的人乞讨？因此，当一个富人向年老贫穷的人献殷勤，好取得他们的财产，而且是在他们有子女的情况下，他又能配得什么宽恕呢？\par

此外：如果你愿意，让我们细查这些话本身：富有的乞丐说些什么，穷苦的乞丐又说些什么。那么穷苦人说什么呢？他说，施舍之人永不会按量施舍（\OriginalTerm{限制}{\GreekText{μετριάσει}}，也许有误：建议作\OriginalTerm{饥饿}{\GreekText{πεινάσει}}，即永不饥饿）；他给的是属于天主的；天主是爱人者，并且更丰盛地回报；这一切都是高尚道德、劝诫和劝告的话。因为他劝你仰望主，并除去你对未来贫穷的恐惧。我们可以在求施舍者的话中看到许多教训；但富人说的话又是何种呢？唉，是猪、狗、狼和一切野兽的话。因为他们中有些人喋喋不休地谈论宴席、菜肴、美味、各种酒、香膏、服饰以及所有其余奢侈之物。另一些人则谈论利息和贷款。他们编造账目，将债务堆得难以忍受，仿佛是从祖辈或父辈时开始累积的，他们抢夺一个人的房屋，另一个人的田地，另一个人的奴隶，以及他的一切。何须提及他们的遗嘱呢？那些是用血而不是用墨水写成的。因为他们或是用某种无法忍受的危险来包围人，或是用些微不足道的承诺来蛊惑人，只要他们看到谁拥有一些微小的财产，就说服那人越过所有的亲属——往往当那些亲属正因贫穷而垂死时——而将自己的名字写在遗嘱上。任何野兽的疯狂和凶残，有什么不被这些事情所超越呢？\par

8. 因此，我恳求你们，让我们逃避所有这类的财富，因为它既可耻又充满死亡；让我们获得那属神的财富，让我们寻求天上的宝藏。因为谁拥有这些，谁就是富人，就是富翁，今世和来世都享受一切，甚至万有。因为谁若按天主的圣言成为贫穷的，所有人家的大门都会向他敞开。因为那为天主的缘故而不再拥有任何事物的人，每个人都会从他自己的所有中分给他。但谁若以不义的方式持有少许，就关闭了所有人对他的大门。为此，为使我们能获得今世和来世的美善，让我们选择那不可挪走的财富，那不朽的丰盈：愿天主藉着我们的主耶稣基督的恩宠和慈爱，赐予我们众人获得这一切。\par

\SourceRecord{Translated by Talbot W. Chambers.；Nicene and Post-Nicene Fathers, First Series；Vol. 12.；Edited by Philip Schaff.；Buffalo, NY: Christian Literature Publishing Co.,；1889.；<http://www.newadvent.org/fathers/220113.htm>.}

% structure: p0001 source=h1 target=section reason=page-title

\section{格林多前书第十四篇讲道}

% structure: p0002 source=h2 target=subsection reason=relative-heading-rank; base=h2

\subsection{格林多前书 4:17}

\BibleQuote{为此，我打发弟茂德到你们那里去，他是主内我所亲爱又忠信的儿子，他必提醒你们我在基督耶稣内所走的道路。}\par

也请思考，我恳求你们，这高贵的灵魂，比火更炽热、更锐利的灵魂：他是多么渴望亲自与格林多人在一起，他们如此混乱和分裂。因为他深知，他的临在对门徒是多么大的帮助，他的不在又是多么大的损失。前者他在致\BibleBook{斐理伯}{人书}中宣称，说：\BibleRef{斐 2:12}（公认文本中省略）\Quote{你们不仅要在我临在时，而且更要在我不在时，以恐惧战栗，努力成就你们的救恩。} 后者他在本书信中表明，说：\EditorialNote{78 18}\Quote{有些人以为我不去你们那里，就傲慢自大；但我将要来。} 他似乎是迫切的，渴望亲自到场。但当时这不可能，于是他以自己将要到来的许诺来纠正他们；不仅如此，还通过派遣他的门徒。他说：\Quote{为此，}他又说：\Quote{我打发弟茂德到你们那里去。} \Quote{为此：}是什么意思？\Quote{因为我关心你们如同子女，并如同生了你们一样。} 并且这信息附带着对他本人的推荐：\Quote{他是主内我所亲爱又忠信的儿子。} 他说这话，既是为了显示对他的爱，也是为了预备他们以尊敬的态度看待他。而且不只是说\Quote{忠信，}而是说\Quote{在主内；}也就是说，在属于主的事上。如果在世俗的事上，一个人被称为忠信已是极高的称赞，那么在属神的事上更是如此。\par

如果他是他的\Quote{亲爱的儿子，}试想\Person{保禄}的爱何其大，他为了格林多人的缘故而选择与他分离。若他也是\Quote{忠信}的，那么他在办理他们的事务时将是无可指摘的。\par

\Quote{他必提醒你们。}他没有说\Quote{要教导，}以免他们因此不快，因为他们习惯从他本人学习。因此他在结尾也说，\BibleRef{格前 16:10}\Quote{因为他做主的工作，如同我一样。所以不要有人轻视他。}因为宗徒们之间没有嫉妒，只着眼于一件事：教会的建立。如果被派遣的人比他们低下，他们仿佛\OriginalTerm{支持}{\GreekText{συνεκρότουν}}他，极其热忱。因此他不满足于说\Quote{他必提醒你们，}而为了更彻底地铲除他们的嫉妒——因为弟茂德年轻——他带着这个意图，加上说：\Quote{我的道路，}不是\Quote{他的，}而是\Quote{我的；}也就是说，\OriginalTerm{他的方法}{\GreekText{τὰς} \GreekText{οἰκονομίας}}、他的危险、他的习惯、他的规矩、他的法令、他的宗徒教规，以及一切其余的。因为他曾说过：\Quote{我们赤身露体，且受殴打，没有一定的住所：这一切，}他说，\Quote{他都要提醒你们；}还有基督的法律，以消灭一切异端。然后，他将论证提升，加上：\Quote{那是在基督内的，}按照他的习惯，将一切都归于主，并据此建立后面内容的可信性。因此他接着说：\Quote{正如我在各教会各处所教导的。} \Quote{我没有向你们说什么新奇的事：我的这些做法，其他各教会也和你们一样知晓。}此外，他称它们为\Quote{在基督内的道路}，表明其中没有任何属人的成分，并且依靠那源头的帮助，他凡事做得完善。\par

2. 说了这些话并如此安抚了他们之后，他正要开始指控那个不洁的人，又说出充满怒气的话；并非他内心真的如此，而是为了纠正他们。他放弃了对那个淫乱者的直接指责，转向其余的人，仿佛认为那人甚至不值得自己对他说话；正如我们对待大大冒犯我们的仆人时所做的那样。\par

接着，在他说了\Quote{我打发弟茂德}之后，免得他们因此掉以轻心，请注意他说的话：\par

% structure: p0009 source=h2 target=subsection reason=relative-heading-rank; base=h2

\subsection{格林多前书 4:18}

\Quote{有些人以为我不去你们那里，就傲慢自大。}因为这里他既针对他们，也针对某些别的人，打击他们的自高自大：因为当人们滥用教师不在的机会而任性妄为时，根源就在于好居首位。因为当他向民众说话时，请注意他是如何以激发羞愧感的方式进行的；当向肇事者说话时，他的态度更为激烈。因此他对前者说：\Quote{我们成了万物中的渣滓，}并安抚他们说：\Quote{我写这些话，不是为羞辱你们；}但对后者说：\Quote{现在有些人以为我不去你们那里，就傲慢自大，}表明他们的任性显露出幼稚的心态。因为正如孩童在教师不在时变得更加疏忽。\par

这里所表明的是一件事；另一件事是他的临在足以纠正他们。因为正如狮子的出现使所有活物畏缩，同样\Person{保禄}的出现也使教会的败坏者畏缩。\par

% structure: p0012 source=h2 target=subsection reason=relative-heading-rank; base=h2

\subsection{格林多前书 4:19}

因此他接着说：\Quote{但若主愿意，不久就要到你们那里去。}仅仅这样说似乎只是威胁。但既应许自己来，又要求他们以行动提供必要的证明，这才是真正高尚精神的举措。因此他也补充说：\par

\BibleQuote{并且我要知道的，不是那些自高自大者的话语，乃是能力。}因为他们这种任性并非出于自身的任何优越，而是由于他们的老师不在。这本身又是他们对他心怀轻蔑的标记。为此，他说了\BibleQuote{我已经打发弟茂德去}之后，并没有立即加上\BibleQuote{我要来}，而是等到他对他们提出\BibleQuote{自高自大}的控诉之后，才说\BibleQuote{我要来}。因为，如果他在这控诉之前就提出来，那倒更像是为自己没有失职而作的辩解，而不是威胁；即使那样（\OriginalTerm{如此}{\GreekText{οὕτως}}：王家抄本如此；公认经文为\OriginalTerm{这}{\GreekText{οῦτος}}），说法也不会令人信服。而现在，他将这话放在控诉之后，便使自己成为他们既相信又畏惧的人。\par

也要注意他如何使自己的立场稳固而安全：因为他并不简单地说\BibleQuote{我要来}，而是说\BibleQuote{主若愿意}，并且不指定明确的时间。因为他可能来迟，而这种不确定性正可以让他们保持热切的期待。而且，恐怕他们因此又松懈，他又加上\BibleQuote{快}，\par

3. \BibleQuote{并且我要知道的，不是那些自高自大者的话语，乃是能力。}他没有说\BibleQuote{我要知道的不是智慧，也不是神迹}，而是什么？\BibleQuote{不是话语}：他通过所用的词语，一方面贬低前者，另一方面抬高后者。此时他是在针对那些袒护淫乱者的多数人。因为如果他是论及那人，他不会说\BibleQuote{能力}，而是\BibleQuote{行为}，即他所行的败坏之事。\par

那么，你为何不追求\BibleQuote{话语}？\BibleQuote{不是因为我在话语上缺乏，而是因为我们的一切作为都'在乎能力'。}因此，正如在战争中，胜利不属于那些夸夸其谈者，而属于那些大有作为者；同样，在这里，得胜的不是说话的人，而是行事的人。\BibleQuote{你，}他说，\BibleQuote{为这漂亮的言词而自夸。好吧，如果这是一个演讲者的竞赛和时机，你为此自鸣得意尚有道理；但如果是一个宣讲真理、并以神迹证实真理的宗徒的竞赛，你何必为一件多余、虚幻、且对当前目的毫无用处的事而自高自大呢？因为夸夸其谈对于复活死人、驱逐邪灵、或行任何其他奇事有何用处？但这些正是我们现在所需要的，我们的立场也由此而立。}因此他又说：\par

% structure: p0018 source=h2 target=subsection reason=relative-heading-rank; base=h2

\subsection{\BibleRef{格前 4:20}}

\BibleQuote{因为天主的国不在乎言语，乃在乎能力。}他说，我们得胜不是靠漂亮的言词，而是靠神迹；并且我们的教导是神圣的，确实宣告天国的来临，我们给出更大的明证，即我们借着圣神的能力所行的神迹。如果那些现在自高自大的人想要成为大人物，那么我一来，就让他们看看他们是否有这样的能力。并且不要让我发现他们躲在夸夸其谈的背后，因为那种技艺对我们毫无用处。\par

% structure: p0020 source=h2 target=subsection reason=relative-heading-rank; base=h2

\subsection{\BibleRef{格前 4:21}}

4. \BibleQuote{你们愿意怎样呢？要我带着棍棒到你们那里去呢，还是要我存着慈爱和温柔的心？}\par

这话既有极大的威吓，又有极大的温柔。因为说\BibleQuote{我要知道}，是自持之词；而说\BibleQuote{你们愿意怎样？要我带着棍棒到你们那里去吗？}则是登上教师的座位，从那里向他们讲话，并拿起他全部权柄的话。\par

\BibleQuote{带着棍棒}是什么意思？带着惩罚、带着报复：即我要毁灭，我要使他们瞎眼——就像伯多禄在撒斐辣身上所做的，以及他自己在厄里玛斯术士身上所做的。因为此后他不再将自己与别的教师作仔细比较来说话，而是带着权柄。在第二封书信中，他好像也说了同样的话，当他写道：\BibleQuote{既然你们寻求基督在我内说话的凭据。}\par

\BibleQuote{要我带着棍棒来，还是存着爱？}既然如此，带着棍棒来难道不是爱的表现吗？当然是爱的表现。但因着他极大的爱，他在惩罚时退缩，所以他这样表达自己。\par

再者，当他论及惩罚时，他没有说\BibleQuote{在温柔的心之中}，而只是说\BibleQuote{带着棍棒}；然而惩罚也是圣神所赐的。因为有温柔之神的恩赐，也有严厉之神的恩赐。但他不如此称呼它，而是从其较和善的一面（\OriginalTerm{从较和善的一面}{\GreekText{ἀπὸ} \GreekText{τῶν} \GreekText{χρηστοτέρων}}）来称呼。同样，天主虽然报仇雪恨，却常常被宣告为\BibleQuote{仁慈、宽仁、富于慈爱和怜悯}；但祂易于惩罚的说法，或许只出现一次或两次，且很节制，只在某些紧急原因下。\par

5. 那么请看保禄的智慧：他将权柄握在自己手中，却把选择和权柄都留给别人，说：\BibleQuote{你们愿意怎样？}\BibleQuote{事情由你们决定。}\par

因为我们也有这两方面的选择：既可能堕入地狱，也可能获得天国；因为天主如此愿意。因为他说：\BibleQuote{看，火和水：你愿往哪一方面，就伸手吧。}\BibleRef{德 15:16}又：\BibleQuote{如果你们愿意，且听从我，就必吃地上的美物；\BibleRef{依 1:19}但若你们不愿意，刀剑必吞灭你们。}\par

但或许有人会说：\Quote{我愿意（没有人愚昧到不愿意的地步）；但愿意对我来说还不够。}不，如果你真正愿意，并做出愿意者所行的事，那就够了。但事实是，你并非真正愿意。\par

让我们也在其他事上试试看，如果这似乎好的话。因为请告诉我，想要结婚的人，仅仅希望就满足了吗？绝不；他反而要寻找为他提亲的女人，请求朋友与他一同看守，并筹集钱财。同样，商人并不满足于坐在家里空想，而是先租船，再挑选水手和划桨者，然后借钱付利息，并打听市场和货物的价格。那么，人们对于世俗之事表现得如此认真，但当他们要为了天堂冒险时，却只满足于愿望，这难道不奇怪吗？还不如说，即使在这方面，他们也没有真正表现出认真。因为一个正当渴望某事的人，也会动手去获取通向他所欲望目标的手段。因此，当饥饿迫使你进食时，你不会等待食物自己送到你面前，而是不遗余力地搜集食物。同样，对于解渴、御寒以及其他类似的事，你都很勤奋，并适时地为照顾身体做准备。现在，也这样对待天主的国吧，你必将得到它。\par

为此，天主使你成为一个自由的主体，免得你日后控告天主，仿佛某种必然性束缚了你；而你，恰恰在那些你受尊荣的事上，竟然抱怨。\par

因为我确实常听人说：\Quote{但祂那时为何使我的善行取决于我呢？} 不，祂怎能把你这个昏睡、贪爱一切不义、生活奢侈、自我放纵的人带到天堂呢？如果祂这样做了，你也不会戒除恶习。因为如果现在，即使面临威胁，你也不肯远离你的邪恶；假如祂把天堂作为你竞赛的终点，那还要到什么时候你才会停止变得更粗心、更糟糕呢？\OriginalTerm{更糟糕得多。许多}{\GreekText{χείρων} \GreekText{πολλῷ}. \GreekText{πολλῶν}}\par

你也不能辩解说：祂确实指示了我什么是善，却没有提供帮助，因为祂对你的帮助也有丰富的应许。\par

6. 你说：\Quote{但德行是沉重而令人不快的；而恶习却混合着极大的快乐；前者宽而广，后者窄而狭。}\par

那么，请告诉我，它们分别从头到尾都是这样，还是仅仅起初如此？因为事实上，你在此所说的话，并非有意，却是在为德行辩护；真理是如此有力。因为假设有两条路，一条通向火炉，另一条通向乐园；一条通往火炉的路是宽的，另一条通往乐园的路是窄的；你会优先选择哪条路？尽管你现在为了反驳而辩驳，但那些显然被所有人认可的事，即使你再无耻，也无法反驳。现在，我将努力从显而易见的事中教导你，那样一条开端艰难、结局却不难的路更值得选择。如果你愿意，让我们先以技艺为例。这些技艺开端充满辛劳，但结局却是有利可图的。你说：\Quote{但没有人会在没有强迫的情况下学习一门技艺；因为，}你补充说，\Quote{只要男孩是自己作主，他宁愿起初偷懒，最终忍受多大的恶，也不愿起初辛苦，后来收获辛劳的果实。} 那么，做出这样的选择，出自\OriginalTerm{孤儿般的心灵}{\GreekText{ὀρφανικῆς} \GreekText{διανοίας}}和幼稚的懒惰；而相反的选择则出自理智和刚毅。对我们也是如此：如果我们不是心智上的儿童，我们就不会像上述那个被遗弃的\OriginalTerm{孤儿}{\GreekText{ὀρφάνῳ}}一样，无忧无虑；而像有父亲的人一样。因此，我们必须摒弃我们幼稚的心智，不要归咎于事物本身；我们必须为我们的良心设置一位车夫，他不允许我们放纵欲望，而是使我们奔跑并奋力拼搏。因为，用起初的痛苦使我们的孩子习惯于那些开端艰辛但结局美好愉快的追求，而我们在属灵之事上却恰恰相反，这难道不是荒谬吗？\par

然而，即使在那些世俗事物中，结局会美好愉快也并不完全确定：因为此前，夭折、贫困、诬告、命运逆转或其他许多此类事情，使人在长期辛劳之后被剥夺了一切成果。此外，从事这些追求的人，即使成功，他们收获的也不是什么大益处。因为一切都会随着今生而消失。但在这里，我们的竞赛不是为了这些无果而短暂的事物，我们对结局也没有恐惧；而是在我们离开此地之后，我们的希望更大且更确定。那么，对于那些不肯为德行之苦而奋力拼搏的人，还有什么宽恕可言，什么借口可言呢？\par

他们还在问：\Quote{为什么道路是窄的？} 怎么，你不认为任何奸淫者、放荡者或\OriginalTerm{酗酒者}{\GreekText{και} \GreekText{τῶν} \GreekText{μεθυόντων}}进入地上君王的宫廷是正当的；而你却要求人们带着放荡、奢侈、酗酒、贪婪以及各种不义进入天堂本身吗？这些事怎能可原谅呢？\par

7. \Quote{不,} 你回答说，我并非说那个，但为何美德没有一条 \Quote{宽路？} 事实上，如果我们愿意，它的路是很容易的。因为请告诉我，哪一样更容易？是凿穿墙壁夺取他人的财物，然后被投入监狱？还是满足于你所有的，并免除一切恐惧？但我还没有说完。因为请告诉我，哪一样更容易？是偷取所有人的财物，在短时间内享受其中很少一部分，然后永远受折磨和鞭打？还是安于正义的贫穷一小段时间，之后永远活在喜乐中？（因为我们暂且不问哪一样更有利，只问现在哪一样更容易。）再者，哪一样更愉快？是做了一个好梦却在现实中受惩罚，还是做了一个不愉快的梦后真正享受快乐？当然是后者。请告诉我，那么你为何称美德为严苛的？我承认，与我们的疏忽相比，它是严苛的。然而，它确实是容易而平坦的，请听基督所说的话：\BibleRef{玛 11:30} \BibleQuote{我的轭是容易的，我的担子是轻省的。} 但如果你感觉不到轻省，显然是因为缺乏勇敢的热忱；因为哪里有热忱，即使沉重的事也变得轻省；同样，没有热忱，即使轻省的事也变得沉重。因为请告诉我，有什么比玛纳的宴席更甜美、更容易获得？然而犹太人不满足，尽管享受着如此愉悦的食物。有什么比饥饿和保禄所忍受的一切其他艰辛更痛苦？然而他雀跃、喜乐，并说：\BibleRef{哥 1:24} \BibleQuote{现在我在苦难中喜乐。} 那么原因是什么？是心态的不同。如果你按照当有的方式调整心态，你就会看到美德的容易。\par

\Quote{那么,} 你说, \Quote{美德是否只通过追求她的人的心态而成为这样？} 她不仅由于他们的心态，而且本性如此。我现在这样证明：如果一种事物始终是痛苦的，另一种始终是相反的，那么某些堕落的人或许会有些可信地说，后者比前者更容易。但如果它们各有开端，一个始于艰辛，一个始于快乐，而各自的结果又恰恰相反；而且这些结果都是无限的，一个带来快乐，另一个带来负担；请告诉我，哪一个更容易选择？\par

\Quote{那么为什么许多人并不选择那容易的呢？} 因为有些人不相信；另一些人虽相信，但判断力败坏，宁可选择短暂的快乐，而不选择永恒的。 \Quote{这难道不容易吗？} 不：这是由于灵魂的病态。正如发烧的人渴望冷饮，并非经过计算认为一时的痛快比从头到尾被烧灼更愉快，而是因为他们无法克制过度的欲望；这些人也是如此。因为如果有人在他们享乐的那一刻就把他们带到惩罚面前，他们肯定永远不会选择它。这样你就看到，恶习在何种意义上不是容易的事。\par

8. 但如果你愿意，让我们在适当的具体事例中再次尝试这一点。请告诉我，举例来说，哪一样更愉快、更容易？（只是我们不要再以多数人的欲望作为这事的标准；因为人应当以健康者而不是病人来判定；正如你可以给我看一万个发烧的人，他们选择不健康的东西，然后为此受苦；但我不会认可这种选择；）我重复一遍，哪一样带来更多的轻松，请告诉我；是渴望大量财富，还是超越那种欲望？就我而言，我认为是后者。如果你不相信，让论证回到事实本身。\par

那么让我们假设一个人渴望许多，另一人什么都不渴望。请告诉我，哪一种状态更好、更可敬？然而，暂且不谈这一点。因为这一点是公认的，即后者比前者有更优质的品格。我们现在并不探究这一点，而是探究谁过着更轻松、更愉快的生活？那么：爱财者连他所拥有的也不能享受：因为他所爱的，他不能选择去花费；甚至宁愿将自己\OriginalTerm{卡特科普西埃}{\GreekText{κατακόψειε}}，舍弃自己的血肉，也不愿舍弃黄金。但那轻视财富的人，同时获得了这样的好处：他安静且极其安全地享受他所拥有的，而且他把自己看得比财富更贵重。那么哪一样更愉快？是自由地享受自己所有的，还是活在主人之下，即财富之下，甚至不敢碰自己的一件东西？啊，在我看来，这很像两个人，各自拥有妻子并极其爱她们，但状况不同：一个人被允许与妻子同住同处，另一个人甚至不被允许靠近她。\par

还有另一件事我想提一提，它显示了一个人的愉悦和另一个人的不适。贪婪者永远不会止息他的贪欲，不仅因为他不可能获得所有人的财物，而且因为他无论得到多少，都认为自己一无所有。但轻视财富的人会认为一切都是多余的，不会用无尽的欲望折磨自己的灵魂。我说\Quote{折磨}，因为没有比无法满足的欲望更符合惩罚的定义了；而这尤其标志着他扭曲的心智。你从这角度看看：贪图财富并增加自己积蓄的人，他的感觉就像是一无所有。那么我问，还有什么比这种病更复杂的？奇怪的不止于此，而是他虽然拥有，却认为自己没有手中所握有的东西，仿佛没有一样，他为自己哀叹。即使他得到了所有人的财物，他的痛苦反而更大。如果获得了一百塔冷通，他因没有收到一千而烦恼；如果收到一千，他又因不是一万而刺痛；如果收到一万，他因不是十倍而\OriginalTerm{痛苦不已}{\GreekText{κατακόπτεται}}。获得更多对他来说反而变成了更大的贫穷；因为他得到的越多，他贪图的就越多。所以，他得到的越多，就变得越穷：因为贪图更多的人，实际上是更穷的人。那么当他有一百塔冷通时，难道不穷吗？因为他贪图一千。当他有了一千时，他变得更穷。因为不再是一千，而是他声称需要一万了。如果你说愿望与得不到是愉悦，那你似乎对愉悦的本质一无所知。\par

9. 为了表明这种事不是愉悦而是惩罚，我们再举一个例子，以此探究。当我们口渴时，我们难道不是因为解渴而在喝水时感到愉悦吗？喝水之所以是愉悦，难道不是因为它在很大程度上解除了我们的一种巨大折磨，即我想说的口渴的欲望？我想每个人都能明白。但如果我们一直处于这种欲望状态，我们就会像比喻中拉匝禄的富人在惩罚方面那样糟糕；因为他的惩罚正是：他强烈渴望一滴水，却得不到。我看所有贪婪者似乎都不断遭受这种痛苦，并与他相似：他乞求得一滴水，却得不到。他们的灵魂比他的更灼热。\par

有人说得对，所有贪爱钱财的人都像是患了水肿病；因为正如他们体内大量积水，反而更加灼热一样，贪婪者拥有巨大财富，却更贪求更多。原因是，前者不能把水保持在身体应有的部位，后者则不能把欲望保持在合理思想的界限内。\par

让我们逃离这种\OriginalTerm{奇怪而空洞}{\GreekText{ξένην} \GreekText{καὶ} \GreekText{κενὴν}}的疾病吧；让我们逃离万恶之根吧；让我们逃离现世的地狱吧；因为对这一切的贪欲就是地狱。你只需敞开每个人的灵魂——轻视财富的人和不轻视的人——你就会看到，前者如同精神错乱者，既不愿听也不愿看任何事物；后者则如同风平浪静的港湾。他是众人的朋友，而前者是众人的敌人。因为如果有人拿了他的东西，他不会烦恼；反之，如果有人给他东西，也不会使他骄傲；他身上有一种完全稳妥的自由。前者不得不向所有人谄媚和伪装；后者则不必向任何人如此。\par

如果贪爱钱财既是贫穷、胆怯、虚伪、假冒为善，又充满恐惧、巨大的惩罚痛苦和责罚；而轻视财富的人却享有相反的一切快乐，那么美德不是更令人愉悦吗？\par

我们本可以也详述其他一切恶事，以表明没有一种恶行具有愉悦可言，但之前我们已经说得够多了。\par

因此，让我们认识这些事，选择美德；这样我们既能在此世享受这样的快乐，也能获得未来的福乐，藉着天主的恩宠和慈爱，等等。\par

\SourceRecord{Translated by Talbot W. Chambers.；Nicene and Post-Nicene Fathers, First Series；Vol. 12.；Edited by Philip Schaff.；Buffalo, NY: Christian Literature Publishing Co.,；1889.；<http://www.newadvent.org/fathers/220114.htm>.}

% structure: p0001 source=h1 target=section reason=page-title

\section{\WorkTitle{格林多前书讲道集第十五篇}}

% structure: p0002 source=h2 target=subsection reason=relative-heading-rank; base=h2

\subsection{格林多前书 5:1-2}

\BibleQuote{确实听说你们中间有淫乱，而且是那样的淫乱，连外邦人中也没有的，就是有人收了他的父亲之妻。你们还自高自大，而不哀痛，好把行这事的人从你们中间除去。}\par

\BibleRef{格前 1:11}\BibleQuote{因为革来氏家里的人曾对我提起你们的事，说你们中间有纷争。}但在这里并非如此；他立即全身投入，使指责的谴责尽可能广泛。因为他没有说，\Quote{为什么这样的人犯了奸淫？}而是说，\Quote{听说你们中间有淫乱；}这样，他们作为完全超出其指控的人可以轻松接受；但却充满自然的焦虑，因为整个身体受伤，教会蒙受羞辱。他说：\Quote{因为没有人会这样说，\InnerQuote{这样的人犯了奸淫}，而是\InnerQuote{在格林多教会中，那罪已被犯下。}}\par

并且他没有说，\Quote{淫乱发生了，}而是说，\Quote{听说——连外邦人中也没有的。}因为他一贯利用外邦人作为谴责信徒的话题。这样，在给得撒洛尼人的信中，他说，\BibleRef{得前 4:4}省略了\OriginalTerm{其余的}{\GreekText{τὰ} \GreekText{λοιπὰ}}插入说：\BibleQuote{各人应以圣洁和尊贵保持自己的器皿，不可放纵情欲，如同其余的外邦人。}给哥罗森人和厄弗所人的信中，\BibleRef{弗 4:17}说：\BibleQuote{你们不要再像其他外邦人那样行事。}如果他们在同样的事上犯罪是不可原谅的，甚至超过了外邦人，我们还能为他们找到什么位置？告诉我：\BibleQuote{因为在外邦人中，}他这样说，不但他们不敢做这种事，甚至连名称都不提。你看到他是如何加重他的指控吗？因为当他们被证实发明了连不信者都不敢尝试、甚至不知道的种种污秽时，这罪必然极大，非言语所能形容。而\Quote{在你们中间}这句话也说得很有分量；也就是说，\Quote{在你们中间，蒙受了如此崇高奥迹的信徒，分享奥秘、被邀请到天上的人。}你看到他的言语充满何等的愤慨？他对所有人的怒气何等之大？因为若不是他心中充满极大的愤怒，若不是他树立自己反对他们所有人，他大概会这样说：\Quote{听说某某人犯了奸淫，我命你们惩罚他。}但事实并非如此；他反而立即挑战所有人。确实，如果他们先写信，他大概会这样说。既然他们非但没有写信，甚至还遮掩了过错，因此他更加激烈地安排了这篇讲论。\par

2. \Quote{你们中间有人收了他父亲的妻。}为什么他不说，\Quote{他玷污了他父亲的妻？}这事的极端污秽使他退缩。因此他带着一种谨慎匆匆带过，好像之前已经明确提及似的。借此他又加重了指控，暗示他们中间竟敢做这样的事，以至于连保禄都无法直言。因此，随着他的讲述，他用同样的措辞说，\Quote{行了这事的人：}并再次羞愧脸红，不敢说出来；这也是我们在对待极其可耻之事时的习惯。他没有说，\Quote{他的继母，}而是说，\Quote{他父亲的妻；}以便更加严厉地打击。因为当词语本身就足以传达指控时，他就简单使用它们，不加任何东西。\par

并且他说，\Quote{不要告诉我，}他说，\Quote{那犯奸淫的只是一个人：指控已变成全体共同的。}因此他立刻加上，\Quote{你们还自高自大：}他没有说，\Quote{因着那罪；}因为这暗示完全缺乏理性：而是因着你们从那人所听的道理。然而他自己并没有写明，而是留着不确定，以便给予更沉重的一击。\par

请注意保禄的明智。他首先推翻了外在的智慧，表明它本身毫无价值，即使没有罪恶与之相连；然后，直到那时，他才论及那罪。因为如果通过与那或许有智慧的犯奸淫者比较，他坚持自己神恩的伟大；他并没有做什么了不起的事：但即使没有罪恶，也要贬低外邦人的智慧，证明它毫无价值，这才是真正表明它的极端无用。因此，正如我所说，他首先作了比较，然后才提到那人的罪。\par

而且他确实不屑与他辩论，从而表明他的极度羞辱。但向其他人说，\Quote{你们应当哭泣哀号，掩面，但现在你们却相反。}这就是下一句的意思，\Quote{你们还自高自大，而不哀痛。}\par

\Quote{我们为什么要哭泣？}有人会说。因为羞辱已蔓延到你们整个教会。\Quote{我们哭泣有什么好处？}\Quote{好叫这样的人从你们中间除去。}他甚至在这里也没有提到他的名字；更准确地说，任何地方都没有；这在一切怪诞之事中是我们惯常的做法。\par

他没有说，\Quote{你们没有反而把他赶出去，}而是，如同对待任何疾病或瘟疫一样，\Quote{需要哀痛，}他说，\Quote{和恳切的祈求，\InnerQuote{好叫他被除去。}你们应当为此祈祷，竭尽一切使他被剪除。}\par

他也不是责备他们没有向他报告，而是责备他们没有为那人的被除去而哀恸；暗示即使没有他们的老师，这件事也应该做，因为那恶行的昭彰。\par

% structure: p0013 source=h2 target=subsection reason=relative-heading-rank; base=h2

\subsection{\BibleRef{格前 5:3-5}}

3. \BibleQuote{我身体虽不在你们那里，但心神却与你们同在。}\par

请留意他的活力。他甚至不容他们等待他的亲临，也不容他们先接待他，然后再通过捆绑的判决；而是仿佛要赶在某种传染蔓延到身体其他部分之前将其清除，他急忙加以制止。因此他又补充了一句：\BibleQuote{我已判决，如同我在那里一样。}他这样说，不仅是为了促使他们宣判，不给他们另作安排的机会，也是为了吓唬他们，因为知道将要做什么并已下定决心。因为\BibleQuote{心神同在}的意思正是如此：就如厄里叟与革哈齐同在时所说：\BibleQuote{我的心岂没有同你一起去吗？}\BibleRef{列下 5:26}何等奇妙！那恩赐的能力何其伟大，竟然能使一切相聚如同一体，并使他们能知晓远方的事。\BibleQuote{我已判决，如同我在那里一样。}\par

他不容许他们有别的打算。\BibleQuote{如今我已说出我的决定，如同我在那里：不要拖延再耽搁，因为决不可有别的做法。}\par

然而，恐怕他显得过于独断，言语听起来自以为是，请看他又如何使他们成为判决的同工。因为说了\BibleQuote{我已判决，}之后，他又加上：\BibleQuote{对那个做了这事的人，因我们的主耶稣基督之名，当你们集合时，连同我的心神，并以我们主耶稣基督的能力，把这样的人交给撒殚。}\par

\BibleQuote{因我们的主耶稣基督之名}是什么意思？\BibleQuote{按照天主，} \BibleQuote{不怀任何人的偏见。}\par

不过，有些人如此读：\BibleQuote{对那个在我们的主耶稣基督之名内做了这事的人，}在这里停一下或稍作停顿，然后接续下文说：\BibleQuote{当你们集合时，连同我的心神，把这样的人交给撒殚。}他们断言这种读法的含义如下：\BibleQuote{对那个在基督之名内做了这事的人，}圣保禄说，\BibleQuote{你们把他交给撒殚；}也就是说：\BibleQuote{对那个侮辱了基督之名的人，那个在成为信徒并被称为基督信徒之后，竟敢做这样的事的人，你们把他交给撒殚。}但在我看来，前一种解释（\OriginalTerm{解释}{\GreekText{ἐκδοσις}}，意为\Quote{论述}）更为真实。\par

那么这是什么意思？\BibleQuote{当你们因主的名而集合时。}也就是说：你们因祂的名而聚集，这名把你们集合在一起。\par

\BibleQuote{连同我的心神。}他再次使自己处于他们的前列，以便当他们宣判时，他们不能以其他方式革除那犯罪者，除非如同他在场一样；也免得有人胆敢认为他可宽恕，因为知道保禄会知晓整个程序。\par

4. 进而使之更加可畏，他说：\BibleQuote{以我们主耶稣基督的能力；}意思是：要么基督能赐予你们如此恩宠，使你们有能力把他交给魔鬼；要么是祂自己与你们一起对他宣判。\par

他没有说：\BibleQuote{把}这样的人交给撒殚，而是说\BibleQuote{交出}；为他开启了悔改之门，如同将这样的人交托给一位教长。另外也说：\BibleQuote{这样的人}；他绝不肯提及他的名字。\par

\BibleQuote{为毁灭肉体。}正如蒙福的约伯所受的，但并非基于同样的理由。因为在约伯的事上是为了更光明的荣冠，而在这里是为了罪过的赦免；使他受责打，患上严重的疮或其它疾病。诚然，在其他地方他说：\BibleQuote{我们受审判，是因主惩戒我们，当我们遭受这些事时。}\BibleRef{格前 11:32}但在这里，他想要让他们更严厉地感受，于是\BibleQuote{交给撒殚}。于是天主所决定的事也发生了：那人的肉体受了惩罚。因为过度的饮食和肉体的纵欲是欲望之母，所以他责罚肉体。\par

\BibleQuote{为使灵魂在主耶稣的日子得以得救；}即心灵。并非只有心灵得救，而是因为若心灵得救，毫无疑问身体也要分享其救恩。正如因为灵魂犯罪，身体成了可朽的；所以若灵魂行义，身体也将享受极大的荣耀。\par

但有些人认为\BibleQuote{神}是指那恩赐，当我们犯罪时它便熄灭。\BibleQuote{为了防止这事发生，}他说，\BibleQuote{让他受罚；好使他借此变得更好，为自己赢得天主的恩宠，并在那日子发现它安然无恙。}所以这一切都出于一位护士或医师的职分，并非仅仅是责打或随意惩罚。因为收益大于惩罚：一个是暂时的，另一个是永恒的。\par

他没有简单地说\BibleQuote{为使灵魂得救}，而是说\BibleQuote{在那日子。}他适时地提醒他们那日子，好使他们更愿意接受治疗，也好使被责备的人不是从愤怒中，而是从忧虑父亲般的远见中领受他的话。为此他也说\BibleQuote{为毁灭肉体}，继续为魔鬼设定界限，不容它越界一步。如同在约伯的事上，天主说\BibleRef{约 2:6}：\BibleQuote{只是不可伤害他的性命。}\par

然后，在结束了判决，简单提及而不赘述之后，他又插入一段责备，指向他们：\par

% structure: p0029 source=h2 target=subsection reason=relative-heading-rank; base=h2

\subsection{\BibleRef{格前 5:6}}

\BibleQuote{你们的夸耀是不好的：}意即正是他们直到如今因以他为荣而阻止了他悔改。接着他指出，他采取这一步，不仅是为了宽恕那人，也是为了他所写信的对象。为此他加上：\par

\BibleQuote{你们不知道，一点酵母能使整个面团发酵么？} \Quote{因为，}他说，\Quote{虽然过失是他的，但若忽视它，它就能败坏教会其余的身体。因为当第一个犯过者逃脱惩罚时，别人也会迅速犯同样的过错。}\par

在这些话中，他还表明他们的挣扎和危险是为整个教会，而不是为任何人。为此他也需要酵母的比喻。因为他说：\Quote{正如酵母，}他说，\Quote{虽小，却使整个面团变成自己的本性；同样，若让这人不受惩罚，这罪不报应，他也会败坏其余所有的人。}\par

% structure: p0033 source=h2 target=subsection reason=relative-heading-rank; base=h2

\subsection{格林多前书 5:7-8}

\BibleQuote{你们要把旧酵母除净，}即这个恶人。不是单指这人；他也旁及与他一起的人。因为\Quote{旧酵母}不仅指私通，也指每种罪过。他没有说\Quote{除净}，而说\BibleQuote{你们要把……除净}：\Quote{仔细洁净，以致不留丝毫或阴影。}因此当他说\BibleQuote{你们要把……除净}时，他表明他们中间仍有不义。但他说\BibleQuote{好使你们成为新面团，正如你们是无酵的一样}时，他肯定并声明罪恶并未在太多人中得势。但他说\Quote{正如你们是无酵的一样}，并非意指所有人都洁净，而是指你们应当成为怎样的人。\par

6. \BibleQuote{因为我们的逾越节羔羊，基督，已被祭杀，所以我们过节：不可用旧酵母，也不可用恶毒和邪恶的酵母，而要用真诚和真理的无酵饼。}同样，基督也称祂的教义为酵母。并且他本人进一步发挥这个比喻，提醒他们一段古代历史、逾越节和无酵饼，以及他们那时和现在的祝福、惩罚和灾祸。\par

因此，我们生活的整个时期都是节期。因为他说\BibleQuote{我们过节}，并不是指逾越节或五旬节的临在，而是指出整个时间对基督徒而言都是节期，因为所赐美物的卓越。有什么好事没有发生呢？天主子为你成了人；祂让你摆脱死亡；并召叫你进入天国。因此，你既已获得且仍在获得如此事物，难道不应一生\BibleQuote{过节}么？那么，不要因贫穷、疾病和仇敌的诡计而沮丧。因为这是节期，甚至我们的整个时间。为此保禄说：\BibleRef{斐 4:4}\BibleQuote{你们要在主内常常喜乐；我再说，喜乐。}在节庆日，无人穿着脏衣。我们也不要这样。因为已经举办了婚宴，一场属神的婚宴。因为祂说：\Quote{天国好比一个国王，}他说，\Quote{为他的儿子摆设婚宴，\BibleRef{玛 22:1}，修订文本 \OriginalTerm{}{\GreekText{ἐποίησε}}}如今，既是国王摆设婚宴，且是为他儿子的婚宴，还有什么比这筵席更伟大的？那么，没有人穿着破衣进来。我们讲的不是衣服，而是不洁的行为。因为如果所有人都穿着亮丽，只有一人被发现穿着脏衣赴宴，而被羞辱地赶出去，那么进入那婚宴所需要的是怎样的严谨和纯洁啊！\par

7. 然而，他不仅为此提醒他们\Quote{无酵饼}，也为了指出旧约与新约的相似性，并指出在\Quote{无酵饼}之后不能再进入埃及；但若有人选择回去，他将遭受与他们同样的苦难。因为那些事是这些事的预像，无论犹太人如何顽固。因此你若问他，他说话并不伟大，而是他说的话虽伟大，却不像我们所说的：因为他不知真理。因为他会说：\Quote{扣押我们的埃及人被全能者改变，以致他们自己催促并赶我们出来，他们先前强行拘留我们；他们甚至不许我们让面团发酵。}但若有人问我，他听到的将不是埃及或法郎，而是我们脱离魔鬼的欺骗和黑暗的拯救；不是梅瑟，而是天主子；不是红海，而是充满万千祝福的圣洗圣事，在那里\Quote{旧人}被淹死。\par

再者，你若问犹太人为何从他所有的边界中驱逐一切酵母；他在这里甚至闭口不言，不说任何理由。这是因为，虽然有些情况既是将来之事的预像，也因当时的事而发生；但另一些却不是，免得犹太人诡诈行事；免得他们停留在预像中。因为告诉我，羔羊是\Quote{公的}、\Quote{无瑕疵的}、\Quote{一岁的}，以及\Quote{不可折断一根骨头}是什么意思？召叫邻居的命令，\BibleRef{出 12:4}以及要\Quote{站着}\Quote{在晚上}吃，或用血巩固房屋，又是什么意思？他将除了反复谈论埃及外无话可说。但我能告诉你血的含意、晚上的含意、一起吃的规定，以及所有人都要站立的规矩。\par

8. 但首先让我们解释为何将酵母从他们的全境内除去。隐藏的意义是什么？信徒必须摆脱一切不义。因为如同在他们中，凡发现旧酵母的人都将灭亡，同样在我们中，凡发现不义之处也是如此：既然在预像中惩罚如此之大，在我们这里必然更大。因为他们如此仔细地清除家中的酵母，甚至探察鼠穴；我们更当搜查灵魂，以除去一切不洁的思念。\par

然而，这是他们后来所做的；但如今不再如此。因为凡有犹太人的地方，到处都有酵母。因为无酵节是在城市中举行的；这如今与其说是法律，不如说是游戏。因为既然真理已经来到，预像不再有任何地位。\par

因此，藉着这个例子，他也大力地将行淫者逐出教会。因为他说，他的在场不仅无益，反而有害，损害了身体的共同状态。因为当腐败的部分被隐藏时，无人知道恶臭从何而来，于是就归咎于全体。因此，他强烈敦促他们\BibleQuote{把旧酵母除净，好使你们成为新面团，如同你们原是无酵的一样。}\par

\BibleQuote{因为我们的逾越节羔羊基督，已经被祭献了。}他没有说死了，而是更切合主题地说\BibleQuote{被祭献了}。那么，不要寻求这种无酵饼，因为你们也没有同类的羔羊。不要寻求这种酵母，因为你们的无酵饼并非如此。\par

9. 因此，在物质的酵母上，无酵可以变成有酵，但绝不会相反；而在这里，却有发生完全相反情况的机会。然而他没有明确宣告这一点：请观察他的明智。在前一封书信中，他没有给行淫者任何回来的希望，而是命令他整个生命都应在忏悔中度过，以免因应许而使他不那么热忱。因为他没有说\BibleQuote{把他交给撒殚,}以便他悔改后可以再次被接纳到教会中。而是说什么？\BibleQuote{使他在末日可以得救。}因为他带领他直到那时，使他充满焦虑。至于悔改后他预备赐予什么恩惠，他没有透露，效法他的师傅。因为正如天主所说：\BibleRef{纳 3:4}\BibleQuote{还有三天，尼尼微就要毁灭了,}并没有加上\BibleQuote{但如果她悔改，她将得救；}同样，他在这里也没有说\BibleQuote{如果他相称地悔改，我们就要\InnerQuote{向他重申我们的爱}}\BibleRef{格后 2:8}。他等待他完成这工作，以便那时他再接受恩惠。因为如果他在开始时说了这话，可能会使他从恐惧中解脱。因此，他不仅没有这样做，反而藉着酵母的例子不给他任何回归的希望，将他保留到那日：\BibleQuote{把旧酵母除净（他说）}；并\BibleQuote{我们不可用旧酵母过节。}但一旦他悔改了，他就以极大的热忱把他重新接纳进来。\par

10. 但为什么他称它为\Quote{旧}？要么因为我们先前的生活属于这种，要么因为旧的东西\BibleQuote{快要消逝}\BibleRef{希 8:13}，且无味而污秽；这正是罪的本性。因为他既不是单纯地谴责旧的，也不是单纯地赞美新的，而是相对于主题。因此，在别处他说：\BibleRef{德 9:15}\BibleQuote{新酒如同新朋友；但若变成陈酒，你就会喜欢喝它：}在友谊上，他更赞美旧的而非新的。同样，\BibleQuote{万古常存者坐下了}\BibleRef{达 7:9}，这里又取\Quote{古}作为那些赋予最高荣耀的赞美词。别处圣经以贬义取\Quote{旧}一词；因为看到事物具有多种面貌，由于由许多部分构成，它用相同的词语表达好和坏的含义，并非根据相同的意义。关于这一点，你可以看到在别处对旧的谴责：\BibleRef{咏 16:46}\BibleQuote{他们变老了，从他们的路上蹒跚而行。}又\BibleRef{咏 6:7}\BibleQuote{我在一切仇敌中变老了。}又\BibleRef{达 13:52}\BibleQuote{你这位在恶日子中变老的人。}同样，\Quote{酵母}也常被用来指天国，尽管在这里被指责。但在那里它有一种含义，在这里有另一种。\par

11. 但我强烈确信，关于酵母的这番话也指向那些司铎，他们容忍大量的旧酵母留在里面，没有从他们的境内，即从教会中，除去贪婪者、敲诈者以及一切将人排除于天国之外的事物。因为贪婪确实是\Quote{旧酵母}；无论它落在何处，进入哪家，都使之不洁：即使你因不义只获得少许，它也发酵你全部的家产。因此，不义之财虽少，却常常驱逐诚实的储备，无论后者多么丰盈。因为没有什么比贪婪更腐败。你可以用钥匙、门和闩锁住那人的橱柜：但你把贪婪这最坏的盗贼关在里面，它能够带走一切，你这是徒劳的。\par

\Quote{但是，}你说，\Quote{如果有许多贪婪的人没有经历这种事呢？}首先，他们将会经历，即使不是立即。如果他们现在逃脱了，那就更要害怕：因为他们被保留以受更大的惩罚。此外，即使他们自己逃脱了，那些继承他们财富的人也将承受同样的事。\Quote{但这怎么可能是正义的呢？}你会说。这是十分正义的。因为继承遗产的人，尽管他本人没有抢劫，却仍然扣留他人的财产；他完全知道这一点；他为此受苦是公平的。因为如果某人抢劫，而你收到了东西，然后物主来讨回，你辩解说你没有抢夺，这对你有用吗？绝不。因为当你被指控时，你的辩词是什么？告诉我。是另一个人抢夺的？好吧，但你在占有。是他抢劫的？但你正在享用。即使是外邦人的法律也承认这些规则，它们赦免了抢夺和偷窃的人，却命令你从那些你发现你的东西被存放在那里的人那里索取赔偿。\par

如果你知道谁是受害者，就归还，并像匝凯所做的那样，加上许多增加。但如果你不知道，我给你另一条路；我并不排除你的补救方法。将所有这些东西分给穷人：这样你就会减轻邪恶。\par

但若有人将这些事物传给了子女和后代，他们在报应中仍然遭受了其他灾难。\par

12. 我何必提今生的事呢？无论如何，在那一天，这些话都不会再被说，那时两者都赤裸地出现，被抢夺者和抢夺者。或者更确切地说，不是同样赤裸。的确，两者都同样被剥夺了财富；但前者将充满他们所引发的指控。那么那天我们该怎么办，当在可怕的审判台前，那受虐待并失去一切的人被带到中间，而你没有人为你说一句话？你将对法官说什么？现在你或许还能败坏判决，因为你只是人；但在那个法庭和那时，情况不再如此：不，现在你也不能。因为即使在此刻，那个审判台也是临在的：因为天主看见我们的行为，并且接近受委屈的人，即使没有被呼求；肯定的是，无论谁受委屈，即使他自己不配得到任何补偿，然而既然所做的事不中悦天主，他必定有一位复仇者。\par

\Quote{那么，}你会说，\Quote{这样一个邪恶的人怎么会过得好呢？}不，结局不会是这样。听先知所说：\BibleRef{咏 37:1}\BibleQuote{不要因作恶的人而烦恼，因为他们像草一样必要迅速枯萎。}因为请告诉我，那个施行抢夺的人，在离开今世之后，在哪里？他的光明希望在哪里？他显赫的名声在哪里？难道不是一切都过去了，消失了吗？他的一切难道不是一场梦和一个影子吗？你必须对每一个这样的人都期待这样的事，无论是本人活着的时候，还是后来者。但圣人们的状态并非如此，你也不可能对他们说同样的话，说他们的所有是影子、梦和虚构的故事。\par

13. 如果你愿意，让这个说这些话的人，那帐棚匠、基里基雅人、出身不明的人，成为我们提出的榜样。你会说：\Quote{怎么可能变得像他那样？}你真心渴望吗？你迫切想要成为那样吗？\Quote{是的，}你会说。那么，走他所走的路，以及与他同行的人所走的路。他走了什么路？一个人说：\BibleRef{格后 11:27}\BibleQuote{在饥饿、口渴和赤身露体中。}另一个人说：\BibleRef{宗 3:6}\BibleQuote{银子和金子我没有。}因此他们\BibleRef{格后 6:10}\BibleQuote{一无所有，却拥有万有。}有什么比这句话更高贵？什么比这更蒙福、更富足？其他人确实以相反的事为荣，说：\Quote{我有这么多金塔冷，无边的土地，房屋和奴隶}；但这个人以自己一无所有为荣；他不因贫穷而退缩（这是不智之人的感觉），也不遮盖自己的脸，反而将其作为装饰。\par

如今富人在哪里？那些计算单利和复利的人，那些向所有人索取却永不满足的人？你们听到了伯多禄的声音，那声音称贫穷为财富之母？那声音一无所有，却比戴王冠的人更富有？因为这就是那声音，它一无所有，却复活了死人，使瘸子站起来，驱逐了魔鬼，并赐予了如此恩惠的恩赐，以至于那些穿着紫袍、率领强大可怕军团的人从未能赐予。这是那些如今被接到天上、达到那高度的人的声音。\par

14. 因此，一无所有的人有可能拥有所有人的财物。因此，一无所获的人可以获得所有人的财物：而倘若我们得到所有人的财物，我们反而丧失一切。这话听起来似乎矛盾，但并非如此。你将会说：\Quote{但是，}你将会说，\Quote{什么都没有的人怎能拥有所有人的财物？难道不是拥有众人之物的人更富有吗？}绝不是，恰恰相反。因为一无所有的人支配一切，正如宗徒们那样。全世界所有的房屋都向他们敞开，收留他们的人视他们的到来为恩惠，他们来到众人面前如同面对朋友和亲人。因为他们来到那个卖紫布的妇女那里\BibleRef{宗 16:14}，她像仆人一样把他们所需要的摆在他们面前。又来到狱卒那里，他向他们敞开了整个家。还有无数其他人。这样，他们拥有一切又一无所有：因为\BibleRef{宗 4:32}\BibleQuote{他们说，他们所拥有的没有一样是自己的；}所以一切都是他们的。因为凡以为万物都是公有的人，不仅使用自己的，也使用他人的如同自己的。但那划分界限、只做自己东西的主人的人，连自己的东西也做不了主。这一点通过一个例子很清楚。那完全一无所有的人，既无房屋，也无桌子，也无多余的衣服，但为了天主的缘故舍弃一切，将公共之物视为己有；他将从所有人那里得到他所想要的一切，这样，一无所有的人就拥有了众人的财物。但拥有一些东西的人，连这些也做不了主。首先，没有人会给拥有财产的人；其次，他的财产将属于盗贼、小偷、告密者和变化无常的事件，属于任何人而不是他自己。例如，\Person{保禄}走遍全世界，什么也没有带，虽然他既不是去朋友也不是去亲人那里。起初，他是所有人的公敌，但在他成功地进入后，他拥有了所有人的财物。但是\Person{阿纳尼雅}和\Person{撒斐辣}，急于多得一点自己的东西，却失去了所有连同生命本身。那么，放弃你的东西，这样你才能把他人的东西当作自己的来使用。\par

15. 但我必须停下来了：我不知道自己怎么会被引到这样的高潮，对这些人说出这些话，他们认为分享哪怕一点点自己的东西也是一件大事。因此，愿我这些话是对成全的人说的。但对较为不全的人，我们可以这么说：把你所拥有的给予穷人。增加你的财富。因为，主说\BibleRef{箴 19:17}\BibleQuote{周济贫穷的，就是借给上主。}但如果你着急，不等得赏报的时候，想想那些借钱给人的人：连他们也不希望立刻得到利息；他们倒希望本金尽可能长时间地留在借贷人手里，只要还款有保障，他们不怀疑借贷人。现在也当如此。把它们留给天主，让祂加倍偿还你。不要想着在这里全数收回；因为你若在这里全数收回，那里又怎能再得到呢？正是为此，天主将它们储存在那里，因为现世生命充满腐朽。但祂甚至在此世也给予；因为，\Quote{你们寻求，}祂说，\Quote{天国，这一切都将加给你们。}\BibleRef{玛 6:33}那么，让我们向往天国，不要急于全部偿还，免得减少我们的赏报。让我们等待合适的时机。因为这些情况下的利息不是那种，而是适合献给天主的。那么，让我们多多收集这些，这样离开，好获得现世和将来的祝福；藉着我们的主\Person{耶稣基督}的恩宠和慈爱，愿光荣、权能、尊崇归于祂及圣父和圣神，从今世直到永远。阿们。\par

\SourceRecord{Translated by Talbot W. Chambers.；Nicene and Post-Nicene Fathers, First Series；Vol. 12.；Edited by Philip Schaff.；Buffalo, NY: Christian Literature Publishing Co.,；1889.；<http://www.newadvent.org/fathers/220115.htm>.}

% structure: p0001 source=h1 target=section reason=page-title

\section{格林多前书第十六篇讲道}

% structure: p0002 source=h2 target=subsection reason=relative-heading-rank; base=h2

\subsection{格林多前书 5:9-11}

\BibleQuote{我先前在信上写过，不要与淫乱的人交往，但不是指这世上所有淫乱的人，或贪婪的，或勒索的，或拜偶像的，因为这样你们就必须离开这世界；但现在我写信给你们，若有人称为弟兄却是淫乱的人，或贪婪的，或拜偶像的，或酗酒的，或辱骂的，或勒索的，这样的人你们不可与他交往，甚至连吃饭也不可。}\par

因为他曾说：\BibleQuote{你们反倒不哀痛，使那作这事的人从你们中间除去？}并且说：\BibleQuote{你们应把旧酵母除净。}他们很可能推断出他们的责任是避开所有淫乱的人：因为如果犯罪的人把他的一些邪恶传染给没有犯罪的人，那么离开那些外教人更是应当的；（因为如果因这种从朋友而来的邪恶，人不应饶恕朋友，那么对其他人更该如此；）基于这种印象，他们可能也会与希腊人中的淫乱者分开，而这件事既然不可能做到，他们会更放在心上：于是他用这种方法来纠正，说：\BibleQuote{我写信给你们不要与淫乱的人交往，但不是指这世上所有淫乱的人：}他用\Quote{完全}这个词，好像是一件公认的事。为了不让他们以为他并没有把这件事指给他们，因为这是不完善的，并且错误地认为他们是完善的而试图去做，他表明这件事即使他们非常愿意也不可能做到。因为必须寻找另一个世界。因此他加上：\BibleQuote{因为你们必须离开这世界。}你看见他不是严厉的主人，在他的立法中，他不仅考虑什么是可做的，也考虑什么是容易做的。因为他说，一个人照顾家庭、孩子，参与城市事务，或是工匠或士兵（人类的大部分是希腊人），如何能避免到处都是的不洁之人？因为\Quote{这世上的淫乱者}他指的是希腊人中的那些人。\BibleQuote{但现在我写信给你们，若有称为弟兄的是这样的人，}这样的人甚至不可与他吃饭。\par

但一个人如何能称为\Quote{弟兄}却又是拜偶像者？正如从前发生在撒玛黎雅人身上的情况，他们只选择了一半的虔诚。此外，他也预先为以后要论述的关于祭偶像之物的论述打下基础。\par

\Quote{或贪婪的。}他也与这些人冲突。因此他也说：\BibleQuote{为什么不情愿受欺？为什么不情愿被亏待？然而你们反倒欺压、亏待。}\par

\Quote{或酗酒的。}这一点他后来也指责他们；如他所说：\BibleQuote{有人饥饿，有人醉酒。}以及：\BibleQuote{食物是为肚腹，肚腹是为食物。}\par

\Quote{或辱骂的，或勒索的：}这些他先前也已责备过。\par

2. 接着他也加上禁止他们与那种品性的外教人交往的原因，暗示这不仅是不可能的，也是多余的。\par

% structure: p0010 source=h2 target=subsection reason=relative-heading-rank; base=h2

\subsection{格林多前书 5:12-13}

\BibleQuote{与我何干？审判那些外教人？}他称基督徒为\Quote{内}，希腊人为\Quote{外}，正如他在别处也说：\BibleRef{弟前 3:7} \BibleQuote{在外教人也必须有好的声望。}在给得撒洛尼人的书信中他也说了同样的话：\BibleRef{得后 3:14} \BibleQuote{不要和他来往，使他羞愧。}以及：\BibleQuote{不要以他为仇人，但要劝戒他如弟兄。}但在这里，他没有加上理由。为什么？因为在另一情形中他愿意安抚他们，但在这里却不是。因为这里的过失和那里的过失并不相同，在得撒洛尼人的情形中过失较轻。因为在那里他指责的是懒惰；但在这里却是淫乱和其他最严重的罪。如果有人愿意去希腊人那里，他不禁止他与这样的人吃饭；这也是出于同样的理由。我们也如此行：为了我们的儿女和弟兄，我们无所不为，但对陌生人我们却不怎么关心。那么怎样？保禄不也关心那些外人吗？是的，他关心他们；但要等到他们接受了福音，并使他们都服从基督的道理之后，他才为他们制定法律。但只要他们藐视，向不认识基督本身的人宣讲基督的诫命就是多余的。\par

\BibleQuote{你们不是审判内在的人吗？至于外在的人，由天主审判。}因为他说了：\BibleQuote{与我何干？审判那些外人？}免得有人以为这些人会不受惩罚，他为他们设立了另一个审判庭，而且是可怕的。他说这个，既是为了吓唬他们，也是为了安慰这些；他也暗示，这暂时的惩罚把他们从永久的、不停的惩罚中解救出来：这他在别处也清楚地表明了，说：\BibleRef{格前 11:32} \BibleQuote{我们受审判，正是被主惩诫，免得我们和世界一同被定罪。}\par

3. \Quote{你们应当把那恶人从你们中间除掉。}他采用了旧约中的一个表达\BibleRef{申 17:7}，部分地暗示他们也将大有所获，如同从某种严重的瘟疫中获得释放；部分地表明这种做法并非新事，而是从一开始立法者就认为这样的人应当被剪除。但在旧约中，这样做更为严厉；在此则更为温和。因此，人们有理由质疑：为什么在那情况下他允许罪人受严厉惩罚并被石头砸死，而在当前情况下却不如此；相反，他引导他悔改。那么，为什么前一种情况下的处理方式与后一种不同呢？原因有二：其一，因为这些人面临更大的考验，需要更多的忍耐；其二，也是更真实的理由，因为这些人若不受惩罚，更容易被纠正，从而可能悔改；而另一些人则可能继续作恶。因为如果当他们看到最先受惩罚的人仍坚持同样的事，而若完全无人受惩罚，他们的感受会更甚。因此，在那旧约中，奸夫和杀人者立刻被处死；但在新约中，若他们通过悔改获得赦免，就免去了惩罚。不过，在此我们也可以看到一些更重惩罚的例子，在旧约中也有一些较轻的，以便从各方面表明两约是彼此相近的，出自同一位立法者：你会看到惩罚立刻降临，既在这约中也在那约中，并且常常是在长时间之后。不仅如此，甚至并非总是长时间之后，全能者仅以悔改作为满足。例如在旧约中，犯奸淫和杀人的达味因悔改而得救；在新约中，阿纳尼雅因私留了田地价款的一部分而与妻子一同丧亡。如果这些例子在旧约中更常见，而相反的例子在新约中更常见，那么人物的不同导致了在这些事情上处理方式的差异。\par

% structure: p0014 source=h2 target=subsection reason=relative-heading-rank; base=h2

\subsection{格林多前书 6:1-2}

4. \Quote{你们中间有人与弟兄（\OriginalTerm{弟兄}{\GreekText{τὸν} \GreekText{ἀδελφόν}}，公认文本作\OriginalTerm{另一人}{\GreekText{τὁν} \GreekText{ἓτερον}}）相争，竟敢在不义的人面前，而不在圣人面前求审吗？}\par

在此，他再次基于公认的事实提出指责；因为在另一处他说：\Quote{我确实听说你们中间有淫乱的事。}而在此处他说：\Quote{你们中间竟有人敢？}从最初开始就表现出他的愤怒，并暗示所说之事源于一种大胆和无法无天的精神。\par

那么，他为什么顺便插入关于贪婪和不应在教会外诉讼的论述呢？是为了履行他自己的规则。因为他有一个习惯，即当遇到问题时便加以纠正；正如当他谈论他们共同使用的餐桌时，便引出了关于奥秘的论述。同样，你在这里看到，既然他提到了贪婪的弟兄，急于纠正犯罪者，便不严格遵守顺序；而是再次纠正了在常规之外引入的罪，然后回到原来的主题。\par

让我们听听他对此还说了什么。\Quote{你们中间有人相争，竟敢在不义的人面前，而不在圣人面前求审吗？}他暂时使用这些个人化的措辞来揭露、指责和责备他们的行为；他并没有一开始就完全推翻在信徒面前寻求审判的习俗；但当他用许多话打击他们之后，才完全取消了所有诉讼。他说：\Quote{首先，如果必须诉讼，在不义的人面前诉讼是错的。但你们根本不应该诉讼。}然而，这是他后来才补充的。目前他彻底审查了前一个主题，即他们不应将事情提交给外界仲裁。他说：\Quote{一个与朋友不和（\OriginalTerm{心怀不满}{\GreekText{μικροψυχοῦτα}}）的人，怎能不荒谬地让他的仇敌作他们之间的调解者？当一个希腊人坐在那里审判一个基督徒时，你怎能不感到羞耻和脸红？如果连私事都不应在希腊人面前诉讼，我们又怎能将其他更重要的事情交由他们裁决呢？}\par

再者，请注意他是如何说话的。他没有说\Quote{在不相信的人面前}，而是说\Quote{在不义的人面前}；他使用了对他当前问题最为必要的表达，以阻止和远离他们。因为你看，他的论述是关于诉讼，而那些参与诉讼的人最渴望的莫过于法官对正义有极大的关注；他以此作为劝诫的理由，几乎是在说：\Quote{你要去哪里？你在做什么，人啊，让自己遭遇与你愿望相反的事，为了获得正义而将自己交托给不义的人？}因为立刻被告知不要诉讼是难以忍受的，所以他并没有立即加上这一点，而只是改变了法官，将诉讼的一方从外面带到教会内。\par

5. 然后，由于被教会内的人审判似乎容易遭到轻视，特别是在当时（因为他们或许不能理解某个论点，也不像外教法官那样精通法律和修辞学，因为其中大部分是未受教育的人），请注意他是如何使他们值得信任的：首先称他们为\Quote{圣人}。\par

但既然这事见证的是生活的纯洁，而非听案的精确，请注意他如何有条理地处理这部分，这样说：\BibleQuote{你们不知道圣人们将要审判世界吗？} 那麼，你们现今在审判他们时，怎能容忍被他们审判呢？他们确实不会亲自坐堂、要求对簿来审判，但他们会定罪。至少他明确说了：\BibleQuote{如果这世界在你们中间受审判，难道你们不配判断最小的事吗？} 他不是说\BibleQuote{由你们}，而是说\BibleQuote{在你们中间}：正如主曾说：\BibleRef{玛 12:42} \BibleQuote{南方的女王将起来定这世代的罪}，和\BibleQuote{尼尼微人要起来定这世代的罪}。因为当看见同一个太阳、分享所有同样的事物时，我们被发现是信徒，而他们是不信者，他们将无法以无知为托辞。因为我们单凭所行的事就要控告他们。在那里，你将发现许多这样的审判方式。\par

然后，为免有人以为他指的是别人，请注意他如何概括他的话：\BibleQuote{如果这世界在你们中间受审判，难道你们不配判断最小的事吗？}\par

他说，这事对你们是耻辱，是说不出的谴责。因为既然他们很可能会因为在内部的人面前受审判而难堪；他说：\Quote{不}，\Quote{相反，耻辱是在你们被外人审判时：因为那些是极小的事，而非这些。}\par

% structure: p0024 source=h2 target=subsection reason=relative-heading-rank; base=h2

\subsection{格林多前书 6:3}

\BibleQuote{你们不知道我们将审判天使吗？何况是今生的事呢？}\par

有人说这里暗示的是司铎，但不要这样想。他的话是关于魔鬼的。因为如果他是在说腐败的司铎，那么他刚才说\BibleQuote{世界在你们中间受审判}时就会指他们（因为圣经也惯常称恶人为\Quote{世界}），并且他不会说同样的话两次，也不会好像是在讲更重要的事而将其放在后面。但他指的是那些天使，基督关于他们说：\BibleQuote{你们到那为魔鬼和他的天使预备的火里去}。\BibleRef{玛 25:41} 保禄也说：\BibleQuote{他的天使装作义德的仆役}。\BibleRef{格后 11:15} 因为当那些无形的权能竟被我们这些穿着肉身的所超越时，他们将受更重的惩罚。\par

但如果还有人争辩说他指的是司铎，让我们问：\Quote{哪种司铎？} 当然是那些生活行径属世俗的。那么他说\BibleQuote{我们将审判天使，何况是今生的事呢？} 是什么意思？他提到天使，是与\BibleQuote{今生的事}相对照：确实如此；因为他们远不需要这些东西，因其本性的卓越超凡。\par

% structure: p0028 source=h2 target=subsection reason=relative-heading-rank; base=h2

\subsection{格林多前书 6:4}

6. \BibleQuote{如果你们需要判断今生的事，就当在教会中设置那些微不足道的人来审判。}\par

为尽可能有力地教训我们，他们不应将任何事交给外人；他先提出一个似乎的反对，然后解答。他所说的与此类似：也许有人会说：\Quote{你们中没有一个智慧的人，也没有一个能判决的；都是可鄙的。} 那么结果如何？他说：\Quote{即使没有一个智慧的人，我仍吩咐你们把事务托付给那些最微不足道的人。}\par

% structure: p0031 source=h2 target=subsection reason=relative-heading-rank; base=h2

\subsection{格林多前书 6:5}

\BibleQuote{但我说这话，是为叫你们羞愧。} 这是揭露他们的反对是空托辞的话；因此他加上：\BibleQuote{难道你们中没有一个智慧的人，甚至一个也没有吗？} 他说，匮乏如此之大吗？你们中明智的人如此缺乏吗？他接着说的更严厉。因为说了\BibleQuote{难道你们中没有一个智慧的人，甚至一个也没有吗？} 之后，他加上：\BibleQuote{谁能在自己兄弟的案件中判断呢？} 因为当兄弟与兄弟对簿公堂时，调解此事的人根本不需要聪明才智，感情和关系大大有助于解决这种争吵。\par

\BibleQuote{但弟兄与弟兄告状，而且是在不信者面前。} 你注意到他最初是如何通过称审判者为不义来贬低他们，而在此，为叫他们羞愧，他称他们为不信者？因为如果司铎甚至不能成为弟兄间的和好媒介，而必须求助于外人，那真是极其可耻。因此，当他说\Quote{那些微不足道的人}时，他的主要意思（\OriginalTerm{这不是他主要的意思}{\GreekText{οὐ} \GreekText{τοῦτο} \GreekText{εἰπε} \GreekText{προηγουμένως}}）并不是主张教会的弃民应被任命为审判者，而是为了责备他们。因为他已通过说\BibleQuote{难道你们中没有一个智慧的人，甚至一个也没有吗？} 表明，理应推荐那些有能力裁判的人。并且他以极大的说服力堵住他们的口，说：\Quote{即使没有一位智慧的人，听证也应该留给你们这些不智慧的人，而不是让外人审判。} 因为还有什么比这更荒谬：在家中发生争吵时，我们不请外人介入，并且如果内情传到陌生人耳中，我们感到羞愧；而教会，这不可言喻的奥秘之宝库，竟让一切外传呢？\par

% structure: p0034 source=h2 target=subsection reason=relative-heading-rank; base=h2

\subsection{格林多前书 6:6}

\BibleQuote{但弟兄与弟兄告状，而且是在不信者面前。}\par

指控是双重的：既因他\Quote{告状}，又因\Quote{在不信者面前}。因为即使告状本身——与弟兄告状——已是过错，何况还在外人面前做这事，这还能得到什么宽恕呢？\par

% structure: p0037 source=h2 target=subsection reason=relative-heading-rank; base=h2

\subsection{格林多前书 6:7}

7. \BibleQuote{不，你们彼此诉讼，这已经是你们的一个缺陷了。}\par

你们看，他把这一点保留在什么地方？又是多么适时地把对它的讨论澄清了？因为\Quote{我尚未谈论，}他说，\Quote{谁损害人，或被损害。}至此，诉讼行为本身使双方都受他的责备；在这点上，一方并不比另一方好。但一个人诉讼是公正还是不公正，那是另一件事。所以不要说：\Quote{谁做了错事？}因为单凭诉讼行为本身，我就立刻定你的罪。\par

如果忍受不了作恶的人算是一种过错，那么有什么控告能比得上实际的恶行呢？\Quote{为什么不情愿受冤？为什么不情愿被欺骗？}\par

% structure: p0041 source=h2 target=subsection reason=relative-heading-rank; base=h2

\subsection{格前 6:8-10}

\Quote{不，你们自己反而做错事，欺骗人，而且是对弟兄做的。}\par

这又是双重罪行，甚至可能是三重或四重。其一，不知道怎样忍受冤屈。其二，实际去做错事。其三，竟将这些事务的解决委托给不义的人。其四，这一切竟是对一位弟兄做的。因为人的过犯，当它们是对任何偶然的人所犯，与对自身的肢体所犯时，判断的标准并不相同。因为冒犯后者，必定是更大程度的鲁莽。在前一种情形，只践踏了事情的本质；但在后一种，也践踏了人的身份。\par

正如你们所见，他已经从普遍原则的论证使他们羞愧，而在那之前，又从所提出的赏报；他用威胁来结束劝诫，使他的讲话更严厉，这样说：\BibleRef{格前 6:9}\Quote{你们岂不知道不义的人不得继承天主的国吗？不要被迷惑：无论是淫荡的、拜偶像的、奸淫的、柔弱的、与男人苟合的、\BibleRef{格前 6:10}贪婪的、偷窃的、醉酒的、辱骂的、勒索的，都不得继承天主的国。}你说什么？论及贪婪之人时，你竟把如此庞大的一群不法之徒带到了我们面前？\Quote{是的，}他说，\Quote{但我这样做，并没有混淆我的论述，而是按着次序进行。}因为正如论及不洁之人时，他同时提到了所有罪；同样，再次提到贪婪之人时，他带出所有罪，这样使责备对那些良心有这些事的人变得熟悉。因为不断提及为他人积存的惩罚，使得责备在面对我们自己的罪时容易被接受。所以在此例中，他发出威胁，绝无意味他知道他们做这些事，也不是在责问他们——这是一个特别有力的方式，能抓住听者并防止他脱逃：即，说话并不针对他，而是无指定地发言，从而暗中刺伤他的良知。\par

\Quote{不要被迷惑。}在这里，他暗示某些人（正如现在大多数人所宣称的那样）主张天主既然是良善仁慈的，就必不会报复我们的恶行：\Quote{那么我们就不必害怕吧。}因为他决不会向任何人追讨任何事的正义。正是为了这些人，他说：\Quote{不要被迷惑。}因为依赖于善却遇到相反的事，这是极端的错误和迷妄；并且猜想关于天主的这样的事，是人所不能想到的。因此先知以天主的口气说：\BibleRef{咏 49:21}\Quote{你怀了恶念，以为我像你一样：我要责斥你，将你的罪孽摆在你面前。}而保禄在此说：\Quote{不要被迷惑；无论是淫荡的}（他先把已经被定罪的放在前面）\Quote{奸淫的、柔弱的、醉酒的、辱骂的，都不能继承天主的国。}\par

许多人攻击这段经文极为严厉，因为他将醉酒的与辱骂的放在与奸淫的、可憎的以及与男人苟合者同等的位置。然而罪行并不相等：那么惩罚如何相同？那么我们该说什么？首先，醉酒与辱骂并非小事，因为基督亲自将称弟兄为\Quote{愚痴}的人交付地狱。而且那罪常常带来死亡。再者，犹太人民也因醉酒犯下他们最大的罪。其次，他至今所谈论的不是惩罚，而是从国度的排除。从国度中，两者都被同样地驱逐出去；但在地狱中他们是否会感到差别，不属于此时此地要探讨的。因为那个话题现在不在我们面前。\par

% structure: p0047 source=h2 target=subsection reason=relative-heading-rank; base=h2

\subsection{格前 6:11}

9. \Quote{你们中从前也有人是这样：但你们被洗涤了，但你们被圣化了。}\par

为了使你们极其羞愧，他添加了这一点：好像他说，想一想天主将我们从何等邪恶中拯救出来；祂赐予我们多么伟大的慈爱实验和证明！祂并没有将祂的救赎局限于单纯的拯救，而是大大扩展了恩惠：因为祂也使你们洁净。难道就这样吗？不：祂还\Quote{圣化了。}即便如此也不是全部：祂还\Quote{成义了。}然而，即使是单纯地从我们的罪中拯救出来也是一份大礼：但如今祂还用无数的祝福充满你们。祂所做这一切，是\Quote{因我们的主耶稣基督之名，}不是这个或那个名：是的，也\Quote{因我们天主的圣神。}\par

因此，亲爱的，知道这些事情，并牢记所成就的祝福的伟大，让我们一方面继续清醒地生活，远离所列举的一切事物；另一方面避免在外邦人的广场上的法庭；并保全天主白白赐给我们的高贵出生，直到终了。因为想一想，一个希腊人竟坐下来审判你，是多么羞耻的事。\par

但你会说：倘若内中的人按相反法律判决，又当如何？他为何要这样？请告诉我。因为我想知道，希腊人依据何种法律施行审判，基督徒又依据何种？难道不是明摆着，人的法律是希腊人的准则，而天主的法律是基督徒的准则吗？那么后者显然有更大的公正机会，因为这些律法甚至是从天上而来的。至于那些外在的人，除了所说的之外，还有许多别的事也值得怀疑：演说者的口才、长官的腐败，以及许多其他破坏公正的事。但在我们这里，却没有这类事。\par

\Quote{那么，}你会说，\Quote{倘若仇敌是位高权重的人呢？}这恰恰是最应该到基督徒法庭去诉讼的理由：因为在世俗法庭，他无论如何都会胜过你。\Quote{但他若不听从，既藐视我们内部的人，又强行拖到外面的法庭，又当如何？}那就不如甘愿忍受你可能会被迫承受的事，不去诉讼，这样你还能获得赏报。因为，\BibleRef{玛 5:40}\Quote{若有人愿与你诉讼，拿你的内衣，你连外衣也让他拿去；}并且\BibleRef{玛 5:25}\Quote{当你同仇敌还在路上时，赶快与他和解。}我何必谈我们的规则呢？就连外邦法庭上的律师也常告诉我们：\Quote{庭外和解更好。}但是，财富啊，更不如说是荒谬的爱财之心！它颠覆一切，败坏一切；由于金钱，许多事物对众人而言不过是空谈和寓言！现在，给法庭带来麻烦的是世俗之人，这并不奇怪；但许多告别了世俗的人竟也做同样的事，这是完全不可原谅的。因为如果你愿意按照圣经的规则，看看自己该远离这种需要——我指法庭的需要——多远，并要知道法律是为谁设立的，请听保禄所说：\BibleRef{弟前 1:9}\Quote{因为法律不是为义人设立的，乃是为不法与悖逆的人设立的。}如果他说这些是关于梅瑟法律，那么关于外邦人的法律更是如此。\par

10. 如今，你若行了不义，显然不能是义人；但若你受损害并忍受之（这正是义人的特殊标志），你就不需要外在的法律。\Quote{那么，}你说，\Quote{我受损害时怎能忍受呢？}然而，基督竟命令了比这更多的事。因为祂不仅命令你受损害时要忍受，更要在甘愿受苦的热心上超过作恶者的急切。祂没有说：\Quote{若有人愿与你诉讼，拿你的内衣，你就把内衣给他，}而是说：\Quote{连外衣也一并给他。}祂说：我命令你以受苦而非作恶来胜过他；因为这是确定而辉煌的胜利。因此保禄继续说：\Quote{你们彼此有诉讼，这已经是你们的一个\OriginalTerm{缺陷}{\GreekText{ἥττημα}}（修订本：\InnerQuote{过失}）了。}又说：\Quote{为什么不宁愿受亏负呢？}因为受损害的那个人胜过不能忍受损害的人，这一点我将向你们证明。不能忍受损害的人，即使强迫对方上法庭并赢得判决，那时他反而是最失败的。因为他遭受了他所不愿的事：仇敌迫使他既感到痛苦又去诉讼。你胜诉了又怎样？你收回了所有钱又怎样？你暂时承受了你不愿意的事，被迫以法律解决争端。但如果你忍受不义，你就得胜了；虽然失去了金钱，但绝没有失去因这种自制而来的胜利。因为对方没有能力强迫你做你不喜欢的事。\par

为证明这是真的，请告诉我：在粪堆上，谁胜了？谁败了？是被剥夺了一切的约伯，还是剥夺了他一切的魔鬼？显然是剥夺了他一切的魔鬼败了。我们钦佩谁得胜？是击打的魔鬼，还是被击打的约伯？显然是约伯。他虽然不能保留他易逝的财富，也不能救自己的孩子。我何必提财富和儿女呢？他连身体的健康都不能保证。然而尽管如此，他仍是得胜者，他失去了所有的一切。他的财富的确不能守住，但他的虔诚却严谨地持守了。\Quote{但他不能帮助他的孩子免于死亡。}那又怎样？所发生的事既使孩子们更光荣，另一方面，他也以这种方式保护了自己免受恶意对待。倘若他没有受苦、没有受魔鬼的亏负，他就不会获得那显著的胜利。若受苦是恶事，天主就不会命令我们这样做；因为天主不命令恶事。难道你们不知道祂是荣耀的天主吗？祂不愿意使我们蒙受羞耻、嘲笑和损失，而是愿意把我们\OriginalTerm{引入}{\GreekText{προξενῆσαι}}与这些相反的事。因此祂命令我们受亏负，并且做一切事使我们从世俗事物中抽离，使我们确知什么是荣耀，什么是羞耻；什么是损失，什么是收获。\par

\Quote{但受亏负和被恶意对待是难的。}不，人啊，并不难，并不难。你的心要到几时还为眼前的事不安呢？因为天主必不至于命令难做的事，你可以确信。请思考：作恶者带着钱走了，但同时也带着邪恶的良心；受亏负者虽然确实被骗走了某些钱，却因对天主的信赖而富有；这种获得比无数财宝更宝贵。\par

11. 既然知道这些事，我们就要出于自由意志，持守严格原则，不要像那些不明智的人，以为当自己因试炼而受苦时，就没有受委屈。但恰恰相反，那是最严重的伤害；所以每当我们在这些事上自我克制时，并非出于自愿，而是在那另一方面失利之后。因为一个人在被判决中忍受它，并无益处；因为从此事成了必然。那么什么才是辉煌的胜利？当你藐视它：当你拒绝诉讼时。\par

\Quote{你怎么说？我被剥夺了一切吗？} 一人说，\Quote{你还要我保持沉默？我遭受了羞辱，你还劝我温顺地忍受？我怎能做到呢？} 不，但这是最容易的，只要你仰望天上；只要你看见眼前的美丽；并且天主应许在你高尚地忍受冤屈时接纳你。那么就这样做；仰望天上，思想你被造成相似那坐在革鲁宾上的那一位。因为祂也受了伤害，祂忍受了；祂被辱骂，却没有报复；祂被打，却没有辩护。相反，祂以相反的方式回报那些行这种事的人，甚至赐予无数的恩惠；并且祂命令我们效法祂。想一想你从母胎赤身而来，那得罪你的人也将赤身而去；而他呢，带着无数的伤口，生虫。想一想现在的事不过是暂时的；数算你祖先的坟墓；精确地知晓过去的事；你将会看到那作恶者使你更坚强。因为他加重了自己的情欲，我指的是贪婪；但他减轻了你的（情欲），拿走了野兽的食物。除此之外，他使你摆脱了忧虑、痛苦、嫉妒、告密、麻烦、担忧、持久的恐惧；他将那污秽的恶行堆在自己头上。\par

\Quote{那么，} 一人说，\Quote{如果我必须与饥饿斗争呢？} 你与保禄一同忍受它，保禄说，\BibleRef{格前 4:10} \BibleQuote{直到此时此刻，我们仍是又饥又渴，又赤身露体。} 但他是为了天主的缘故，你会说，\Quote{为了天主的缘故：} 你也为天主的缘故这样做吧。因为当你放弃报复时，你是为天主的缘故这样做。\par

\Quote{但那得罪我的人，与富人一同享乐。} 不，不如说与魔鬼一同享乐。但你与保禄一起接受冠冕。\par

因此，不要害怕饥饿，因为\BibleRef{箴 10:3} \BibleQuote{上主不让义人受饥挨饿。} 此外，另一位说，\BibleRef{咏 54:23} \BibleQuote{将你的挂虑托给上主，祂必供养你。} 因为田野的麻雀都蒙祂饲养，何况你呢？现在，我亲爱的，让我们不要小信德，也不要小气！因为祂应许了天国和如此大的祝福，怎能不赐予现世之物呢？让我们不要贪求多余之物，而要持守知足，我们就总是富足。让我们寻求居所和食物，我们必将得到一切；既得到这些，也得到那远为更大的。\par

但如果你仍在悲伤和沮丧，我想向你展示作恶者在得胜之后的灵魂，它如何变成了灰烬。因为罪确实是这样的：当人犯罪时，它提供某种快乐；但一旦结束，那微小的快乐就消失得无影无踪，取而代之的是沮丧。我们伤害别人时也是这样：之后，我们至少会谴责自己。同样，当我们欺诈时，我们也有快乐；但之后我们会受到良心的刺痛。你看到某人的财产中有穷人的家吗？不要为那被掠夺的人哭泣，而要哭那掠夺者：因为他不是施加了恶，而是承受了恶。因为他抢夺了别人的现世之物；但他自己却被逐出那不可言宣的祝福。因为凡不给穷人的人要下地狱；那抢夺穷人财物的人将受什么苦呢？\par

\Quote{可是，} 一人说，\Quote{如果我受害，有什么益处呢？} 确实，益处是大的。因为天主并非按那害你之人的惩罚来为你安排补偿：因为那算不了什么大事。如果我受害而他受害，那有什么大好处呢？然而我知道许多人，他们认为这是最大的安慰，当他们看到那些曾侮辱他们的人受到惩罚时，他们认为自己已得回了所有。但天主并不将他的补偿限于此。\par

那么，你真心想知道等待你的祝福有多大吗？祂为你敞开整个天空；祂使你成为圣人的同胞；祂使你参与他们的歌咏团：祂赦免了你的罪；以义德为你加冕。因为如果那些宽恕冒犯者的人将获得宽恕，那些不仅宽恕而且还慷慨施舍的人，将继承什么样的祝福呢？\par

因此，不要以小气的心态忍受，反而要为那得罪你的人祈祷。你这样做是为了你自己。他拿走了你的钱财？好吧：他也拿走了你的罪：那正是纳阿曼和革哈齐的情况。你愿意付出多少财富来使你的罪孽得到赦免？相信我，现在正是如此。因为如果你高尚地忍受而不咒骂，你就在自己身上绑了一顶光荣的冠冕。这不是我的话，而是你听到基督说的：\BibleQuote{为那迫害你们的人祈祷。} 而且考虑奖赏有多大！\BibleQuote{好使你们成为你们在天之父的子女。} 如此，你什么也没有失去，反而有所得：你没有受害，反而被加冕；因为你在灵魂中得到了更好的操练；变得相似天主；从对钱财的忧虑中解脱出来；成为天国的拥有者。\par

因此，亲爱的，考虑到这一切，让我们在受伤害时克制自己，好使我们既能摆脱现世生活的纷扰，又能驱除一切无益的灵性悲伤，并获得未来的喜乐；藉着我们的主\Person{耶稣基督}的恩宠与慈爱，愿光荣、权能、尊崇归于祂，及\Person{圣父}与\Person{圣神}，自今而始，及至永远。阿们。\par

\SourceRecord{Translated by Talbot W. Chambers.；Nicene and Post-Nicene Fathers, First Series；Vol. 12.；Edited by Philip Schaff.；Buffalo, NY: Christian Literature Publishing Co.,；1889.；<http://www.newadvent.org/fathers/220116.htm>.}

% structure: p0001 source=h1 target=section reason=page-title

\section{格林多前书第十七篇讲道}

% structure: p0002 source=h2 target=subsection reason=relative-heading-rank; base=h2

\subsection{格前 6:12}

\BibleQuote{凡事我都可行，但不全有益；凡事我都可行，但我却不受任何事物的辖制。}\par

此处他暗示贪食者。因为他打算再次攻击行淫者，而淫乱源自奢侈和缺乏节制，所以他严厉地谴责这种激情。他不可能就这样谈论禁止的事，那些事不是\Quote{合法的：}而是谈论那些看似无关紧要的事。为了说明我的意思：\Quote{合法，}他说，\Quote{可以吃喝；但过度则无益。}因此，他那奇妙的、出人意料的转折，是他经常采用的习惯；\BibleRef{罗 12:21}将他的论证完全转到相反的方面，他在这里也设法引入；并指出，做自己有权能做的事不仅无益，甚至不是权能的一部分，而是奴役的一部分。\par

首先，他以事情无益为由劝阻他们，说：\Quote{它们并无益}；其次，以事情自相矛盾为由，说：\BibleQuote{我绝不受任何事物的辖制}。他的意思是：\Quote{你有吃的自由，}他说，\Quote{那么，就保持自由，注意不要成为这欲望的奴隶：因为适当使用的人，是它的主人；但过度的人不再是它的主人，而是它的奴隶，因为贪食在他内作主。}你们看出，人以为自己有权柄的地方，保禄却指出他是在权柄之下吗？因为这是他的习惯，正如我前面所说，把所有反对意见转向相反的方向。他在这里正是这样做的。请注意：他们每个人都说：\Quote{我有权过奢侈的生活。}他回答说：\Quote{这样做，你不像是对事物有权柄，反而像是你自己受某种权柄的支配。因为只要你是放荡的，你甚至连自己的肚腹也没有权柄，而是它有权柄支配你。}对于财富和其他事物，我们也可以这样说。\par

% structure: p0006 source=h2 target=subsection reason=relative-heading-rank; base=h2

\subsection{格前 6:13}

\BibleQuote{食物是为肚腹。}这里他所说的\Quote{肚腹}，不是指胃，而是指胃的贪饕。正如他说的：\BibleRef{斐 3:19}\BibleQuote{他们的神是肚腹：}不是谈论身体的那部分，而是贪饕。为证明就是这样，请听接下来：\BibleQuote{肚腹是为食物；但身体不是为淫乱，而是为主。}然而\Quote{肚腹}也是\Quote{身体}的一部分。但他列出两对：\Quote{食物}和贪食（他称之为\Quote{肚腹}）；\Quote{基督}和\Quote{身体}。\par

那么，\Quote{食物是为肚腹？}是什么意思？他说，\Quote{食物}与贪食关系良好，贪食也与食物关系良好。因此它不能引领我们到基督那里，反而把我们拖向这些事物。因为这是一种强烈而野蛮的激情，使我们成为奴隶，让我们服事肚腹。那么，人哪，你为何激动地张大口渴望食物呢？因为那种服事的结局就是这样，此外再也看不到什么：但就像有人服侍一位女主人，它始终维持着这种奴役，不再前进，除了这无果的服事之外，别无作为。两者相互关联，也一同消亡；\Quote{肚腹}与\Quote{食物}，\Quote{食物}与\Quote{肚腹}；编织出一种无尽的循环；正如从腐败的身体生出虫子，虫子又消耗身体；或如涌起的高浪破碎，再没有任何效果。但他说的这些不是关于食物和身体，而是谴责贪食和饮食过度的激情：下面的话也显示这一点。因为他继续：\par

\BibleQuote{但天主必要使二者都归于乌有：}这不是指胃，而是指不节制的欲望：不是指食物，而是指过度的饮食。因为对于前者（食物），他并不生气，甚至为他们制定规则，说：\BibleRef{弟前 6:8}\BibleQuote{我们有吃有穿，就当知足。}然而，他这样谴责整个事情；它的改正（在劝告之后）由他留给了祈祷。\par

但有些人说这些话是预言，宣布将来生命中的状态，那里没有吃喝。如果适度的饮食也有终结，我们更应当戒绝过度。\par

然后，免得有人以为身体是他指责的对象，并且怀疑他由部分责怪全体，说身体的本性是贪食或行淫的原因，请听下面的话。\Quote{我不责备，}他说，\Quote{身体的本性，而是思想的过度放纵。}因此他又加上：\BibleQuote{现在身体不是为淫乱，而是为主；}因为它不是为了这个目的而形成的，即放荡生活和行淫，正如肚腹也不是为了贪婪；而是为了使它跟随基督为元首，并使主作身体的元首。让我们深感羞愧，让我们惊恐，因为我们被算为配得上如此大的荣誉，成为那坐在至高者身上的肢体，却用如此大的邪恶玷污自己。\par

2. 在充分谴责了贪食者之后，他也利用对将来事物的希望来使我们远离这邪恶：说，\par

% structure: p0013 source=h2 target=subsection reason=relative-heading-rank; base=h2

\subsection{格前 6:14}

\BibleQuote{天主既使主复活了，他也要藉自己的能力使我们复活。}\par

你看出他宗徒的智慧了吗？因为他总是藉基督建立复活的可靠性，尤其是现在。因为如果我们的身体是基督的肢体，而基督已经复活，身体也必定跟随元首。\par

\Quote{藉他的德能。}因为他既然提到一件不被相信且无法用理性理解的事，就完全将其交给基督自己复活的那不可理解的能力，也以此作为反对他们的一个不小的论证。关于基督的复活，他没有插入这一点：因为他没有说，\Quote{天主也将复活主；}——因为事情已经过去了——而是怎样呢？\Quote{天主既已复活了主；}并不需要任何证明。但关于我们的复活，因为它尚未发生，他没有这样说，而是怎样呢？\Quote{也要藉他的德能复活我们；}凭着对工作者德能的信赖，他堵住了反对者的口。\par

再者：如果他将基督的复活归于圣父，切莫为此困扰。因为他记下这一点，并非以为基督没有能力，因为正是他自己说：\BibleRef{若 2:19}\Quote{你们拆毁这座圣殿，三天内我要把它重建起来。}又说：\BibleRef{若 10:18}\Quote{我有权舍去我的生命，也有权再取回它。} 路加在\BibleBook{宗徒}{大事录}中也说：\BibleRef{宗 1:3}\Quote{他向他们显示自己活着。} 那么，保禄为何这样说呢？因为圣子的作为都归于圣父，圣父的作为也归于圣子。因为他说：\BibleRef{若 5:19}\Quote{圣父所做的，圣子也同样做。}\par

他很合时宜地在此提到了复活，用这些希望来压制贪饕欲望的暴政；他几乎是在说：你吃了，你喝得过量了；结果如何呢？无非是毁灭。你与基督结合了；结果如何呢？一件伟大而奇妙的事：将来的复活，那光荣的、远超一切言语的复活！\par

3. 所以，谁也不要再不信复活了；但若有人不信，让他想想天主如何从无中创造了万物，并且以此为其他事的证明。因为已经过去的事，其不可思议远超一切，充满了令人惊奇的奥秘。试想：他取土和成泥，造了人；这土先前并不存在。那么，土如何变成了人？它又是如何从无中产生的？还有，由土所造的一切事物呢？那无数种类的无理性动物、种子、植物，在没有产痛的情况下，没有雨水降下，没有耕作，没有牛，没有犁，也没有任何其他因素帮助它们产生？正因如此，那无生命、无感觉的东西在起初就产生了如此多种类的植物和无理性动物，为的是从一开始就教导你复活的道理。因为复活比这更难以解释。重新点燃熄灭的灯，与展示从未出现过的火，是不同的。重建倒塌的房屋，与建造从未存在的房屋，是不同的。因为在前一种情况下，即使没有别的，至少还有材料可用；但在后一种情况下，连实体都不存在。所以，他先做了那看似更难的事，好让你由此承认那更容易的事——我说更难，不是对天主而言，而是就我们的理智所能追随的程度而言。因为对天主来说，没有难事；正如画家画了一幅肖像，就能轻易画出千万幅，同样，对天主来说，创造无数无量的世界也是容易的。甚至，正如你轻易就能想象出无止境的城市和世界，同样，对天主来说，创造它们也很容易；或者更确切地说，容易得多。因为你在想象中还要花费时间，尽管很短；但天主连这点时间都不用，正如石头比任何最轻的东西都重，甚至比我们的思想还重，同样，我们的思想在天主创造的迅捷面前也望尘莫及。\par

你惊叹他在大地上的德能吗？再想想天是如何造成的，它先前并不存在；那无数的星辰，太阳，月亮，它们先前都不存在。再说，它们被造之后是如何保持稳定的？在什么之上？它们有什么根基？大地呢？大地的后续是什么？再后来，大地后续之后又是什么？你看见吗？除非你迅速投奔信德和造物主那不可理解的德能，你的理智之眼将陷入怎样的漩涡？\par

但如果你也选择从人类的事物中推测，你将能逐渐为你的悟性添上翅膀。\Quote{什么样的人类事物？}有人会问。难道你没看见陶工如何将那摔碎、不成形的器皿重新塑造？熔炼矿石的人如何将土在手中变成\OriginalTerm{金、银、铜}{\GreekText{τὴν} \GreekText{γῆν} \GreekText{χρύσιον} \GreekText{ἀποφαίνουσι}}？其他加工玻璃的人，如何将沙变成一种紧密透明的物质？难道还要我说皮革匠、染紫袍的工匠，如何使那曾接受染色的东西显示出另一种颜色？难道还要我说我们人类的生育？一粒小小的种子，起初无形无状，进入接受它的子宫，然后那活物如此精细的构造从何而来？麦子是什么？不是一粒赤裸的种子被撒进地里吗？撒进之后，它不是腐烂了吗？穗、芒、秆和所有其他部分从何而来？难道不是常常有一粒小小的无花果籽落在地上，就生出根、枝和果实？对此你每件都承认，不加追究；唯独对天主改变我们身体形态的工作，你却要求解释？这类事怎能得到宽恕呢？\par

我们向希腊人说的是这些和类似的话。因为对于那些顺从圣经的人，我根本没有必要说什么。\par

我說，你若要好奇地追究祂的一切作為，那麼天主比人有什麼更多呢？然而即使是人，也有許多事我們是不這樣追問的。更何況我們更應當避免對天主的智慧作不相干的探究，並避免要求說明它的理由：首先，因為那作證者是可信的；其次，因為這事本身不能藉推理來考察。因為天主並非如此貧乏，以致只做那些能被你推理的軟弱所理解的事。如果你不能理解一個工匠的工作，更何況是眾工匠中最好的天主呢？那麼，不要懷疑復活，因為你將遠遠離開對未來之事的希望。\par

但反駁者們的明智論點是什麼？我該說，是他們極其愚蠢的論點？\Quote{為什麼，當身體與土混合而成了土，而這土又被移到別處時，}他們說，\Quote{它將如何復活？}對你來說，這似乎不可能，但對那不眠之眼卻不然。因為在祂眼中，萬事都是清楚的。而你在那混亂中看不到各部分的區別；但祂卻知道它們全部。因為你也不知道你鄰居的心，以及其中的事；但祂知道一切。那麼，因為你不知道天主如何使人復活，你就不相信祂使人復活，你也會不相信祂知道你在想什麼嗎？因為那也不是顯而易見的。然而在身體中，雖然它被分解，卻仍是可見的物質；但那些思想是不可見的。那麼，那確切知道不可見之事的那一位，難道看不見那些可見之事，並輕易分辨身體分散的各部分嗎？我想這是每個人都明白的。\par

那麼，不要懷疑復活；因為這是魔鬼的教義。魔鬼所極力追求的，不僅是使復活遭到懷疑，而且也使善行被廢除。因為那不期望自己將要復活並為自己所做的事交帳的人，不會很快致力於德行；他反過來會完全不相信復活：因為這兩者是互相鞏固的：惡行由不信鞏固，不信由惡行鞏固。因為充滿許多邪惡的良心，因害怕和戰慄將來的報應，又不願藉轉向至善來為自己提供安慰，便安於不信。因此，當你否認復活和審判時，另一方便會說：\Quote{那麼，我也不必為我的大膽行為交帳了。}\par

4. 但基督為什麼說？\BibleQuote{你們錯了，因為不認識聖經，也不認識天主的大能。}\BibleRef{玛 22:29}因為天主若無意要使我們復活，而只是要分解並消滅我們歸於虛無，祂就不會行這樣多的事。祂不會鋪開這諸天，不會伸展這大地，不會為了這短暫的生命而創造宇宙的其餘部分。但如果這一切都是為了現在，那麼為了將來，祂還會不做什麼呢？相反地，如果沒有未來的生命，那麼在這一點上，我們比那些為我們而造的事物更為低劣。因為諸天、大地、海洋、河流都比我們更持久；甚至有些禽獸也是如此：因為烏鴉、象族和許多其他動物，享壽更長於我們今生。對我們來說，生命既短暫又勞苦，但對牠們卻不然。牠們的生命既長久，又較少憂愁和掛慮。\par

\Quote{那麼，告訴我：祂使奴僕比主人更好嗎？}我懇求你，不要，不要這樣推理，人啊，也不要這樣思想貧乏，也不要無知於天主的富饒，因為你有如此偉大的主人。因為從起初天主就願意使你成為不朽的，但你卻不願意。因為那時候的事也是不朽的暗示：與天主來往；生活中沒有不安；沒有憂愁、掛慮和勞苦，以及其他屬於暫時存在的事物。因為亞當不需要衣服或住所，或任何這類供應；反而像天使一樣；而且他預知許多將來的事，充滿了大智慧。甚至天主在隱密中所做的事，他也知道——我指的是關於女人的事；因此他也說：\BibleQuote{這是我骨中的骨，肉中的肉。}\BibleRef{创 2:23}勞苦後來才出現；汗也一樣，羞恥、膽怯、缺乏自信也是如此。但在那日子，沒有憂愁、痛苦、哀嘆。然而他沒有持守在那尊位中。\par

那麼，有人說，我該怎麼辦？我必須因他而滅亡嗎？我回答說，首先，這不是因他：因為你也沒有保持無罪；就算不是同樣的罪，至少你犯了別的罪。再者，你也沒有因他的懲罰而受損，反而有所得。因為如果你要完全成為可朽的，那麼所說的話或許有些道理。但現在你是不朽的，而且如果你願意，你可以比太陽本身更明亮。\par

5.\Quote{但是，}有人說，\Quote{如果我沒有得到可朽的身體，我就不會犯罪。}那麼告訴我，他犯罪時有可朽的身體嗎？當然沒有：因為如果身體在此之前已是可朽的，它就不會後來因懲罰而經歷死亡。而且可朽的身體不是德行的障礙，反而使人遵守秩序並有極大的益處，從以下的事可以明白。因為如果單單不朽的期望就使亞當如此驕傲；那麼如果他實際上是不朽的，他會狂妄到什麼地步呢？而事實上，你在犯罪之後可以消除你的罪，因為身體是卑微的、衰敗的、易於分解的：因為這些想法足以使人清醒。但如果你在不朽的身體中犯罪，你的罪很可能會更持久。\par

那么，必朽并非罪的成因：不要归咎于它；邪恶的意志才是一切祸患的根源。因为为何亚伯尔丝毫不因身体而受害？为何魔鬼虽无形体却毫无改善？你愿意听听身体成为可朽的，非但无害，反而有益吗？请注意，你若清醒，会从中获得多少益处。它借着忧伤、痛苦、辛劳以及其他类似的事，将你拖回、拉离邪恶。\Quote{但它诱惑人行不洁之事，}也许你会说。这不是身体做的，而是不节制做的。因为我所提到的这一切确实属于身体：因此，进入此生的人不可能免于疾病、痛苦和沮丧；但不犯不洁之事却是可能的。由此可见，若邪恶的情欲是身体本性的一部分，它们就会是普遍的：因为一切本性之事都是如此；但行邪淫却并非如此。痛苦确实出自本性；但行邪淫却出于故意。\par

因此，不要责备身体；不要让魔鬼夺去天主赐予你的尊荣。因为，若我们愿意，身体是一副极好的缰绳，可抑制灵魂放纵的冲动，压下高傲，遏止骄横，并在最大的德行成就中帮助我们。请勿对我提及那些失去理智的人；因为我们常见马匹，在甩掉骑手后，拖着缰绳冲下悬崖，但我们并不责备缰绳。因为造成这一切的不是缰绳断裂，而是骑手未能勒住它们，才毁了一切。你在此事上也应这样推理。若你看见一个年轻人生活在孤苦中，行无数恶事，不要责备身体，而要责备被拖曳的御者，即人的理性官能。正如缰绳不为难御者，是御者因不善用而成为一切祸患的根源（因而缰绳常惩罚他，缠住他，拖曳他，迫使他参与自身的不幸），在我们眼前的情形也是如此。\Quote{我，}缰绳说，\Quote{曾使马口流血，只要你握着我；但自从你抛弃我，我要求你为藐视我而赔偿，我缠绕你，拖曳你，以免再受同样的待遇。}因此，没人应责备缰绳，而应责备他自己和他自己败坏的心思。因为在我们之上也有一位御者，即理性；缰绳是身体，连接马匹与御者；若它们状态良好，你便不受伤害；但若你松开它们，你就毁灭并败坏了一切。因此让我们节制，将一切责备不归于身体，而归于邪恶的心思。因为这是魔鬼的特殊工作：使愚昧人责备身体、责备天主、责备邻人，而不责备自己败坏的心思，免得他们发现原因后，从祸根中解脱出来。\par

但你们，既认清他的计谋，就当将愤怒转向他；并将御者安置在车上，使你们的心灵之眼朝向天主。因为在其他一切比赛中，设立竞赛者并不贡献什么，只是等候结局。但在这一比赛中，设立竞赛者——就是天主——是一切的一切。因此，让我们使他仁慈，我们就必获得预许的福分；借着我们的主耶稣基督的恩宠和慈爱，愿光荣、权能、尊崇归于他，及圣父及圣神，自今，及世之世。阿们。\par

\SourceRecord{Translated by Talbot W. Chambers.；Nicene and Post-Nicene Fathers, First Series；Vol. 12.；Edited by Philip Schaff.；Buffalo, NY: Christian Literature Publishing Co.,；1889.；<http://www.newadvent.org/fathers/220117.htm>.}

% structure: p0001 source=h1 target=section reason=page-title

\section{\WorkTitle{《格前》第十八篇讲道}}

% structure: p0002 source=h2 target=subsection reason=relative-heading-rank; base=h2

\subsection{格前 6:15-16}

\BibleQuote{你们不知道你们的身体是基督的肢体吗？那么，我可以把基督的肢体拿去，作为娼妓的肢体吗？断乎不可！}\par

在从行淫者转到贪财者之后，他又从后者回到前者，从此不再与那人谈论，而是与那些没有犯奸淫的其他人交谈。在保护他们以免陷入同样罪过的行动中，他又再次攻击那人。因为那犯了罪的人，即使你将话转向别人，他也会被刺痛；他的良心被彻底唤醒并鞭策他。\par

对惩罚的恐惧确实足以使他们保持贞洁。但他不愿单凭恐惧来纠正这些事，而是既用威胁也用理由。\par

在先前那场合，他陈述了罪过，规定了惩罚，并指出了与行淫者交往给众人带来的伤害，然后他停止并转到贪财的话题上：威胁贪财者以及他所提的其余所有人，将他们逐出天国，他就这样结束了讲论。但在这里，他以更可怕的方式着手劝诫。因为只惩罚罪过而不指出其最极端的不义，这样的惩罚不会产生如此大的效果；同样，只使人羞愧而不以惩罚方式恐吓，也不会对刚硬之心的人造成深刻打击。因此保禄两者兼用：在这里他使人羞愧，说：\BibleQuote{你们不知道我们要审判天使吗？}在那里他又恐吓，说：\BibleQuote{你们不知道贪财者不能继承天主的国吗？}\par

关于行淫者，他又用了这种论述顺序。因为他先用之前所说的话恐吓他：首先把他割除并交给撒殚，然后提醒他那即将来临的日子；他又用\BibleQuote{你们不知道你们的身体是基督的肢体吗？}来使他羞愧，从此以后就仿佛对高贵的孩子说话。因为他说过\BibleQuote{身体是为主的}，现在则更清楚地指明这一点。他在另一处也做了同样的事，说：\BibleRef{格前 12:27}\BibleQuote{你们是基督的身体，各自都是肢体。}他经常使用这个比喻，但目的不同：有时显示祂的爱，有时增加他们的恐惧。但在这里他用来使他们震惊和充满惊慌。\BibleQuote{那么，我可以把基督的肢体拿去，作为娼妓的肢体吗？断乎不可！}没有什么比这句话更令人恐惧了。他没有说：\Quote{我可以把基督的肢体拿去，连接到娼妓身上吗？}而是说什么？\Quote{作为娼妓的肢体吗？}这无疑更尖锐地刺痛人。\par

然后他说明行淫者如何成为这样，这样说：\BibleQuote{你们不知道那与娼妓结合的，成为一体吗？}这如何显明？\BibleQuote{因为祂说：二人成为一体。}\par

% structure: p0009 source=h2 target=subsection reason=relative-heading-rank; base=h2

\subsection{格前 6:17}

\BibleQuote{但与主结合的，成为一神。}\par

因为结合不再让两者成为两个，而是使两者都成为一。\par

2. 现在请注意，他如何仅仅通过术语来推进，以娼妓和基督的名字进行指控。\par

% structure: p0013 source=h2 target=subsection reason=relative-heading-rank; base=h2

\subsection{格前 6:18}

\BibleQuote{你们要逃避奸淫。}\par

他没有说：\Quote{你们要戒避奸淫}，而是\Quote{你们要逃避奸淫。}\BibleQuote{凡人所犯的罪，都在身体以外；但那犯奸淫的，却是得罪自己的身体。}这比之前的话轻一些；但由于他必须对行淫者说话，他从各方面、从更大的和更小的事物来扩大那罪过的严重性，使之显得可憎。事实上，前一个话题是对更虔诚的人说的，但这一话题是对较弱的人说的。因为这也是保禄智慧的特点：不仅用伟大的事物使人羞愧，也用较小的事物以及可耻和不体面的考量。\par

\Quote{那么，}你说，\Quote{杀人者岂不玷污他的手吗？贪财者和勒索者呢？}我想每个人都清楚。但由于无法说出比行淫者更糟糕的事，他以另一种方式扩大罪行，说在行淫者身上，整个身体都被玷污。因为它就像掉进污秽的容器里并被浸没在污秽中一样污染。我们也是如此。因为从贪财和勒索，没有人会急忙去洗澡，而是好像什么也没发生就回家。但从与娼妓的交往，如同变得完全不洁，他去洗澡。因此良心从这罪中保留了一种不寻常的羞耻感。然而两者都是坏的，贪财和奸淫，都使人堕入地狱。但保禄以良好的管理做每一件事，所以他用他手头的所有话题来夸大奸淫的罪。\par

% structure: p0017 source=h2 target=subsection reason=relative-heading-rank; base=h2

\subsection{格前 6:19-20}

\BibleQuote{你们不知道你们的身体是住在你们内的圣神的宫殿吗？}他没有仅仅说\Quote{圣神的}，而是\Quote{在你们内的}，这也是一种安慰。然后，他进一步解释自己，补充说\BibleQuote{你们从天主那里领受了。}他也提到了施予者，既提升听者又使他恐惧，既通过那寄存物的伟大，也通过那施予者的慷慨。\par

\BibleQuote{你们不是自己的主人。}这不仅使人羞愧，甚至迫使人走向德行。\Quote{为什么，}他说，\Quote{你照你所愿的行吗？你不是自己的主人。}他说这些话，并不是要剥夺自由意志。因为他说\BibleQuote{凡事都可行，但不都有益处}时，他没有剥夺我们的自由。这里再次写道\BibleQuote{你们不是自己的主人}，他并没有侵犯选择的自由，而是引导人远离邪恶并指明主的眷顾。因此他补充说：\BibleQuote{因为你们是用重价买来的。}\par

但若我不是属于我自己的，你们凭什么向我要求应尽的本分呢？你们为何又接着说：\Quote{所以你们要在你们身上和你们的心灵上光荣天主，因为二者都属于天主。}那么，\Quote{你们不是属于你们自己的}是什么意思？他借此想要证明什么？为使他们在抗拒罪恶和顺从不当的思想欲望上处于稳固的状态。因为我们确实有许多不当的意愿，但我们必须抑制它们，因为我们能够做到。如果我们做不到，劝诫就是徒然的。因此，请注意他如何巩固自己的立场。因为说了\Quote{你们不是属于你们自己的}之后，他没有接着说\Quote{而是被迫的}，而是说\Quote{你们是用代价买来的。}你为何这样说？当然，从另一个角度来说，也许有人会说，你应该说服人，指出我们有一位主人。但这对于希腊人也与我们一样是共同的；而\Quote{你们是用代价买来的}这句话却是我们特有的。因为他提醒我们恩惠的伟大和救恩的方式，表明当我们成为外人时，是被\Quote{买来}的；而且不单单是\Quote{买来}，而是\Quote{用代价}。\par

\Quote{所以，你们要承担并在你们身上和你们的心灵上光荣天主。}他说这些话，是要我们不仅在身体上逃避淫乱，而且也要在心灵的精神上戒除一切邪恶的念头，不把恩宠赶走。\par

\Quote{二者都属于天主。}因为他说了\Quote{你们的}，所以又加上\Quote{属于天主的}，不断提醒我们：所有的一切，身体、灵魂和精神都属于主。因为有些人说，\Quote{在精神上}这句话是指神恩；如果那神恩在我们内，天主就受到了光荣。这将会实现，如果我们有一颗洁净的心。\par

祂之所以称这些东西是天主的，不仅因为祂创造了它们，而且因为当它们疏离时，祂第二次又把它们赢回，付出了圣子的血作为代价。请注意祂如何在基督内使一切达于圆满，如何把我们提升到天上。他说：\Quote{你们是基督的肢体，你们是圣神的宫殿。}不要成为\Quote{娼妓的肢体}，因为受侮辱的不是你们的身体；它根本不是你们的身体，而是基督的。他说这些话，既是为了显明祂的慈爱——我们的身体是属于祂的，也是为了使我们脱离一切邪恶的放纵。因为如果身体是属于别人的，他说：\Quote{你无权侮辱别人的身体，尤其是当它是主的时候；也不可玷污圣神的宫殿。}因为如果有人侵入私宅，狂欢乱闯，他必定要受到最严厉的惩罚；那么，若有人使君王的宫殿成为盗贼的巢穴，他将会遭受何等可怕的事呢？\par

因此，你当思考这些事，敬畏那住在你内的那一位。因为祂是\OriginalTerm{护慰者}{Paraclete}。当为那拥抱你、与你结合的那一位而战栗；因为祂是基督。你果真使自己成为基督的肢体了吗？你当这样思想，并保持贞洁：它们曾是某某人的肢体，现在成了谁的肢体。从前它们是娼妓的肢体，基督却使它们成为祂自己身体的肢体。因此，你从今以后对它们不再有主权。要为那释放了你的一位服务。\par

譬如你有一个女儿，你在极端疯狂中把她交给拉皮条的人赚钱，使她过娼妓的生活；然后有一位王子经过，把她从那奴役中解救出来，并娶她为妻；那么你从此不能再把她带回妓院。因为你已经一次把她交出去了，卖掉了。我们的情况也是如此。我们把自己的肉体出租给魔鬼——那个可恶的拉皮条者——以换取利益；基督看见了，便释放了它，把它从那邪恶的暴政下解救出来；它不再属于我们，而是属于拯救它的一位。如果你愿意把它当作王子的新娘来使用，没有人会阻止；但如果你把它带回到原来的地方，你必将遭受那些犯如此暴行的人应得的惩罚。因此，你更应该装饰它，而不是玷污它。因为对于邪恶的情欲，你对肉体没有权柄，只有在天主所吩咐的事情上才有权柄。至少要让这样的思想进入你心：天主从多么大的暴行中拯救了它。事实上，从来没有一个娼妓像我们的本性以前那样无耻地暴露自己。因为抢劫、谋杀和一切邪恶的思想进入并与灵魂同寝，而报酬仅仅是一点微小的、庸俗的眼前快乐。因为灵魂与一切邪恶的计谋和行为混杂在一起，所得的回报仅此而已。\par

然而，在此以前，虽然成为这样的人是坏的，但还不至于如此之坏；但是在天堂、在君王的宫廷、在分享了可敬畏的\OriginalTerm{圣事}{Mysteries}之后，再受玷污，这还有什么可赦免的呢？或者，你不认为那些贪婪的人，以及他先前列举的所有那些人，都有魔鬼与他们同寝吗？你不认为那些为了污秽而打扮自己的女人与魔鬼交往吗？谁能反驳这句话呢？如果有人争辩，让他揭开那些行为不当的女人的灵魂，他必然会看见邪恶的恶魔紧贴着她们。因为弟兄们，这是难的，也许是甚至不可能的：当身体这样打扮时，灵魂同时也得到装扮；而必定是一个被忽略，另一个被照顾。因为本性不允许这两件事同时发生。\par

4. 故此他说：\Quote{那与娼妓结合的，成为一体；但那与主结合的，成为一神。}因为这样的人从那时起就成为神，虽然身体仍包围着他。因为当没有任何属肉体、粗俗或属世的东西在他周围时，身体仅仅包围着他；因为他的整个治理在于灵魂和神。这样，天主就受到了光荣。因此，在祈祷中我们受命说：\Quote{愿你的名被尊为圣；}基督也说：\Quote{你们的光当在人前照耀，好使他们看见你们的善行，光荣你们在天圣父。}\par

诸天也不出声，却借其景象引人赞叹，将光荣归于伟大工程师，以此光荣祂。我们也当如此光荣祂，甚至比诸天更甚，因为我们若愿意，就能做到。因为诸天、白昼、黑夜光荣天主，远不如一个圣善的灵魂。正如有人仰观天空之美，说道：\Quote{愿光荣归于你，天主！你造了何等美妙的化工！}同样，当人目睹任何人的德行时，也会发出赞叹，而且后者的赞叹更甚。因为并非所有人都从受造物光荣天主，许多人甚至断言万物是自行运动的，又有人将世界和眷顾的工程归咎于魔鬼；这些人确实大大地、不可原谅地错了。但论到人的德行，没有人能持这些无耻的意见，反而当看见事奉主的人生活在良善中时，必定光荣天主。因为当一个人具有人的本性，分有我们共同的本性，生活在众人之中，却如金刚石般不为情欲的群集所动，谁不惊讶？当他身处火、铁、野兽之中，却比金刚石更坚硬，为了虔诚之道的缘故征服一切？当他受伤害时祝福，受诽谤时赞美，受恶待时为伤害他的人祈祷，遭人暗算时却善待那些与他争斗、设网罗的人？这些事以及类似的事，远比诸天更能光荣天主。因为外邦人观看诸天时不觉敬畏，但看见一个圣人极其严厉地过着严谨的生活时，便畏缩退避，自责不已。既然那与己同性的圣者远远高出他们（其差别远甚于天地之别），他们便不由自主地承认，这是神圣的力量在运作。因此祂说：\Quote{光荣你们在天的父。}\par

5. 你也愿意从另一处学习如何通过祂仆人的生活以及如何通过奇迹光荣天主吗？拿步高曾把三位青年投入火窑。当他看见火焰没有制胜他们时，他说道（正如\BibleRef{达 3:28}所记载的）：\Quote{愿天主受颂扬，他派遣了他的天使，拯救了他的仆人们脱离火窑，因为他们信赖他，并且改变了君王的谕旨。}他问道：\Quote{你怎么说？你被轻视，却赞美那些唾弃你的人吗？}他回答说：\Quote{是的，正因我被轻视。}对这件奇事他给出了这个理由。所以当时光荣归于天主，不仅因为奇迹，也因为这些投入火窑之人的意向。若有人按事物本身来考察这一点和那一点，就会看出前者并不亚于后者：因为说服灵魂勇敢面对火窑，在奇迹方面并不亚于从火窑中拯救出来。因为世界的皇帝，周围有如此多的武装、军团、将军、总督、执政官，海陆均受其统治，竟被被掳的青年轻视；受捆绑的竟战胜了捆绑者，征服了全部军队，这怎能不令人惊奇呢？国王及其随从没有任何力量达成所愿，即使以火窑为同盟也不成。但那些赤身露体、身为奴隶、异乡人、人数稀少（有什么比三人更可轻看呢？）的人，虽带着锁链，却战胜了无数的军队。因为死亡此时已被轻视，基督即将开始住在世界。正如太阳即将升起时，即使光芒尚未显现，白昼之光已渐明亮；同样，当正义的太阳即将来临时，死亡从那时起便开始退缩。有什么比那个场面更辉煌？有什么比那胜利更引人注目？有什么比他们的那些新战利品更显著？\par

同样的事在我们的时代也发生了。如今也有一个巴比伦火窑的君王，他也点燃比那更猛烈的火焰。如今也有这样一个像，有人命令人崇拜它。他身旁有总督、士兵和魅惑的音乐。许多人赞赏这像，它是如此多样、如此巨大。因为那像与贪婪有几分相似，贪婪连铁也不放过，但与构成它的材料不同，它命令人赞赏一切，包括铜、铁，以及比它们更普通的东西。\par

但正如这些事一样，如今也有一些效法这些青年的人，他们说：\Quote{你的神祇我们不事奉，你的偶像我们不崇拜；但为了天主的法律，我们忍受贫穷的火窑和其他一切困苦。}而富人方面，正如当时的人那样，常常也崇拜这像，并自身被焚烧。但那些一无所有、虽处贫穷却仍在这像面前的人，比生活在富足中的人更如在甘露之中。正如当时那些把别人投入火中的人被烧毁，而处于火中的人却如同在雨露之中。那时那暴君也比那些青年更被火焰焚烧，他的怒气猛烈焚烧着他。至于他们，火焰甚至无力触及他们的发梢；但比那火焰更猛烈的是，怒气焚烧了他的心神。请想一想，在众人面前，他竟被被掳的青年藐视，这是何等的事。这是一个征兆，表明他攻占他们的城并非凭自己的大能，而是由于他们中间众人的罪恶。因为若他不能制伏这些带着锁链、被投入火窑的人，那么如果犹太人都是这样的人，他又怎能通过常规战争制伏他们呢？由此清楚看出，众人的罪恶出卖了城。\par

6. 但也请注意孩子们不贪图虚浮的荣耀。因为他们并没有跳进火窑，而是事先遵守了基督的诫命，祂说：\BibleRef{玛 26:41} \Quote{祈祷，免得陷于诱惑。}他们被带到火窑时也不退缩，反而英勇地站在当中，既未受召就不争战，受召时也不怯懦：他们准备好一切，高贵而充满全部放胆直言的精神。\par

但让我们也听听他们说什么，好从中学到他们的崇高精神。\BibleRef{达 3:17} \Quote{有一位天主在天上，能拯救我们：}他们毫不为自己着想，即使即将被烧，所想的全是天主的荣耀。因为他们的话意思是：恐怕我们若被焚烧，你们会归咎于天主的无能，我们现在向你们准确宣告我们全部的教义。\Quote{有一位天主在天上，}不像这地上的像，这无生命且哑然的东西，却能拯救人脱离燃烧的火窑。那么，不要因祂容许我们落入火窑而谴责祂的无能。祂如此强大，以致我们堕落后，仍能将我们从火焰中救出。\Quote{即便不然，大王，你应当知道，我们决不侍奉你的神，也不朝拜你所立的金像。}请注意，他们因特殊的安排而不知道未来：因为若他们预知，他们的行为就毫无惊人之处。若他们已有安全的保证，又能藐视一切恐怖，那有何奇妙？那时天主固然因能从火窑中拯救而受光荣，但他们却不会受赞叹，因为他们不会投身于危险之中。因此，祂容许他们不知道未来，好更光荣他们。正如他们\OriginalTerm{告诫}{\GreekText{ἠσψαλίξοντο}}国王不要因他们可能被烧而谴责天主无能，天主也实现了两个目的：显示自己的权能，并使孩子们的热忱更显突出。\par

那么，他们的疑惑和没有把握自己必得保全从何而来？因为他们自认为太卑微，不配得这样的恩惠。为证明我并非猜测：当他们落入火窑时，他们这样哀叹自己，说：\BibleRef{达 4:6} \Quote{我们犯了罪，行了不义，我们不能开口。}因此他们说：\Quote{即便不然。}但如果他们没有明白说出这句话，即\Quote{天主能拯救我们；但若祂不拯救我们，是因为我们的罪，祂不拯救我们；}，不要奇怪。因为在异邦人看来，他们似乎是在用自己的罪来掩饰天主的无能。因此他们只提到祂的权能，不说原因。此外，他们受过良好训练，不过分追究天主的判断。\par

他们带着这些话进入了火焰；既没有侮辱国王，也没有推倒雕像。因为勇敢的人应当如此，节制而温和；尤其在危险中，以免显得是出于愤怒和虚浮的荣耀而投入这样的争战；而是带着刚毅和镇定。因为行事傲慢的人会招致这些过错的嫌疑；但忍耐、被迫争斗、并以温良经受考验的人，不仅因其勇敢而受钦佩，其镇定和体贴也使他受同样赞扬。这正是他们那时所做的：展现出全部刚毅和温良，不为酬报、赏报或回报而做任何事。\Quote{即使祂不愿意拯救我们——这话确实如此——我们也不侍奉你的神。}因为我们已有酬报：被算作配得远离一切不敬，并为此献出身体被焚烧。\par

我们既已有了酬报（因为的确我们因蒙受了对祂的圆满认识，蒙受作了基督的肢体，已得了酬报），就应谨慎，不使这些肢体成为娼妓的肢体。因为我们必须以这句极其可畏的话来结束我们的讲道，以便在威胁的恐惧充分生效下，我们比黄金更纯洁，这恐惧帮助我们如此。因为这样我们才能脱离一切邪淫，得以看见基督。愿天主恩赐我们众人在那日得以放胆瞻仰祂，藉着我们的主耶稣基督的恩宠和慈爱；愿光荣归于祂，至于无穷之世。阿们。\par

\SourceRecord{Translated by Talbot W. Chambers.；Nicene and Post-Nicene Fathers, First Series；Vol. 12.；Edited by Philip Schaff.；Buffalo, NY: Christian Literature Publishing Co.,；1889.；<http://www.newadvent.org/fathers/220118.htm>.}

% structure: p0001 source=h1 target=section reason=page-title

\section{\WorkTitle{论格林多前书讲道集第十九篇}}

% structure: p0002 source=h2 target=subsection reason=relative-heading-rank; base=h2

\subsection{\BibleRef{格前 7:1-2}}

\BibleQuote{论到你们信上所提的事：男人不亲近女人，原是好。但为了避免淫乱，男人当各有自己的妻子，女人当各有自己的丈夫。}\par

在纠正了他们所受指控的三件最严重之事——其一，教会的分裂；其二，关于淫乱者；其三，关于贪婪者——之后，他转而采用较为温和的措辞。他插入一些关于婚姻与贞洁的劝勉与建议，使听众从较不愉快的话题中得到些许喘息。但在第二封书信中，他做了相反的事：他从较温和的话题开始，以更令人痛苦的话题结束。同样在此处，当他结束了关于贞洁的讨论后，再次转向更接近责备的内容；但他并非按固定顺序逐一列出，而是根据时机与所涉事宜的紧迫性，以两种方式交替变化其论述。\par

因此他说：\BibleQuote{论到你们信上所提的事。}因为他们曾写信问他：\Quote{是否应当远离自己的妻子？}他在回信中答复了此事，并制定了关于婚姻的规则，也引入了关于贞洁的论述：\BibleQuote{男人不亲近女人，原是好。}他说：\Quote{你若探问何者为卓越且远胜之道，那就是根本不与女人有任何交往；但若你询问何者为安全且有助于自身软弱之道，那就通过婚姻结合。}\par

但由于很可能——正如现在也常发生的那样——丈夫愿意而妻子不愿意，或反之，请注意他如何分别讨论每种情况。有些人认为，这番话是向司铎们说的。但根据下文判断，我不能肯定如此；因为他不会以一般性措辞给予劝告。若他仅为司铎写这些，他定会说：\Quote{教师不亲近女人，原是好。}但他现在使此话具有普遍适用性，说：\Quote{男人；}不仅是司铎。又说：\BibleQuote{你脱离了妻子吗？不要寻求妻子。}他没有说：\Quote{你身为司铎与教师，}而是不加限定。整篇讲话完全以同样的语调进行。在说\BibleQuote{为了避免淫乱，男人当各有自己的妻子}时，借由这一让步所依据的原因，他引导人们走向节欲。\par

% structure: p0007 source=h2 target=subsection reason=relative-heading-rank; base=h2

\subsection{\BibleRef{格前 7:3-4}}

2. \BibleQuote{丈夫当还妻子应尽的义务；妻子也同样还丈夫应尽的义务。}\par

那么\Quote{应尽的义务}是什么意思？\BibleQuote{妻子没有权力掌管自己的身体；}她既是丈夫的奴仆，也是他的主人。你若回避应尽的义务，便得罪了天主。你若想暂时退避，必须得到丈夫的许可，即使只是短暂的时间。正因如此，他称此事为债务，以表明没有人是自己身体的主人，而是彼此互为仆人。\par

因此，当你看见妓女诱惑你时，你要说：\Quote{我的身体不属于我，而属于我的妻子。}同样，妻子也当对企图破坏她贞洁的人说：\Quote{我的身体不属于我，而属于我的丈夫。}\par

若夫妻连自己的身体都没有权柄，何况财产呢？所有有丈夫和所有有妻子的人，请听：若你们连自己的身体都不视为己有，何况钱财呢？\par

在别处，我承认，无论在\WorkTitle{新约}还是\WorkTitle{旧约}中，他都赋予丈夫较多的优先地位，例如\WorkTitle{旧约}中说（\OriginalTerm{你的归向}{\GreekText{ἡ} \GreekText{ἀποστροφή} \GreekText{σου}}，\BibleRef{创 3:16}）\BibleQuote{你的归向必指向你的丈夫，他要管辖你。}保禄也是如此，他作出区分并写道：\BibleRef{弗 5:25}\BibleQuote{作丈夫的，要爱你们的妻子；并要妻子敬畏丈夫。}但在此处，我们不再听到\Quote{更大}与\Quote{更小}之分，而是同一权利。这是为何？因为他的讲话是关于贞洁的。\Quote{在其他一切事上，}他说，\Quote{让丈夫享有特权；但在涉及贞洁的问题上，不可如此。}\BibleQuote{丈夫没有权力掌管自己的身体，妻子也没有。}这是极大的平等，没有特权。\par

% structure: p0013 source=h2 target=subsection reason=relative-heading-rank; base=h2

\subsection{\BibleRef{格前 7:5-9}}

3. \BibleQuote{你们不可彼此亏负，除非两相情愿。}\par

那么这是什么意思？他说：\Quote{如果丈夫不愿意，妻子不可行节欲；同样，若妻子不同意，丈夫也不可行。}为何如此？因为这种节欲往往产生大恶。奸淫、淫乱和家庭的破裂常由此而起。如果男人有自己的妻子尚且犯奸淫，那么你若剥夺他们这慰藉，岂不更甚？他恰当地说：\BibleQuote{不可亏负}；此处用\Quote{亏负}，上文用\Quote{债务}，为表明所涉支配权的严格性。因为未经对方同意而行节欲，就是\Quote{亏负}；但若对方同意，则不然：正如我若说服你，然后拿走你任何东西，我不认为这是亏负。因为只有违反他人意愿、强行取走者才是亏负。许多妇女正是如此，她们所行的是罪而非义，进而对丈夫的不洁承担责任，并使一切分裂。其实她们应将和谐置于一切之上，因为这比所有其他事都更重要。\par

若你们愿意，我们不妨用实际例子来思考。例如，假设有夫妻二人，妻子在没有丈夫同意的情况下守节；那么，如果丈夫因此犯奸淫，或者虽未犯奸淫却烦躁不安、焦躁、争吵，给妻子带来各种烦恼；那么，在爱情破裂的情况下，禁食和节欲的益处在哪里呢？毫无益处。因为必然会生出多么奇怪的指责、多少麻烦、多大的争战！因为当屋内的夫妻不和时，这屋子不会比暴风雨中的船更好，那时主人与舵手关系不睦。因此他说：\Quote{不要彼此亏负，除非两相情愿，暂时分居，为要专心祈祷。} 他这里指的是极其热切的祈祷。因为如果他禁止有夫妻关系的人祈祷，那么\Quote{不断祈祷}又如何能实现呢？可见，与妻子同住仍可专心祈祷。但节欲使祈祷更完美。因为他没有仅仅说：\Quote{为要祈祷}，而是说：\Quote{为要专心于此}；好像他所说的不是引起不洁，而是许多事务。\par

\Quote{然后再在一起，免得撒殚诱惑你们。} 因此，为使这看起来不是明令，他补充了理由。这理由是什么？\Quote{免得撒殚诱惑你们。} 为使你明白，造成这种罪的不仅是魔鬼，我说的是奸淫，他接着说：\Quote{因为你们的不能自持。}\par

\Quote{但这是我说的，是许可的话，不是命令。因为我愿意众人像我一样，} 处于节欲的状态。他在许多地方遇到难以劝告的事就这样做：他提出自己为例，说：\Quote{你们该效法我。}\par

\Quote{然而每人从天主领受了自己的恩赐，有人这样，有人那样。} 既然他曾严厉斥责他们说：\Quote{由于你们的不能自持}，他再次安慰他们说：\Quote{每人从天主领受了自己的恩赐}；这不是说在这德行上我们不需要努力，而是如我之前所说，为了安慰他们。因为如果这是一项\Quote{恩赐}，而且人对此毫无贡献，那么你为何说：\Quote{但是，\BibleRef{格前 7:8} 我对未婚和寡妇说：如果他们像我一样，对他们就有益；\BibleRef{格前 7:9} 但如果他们不能自持，就让他们结婚？} 你看见保禄的强烈感觉了吗？他既表明守节更好，又不向不能达到的人施加任何压力，恐怕引起什么反感？\par

\Quote{因为结婚比欲火中烧更好。} 他指出了情欲的肆虐是何等巨大。他的意思大致是：\Quote{如果你要忍受极大的暴力和欲火，就从痛苦和辛劳中退出来，免得你被倾覆。}\par

% structure: p0021 source=h2 target=subsection reason=relative-heading-rank; base=h2

\subsection{格前 7:10-11}

4. \Quote{至于那些已经结婚的，我吩咐他们——其实不是我，而是主吩咐——}\par

因为基督亲口明确制定的法律，他即将读给他们听，关于除非因奸淫不得休妻的事；\BibleRef{玛 5:32}因此他说：\Quote{不是我。} 诚然，之前所说的话虽未明确陈述，却也是祂的旨意。但请看，这一条祂已明显宣示了。所以\Quote{我}和\Quote{不是我}的区别就在这里。为使你不会认为他自己的话也是人的话，他因此补充说：\Quote{因为我想我也有天主的圣神。}\par

那么，什么是他所说的\Quote{主对已婚者的命令？妻子不可离开丈夫：\BibleRef{格前 7:11}但如果她离开了，就当独身，或与丈夫和好。} 在这里，他看到由于守节和其他借口，以及因为\OriginalTerm{心窄}{\GreekText{μικροψυχίας}}的软弱，离异便发生了；他说，这样的事最好根本不要发生；但若发生，妻子应与丈夫同住，即使不和好同居，至少不要引进另一人为她丈夫。\par

% structure: p0025 source=h2 target=subsection reason=relative-heading-rank; base=h2

\subsection{格前 7:12-14}

\Quote{至于其余的人，我说，不是主说：若某弟兄有不信的妻子，而她同意与他同住，就不要离弃她；若某妇人有不信的丈夫，而他同意与她同住，就不要离弃他。}\par

因为当谈论与淫乱者分离时，他藉对言语加以修正，使事情变得容易，说：\Quote{然而，不是绝对不与这世界上的淫乱者；} 同样在此，他也为这责任的极大轻松预备了条件，说：\Quote{若某妻子有不信的丈夫，或丈夫有不信的妻子，不要离弃她。} 你说什么？\Quote{若他是不信者，就让他与妻子同住，但若是淫乱者，就不吗？而奸淫是比不信更轻的罪。} 我承认，奸淫是更轻的罪；但天主极其体恤你们的软弱。这正是他关于祭献所作的，说：\BibleRef{玛 5:24} \Quote{放下你的祭品，先与兄弟和好。} 关于欠一万塔冷通的仆人也是如此。他并没有因为欠他一万塔冷通而惩罚他，而是因为向他同伴追讨一百银钱而报复了他。\par

然后，免得妻子害怕，以为因与丈夫同居而不洁，他说：\Quote{因为不信的丈夫因妻子而被圣化，不信的妻子因丈夫而被圣化。} 然而，如果\Quote{与娼妓结合的是成为一体}，很明显，与偶像崇拜者结合的女人也是一体。诚然是一体；但她并未因此不洁，反而妻子的洁净胜过了丈夫的不洁；同样，相信的丈夫的洁净胜过了不信的妻子的不洁。\par

那么，在这种情况下，不洁是如何被克服的，因此婚姻关系被允许；而在卖淫的妇人身上，丈夫将她赶出并不受谴责？因为在前者中，有希望失落的肢体可通过婚姻得救；但在后者中，婚姻已然破裂；而在那里，双方都败坏了；但在这里，过错只在于二者之一。我的意思是这样的：犯了淫乱的妇人完全可憎；如果\BibleQuote{与娼妇结合的，成为一体}\BibleRef{格前 6:16}，那么与她结合的人也因与娼妇交合而变得可憎；因此一切纯洁都飞逝了。但在我们面前的情况并非如此。那么是怎样的呢？拜偶像者是不洁的，但妇人并不不洁。如果她确实与他在他不洁的那方面，即他的不敬虔上合伙，那么她自己也会变得不洁。但现在拜偶像者在不洁的某一方面是，而妻子与他交往的另一方面他并不不洁。因为婚姻和身体的结合正是交往所在。\par

再者，这个人有希望被他的妻子挽回，因为她完全属于他；但对另一个人来说，这并不十分容易。因为那个先前羞辱他、成为别人的、并破坏了婚姻权利的妇人，如何能挽回她曾错待的他呢？况且他仍然对她如同外人？\par

再者，在那种情况下，淫乱之后丈夫不再是丈夫；但在这里，即使妻子是拜偶像者，丈夫的权利并未被破坏。\par

然而，他并非简单地推荐与不信者同居，而是有条件地，即他愿意如此。因此他说：\BibleQuote{他自己也情愿与她同住。}\BibleRef{格前 7:12}请告诉我，当敬虔的职责不受损害，且对不信者有良好希望时，已经结合的人如此继续，不引起不必要的争战，有什么害处呢？因为现在的问题不是关于从未结合的人，而是关于已经结合的人。他没有说：\Quote{如果有人想娶不信的妻子}，而是说：\BibleQuote{如果有人有不信的妻子。}\BibleRef{格前 7:12}意思是：如果有人在结婚后接受了信仰之言，而另一方继续不信，仍然渴望与他们同住，就不要中断婚姻。因为他说：\BibleQuote{不信的丈夫因妻子而成圣。}\BibleRef{格前 7:14}你们的纯洁是如此超量。\par

那么，希腊人是圣的吗？当然不；因为他不是说\Quote{他是圣的}，而是说：\BibleQuote{他因妻子而成圣。}\BibleRef{格前 7:14}他说这话并非表明他是圣的，而是要尽可能彻底地将妇人从恐惧中释放出来，并引导那人渴望真理。因为不洁不在发生共融的身体里，而在心思和意念中。接下来就是证据：因为如果你继续不洁而有后代，那孩子不仅仅出于你，当然就是不洁或半洁。但现在它并非不洁。为此他补充说：\BibleQuote{不然，你们的儿女就是不洁的；但如今他们是圣洁的。}\BibleRef{格前 7:14}即不洁的。但宗徒称他们为\Quote{圣洁}，通过强调的表达再次消除由那种怀疑产生的恐惧。\par

% structure: p0034 source=h2 target=subsection reason=relative-heading-rank; base=h2

\subsection{格林多前书 7:15}

\BibleQuote{倘若那不信的人离开，就由他离开吧。}\BibleRef{格前 7:15}因为在这种情况下，事情不再属于淫乱。但\Quote{倘若那不信的人离开}是什么意思？例如，如果他命令你因婚姻而献祭并参与他的不敬虔，否则分道扬镳；那么婚姻被废止，不在敬虔上造成裂痕，这样更好。因此他补充说：\BibleQuote{在这种情况下，兄弟或姐妹都不受束缚。}\BibleRef{格前 7:15}如果他每天为此事烦扰你、持续争斗，那么最好分离。因为他暗示这一点，说：\BibleQuote{但天主召叫我们原是召叫我们平安。}\BibleRef{格前 7:15}因为分离的理由是对方提供的，如同犯了不洁的人一样。\par

% structure: p0036 source=h2 target=subsection reason=relative-heading-rank; base=h2

\subsection{格林多前书 7:16}

\BibleQuote{因为妻子，你哪里知道能救你的丈夫呢？}\BibleRef{格前 7:16}这再次指向那话：\BibleQuote{她不可离开他。}\BibleRef{格前 7:12}也就是说：\Quote{如果他不制造混乱，就留下，}他说，\Quote{因为这甚至有好处；留下，给予建议和劝告，说服。}因为没有任何教师能像妻子那样有能力说服（Reg. \OriginalTerm{说服}{\GreekText{πεῖσαι}}，Bened. \OriginalTerm{得胜}{\GreekText{ἰσχῦσαι}}）。而且，一方面，他没有将任何必须强加给她，绝对地要求她达到这一点，以免再次做出过于痛苦的事；另一方面，他也没有叫她绝望；但他通过未来的不确定性使事情悬而未决，说：\Quote{因为妻子，你哪里知道能救你的丈夫呢？或者丈夫，你哪里知道能救你的妻子呢？}\par

% structure: p0038 source=h2 target=subsection reason=relative-heading-rank; base=h2

\subsection{格林多前书 7:17-22}

5. 再者，\BibleQuote{只要照主所分给各人的，和天主所召各人的而行。有人蒙召时是受割礼的吗？就不要变成未受割礼的。有人蒙召时是未受割礼的吗？就不要受割礼。受割礼算不得什么，不受割礼也算不得什么；只有遵守天主的诫命才是要紧的。各人蒙召的时候是什么身份，仍要守住这身份。你是作奴隶蒙召的吗？不要因此忧虑。}\BibleRef{格前 7:17-22}他说，这些事对信仰毫无贡献。所以不要争竞，也不要烦恼；因为信仰已经废去了这一切。\par

各人蒙召的时候是什么身份，仍要守住这身份。你蒙召时有不信的妻子吗？继续拥有她。不要为信仰的缘故赶出你的妻子。你蒙召时是奴隶吗？不要为此忧虑。继续作奴隶。你蒙召时是未受割礼的吗？保持未受割礼。你蒙召时受过割礼了吗？继续受割礼。因为这就是\BibleQuote{照主所分给各人的}\BibleRef{格前 7:17}的意思。因为这些东西不是敬虔的阻碍。你蒙召时是奴隶；另一个人有不信的妻子；另一个人受过割礼。\par

令人惊奇！他将奴役置于何处？正如割礼毫无益处，未受割礼也无损害；同样，奴役也无害处，自由也无益处。为了极其清楚地指出这一点，他说：\Quote{但即使（\OriginalTerm{但即使你能}{\GreekText{Αλλ}' \GreekText{εὶ} \GreekText{καὶ} \GreekText{δυνάσαι}}）你能获得自由，还是保持原状吧：}也就是说，宁愿继续为奴。现在，他凭什么理由告诉那个可能获得自由的人继续为奴？他的意思是要指出，奴役并无害处，反而是一种益处。\par

现在我们并非不知道，有些人说，\Quote{还是保持原状吧}这句话是针对自由而言的，将其解释为\Quote{如果你能获得自由，就获得自由。}但如果他想要表达这个意思，这种说法就非常违背保禄的作风。因为当他安慰奴隶并表明他毫发无损时，他不会叫他去获得自由。因为也许有人会说，\Quote{那么，如果我无法获得自由呢？我是一个受损害、被贬低的人。}所以，他说的不是这个意思；而是如我所说，要指出一个人获得自由毫无益处，他说：\Quote{即使你有能力获得自由，还是宁愿为奴。}\par

接着他又补充原因；\Quote{因为那在主里蒙召的奴仆，是主所释放的人；同样，那蒙召的自由人，是基督的奴仆。}\Quote{因为，}他说，在关于基督的事情上，两者都是平等的。那么，奴仆如何是自由人呢？因为祂不仅将你从罪恶中释放，也当你继续为奴时，将你从外在的奴役中释放。因为祂不允许奴仆为奴，即使他仍处于奴役状态，这正是伟大的奇迹。\par

但奴仆如何一面继续为奴，一面又是自由人呢？当他从情欲和心灵的疾病中释放出来时；当他轻视财富、愤怒和所有其他类似的情欲时。\par

% structure: p0045 source=h2 target=subsection reason=relative-heading-rank; base=h2

\subsection{\BibleRef{格前 7:23-24}}

\Quote{你们是用代价买来的；不要成为人的奴仆。}这句话不仅是对奴隶说的，也是对自由人说的。因为一个为奴的人可能不是奴隶；一个自由人可能成为奴隶。\Quote{一个人如何是奴隶又不是奴隶？}当他为天主做一切时：当他毫无虚饰，也不出于对人的眼前服侍而做任何事时：这就是一个为奴于人的如何能自由。或者，一个自由人如何成为奴隶？当他服务于人的任何邪恶差事，无论出于贪食、贪财或官职的缘故。这样的人，尽管他是自由的，却比任何人都更像奴隶。\par

请思考这两点。\Person{若瑟}是奴隶，但不是人的奴隶；因此，即使在奴役中，他比所有自由人更自由。例如，他没有屈服于他的主母；没有屈服于拥有他的她所渴望的目的。再者，她是自由的；然而，没有人像她那样像一个奴隶，讨好乞求她自己的仆人。但她无法说服这个自由的人去做他不愿做的事。所以，这不是奴役；而是最高超的自由。因为他的奴役对他追求德行有何阻碍？让众人听，无论奴隶还是自由人。哪一个是奴隶？是被恳求的那位，还是恳求的那位？是祈求的那位，还是轻视她恳求的那位？\par

事实上，天主亲自为奴隶设定了界限；并且一个人应当遵守到何种程度也已确定，逾越它便是错误。也就是说，当你的主人命令任何不悦于天主的事时，你应当追随并服从；但仅此而已。因为这样奴隶便成为自由人。但你若再进一步，即使你是自由的，你已成为奴隶。至少他暗示了这一点，说：\Quote{不要成为人的奴仆。}\par

但如果这不是他的意思，如果他命令他们离弃主人并分争地争取自由，那么他劝诫他们时说\Quote{各人应留在自己蒙召时的身份中}是什么意思？在另一处，\BibleRef{弟前 6:1}\Quote{凡在轭下做奴仆的，应认为自己的主人配受一切尊敬；那些有信主的主人的，不可因他们是弟兄而轻视他们，因为他们分享益处。}并且写给厄弗所人和哥罗森人时，他也制定并要求同样的规则。由此清楚可见，他废除的不是这种奴役，而是那种由恶习引起的、也发生在自由人身上的奴役；这是最坏的一种奴役，即使一个自由人受其束缚。\Person{若瑟}的弟兄们从他们的自由中得到什么益处？他们岂不是比所有奴隶更为奴性：既对父亲说谎，又对商人用假话，还对他们的兄弟如此？但自由人却不是这样：相反，在任何地方、任何事上，他都是真实的。没有任何事物有权奴役他，无论是锁链、束缚、主母的爱，还是身处异乡。但他在任何地方都保持自由。因为当自由甚至在束缚中闪耀时，这才是最真实的自由。\par

6. 基督教就是如此：在奴役中赐予自由。正如那生来无懈可击的身体，在受到箭矢却毫发无伤时，才显出它的无懈可击；同样，那严格意义上的自由人，在主人之下仍不受奴役时，才显出他的自由。因此，他的命令是：为奴。但若一个奴隶不可能成为他所应有的基督徒，那么希腊人将谴责真宗教的极其软弱；然而，如果他们能被教导奴役丝毫无损于虔敬，他们将钦佩我们的教义。因为，若死亡、鞭打、锁链都不能伤害我们，那么奴役就更不能了。火、铁、无数的暴政、疾病、贫穷、野兽以及比这些更可怕的事物，都不能伤害信者；反而使他们更坚强。那么，奴役如何能伤害人呢？亲爱的，不是奴役本身伤害人；真正的奴役是罪恶的奴役。如果你不是这种意义上的奴隶，就勇敢并喜乐吧。没有人能有权做任何对你不利的事，因为你拥有不可奴役的情操。但如果你成为罪恶的奴隶，即使你一万次自由，你也白白拥有你的自由。\par

因为，请告诉我，你虽不受制于人，却屈服于自己的情欲，这有什么益处呢？人常常知道如何宽恕；但那些主人对你的毁灭永不知足。你受制于人吗？其实，你的主人也为你作奴仆：安排你的饮食，照顾你的健康，照料你的鞋子以及一切其他事物。你怕触怒主人，还不如他怕你缺少那些必需品。\Quote{但他坐着，你却站着。}那又怎样？这句话也可以说你和说他。常常，至少当你躺下安睡时，他不但站着，而且在市场上忍受无尽的痛苦；他比你更痛苦地失眠。\par

例如，若瑟从他的主母所受的苦，与她因邪恶欲望所受的苦相比，又算得了什么呢？因为他确实没有做她想要强加于他的事；但她却执行了她的主人——我指的是她的不洁之灵——所命令的一切：它一直不罢休，直到使她公开蒙羞。什么主人会命令这样的事？什么野蛮的暴君？\Quote{恳求你的奴仆，}就是这话：\Quote{讨好你用钱买来的人，哀求俘虏；即使他厌恶地拒绝你，再次围困他：即使你多次对他说话，他不同意，也要等他独自一人，强迫他，成为笑柄。}还有什么比这些话更不光彩、更可耻的呢？\Quote{如果通过这些方法仍无进展，那么，就诬告他，欺骗你的丈夫。}请看这些命令多么卑劣、可耻，多么无情、野蛮、疯狂。主人曾向奴仆下达过什么命令，如同她的淫欲当时向那位王后所下的命令那样？而她竟不敢违抗。但若瑟却没有遭受任何这类事，反而一切都带来光荣和尊荣。\par

你愿意再看另一个人受到严厉女主人命令的束缚，却无力违抗其中任何一条吗？请看加音，他的嫉妒给他下了什么命令。它命令他杀害兄弟，对天主撒谎，使父亲忧伤，抛弃羞耻；他都做了，毫不拒绝服从。为什么惊奇于这女主人在一个人身上有如此大的力量？她常常毁灭整个民族。例如，米德杨妇女擒获了犹太人，几乎将他们掳去；她们自己的美貌燃起情欲，成为她们征服那整个民族的手段。因此，保禄为了驱除这种奴役，说：\Quote{不要成为人的奴仆；}这就是说：\Quote{不要服从那些命令不合理之事的人：不，不要服从你们自己。}然后，他提高了他们的心思，使之升到高处，说：\par

% structure: p0054 source=h2 target=subsection reason=relative-heading-rank; base=h2

\subsection{\BibleRef{格前 7:25-26}}

7. \Quote{论童贞，我没有主的命令，但我将自己的意见，作为蒙主怜悯而成为忠信的人。}\par

他按着次序前进，接着谈到童贞。因为在论节制的话中磨练他们之后，他迈向更高的事，说：\Quote{我没有命令，但我认为这是好的。}什么原因呢？就是论节制时他所提过的同一原因。\par

% structure: p0057 source=h2 target=subsection reason=relative-heading-rank; base=h2

\subsection{\BibleRef{格前 7:27-34}}

\Quote{你已与妻子结合吗？不要寻求分离。你已脱离妻子吗？不要寻求妻子。}\par

这些话与前面所说的并不矛盾，反而是完全一致。因为他在那里也说：\Quote{除非双方同意，}正如这里说：\Quote{你已与妻子结合吗？不要寻求分离。}这并不矛盾。因为违反同意才会构成解除；但如果双方同意，且二人都过节制的生活，那便不是解除。\par

然后，免得这好像是制定法律，他接着引用\BibleRef{格前 7:28}说：\Quote{但如果你结婚，你没有犯罪。}接着他引证现状：\Quote{现在的困难，时间的短促，}和\Quote{痛苦。}因为婚姻带来许多事物，他在此及论述节制时都曾提及：那里他说：\Quote{妻子没有权力掌管自己的身体；}这里他说：\Quote{你已被结合。}\par

\Quote{但若你结婚，你没有犯罪。}他并不是在谈论选择了守童贞的人，因为若那样，她就犯罪了。因为如果寡妇在选择了寡居之后再缔结第二次婚姻要受谴责，何况童贞女呢？\par

\Quote{但这样的人肉身必有苦难。}\Quote{也有喜乐，}你可能会说：但请看他如何以时间的短促来限制它，引用\BibleRef{格前 7:28}说：\Quote{时间是短促的；}意思是：\Quote{我们现在被劝勉要离开前进，你们却更深入。}即使婚姻没有苦难，我们也应该朝着将来的事物努力。但当它还有痛苦时，何必给自己加上额外的负担呢？何必担负这样的重担，而背负它之后还得如同没有一样使用它呢？因为他说：\Quote{那些有妻子的，}该\Quote{像没有一样。}\par

然后，在插入一些关于将来的话题之后，他将话题转回现在。因为他的一些论点是属神的：如\Quote{一个关心丈夫的事，另一个关心天主的事。}另一些论及今世生活：如\Quote{我愿你们无挂虑。}但尽管如此，他仍将选择留给他们：因为若是证明了什么是最好的之后又回过头来强迫，就好像不信任自己的话一样。因此，他倒是以让步来吸引他们，并如下阻止他们：\par

% structure: p0064 source=h2 target=subsection reason=relative-heading-rank; base=h2

\subsection{\BibleRef{格前 7:35}}

我这样说，是为你们自己的益处，并不是要给你们设下圈套，而是要你们做合宜的事，并得以专心事主，毫无分心。愿童贞女们听见：童贞并不只凭这一点来界定；因为那挂虑世俗之事的，不能算是童贞女，也不能算是合宜的。因此，当他说：\BibleQuote{妻子与童贞女之间是有区别的}，他加上这一点作为区别，以及她们彼此区分之处。并且，在界定童贞女与非童贞女之时，他所指出的并非婚姻或节制，而是摆脱事务的闲暇与事务繁多。因为邪恶不在于同居，而在于对生活严谨的妨碍。\par

% structure: p0066 source=h2 target=subsection reason=relative-heading-rank; base=h2

\subsection{格前 7:36-40}

\BibleQuote{若有人以为对待他的童贞女不合宜}\par

他似乎在这里谈论婚姻；但他所说的一切都与童贞有关；因为他甚至允许第二次婚姻，说：\BibleQuote{只可在主内}。那么\Quote{在主内}是什么意思？就是带着贞洁与尊重：因为这在各处都是必要的，必须追求，否则我们不能看见天主。\par

现在我们若轻率地掠过他论童贞的话，愿无人指责我们疏忽；因为我们已就此主题写了一整本书，并在其中尽我们所能的精确，研究了此主题的每一个分支，我们认为在此重复是浪费言辞。因此，关于这些事情，我们请听众参考那部著作，在此只说这一件事：我们必须追求节制。因为，他说：\BibleQuote{你们要追求和平与圣化，没有圣化，谁也不能见主}。因此，为使我们被堪当见到祂，无论我们处于童贞、初婚或再婚，让我们追求这圣化，好能获得天国，藉着我们的主耶稣基督的恩宠和慈爱；愿光荣、权能、尊崇归于祂，偕同圣父及圣神，自今而始，直至永远万世。阿们。\par

\SourceRecord{Translated by Talbot W. Chambers.；Nicene and Post-Nicene Fathers, First Series；Vol. 12.；Edited by Philip Schaff.；Buffalo, NY: Christian Literature Publishing Co.,；1889.；<http://www.newadvent.org/fathers/220119.htm>.}

% structure: p0001 source=h1 target=section reason=page-title

\section{格林多前书第二十篇讲道}

% structure: p0002 source=h2 target=subsection reason=relative-heading-rank; base=h2

\subsection{格林多前书 8:1}

论及祭偶像之物：我们知道我们众人都有\OriginalTerm{知识}{knowledge}。\OriginalTerm{知识}{knowledge}使人自大，\OriginalTerm{爱德}{love}却立人。\par

首先有必要说明这段经文的意思：因为这样我们就能轻易理解宗徒的论述。因为凡看到有人被控诉，若不先明白过犯的性质，就不会理解所说的话。那么，他当时控诉格林多人什么呢？一个沉重的罪过，是许多恶的根源。嗯，是什么呢？他们中有许多人，得知\BibleRef{玛 15:11} \BibleQuote{不是入口的使人污秽，而是出口的，}且木石之偶像和魔鬼无力加害或帮助，便过度利用对此的完美知识，既损害别人，也损害自己。他们既进入偶像所在之处，并分享那里的筵席，从而产生了巨大且毁灭性的恶。一方面，那些仍保留对偶像的畏惧、不知如何藐视它们的人，因看见更完美的人这样做，也参与了那些宴席；因而受到了最大的伤害：因为他们不像其他人那样以同样的心思接触面前的食物，而是当作祭偶像之物；这事竟成了通向偶像崇拜的道路。另一方面，这些自命更完美的人，因参与魔鬼的筵席，也受到非同寻常的伤害。\par

这就是抱怨的缘由。现在这位有福之人正要纠正此事，却没有立即严厉发言；因为所行之事更多出于愚蠢而非邪恶：因此首先需要的不是严厉斥责和愤怒，而是劝诫。请在此留意他的明智，看他如何立即开始告诫。\par

\Quote{论及祭偶像之物，我们知道我们众人都有知识。}他撇下软弱的人——他常这样做——先与强壮的人交谈。他在\WorkTitle{罗马书}中也这样做，说：\BibleRef{罗 14:10} \BibleQuote{可是你为什么判断你的弟兄呢？}因为这样的人也能欣然接受责备。他在这里也完全一样。\par

首先，他通过宣告他们认为独有之事——拥有完美的知识——是众人共有的，来破除他们的自负。因此，他说：\BibleQuote{我们知道，我们众人都有知识。}因为如果他允许他们高抬自己，首先指出这事对他人有害，那么他对他们带来的害处会大于益处。因为野心勃勃的灵魂，当它自恃于某事时，即使这事伤害他人，它也因虚荣的暴政而固执己见。所以\Person{保禄}首先单独审查事物本身：正如他之前在论及外来的智慧时，以高压手段将其摧毁。但在那件事上，他按我们所期望的做了：因为那件事完全是可责难的，他的任务很容易。因此他指明它不仅无用，甚至与福音相悖。但在当前的事上，他不能这样做。因为所行的是关于知识，而且完美的知识。推翻它并不安全，然而不这样做就不可能除掉由此产生的自负。那他怎么办呢？首先，通过指出它是共有的，他抑制了他们那种膨胀的骄傲。因为拥有伟大卓越之物的人，当只有他们拥有时，会更加得意；但若证明他们与他人共同拥有，他们就不再有那么多这种感受了。因此，他首先使之成为公有物，因为他们认为只属于自己。\par

其次，在使之成为共有后，他并没有让自己单独与他们分享；因为这样他反而会抬举他们；正如唯一拥有者会使人得意，有一两个同侪（或许是领袖人物）也有同样的效果。为此他没有提到自己，而是提到众人：他没有说\Quote{我也有知识，}而是说\BibleQuote{我们知道我们众人都有知识。}\par

2. 这是一种方式，而且是第一种，他以此打击了他们的骄傲；接下来的方式更有力。那是什么呢？他表明，即使这事本身也不是完全完整的，而是不完美的，极其不完美。不仅不完美，而且有害，除非有另一物与之结合。因为说了\BibleQuote{我们有知识}之后，他加上了\BibleQuote{知识使人自大，爱德却立人：}因此，当知识没有爱德时，它便将人提升到绝对的傲慢。\par

\Quote{然而，即使爱德，}你会说，\Quote{没有知识也没有益处。}嗯：他没有说这话；而是把它当作众人公认之事略过，他表明知识极需爱德。因为爱德的人，既然履行了最绝对的诫命，即使有些缺陷，也会因他的爱德很快蒙受知识的恩赐；如\Person{科尔乃略}和许多其他人。但那有知识却没有爱德的人，不但一无所获，反而会被剥夺所拥有的，在许多情况下陷入傲慢。看来知识并不产生爱德，反而阻止那不加防备的人得到爱德，使他自高自大。因为傲慢惯于造成分裂，而爱德既凝聚人，又引人走向知识。为阐明这一点，他说：\BibleQuote{但如果有人爱天主，这人是蒙祂所知的。}因此他说：\Quote{我不是禁止这个，即你们拥有完美的知识；而是命令你们要有爱德地拥有它；否则不是收益，而是损失。}\par

你看他是如何开始为有关爱德的讲论奏出先声？因为既然这一切邪恶都源自下列根源，即：不是出于完美的知识，而是因为他们没有大大地爱邻人，也没有顾惜邻人；由此引出了他们的分歧和自满，以及他之前指控他们的所有其他事；在此之前和之后，他都不断地为爱德作预备；如此纠正一切美善的泉源。\Quote{现在，}他说，\Quote{你为什么因知识而自夸呢？如果你没有爱德，你甚至会因此受损。因为有什么比自夸更糟呢？但如果再加上另一者，前者也将安全。因为即使你比你的邻人知道得更多，如果你爱他，你就不会自高，反而会引导他也达到同样的境界。}因此，在说了\BibleQuote{知识使人自高自大}之后，他又加上了\BibleQuote{然而爱德却立人}。他没有说\Quote{行为谦逊}，而是说了更重大、更有益处的话。因为他们的知识不仅使他们自大，还使他们分心。为此，他将二者对立起来。\par

3. 然后他加上第三个考量，这足以使他们谦抑。这考量是什么呢？就是：即使爱德与之结合，我们的知识在那时也不完美。因此他加上：\par

% structure: p0013 source=h2 target=subsection reason=relative-heading-rank; base=h2

\subsection{\BibleRef{格前 8:2}}

\BibleQuote{但如果有人自以为知道什么，他还不知道他应当知道的事。}这是一记致命打击。\Quote{我不纠缠于知识是众人皆有的这一点，}他说，\Quote{我不说因你憎恨邻人、因傲慢而最伤害自己。但即使你独自拥有它，即使你谦逊，即使你爱弟兄，在这种情况下，你在知识上仍不完全。\InnerQuote{因为你们至今还不知道你们应当知道的事。}}如果我们至今对任何事尚无确切知识，怎么有些人竟疯狂到说他们确切地认识天主呢？而虽然我们对其他一切事有确切知识，即便如此，我们也不可能拥有这种程度的认识。因为天主超越万物的距离，是无法言说的。\par

请注意他是如何压下他们膨胀的骄傲：因为他没有说\Quote{对于面前的事，你们没有恰当的知识}，而是说\Quote{关于一切}。他也没有说\Quote{你们}，而是说\Quote{无论什么人}，无论是\Person{伯多禄}、\Person{保禄}，或任何其他人。因为借此他既安抚了他们，又小心地使他们谦抑。\par

% structure: p0016 source=h2 target=subsection reason=relative-heading-rank; base=h2

\subsection{\BibleRef{格前 8:3}}

\BibleQuote{但如果有人爱慕天主，那同一人，}他没有说\Quote{认识祂}，而是说\Quote{被祂所认识}。因为我们没有认识祂，而是祂认识了我们。因此基督说：\BibleQuote{不是你们拣选了我，而是我拣选了你们。}保禄在别处说：\BibleQuote{到那时我将完全认识，如同我也被完全认识一样。}\par

现在，我恳请你们注意，他是通过什么方式使他们的高傲降下。首先，他指出他们所知道的事并非唯独他们知道；因为他说：\BibleQuote{我们众人都有知识。}其次，这事物本身在没有爱德的情况下是有害的；因为他说：\BibleQuote{知识使人自高自大。}第三，即使与爱德结合，它也不完全、不完美。他说：\BibleQuote{因为如果有人自以为知道什么，他还不知道他应当知道的事。}此外，他们连这一点也不是出于自己，而是来自天主的恩赐。因为他没有说\Quote{认识天主}，而是说\Quote{被祂所认识}。再者，正是这件事来自于爱德，而他们所拥有的爱德并不相称。因为他说：\BibleQuote{如果有人爱慕天主，那同一人就被祂所认识。}如此大大平息了他们的恼怒后，他开始以教义的方式讲话，如此说：\par

% structure: p0019 source=h2 target=subsection reason=relative-heading-rank; base=h2

\subsection{\BibleRef{格前 8:4}}

4. \BibleQuote{因此，关于吃祭偶像之物，我们知道：世上没有什么偶像，也只有一个天主。}看他陷入了怎样的困境！因为他既要证明：应当戒避这类宴席，又要证明它无力伤害那些参与的人：这两件事彼此不太一致。因为当他们被告知这些东西对他们无害时，他们自然会将它们当作无关紧要的事物去参与。但当被禁止接触时，他们会反过来怀疑禁令是因为它们能造成伤害。因此，你看，他首先放下他们对偶像的看法，然后陈述他们应当戒避的首要理由：他们在弟兄们面前放置的绊脚石；他说：\BibleQuote{关于吃祭偶像之物，我们知道：世上没有什么偶像。}他又使这事成为公共财产，不允许唯他们独有，而是将知识扩展到整个世界。因为他说：\Quote{不但在你们中间，而且这道理在世上任何地方都通行。}那么是什么呢？\BibleQuote{世上没有什么偶像，只有一个天主。}那么怎样？难道没有偶像吗？没有雕像吗？确实有，但它们没有能力：它们不是神，而是石头和魔鬼。因为他现在正反对两派人：一派是其中较粗俗的人，另一派是那些被视为爱好智慧的人。因此，前者只知道不过是石头，后者则断言某些能力寓于其中，他们也称之为神；因此他对前者说：\BibleQuote{世上没有什么偶像}，对后者说：\BibleQuote{只有一个天主}。\par

你是否注意到他写这些事，并非单纯地立下教义，而是为了反对那些外人？这件事我们确实必须时时密切观察：他说话是抽象地，还是针对某些人？因为这对我们教义观点的准确性和对他表达的精确理解大有裨益。\par

% structure: p0022 source=h2 target=subsection reason=relative-heading-rank; base=h2

\subsection{\BibleRef{格前 8:5-6}}

5. \Quote{因为虽有称为神的，或在天上，或在地上，就如那许多的神，许多的主；可是我们只有一位天主，就是圣父，万物都出于祂，我们也归于祂；也只有一位主耶稣基督，万物都借着祂，我们也借着祂。}既然他曾说，\Quote{偶像算不得什么}且\Quote{没有别的神；}可是仍有偶像，仍有那些被称为神的；为免他似乎与明显的事实相矛盾，他便接着说，\Quote{因为虽有称为神的，的确也有；}并非绝对地\Quote{有；}而是\Quote{称为，}并非实际拥有此名，而是名义上：\Quote{或在天上，或在地上：——在天上，}指太阳、月亮和其余众星；因为这些也是希腊人所崇拜的；在地上则有魔鬼和所有那些从人中被奉为神的。——\Quote{可是我们只有一位天主，圣父。}最初他表达时没有用\Quote{圣父}一词，只说\Quote{除了一位天主外，没有别的神，}现在他加上这一句，当他完全摒弃了其他神之时。\par

接着，他引证那确实是神性的最大标志：\Quote{万物都出于祂。}因为这也暗示那些别的不是神。因为经上说\BibleRef{耶 10:11}：\Quote{那没有创造天地的神，必从地上灭亡。}然后他又加上同样重要的一句：\Quote{我们也归于祂。}因为当他说\Quote{万物都出于祂}时，他指的是受造和从无中产生万有；但他说\Quote{我们也归于祂}时，他是指信仰之言和互相的\OriginalTerm{私有关系}{\GreekText{οἰκειώσεως}}，正如他以前在\BibleRef{格前 1:30}说的：\Quote{但你们在基督耶稣内也是出于祂。}我们在两方面出于祂：因从无中被造，又因成为信徒。因为这也是一个受造：这也是他在别处所宣告的：\BibleRef{弗 2:15}\Quote{为把双方在自己内造成一个新人。}\par

\Quote{并且只有一位主，耶稣基督，万物都借着祂，我们也借着祂。}关于基督，我们又当同样理解。因为借着祂，人类既从无中被造而存在，又从错误返回真理。因此，对于\Quote{出于祂}这一说法，不可离了基督来理解。因为我们是出于祂，借着基督被造的。\par

6. 然而，你若注意，他并没有将名称分配为专属，将主之名归于圣子，将天主之名归于圣父。因为圣经也常互换它们；正如经上说\BibleRef{咏 109:1}：\Quote{上主对我主说；}又\BibleRef{咏 64:8}：\Quote{因此天主，你的天主，给你傅油；}以及\BibleRef{罗 9:5}：\Quote{按肉身说，基督也是出于他们，祂是超越万物的天主。}在许多实例中，你可看到这些名称互换位置。此外，如果它们被分别归于每个本性，并且如果圣子不是天主，不像圣父一样是天主，却仍是一个儿子：那么，在说了\Quote{但对我们来说，只有一位天主}之后，再加上\Quote{圣父}一词以表明那无始者，就是多余的。因为天主的话本身就足以说明这一点，如果这话是指祂独一的话。\par

这还不止，还有另一项评论：如果你说，\Quote{因为经上说\InnerQuote{一位天主}，所以天主一词不应用于圣子；}请注意，同样的道理也适用于圣子。因为圣子也被称为\Quote{一位主}，然而我们并不因此主张主一词只限于祂。所以，\Quote{一}这一表述应用于圣子时具有的力量，应用于圣父时也同样具有。并且，正如圣父并不因圣子被称为一位主就被排除为与圣子相同意义上的主；同样，圣子也不因圣父被称为一位天主就被排除为与圣父相同意义上的天主。\par

7. 如果有人问：\Quote{为何他未提及圣神？}我们可这样回答：他的论敌是崇拜偶像者，争论焦点在于\Quote{诸神和诸主}。因此，他称圣父为天主，称圣子为主。若他不敢与圣子同称圣父为主，以免被人怀疑谈论两位主；也不敢与圣父同称圣子为天主，以免被人认为谈论两位天主：那么，他未提及圣神又有什么可惊奇的呢？目前他的争辩对象是外邦人，目的在于表明在我们中间并无多位神祇。因此他紧抓\Quote{一}这个词，说：\Quote{除了一位天主，没有别的神；我们只有一位天主，一位主。}由此可见，他迁就听者的软弱而采用这种表达方式，也因此完全没有提及圣神。若不是这样，他就不应在别处提及圣神，也不应将祂与圣父和圣子并列。因为，如果祂被排除于圣父和圣子之外，那么在圣洗圣事中更不应将祂置于同等地位；而正是在圣洗圣事中，天主性的尊荣尤为彰显，且施予唯独天主才能赐予的恩赐。如此，我指出了祂在这段经文中被沉默的原因。如果你认为这不是真正的原因，请告诉我，为何在圣洗中祂与圣父、圣子同列？你只能给出一个原因，即祂享有同等的尊荣。无论如何，当他没有这种限制时，便将圣神置于同等地位，如此写道：\BibleRef{格后 13:14}\Quote{愿主耶稣基督的恩宠，和天主的爱，及圣神的相通，与你们众人同在。}又说：\BibleRef{格后 12:4}\Quote{恩赐虽有区别，却是同一的圣神；职务虽有区别，却是同一的主；功效虽有区别，却是同一的天主。}但因现在他的对话者是希腊人及希腊人归化者中较软弱的，所以他这样\OriginalTerm{谨慎处理}{\GreekText{ταμιεύεται}}。这正是先知们在论及圣子时所行的：因听者的软弱，从未明确提及祂。\par

% structure: p0029 source=h2 target=subsection reason=relative-heading-rank; base=h2

\subsection{格前 8:7}

他说：\Quote{但所有人没有这种知识。}这是指什么知识？关于天主的，还是关于祭过偶像之物的？因为他此处要么暗指那些说有多神多主、不认识真天主的希腊人，要么暗指那些仍较为软弱的希腊归化者，他们尚未清楚知道不应害怕偶像，且\Quote{世上偶像算不得什么}。但他说这话，是在温和地安抚和鼓励后者。因为他无需一一提及他所要责斥的一切，尤其是他打算以后用更严厉的态度再次访问他们。\par

8. \Quote{但有些人因惯于偶像，吃祭物时仍视为祭过偶像之物，他们的良心既然软弱，就受到了玷污。}他说：他们仍在偶像前战栗。不要对我说现在的安定，以及你们从祖先接受了真宗教。要回溯到那时代，想想福音刚开始，不敬神仍猖獗，祭坛燃烧，祭品和奠酒献上，且大部分人都是外邦人；想想那些从祖先接受了不敬神的人，他们的父亲、祖父、曾祖父都如此，他们曾受魔鬼的极大痛苦。他们突然改变后，内心会怎样！面对魔鬼的攻击，他们会怎样战栗！为了他们，他也稍加保留，说：\Quote{但有些人因祭过偶像之物的良心。}他既未公开揭露他们，也未严厉打击他们；也未完全忽略他们，而是模糊地提及他们，说：\Quote{现在有些人因偶像的良心，仍视之为祭过偶像之物而吃}；即带着与先前相同的想法；\Quote{他们的良心既然软弱，就受到了玷污}；尚未能鄙视并彻底嘲笑它们，仍有些疑虑。好比有人按犹太习俗认为摸死尸会玷污自己，然后看见别人以清白的良心触摸，自己却未以同样心意去摸，就会受到玷污。这就是他们当时的心态。他说：\Quote{有些人因偶像的良心，仍这样做直到如今。}他不是无故加上\Quote{直到如今}，而是表明他们因不肯迁就，并未取得进展。因为这不是带他们归向的方法，而是要通过其他方式，用言语和教导去说服他们。\par

\Quote{他们的良心既然软弱，就受到了玷污。}他至今未从事物本身的性质论述，而是转而围绕着领受者的良心。因为他担心，在纠正软弱者的同时，会给强壮者造成重击，使他也变得软弱。因此他同样顾惜两者。他也不认为事情本身有何重大意义，而是使他的论证十分充分，以防止任何这类猜疑。\par

% structure: p0033 source=h2 target=subsection reason=relative-heading-rank; base=h2

\subsection{格前 8:8}

9. \Quote{但食物不能使我们得到天主的赞许。因为我们吃也无益，不吃也无损。}你看他如何再次打压他们的高傲？他说，不但他们，我们众人都有知识，而且没有人按应当知道的知道，知识使人自大，然后他安抚他们，说这知识并非人人都有，软弱是这些人被玷污的原因，免得他们说：这于我们何干？为什么那人没有知识？为什么他软弱？——我说，免得他们这样反驳，他没有立刻明确指出因怕伤害别人而应当禁戒，而是先略作试探，然后指出比这更重要的事。这是什么？就是即使没有人受伤，也没有导致别人的错谬，也不应当这样做。因为前一个论点本身是徒劳的。因为听到别人受伤害而自己有益的人，不太容易禁戒；只有当他发现自己毫无得益时才会这样做。所以他把这一点先写下，说：\Quote{但食物不能使我们得到天主的赞许。}看他把那被认为出自完全知识的事看得多么轻！\Quote{因为我们吃也无益}（即在天主眼中并不更高，好像做了什么好事或大事），\Quote{不吃也无损}，即丝毫不亚于别人。至此他已表明这事本身是多余的，如同无物。因为做了无益、不做无损的事，必是多余的。\par

10. 但他在下文揭示了此事可能引起的一切害处。目前，他所谈论的是弟兄们所遭遇的事。\par

% structure: p0036 source=h2 target=subsection reason=relative-heading-rank; base=h2

\subsection{\BibleRef{格前 8:9}}

\Quote{你们要谨慎，}他说，\Quote{恐怕你们这自由竟成了那软弱弟兄的绊脚石。}（\OriginalTerm{弟兄们}{\GreekText{τῶν} \GreekText{ἀδελφῶν}} 不在公认文本中。）\par

他没有说：\Quote{你们的自由成了绊脚石，}也没有肯定地说，免得使他们更加无耻；而是用：\Quote{你们要谨慎，}来恐吓他们，使他们羞愧，引导他们否定这样的行为。他也没有说：\Quote{你们这知识，}那样听起来更像赞美；也没有说：\Quote{你们这完全，}而是说：\Quote{你们的自由，}这似乎更带有鲁莽、固执和傲慢的意味。他也没有说：\Quote{对弟兄们，}而是说：\Quote{对那些软弱的弟兄，}以他们连软弱的人、而且又是他们的弟兄都不放过，来加重对他的指控。就算你不纠正他们，不激励他们，难道还要绊倒他们，使他们跌倒吗？你本应伸出援手。既然你不愿意，那至少不要推倒他们。因为如果一个人是邪恶的，他需要惩罚；如果是软弱的，他需要医治；但现在他不仅是软弱的，而且是弟兄。\par

% structure: p0039 source=h2 target=subsection reason=relative-heading-rank; base=h2

\subsection{\BibleRef{格前 8:10}}

\Quote{若有人看见你这有知识的，在偶像的庙里坐席，这人既是软弱的，他的良心岂不放胆去吃那祭偶像之物吗？}\par

在说了\Quote{你们要谨慎，恐怕你们这自由成了绊脚石}之后，他解释它是如何以及以何种方式成为绊脚石的。他不断使用\Quote{软弱}一词，免得人以为祸患出于事物的本性，或以为魔鬼可怕。他说：\Quote{此刻，一个人正要完全远离一切偶像；但当他看见你喜爱在偶像庙里流连，他就以此为示范，自己也留在那里。因此，不仅他的软弱，而且你不合时宜的行为，都助长了针对他的阴谋；因为是你使他变得更软弱。}\par

% structure: p0042 source=h2 target=subsection reason=relative-heading-rank; base=h2

\subsection{\BibleRef{格前 8:11}}

\Quote{因此，基督为他死的那软弱弟兄，也就因你的食物而丧亡。}\par

因为在这恶事上，有两件事使你无可推诿：一是他是软弱的，二是他是你的弟兄；其实还有第三件，比这一切更可怕。是什么？就是基督甚至不拒绝为他而死，而你不能容忍自己迁就他。你看，他也提醒那完全的人，他从前是什么，以及基督为他而死。他没有说：\Quote{你应当为他而死，}也没有说：\Quote{为你的完全，}也没有说：\Quote{为你的知识，}而是说：\Quote{为你的食物。}因此，指控有四条，且极其沉重：他是弟兄，他是软弱的，基督如此看重他甚至为他而死，而这一切之后，他竟为了\Quote{为了一口食物}而丧亡。\par

% structure: p0045 source=h2 target=subsection reason=relative-heading-rank; base=h2

\subsection{\BibleRef{格前 8:12}}

\Quote{你们这样得罪弟兄们，伤了他们软弱的良心，就是得罪基督。}\par

你看到他是如何平缓而逐步地将他们的过犯带到极恶的顶峰？他又提到另一类人的软弱，于是这些自认为对他们有利的事，他处处转过来压在他们自己头上。他没有说：\Quote{把绊脚石放在他们路上，}而是说：\Quote{伤害，}以他言辞的力量表明他们的残忍。因为有什么比伤害病人更野蛮的呢？然而，没有一种创伤比使人跌倒更严重。事实上，这也常常是死亡的原因。\par

但他们如何\Quote{得罪基督}？一方面，因为祂视祂仆人的事如己事；另一方面，因为受伤的人构成了祂的身体，而他们是祂的一部分；第三方面，因为祂用自己的血所建造的工程，这些人为了自己的野心而在摧毁。\par

% structure: p0049 source=h2 target=subsection reason=relative-heading-rank; base=h2

\subsection{\BibleRef{格前 8:13}}

11. \BibleQuote{因此，如果食物使我的弟兄跌倒，我就永远不吃肉。}这正如最好的教师，以身作则来教导他所言之事。他并没有说是合理还是不合理；但无论如何。\Quote{我不是说，}（这是他说话的语气）\Quote{祭偶像的肉，这已经因另一个理由被禁止；但如果即使是许可的和允许的事物中，有任何导致跌倒的，连这些我也要戒绝：不是一两天，而是我一生之久。}因为他说：\BibleQuote{我永远不吃肉。}他又不是说：\Quote{免得我毁灭我的弟兄，}而只是说：\Quote{免得我使我的弟兄跌倒。}因为这实在是极端的愚昧：基督所极其关心的事物，甚至祂选择了为他们而死，而我们将它们视为完全不屑一顾，以致连为了他们而戒绝肉食都不愿意。\par

这些话不仅当时对他们说，也对我们是合宜的，因为我们是这样容易轻视邻人的救恩，并说出那些撒殚般的话。我说撒殚般的：因为这样的说法，\Quote{我何在乎，虽然这个人跌倒，那个人丧亡？}散发出他的残忍和没有人性的思想。然而在那种情况下，被得罪者的软弱也影响了结果；但在我们的情况下则不然，我们犯罪的方式甚至足以使强壮的人跌倒。因为当我们打击、抢夺、欺诈、把自由人当作奴隶使用时，这难道不足以使任何人跌倒吗？不要对我说这个人是个鞋匠，另一个人是个染工，再一个人是个铜匠；但请记住，他是一位信徒，是一位弟兄。这些人是我们的门徒来源：渔夫、税吏、造帐篷的，属于那位在木匠家长大、屈尊以木匠的未婚妻为母亲、在襁褓之后被放在马槽里、没有枕头的地方、旅途如此漫长以至于奔波足以使祂疲惫、被人供养的主。\par

12. 要思想这些事，视人的骄傲为无有。但要把造帐篷的与你那乘战车、有无数的仆从、在市场上趾高气扬的人一样当作兄弟：不，宁可是前者而非后者；因为兄弟一词更自然地用于有更大相似之处的人。那么，谁更像渔夫呢？是那个靠每日劳作维生、既无仆从也无居所、完全被匮乏所包围的人；还是那个被如此巨大的排场所包围、并违反天主法律的人？那么，不要轻视那两人中更像是你兄弟的人，因为他更接近宗徒的样式。\par

\Quote{然而，}你说，\Quote{并非出于自愿，而是出于被迫；因为他不是心甘情愿做这事。}这从何说起？你难道没有听过\BibleQuote{你们不要判断，免得你们受判断}？但是，为说服自己他不是违背意愿，你走上前，给他一万塔冷通的金子，你必看见他推开它。这样，即使他没有从祖先那里继承任何财富，但当他有能力取用时，他却不让它接近自己，也不增加自己的财物，就显示了他对财富的极大蔑视。因为若望就是载伯德的儿子，那个极其贫穷的人：但我想我们因此不能说他贫穷是迫于无奈。\par

所以，每当你看见有人钉钉子、用铁锤敲打、满身煤灰，不要因此轻视他，反而要因此而钦佩他。因为连伯多禄也曾束腰，收拾拖网，在主的复活后去打鱼。\par

我为何说伯多禄？因为这位保禄自己，在不断的奔波和那些伟大的奇迹之后，站在造帐篷的店铺里，缝制皮革：当时天使尊敬他，恶魔战栗。他不羞于说，\BibleRef{宗 20:34} \BibleQuote{这双手供应了我的需要，以及同我在一起的人的需要。}我说什么？他不羞于？是的，他反而以此夸耀。\par

但你会说：\Quote{现今有谁可与保禄的德行相比？}我也知道没有，但因此不应轻视今世活着的人：因为如果你为基督的缘故给予尊荣，即使他是最末等的，但只要他是信徒，就理应受尊荣。设想一位将军和一名普通士兵都来到你面前，他们都是国王的朋友，你向两人敞开家门，你在哪个人的身上显得最尊敬国王？显然是在士兵身上。因为在将军身上，除了对国王的忠诚外，还有许多别的足以赢得你如此尊重的原因；但士兵除了对国王的忠诚外，别无他物。\par

因此，天主命令我们请瘸腿的、残疾的、无力回报我们的人来参加我们的晚餐和宴席；因为这些才是真正为天主而行善事。然而，如果你款待某个伟大显赫的人，你所做的就不是如此纯粹的慈悲，而常常有一部分归于你自己，既出于虚荣，也出于回报的期待，还由于因你的客人而使你在许多人眼中的地位提升。无论如何，我想我能指出很多人，怀此目的巴结圣人中更显赫者，即借此得以与统治者更亲近，并使他们此后在自身事务和家庭中更为有用。他们向那些圣人要求许多这样的回报作为酬谢；这种行为破坏了款待的回报，因为他们心怀这种意图。\par

我何必提及圣人的事呢？因为凡在今世寻求甚至从天主获得其劳苦的赏报，并为现世的利益而追随德行的人，必定会减少他的报酬。但那完全在那里要求他全部荣冠的人，更令人钦佩：像那位拉匝禄，他现在就在那里\Quote{领受}\BibleRef{路 16:25}他一切的\Quote{他的福乐}；像那三位青年，当他们即将被投入火窑时，说：\BibleRef{达 17:17}\Quote{有一位天主在天上，能拯救我们；即使不，大王，愿你知道，我们不事奉你的神，也不朝拜你所立的金像。}像亚巴郎，他甚至献上自己的儿子并杀了他；他这样做，不是为任何赏报，而是认为这一件事就是最大的报酬：服从主。\par

让我们也效法这些人。因为这样，我们必因我们一切善行而得加倍的回报，因为我们怀有这样的心思行一切事：这样我们也必获得更明亮的荣冠。愿天主恩赐我们众人获得这些荣冠，借着我们的主耶稣基督的恩宠和慈爱，愿光荣、权能、尊崇归于祂，及圣父与圣神，从今直到永远。阿们。\par

\SourceRecord{Translated by Talbot W. Chambers.；Nicene and Post-Nicene Fathers, First Series；Vol. 12.；Edited by Philip Schaff.；Buffalo, NY: Christian Literature Publishing Co.,；1889.；<http://www.newadvent.org/fathers/220120.htm>.}

% structure: p0001 source=h1 target=section reason=page-title

\section{\WorkTitle{格林多前书讲道 第廿一篇}}

% structure: p0002 source=h2 target=subsection reason=relative-heading-rank; base=h2

\subsection{\BibleRef{格前 9:1}}

\Quote{我岂不是宗徒么？我岂不是自由么？我岂不是见过我们的主耶稣基督么？你们岂不是我在主内的工作么？}\par

正因为他说过：\Quote{若食物使我的弟兄跌倒，我就永远不吃肉；}这是一件他尚未做到的事，却宣称若有必要他会做。免得有人会说：\Quote{你妄自夸口，说话严厉，满口承诺，这事对我来说或对任何人都容易；但若这些言语出自信，就请通过行为证明你为了不使弟兄跌倒而舍弃了某事：}因此，在接下来的内容中，他不得不也对此加以证明，指出他如何常常舍弃甚至合法的事，以免冒犯他人，尽管没有法律强迫他这样做。\par

我们尚未触及事情的可钦赞之处：尽管他为了避免冒犯而不仅从合法的事上自制，这已是可钦赞的；但他这样做所付出的辛劳和危险更甚。\Quote{何必，}他说，\Quote{谈那些祭偶像之物呢？因为虽然基督命令那些传福音的人应靠门徒的供给生活，我却没有这样做，反而选择了——若有必要——以饥饿结束生命，忍受最痛苦的死亡，以便避免从我所教导的人那里接受供给。}\par

并非因为他们会因此跌倒，而是因为他不接受供给会建立他们：这是他做的一件更伟大的事。为此他召唤他们自己作证，他在他们中间常常劳苦、饥饿，靠别人供养，忍受窘迫，以免冒犯他们。但他们并没有理由冒犯，因为这不过是他所遵守的法律罢了。尽管如此，他仍以某种超额善功的方式宽容他们。\par

既然他做了超出法律要求的事，以免他们冒犯，并从合法的事上自制以建立他人；那么，那些不从祭偶像之物上自制的人该当何罪？尤其当许多人因此丧亡时？这是一件即便没有冒犯的事也应当远离的，因为是\Quote{魔鬼的筵席}。\par

因此，整个话题的要点就是他在许多经文中阐述的内容。但我们必须重新审视，并重新进入他所提出的话题。因为他既没有像我这样明确地写下来，也没有立刻切入主题；而是从另一个话题开始，这样说：\par

2. \Quote{我岂不是宗徒么？}因为除了已经说过的一切，这一点也带来不小的差异：即保禄本人就是如此行事的人。例如：为了防止他们声称：\Quote{你可以尝祭物，同时封印；}他暂时不反驳那说法，而是论证：\Quote{即使那是合法的，你弟兄的损伤也应阻止你这样做；}随后他又证明那甚至是不合法的。不过在这一点上，他通过与自己相关的情况确立了前者。而当他即将说他没有从他们那里接受任何东西时，并非立即摆出来，而是首先肯定自己的尊荣：\Quote{我岂不是宗徒么？我岂不是自由么？}\par

因此，为了阻止他们说：\Quote{真的，你没有接受，但你没有接受的原因是因为那是不合法的；}于是他首先列出他本可以合理接受的原因，如果他愿意的话。\par

此外，为了不使他在说这些话时显得对伯多禄及类似伯多禄的人有所嫉妒（因为他们并不拒绝接受），他首先表明他们有权接受，然后为了不让任何人说：\Quote{伯多禄有权接受，但你却没有，}他先用这些对自己的称颂来预告听者。他意识到必须赞美自己（因为那是纠正格林多人的方法），却又厌恶言及自己的伟大之处，请看他是如何根据时宜调和这两种情感的：他限制自己的颂词，不是根据他对自己所知，而是根据主题的必要性。他本可以说：\Quote{我比任何人更有权接受，甚至超过他们，因为\InnerQuote{我比他们更劳碌}。}但他省略了这一点，因为那是他超越他们的方面；他只列出那些使他们伟大的、他们接受供给的正当理由，如下：\par

\Quote{我岂不是宗徒么？我岂不是自由么？}即\Quote{我岂没有管理自己的权柄？我岂在任何人之下，可以制止我、禁止我接受么？}\par

\Quote{但他们比你有优势，因为他们曾与基督同在。}\par

\Quote{不，这一点也没有被拒绝给我。}为此他说：\par

\Quote{末了，} \BibleRef{格前 15:8} 他说，\Quote{也显现给我这个像流产儿一样的人。} 这一点也是不小的尊荣：因为\Quote{许多先知，} \BibleRef{玛 13:17} 祂说，\Quote{和义人曾渴望看见你们所看见的，却没有看见；}并且\Quote{日子将到，你们渴望看见这日子中的一天。} \BibleRef{路 17:22}\par

\Quote{那么，即使你做了\InnerQuote{宗徒}，\InnerQuote{自由}，并且\InnerQuote{看见了基督}，若你没有显露出任何宗徒的工作，你怎么能有权接受呢？}因此在此之后他补充说：\par

\Quote{你们岂不是我在主内的工作么？}因为这是大事；没有它，其他一切全无益处。连犹达斯自己也是\Quote{宗徒}，\Quote{自由}，并\Quote{看见了基督}；但因他没有\Quote{宗徒的工作}，这一切对他都无益。你看他为什么也加上这一点，并召唤他们自己作证。\par

此外，因为他所说的是大事，请看他是如何节制它的，加上\Quote{在主内：}即\Quote{这工作是天主的，不是我的。}\par

% structure: p0019 source=h2 target=subsection reason=relative-heading-rank; base=h2

\subsection{\BibleRef{格前 9:2}}

\Quote{如果我为别人不是宗徒，至少我为你们是。}\par

你看到他何等无必要地在这里展开？然而他本可论及全世界、蛮族、海洋和陆地。但他完全不提这些，反而以让步的方式证明他的论点，甚至给予的比他需要的更多。好像他说：\Quote{为何我要详论额外的事，既然这些本身就足以达成我当前的目的？你会注意到，我所说的不是在别处的成就，而是你们作为见证的那些事。由此得出的结论是，即使从别处得不到，但从你们那里我有权利接受。然而，从那些我最有权接受的人——你们，我的学生——我却没有接受。}\par

\Quote{如果对别人我不是宗徒，至少对你们我是宗徒。}再次，他以让步的方式陈述他的论点。因为全世界都视他为宗徒。\Quote{不过，}他说，\Quote{我不说那个，我不争辩，但把关于你们的事讲明。\InnerQuote{因为你们是我宗徒职分的印证。}}即其证明。\Quote{此外，若有人想知道我从哪里成为宗徒，你们就是我提出的人证：因为我在你们中显示了宗徒的一切标记，没有一样缺失。}正如他在第二封书信中所说，\BibleRef{格后 12:12}\Quote{我虽是虚无，却在你们中，以各种坚忍，藉着征兆、奇迹和异能，真正显示了宗徒的标记。因为在什么事上，你们不如其他教会呢？}所以他说，\Quote{你们是我宗徒职分的印证。}\Quote{因为我既行了奇迹，又以言语教导，经历了危险，并显示了无可指责的生活。}这些主题你们可以从这两封书信中看到充分阐述，他如何精确地向他们展示每一点的证明。\par

% structure: p0023 source=h2 target=subsection reason=relative-heading-rank; base=h2

\subsection{格林多前书 9:3}

3. \Quote{我的辩白对查问我的人就是这。}什么是\Quote{我的辩白对查问我的人就是这？} \Quote{对那些想知道凭什么证明我是宗徒的人，或指责我收钱的人，或询问我为何不收钱的人，或想要表明我不是宗徒的人：对所有这些人，我给你们的教导以及我即将说的话，可以构成一个充分的说明和辩白。}那么这些是什么呢？\par

% structure: p0025 source=h2 target=subsection reason=relative-heading-rank; base=h2

\subsection{格林多前书 9:4-5}

\Quote{难道我们没有权利吃喝吗？难道我们没有权利带一位女信友到处旅行吗？}为什么，这些话怎么是辩白？\Quote{因为，当它显得我连允许的事都节制时，就不该怀疑我是骗子或为利益而行动。}\par

因此，从先前所举的和我曾教导你们的事，以及我刚刚所说的，我有足够的材料对你们进行辩白：所有查问我的人，我都以此为依据，既提出前面的话，也提出以下的话：难道我们没有权利吃喝吗？难道我们没有权利带一位女信友到处旅行吗？\Quote{然而，尽管有权利，我却节制不用？}\par

那么怎样？他平时不吃不喝吗？最真实地说，在许多地方他确实不吃不喝：因为\BibleRef{格前 4:11}\Quote{我们正挨饿、受渴、赤身露体。}然而，这里不是他的意思；而是什么？\Quote{我们不吃不喝，不向我们教导的人收取费用，尽管我们有权利收取。}\par

\Quote{难道我们没有权利带一位女信友到处旅行，如同其余的宗徒、主的弟兄和刻法一样吗？}请观察他的巧妙。领唱者被放在排列的最后：因为那是提出最强论据的时机。能在其他人中指出这种行为的榜样并不那么了不起，而在最重要的斗士和受托管理天钥匙的人身上指出，才是了不起的。但他不只是提到伯多禄一人，而是他们全体：好像他说，无论你们寻找较低级的还是更杰出的，在所有人中你们都会找到这种榜样。\par

因为主的弟兄，从他们最初的不信\BibleRef{若 7:5}中释放出来，也已成为那些受称赞的人，尽管他们未达到宗徒的地位。相应地，他给他们安排了中间的位置，把极端的人放在前后。\par

% structure: p0031 source=h2 target=subsection reason=relative-heading-rank; base=h2

\subsection{格林多前书 9:6}

\Quote{难道只有我和巴尔纳伯没有权利不做工吗？}\par

（请看他的谦逊心怀和他那无嫉妒的灵魂，他如何注意不隐瞒他所知道与他共享这一完美的人。）因为如果其他事是共同的，怎能不是共同的？他们和我们都是宗徒，都是自由的，都见过基督，都行了宗徒的工作。因此，我们同样有权利不工作而过活，并由我们的门徒供养。\par

% structure: p0034 source=h2 target=subsection reason=relative-heading-rank; base=h2

\subsection{格林多前书 9:7}

4. \Quote{有谁当兵而自备粮饷呢？}因为，他既然从宗徒证明了这样做是合法的，即最有力的一点，接着他就来到例子和普通做法；正如他惯常做的：\Quote{有谁当兵而自备粮饷呢？}他说。但请你思考，我祈祷，这些例子与他的主题多么吻合，他又如何首先提到伴随着危险的那个；即兵役、武器和战争。因为宗徒之职就是这样的事，甚至比这些更危险。因为他们的战斗不仅与人，也与魔鬼，并反抗那些存在者的首领。因此他所说的是：\Quote{连异教统治者，残忍和不义，都不要求他们的士兵忍受劳苦和危险而自谋生计。基督怎么会要求这个呢？}\par

他不仅满足于一个例子。因为对于较简单和迟钝的人，看到平常的习惯也随从天主的法律，往往是一种大安慰。所以他又转向另一个话题，说道：\Quote{谁种葡萄园，却不吃它的果实呢？}因为正如前者表明了他的危险，后者表明了他的劳苦、大量劳作和关怀。\par

他同样添加了第三个例子，说：\Quote{谁牧养羊群，而不吃羊的奶呢？}他正在展示教师对他管辖之下的那些人应当表示的巨大关怀。事实上，宗徒们既是士兵，又是农夫和牧人，不是在地上、也不是对无理性的动物、也不是在感官可察觉的战争中，而是在理性的灵魂中并与魔鬼作战。\par

也必须注意到，他在各处都保持节制，只寻求有用的，而非非凡的。因为他说：\Quote{什么士兵服役而不致富？}而是：\Quote{什么士兵曾经自备粮饷服役？}也没有说：\Quote{谁栽种葡萄园，却不收取金子，或吝于收取全部果实？}而是：\Quote{谁不吃其中的果子？}也没有说：\Quote{谁牧养羊群，而不贩卖羊羔？}而是什么？\Quote{而不吃羊的奶？}不是羊羔，而是奶；表明教师应得的微薄补偿就够了，甚至仅是他必需的食物。（这是指那些会吞没一切并收取全部果实的人。）\Quote{主也是这样规定了}，说：\BibleQuote{工人配得他的食物。} \BibleRef{玛 10:10}\par

他不仅通过他的比喻确立了这一点，而且还展示了一位司铎应当是怎样的人。他应当拥有士兵的勇气、农夫的勤勉和牧人的细心，在这一切之后，只寻求必需品。\par

你们看，他既从宗徒们那里表明教师接受供应并非被禁止，也从日常生活的比喻中表明这一点，又继续到第三个主题，这样说：\par

% structure: p0041 source=h2 target=subsection reason=relative-heading-rank; base=h2

\subsection{格前 9:8}

\Quote{是我说这些合乎人的想法吗？或者法律不也说同样的话吗？}\par

因为他迄今为止尚未引用任何圣经，只提出了普通习惯；\Quote{不要以为}，他说，\Quote{我只依靠这些，也不是以人的意见作为这些规定的依据。因为我还能表明这些事也是天主所喜悦的，并且我读到一条古老的法律命令它们。}因此他也以提问的形式进行论述，这在完全承认的事情上常会这样；因此他说：\Quote{我说这些是合乎人的想法吗？}即\Quote{我只用人的例子来加强自己吗？}\Quote{或者法律不也说同样的话吗？}\par

% structure: p0044 source=h2 target=subsection reason=relative-heading-rank; base=h2

\subsection{格前 9:9-10}

\Quote{因为在梅瑟法律上记载说：\BibleQuote{牛在场上踹谷的时候，不可笼住它的嘴。}}\par

他为什么提到这个，有了司铎的例子？他想确立它远远超出所要求的范围。此外，免得有人说：\Quote{这关于牛的话与我们有什么关系？}他更精确地阐述说：\Quote{天主岂关心牛呢？}那么，天主，请告诉我，不关心牛吗？他确实关心它们，但并不至于为此制定一条法律。所以，如果他没有暗示某些重要的事，训练犹太人对待牲畜要有仁慈，并透过这些事也向他们谈论教师，他就不会如此关心，以至于甚至制定法律禁止笼住牛的嘴。\par

在此他也指出另一件事：教师的劳动既伟大也应当是伟大的。\par

还有另一件事。那么这是什么？旧约中关于关心牲畜的言论，其主要意义在于教导人类：事实上所有其他也是如此；例如关于各种衣服的诫命；关于葡萄园和种子，不使土地混杂各种作物；关于癞病；总之，所有其余：因为他们比较迟钝，他从这些题目与他们谈论，逐步提升他们。\par

并且看他在接下来的部分甚至没有确认它，认为它是清楚自明的。因为说了\Quote{天主岂关心牛呢？}他补充说：\Quote{还是全为我们说的？}甚至不是随意加上\Quote{全}字，而是免得给听众留下任何可以反驳的余地。\par

他详细阐述了比喻，说并宣告：\Quote{的确，这是为我们写的，因为耕种的当存着希望去耕种；}即教师应当享受他劳动的回报；\Quote{打场的也当存有分享的希望去打场。}注意他的智慧，他从种子转到打谷场；在此又一次表明教师的许多辛劳，他们亲身耕种和打场。关于耕种，因为没有什么可收成，只有劳动，他使用了\Quote{希望}一词；但关于打场，他立即允许果实，说：\Quote{打场的是他希望的分享者。}\par

此外，免得有人说：\Quote{这就是如此多辛劳的回报吗？}他补充说\Quote{在希望中}，即\Quote{未来的。}因此，这不笼嘴的动物的嘴所宣告的无非就是：劳动的教师也应当享受一些回报。\par

% structure: p0052 source=h2 target=subsection reason=relative-heading-rank; base=h2

\subsection{格前 9:11}

6. \BibleQuote{如果我们给你们播种了属神的事，我们从你们收获属血肉的事，还算大事吗？}\par

看，他又增加了第四个论据，说明供养的义务。因为他既已说了：\Quote{哪个士兵曾自备粮饷？}和\Quote{谁栽种葡萄园？}和\Quote{谁牧养羊群？}并引用了踹谷的牛；他同样指出另一个非常合理的理由，使他们可以正当地接受；即，他们已赐予更大的恩赐，不仅仅是劳苦的回报。那是什么？\Quote{我们若给你们播种了神性的事物，收割你们的世俗事物，还算大事吗？}你看见一个最公道的理由，比先前所有更充满理性吗？因为他说：\Quote{在那些例子中，种子是世俗的，果实也是世俗的；但这里却不是这样，种子是神性的，收成是世俗的。}如此，为阻止那些供养教师的人产生高傲思想，他表明他们收获得比给予的多。好像他曾说：\Quote{农夫，不管他们播种什么，也收获什么；但我们，在你们的灵魂中播种神性的事物，却收割世俗的事物。}因为他们所给的供养就是这样的。此外，更要使他们脸红。\par

% structure: p0055 source=h2 target=subsection reason=relative-heading-rank; base=h2

\subsection{格林多前书 9:12}

\Quote{如果别人分享对你们的这权利，我们岂不更该吗？}\par

请再看另一个论点，这也是来自例子，尽管不是同一类。因为他在这里提到的不是伯多禄，也不是宗徒，而是某些其他虚伪的人，他后来与他们争战，并且论到他们，他说：\BibleRef{格后 11:20}\BibleQuote{若有人吞吃你们，若有人俘虏你们，若有人自高自大，若有人打你们的脸，}并且他已经奏响了与他们对阵的前奏。因此，他也没有说：\Quote{如果别人从你们收取}，而是指出他们的傲慢、暴虐和贩卖，他说：\Quote{如果别人分享对你们的这权利}，即\Quote{统治你们，行使权威，把你们当作仆人使用，不仅俘虏你们，而且大有权威。}因此他加上：\Quote{我们岂不更该吗？}如果这段话是关于宗徒的，他就不会这么说。但显然他暗示的是某些有害之人，以及欺骗他们的人。\Quote{因此，除了梅瑟法律，连你们自己也制定了一条关于供养义务的法律。}\par

并且说了\Quote{我们岂不更该吗？}之后，他没有证明为什么更该，而是让他们的良心去说服他们这一点，同时希望既警告他们，也更彻底地使他们羞愧。\par

7. 然而，我们没有使用这权利；即\Quote{没有接受。}你们看见吗？当他用如此多的理由证明接受并非不合法之后，他接着说：\Quote{我们不接受，}免得他好像因禁止之事而节制？他说：\Quote{因为并非不合法，}他说，\Quote{我仍不接受；因为这是合法的，并且我们已用许多方式表明：从宗徒；从生活的事例，如士兵、农夫和牧人；从梅瑟法律；从事情的本质，即我们给你们播种了神性事物；从你们自己向别人所行的。}但是当他提出了这些，免得他好像羞辱那些习惯接受的宗徒，使他们羞愧并表明他不是因禁止之事而节制：同样，又免得他因大量的证据和例子以及推理（他藉此指出了接受的正当性），显得他本人渴望接受，因此说这些话；他现在纠正了这一点。后来他更清楚地说明，他说：\Quote{我写这些事，并不是为叫在我身上这样行}；但这里他的话是：\Quote{我们没有使用这权利。}\par

还有更重大的事，没有人能这样说：我们因富足而拒绝使用这权利；反而，当迫切的需要压迫我们时，我们也不屈服于需要。这事他在第二封书信中也说到：\BibleQuote{我抢夺了别的教会，取了他们的报酬，为能服事你们；当我和你们同在，并且匮乏时，我没有成为任何人的负担。}\BibleRef{格后 11:8}在这封书信中又说：\BibleQuote{我们又饥又渴，又赤身露体，又挨打。}\BibleRef{格前 4:11}这里他又暗示同一件事，说：\Quote{但我们忍受一切。}因为藉着说\Quote{我们忍受一切}，他暗示了饥饿、极大的困苦和所有其他事情。他说：\Quote{但即使这样，我们也没有被迫}破坏我们为自己制定的法则。为什么？\Quote{为叫我们对基督的福音不造成任何阻碍。}因为格林多人是心思软弱的人，他说：\Quote{免得我们因接受而伤害你们，我们选择做比命令更多的，而不妨碍福音}，即你们的教导。现在，如果我们在一个自由的事上，并且当我们既忍受许多艰苦，又有宗徒作我们的榜样时，我们节制以免造成阻碍（他没有说\Quote{颠覆}，而是\Quote{阻碍}；也不仅仅说\Quote{阻碍}，而是\Quote{任何}阻碍），好使我们，可以说，不至于引起圣言历程的丝毫延迟和停顿：\Quote{如果}他说，\Quote{我们用了如此大的谨慎，你们岂不更该节制？你们远不及宗徒，也没有任何法律提出来允许你们；反而，你们正着手做禁止的事，以及那些倾向于大损害福音的事，不仅仅是阻碍而已；而且甚至没有任何迫切的需要作为理由。}他这一切讨论都是因这些格林多人而起，他们因吃祭偶像之物而绊倒了软弱的弟兄。\par

8. 亲爱的，也让我们听这些事，好使我们不轻视那些跌倒的人，也不\Quote{阻碍基督的福音}；好使我们不背叛自己的救恩。当你的弟兄跌倒时，不要对我说：\Quote{这件事或那件事，使他跌倒的，并没有被禁止；它是许可的。}因为我有更重大的事要对你说：\Quote{即使基督自己许可了它，但如果你看见有人受伤，就要停止，不要使用那许可。}这也是保禄所做的；当他本可以接受时，基督已授予许可，他却没有接受。同样，我们的主以他的仁慈，在他的诫命中掺和了许多温和，使得它不全是出于命令，而是我们也可以出于自己的心意多做许多。因为如果他不愿意，他本可以扩展诫命，说：\Quote{不持续禁食的，要受责罚；不保持童贞的，要受惩罚；不放弃一切所有的，要受最严厉的刑罚。}但他没有这样做，而是给你机会，如果你愿意，可以主动多做。因此，当他谈论贞洁时，他说：\Quote{能领受的，就领受吧}；关于那位富人，他命令了一些事，但有些则留给他自己决定。因为他没有说：\Quote{变卖你所有的}，而是说：\Quote{你若愿意做成全的，就去变卖。}\par

但我们不仅不主动多做、超越诫命，反而在命令的事上远远不足。保禄忍受饥饿以免阻碍福音，而我们却不愿触动自己仓库里的东西，尽管看到无数灵魂跌倒。\Quote{是的}，有人说，\Quote{让蛀虫吃，不要让穷人吃；让虫子吞，不要让赤身者穿衣；让一切被时间朽坏，不要让基督得喂养，而当他饥饿时。}\Quote{咦，谁说了这话？}有人会问。然而，这正是可悲之处，因为这些话不是以言语，而是以行为说出；因为用言语说出比用行为做出更轻。那位残忍冷酷的暴君贪吝，岂不是一个喊着，对那些被她俘虏的人说：\Quote{你们的财物要摆在告密者、强盗和叛徒面前供他们奢靡，而不是摆在饥饿和穷乏者面前供他们度日？}难道不是你们制造了强盗？难道不是你们给嫉妒之火添柴？难道不是你们制造了流浪者和叛徒，把你们的财富放在他们面前作诱饵？这是何等的疯狂？（因为这确实是疯狂和明显的颠狂）将衣服装满你们的箱子，却忽视按天主的肖像和模样所造的人赤裸着、冷得发抖、艰难地站着。\par

\Quote{但他假装}，有人说，\Quote{颤抖和虚弱。}你难道不怕因这话而引发天主的雷霆落在你身上吗？（因为我充满愤怒，请容忍我。）我说你啊，养尊处优、饱食终日、饮酒至深夜、躺在柔软的被褥中，却认为自己不受审判，如此肆意妄为地使用天主的恩赐（因为酒不是为让我们醉酒，食物不是为让我们满足口腹，肉类不是为让我们撑大肚子）？而对于穷人、可怜人、近乎死亡的人，你却严格要求，难道你不怕基督那充满威严和恐怖的审判台吗？即使他是假装的，那也是出于需要和匮乏，因为你的残忍和无情，迫使他使用这样的面具，拒绝一切仁慈的倾向。有谁如此悲惨，在没有紧迫必要的情况下，为了一块面包而忍受这样的羞辱，哀哭自己，忍受如此严厉的惩罚？因此，他的这种假装四处传播，正是你无情的宣告。因为他虽恳求、哀求、说可怜话、哭泣、整天奔走，却得不到必需的食物，也许他设计了这种诡计，其耻辱和责备主要落在你身上，而非他自己身上：因为他值得怜悯，因为他陷入了如此大的必要；而我们却应受无数的惩罚，因为我们迫使穷人遭受这些事。如果我们轻易让步，他绝不会选择忍受这些。\par

我为何只提赤裸和颤抖？我还要说一件更令人战栗的事：有些人甚至被迫在幼年时就弄瞎自己孩子的眼睛，为了触动我们的麻木。因为当他们能看见并赤裸行走时，既不能因年龄也不能因不幸而赢得无情者的怜悯，他们便在如此巨大的灾祸上再加上另一出更严厉的悲剧，以期消除饥饿；他们以为被剥夺了这公共的光和那赐给所有人的阳光，比忍受连续饥饿和经历最悲惨的死亡要轻一些。这样，既然你们没有学会同情贫困，反而以不幸为乐，他们就满足你们永不知足的欲望，并且为他自己又为我们在地狱中燃起更猛烈的火焰。\par

9. 为叫你相信，正因这个缘故，才发生这些以及类似的事，我要告诉你一个无人能反驳的公认证据。另有一些穷人，心思轻浮不稳，不知如何忍受饥饿，反而什么都忍受，唯独饥饿不能忍。他们屡次试图用哀怜的姿态和言语来打动我们，发现毫无效果，便放弃了那些哀求；从此以后，连我们那些奇人异士也被他们超过了：有的嚼烂旧鞋的皮，有的把尖钉钉进自己头上，有的赤着肚子躺在结冰的水池里，另一些则忍受比这些更可怕的事，为要吸引那些不虔敬的围观者。而你，当这些事发生时，你却站在那里又笑又惊叹，把别人的苦难——我们共同本性在羞辱自己——当作一场好戏。一个凶恶的魔鬼还能做什么更多的事？其次，你大量给他钱，为让他更起劲地做这些事。而对于那祈祷、呼求天主、谦逊前来的人，你既不予回答，也不屑一顾；反而对他不断纠缠你的话，说出那些可恶的言辞：\Quote{这家伙该活着吗？或者还该呼吸、看见这太阳吗？}而对另一类人，你却又高兴又慷慨，好像你被指定去分发那种可笑而撒殚式的丑态的奖品。因此，对那些设立这些游戏、非要等到别人自我惩罚才肯施舍的人，更恰当地应当对他们说：\Quote{这些人该活着、该呼吸、或看见太阳吗？他们冒犯我们共同的本性，侮辱天主！}因为天主说：\Quote{施舍吧，我给你天国，}你却不听；但魔鬼把一个钉满钉子的头指给你看时，你立刻变得慷慨了。那充满如此多祸患的恶灵的计谋，对你的影响超过那带来无数祝福的天主的应许。即使有人拿出金子来阻止这些事发生或观看它们，你本应不惜一切努力和忍耐，为摆脱这极度的疯狂；但你却千方百计让它们发生，并且观看它们发生。那么，请问，你还问地狱之火有什么用吗？不，你再也别问那个了，而该问：为何只有一个地狱？因为这些搞出这个残忍无情把戏的人，以及那些为应当哭泣的事而发笑的人——甚至你们当中那些强迫他们如此丑行的人——该受多少惩罚？\par

\Quote{但我没有强迫他们，}你说。请问，那算什么，如果不是强迫？对那些更谦逊、流泪呼求天主的人，你连听都不耐烦；但对这些人，你却找到大量的银子，并招来许多人围观赏奇。\par

\Quote{好吧，让我们停止可怜他们。}你说，\Quote{你也命令这样做吗？}不，这不是怜悯，人啊，为了一点钱就要求这样严厉的惩罚，叫人为了必需的食物而自残，如此残忍可怜地把他们头上的皮肤切成许多片。\Quote{慢着，}你说，\Quote{又不是我们钉那些头。}我倒希望是你钉的，那样恐怖还不至于如此恐怖。因为杀一个人，比命令他自杀更为严重——而这正是这些人的情况。当他们被命令亲自执行这些邪恶命令时，他们忍受着更痛苦的折磨。\par

而这一切发生在安提约基雅，那里的人首先被称为基督徒，那里生长着最文明的人类，那里往昔爱德之果如此丰盛。因为他们不仅向近处的人，也向远方的人施舍，而且是在饥荒将到之时。\par

10. 那么，我们该做什么？你说。停止这种野蛮行径：并叫所有有需要的人相信，做这些事一无所获；如果他们谦逊地前来，必会得到你的慷慨。一旦他们知道这点，即使他们是最可怜的人，我保证，他们绝不会选择如此严厉地惩罚自己；不，他们甚至会感谢你，因为使他们摆脱了那种生活的嘲弄和痛苦。但如今，为了车夫你甚至肯典卖自己的儿女，为了舞者你肯丢弃你的灵魂，而为那饥饿的基督，你却连最小的零碎都不肯省下。但即使你给了一点银子，你就觉得好像已经倾尽所有，不知道：不是施舍本身，而是慷慨的施舍，才是真正的施舍。因此，先知宣告并称为有福的，不是那些简单施舍的人，而是那些慷慨施舍的人。因为他不是简单地说\Quote{他施舍了}，而是说什么？\BibleRef{咏 111:8}\BibleQuote{他散施了，他给了穷人。}因为，你若从大海中只给出一杯水，且连一个寡妇的慷慨你都不愿效法，又有什么益处？而你将如何说：\Quote{上主，求你按你的大仁慈怜悯我，按你丰厚的慈悲涂抹我的过犯，}而你自己却不按任何大仁慈去怜悯，甚至可能连一点点也吝于怜悯？因为我非常羞愧，我承认，当我看到许多富人骑着金辔的骏马，带着穿金戴银的仆从，有银制的卧榻以及更多其他奢华，而当需要给穷人施舍时，却变得比最穷的人还穷酸。\par

11. 但他们常说什么？\Quote{他，}他们说，\Quote{有教会的公共津贴。}那与你何干？因为就算我施舍，你也不会得救；即使教会施舍，你也不能涂去自己的罪。难道因此你就不施舍，因为教会应该给穷人？因为司铎们祈祷，你就从不祈祷吗？因为别人禁食，你就终日酗酒？难道你不知道，天主订立施舍，与其说是为穷人的缘故，不如说是为施舍者本身的益处？\par

但是，你怀疑司铎吗？首先，这件事本身就是一项大罪。不过，我不会追究得太过仔细。你要亲自去做这一切，这样你将获得双倍的赏报。事实上，我们为施舍所说的话，并不是要你献给我们，而是要你亲手去施予。因为如果你把施舍带给我，也许你还会被虚荣所俘获，而且常常也会因怀疑某些坏事而生气离去；但如果你凡事由自己去做，你就能免于冒犯和不合理的怀疑，并且你的赏报也更大。我说这些话，不是为了强迫你们把钱带到这里来；也不是出于对司铎被诽谤的愤慨。因为如果有人应当愤慨和悲伤，那你们才应当是我们的悲伤，你们说了这些坏话。因为那些被虚假而徒然地诽谤的人，赏报更大；而诽谤者所受的谴责和惩罚更重。因此，我说这些话不是为了他们，而是出于对你们的关心和挂念。在我们这一代人中，有人被怀疑，这有什么稀奇呢？就连那些效法天使、一无所有的圣人，我说的是宗徒们，在照顾寡妇的事上还发生了怨言\BibleRef{宗 6:1}，说穷人被忽视了？当时\Quote{没有人说他所拥有的哪一样东西是自己的，而是大家公用一切？}\BibleRef{宗 4:32}\par

那么，我们不要提出这些借口，也不要以为教会富有就可以推卸责任。但是，当你看到她财富的巨大时，也要记住她名册上穷人的众多、病人的众多、以及她无休止的开支。查问吧，审查吧，没有人禁止，甚至他们乐意给你一份账目。但我希望走得更远。就是说，当我们交出账目，证明我们的支出不少于收入，甚至有时更多时，我很乐意再问你们这个问题：当我们离开这里，听见基督说：\Quote{你看见我饿了，却没有给我吃的；我赤身露体，却没有给我穿的}，那时我们将说什么？我们将如何辩护？我们能抬出某个违背这些命令的人吗？或者某些被怀疑的司铎？祂说：\Quote{不，这与你何干？因为我控告你，}祂说，\Quote{那些你自己犯罪的事。这些事的辩护就是洗清你自己的过犯，而不是指出别人犯了一样的错误。}\par

事实上，教会因你们的吝啬才被迫拥有现在这样的财产。因为，如果人们都按法律行事，教会的收入本应是你们的善意，那既是安全之箱，又是用之不竭的宝库。但现在，当你们为自己积蓄财宝在地上，把一切都锁在自己的仓库里，而教会却被迫承担着成群的寡妇、贞女的歌咏团、旅客的寄居、外乡人的困苦、囚犯的不幸、病者和残疾者的需要，以及其他类似的原因，那该怎么办呢？难道要拒绝这一切，堵住这么多港口吗？那么谁能承受由此而来的沉船？以及从四面八方传来的哭泣、哀叹和悲号？\par

那么，我们不要随口说出我们心里想的话。因为现在，正如我刚刚所说，我们确实准备好向你们交出我们的账目。但即使情况相反，你们有腐化的教师掠夺和攫取一切，他们的邪恶也绝不能成为你们的借口。因为爱世人者和全智者、天主的独生子，看见一切，知道在如此漫长的时间和在如此广阔的世界里，会有许多腐化的司铎；祂为了防止下属因他们的疏忽而变得更粗心大意，除掉了所有漠不关心的借口；祂说：\Quote{在\Person{梅瑟}的座位上，}祂说，\Quote{坐着经师和法利塞人；所以，凡他们所吩咐你们的，你们都要遵行，但不要效法他们的行为。}意思是，即使你们有一个坏老师，如果你们不听从所说的话，这对你们也没有帮助。因为天主审判你们，不是看你们的老师做了什么，而是看你们听到了什么却不服从。所以，如果你们遵行所吩咐的事，你们就会大有胆量地站立；但如果你们不服从所说的话，即使你们指出一万个腐化的司铎，这也绝不能为你们辩护。因为\Person{犹达斯}也是宗徒，但这绝不能成为亵渎者和贪婪者的借口。也没有人在受控告时能说：\Quote{为什么宗徒是个贼，是亵渎者，是叛徒？}是的，这件事最将成为我们的惩罚和定罪，因为我们连别人的恶行都不能使我们改正。为此，这些事也被记录下来，好叫我们避免效法这类事情。\par

因此，让我们放下这个人和那个人，注意自己。因为\Quote{我们每个人都要向天主交自己的账。}为了能带着好的辩护交出这个账，让我们好好管理自己的生活，向穷人大方地伸手，知道这是我们的唯一辩护：表明我们正确地遵行了所吩咐的事；没有别的辩护。如果我们能提出这个，我们将逃脱地狱那些不可忍受的痛苦，并获得将来的美善事物；愿我们众人靠着我们的主耶稣基督的恩宠和怜悯，得以到达那里。愿光荣、权能和尊崇归于祂，偕同圣父和圣神，从今世直到永远，世世无穷。阿们。\par

\SourceRecord{Translated by Talbot W. Chambers.；Nicene and Post-Nicene Fathers, First Series；Vol. 12.；Edited by Philip Schaff.；Buffalo, NY: Christian Literature Publishing Co.,；1889.；<http://www.newadvent.org/fathers/220121.htm>.}

% structure: p0001 source=h1 target=section reason=page-title

\section{\WorkTitle{论格林多前书——第22讲}}

% structure: p0002 source=h2 target=subsection reason=relative-heading-rank; base=h2

\subsection{\BibleBook{格林多}{前书}9:13-14}

\BibleQuote{你们岂不知道，在圣物上供职的，就吃圣殿之物？侍候祭台的，就与祭台分有分？主也这样规定了：传福音的人，应靠福音生活。}\par

他极其小心地表明，接受并非被禁止。因此，在讲过这许多之后，他仍不满足，进而引用法律，举出一个比前一个更贴切的例证。因为提起牛，与明确引用关于司祭的法律，不是同一回事。\par

但请思量，我恳求，保禄在此的智慧，看他如何以赋予尊严的方式提及此事。因为他没有说：\Quote{在圣物上供职的，从奉献者那里接受。}而是说什么？\Quote{他们吃圣殿之物：}如此，接受者不受指责，给予者也不被高举。因此，接下来的话他也以同样的方式记下。\par

他也没有说：\Quote{侍候祭台的，从献祭者那里接受，}而是\Quote{与祭台分有分。}因为所献之物如今不再属于奉献者，而是属于圣殿和祭台。他没有说\Quote{他们接受圣物，}而是\Quote{吃圣殿之物，}再次表明他们的节制，且他们不应敛财致富。尽管他说他们\Quote{与祭台}分有分，他并非指平均分配，而是指给予他们当得的供养。然而宗徒的情况更要紧得多。因为前者中，司祭职是一种荣誉，而后者中，却是危险、杀戮和暴死。因此，所有其他例证加在一起，也不及这句话：\BibleQuote{我们若给你们播下了属神之物：}因为说\BibleQuote{我们播下，}他指出了风暴、危险、罗网和说不尽的患难，就是他们在宣讲中忍受的。尽管优势如此之大，他却不愿贬低旧约，也不愿抬高属于己身之事：他甚至收敛自己的事，不按危险计算优势，而是按恩赐的伟大。因为他没有说\Quote{若我们冒了生命危险}或\Quote{暴露于罗网，}而是\BibleQuote{若我们给你们播下了属神之物。}\par

至于司祭的部分，他尽可能地高举，说\Quote{在圣物上供职的}和\Quote{侍候祭台的，}借此意在指出他们持续的服事和忍耐。再者，正如他论及犹太人中的司祭，即肋未人和大司祭，他也表达了每一等级，包括较低和较高的：一个通过说\Quote{在圣物上供职的，}另一个通过说\Quote{侍候祭台的。}因为并不是所有人都被命令做同样的工作；有些受托粗重之职，有些受托更崇高的职任。因此，概括这一切，免得有人说：\Quote{为何与我们谈论旧约？你们不知道如今是更圆满诫命的时候吗？}在所有这些话题之后，他放置了最强有力的话，说：\par

\BibleQuote{主也这样规定了：传福音的人，应靠福音生活。}\newline{}\BibleRef{格前 9:14}\par

甚至在此他也没有说他们由人供养，而是如司祭的例子中\Quote{圣殿}和\Quote{祭台}一般，同样在此是\Quote{福音；}又如同那里他说\Quote{吃，}这里说\Quote{生活，}不是做买卖或积蓄财宝。因为祂说：\BibleQuote{工人当得他的工价。}\par

% structure: p0010 source=h2 target=subsection reason=relative-heading-rank; base=h2

\subsection{\BibleBook{格林多}{前书}9:15}

2. \BibleQuote{但我没有用过这些：}\par

那么，有人说，如果你现在没有用，却打算将来用，并为此说这些话，那又如何？绝非如此；因为他迅速纠正了这想法，如此说：\par

\BibleQuote{我写这些话，并非为在我身上也是这样行。}\par

请看他是以何等激烈的方式否认并拒绝此事：\par

\BibleQuote{我宁愿死，也不愿有人使我的夸耀落空。}\par

他不只一次两次，而是多次使用这话。因为上面他说道：\BibleQuote{我们没有用这权利：}在此之后又说：\BibleQuote{免得我滥用我的权利：}这里则说：\BibleQuote{但我没有用过这些。}\Quote{这些；}什么？许多的例证。也就是说，有许多事给予我许可：兵士、农夫、牧人、宗徒、法律、我们对你们所行的、你们对他人所行的、司祭、基督的规定；我未被这些中的任何一个诱导废除我自己的规则而接受。不要对我说过去的事：（虽然我能说，为此我在过去也忍受了许多，）然而我不仅仅依靠这一点，并且关于将来我也保证：我宁愿饿死，也不愿被剥夺这些荣冠。\par

\BibleQuote{我宁愿死，}他说，\Quote{也不愿有人使我的夸耀落空。}\par

他没有说\Quote{有人废除我的规则，}而是\Quote{我的夸耀。}免得有人说：\Quote{他确实行了，但不是欣然，而是带着哀叹和忧伤，}为了显示他喜乐之极致和热忱之丰盈，他甚至将此称为\Quote{夸耀。}他不但不使自己苦恼，反而夸耀，宁愿死也不愿失去这\Quote{夸耀。}对他而言，那项举动甚至比生命本身更宝贵。\par

3. 接着，他又从另一考虑高举它，并表明这是一件大事，并非为显扬自己（他远非那种性情），而是为表明他喜乐，并更充份地除去一切猜疑。因为为此，如我先前所说，他也称之为夸耀：他说了什么？\par

% structure: p0020 source=h2 target=subsection reason=relative-heading-rank; base=h2

\subsection{\BibleBook{格林多}{前书}9:16-18}

\Quote{因爲我若傳福音，並無可誇；因爲這是不得已的事；若不傳福音，我便有禍了！我若自願作這事，便有賞報；若不情願，我卻受託付了這責任。那麼，我的賞報是甚麼呢？就是當我傳福音時，我能白傳福音，而不享用我在福音中的權利。}\par

你說甚麼？請告訴我。\Quote{你若傳福音，並無可誇；但若白傳福音，就有可誇？}難道這比那更大嗎？決不是；但從另一角度看，它有一些優越之處，因爲一個是命令，另一個是出自自由意志的善行：凡超出命令的事，在這方面有大的賞報；但遵從命令的事，則不那麼大。所以在這方面，他說一個比另一個大；並非在事物本身的性質上。因爲有甚麼能與傳福音相比呢？它使人甚至與天使競逐。然而，既然一個是命令和債務，另一個是自由意志的熱忱，在這方面這個比那個大。因此，他解釋同一點時，說了我剛才提到的那段話：\par

\Quote{我若自願作這事，便有賞報；若不情願，我卻受託付了這責任；}這裏的\Quote{自願}和\Quote{不情願}是指這事是否受託付於他。我們必須這樣理解\Quote{不得已的事}這句話，並非指他違背自己的意願作了這些事，絕對不是，而是如同他被命令所束縛，與前面提到的接受報酬的自由相對。所以基督也對門徒說：\BibleRef{路 17:10}\Quote{你們作完了一切，要說：我們是無用的僕人，因爲我們只作了我們應作的事。}\par

\Quote{那麼，我的賞報是甚麼呢？就是當我傳福音時，我能白傳福音。}那麼，請問，伯多祿沒有賞報嗎？誰能有像他那樣大的賞報呢？我們又該說其他宗徒呢？他怎能說：\Quote{我若自願作這事，便有賞報；若不情願，我卻受託付了這責任？}你看見他的智慧嗎？他沒有說：\Quote{若不情願，}我沒有賞報，而是說：\Quote{我受託付了這責任：}意思是即使這樣，他也有賞報，但如同一個完成命令的人所得的賞報，而非出於自己慷慨、超過命令的人所得的那種賞報。\par

\Quote{那麼，賞報是甚麼呢？}他說：\Quote{當我傳福音時，我能白傳福音，而不享用我在福音中的\OriginalTerm{權利}{right}。}請看，他自始至終使用\Quote{權利}一詞，暗示這點，如我常說的：即那些接受的人也不應受責備。但他加上\Quote{在福音中}，一方面是爲了顯示其合理性，另一方面也是禁止我們將此事擴展到一切情況。因爲教師應得到供養，但純粹的懶惰者則不該。\par

% structure: p0026 source=h2 target=subsection reason=relative-heading-rank; base=h2

\subsection{\BibleRef{格前 9:19}}

4. \Quote{我雖然是自由的，不受任何人管轄，卻使自己成了眾人的奴僕，爲要贏得更多的人。}\par

在這裏，他又引進另一個更高的進展。因爲即使不接受報酬已是一件大事，但他將要說的遠超過這個。他說的是甚麼呢？\Quote{我不僅沒有接受，}他說，不僅沒有使用這權利，反而使自己成爲奴僕，而且是多樣普遍的奴役。因爲不僅在金錢上，而且在許多不同的工作上，我堅持了同樣的原則：我使自己成爲奴僕，雖然本不受任何人管轄，在任何方面都沒有必要（這就是\Quote{雖然是自由的，不受任何人管轄}的意思）；我不僅向一個人，而是向全世界成爲奴僕。\par

因此他又加上：\Quote{我使自己成了眾人的奴僕。}就是說：\Quote{傳福音是命令，宣講所託付的事；但計劃和設計無數其他事情，全是出於我的熱忱。因爲我只負有投資銀子的義務，而我卻爲獲得回報而作了一切，嘗試了超出命令的事。}這樣，他出於自由選擇、熱忱和對基督的愛而作一切事，對人類的救恩懷有無饜的渴望。因此，他常常遠遠超越劃定的界線，處處跳得比天還高。\par

5. 接着，他提到了他的奴役，在隨後的部分描述了其各種方式。這些方式是甚麼呢？\par

% structure: p0031 source=h2 target=subsection reason=relative-heading-rank; base=h2

\subsection{\BibleRef{格前 9:20}}

\Quote{我對猶太人，就成爲猶太人，爲要贏得猶太人。}這事如何發生呢？當他行割損禮時，是爲要廢除割損禮。所以他不說\Quote{猶太人}，而說\Quote{如同猶太人}，這是明智的安排。你說甚麼？世界的使者，觸及高天、在恩寵中如此光輝奪目的人，竟一下子降到如此低微的地步？是的，因爲這是上升。你不僅要看他的下降，也要看他扶起那彎腰的人，並將他提昇到自己那裏。\par

\BibleQuote{对于在法律之下的，我如同在法律之下的——我自己并不在法律之下——为要赢得那些在法律之下的。} 这要么是对前述内容的解释，要么是暗示了除前述之外的另一件事：他称那些原本从起初就是犹太人为犹太人；但\Quote{在法律之下的}指的是归化者，或那些成了信徒却仍固守法律的人。因为他们不再是犹太人，却仍\Quote{在法律之下}。他何时在法律之下？当他剃头时；当他献祭时。这些事之所以做，并非因为他心意改变——因为这样的行为本是恶行——而是因为他的爱德屈尊俯就。为要使那些真正是犹太人的人归向这信仰，他亲自成为这样的人，并非真正成为，只是表面上如此，实际上并非如此，也不是出于如此心意而行这些事。确实，他既如此热心劝化别人，且行这些事只为使那些行这些事的人摆脱那种卑劣，怎么可能心怀恶意呢？\par

% structure: p0034 source=h2 target=subsection reason=relative-heading-rank; base=h2

\subsection{\BibleRef{格前 9:21}}

\BibleQuote{对于没有法律的，我如同没有法律的。} 这些人既不是犹太人，也不是基督徒，更不是希腊人；而是\Quote{在法律之外}，正如\Person{科尔乃略}，以及若有其他像他的人。因为他也出现在这些人之列，曾采纳他们许多方式。但有人说，他暗指他与雅典人从祭坛上的铭文所作的谈论，因此他说：\Quote{对于没有法律的，我如同没有法律的。}\par

然后，免得有人认为这是心意的改变，他补充道：\BibleQuote{我不是在天主面前没有法律，而是在基督的法律之下}；也就是说，\Quote{我绝非没有法律，我不单是在法律之下，而是拥有那远超旧法律的法律，即圣神和恩宠的法律。}因此他又补充说：\Quote{对基督}。然后，在使他们确信他的判断之后，他也说明了这种屈尊俯就的益处，说道：\BibleQuote{为要赢得那些没有法律的。}而且他处处提出他屈尊俯就的原因，甚至不止于此，而是说：\par

% structure: p0037 source=h2 target=subsection reason=relative-heading-rank; base=h2

\subsection{\BibleRef{格前 9:22}}

\BibleQuote{对软弱的人，我就成为软弱的，为要赢得软弱的人。} 在此他正论及他们的情况，也正是为这情况说了这一切。然而，那些是更为重大的事，但这更切题；因此他也把它放在它们之后。确实，他在致\WorkTitle{罗马书}中也同样做了，当他责备关于食物的问题时；在其它许多地方也是如此。\par

接下来，为免逐一列举浪费时间，他说：\BibleQuote{我向众人成为一切，为能藉一切方式拯救一些人。}\par

你看到这被推行到多远吗？\Quote{我向众人成为一切}，但他并不指望拯救所有人，而是为能拯救即使只是少数人。我承受了如此巨大的关怀和服事，如同一个自然想要拯救所有人的人，却远不能希望赢得所有人：这实在是宽宏大量，也是炽热热忱的证明。同样，播种者随处播种，并非所有种子都得救，尽管如此，他尽了本分。而且，在提到得救者人数稀少之后，他又加上\Quote{藉一切方式}，安慰了那些为此忧伤的人。因为尽管不可能所有种子都得救，但也不能所有都灭亡。因此他说\Quote{藉一切方式}，因为一个如此热切热心的人必定会有所成功。\par

% structure: p0041 source=h2 target=subsection reason=relative-heading-rank; base=h2

\subsection{\BibleRef{格前 9:23}}

\BibleQuote{我作一切事，都是为福音的缘故，为能同享福音的恩惠。}\par

也就是说，为使我自己也似乎添加了一些自己的贡献，并能有分于为信者存留的冠冕。因为正如他提到\Quote{靠福音生活}，即靠信徒生活；同样在这里，\Quote{为能同享福音，能与那些信从福音的人一同有分}。你难道看不到他的谦卑吗？在赏报的补偿中，他把自己放在众人之中，尽管他在劳苦上超过众人！由此也显明他也会在赏报上如此。然而，他不要求享受头奖，而是满足于若能与他人同享为他们存留的冠冕。但他说这些话，并非因为他为任何赏报而行，而是借此至少能吸引他们，藉这些希望劝诱他们为弟兄的缘故作一切事。你看到他的智慧吗！你看到他完美的卓越吗？他怎样作了超出命令的事，在合法收取时却不收取。你看到他屈尊俯就的无比伟大吗？那\BibleQuote{在基督的法律之下}并遵守那最高法律的人，对\BibleQuote{那些没有法律的}，却\BibleQuote{如同没有法律的}；对犹太人，如同犹太人，在两方面都显出超群出众，超越一切。\par

6. 你也要这样做，不要以为你既然卓越，当你为弟兄的缘故屈就一些卑下时，你就降低了自己。因为这不是跌倒，而是下降。因为跌倒的人匍匐在地，几乎无法再起来；但下降的人必要带着极大益处再升起来。正如保禄确实独自下降，却与整个世界一同上升：他不是在扮演角色，因为如果他在扮演，就不会寻求得救者的益处。因为\OriginalTerm{伪善者}{hypocrite}寻求人的丧亡，他伪装是为了收取，而不是为了付出。但宗徒并非如此：他更像医生、教师、父亲：医生对病人、教师对学生、父亲对儿子，屈尊俯就成为矫正，而不是伤害；他自己也是如此行。\par

为表明所说之事并非虚饰，在一个他并非被迫行或说此类事情、而意在表达其深情与信赖的事例中，请听他说道：\BibleRef{罗 8:39} \Quote{是生命，是死亡，是天使，是掌权者，是有能的，是现在的事，是将来的事，是高天，是深渊，或是任何受造之物，都不能使我们与天主的爱相隔绝，这爱是在我们的主基督耶稣内的。}你可见一种比火更炽烈的爱吗？让我们也这样爱基督罢。因为实在说来，只要我们愿意，这是容易的。宗徒并非生来如此。你瞧，正因如此，他早年的生活才被记录下来，是与后来如此相反，好叫我们明白这事是出于选择，并且对于愿意的人，一切都是容易的。\par

因此，我们不要失望。即使你是辱骂者、贪婪者，或无论你是什么样的人，请想想保禄曾是\Quote{亵渎者、迫害者、凌辱者、罪人之魁首}\EditorialNote{弟茂德前书 1:13-16}，而他却突然上升到德行的至极顶峰，他先前的生活并未成为他的障碍。没有人像他先前那样疯狂地执着于恶习，正如他执着于攻击教会。那时他简直将性命投入其中；因为他没有一万只手可用以一同砸死斯德望，他便为此烦恼。然而，即便如此，他仍设法用更多的手来砸死他——即那些他看守其衣服的假证人的手。又如，他像野兽一样闯入房屋，毫不留情地冲进去，拖曳撕裂男女，使一切充满骚动、混乱与数不清的争斗。举例来说，他是如此可怕，以至宗徒们\BibleRef{宗 9:26}甚至在他那极为荣耀的改变之后，仍不敢与他交往。然而，在这一切之后，他成了那样的人：因为我无需多言。\par

7. 如今那些主张命运必然性、反对自由意志的人在哪里？让他们听听这些事，让他们住口。因为没有什么能阻止愿意的人成为善人，即使他先前是极卑劣的人。事实上，我们在这方面更有倾向，因为德行合乎我们的本性，而恶习违反本性，如同疾病与健康。因为天主赐给我们眼睛，不是叫我们淫荡地观看，而是叫我们赞叹祂的工程，敬拜造物主。眼睛的用途从所见之物即可显明。因为太阳和天空的光辉，我们从不可估量的远处就能看见，而女子的美貌却不能从那么远辨明。你看见吗？我们的眼睛主要为此而被赐予。同样，祂造了耳朵，叫我们领受的不是亵渎的话语，而是救恩的道理。因此你瞧，当耳朵听到任何不和谐的声音时，我们的灵魂甚至身体都颤抖。因为经上说：\BibleRef{德 27:5} \Quote{多发誓者的谈论使人毛发竖立。}若我们听见任何残酷无情之事，我们的肉体也要战栗；但若听见庄重仁慈之事，我们就甚至欢欣踊跃。再者，若我们的口说出卑贱的话，就会使我们羞愧躲藏；若说出庄重的话，就会轻松自在、毫无拘束地说出。因为没有人会为合乎本性的事羞愧，却会为违反本性的事羞愧。手在偷窃时隐藏自己，寻求借口；但在施舍时，它们甚至夸耀。所以，只要我们愿意，我们从各方面都有强烈的向善倾向。但若你对我谈论来自恶习的快乐，请想想这一点：我们从德行中收获的快乐更多。因为拥有良好的良心、受众人景仰、怀有美好的希望，对于懂得快乐本质的人而言，是万物中最愉快的；相反，对于懂得痛苦本质的人而言，最痛苦的事莫过于受众人责备、受自己良心控告、在将来和现时都战栗恐惧。\par

为使我的话更显明，让我们假设一种情况：有一人拥有妻子，却玷污邻人的婚床，并从这邪恶的掠夺中取乐，享受他的情妇。然后再将另一个爱自己配偶的人与他对立。为使胜利更大且更明显，让那仅享受自己妻子的人，也对那个淫妇有贪恋，却克制自己的情欲，不作任何恶事（虽然这也不算纯粹的贞洁）。然而，我们姑且作出超过必要限度的假设，好叫你确信德行的快乐何其大，为此我们如此编排故事。\par

如今，我们把他们召集在一起，不妨问一问他们：谁的生活更愉快？你将听到一人因战胜情欲而夸耀欢悦；而另一人——更确切地说，无需等待他告知任何事。因为你会看到他，即使他无数次否认，仍然比监牢中的人更痛苦。他畏惧并猜疑一切：自己的妻子、奸妇的丈夫、奸妇本人、仆人、朋友、亲属、墙壁、阴影、自己，而最糟糕的是，他的良心日夜高声责斥他。若他还想起天主的审判宝座，他甚至站立不住。快感短暂，但随之而来的痛苦却无休止。因为无论黄昏还是黑夜，在旷野、城市或任何地方，控告者都萦绕着他，指着一把磨亮的利剑和无法忍受的刑罚，以那恐惧消耗并折磨他。而另一人，那贞洁的人，却免于这一切，自由自在，安心地面对妻子、儿女、朋友，以坦然的目光迎接众人。若那爱慕却自制的人能享受如此大的快乐，那么那不起情欲、真正贞洁的人，将获得何等甜美平静的港湾和心境？正因如此，你能见到少数奸淫者，而许多贞洁者。若前者更愉快，大多数人自会偏好它。不要向我提及法律的恐惧，因为约束他们的并非此，而是过度的不合理，以及其中的痛苦多于快乐，还有良心的判决。\par

8. 奸淫者便是如此。现在，若你乐意，让我们将贪婪者带到你面前，再次揭示另一种不法情欲。我们也必看见他畏惧同样的事，无法享受真正的快乐：因为他既想起自己所亏负的人，又想起同情他们的人，以及众人对他自己的公论，心中便有无数的烦乱。\par

这并非他唯一的烦恼，甚至他所钟爱的对象他也无法享用。因为贪婪者便是如此：他们拥有不为享用，而为不享用。若这对你而言似是谜语，那么请听接下来更糟糕、更令人困惑的事：他们不仅因不敢随心所欲使用财物而失去享用的快乐，还因永不满足、恒常处于干渴之中而受苦——还有什么比这更痛苦？但义人并非如此，他既免于战栗、仇恨、恐惧和这不治的干渴；所有人都咒骂前者，同样所有人都合力祝福后者；前者无友，后者无仇。\par

既然如此，这些事既已公认，还有什么比恶更令人不悦，比德更令人愉悦？不，即使我们永远谈论，无人能在言语中完全表达恶的痛苦或德的喜乐，直至我们亲身经历。那时，当我们充分品尝了德行的蜜，才会发现恶比胆汁更苦。恶现在也绝无不令人不悦、可憎、沉重，连其追随者也不否认；但当我们远离它时，才更清楚地分辨它命令的苦涩。若众人奔向它，这并不奇怪；因为孩童也常选择不甚愉快的事，轻视更令人愉悦的；病人为片刻满足而失去持久且更确定的喜乐。但这出于陷于任何偏好之人的软弱和愚昧，而非事物本身的性质。因为正是有德之人生活在喜乐中；他才是真正富有、真正自由的。\par

若有人将其他一切归于德行——自由、安全、无忧、不惧人、不疑人——却不愿将快乐归于它，我承认，我只能付之一笑。因为快乐若不是无忧、无惧、无沮丧、不受任何人支配，又是什么？那么告诉我，谁在快乐？是那疯狂而抽搐、被各种欲望驱使、甚至不能自主的人，还是那摆脱这一切波涛、安顿于智慧之爱中如同在港湾的人？岂非后者显然？而这似乎是德行独有的事。所以恶只有快乐之名，却没有其实。在享受之前，它是疯狂而非快乐；享受之后，它随即熄灭。若开始时和之后都觉察不到它的快乐，那它何时何地出现？\par

为使你更清楚理解我所说的，让我们用一个例子来检验论证的力量。请想：一人迷恋一位美丽可爱的女子；这人只要得不到他的欲望，就如癫狂失心的人；但一旦得到，他便熄灭了欲火。若在开始他并无快乐（因为此事乃是疯狂），结束时也无快乐（因情欲的满足冷却了他的狂想），那么快乐究竟在何处？但我们的行为并非如此：我们在开始就摆脱了一切烦扰，快乐至终保持鲜活；不，我们的快乐没有终结，我们的福乐没有界限，这快乐永不消失。\par

因此，基于这一切考量，若我们爱快乐，就让我们抓紧德行，以便赢得现今和将来的福善；愿我们众人赖恩宠和仁慈等而达到。\par

\SourceRecord{Translated by Talbot W. Chambers.；Nicene and Post-Nicene Fathers, First Series；Vol. 12.；Edited by Philip Schaff.；Buffalo, NY: Christian Literature Publishing Co.,；1889.；<http://www.newadvent.org/fathers/220122.htm>.}

% structure: p0001 source=h1 target=section reason=page-title

\section{格林多前书讲道集第廿三篇}

% structure: p0002 source=h2 target=subsection reason=relative-heading-rank; base=h2

\subsection{格林多前书 9:24}

\BibleQuote{你们岂不知道，在赛跑中，跑众人都跑，但只有一人得奖赏吗？}\par

具有指出 \OriginalTerm{迁就}{condescension} 的多方面益处，并且这是最高的 \OriginalTerm{成全}{perfectness}，而他自己既已超越众人达到成全，或更确切地说，因拒绝接受而超越了成全，又再次降到众人之下；并让我们知道这两者各自的时候，既为成全也为迁就；他在下面更严厉地触动他们，暗中暗示，他们所作且被视为成全的事，其实是一种多余且无用的劳苦。他没有如此明说，以免他们变得傲慢；但他所采用的论证方法使这一点变得明显。\par

说了这些话，说他们得罪基督，毁灭弟兄，并且除了加上爱德，这种成全的知识毫无益处之后，他又举出一个常见的例子，说：\par

\BibleQuote{你们岂不知道，在赛跑中，跑众人都跑，但只有一人得奖赏吗？}他说这话，并非在这里也只有一人从众人中得救；绝非如此；而是为了表明我们应尽的最大努力。因为正如在那里，虽然许多人进入跑道，但得冠冕的并不多，而只临到一人；同样，仅仅进入竞赛，并涂抹自己或摔跤，是不够的。同样，在这里，相信并随意争斗是不够的；除非我们如此奔跑，以致最终显明自己是无可指责的，并接近奖赏，否则对我们毫无益处。因为即使你自以为按照知识是完全的，你还没有达到全部；他暗示这一点，说：\BibleQuote{你们要这样跑，好使你们获得奖赏。}他们当时似乎还没有达到。\par

% structure: p0007 source=h2 target=subsection reason=relative-heading-rank; base=h2

\subsection{格林多前书 9:25}

\BibleQuote{凡在竞技场上较力的，在一切事上都有节制。}\par

什么是\InnerQuote{一切事}？他并非在一件事上禁戒，在另一件事上犯罪，而是完全制服贪食、纵欲、醉酒和所有情欲。\Quote{为此，}他说，\Quote{甚至在异教赛会中也是如此。因为……}如果在那里，冠冕只落在一个人身上，情况尚且如此，那么在这里，竞赛的激励更为充足，岂不更应如此？因为在这里，并非只有一人得冠冕，而且赏报也远远超过劳苦。因此他也以此使他们羞愧，说：\BibleQuote{他们这样做，是为了能朽坏的冠冕，而我们是为了不朽坏的。}\par

% structure: p0010 source=h2 target=subsection reason=relative-heading-rank; base=h2

\subsection{格林多前书 9:26}

2. \BibleQuote{我这样跑，不是无定向的。}\par

这样，他用外人来使他们羞愧之后，接着也把自己摆出来，这是一种极好的教导方法；我们随处可见他这样做。\par

但什么是\Quote{无定向}？\Quote{瞄准某个目标，}他说，不是随意和徒然，像你们那样。你们进入偶像庙宇，为真理展示那种成全，有什么益处呢？毫无益处。但我并非如此，我所作的一切，都是为邻人的得救。无论我显示成全，是为他们；或显示迁就，也是为他们：无论我超越伯多禄而拒收报酬，是为使他们不致跌倒；或降到众人之下，受割礼并剃头，是为使他们不致被颠覆。这便是\Quote{无定向}。但你们，为什么在偶像庙宇里吃，告诉我？你们不能提出任何合理的理由。因为\BibleQuote{食物并不能使我们推荐给天主；我们若不吃，也不亏损，若吃，也不增添。}\BibleRef{格前 8:8}显然你们在随意奔跑：因为这就是\Quote{无定向}。\par

\BibleQuote{我这样斗拳，不像打空气。}他说这话，再次暗示他的行为不是随意或徒然的。\Quote{因为我有一个可以打击的对象，即魔鬼。但你却没有打击他，只是白费力气。}\par

至此，他完全忍耐他们，这样说话。因为在前文他对他们颇为严厉，现在反而保留责备，把最深的创伤留到讲论的结尾。因为这里他说他们的行为是随意和徒然的；但后来他又表明，这至少会使他们自己的灵魂遭致完全的毁灭，而且即使不伤害他们的弟兄，他们自己胆敢这样做也不是无罪。\par

% structure: p0016 source=h2 target=subsection reason=relative-heading-rank; base=h2

\subsection{格林多前书 9:27}

\BibleQuote{我痛击我的身体，使它成为奴隶，免得我给别人传报了，自己反被淘汰。}\par

这里他暗示他们受肚腹的贪欲支配，并放松缰绳给它，在成全的伪装下满足自己的贪婪；这个思想从前他也曾竭力表达，说过：\BibleQuote{食物是为肚腹，肚腹是为食物。}\BibleRef{格前 6:13}既然奸淫由奢侈引起，它也产生了偶像崇拜，他自然常常抨击这种病；并指出他为福音受了多大的苦，他也把这列入其中。\Quote{我既已}他说，\Quote{超越了诫命，而这对我也不是轻省的事：}（\Quote{我们忍受一切，}据说，）\Quote{所以在这里，我也要付出许多辛劳，以便节制地生活。虽然欲望顽固，肚腹的暴政，我仍束缚它，不让自己屈服于激情，而是忍受一切劳苦，以免被它拉走。}\par

\Quote{因为，我恳求你们，不要以为我轻易达到了这理想的结果。因为这是一场比赛，一场多重的搏斗，一种暴虐的本性不断起来反抗我，试图挣脱。但我并不容忍它，而是压制它，用许多斗争使它服从。} 他说这话，是免得有人因行善的艰辛而绝望地退出争战。因此他说：\Quote{我痛击它，使它成为奴隶。} 他没有说：\Quote{我杀死它；} 也没有说：\Quote{我惩罚它}，因为肉身不是该恨恶的，而是说\Quote{我痛击它，使它成为奴隶；} 这是主人而非仇敌的作为，是教师而非对手的作为，是运动教练而非对手的作为。\par

\Quote{免得我给别人传了福音，自己反而被弃绝。}\par

如果\Person{保禄}——他教导了这么多人，并且在传道之后成了天使，担负起全世界的领导——还惧怕这一点，那么我们能说什么呢？\par

因为，\Quote{不要以为，} 他说，\Quote{你们相信了，这就足以使你们获得救恩。因为如果我——传道、教导、使无数人皈依——都不足以获得救恩，除非我也表现出自己无可指责的行为，那么对你们来说，就更不够了。}\par

3. 然后他再次用其他比喻。正如上面他引用了宗徒的例子、普遍习俗的例子、司铎的例子和他自己的例子，同样在这里，他先列举了奥林匹克竞技和他自己赛跑的例子，然后继续引用旧约的历史。因为他要说的话可能有些令人不悦，所以他把劝诫泛化，不仅谈论眼前的话题，也普遍地谈论\Place{格林多}人中的一切恶行。在异教竞技的例子中，他说\Quote{你们不知道吗？}，但在这里：\par

% structure: p0024 source=h2 target=subsection reason=relative-heading-rank; base=h2

\subsection{\BibleRef{格前 10:1-4}}

\Quote{弟兄们，我不愿意你们不知道。}\par

他说这话，暗示他们对于这些事情并不十分明白。而你不想让我们不知道的是什么呢？\par

\BibleQuote{我们的祖先，} 他说，\BibleQuote{都曾在云彩下，都曾从海中经过，都曾在云彩和海中受圣洗归于梅瑟，并且都吃过同样的神粮，都喝过同样的神饮；因为他们饮自伴随他们的神石，而那神石是基督。然而，他们大多数人不为天主所喜悦。}\par

他说这些话为了什么？为了指出，正如他们享用如此大恩毫无获益，同样，这些人在获得圣洗、参与神妙的奥迹之后，除非他们继续活出与这恩宠相称的生命，也毫无益处。因此，他也引用了圣洗和奥迹的预象。\par

但\Quote{他们受圣洗归于梅瑟}是什么意思？就像我们相信基督及其复活，并受圣洗，因为我们自己也将分享同样的奥迹；因为，\Quote{我们受圣洗，} 他说，\Quote{为死人，} 即为了我们自己的身体；同样，他们信任梅瑟，即看见他先渡过，也冒险进入水中。但他为了使预象接近真理，并没有这样表述，而是用了真理的术语来描述预象。\par

此外，这是水泉的象征，而接下来的则是圣桌的象征。因为正如你吃主的身体，他们吃玛纳；正如你饮血，他们饮从石中流出的水。因为所产生的是可感觉的事物，但它们却是以属灵的方式展现的，不是按照本性的秩序，而是按照恩赐的慈意，并且与身体一同滋养灵魂，引导它走向信德。因此，你看，关于食物他没有作评论，因为它完全不同，不仅在于方式，也在于本性（因为那是玛纳）；但关于饮品，因为供应方式本身非凡且需要证明，所以他说\Quote{他们喝了同样的神饮}后，补充道：\Quote{因为他们饮自伴随他们的神石}，然后又加上：\Quote{而那神石是基督。} 因为不是那石头的本性发水（这是他意思），否则它在此之前也会涌出水来；而是另一种神石，属灵的，完成了整个工作，即基督，他处处与他们同在，行了一切奇迹。因此他说\Quote{伴随他们的}。\par

你体会到\Person{保禄}的智慧吗？他如何在这两种情况下都指出他是施予者，从而使预象接近真理？\Quote{那位把这些东西摆在他们面前的那一位，} 他说，\Quote{同样的那位也为我们预备了这圣桌；同一位带领他们穿越红海，也带领你们穿越圣洗；在他们面前摆上玛纳，在你们面前摆上他的体血。}\par

4. 因此，就他的恩赐而言，情况如此；现在让我们也观察后续，并考虑：当他们表现出不配得恩赐时，他是否宽恕了他们？不，你不能这样说。因此他也补充道：\Quote{然而，他们大多数人不为天主所喜悦；} 尽管他以如此大的尊荣荣耀了他们。是的，这对他们毫无益处，大多数人都丧亡了。事实上，他们都丧亡了，但他为了不显得也预言这些人的完全毁灭，因此他说\Quote{他们大多数}。尽管如此，他们是无数的人，但人数对他们毫无益处；所有这些也是那么多的爱的标记；但这也没有益处，因为他们自己没有结出爱的果实。\par

如此，既然多数人不相信关于地狱所说的事，因为它既不在眼前也不可见，他便用过去所发生的事作为证据，表明天主确实惩罚一切犯罪的人，即使祂曾赐予他们无数的恩惠：\Quote{因为若你们不相信将来的事，}他是这样说的，\Quote{那么过去的事你们总不会不相信吧。}试想，例如祂赐予他们多么大的恩惠：从埃及和那里的奴役中释放了他们，使大海成为他们的道路，从天上降下玛纳，从地下涌出奇异而惊人的泉水；祂处处与他们同在，行奇事并四面护卫他们；然而，既然他们丝毫没有显示出与此恩赐相称的行为，祂并没有宽恕他们，反而将他们全部毁灭。\par

% structure: p0034 source=h2 target=subsection reason=relative-heading-rank; base=h2

\subsection{格前 10:5}

\Quote{因为他们都倒毙在旷野里。}他说这话，既表明全面的毁灭，也表明天主所施的惩罚和报复，并且他们甚至没有达到为他们所预备的赏报。当天主向他们行这些事时，他们并不在应许之地，而是在外面和某处远离那地方；祂以此双重惩罚他们：既不许他们看见那地（尽管这地曾应许给他们），又以实际的严厉惩罚待他们。\par

你说：这些与我们何干？的确，它们与你们有关。因此他又补充说：\par

% structure: p0037 source=h2 target=subsection reason=relative-heading-rank; base=h2

\subsection{格前 10:6}

\Quote{这些事是我们的预像。}\par

因为正如恩赐是预像，惩罚也是预像；正如圣洗和筵席曾被预言式地描绘，同样，通过所发生的事，那临到不配此恩赐之人身上的惩罚的确定性，为了我们的缘故而被预先宣告，好使我们藉这些榜样学会节制。因此他又补充说：\par

\Quote{叫我们不要贪恋恶事，像他们那样贪恋。}因为在恩赐中，预像先行，实体随後，在惩罚中也将如此。你看见他如何不仅表明这些人将受惩罚，而且表明程度比那些古人更严厉吗？因为如果一个是预像，另一个是实体，那么惩罚就必然如同恩赐一样超越前者。\par

看他首先处理的是谁：那些在偶像庙里吃祭物的人。因为他说了\Quote{叫我们不要贪恋恶事}，这是笼统的说法，接着他提出具体的，暗示他们每一项罪恶都来自恶的贪恋。首先他说了：\par

% structure: p0042 source=h2 target=subsection reason=relative-heading-rank; base=h2

\subsection{格前 10:7-8}

\Quote{也不要拜偶像，像他们有些人一样；正如经上所记载：\InnerQuote{百姓坐下吃喝，起来玩耍。}}\par

你听见他如何甚至称他们为\Quote{拜偶像的}吗？这里他确实宣告了，但随后提出了证明。他也指出了他们跑向那些筵席的原因，就是贪食。因此，他说了\Quote{叫我们不要贪恋恶事}，又加上不要\Quote{拜偶像}之后，他指出了这种过犯的原因，就是贪食。\Quote{因为百姓坐下，}他说，\Quote{吃喝，}又加上其结果：\Quote{起来玩耍。}\Quote{因为他们，}他说，\Quote{从放纵情欲堕入偶像崇拜；同样，恐怕你们也会从一个堕入另一个。}你看见他如何表示这些自以为是成全的人，比他们所指责的其他人更不成全吗？不仅在这一方面，即他们不肯容忍弟兄们，而且在于一方因无知而犯罪，另一方却因贪食而犯罪。并且从前者的灭亡，他推算出对后者的惩罚，但他不允许后者将自己的罪因推给别人，反而宣告他们既为自己的伤害负责，也为自己负责。\par

\Quote{也不要行淫，像他们有些人行淫。}为何他在这里再次提到淫乱，之前已经详细讨论过了？保禄的习惯是，当他指控多种罪行时，既逐一列举，分别继续他提出的主题，又在论及其他事时再次提及前者；同样，天主在旧约中也常常这样做，针对每一次过犯，提醒犹太人关于牛犊的事，将那罪摆在他们面前。因此保禄在此也这样做，同时提醒他们那罪，并教导这恶的根源也是安逸和贪食。因此他又加上：\Quote{也不要行淫，像他们有些人行淫，一天就倒毙了二万三千人。}\par

他为何没有同样提到他们拜偶像的惩罚呢？要么是因为那惩罚显而易见，更广为人知；要么是因为当时的瘟疫不如巴郎事件那样大，当他们依附巴耳培敖尔，米德杨妇女出现在营中，按照巴郎的计谋引诱他们放荡的时候。因为摩西在《户籍纪》末尾的陈述中后来指出这邪恶的计谋是巴郎的。\BibleRef{户 31:8}\Quote{他们又在战争中用刀杀了米德杨的王巴郎……并掳掠了所有财物。……摩西就发怒，说：\InnerQuote{为什么留下这些妇女活着呢？这些妇女正是照巴郎的话，引诱以色列人触犯上主，在培敖尔的事上得罪上主。}}\par

% structure: p0047 source=h2 target=subsection reason=relative-heading-rank; base=h2

\subsection{格前 10:9}

\Quote{也不要试探基督，像他们有些人试探过，就被蛇所灭。}\par

由此他再次暗示另一项指控，这他也在最后陈述了，责备他们因为争论迹象而犯罪。事实上，他们因试探而灭亡，说：\Quote{好事何时到来？赏报何时到来？}因此他又加上，为纠正并警戒他们：\par

% structure: p0050 source=h2 target=subsection reason=relative-heading-rank; base=h2

\subsection{格前 10:10}

\Quote{也不要发怨言，像他们有人发怨言，就被灭命的所灭。}\par

因为所需要的不仅是为基督受苦，而且要高尚地、并满怀喜乐地承受临到我们的事：因为这是每一顶冠冕的本性。诚然，若非如此，那些以不情愿的态度承受灾祸的人，反要受到惩罚。因此，宗徒们被打时也欢喜，保禄更以他的受苦为荣。\par

% structure: p0053 source=h2 target=subsection reason=relative-heading-rank; base=h2

\subsection{格前 10:11}

\BibleQuote{所有这些事发生在他们身上，作为预像；它们写了下来，为劝戒我们这些世代已到末期的人。}\par

他又以提及\Quote{终结}来恐吓他们，并预备他们期待比已发生之事更大的事。他说：\Quote{因为我们必将受到惩罚，这是显而易见的，即使对那些不相信有关地狱之火的陈述的人也是如此；但惩罚也将是最严厉的，这从我们所享有的更多祝福以及那些只是预像的事物中便可明见。因为，如果在恩赐上一样超越另一样，那么很明显，在惩罚上也会同样。}因此，他既称它们为预像，又说它们是\Quote{为我们写的}，并提及\Quote{终结}，以便提醒他们万物的结局。因为那时的惩罚并非那种可终止并消除的，而是永久的；因为正如今世的惩罚随现世生命结束而告终，来世的惩罚却永远存在。但当他说\Quote{时代的终结}时，无非是指那可怕的审判从此近在眼前。\par

% structure: p0056 source=h2 target=subsection reason=relative-heading-rank; base=h2

\subsection{格前 10:12}

\BibleQuote{因此，那自以为站得稳的，要小心，免得跌倒。}\par

再者，他打击那些自恃知识高超之人的骄傲。因为，如果他们享有如此大的特惠却遭受了这些事；有人仅因抱怨就受到如此惩罚，有人因试探，无论他们的数量如何众多，还是他们获得如此大的恩惠，都不能使天主回心转意；那么，在我们身上，若我们不谨慎，就更当如此。他说得好：\Quote{那自以为站得稳的}，因为依靠自己甚至算不得站得稳；这样的人很快就会跌倒。因为他们若不是心高气傲、自信满满，而是谦卑自下，就不会遭受这些事。由此可知，主要是骄傲和漫不经心（贪饕也由此而来）是这些邪恶的根源。因此，即使你站得稳，也要小心免得跌倒。因为我们在此世的站立并不稳固，直到我们脱离今生的风浪，驶入平静的港湾。所以，不要因你站得稳而自高自大，却要提防跌倒；因为如果\Person{保禄}（他比任何人都坚定）尚且惧怕，我们更当惧怕。\par

现在，正如我们所见，宗徒的话是：\BibleQuote{因此，那自以为站得稳的，要小心，免得跌倒。}但我们甚至不能说这话；可以说，我们所有人都跌倒了，卧倒在地。因为我能对谁说这话呢？对那每天勒索的人？不，他跌得极重。对那行奸淫的人？他也倒在地上。对那醉汉？他也跌倒了，甚至不知自己跌倒了。所以，现在不是这话的时候，而是先知对犹太人所说的那句话：\BibleRef{耶 8:4}——\BibleQuote{跌倒的，难道不能起来吗？}因为所有人都跌倒了，却无意起来。所以我们的劝勉不是关于不跌倒，而是关于跌倒者有能力起来。那么，让我们起来吧，虽已迟晚，亲爱的，让我们起来，让我们高尚地站立。我们要躺卧多久？我们要醉酒多久？被今世事物的过度欲望迷醉多久？现在正是说这话的时候：\BibleRef{耶 6:10}——\BibleQuote{我向谁说话，向谁作证呢？}所有人都聋了，连对德行的教导也听不见，因而充满种种邪恶。若能看见他们赤裸的灵魂，就会像在战场上，战争结束后，可以看到一些死者，一些伤者；同样在教会里我们也能看到。因此，我恳求并哀求你们，让我们彼此伸出手来，彻底振作起来。因为我自己也是受伤者之一，需要有人来敷一些药。\par

不过，不要因此绝望。因为即使伤口严重，难道它们就不能医治吗？我们的医生就是这样：只要我们感觉到自己的伤口。即使我们达到了邪恶的极致，祂仍为我们开辟了许多安全的道路。这样，如果你不向邻人发怒，你自己的罪就会得到赦免。因为祂说：\BibleQuote{因为如果你们宽恕别人，}祂说，\BibleQuote{你们天上的父也要宽恕你们。}\BibleRef{玛 6:14}如果你施舍，祂会赦免你的罪；因为祂说：\BibleQuote{藉施舍断绝你的罪过。}\BibleRef{达 4:24}如果你恳切祈祷，你必获得宽恕：那寡妇凭着她祈祷的恒心打动了那位残酷的判官，正象征了这一点。如果你控告自己的罪，你便得解脱；因为\BibleQuote{你先宣告你的过犯，好使你成义。}\BibleRef{依 47:26}如果你为此而忧伤，这也会成为你有效的良药：因为祂说：\BibleQuote{我看见他忧伤痛苦，我便医治了他的道路。}\BibleRef{依 57:17}而且，当你遭受任何不幸时，如果你高尚地忍受，你便除去了全部。因为\Person{亚巴郎}也对富翁说过，\Quote{拉匝禄曾享过福，现在他在这里得安慰。}如果你们怜悯寡妇，你们的罪就被洗净了。因为祂说：\BibleQuote{审判孤儿，为寡妇辩护，然后来，我们彼此辩论——上主说。你们的罪虽似朱红，将变成雪一样的洁白；虽红得发紫，仍能变成羊毛一样的皎洁。}\BibleRef{依 1:17}因为祂不容许伤口留下任何疤痕。是的，即使我们陷入那种极深的痛苦，就像那耗尽父亲家产、靠豆荚充饥的人，如果我们悔改，我们无疑会得救。即使我们欠了一万塔冷通，如果我们俯伏在天主面前，不怀怨恨，一切都会被宽恕。即使我们迷失到牧人丢失羊的地方，祂也能从那里把我们找回；只要我们愿意，亲爱的。因为天主是慈悲的。因此，对于那欠一万塔冷通的人，祂只满足于他俯伏在祂面前；对于那耗尽父亲家产的人，只满足于他的回头；对于那只羊，只满足于它愿意被背负。\par

7. 因此，考虑到祂仁慈的伟大，让我们在此使祂对我们仁慈，并\BibleQuote{以充分的告解来到祂面前，}\BibleRef{咏 94:2}以免我们毫无借口地离去，而遭受极端的惩罚。因为如果在今生我们表现出哪怕普通的勤勉，我们也将获得最大的赏报；但如果我们在离去时在此地没有变得更好，即使我们在那里极其诚恳地悔改，也毫无益处。因为我们的责任是在仍留在赛场内时奋斗，而不是在集会散后徒然哀叹哭泣：正如那个富翁所做的，他悲叹痛惜，却毫无用处，因为他在应该做这些事的时候忽略了时机。不单是他，现在许多富翁都像他一样：不愿意轻视财富，却为了财富而轻视自己的灵魂：我不禁对此感到惊奇，因为我看到人们不断向天主祈求怜悯，同时却给自己造成不可救药的伤害，毫不爱惜自己的灵魂，视之如仇敌。因此，亲爱的，我们不要戏弄，不要戏弄也不要自欺，恳求天主怜悯我们，而我们自己却偏爱金钱和奢侈，事实上，把一切都置于这怜悯之上。因为如果有人在你面前提出一个案件，指控某人：他本该受千万次死亡，却有能力用一点点钱免除判决，但他宁愿死也不愿放弃任何财产，你会说他值得任何怜悯或同情吗？现在，你也要以同样的方式推断自己。因为我们也这样做，看轻自己的救恩，却吝惜金钱。那么，当你自己如此不爱惜自己，尊崇金钱胜过你的灵魂时，你怎么能恳求天主爱惜你呢？\par

因此我也大为惊奇，看到在财富中隐藏着多么大的迷惑力，或者不如说不在财富中，而在那些被迷惑的灵魂中。因为有些人，确实有些人彻底嘲笑了这种巫术。因为其中哪一样东西真的能够迷惑我们呢？难道不是无生命的物质？难道不是短暂的？难道它的拥有不值信赖？难道它不充满恐惧和危险？是的，还有凶杀和阴谋？仇恨和敌意？疏忽和许多恶习？难道不是尘土和灰烬？我们这是什么样的疯狂？什么样的疾病？\par

\Quote{但是，}你说，\Quote{我们不仅应该对染上这种病的人提出指责，还应该摧毁这种激情。}此外，除了指出它的卑劣以及它如何充满无数邪恶之外，我们还能用什么方式来摧毁它呢？\par

但若要说服一个恋慕其爱情对象的人，并非易事。那么，我们必须向他展示另一种美。然而，他因仍在病中，看不见无形体的美。那么，让我们向他展示某种有形体的美，并对他说：请看草地和其中的花朵，它们比任何黄金更璀璨，比各种宝石更优雅、更透明。请看从源头流出的清澈溪流，那些如油般无声地涌出地面的河流。升到天上，观看太阳的光辉、月亮的美貌、繁星如花朵般簇拥。\Quote{啊，这是什么？}你说，\Quote{因为我们大概并不像使用财富那样使用它们？}不，我们使用它们多于财富，因为它们的用途更为必要，享受更为稳妥。因为你不会害怕，如怕金钱一样，有人会将它们取走而离开；但你总能安心地拥有它们，且无需焦虑或操心。但若你因与他人共享而不快，且不像金钱那样独占而感到悲伤；那并非金钱，而是纯粹的贪婪，在我看来，你所恋慕的正是这贪婪；即使金钱，如果它被置于众人共享的范围内，也不会成为你欲望的对象。\par

8. 因此，既然我们找到了那被爱的对象——我指的是贪婪——，来吧，让我向你展示她如何恨你、厌恶你，她磨利了多少剑刃指向你，挖了多少陷阱，系了多少套索，预备了多少悬崖；这样至少你可以解除这魅惑。我们从哪里获得这认识呢？从大道，从战争，从海洋，从法庭。因为她使海洋充满鲜血，使法官的剑常常违背法律而染红，使那些昼夜埋伏在大道上的人武装起来，并劝说人忘记本性，制造弑父者和弑母者，将各种邪恶引入人的生命。这就是为什么\Person{保禄}称她为\Quote{万恶之根}。\BibleRef{弟前 6:10}她不让她的恋慕者比那些在矿场中劳作的人处境更好。因为正如他们永远被关在黑暗和锁链中，徒劳地劳作；同样，这些人埋葬在贪婪的洞穴中，无人强迫他们，却自愿承受惩罚，给自己绑上无法挣断的锁链。那些被判在矿场劳作的人，至少当夜晚来临时，可从苦工中解脱；但这些人无论昼夜都在这些悲惨的矿场中挖掘。对于那些人，那艰苦的劳作有明确的限度，但这些人的劳作不知限度，他们越挖就越渴望更大的艰辛。那些人是被迫的，而这些人却是自愿的，这又如何？你告诉我的是这疾病的可怕之处：他们甚至无法摆脱它，因为他们并不憎恨自己的悲惨。但正如猪在泥中，这些人也喜欢在贪婪的恶臭泥沼中打滚，遭受比那些被判刑者更糟的境遇。至于他们处境更糟的事实，请听一听一方的情形，然后你就会知道另一方的状况。\par

据说那富含金矿的土壤，在那些阴暗的洞穴中有某些裂缝和凹陷。被判在该处劳作的犯人，带着一盏灯和一把镐，这样，我们听说，他进入其中，并带一个油壶，以便从中滴油到灯里，因为那里即使白天也黑暗，没有一丝光线，如我之前所说。然后，当白天的时间召唤他吃那可怜的一餐时，他自己，据说，不知道时间，但他的狱卒从上方猛烈敲击洞穴，用那敲击声向下面劳作的人宣告一天的结束。\par

你们听到这些时，难道不战栗吗？现在让我们看看，在贪婪者身上是否没有比这些更可怕的事。首先，这些人也有一个更严厉的狱卒，即贪婪，而且严厉得多：不仅捆绑他们的身体，还捆绑他们的灵魂。这黑暗也比那更可怕。因为它不受感官的影响，而他们从内心产生它，无论走到哪里，都随身携带。因为他们灵魂的眼睛被弄瞎了：正因如此，\Person{基督}特别称他们为有祸的，说：\Quote{但若你内的光成了黑暗，那黑暗是何等大呢！}\BibleRef{玛 6:23}而那些人至少有一盏灯照亮，但这些人甚至被剥夺了这一线光；因此他们每天落入无数的陷阱。被判刑的人在黑夜来临时得到喘息，驶入那个所有不幸者共有的平静港湾，即夜晚；但对贪婪者来说，就连这个港湾也被他们自己的贪婪堵塞了：因为即使在夜晚，他们也有如此痛苦的念头，那时无人打扰，他们完全闲暇地自相残杀。\par

他们在今世的情形如此；但来世的情形，什么言词能展现呢？那无法忍受的火炉、燃烧着火的河流、咬牙切齿、永不解开的锁链、有毒的虫豸、无光的幽暗、无尽的痛苦。我们当畏惧这些，亲爱的，当畏惧如此大惩罚的源头、那疯狂的贪得无厌、我们救恩的毁灭者。因为不可能同时爱金钱又爱你的灵魂。让我们深信财富是尘土和灰烬，它在我们离开这里时就离开我们，或者不如说，在我们离开之前，它就常常从我们身边飞走，并在未来和现世两方面伤害我们。因为在地狱之火和那惩罚之前，甚至在这里，它就用无数的战争包围我们，激起纷争和争斗。没有什么比贪财更容易引起战争，没有什么比它更容易产生乞讨，无论它表现在财富中还是在贫穷中。因为在穷人的灵魂里，这严重的疾病也会出现，并且更加剧了他们的贫穷。如果有一个贫穷的贪婪者，这样的人所遭受的惩罚不在于金钱，而在于饥饿。因为他不能使自己舒心地享用那微薄的收入，反而用饥饿折磨自己的肚腹，用赤身和寒冷惩罚自己的全身，到处都显得比任何囚犯更褴褛和肮脏；他总是哀哭悲叹，仿佛他比所有人更不幸，虽有上万的人比他更贫穷。这人若进入市场，便带着无数鞭伤离开；若进入浴场或剧院，他仍会遭受更多的伤害，不仅来自观众，也来自舞台上的那些人，他在那里看见不少不贞的妇女以黄金闪耀。这人若航行海上，看着商人和他们满载货物的船只及巨大的利润，他甚至不认为自己还活着；或若在陆地上旅行，计算着田地、城郊的庄园、客栈、浴场、从中所得的收益，便从此认为自己的生命不值得活；或若你把他关在家里，他只会摩擦和磨损在市场上所受的伤口，从而更虐待自己的灵魂：他只知道一种缓解压迫他的痛苦的方法：死亡和脱离此生。\par

这些事不仅穷人会遭受，富人也一样，只要他染上这疾病，而且比穷人更甚，因为暴政在他身上压得更猛烈，醉意也更浓。因此他也将认为自己比所有人更贫穷；或者不如说，他确实是更贫穷的。因为财富和贫穷不是由物质的多少来决定的，而是由心灵的倾向决定的；那总是贪求更多、永远不能抑制这邪恶欲望的人，才是最贫穷的。\par

因此，为所有这些缘故，让我们逃避贪财——它是乞儿的制造者、灵魂的毁灭者、地狱的朋友、天国的敌人、一切邪恶之母；让我们藐视财富，好能享受财富，并随同财富也享受为我们预备的福乐；愿我们都能达到那福乐，等等。\par

\SourceRecord{Translated by Talbot W. Chambers.；Nicene and Post-Nicene Fathers, First Series；Vol. 12.；Edited by Philip Schaff.；Buffalo, NY: Christian Literature Publishing Co.,；1889.；<http://www.newadvent.org/fathers/220123.htm>.}

% structure: p0001 source=h1 target=section reason=page-title

\section{《格前》第二十四篇讲道}

% structure: p0002 source=h2 target=subsection reason=relative-heading-rank; base=h2

\subsection{格前10:13}

\BibleQuote{你們所受的試探，無非是普通人所能承受的；天主是忠信的，他決不許你們受試探超過你們的能力，而且在受試探的時候，他必也開一條出路，叫你們能夠擔當。}\par

這樣，因為他用古時的例證大大恐嚇了他們，使他們陷入痛苦，說：\BibleQuote{所以，自己以為站得穩的，要謹慎，免得跌倒。}雖然他們曾受過許多試探，並在其中多次磨練自己；因為他說：\BibleQuote{我從前在你們那裡，又軟弱，又恐懼，又戰戰兢兢。}\BibleRef{格前 2:3} 恐怕他們說：\Quote{你為什麼恐嚇我們、驚嚇我們？我們在這些患難中並非沒有操練，因為我們曾被驅逐、被迫害，經歷了許多持續的危險。}他再次壓制他們的驕傲，說：\BibleQuote{你們所受的試探，無非是普通人所能承受的。}意思是微小、短暫、適度的。因為他用了\Quote{普通人所能承受的}來表示微小的事，就如他所說：\BibleQuote{我因你們肉體的軟弱，照人的常態說話。}\BibleRef{羅 6:19} \Quote{所以，不要想得太高，好像你們已經勝過風暴一樣。因為你們從未見過死亡的危險，也未見過殺戮的試探。}這話他也對希伯來人說過：\BibleQuote{你們還沒有抵抗到流血的地步，與罪惡爭鬥。}\BibleRef{希 12:4}\par

然後，因為他恐嚇了他們，請看他如何又提振他們，同時推薦節制，說：\BibleQuote{天主是忠信的，他決不許你們受試探超過你們的能力。}所以，有些試探是我們不能承受的。這些是什麼呢？可以說，全都是。因為能力在於天主的恩寵影響；這是我們靠自己的意志所獲取的力量。因此，為了讓你們知道並看到，不僅那些超過我們能力的試探，就連這些\Quote{普通人所共有的}試探，若沒有來自天主的幫助，也不可能輕易承受，他補充說：\par

\BibleQuote{而且在受試探的時候，他必也開一條出路，叫你們能夠擔當。}\par

因為他說，正如我所說的，就連那些適度的試探，我們也不能靠自己的力量承受；即使在這些試探中，我們也需要他的幫助，在奮戰中得勝，直到我們通過它們並承擔它們。因為他賜予忍耐，並帶來迅速的解放；這樣，試探也變得可以承受。他隱晦地暗示這點，說：\BibleQuote{也開一條出路，叫你們能夠擔當。}並將一切歸於他。\par

% structure: p0008 source=h2 target=subsection reason=relative-heading-rank; base=h2

\subsection{格前10:14}

2. \BibleQuote{所以，我親愛的諸位，你們要逃避崇拜偶像的事。}\par

他又用\Quote{諸位}這個親屬的稱呼來討好他們，敦促他們趕快擺脫這罪。因為他沒有簡單地說\Quote{離開}，而是說\Quote{逃避}；他稱這事情為\Quote{崇拜偶像}，不再僅僅因為對近人的傷害而叫他們離棄它，而是表明這事情本身足以帶來極大的毀滅。\par

% structure: p0011 source=h2 target=subsection reason=relative-heading-rank; base=h2

\subsection{格前10:15}

\BibleQuote{我好像對明白人說話；你們自己審斷我所說的吧。}\par

因為他大聲呼喊，加重了指控，稱之為崇拜偶像；為了不使他們激怒，令他的言論令人厭惡，在接下來的內容中，他將決定權交給他們，並以讚詞將他的審判者安置在他們的裁判席上。他說：\Quote{我好像對明白人說話。}這是一個對自己的立場非常有信心的人的特徵：他讓被告自己判斷他的論點。\par

這樣，當他說話不是命令，也不是定規，而是與他們商量、實際上在他們面前辯護時，他更能提升聽眾。因為對猶太人，他們比較愚昧幼稚，天主沒有這樣說話，也沒有每次告知他們命令的理由，而只是命令；但在這裡，由於我們擁有極大的自由，甚至被允許成為參謀。他如同對朋友說話，說：\Quote{我不需要其他審判者，你們自己對我做出判斷，我以你們為仲裁者。}\par

% structure: p0015 source=h2 target=subsection reason=relative-heading-rank; base=h2

\subsection{格前10:16}

3. \BibleQuote{我們所祝福的祝福之杯，豈不是共結合於基督的血嗎？}\par

有福的保祿，你說什麼？當你要呼喚聽眾的敬畏，當你提到可畏的奧蹟時，你竟將那可怕而極其驚恐的杯稱為\Quote{祝福之杯}嗎？他說：\Quote{是的；而且這話並非平凡的稱呼。因為當我稱之為\InnerQuote{祝福}時，我指的是感恩；當我稱之為感恩時，我揭示了天主所有恩惠的寶藏，並記起了那些偉大的恩賜。}因為我們在杯中數算天主無可言喻的慈悲，以及我們所分享的一切，如此走近他並共融；感謝他將整個人類從錯誤中解救出來；使我們這些遠離的人得以靠近；使我們這些在世上沒有盼望、沒有天主的人，成為他自己的兄弟和共同繼承人。為這一切和所有類似的事，我們感謝著，如此走近。\Quote{那麼}，他說，\Quote{你們格林多人啊，你們的行為豈不是不一致嗎？你們感謝天主將你們從偶像中解救出來，卻又跑回他們的筵席？}\par

\Quote{我们所祝福的祝福之杯，岂不是基督之血的共融吗？} 他讲得极具说服力，且令人敬畏。因为他所说的是：\Quote{杯中所盛的，是从祂肋旁流出的，我们正分受那血。} 但他称其为祝福之杯，因为我们手中高举它，在赞歌中如此举扬祂，惊叹于祂难以言喻的恩赐，祝福祂，特别是为倾倒这同一饮剂，使我们不留在错谬中；不仅为倾倒，也为将它分给我们众人。\Quote{所以，如果你渴望血，} 祂说，\Quote{不要用牲畜的宰杀染红偶像的祭台，而是用我的血染红我的祭台。} 告诉我，有什么比这更可畏？有什么比这更温柔慈爱？恋人们也有这样的行为。当他们看到自己所爱的人渴望属于外人的东西而藐视自己的东西时，他们就把自己所有的给他们，以此劝服他们远离别人的礼物。然而，恋人们在财物、金钱和衣裳上显示这种慷慨，但从未有人以血如此行。而基督就在这件事上显示了对我们的关怀和炽热的爱。在旧盟约中，因为他们处于不完美的状态，他们通常用来祭偶像的血，祂自己竟接受，以便将他们与那些偶像分开；这本身也是祂不可言喻之爱的明证：但在这里，祂将敬礼转移到更可畏、更荣耀的事上，改变了祭献本身，命令献上自己，而不是宰杀无理性的受造物。\par

4. \Quote{我们所擘开的饼，岂不是基督身体的共融吗？} 他为何不说\Quote{分受}？因为他想表达更多，并指出这种结合何等密切：我们不仅通过参与和分受而共融，也因结合而共融。因为正如那身体与基督结合，我们也藉着这饼与祂结合。\par

但他为何又加上\Quote{我们所擘开的？} 因为在圣体圣事中我们确实看到这是可行的，但在十字架上却不是如此，而是恰恰相反。因为有人说：\Quote{祂的骨头，一根也不可折断。} 但祂在十字架上所受的（不被折断），在祭献中为你们而受，并甘愿被擘开，为要充满众人。\par

此外，因为他说了\Quote{身体的共融}，而共融者与所共融的对象本是不同的；就连这看似微小的差别，他也消除了。因为说了\Quote{身体的共融}之后，他又试图表达更亲密的关系。因此他又加上：\par

% structure: p0022 source=h2 target=subsection reason=relative-heading-rank; base=h2

\subsection{\BibleBook{格林多}{前书} 10:17}

\Quote{我们众人都是一饼，一个身体。} \Quote{为什么要说共融？} 他说，\Quote{我们就是那个同一的身体。} 因为饼是什么？基督的身体。而分受它的人成为什么：基督的身体：不是许多身体，而是一个身体。正如饼由许多麦粒合成一，以致麦粒无处可见；它们确实存在，但由于结合，它们的差异看不见了；同样，我们彼此结合并与基督结合：没有为你预备一个身体，为邻人预备另一个，而是众人共享同一个。因此他又加上：\par

\Quote{因为我们众人都分受同一饼。} 现在，如果我们都以同样的食物滋养，都成为同样的，为什么我们不也显明同样的爱，并在这点上成为一呢？因为在我们祖先的时代也有此风：圣经记述：\Quote{那众多信众，} 经文说，\Quote{都一心一意。} \BibleRef{宗 4:32} 然而，如今并非如此，而是全然相反。众人之间多有纷争，我们对待彼此的肢体比野兽还凶残。基督早已使你远离至此，使你与祂合而为一；但你却不屑于以应有的精确与你的兄弟联合，反而自相分离，你既然从主领受了如此伟大的爱与生命。因为祂不仅简单赐下祂自己的身体；而且因为那从土而造的肉身本性，先已被罪麻木且失去生命；祂引进，可谓另一种面团和酵母，祂自己的肉身，本性相同，但无罪恶且充满生命；并赐给众人分受，使我们藉此滋养，脱去那旧的死物质，藉着这筵席，融合于那永生不朽的。\par

% structure: p0025 source=h2 target=subsection reason=relative-heading-rank; base=h2

\subsection{\BibleBook{格林多}{前书} 10:18}

5. \Quote{你们看那按血肉的以色列人：那些吃祭物的，岂不是与祭台共融吗？}\par

他又从旧盟约引导他们至此。因为他们远不及所言之事的伟大，他既以过往的事，又以他们习以为常的事来说服他们。他恰当地说\Quote{按血肉的}，仿佛他们自己按圣神行事。他所说的大意是：\Quote{甚至从更粗俗之人，你们可得知，那些吃祭物的，是与祭台共融的。} 你看见他如何暗示，那些看似完美的人并无完美的知识，如果他们连这点都不知道：这些祭献的结果往往是与魔鬼的某种共融和友谊，而实践逐渐牵引他们？因为如果在人中间，盐和餐桌的共食成为友谊的契机和标记，那么在魔鬼身上也可能发生同样的事。\par

但请你考虑，论到犹太人，他并没有说\Quote{他们与天主分受}，而是说\Quote{他们与祭台共融}；因为摆在祭台上的被焚烧了；但论到基督的身体，却不是这样。而是如何？乃是\Quote{主的身体的共融}。因为我们不是与祭台共融，而是与基督自己共融。\par

但他说了他们与祭台共融之后，后来恐怕他像是在说偶像有任何能力并造成伤害，请看他又如何推翻他们，说道：\par

% structure: p0030 source=h2 target=subsection reason=relative-heading-rank; base=h2

\subsection{\BibleRef{格前 10:19}}

\BibleQuote{那么我说什么呢？难道说偶像算得什么？或者祭偶像之物算得什么？}\par

\Quote{就好像他说：\InnerQuote{现在我所肯定的这些事，就是设法使你们远离偶像，并非因为它们能造成什么伤害或有什么力量：因为偶像算不了什么；但我希望你们藐视它们。}} \Quote{既然你要我们藐视它们，}有人说，\Quote{那你为什么仔细地使我们远离它们？}因为它们不是奉献给你们的主的。\par

% structure: p0033 source=h2 target=subsection reason=relative-heading-rank; base=h2

\subsection{\BibleRef{格前 10:20}}

\BibleQuote{因为外邦人所祭献的，是祭献给魔鬼，而不是祭献给天主。}\par

那么，不要跑到相反的事情上去。因为假如你是国王的儿子，享有父王的餐桌特权，却离开它，选择去分享地牢中那些被定罪者和囚犯的餐桌，你的父亲不会允许，反而会极其强烈地阻止你；并非因为那张餐桌能伤害你，而是因为它羞辱了你的尊贵和皇家餐桌。因为这些人确实也是冒犯的仆人；受羞辱、被定罪、被囚禁、等待不可忍受的惩罚、犯下万千罪行。那么，你怎么还不羞于模仿那些贪吃粗俗的人，当这些被定罪者摆出餐桌时，你跑到那里去享用食物呢？这就是我试图使你们远离的原因。因为献祭者的意图和接受者的位格，使摆在你们面前的东西成为不洁的。\par

\BibleQuote{我不愿你们与魔鬼有来往。}你看出一位细心父亲的慈爱吗？你也看出这话本身表达他感情的力量吗？\Quote{因为我的愿望是你们与他们毫无共同之处。}\par

然后，因为他借着劝诫的方式说出这话，免得比较粗鲁的人因着他说了\Quote{我不愿意}和\Quote{你们判断}而轻视它，以为享有自由；他在后面明确断言并立下规矩，说：\par

% structure: p0038 source=h2 target=subsection reason=relative-heading-rank; base=h2

\subsection{\BibleRef{格前 10:21}}

\BibleQuote{你们不能喝主的杯，又喝魔鬼的杯；你们不能参加主的筵席，又参加魔鬼的筵席。}\par

他只满足于用这些词语，为了使他们远离。然后，也针对他们的羞耻感说：\par

% structure: p0041 source=h2 target=subsection reason=relative-heading-rank; base=h2

\subsection{\BibleRef{格前 10:22}}

\BibleQuote{我们惹主嫉妒吗？我们比祂更强吗？}也就是说，\Quote{我们是在试探祂，看祂能否惩罚我们，并且因投奔敌对者、站在祂仇敌一边而激怒祂吗？}他说这话，是要提醒他们一段古代历史和祖先的过犯。因此他也使用了这个说法，就是\Person{梅瑟}昔日同样用来指责犹太人、藉天主的口斥责他们拜偶像的话：\BibleQuote{他们以不是神的，惹动我的嫉妒；以他们的偶像，激怒我的怒气。}\BibleRef{申 32:21}\par

\BibleQuote{我们比祂更强吗？}你们看见他是多么严厉、多么可怕地责备他们，彻底震动他们的神经，并以归谬法触及其痛处、挫其傲气吗？\Quote{哦，但为什么，}有人会说，\Quote{他一开始不把这些最有效使他们远离的事写出来呢？}因为他的习惯是用许多细节来证明他的论点，把最强的放在最后，并通过证明超乎必要来取胜。因此，他从较小的主题开始，这样一路走到万恶的总汇：因为这样最后一点也更容易被接受，他们的心已被之前所说的话所软化。\par

% structure: p0044 source=h2 target=subsection reason=relative-heading-rank; base=h2

\subsection{\BibleRef{格前 10:23-24}}

\BibleQuote{凡事我都可行，但不都有益处；凡事我都可行，但不都造就人。各人不要寻求自己的事，但各人要寻求邻人的好处。}\par

你看见他那精确的智慧吗？因为他们很可能说：\Quote{我是完美的，能掌管自己，吃摆在面前的食物对我无害；} \Quote{即便如此，}他说，\Quote{你是完美的，能掌管自己；然而不要看这个，而要看结果是否涉及伤害，甚至是颠覆。}因为他提到这两点，说：\BibleQuote{凡事不都有益处，凡事不都造就人；}并且前者针对自己，后者针对弟兄：因为\Quote{不都有益处}这句话，暗中暗示了他所说话的人的毁灭；而\Quote{不造就人}这句话，则暗示了弟兄的绊脚石。\par

因此他也加上：\BibleQuote{不要寻求自己的事；}这在他整封书信中和在\WorkTitle{致罗马人书}中都一再强调；他说：\BibleQuote{因为连基督也没有寻求自己的喜悦：}\BibleRef{罗 15:3}；又说：\BibleQuote{就像我在一切事上使众人喜悦，不求自己的益处。}\BibleRef{格前 10:33}。在这里他又一次这么做；不过他没有完全展开。也就是说，因为在前文中他已经详细确立了这一点，并表明他哪里也不\Quote{寻求自己的事}，反而\BibleQuote{对犹太人，就成为犹太人；对没有法律的人，就成为没有法律的人}，并且他没有随意使用自己的\Quote{自由}和\Quote{权利}，而是为了众人的益处，服事众人；他在这里就停住了，满足于几句话，借着这几句话引导他们回想所有已经说过的话。\par

7. 知道这一切后，亲爱的，让我们也顾念弟兄们的益处，与他们保持合一。因为那令人敬畏的、可怖的祭献正引导我们走向此事，警告我们尤其要同心合意、以炽热的爱德接近它，并由此成为鹰，如此翱翔直至高天，甚至超越诸天。因为祂说：\Quote{因为无论那里有尸体，鹰也必聚集在那里。}\BibleRef{玛 24:28}祂称自己的身体为尸体，是由于祂的死亡。除非祂倒下，我们不会复活。但祂称我们为鹰，暗示那接近这身体的人必须居于高处，与尘世无涉，不可向下蜷曲爬行；而应永远向天飞翔，仰望正义的太阳，并拥有敏锐的心灵之眼。因为鹰，而非寒鸦，才有权享用这筵席。那些如今相称地享有这一特权的人，将来也要迎接从天而降的祂；而那些不相称地如此行的人，将遭受最极端的折磨。\par

因为若有人不会轻率地接待一位君王——（我说君王何意？不，即便是王袍，也不可轻率地用污秽的手去触碰；）——即使他独处、孤单、无人近旁：然而那袍子不过是虫所吐的丝线；你若赞叹那染料，那也是死鱼的血；尽管如此，人也不愿用污染的手去冒险接触它。我说，若连一个人的衣服也不可轻率地触碰，那么对于那统御万有的天主的身体，无玷、纯净、与天主性结合的身体，我们藉以生存、生活的身体，地狱之门因它被打破、天堂圣所因它被打开的身体，我们又当如何呢？我们怎能如此无礼地领受它？我恳求你们，不要，不要因我们的不敬而自取灭亡，而要以全部的敬畏和纯洁接近它；当你看到它被摆在你面前时，你要对自己说：\Quote{因这身体，我不再是尘土和灰烬，不再是囚犯，而是自由；因这身体，我盼望天堂，并领受其中的美善，不朽的生命，天使的份位，与基督的交往；这被钉、被鞭打的身体，死亡无法抵挡；这身体被献祭时，太阳曾目睹并收回它的光芒；为此，幔子在那一刻撕裂，岩石崩裂，整个大地震动。这就是那身体，染血的、被刺透的，从中涌出救恩的泉源，一为血，一为水，为整个世界。}\par

你还想从另一个来源学习它的能力吗？问问那患血漏的妇人，她所抓的不是身体本身，而是包裹它的衣服；不，甚至不是整件衣服，只是衣边；问问那承载它的海；甚至问问魔鬼本身，说：\Quote{你那不可治愈的创伤从何而来？你为何不再有任何能力？你为何被俘虏？谁在逃亡中擒获了你？}他只会这样回答：\Quote{那被钉死的身体。}藉此，它的刺被折断；藉此，它的头被击碎；藉此，那些执政的、掌权的被公开示众。因为祂说：\Quote{祂解除了执政者和掌权者的武装，公开地羞辱他们，藉着十字架战胜了他们。}\BibleRef{哥 2:15}\par

也要问问死亡，说：\Quote{你的毒刺为何被除去？你的胜利为何被废除？你的筋腱为何被切断？你从前连君王和一切义人都惧怕，如今却成为孩童和少女的笑柄？}它会将这归功于这身体。因为这身体被钉时，死人复活，那监狱被冲破，铜门被打破，死人得释放，地狱门口的看守都惊惶失措。然而，若祂是众人之一，死亡反而会更强；但并非如此。因为祂不是众人之一。因此死亡被瓦解。正如那些吃了无法消化的食物的人，因反胃而把先前存留的也吐出；死亡也是如此。它吞下了那无法消化的身体，因此不得不吐出它里面所存的。是的，它怀胎痛苦，被拘束，直到将它呕出。故此宗徒说：\Quote{解除了死亡的痛苦。}\BibleRef{宗 11:24}因为从未有妇人分娩时像它在怀着主的身体时那样被撕裂和折磨。那发生在巴比伦龙身上的事，当它吞下食物后肚腹爆裂，也同样发生在它身上。因为基督不是从死亡的口中再次出来，而是像从内室中爆裂、撕开龙的肚腹，从祂的隐秘处\BibleRef{咏 18:5}极其荣耀地出来，将祂的光芒不仅洒向这天空，更洒向至高的宝座。因为祂甚至将它携带上去。\par

祂将这身体赐给我们，不仅要我们持守，更要我们吃；这是与炽热之爱相称的事。因为那些我们深吻的人，我们有时甚至会咬他们。因此约伯在表明仆人对他的爱时说，他们因对他的大爱常常说：\Quote{哦！愿我们饱尝他的肉！}\BibleRef{约 31:31}同样，基督赐我们饱尝祂的肉，引领我们走向更大的爱。\par

8. 让我们殷勤地、以炽热的爱走近祂，免得我们受到惩罚。因为所赐予我们的恩惠越大，当我们显示自己不配这慷慨时，我们受的责罚就越大。这圣体，即使躺在马槽里，贤士们也朝拜了。是的，那些亵渎和野蛮的人，离开他们的国家和家园，踏上长途旅程，来的时候，带着敬畏和巨大的战栗朝拜了祂。那么，我们这些天国的子民，至少该效仿那些外邦人。因为他们看见祂在马槽里、在茅屋里，没有你现在所见的任何事物，却带着极大的敬畏走近；但你看见祂不是在马槽里，而是在祭台上，不是一个女人抱着祂，而是司铎站在旁边，圣神极其慷慨地盘旋在我们面前的供品上。你不仅仅像他们那样看见这圣体本身，你还知道祂的德能，以及整个救恩计划，并且不忽视任何通过祂成就的神圣事物，因为你已完全领受了这一切。\par

因此，让我们振奋起来，充满敬畏，并显示出远超过那些外邦人的崇敬；免得我们因随意和粗心的接近，而在自己头上堆起烈火。但我说这些话，不是要阻止我们接近，而是阻止我们不加考虑地接近。因为随意接近是危险的，同样，不参与那些奥秘的晚餐就是饥荒和死亡。因为这祭台是我们灵魂的筋腱、思想的纽带、信心的根基、我们的希望、救恩、光明、生命。当我们带着这祭献离开进入外世时，我们将怀着极大的信心踏过神圣的门槛，四面都被一种金甲所环绕。\par

那么，我何必谈论将来世界？因为在这里，这奥秘使大地在你眼中成为天堂。你只需打开天堂的门看一眼；不，不是天堂，而是诸天之天；然后你就会看到我所说的一切。因为最珍贵的一切，我将指给你看它躺在大地上。因为在王宫中，最荣耀的不是墙壁，也不是黄金屋顶，而是坐在宝座上的君王本人；同样，在天上，就是君王的圣体。但现在，你被允许在地上看见这圣体。因为我所展现给你的，不是天使，也不是总领天使，也不是天和诸天之天，而是这一切的主与拥有者。你岂不明白，那比万物更珍贵的，你在地上就能看见；不仅看见，还能触摸；不仅触摸，还能吃；并且领受后你就回家？\par

那么，请洁净你的灵魂，预备你的心灵以领受这些奥秘。因为如果你受托抱着一个君王的儿子，穿着袍子、紫袍和冠冕，你会抛弃地上的一切。但现在你所领受的不是人的孩子，无论多么尊贵，而是天主自己的独生子，你难道不战栗，告诉我，并且抛弃了一切世俗之物的爱，而除了用以装饰自己的勇气之外，没有别的勇气吗？或者你仍注目于地上，爱钱财，渴求黄金？那么你还能得到什么宽恕？什么借口？你难道不知道，所有这些世俗的奢华在主眼中是可憎的吗？难道不是为了这个，祂在诞生时被放在马槽里，并使自己有一位卑微的母亲？难道不是为了这个，祂对那个追求利益的人说：\Quote{但是人子却没有枕头的地方？} \BibleRef{玛 8:20}\par

那么门徒们做了什么？他们不是遵守了同样的法则，被带到穷人的家里并留宿，一个与皮匠在一起，另一个与帐棚匠在一起，还有与卖紫布的人在一起？因为他们不追求房屋的华丽，而是追求人们灵魂的美德。\par

因此，我们也要效仿他们，匆匆走过柱子和大理石的美景，寻求天上的居所；让我们践踏一切地上的骄傲和一切对金钱的爱，并获得崇高的思想。因为如果我们清醒自持，连整个世界都不配我们，更不用说门廊和拱廊了。因此，我恳求你们，让我们装饰我们的灵魂，让我们整理这所房屋，就是我们在离开时也要带走的；这样我们或许能获得永恒的福乐，藉着恩宠和仁慈，等等。\par

\SourceRecord{Translated by Talbot W. Chambers.；Nicene and Post-Nicene Fathers, First Series；Vol. 12.；Edited by Philip Schaff.；Buffalo, NY: Christian Literature Publishing Co.,；1889.；<http://www.newadvent.org/fathers/220124.htm>.}

% structure: p0001 source=h1 target=section reason=page-title

\section{关于《格林多前书》的讲道第二十五篇}

% structure: p0002 source=h2 target=subsection reason=relative-heading-rank; base=h2

\subsection{格林多前书 10:25}

凡在肉市上售卖的，你们只管吃，不要为良心的缘故查问什么。\par

他说了\BibleQuote{你们不能喝主的杯，又喝魔鬼的杯}，并且藉犹太人的事例、人的理智、那可怕的奥迹\OriginalTerm{奥迹}{Mysteries}和偶像中所举行的礼仪，一次而永远地把他们从那宴席上引开，又使他们充满极大的恐惧；但为了不因这恐惧而把他们驱向另一个极端，叫他们被迫过度谨慎，担心可能连自己也不知道就从市场或其他地方沾上这类东西，为了使他们从这种困境中解脱，他说：\Quote{因为如果你们是在不知情而非故意的情况下吃，就不会受到惩罚：此后便不再是贪吃的问题，而是无知的问题了。}\par

他不仅让人摆脱这种忧虑，还从另一种忧虑中释放他们，使他们处于彻底的安全和自由之中。因为他甚至不让他们\BibleQuote{查问}，即搜索和探询那是否是祭过偶像的肉或并非如此；而只是简单地吃一切来自市场的东西，甚至不让自己知道摆在我们面前的是什么。这样，即使是吃的人，若处在无知中，也能免于焦虑。因为那些本质并非邪恶、而因人的意图使人不洁的事物，其本性就是如此。因此他说：\BibleQuote{不要查问。}\par

% structure: p0006 source=h2 target=subsection reason=relative-heading-rank; base=h2

\subsection{格林多前书 10:26}

\BibleQuote{因为大地和其中的万物属于主。}不属于魔鬼。如果大地、果实和牲畜都属于祂，那么就没有什么是不洁的；但它会因我们的意图和悖逆而以另一种方式成为不洁。因此他不仅允许，而且，\par

% structure: p0008 source=h2 target=subsection reason=relative-heading-rank; base=h2

\subsection{格林多前书 10:27}

\BibleQuote{若有不信的人请你们，赴宴，你们愿意去；凡摆在你们面前的，只管吃，不要为良心的缘故查问什么。}\par

请看他的节制。因为他没有命令或制定法律要他们退出，但也没有禁止。而且，如果他们离开，他就免去他们一切猜疑。这是什么原因呢？是为了使这样大的挑剔不显得出自恐惧和怯懦。因为谨慎查问的人是由于害怕；而听到事实后戒绝的人，是出于轻视、憎恨和厌恶。因此，保禄为了确立这两点，说：\BibleQuote{凡摆在你们面前的，只管吃。}\par

% structure: p0011 source=h2 target=subsection reason=relative-heading-rank; base=h2

\subsection{格林多前书 10:28-29}

\BibleQuote{但若有人对你们说：\InnerQuote{这是祭献给偶像的，}你们就不要吃，为了那指明这事的人的缘故。}\par

因此，他命令戒绝它们，完全不是因为它们有什么能力，而是因为它们是被咒诅的。所以，既不要像它们能伤害你那样逃避（因为它们毫无力量），也不要因为它们没有力量就漠然地参与：因为那是敌对和堕落之灵的宴席。因此他说：\BibleQuote{不要吃，为了那指明这事的人和良心的缘故。因为大地和其中的万物属于主。}\par

\Quote{我并非禁止，不是因为这些东西属于别人（因为大地属于主），而是出于我所说的理由，为了良心的缘故；}即，免得良心受伤。那么人应该谨慎查问吗？他说：\Quote{不，因为我说的不是你的良心，而是他的。因为我已经说过，\InnerQuote{为了那指明这事的人的缘故。}}接着又说：\BibleQuote{我说的是良心，不是你自己的，而是别人的。}\BibleRef{格前 10:29}\par

2. 但也许有人会说：\Quote{你固然出于自然照顾弟兄们，为了他们的缘故不许我们尝，免得他们软弱的良心被怂恿去吃祭偶像之物。但如果是外邦人，这与你有何相干？岂不正是你自己说的：\BibleQuote{关我何事，评断教外之人？}\BibleRef{格前 5:12} 那么你为什么反倒关心他们呢？}他回答说：\Quote{我不是关心他，而是在这种情况下也关心你们。}\par

\BibleQuote{为什么我的自由要受别人的良心论断？}这里\Quote{自由}\OriginalTerm{自由}{liberty}指的是那不加谨慎或禁止而留下的。因为这是自由，从犹太人的束缚中获得释放。他的意思是：\Quote{天主使我自由，完全不受伤害，但外邦人不知道如何判断我的生活准则，也看不到我主人的宽宏，反而会谴责并对自己说：\InnerQuote{基督教是神话；他们戒绝偶像，逃避魔鬼，却又依恋献给他们的事物：他们的贪食多么严重。}}\Quote{那又怎样？}也许有人说。\Quote{他们不公正地论断我们，对我们有何损害？}但完全不给他们论断的余地岂不更好！因为如果你戒绝，他们甚至不会说这话。\Quote{怎么？}你说，\Quote{他们不会说吗？因为当他们看到我在市场或宴席上不作这些查问时，什么能阻止他们这样说，责备我不加区分地吃呢？}完全不是这样。因为你吃，不是像吃祭偶像之物，而是像吃洁净之物。而你不作仔细查问，是为了表示你不害怕摆在面前的东西；这就是为什么无论是进入外邦人的家还是去市场，我都不让你们查问；免得你们变得胆怯和困惑，给自己招来不必要的麻烦。\par

% structure: p0017 source=h2 target=subsection reason=relative-heading-rank; base=h2

\subsection{格林多前书 10:30}

\Quote{若我赖恩宠而领受，为何我所感谢的倒被人毁谤呢？} \Quote{你\Quote{赖恩宠领受}的是什么，请告诉我。} 天主的恩赐。因祂的恩宠如此伟大，竟使我的灵魂纯洁无瑕，超越一切玷污。犹如太阳将光芒照射在许多污秽之处，又纯洁地收回；同样，我们生活在世界中，若愿意，也能保持纯洁，因为我们的能力甚至比太阳更大。\Quote{那么为何要禁戒呢？} 你说。并非因为我会成为不洁，绝不然；而是为了我的弟兄，并且为了不让我与魔鬼有份，也为了不被不信者论断。因为在此情况下，不再是物本身的性质，而是不顺从以及与魔鬼的友谊使我成为不洁，心意的倾向造成了玷污。\par

但所谓\Quote{为何我所感谢的倒被人毁谤呢？}是什么意思？\Quote{我，就我而言，} 他说，\Quote{感谢天主，因祂如此高举我，使我超越犹太人卑微的境地，以致丝毫不受损害。但外邦人不了解我高尚的生活准则，会怀疑相反的事，他们说：\InnerQuote{基督徒竟喜好我们的习俗；他们是一类伪善者，既辱骂并厌弃魔鬼，却又奔向他们的筵席，还有什么比这更愚蠢的呢？我们断定，他们并非为真理，而是出于野心和权力欲才皈依这教义。}那么，对于我因之如此获益，甚至庄重感谢的事，竟成为毁谤的缘由，这是何等愚昧呢？} \Quote{然而，即使现在，} 你说，\Quote{当外邦人见我不加查问时，也会如此猜测。} 绝非如此。因为并非所有事物都充满祭偶像之物，以致他会如此揣测；你也不会把它们当作祭偶像之物来吃。但既不要过分挑剔，另一方面，当有人说那是祭偶像之物时，你也不可吃。因为基督赐给你恩宠，使你超越那方面的一切损害，不是要你被人毁谤，也不是要让那对你如此有益、甚至成为特别感恩之事的境况，反而伤害他人，使他们亵渎。\Quote{那么，为何我不可对外邦人说：\InnerQuote{我吃，我毫不受损，我这样做并非与魔鬼友好}呢？} 因为你无法说服他，即使你说千万次也不行：他既软弱又有敌意。如果你的弟兄尚未被你说服，何况敌人和外邦人呢？若他被祭偶像的意识所控制，何况不信者呢？此外，我们何需如此大的麻烦？\par

\Quote{那么，我们既认识基督并感谢，而他们亵渎，我们因此也要放弃这习俗吗？} 绝不。因为事情并不相同。在前者，忍受羞辱大有裨益；但在后者，却毫无益处。所以之前他也说：\Quote{因为无论我们吃，也不更好；不吃，也不更坏。} \BibleRef{格前 8:8} 此外，他也表明应当避免此事，所以不仅要为此缘故，也要为他所指出的其他理由而禁戒。\par

% structure: p0021 source=h2 target=subsection reason=relative-heading-rank; base=h2

\subsection{格前 10:31}

3. \Quote{所以，你们或吃或喝，或无论做什么，一切都要为光荣天主而做。}\par

你看出他如何从眼前的主题引出普遍的劝诫，给我们一个最卓越的目标，即在一切事上光荣天主吗？\par

% structure: p0024 source=h2 target=subsection reason=relative-heading-rank; base=h2

\subsection{格前 10:32}

\Quote{不可成为绊脚石，或为犹太人，或为希腊人，或为天主的教会：} 即不可给任何人留下把柄；因为在所假设的情况下，你的弟兄被绊倒，犹太人会更加憎恨并责难你，外邦人同样会嘲笑你，如同贪食者和伪善者。\par

然而，不仅弟兄不应因我们受到伤害，就连外人也当竭力不使他们受伤害。因为如果我们是\Quote{光}、\Quote{酵母}、\Quote{明灯}和\Quote{盐}，我们就该照亮，而非昏暗；该约束，而非放松；该吸引不信者归向我们，而非驱散他们。那么，你为何要赶走那些你本该吸引的人呢？因为连外邦人看到我们回归这类事，也会受伤害：他们不了解我们的心意，也不明白我们的灵魂已超乎一切感官的玷污。同样，犹太人和软弱的弟兄也会遭受同样的伤害。\par

你看到他提出了多少理由，说明我们应当禁戒祭偶像之物吗？因为它们的无益，因为它们的非必需，因为对弟兄的伤害，因为犹太人的毁谤，因为外邦人的辱骂，因为我们不应与魔鬼有份，因为这事是一种偶像崇拜。\par

此外，因为他曾说\Quote{不可成为绊脚石}，并且他使他们对外邦人和犹太人所受的伤害负有责任；这话本是严厉的；请看他是如何使之变得可接受和轻松，自己挺身而出，说道：\par

% structure: p0029 source=h2 target=subsection reason=relative-heading-rank; base=h2

\subsection{格前 10:33}

\Quote{就如我在一切事上使众人喜悦，不求自己的利益，只求众人的利益，为使他们得救。}\par

% structure: p0031 source=h2 target=subsection reason=relative-heading-rank; base=h2

\subsection{格前 11:1}

\Quote{你们该效法我，如同我效法基督一样。}\par

这是最成全的基督信仰的一条规则，这是精确划定的界标，这是所有标准中的最高点：即寻求那些为公共利益的事。保禄自己也说明了这一点，他加上说：\Quote{如同我属于基督一样。}因为没有什么比关心邻人更能使人成为基督的效仿者。即使你守斋、睡在地上、甚至折磨自己，却不顾念你的邻人，你也没有成就什么大事，你仍然远远地远离这肖像。然而，在目前的情况下，连那事本身自然也是有益的，即禁戒祭偶像之物。但他说：\Quote{我}也做了许多无益的事。例如：当我行割损礼时，当我献祭时；若有人单独审视这些，它们反而毁灭那些追随它们的人，使他们从救恩中坠落。然而，为了从中得益，我甚至忍受了这些。但这里没有这样的事。因为在那情况下，除非有某种益处，除非为他人而行，否则那事便成为有害的；但在这种情况下，即使没有人跌倒，人也应当禁戒那些被禁止的事。\par

我不但忍受了有害的事，也忍受了辛苦的事。因为他说：\BibleQuote{我抢夺了别的教会，向他们取了工资；并且虽然可以吃饭而不工作，我却不寻求这事，而宁愿饿死也不愿得罪别人。}\BibleRef{格后 11:8}因此他说：\Quote{我在一切事上使众人喜悦。}\Quote{即使要违反律法，即使要劳苦和冒险，为了他人的益处，我忍受一切。所以，既然他在成全上超乎众人，他在迁就上就成了低于众人。}\par

4. 因为没有一种德行行为，当它不将其利益分施给他人时，能是非常崇高的：正如那带来一塔冷通的人所显示的，因为他没有使之更多，就被砍碎了。所以，兄弟，即使你不进食，即使你睡在地上，即使你吃灰并不断哀哭，却不帮助别人，你也不会成就大事。因为那些起初的伟大和高贵的人也是这样以此为首要的关心：仔细查考他们的生活，你就会清楚看到，他们中没有一人寻求自己的事，而是每个人都寻求邻人的事，因此他们也更加闪耀。例如梅瑟（首先提到他），行了许多大奇事和征兆；但没有任何事使他如此伟大，如同他向天主说的那句有福的话：\BibleQuote{你若赦免他们的罪，就赦免；若不然，就把我从你所写的册上涂抹。}\BibleRef{出 32:32}达味也是如此：因此他也说：\BibleQuote{是我这牧者犯了罪，行了恶；但这群羊，他们作了什么呢？愿你的手攻击我和我父的家。}\BibleRef{撒下 24:17}同样，亚巴郎不寻求自己的利益，而是寻求许多人的利益。因此他既使自己置身危险，也为那些丝毫不属于他的人向天主求情。\par

是的，这些人如此变得光荣。至于那些寻求自己利益的人，也要想想他们受到了什么损害。例如，刚才提到的那位的侄儿，因为他听从了这句话：\BibleQuote{你若向右，我就向左；}\BibleRef{创 13:9}并接受了选择，寻求自己的利益，却连自己的利益也没有找到：这个地区被烧毁了，而那个地区却保持无恙。约纳再次不寻求许多人的利益，而是寻求自己的，甚至处于丧亡的危险中：当那城屹立不动时，他自己却在海上被抛掷和淹没。但当他寻求许多人的利益时，他也找到了自己的。同样，雅各伯在羊群中不寻求自己的利益，却得到了极大的财富为他的份。若瑟也寻求他兄弟们的利益，找到了自己的。至少，当他被父亲派遣时，\BibleRef{创 37:14}他没有说：\Quote{这是什么事？难道你们没有听见，因为一个异梦和某些梦，他们甚至企图把我撕碎，我因我的梦而受责难，因受你宠爱而受惩罚？那么当他们把我放在他们中间时，他们会做什么呢？}他没有说这些事，他没有想这些，而是将关切他的弟兄置于一切之上。因此他也享受了随后的一切美好事物，这使他非常卓越并宣告了他的光荣。同样，梅瑟——没有什么阻止我们再次提到他，并看他如何忽略自己的事而寻求他人的事：——我说这梅瑟，在君王的宫廷中生活，因为他\BibleQuote{看基督的凌辱比埃及的财宝更宝贵；}\BibleRef{希 11:26}并将这一切都从他手中抛出，成了希伯来人苦难的分享者；——他自己不但没有受奴役，反而也解放了他们。\par

是的，这些当然是伟大的事，值得天使般的生活。但保禄的行为远超于此。因为所有其余的人离开自己的福乐，选择分享他人的苦难；但保禄做了一件大得多的事。因为他不是同意分享他人的不幸，而是他自己选择处于一切绝境，好使他人能享受福乐。现在，一个生活在奢华中的人抛弃奢华而受苦，与一个独自受苦而使他人享受安全与尊荣，是不同的。因为在前一种情况下，虽然为了邻人的益处而以顺境交换逆境是一件伟大的事，但有受苦的同伴带来一些安慰。但同意自己独处困境以便他人享受他们的美好事物——这属于一个更有力的灵魂，属于保禄自己的精神。\par

不应仅凭这一点，而且凭另一种更伟大的卓越，他超越了先前提到的一切。即\Person{亚巴郎}和其余众人都在现世生活中置身于危险，而所有这些人都只是要求一次性地丧命；但\Person{保禄}却祈祷\BibleRef{罗 9:3}，说他宁愿从将来世界的荣耀中失落，为的是他人的救恩。\par

我还可以提及第三个优越之处。是什么呢？即有些先祖虽然为他们所反对之人代祷，但毕竟是为那些托付给他们引导的人；这就好像一个人为一个蛮横无法的儿子辩护，但终究是儿子；而\Person{保禄}却愿为那些没有被托付给他监护的人而被诅咒。因为他是被派遣到外邦人那里的。你注意到他灵魂的伟大和精神的崇高，超越天际吗？你要效法这人；但若不能，至少跟随那些在旧盟约中闪耀的人。因为这样，你若是寻求邻人的益处，就会找到自己的益处。因此，当你对照顾弟兄感到踌躇时，要想到别无他法可以得救，至少为了你自己，要为他及其利益挺身而出。\par

5. 虽然以上所说的足以说服你，别无他法可以确保我们自己的利益；然而，如果你还想通过日常生活的例子来确认这一点，设想一个房屋某处起火，然后一些邻居为了自己的利益而拒绝面对危险，反而闭门不出，留在家中，生怕有人闯入偷走一些家什；他们将受到多大的惩罚呢？因为火势会蔓延，同样烧毁他们的一切；由于他们不顾邻人的益处，连自己的也丧失了。因为天主，愿意将我们彼此连结，就把事物安排得如此必要，以至于一个邻人的利益与另一个人的利益交织在一起；整个世界就是这样构成的。因此，在船上也是如此，如果风暴来临，舵手不顾众人的利益，只寻求自己的利益，他会很快使自己和大家都沉没。我们也可以说，每一种技艺，如果只寻求自己的利益，生活就永远无法维持，连那寻求自身利益的技艺也无法存在。因此，农夫播种的谷物不是只够自己吃，否则他早就饿死了自己和别人；而是寻求许多人的利益：士兵上阵冒险，不是为了救自己，而是为了使他的城市安全；商人带回的货物不是只够自己用，而是也为许多其他人。\par

若有人说：\Quote{各人这样做，并不是着眼我的利益，而是他自己的利益，因为他从事这一切是为了为自己获得金钱、荣耀和安全，以致于寻求我的利益时，他是在寻求自己的利益。}这也是我所说的，并且早就希望从你们口中听到，为此我构想了全部讲论；即是要说明，你的邻人只有在顾及你的利益时，才是在寻求他自己的利益。因为既然人们除非被逼到这种必要，否则不会下定决心寻求邻人的事，因此天主这样将事物连结起来，不允许他们达到自己的利益，除非他们先经过别人的利益。\par

那么，追求邻人的利益是人的天性；但人不应因这个理由而被说服，而应因天主的喜悦。因为缺少这点，就不可能得救；即使你实践了最高境界的工作，却忽视了其他正在丧亡的人，你在天主面前也不会有信心。这从何显明呢？从真福\Person{保禄}所宣称的：\Quote{若是我把我的财物分给穷人，并舍身被火焚烧，却没有爱德，于我毫无益处。}\BibleRef{格前 13:3}他说。你看见\Person{保禄}向我们要求多少吗？然而，那将财物分给穷人的人，寻求的不是自己的利益，而是邻人的。但仅此还不够，他说。因为他要我们怀着真诚和极大的同情去做。因此，天主也立了这诫命，为使我们进入爱的联系。当祂要求如此大的尺度，而我们连那较小的也不付出时，我们配得怎样的宽容呢？\par

\Quote{怎么，}有人说，\Quote{天主藉天使对\Person{罗特}说：\InnerQuote{逃命吧？}}\BibleRef{创 19:17} 请问，什么时候，为什么？当惩罚临近时，不是在有改正机会的时候，而是当人们已被定罪、病入膏肓，老老少少都陷于同样的情欲，从此必然被烧毁，就在那雷电将要发出之日。此外，这话并非关乎恶习与美德，而是关乎天主所加的惩戒。告诉我，他该做什么？坐着不动，等待惩罚，并不丝毫有益于他们，而被烧死吗？不，这是极端的愚昧。\par

因为我并不断言，人应当不加分辨、随意地给自己招致惩罚，违背天主的旨意。但当一个人长期陷于罪恶时，那时我劝你挺身而出，纠正他：若你愿意，为了邻人的缘故；若不愿意，至少为了你自己的利益。的确，前者是更好的途径；但你若尚未达到那高度，即使为这个原因也去做。并且，不要让人寻求自己的利益以找到自己的利益；要牢记，除非我们拥有爱这至高的德行，否则自愿贫穷、殉道或其他任何事都不能为我们作证；让我们保有这超越其余一切的德行，好使我们藉着它获得所有其他事物，无论是现世的还是预许的福分；愿我们都能因我们的主\Person{耶稣基督}的恩宠和仁慈而到达那里；愿荣耀归于祂，永世无穷。阿们。\par

\SourceRecord{Translated by Talbot W. Chambers.；Nicene and Post-Nicene Fathers, First Series；Vol. 12.；Edited by Philip Schaff.；Buffalo, NY: Christian Literature Publishing Co.,；1889.；<http://www.newadvent.org/fathers/220125.htm>.}

% structure: p0001 source=h1 target=section reason=page-title

\section{论格林多前书第二十六篇讲道}

% structure: p0002 source=h2 target=subsection reason=relative-heading-rank; base=h2

\subsection{格林多前书 11:2}

我称赞你们在一切事上记念我，并持守我所传授给你们的传统。\par

他既完成了关于祭偶像之物的论述，恰如其分，且在各点上极其完美，便转而处理另一件事，这件事本身也是一项责备，但并非那样严重。因为正如我先前所说，我现在也要说：他并非连续列出所有重大指控，而是按适当顺序排列，在其中插入较轻的事项，以缓和读者因他不断责备而可能感到的反感。\par

因此，他也把最严重的一项——即有关复活的事——放在最后。但目前他转而处理另一件较轻的事，说：\Quote{我称赞你们在一切事上记念我。} 这样，当过错被承认时，他既强烈指责，又加以威胁；但当过错被质疑时，他先证明它，然后责备。对于已被承认的事，他加重其严重性；但对于可能引起争议的事，他表明它是被承认的。例如，他们的淫乱是公认的事。因此，无需证明那是一种过错；但在那种情况下，他证明了过犯的重大程度，并用比较的方式展开论述。再者，他们在外邦人面前诉讼是一种过错，但并非那样严重。因此，他顺便考虑了此事，并加以证明。至于祭偶像之物，则是有争议的。然而，那是一种极其严重的恶。因此，他既表明它是一种过错，又通过论述加以扩大。但当他这样做时，他不仅使他们远离各罪行，还邀请他们转向相反的美德。例如，他不仅说不可犯奸淫，同样也说应当表现出极大的圣洁。因此他补充说：\Quote{所以要在你们的身子和心灵上光荣天主。} \BibleRef{格前 6:20} 此外，他又说不可凭外邦的智慧而自夸，他并不满足于此，还吩咐他要\Quote{成为一个愚拙的人。} \BibleRef{格前 3:18} 在他劝告他们不要在外邦人面前诉讼，也不要做错事的地方，他更进一步，甚至完全禁止诉讼，并劝诫他们不仅不做错事，还要甘愿受屈。\EditorialNote{参阅：六章7-8节}\par

在论及祭偶像之物时，他不仅说应当戒绝被禁止的事，还应当在引人跌倒时也戒绝允许的事；并且不仅不伤害弟兄，甚至不伤害外邦人或犹太人。因此他说：\Quote{不可成为犹太人、希腊人或天主的教会跌倒的原因。} \BibleRef{格前 10:32}\par

2. 他既完成了关于这一切事的所有论述，便接着处理另一项责备。这是什么？他们的妇女在祈祷和说预言时不蒙头，露着头（因为那时妇女也说预言）；而男人却留长发，仿佛以哲学为业，并在祈祷和说预言时蒙头，这两者都是希腊人的习俗。既然他曾在场时已就这些事告诫过他们，也许有人听从了他，有人没有听从；因此在他的书信中，他又像医生一样，通过他说话的方式敷上药膏，从而纠正了过错。他从前曾亲自告诫过他们，这从他开头的话可以明显看出。否则，他在书信中此前从未提及此事，而是从其他责备转而直接说：\Quote{我称赞你们在一切事上记念我，并持守我所传授给你们的传统？}\par

你看，有些人听从了，他称赞他们；另一些人没有听从，他在后面纠正他们，说：\Quote{若有人似乎好争辩，我们没有这样的习俗。} \BibleRef{格前 11:16} 因为如果有些做得好，另一些人没有听从，而他把所有人都包括在责备中，那么他既会使前者更大胆，也会使后者更松懈；而如今他通过称赞和认可前者，责备后者，既更有效地激励前者，也使后者在他面前畏缩。因为责备本身就足以伤害他们；但现在它发生在与那些做得好并受称赞的人对比中，更加刺人。然而，目前他并不以责备开始，而是以赞扬和极大的赞扬开始，说：\Quote{我称赞你们在一切事上记念我。} 因为保禄的品格就是如此；即使是为了小事，他也编织出高度的赞扬；他这样做并非出于奉承；绝不是，他怎会如此行事呢？金钱、荣耀或其他任何东西对他来说都不可欲？他为他们的得救而安排一切。这就是他扩大赞扬的原因，说：\Quote{我称赞你们在一切事上记念我。}\par

哪些事？因为直到现在，他的论述仅仅关于他们不留长发和不蒙头；但正如我所说，他在赞扬中非常慷慨，使他们更加积极。因此他说：\par

\Quote{你们在一切事上记念我，并持守我所传授给你们的传统。} 由此可见，他当时也口头传授了许多未写下来的事，他在其他地方也多次表明这一点。但那时他只是传授，而现在则添加了对理由的解释：这样既使顺服的人更加坚定，也挫败了那些反对之人的骄傲。此外，他没有说：\Quote{你们听从了，而别人没有听从，} 而是通过随后的教导方式不引人怀疑地暗示了这一点，他说：\par

% structure: p0011 source=h2 target=subsection reason=relative-heading-rank; base=h2

\subsection{格林多前书 11:3}

\Quote{但我愿意你们知道：男人的头是基督，女人的头是男人，而基督的头是天主。}\par

这是他说明事情理由的论述，他陈述理由是为了使软弱者更加留心。确实，那信德坚固、坚定不移的人，对于所命令之事并不需要任何理由或原因，只满足于诫命本身。但软弱者若也得知原因，就会更谨慎地记住所说的话，并且更乐意地服从。\par

因此，直到他看见诫命被违反，他才陈述原因。那么原因是什么呢？\Quote{一切男人的头是基督。}那么祂也是外邦人的头吗？绝对不。因为\Quote{我们是基督的身体，各自是肢体}，\BibleRef{格前 12:27}，这样祂是我们的头，祂不能是那些不在身体内、不列在肢体中的人的头。所以，当他说\Quote{一切男人}时，必须理解为信徒。你看出他如何处处从高处论证，以唤起听者的羞愧吗？如此，当他论及爱德、谦卑和施舍时，也都是从那里汲取他的例证。\par

3.\Quote{但女人的头是男人；而基督的头是天主。}在此异端者向我们冲来，用某些从这些话中针对圣子捏造的卑下宣言。但他们自己绊倒自己。因为如果\Quote{男人是女人的头}，而头与身体是同质的，又\Quote{基督的头是天主}，那么圣子便与圣父同质。他们说：\Quote{不，我们并非要从此指出祂是不同本质，而是说祂处于隶属地位。}对此我们该说什么呢？首先，当任何低微之事论及祂与肉体结合时，所说的话并不贬损天主性，因为降生奥迹容许这样的表达。然而，请告诉我，你打算如何从这段经文证明这一点？为什么？正如男人统治妻子，他说，\Quote{同样，圣父也统治基督。}因此，正如基督统治男人，同样，圣父也统治圣子。我们读到：\Quote{一切男人的头是基督。}谁能承认这一点？因为如果圣子相对于我们的优越性，成为圣父相对于圣子的尺度，那么你要将祂贬低到何等卑微的地步！因此，在关于我们自己和关于天主的事上，我们不可用同样的尺度衡量一切，尽管所用的语言相似；而必须将一种合适的卓越性归于天主，且是相称于天主的卓越性。因为如果他们不承认这一点，就会产生许多荒谬之处。例如：\Quote{基督的头是天主；}而\Quote{基督是男人的头，男人是女人的头。}因此，如果我们选择在所有从句中取\Quote{头}一词的相同含义，那么圣子将如同我们远离祂一样远离圣父。不仅如此，女人也将如同我们远离天主的圣言一样远离我们。而圣子之于圣父的关系，就是我们之于圣子的关系，以及女人之于男人的关系。谁能容忍这个？\par

但你所理解的\Quote{头}一词，在男人和女人的情况下，与在基督的情况下不同吗？因此，在圣父和圣子的情况下，我们也必须不同地理解它。反对者说：\Quote{如何不同地理解？}根据场合。因为如果保禄要谈论统治和隶属，如你所说，他不会提出妻子的例子，而是奴隶和主人的例子。因为即使妻子隶属我们，她是作为妻子、作为自由者、作为同等尊荣。而圣子虽然确实服从了圣父，却是作为天主子、作为天主。因为如同圣子对圣父的服从大于人类对生养者的服从，同样祂的自由也更大。当然，不能说圣子与圣父的关系比人类之间的关系更大更密切，而圣父与圣子的关系则更小。因为如果我们钦佩圣子服从以至于死，且死在十字架上，并视此为关于祂的伟大奇事，那么我们也应当钦佩圣父，因为祂生了这样的儿子，不是作为受命令的奴隶，而是作为自由者，既服从又赐予劝言。因为劝言者不是奴隶。但是，当你听到劝言者时，不要理解为圣父有所需要，而是圣子与生祂者有同等尊荣。因此，不要将男人和女人的例子完全应用。\par

因为在我们中间，女人确实合理地隶属于男人：因为尊荣平等会引起争端。不仅为此，也因为起初发生的欺骗\BibleRef{弟前 2:14}。所以你看到，她并非一被造就受隶属；当天主将她带到男人面前时，既没有从天主听到任何这样的事，男人也没有对她说这样的话：他确实说她是\Quote{他骨中的骨，肉中的肉}，\BibleRef{创 2:23}但统治或隶属，他从未对她提及。然而，当她滥用了她的特权，那被造为助手的竟成为诱捕者并毁灭一切时，她才公正地为将来被告知：\Quote{你的归向将归于你的丈夫。}\BibleRef{创 3:16}\par

为此解释：这罪很可能使我们的种族陷入战争状态；（因为既然发生了这事，她由他而出并不有助于和平，反而这本身就使男人更加严厉，即她虽由他而出，却连自己身上的肢体都没有顾惜：）因此，天主考虑到魔鬼的恶意，就用这话立起了屏障，并通过这判决和植入我们心中的欲望，消除了由魔鬼的诡计可能产生的敌意：从而拆毁了隔墙，即那罪引起的愤恨。但在天主和那无玷的本体里，我们不可设想任何如此的事。\par

所以，不要将比喻应用于一切情况，因为从这同一个来源，在其他地方也会产生许多严重的错误。例如，就在这封书信的开头，他说：\BibleRef{格前 3:22} \BibleQuote{一切全是你们的，你们是基督的，基督是天主的。}那么，怎样呢？难道一切同样属于我们，如同\BibleQuote{我们是基督的，基督是天主的}吗？绝非如此。即使对最单纯的人来说，差别也是明显的，尽管同样的表达用于天主、基督和我们。又在别处，称丈夫为\Quote{妻子的头}后，他加上：\BibleRef{弗 5:23} \BibleQuote{如同基督是教会的头、救主和保护者，丈夫也当如此对待自己的妻子。}那么，我们是否应以同样的方式理解经文中的这句话，以及此后写给厄弗所人关于此主题的所有话呢？绝非如此。不可能。因为尽管同样的话用于天主和人，它们在天主和人身上并非具有同等效力，而必须以一种方式理解前者，以另一种方式理解后者。然而，并非一切事物都截然不同：否则，它们似乎是被随意且徒然地引入，我们从中一无所获。但正如我们不可完全等同地接受一切，同样也不可完全拒绝一切。\par

为了使我的话更清晰，我将努力用一个例子来说明。基督被称为\BibleQuote{教会的头}。如果我完全从人的观念中不取任何东西，那么，请问，这个表达为何被使用？另一方面，如果我完全照此理解，就会产生极端的荒谬。因为头与身体有相同的情感，并处于同样的境况。那么，我们应当放弃什么，接受什么？我们应当放弃我刚才提到的那些特殊之处，但接受完全合一的观念和元始；甚至这些观念也不可绝对化，而在这里，我们也必须自己形成一种概念，以领会那过于我们、适合于天主性的事物：因为合一更可靠，元始更尊贵。\par

同样，当你听到\Quote{圣子}这个词时，不要在此情况下接受所有细节；但也不可完全拒绝。要接受一切适合天主的，如：祂是同一性体，祂出自天主；那些不相称的、属于人性软弱的事物，你留在地上吧。\par

同样，天主被称为\BibleQuote{光}。那么我们是否要接受属于自然光的一切情况呢？绝非如此。因为这光屈服于黑暗，被空间所限，受另一种力量推动，能被遮蔽；这些对那本体连想都不可想。然而，我们不应因此拒绝一切，而应从比喻中汲取一些有用的东西。那就是从天主而来的光照，脱离黑暗，这些是我们要从中领会的。\par

4. 以上是对异端者的回答：但我们也必须有序地讲解整段经文。因为或许有人在此也会产生疑惑，自问：女人不蒙头，男人蒙头，究竟是什么样的罪过？是什么样的罪过，请从此处学习。\par

许多不同的象征被赐给了男人和女人：对男人是统治的象征，对女人是服从的象征；其中也包括这一点：女人应当蒙头，而男人头不蒙。如果这些是象征，那么你就会看到，当两者打乱正当秩序、违反天主的安排和各自本分时——男人堕入女人的低下地位，女人藉着外在服饰反抗男人——两者都犯了错误。\par

因为如果交换服装是不合法的——既不可让她穿外袍，也不可让他穿斗篷或面纱（\BibleQuote{女人不可穿戴男人的衣物，男人也不可穿女人的衣服}\BibleRef{申 22:5}）——那么，这些事互换就更加不体面了。前者（服装）确实是由人制定的，尽管后来天主予以批准；但后者——即蒙头或不蒙头——却是出于本性。但我说本性，意即天主。因为祂是创造本性的。因此，当你推翻这些界限时，请看产生了多大的损害。\par

不要告诉我，这只是小错。首先，它本身是大的，因为是悖逆。其次，即使它原本微小，但因它所象征的事物伟大而变得重大。然而，它是一件大事，这从它如此有效地维护人类秩序——统治者和被统治者各安其位——即可看出。\par

因此，违反者扰乱了万事，背叛了天主的恩赐，将从上而来的尊荣抛在地上；不仅男人如此，女人亦然。因为对她来说，最大的尊荣就是保守自己的本分；正如反叛的行为是最大的耻辱。因此，他针对两者制定了规定，这样说道：\par

% structure: p0028 source=h2 target=subsection reason=relative-heading-rank; base=h2

\subsection{格前 11:4-5}

\BibleQuote{凡男人祈祷或说预言，若蒙着头，就是羞辱自己的头。但凡女人祈祷或说预言，若不蒙头，就是羞辱自己的头。}\par

因为，如我所说，那时有男人说预言，也有女人拥有这神恩，如斐理伯的女儿们，\BibleRef{宗 21:9} 以及她们之前和之后的其他妇女；关于她们，先知早先也曾说过：\BibleQuote{你们的儿子要说预言，你们的女儿要见异象。}\BibleRef{岳 2:28}\par

那么：男人，他并非强迫他总不蒙头，而只是在祈祷时。\BibleQuote{因为凡男人}，他说，\BibleQuote{祈祷或说预言，若蒙着头，就是羞辱自己的头。}但女人，他命令她总该蒙头。因此，说了\BibleQuote{凡女人祈祷或说预言，若不蒙头，就是羞辱自己的头}之后，他并未停留于此，还继续说道，\BibleQuote{因为她与剃了头发一样。}倘若剃发总是可耻的，显然不蒙头也常是羞辱。他不仅满足于此，又加上说，\BibleQuote{女人应该在头上有权柄的记号，为了天使的缘故。}他表明，不仅祈祷时，而且一直，她都应该蒙头。至于男人，他不再论及蒙头，而是关于留长发，如此形成他的论述。他只在男人祈祷时禁止蒙头；但留长发则始终不鼓励。因此，关于女人，他说了，\BibleQuote{若不蒙头，就剪去头发；}同样关于男人，\BibleQuote{若有长头发，便是他的羞辱。}他并未说\BibleQuote{若蒙头}，而是\BibleQuote{若有长头发。}因此他在开头也说，\BibleQuote{凡男人祈祷或说预言，若头上有什么，就是羞辱自己的头。}他未说\BibleQuote{蒙头}，而说\BibleQuote{头上有什么}；表明即使他光头祈祷，若有长头发，也如同蒙着头。\BibleQuote{因为头发}，他说，\BibleQuote{是用来作遮盖的。}\par

% structure: p0032 source=h2 target=subsection reason=relative-heading-rank; base=h2

\subsection{\BibleRef{格前 11:6}}

\BibleQuote{若不蒙头，就该剪去头发；但若女人剪发或剃发为可耻，就该蒙头。}\par

如此，起初他只要求头不要裸露；但继续下去，他暗示这规则的持续性，说道：\BibleQuote{因为她与剃了头发一样，}并且要以一切谨慎和勤奋遵守它。因为他不仅说蒙头，而是\Quote{完全蒙上}，意思是要她每一边都仔细包好。并且他以荒谬的方式，诉诸他们的羞耻，严厉责备说：\BibleQuote{若不蒙头，就剪去头发。}仿佛他说：\Quote{你若抛掉天主法律所定的遮盖，也抛掉本性所定的遮盖吧。}\par

若有人说：\Quote{怎么，女人若上升到男人的光荣，这怎能是她的羞耻呢？}我们可以这样回答：\Quote{她并非上升，反而从自己应有的尊荣中坠落。}因为不守自己的界限和天主制定的法律，而越界，不是增加，而是减少。正如那些贪图别人的财物、夺取非己之物的人，并未获得更多，反而亏损，连自己原有的也失去了（这也发生在乐园里），同样，女人并未获得男人的尊严，反而失去自己原有的女性端庄。不仅从此处，而且是因她的贪心，才有她的羞耻和责备。\par

他取了公认可耻的事，说道：\BibleQuote{但若女人剪发或剃发为可耻，}接着陈述自己的结论，说：\BibleQuote{就该蒙头。}他未说\Quote{该有长头发，}而说\BibleQuote{该蒙头，}将这两者视为同一，并从习惯及其反面两方面确立：因为他既肯定蒙头和头发是一回事，又肯定剃发的女人与光头女人相同。\BibleQuote{因为这是一样的事，}他说，\BibleQuote{如同她剃了头发。}若有人说：\Quote{怎么一样呢？这女人有本性的遮盖，而剃发的女人连这个也没有？}我们回答：就她的意愿而言，她也丢弃了那遮盖，因为头上赤裸。若她并非缺发，那是本性的作为，不是她自己的。因此，如同剃发的女人头是赤裸的，这女人也一样。为此，他让本性为她提供遮盖，好使她从本性也学到这一课而蒙头。\par

随后他也说明一个原因，如同对自由人说话：这是我在许多地方已经注意到的。那么原因是什么？\par

% structure: p0038 source=h2 target=subsection reason=relative-heading-rank; base=h2

\subsection{\BibleRef{格前 11:7}}

\BibleQuote{男人本不该蒙着头，因为他是天主的肖像和光荣。}\par

这又是另一个原因。\Quote{不仅，}他如此说，\Quote{因为基督是他的头，他不该蒙头，而且因为他管辖女人。}因为统治者来到君王面前，应当有他统治的象征。因此，正如没有统治者敢不带武带和披风出现在戴冠冕者面前，同样，你也不可没有你统治的象征（其中之一就是不被蒙头），在天主面前祈祷，免得你羞辱自己和那尊荣你的主。\par

关于女人，同样也可以这么说。因为对她而言，没有她服从的象征也是羞辱。\BibleQuote{女人是男人的光荣。}因此男人的管辖是自然的。\par

5. 然后，在肯定了论点之后，他又陈述了其他理由和原因，引导你回到最初的创造，这样说：\par

% structure: p0043 source=h2 target=subsection reason=relative-heading-rank; base=h2

\subsection{\BibleRef{格前 11:8}}

\BibleQuote{因为男人不是出于女人，而是女人出于男人。}\par

若出于某人，是那人的光荣，那么作他肖像更是如此。\par

% structure: p0046 source=h2 target=subsection reason=relative-heading-rank; base=h2

\subsection{\BibleRef{格前 11:9}}

\BibleQuote{因为男人不是为女人造的，而女人是为男人造的。}\par

这又是第二个优越性，甚至也是第三个和第四个：第一，基督是我们的头，我们是女人的头；第二，我们是天主的荣耀，而女人是我们的荣耀；第三，我们不是出于女人，而她出于我们；第四，我们不是为了她，而她是为了我们。\par

% structure: p0049 source=h2 target=subsection reason=relative-heading-rank; base=h2

\subsection{\BibleRef{格前 11:10}}

\Quote{因此，女人应当在头上有一个权柄的记号。}\par

\Quote{因此：}是什么原因，请告诉我？他说：\Quote{为了上述这一切；}或者不如说，不仅为了这些，也\Quote{因了天使的缘故。}他又说：\Quote{纵然你轻视你的丈夫，}但仍要\Quote{尊敬天使。}\par

从而可知，蒙头是服从和权柄的记号。因为这使她低头羞愧，并完全保持自己应有的美德。因为被管辖者的美德与光荣在于坚守自己的服从。\par

再者：男人并不被迫这样做；因为他是他主的肖像；但女人却被迫，且合情合理。那么，试想这过犯之深：你既被授予如此崇高的特权，却自取其辱，夺去女人的服饰。你所作的正如同领受了一顶王冠，却将王冠从头上丢弃，反而取了一件奴隶的衣服。\par

% structure: p0054 source=h2 target=subsection reason=relative-heading-rank; base=h2

\subsection{格前 11:11}

\Quote{然而，在主内，男人不能没有女人，女人也不能没有男人。}\par

如此，因为他赋予男人极大的优越，说女人出于他、为他、且在他之下；为了既不使男人被抬举得过高，也不使女人被压抑得过低，请看他是如何引入修正的，说：\Quote{然而，在主内，男人不能没有女人，女人也不能没有男人。}他说：\Quote{我请求你们，不要只考察起初的事和那造化。因为如果你们考察后来之事，两者互为因果；或者甚至也不能说是彼此为因，而是天主为万有之因。}因此他说：\Quote{在主内，男人不能没有女人，女人也不能没有男人。}\par

% structure: p0057 source=h2 target=subsection reason=relative-heading-rank; base=h2

\subsection{格前 11:12}

\Quote{因为正如女人是由男人而出，照样男人也是由女人而生。}\par

他没有说\Quote{由女人，}而是重复了那表达，\BibleRef{格前 11:7}说\Quote{由男人。}因为这项特殊的特权仍然完全保留给男人。然而，这些优点并非男人的财产，而是天主的。因此他也加上：\Quote{但一切都是出于天主。}所以，如果一切都是属于天主的，而祂命定了这些事，你就当服从，不可争辩。\par

% structure: p0060 source=h2 target=subsection reason=relative-heading-rank; base=h2

\subsection{格前 11:13}

\Quote{你们自己判断：女人蒙头向天主祈祷，适当吗？}他又让他们对所讲的话作评判，正如他对待祭偶像之物时所做的一样。因为正如在那里他说：\Quote{你们判断我所说的：}\BibleRef{格前 10:15}，同样在这里：\Quote{你们自己判断：}并且他在这里暗示了更可怕的事。因为他说，这侮辱直接上达于天主：虽然他自己并没有这样表达，而是用一种更温和、更隐晦的说法：\Quote{女人不蒙头向天主祈祷，适当吗？}\par

% structure: p0062 source=h2 target=subsection reason=relative-heading-rank; base=h2

\subsection{格前 11:14}

\Quote{难道本性本身没有教导你们，男人若留长发，就是他的羞辱吗？}\par

% structure: p0064 source=h2 target=subsection reason=relative-heading-rank; base=h2

\subsection{格前 11:15}

\Quote{但女人若留长发，却是她的荣耀；因为她的头发是给她作蒙头的。}\par

他常常引用公认的理由，在此处也是如此，求助于普遍习俗，并使那些等待从他这里学习这些事的人大大羞愧，这些事他们本可以从人们的日常实践中获得认知。因为这样的事甚至连外邦人也并非不知：请看他在各处是如何使用尖锐的表达：\Quote{凡男人祷告若蒙着头，就是羞辱自己的头；}并且又说：\Quote{但女人若剪了发或剃了头，就是羞辱，她就该蒙头；}这里又指：\Quote{男人若留长发，就是他的羞辱；但女人若留长发，却是她的荣耀，因为她的头发是给她作蒙头的。}\par

\Quote{若它已给她作蒙头，}你们说，\Quote{为什么还需要再加一个蒙头呢？}为的是不仅本性，而且她自己的意愿也当在她承认服从中有份。因为本性本身早已预先制定了一条法律，你应当蒙头。所以，我请求你，也加上你自己的那一份，免得你似乎颠覆了本性本身的规律；这是极其傲慢鲁莽的证据，不仅与我们一起抗争，也与本性抗争。这就是为什么天主指责犹太人时说，\BibleQuote{你杀了你的儿女：这超过了你所有可憎之事。}\BibleRef{则 16:21}\par

再者，保禄责备罗马人中的不洁者时，如此加重指控说，他们的行为不仅违背天主法律，甚至违背本性。\BibleQuote{因为他们把自然之用转为反乎自然之用。}\BibleRef{罗 1:26}因此，他也在这里使用这个论据，表明这件事本身：他并没有制定任何新奇的法律，而且在外邦人中，他们的发明都会被视作一种违背本性的新奇之物。同样，基督也暗示了同一件事，说：\Quote{凡你们愿意人给你们做的，你们也要照样给人做；}表明祂并没有引进任何新事。\par

% structure: p0069 source=h2 target=subsection reason=relative-heading-rank; base=h2

\subsection{格前 11:16}

\Quote{但若有人似乎是好辩的，我们却没有这样的规矩，天主的众教会也没有。}\par

因此，反对这些事就是好辩，而非任何理性的运用。然而，即便如此，他采用的是一种有节制的斥责，为使他们更加自愧；这实际上使他的话更为严厉。他说：\Quote{因为我们，}没有\Quote{这样的规矩，}好争辩、好争斗、好反对自己。他甚至没有停在这里，还加上了：\Quote{天主的众教会也没有；}表示他们因不肯让步而抵抗并反对全世界。然而，即使格林多人当时好辩，但现在全世界都已接受并遵守了这条法律。被钉十字架者的能力何其伟大！\par

6. 但我恐怕有些妇女虽然穿上了衣服，在行为上却显出不端庄，在其他方面也是裸露的。因为\Person{保禄}在写给\Person{弟茂德}的书信中也不满足于这些事，还加了别的，说：\BibleQuote{她们应以正派服装，以廉耻和庄重装饰自己，不以编发和金饰。} \BibleRef{弟前 2:9} 因为如果人不该露着头，而应到处带着权柄的记号，那么更应在行为上表现出这一点。因此，从前的妇女也常称丈夫为主，\BibleRef{伯前 3:6} 并将首位让给他们。\Quote{因为他们那方面，}你说，\Quote{曾爱自己的妻子。} 你和我知道的一样清楚：我并不无知。但当我们劝告你们关于你们自己的本分时，不要让他们的本分占据你们全部的注意力。同样，当我们劝孩子们服从父母，说：经上记载：\BibleQuote{孝敬你的父亲和你的母亲，}他们回复我们说：\Quote{也提一下后面的话：\InnerQuote{并且你们作父亲的，不要惹你们的子女发怒。}}\BibleRef{弗 6:1} 并且仆人们，当我们告诉他们经上记载他们应当\BibleQuote{服从自己的主人，不要只作眼下的服务}，他们又要求我们提后面的话，叫我们也同样劝告主人。因为他们看见\Person{保禄}也命令他们\BibleQuote{不要威吓。} 但让我们不要这样做，也不要当我们被指控有关自己的事时去追究别人所担负的；因为即使你找到一个分担指控的同伴，也不能使你免于责备：而只应注重一件事，即如何摆脱那些针对你的指控。因为\Person{亚当}也把罪责推给女人，而女人又推给蛇，但这丝毫没有解救他们。因此，你也不要对我说这些，而应小心谨慎，尽力履行你对丈夫应尽的义务；因为当我在与你丈夫谈话，劝他爱惜你、珍视你时，我不容许他提出为女人所定的法律，而要求他履行为他本人所写的规定。所以你也只应忙于属于你自己的事，向你的伴侣表示顺从。因此，如果你真是为了天主的缘故而服从你的丈夫，就不要向我讲述他应该做的事，而只应精确地履行立法者让你负责的那些事。因为这就是特别服从\Person{天主}，即使在遭受相反的事时也不违反法律。同理，那被爱而爱人的，不被算作行了什么大事；但那服事恨他的人，这个人尤其要接受冠冕。同样，你也应认为如果丈夫令你厌恶，而你忍受了，你将得到荣耀的冠冕；但如果他温柔和善，\Person{天主}还有什么可赏报你的呢？我说这些话，不是要丈夫们严厉，而是劝妻子们甚至忍受丈夫的粗暴。因为当每个人都谨慎履行自己的本分时，对方的职责也会很快随之而来：正如妻子准备好甚至忍受丈夫的粗暴行为，而丈夫在她生气时不辱骂她；那么一切都平静如水，是风平浪静的港湾。\par

7. 古时的人也是如此。每个人都履行自己的本分，而不强求邻人的本分。因此，你若留意，\Person{亚巴郎}带了他兄弟的儿子；他的妻子没有责怪他。他命令她长途跋涉；她甚至没有反对，而是跟随。再者，在经历了许多苦难、辛劳和劳碌之后，他成了万物的主宰，却将优先权让给了\Person{罗特}。\Person{撒辣}并未因此生气，她甚至没有开口，也没有说出任何像当今许多妇女所说的话，她们看见自己的丈夫在这样的分配中处于劣势，尤其是在对待下等人时，她们就责备他们，称他们为傻瓜、无知、懦弱、背叛和愚蠢。但她没有说也没有想任何这类事，反而对他所做的一切感到满意。\par

还有另一件事，而且更大：在\Person{罗特}有权选择，并把较差的份给了他的叔父之后，一个大危险临到了他身上。\par

族长听见这事，就武装了他所有的人民，只用自己的家人对抗整个\Place{波斯}军队；那时她也没有阻止他，也没有像可能的那样说：\Quote{人啊，你要往哪里去？投身于悬崖绝壁，暴露于如此大的危险？为了一个得罪了你、夺取了你一切、流你血的人？即使你轻看自己，也请可怜我，我撇下了房屋、家乡、朋友和亲属，跟随你经过了漫长的漂泊；不要使我成为寡妇，陷入寡妇的悲惨境地。}她什么也没有说；她没有想这些，而是默默地忍受了一切。\par

此后，她的子宫仍不生育，她自己并不受妇女的悲伤，也不哀哭；但亚巴郎抱怨，却不是向妻子，而是向天主。请看每人如何持守自己合宜的本分：他并未因撒辣无子而轻视她，也未以此类事责备她；而她却急于想借使女为他无子策划一些安慰。因为那时这些事尚未像现在一样被禁止。因为现在，妇女不许纵容丈夫行此类事，男子也不许，无论妻子知情与否，建立这样的关系，即使无子的痛苦无限地折磨他们也不可；因为他们也要听到那句话：\BibleQuote{他们的虫子不死，他们的火也不灭}。因为现在这是不允许的，但那时尚未被禁止。因此，他的妻子命令此事，他听从了，但即便如此，也不是为了快感。但\Quote{看，}有人会说，\Quote{他如何又因她的吩咐把哈加尔赶出去}。是的，这正是我要指出的：他在一切事上服从她，她也在一切事上服从他。但你，你这以此为由的人，不要只注意这些，也要考察先前所发生的事：哈加尔侮辱主母，自夸反对主母；对于一个自由而尊贵的妇人来说，还有什么比这更令人烦恼的呢？\par

8. 那么，妻子不要等待丈夫的美德，然后才显出自己的美德，因为这不是什么大事；另一方面，丈夫也不要等待妻子的服从，然后才操练自制；因为这样也不再是他自己的善行。但正如我所说，让每人首先尽自己的一份。因为如果对外邦人打我们的右脸，我们须转过另一边脸；那么对于丈夫的粗暴行为，更应忍受。\par

我这样说并不是要妻子挨打；绝不是：因为这是极端的侮辱，不是对被殴打者，而是对打人者。但即使你因某种不幸被分配了这样的伴侣，妇人，不要介意，考虑到为这些事所预备的赏报以及在此生中的称赞。对你们做丈夫的，我也这样说：立下一个规矩，没有任何冒犯可以迫使你必须动手打妻子。为什么我说妻子？因为自由的人甚至不能忍受对使女施加打击和暴力。但如果打使女已是极大的羞耻，那么伸手打自由人就更甚了。这一点甚至可以从异教的立法者看出，他们不再强迫受虐待的妻子与打她的丈夫同住，认为他不配与她共处。因为当你的生活伴侣，那在最亲密关系和最高程度上与你结合的人，像一个卑贱的奴隶一样被你羞辱时，这无疑源于极度的无法无天。因此，这样的人，如果还可以称之为人而不是野兽，我想说，他就像一个弑父者和弑母者。因为如果为了妻子的缘故，我们被命令甚至离开父亲和母亲，不是亏待他们，而是履行天主的法律；并且这法律如此令我们父母喜悦，以至于连他们——那些我们离弃的人——也感激，并非常热切地促成此事；那么，为了她而令天主命令我们离开父母的人，侮辱她，岂不是极度的疯狂？\par

但我们大可以问：这仅仅是疯狂吗？还有羞耻：我倒想知道谁能忍受它。有什么描述能向我们呈现呢？当尖叫声和哀号声沿着小巷传来，邻居和路人纷纷跑向那个如此自辱之人的家，仿佛有野兽在其中肆虐？还不如地裂开吞没这样一个狂人，总好过他在那之后出现在广场上。\par

他说：\Quote{但这妇人傲慢无礼。}尽管如此，请考虑她是个女人，是较弱的器皿，而你是个男人。因此你被设立为统治者，并被赋予她作为头的位置，好使你忍受在你以下的她的软弱。要使你的统治光荣。当你的臣属没有从你那里受到任何羞辱时，你的统治就是光荣的。正如君主越显赫，就越彰显他属下的官吏的威严；但如果他羞辱并贬低那官阶的伟大，也就间接地削减了自己荣耀的不小部分：同样地，你若羞辱那在你以下治理的人，也会极大损害你治理的荣誉。\par

因此，考虑这一切，你要管束自己；同时也想想那个晚上，父亲召叫了你，将他的女儿如同一种寄托物交付给你，并将她与一切分离，与她的母亲、与自己、与家族分离，将她的全部守护托付给了你的右手。想想，你（在天主之下）藉着她有了儿女，成为了父亲，因此你也要对她温柔。\par

你们岂不见农夫，一旦土地接受了种子，他们便用各种栽培方法悉心照料，即使土地有千般不利：例如土质恶劣、杂草丛生，或因地势关系而遭过多雨水侵袭？你也要这样做。因为这样你必先享受果实和安宁。你的妻子对你而言，既是港湾，又是使你心中喜乐的强大疗愈之方。那么，如果你使你的港湾免于风和浪，你从市场归来时便会享有极大的宁静；但如果你以吵闹和骚动充满它，你不过是为自己预备更悲惨的沉船。\newline{}为避免此事，请照我所建议的去做：家中发生任何不快之事时，若她有错，安慰她，不要加重不快。因为即使你丧失一切，也没有比与没有善意的妻子同居更痛苦的了。无论你提起何种冒犯，你告诉我的事都不会像与她争吵那样痛苦。因此，即使只为这些理由，让她的爱比一切更宝贵。因为如果我们要互相担当重担，更当担当我们自己妻子的。\par

即使她贫穷，不要责备她；即使她愚笨，不要践踏她，而应教导她：因为她是你的肢体，你们已成为一体。\Quote{但她轻浮、酗酒、暴躁。}那么你应为这些事忧伤，而非发怒；并祈求天主、劝诫她、给她建议、尽一切努力除去这恶。但若你打她，你便加重了病患：因为凶暴是由温和而非由对抗的凶暴除去。同时，切记来自天主的赏报：当你被允许休弃她，却因敬畏天主而不这样做，并忍受如此大的缺陷，畏惧在此事上规定的法律，即禁止无论她患何病痛而休妻，你必获得不可言喻的赏报。是的，在赏报之前，你已大有收获，既使她更顺从，也使你自己因而变得更温和。\newline{}例如，据说有一位外邦哲学家，他有一个坏妻子，既轻浮又好争吵。当被问及\Quote{为何有这样一个妻子还忍受她}时，他回答说：\Quote{为使我家中有一所哲学学校与训练所。因为每日在这里受训，我对所有其他人将更温和。}你们大声欢呼了吗？唉，此刻我正大大哀伤，当外邦人比我们更热爱智慧时；我们这些被命令效法天使，不，被命令在温良方面追随天主本身的人。\par

进一步说：据说为此原因，那位哲学家虽有坏妻子，却未将她赶出；有人甚至说这正是他娶她的理由。但我，因许多人具有不甚合理的性格，建议他们起初竭尽全力，小心选择一个合适的伴侣，一个充满一切德行的伴侣。然而，若他们未能如愿，带进家中的新娘不善良或不差强人意，那么我仍愿他们至少努力效法那位哲学家，以各种方式训练她，并将此视为最要之事。因为一个商人若未与合伙人达成能带来和平的协议，就不会将船驶入深海，也不会着手进行其余交易。同样，让我们尽力让那与我们同舟共济、同航此世之船的人，能在一切和平中安于内。因为这样，我们其余的事务也会一片平静，我们将在宁静中航行过今生的大洋。与此相比，房屋、奴隶、金钱、田地，甚至国事本身，都当视为次要。在我们眼中，那与我们一起摇橹的人不与我们发生叛乱和分裂，这应比一切更宝贵。因为如此，我们其余的事也会顺风顺水，在属灵之事上我们也将发现自己少受阻碍，同心合意地背负这轭；且完成一切善事后，我们将获得预许的福分。愿我们都能达到，赖我们的主耶稣基督的恩宠和慈爱，祂与圣父及圣神，永享光荣、权能和尊威，自今至永远，世世无尽。阿们。\par

\SourceRecord{Translated by Talbot W. Chambers.；Nicene and Post-Nicene Fathers, First Series；Vol. 12.；Edited by Philip Schaff.；Buffalo, NY: Christian Literature Publishing Co.,；1889.；<http://www.newadvent.org/fathers/220126.htm>.}

% structure: p0001 source=h1 target=section reason=page-title

\section{格林多前书第廿七篇讲道}

% structure: p0002 source=h2 target=subsection reason=relative-heading-rank; base=h2

\subsection{格前 11:17}

\BibleQuote{然而我吩咐这事，我并不称赞你们，因为你们聚在一起，并不是为得益处，而是为受损。}\par

在考虑当前的吩咐时，必须先说明其缘由。因为这样，我们的讲论才会更清晰。那么，这缘由是什么？\par

正如当初那些相信的三千人，大家都共同用饭、一切公用；在宗徒写此信时，也保持着这样的做法。并非完全相同，但可以说，那存于他们中间的共融的某种流露，也传给了后来的他们。因为当然有些人贫穷，另有些人富有，他们并非将全部财物放在中间，而是在指定的日子开放餐桌，似乎是这样的。当庄严的礼仪完成后，在奥迹的共融之后，他们全都赴一场共同的宴席，富人带着自己的食物，穷困者由他们邀请，大家共同宴饮。但后来这习俗也败坏了。原因是他们分裂，有的人依附这一派，有的人依附那一派，并说：\Quote{我是属某人的}，\Quote{我是属某人的}；为纠正此事，他在书信开头说：\BibleQuote{因为关于你们，我的弟兄们，我由\Person{克罗厄}家里的人得知你们中间有纷争。我的意思是：你们各自说：我是属\Person{保禄}的；我是属\Person{阿颇罗}的；我是属\Person{刻法}的。}并非\Person{保禄}是他们所依附的人；因为他不会容忍此事；但他想通过让步从根本上铲除这习俗，便引用了自己，表明如果有人甚至将自己的名字标榜在脱离共同身体的事上，那么所行的事也是亵渎的、极端的邪恶。若在他身上是邪恶，那么在那些不如他的人身上更是如此。\par

2. 既然这习俗——一个极其优秀且极其有益的习俗（因为它是爱德的基础、贫穷的安慰、富贵的纠正、最高哲学的机会和谦卑的教导）——被破坏了，既然他看到如此巨大的益处正在被摧毁，他自然严厉地对他们说话，如此说：\BibleQuote{然而我吩咐这事，我并不称赞你们。}因为在先前的吩咐中，因有许多人遵守（那些规条），他以另一种方式开始，如此说：\BibleQuote{我称赞你们，因你们凡事记念我。}但此处相反：\BibleQuote{然而我吩咐这事，我并不称赞你们。}这就是为什么他没有把它放在责备那些吃祭偶像之物的人之后。但因为那责备异常严厉，他插入了关于留长发的讲论，以免他从一组严厉的谴责又转到另一组令人反感的谴责，显得过于严厉；然后他回到更严厉的语气，说：\BibleQuote{然而我吩咐这事，我并不称赞你们。}这是什么？那就是我要告诉你们的。\BibleQuote{我吩咐这事，我并不称赞你们}是什么意思？他说：\Quote{我不认可你们，因为你们使我不得不给予劝告；我不称赞你们，因为你们在这事上需要教导，因为你们需要我的训诫。}你们看出他如何从一开始就表示所行的事非常亵渎吗？因为当犯错者连一个小小的暗示都不应当需要以避免犯错时，这错似乎就不可原谅了。\par

为什么不称赞呢？因为他说：\BibleQuote{你们聚在一起，并不是为得益处，而是为受损。}就是说，你们没有在德行上前进。因为你们的慷慨本应增长并变得丰富，你们却反而从已有的习俗中后退了，并且后退到甚至需要我从旁警告，以便你们回到原先的秩序。\par

再者，为了不显得他这些话只是为穷人而说，他没有立即切入关于餐桌的讲论，以免他的责备变得容易被轻视，而是寻找一个最令人困窘、最令人畏惧的表达。他说了什么呢？\par

% structure: p0009 source=h2 target=subsection reason=relative-heading-rank; base=h2

\subsection{格前 11:18}

\BibleQuote{首先，你们聚在教会时，我听说你们中间有分裂。}\par

他并没有说：\Quote{恐怕你们没有共同用餐；我听说你们私下宴饮，不与穷人一起。}而是写下最能震撼他们心灵的事，即分裂之名，这也是那祸害的原因。于是他再次提醒他们那在书信开头所说的，并由\Quote{\Person{克罗厄}家里的人}所\Quote{告知}的事。\BibleRef{格前 1:11}\BibleQuote{我也略略相信。}\par

因此，为免他们说：\Quote{如果告状者说谎呢？}他既不说：\Quote{我相信，}以免使他们更放肆；也不说：\Quote{我不相信，}以免显得无故责备；而是说：\BibleQuote{我也略略相信。}即：\Quote{我信一小部分，}\par

% structure: p0013 source=h2 target=subsection reason=relative-heading-rank; base=h2

\subsection{格前 11:19}

3. \BibleQuote{因为在你们中间必须分党结派，好叫那些经得起考验的人得以显明出来。}\par

\Quote{分党}一词，他这里指的是那些不涉及教义，而是这些当前的分裂。但即使他谈论的是教义上的异端，他也并未给他们任何把柄。因为基督亲自说过：\BibleQuote{引人跌倒的事是免不了的，}\BibleRef{玛 18:7}不摧毁意志的自由，也不设定任何必然性和强迫凌驾于人的生命之上，而是预告那必定会从人的邪恶心智中产生的事；这些事会发生，不是因为他的预言，而是因为那些不可救药的人有此心意。并非因为他预言了，这些事才发生；而是因为它们必定要发生，所以他预言了它们。既然引人跌倒的事若出于必然而非出于引入者的心意，那么他说\Quote{祸哉，那引人跌倒的人}便是多余的。但我们在讲解那段经文时已经详细讨论过这些；现在我们必须继续讲面前的内容。\par

既然他说的这些关于这些与宴席有关的分党，以及那种争执和分裂，他也从接着的话显明了这一点。因为他说了\Quote{我听说你们中间有分裂}，他并未就此止步，而是指明他所指的分裂是什么，继续说到\Quote{各人先吃自己的晚餐}；又说：\Quote{什么？你们没有房子可以吃喝吗？还是你们藐视天主的教会？}然而，他说的正是这些，这是明显的。如果他把它们称为分裂，不必惊奇。因为，如我所说，他愿意用此说法触动他们；但如果这是教义上的分裂，他就不会这样温和地与他们谈话了。听他的例子，当他谈到任何此类事情时，他在断言和责备中是多么激烈：在断言中，如他说：\BibleQuote{即使一位天使从天上给你们宣讲别的福音，与你们所领受的不同，当受诅咒；}\BibleRef{迦 1:8}但在责备中，如他说：\BibleQuote{凡你们这些靠法律成义的人，是与恩宠断绝了关系}\BibleRef{迦 5:4}。有一次他称那些败坏者为\Quote{狗}，说：\BibleQuote{你们要提防狗：}\BibleRef{斐 3:2}另一次说：\BibleQuote{他们的良心被烙铁烙过}\BibleRef{弟前 4:2}；又说：\BibleQuote{撒旦的天使：}\BibleRef{格后 11:14}但是此处他没有说这样的话，而是以温和低调的语调说话。\par

但是\Quote{好叫那些经得起考验的人在你们中间显明出来}是什么意思？意思是使他们更加闪耀。他想要说的就是，那些坚定不变的人不仅丝毫不会因此受损害，反而更显明他们，并且使他们更加光荣。因为\Quote{为的是，}一词并不总是指示原因，也常常指示事情的结果。同样，基督自己就这样使用它，当他说：\BibleQuote{为审判我来到这世界；叫那看不见的看见，叫那看见的成为瞎子。}\BibleRef{若 9:39}同样，保禄在别处，当谈论法律时，他写道：\BibleQuote{而且法律也介入，好叫过犯增多。}\BibleRef{罗 5:20}但是，法律既不是为此目的而赐予，即犹太人的过犯增多（尽管确实发生了这样的结果），基督也不是为此目的而来，即叫那看见的成为瞎子，而是为了相反的目的；但结果却是如此。因此，在这里也必须这样理解\Quote{好叫那些经得起考验的人显明出来}这句话。因为异端的出现绝非为了这个目的，即\Quote{好叫那些经得起考验的人显明出来}，而是随着这些异端的出现，产生了这样的结果。现在他说这些是为了安慰那些穷人，那些高贵的忍受了这种轻蔑的人。所以他不是说\Quote{好叫他们成为经得起考验的人}，而是说好叫那些经得起考验的人显明出来；表明在此之前他们也是这样的人，但他们混杂在众人之中，虽然享受着富人给予的这种救济，却并不显眼；而现在这种争执和纷争使他们显明出来，就如同暴风雨显明舵手一样。他也不是说\Quote{好叫你们显为经得起考验的人}，而是\Quote{好叫那些经得起考验的人显明出来，就是你们中间这样的人}。因为他既不在此处责备时把他们都揭发出来，以免使他们更加鲁莽；也不在此处赞美时，以免使他们更加自夸；而是将这两种说法都悬置起来，让各人的良心去应用他所说的话。\par

而且，在我看来，他这里似乎不仅安慰穷人，也安慰那些没有违反惯例的人。因为很可能他们中间也有遵守它的人。\par

这就是为什么他说\Quote{我部分地相信它}。所以他恰当地称这些人为\Quote{经得起考验的人}，他们不仅与其余人一起遵守了惯例，甚至在没有他们的情况下也安然持守了这条好法律。他这样做，是藉着这样的赞美，努力使其他人和他们自己都更进步。\par

4. 最后他加上了具体的冒犯形式。是什么？\par

% structure: p0021 source=h2 target=subsection reason=relative-heading-rank; base=h2

\subsection{\BibleRef{格前 11:20}}

\Quote{当你们聚集在一起时，}他说，\Quote{不可能吃主的晚餐。}\par

你看他是如何有效地诉诸他们的羞愧，甚至通过叙述的方式，设法给他们忠告？他说：\Quote{因为你们聚会的面貌，}他说，\Quote{是不同的。那是爱与兄弟情谊的面貌。至少一个地方接纳你们众人，你们同在一群。但筵席，当你们来到它面前时，与敬拜者的集会毫无相似之处。}他没有说：你们聚集时，这不是共同吃饭；\Quote{这不是彼此宴饮；}而是以另一种更严厉的方式斥责他们说：\Quote{不可能吃主的晚餐}，此时将他们从这一点遣送到那个基督交付可畏奥迹的晚上。因此他也称那早先的筵席为\Quote{晚餐}。因为那晚餐也使所有人一同斜倚用餐：但富人与穷人之间的距离，当然不如师傅与门徒之间的距离那么大。因为那是无限的。我为什么说师傅与门徒呢？试想师傅与叛徒之间的差距：然而，主亲自与他们一同坐席，甚至没有将他赶走，反而给了他一份盐，并使他成为奥迹的分享者。\par

接着他解释如何\BibleQuote{不可能吃主的晚餐}。\par

% structure: p0025 source=h2 target=subsection reason=relative-heading-rank; base=h2

\subsection{\BibleRef{格前 11:21}}

\BibleQuote{因为在你们吃的时候，各人预先取用自己的晚餐}，他说，\BibleQuote{并且一个饥饿，另一个醉酒}。\par

你看出他是如何暗示他们更是在羞辱自己吗？因为那属于主的，他们却使之成为私事：所以他们自己是首先受辱的，剥夺了他们宴席的最大特权。怎样以何种方式？因为主的晚餐，即主的筵席，应当是共有的。因为主人的财产并不属于这个仆人而不属于那个，而是所有人都共有。所以藉着\Quote{主的}晚餐，他表达了这一点，即筵席的\Quote{共有}。好像他曾说：\Quote{如果是你们主人的，确实如此，你们不应该私用，而是作为你们主和主人的，摆在众人面前共有。因为这就是'主的'的意思。但现在你不允许它是主的，不允许它是共有的，而是独自宴饮。}因此他继续说：\par

\BibleQuote{各人预先取用自己的晚餐}。他没有说\Quote{切断}，而是说\Quote{预先取用}，暗中指责他们的贪食和急躁。至少下文也显示出这一点。因为说了这话后，他又加上：\BibleQuote{一个饥饿，另一个醉酒}，两者都表明缺乏节制，无论是贪求还是过度。再看另一个过错，同样使那些人受损害：第一，他们羞辱了自己的晚餐；第二，他们贪食醉酒；更糟的是，当穷人饥饿时。因为本来要摆在众人面前共享的，这些人独享了，并且达到饱胀和醉酒。因此他也没有说：\Quote{一个饥饿，另一个饱足；}而是说：\Quote{醉酒}。这两者每一个本身都该受责备：因为即使不轻视穷人而醉酒，也是过错；而不醉酒却轻视穷人，也是一个指控。当两者同时结合时，试想这罪过是何等重大。\par

接着，在指出他们的亵渎后，他非常愤怒地在以下话语中加上斥责，说：\par

% structure: p0030 source=h2 target=subsection reason=relative-heading-rank; base=h2

\subsection{\BibleRef{格前 11:22}}

\BibleQuote{什么？你们没有房子可以吃喝吗？还是你们藐视天主的教会，使那没有的羞愧呢？}\par

你看见他是如何将指控从对穷人的侮辱转移到教会，使他的话更令人厌恶吗？现在你看这是第四个指控，不仅穷人，连教会也被侮辱。因为正如你使主的晚餐成为私人宴席，同样地，你又使用这个场所，将教会当作房子。因为教会成为教会，不是要我们这些聚集的人被分裂，而是要那些分裂的人被联合：这一点聚集的行动表明了。\par

\Quote{使那没有的羞愧。}他没有说\Quote{饿死那没有的，}而是说\Quote{使他们羞愧}，从而更加令他们脸红；以指明他所关心的不是食物，而是对他们的伤害。看，又是第五个指控，不仅忽视穷人，甚至羞辱他们。他说这话，部分是由于尊重穷人的事，暗示他们不为肚子而悲伤，却为羞辱而悲伤；部分也是引起听者的同情。\par

因此，在指出了如此多的不虔敬——对晚餐的侮辱、对教会的侮辱、对穷人的藐视——之后，他再次缓和了斥责的语调，突然说道：\Quote{我该称赞你们吗？在这事上我不称赞你们。}其中人们尤其会惊奇，当需要更严厉地打击和斥责之后，证明了如此大的过犯，他反而做了相反的事，让步，让他们喘息。那么原因可能是什么？他在加重指控时比平时更痛苦地触及了伤口，而作为最高明的医生，他使切口适应伤口，既不浅割那些需要深刀的地方（因为你们听过他如何切除了那些人中犯奸淫的人），也不把需要更温和治疗的东西交给刀。因此，在这里他也更温和地讲话，而且从另一个角度看，他特别寻求使他们对穷人温和：这就是为什么他用一种更柔和的语气与他们谈论。\par

5. 接着，为了从另一个话题更加羞辱他们，他又取了那些最要点，并藉着它们编织他的讲论。\par

% structure: p0036 source=h2 target=subsection reason=relative-heading-rank; base=h2

\subsection{\BibleRef{格前 11:23}}

\BibleQuote{我原是从主领受的}，他说，\BibleQuote{我也传授给你们了：就是主耶稣在被出卖的那一夜，拿起饼来：}\par

% structure: p0038 source=h2 target=subsection reason=relative-heading-rank; base=h2

\subsection{\BibleRef{格前 11:24}}

\BibleQuote{他祝谢了，就擘开，说：你们拿去吃，这是我的身体，为你们而擘开的；你们应这样做，为纪念我。}\par

所以，他为何在此提及这奥迹？因为这个论点对于他当前的目的十分必要。就像这样：\Quote{你们的主，}他说，\Quote{认为所有人都配享同一宴席，尽管它非常可畏且远超所有人的尊贵；但你们却认为他们连你们自己的宴席都不配，而你们的宴席如我们所见的又小又卑贱；并且，他们在属灵的事上并不比你们优越，你们却在属世的事上掠夺他们。因为连这些也不是你们的。}\par

然而，他并没有这样说，以免他的讲论变得严厉；而是以更温和的方式表述说：\BibleQuote{主耶稣在被出卖的那一夜，拿起饼来。}\par

他为何也要我们记起那个时间、那个晚上和那出卖？并非偶然或无理由，而是为了以时间本身为例，极大地充满他们的痛悔。因为即使是一块石头，当他想到那一夜，耶稣如何与门徒在一起，\Quote{非常忧郁，}如何被出卖，如何被捆绑，如何被带走，如何受审，如何依次承受其余的一切，他就会比蜡还软，并从尘世和这世界的一切浮华中超脱出来。因此，他通过时间、他的餐桌和他的出卖，引导我们想起这一切，使我们羞愧，并说：\Quote{你们的主甚至为你舍弃了自己；而你却连一点食物也不愿为你的兄弟分享，为了你自身的缘故。}\par

但是，他怎会说\BibleQuote{他是从主领受的？}因为他当时并不在场，反而是迫害者之一。这是为使你们知道，第一张桌子并不比后来的有更多优越。因为即使是今天，也是他做这一切，并像那时一样分授。\par

他提醒我们那一夜，不仅为此，也为以另一种方式使我们痛悔。因为正如我们特别记得从临终者耳边听到的最后的话；如果他们的继承人竟敢违背他们的命令，我们想要使他们羞愧时，就说：\Quote{想想这是你父亲对你说的最后一句话，直到他临咽气的晚上，他还在重复这些吩咐。}同样，保禄为了使他的论证更充满威严，说：\Quote{请记住，}他说，\Quote{这是他赐给你们的最后一项奥礼，在他即将为我们被杀的那一夜，他命令了这些事，并在将晚餐传给你们之后，再也没有增添什么。}\par

接着，他继续叙述当时所行的事，说：\BibleQuote{他拿起饼，祝谢了，就擘开，说：你们拿去吃，这是我的身体，为你们而擘开的。}因此，如果你们前来是为感恩的祭献，你们自己就不要做任何与这祭献不相称的事：绝不可侮辱你们的兄弟，也不可忽视他的饥饿；不可醉酒，不可侮辱教会。既然你们来是为你们所享有的而感恩，你们自己也当相应地回报，而不与邻人隔绝。因为基督一方平等地给了众人，说：\BibleQuote{你们拿去吃。}他平等地给了他的身体，但你连普通的饼都不平等地分给吗？是的，它确实为众人一同擘开，并成了众人同等的身体。\par

% structure: p0046 source=h2 target=subsection reason=relative-heading-rank; base=h2

\subsection{\BibleRef{格前 11:25}}

\BibleQuote{晚餐后，他也照样拿起杯来，说：这杯是用我的血所立的新盟约；你们每次喝它时，应这样做，为纪念我。}\par

你说什么？你们在纪念基督，却轻视穷人，并且不战栗吗？为什么，如果一个儿子或兄弟死了，你们在纪念他，如果没有履行习俗并邀请穷人，你们会受良心的谴责；而当你们纪念你们的主时，你们难道连简单分给宴席的一份也不给吗？\par

但他所说\BibleQuote{这杯是新盟约吗？}是什么意思？因为也有旧盟约的杯，就是奠祭和牲畜的血。因为献祭后，他们习惯用圣爵和碗接血，然后倒出。既然他用自己的血代替了牲畜的血；唯恐有人听到这个会不安，他就提醒他们那古老的祭献。\par

6. 接着，在讲论了那晚餐之后，他将现在的事与那时的事联系起来，使得就像在那个晚上，斜靠在那个榻上，并从基督自己手中领受这祭献一样，现在的人们也能同样受感动；于是他说：\par

% structure: p0051 source=h2 target=subsection reason=relative-heading-rank; base=h2

\subsection{\BibleRef{格前 11:26}}

\BibleQuote{因为你们每次吃这饼，喝这杯，就是宣告主的死亡，直到他再来。}\par

因为基督关于饼和杯说过：\BibleQuote{你们这样做，为纪念我。}向我们揭示了赐予这奥迹的原因，并且除了他所说的之外，还宣告这是一个足够的基础来建立我们的宗教敬畏：——（因为当你考虑你的主为你所受的苦时，你会更好地弃绝自己）——同样，保禄在这里也说：\BibleQuote{你们每次吃，就是宣告他的死亡。}而这正是那晚餐。然后，暗示这要持续到终结，他说：\BibleQuote{直到他再来。}\par

% structure: p0054 source=h2 target=subsection reason=relative-heading-rank; base=h2

\subsection{\BibleRef{格前 11:27}}

\BibleQuote{所以，无论何人，若不配地吃这饼，喝这主的杯，就是干犯主的身体和主的血。}\par

何以如此？因为他浇奠了，并使这事看起来像一场屠杀，不再是祭献了。因此，正如当时那些刺透祂的人，他们刺透祂并非为饮，而是为流祂的血：同样，那不相称地前来领受且毫无所获的人也是如此。你们看见祂如何使祂的讲论令人战栗，并极其严厉地谴责他们，表明如果他们这样喝，就是不相称地分受这元素？因为若他不顾饥饿的人，且轻慢他，又怎能不是不相称呢？既然不施舍穷人，即使是一童贞女，也将被逐出天国；或者更确切地说，不慷慨施舍：（因为那些童贞女也有油，只是不充足而已：）那么，行了如此多不虔之事，这恶将何等巨大，试想吧！\par

\Quote{什么不虔之事？}你说。你为何说\Quote{什么不虔之事}？你分受了这样的筵席，理当比任何人都更温柔，如同天使一般，却无人像你这样残忍。你尝了主的血，竟仍不认识你的兄弟。那你还配得什么宽恕呢？倘若此前你还不认识他，你本该从这筵席认识他；但现在你竟侮辱了这筵席本身：他已被认为配得分受这筵席，而你却不认为他配得你的食物。你没有听过那要求一百个银钱的人受了多大的苦吗？他如何使赐予他的恩赐归于无效？你难道不思想你从前是什么、如今成了什么？你不提醒自己：若此人在财物上贫穷，你在善工上岂不更加贫乏，充满万般的罪恶？然而，天主解救你脱离这一切，并使你配得这样的筵席；但你甚至因此仍未变得更加慈悲：因此，当然别无他途，只有你被\Quote{交付给刑役}了。\par

7. 我们众人也要听这些话，凡在这地方同穷人一起亲近这圣筵，但出去时却仿佛没有看见他们，反而醉酒、漠视饥饿者的人，就是格林多人所被控告的。这事何时发生呢？你说。确实时时发生，但在节庆之日尤其如此，那时最不应当如此。难道不是吗？在这些时候，领圣体之后立即继以醉酒和轻视穷人。你既分受了血，本该守斋和警醒，却沉溺于酒和狂欢。你若偶然在早晨吃了什么好东西，便小心翼翼，以免别的难吃菜肴败坏前味；如今你以圣神为筵，却招来了撒殚式的奢侈。试想：宗徒们分受那圣晚餐后做了什么？他们不是去祈祷、唱圣咏吗？不是守神圣的夜课吗？不是从事那漫长而充满克己的教导工作吗？因为那时，当犹达斯出去召唤那些要钉死祂的人时，祂便向他们讲述并交付了那些伟大而奇妙的事。你没有听说那三千人也领了共融，继续祈祷和听道，而非在醉酒和狂欢中吗？但你在领受前守斋，以某种方式显得配得共融；而领受后，你本应更加节制，却反而前功尽弃。然而，在这之前和之后守斋，确实不同。虽然两时都当节制，但尤其在领受了新郎之后更当如此。之前，为能配得领受；之后，为不显得所领受的不配。\par

\Quote{那么，我们该在领受后守斋吗？}我并非说这话，也不施加强制。这固然好；然而，我不强求此，只劝你们不要纵情宴乐。因为若有人永远不该奢侈度日，保禄已表明此点，他说：\Quote{那好宴乐的女人，虽活着也是死的} \BibleRef{弟前 5:6}；那么她那时更将是死的。若奢侈对女人是死亡，对男人更是如此；若在其他时候行此是致命的，那么在领受奥迹之后就更甚了。你既拿了生命之粮，却行死亡之举，难道不战栗吗？你不知奢侈带来多少大恶？不合时宜的欢笑、无礼的言论、充满毁灭的戏谑、无益的闲谈，以及一切不配提及之事。你在享受基督的筵席之日，在那你被堪得以舌头触祂血肉的日子，竟行这些事。那么，如何防止这些事呢？洁净你的右手、你的舌头、你的嘴唇，它们成了基督践踏的门槛。试想你亲近并摆出物质筵席的时刻，将你的心提升到那筵席、主的晚餐、门徒的夜课，那一夜，那圣洁之夜。不，更确切地说，若详察，此刻正是黑夜。那么，让我们与主一同警醒，与门徒一同痛心。这是祈祷而非醉酒的时刻，永远如此，尤其在节庆期间。因为节庆被设立，不是为使我们行为不端，也不是为累积罪恶，而是为涂抹已有的罪恶。\par

我知道，说这些话是徒然的，但我仍将不停地说。因为即使你们不全听从，也必不会全不听从；或者，即使你们都不听从，我的赏报更大，而你们的定罪则更重。然而，为使定罪不更重，为此我将不停地说。因为或许，或许，因着我的坚持，我能触及你们。\par

为此，我恳求你们，不要使这事归于定罪；我们要喂养基督，给祂喝，给祂穿。这些事才配得上那祭台。你们听到了神圣的赞美诗吗？你们看见了神婚吗？你们享受了君王的宴席吗？你们被圣神充满了吗？你们加入了色辣芬的唱诗班吗？你们成为了高天权势的分享者吗？不要抛弃如此大的喜乐，不要浪费宝藏，不要带来醉酒——沮丧之母、魔鬼的欢乐、万恶之源。因为由此而来的睡眠如同死亡，头痛、疾病、健忘，以及死人状态的形象。再者，如果你不愿意醉酒时遇见朋友，那么当你内心有基督时，告诉我，你怎敢将如此过度的东西强加给祂？\par

但是你喜欢享乐吗？那么就因此停止醉酒。因为我也愿意你享受自己，但却是那真正的享受，永不褪色的。那么真正的、常青的享受是什么呢？邀请基督与你共进晚餐\BibleRef{默 2:20}；让祂分享你的，更确切地说，分享祂自己的。这带来无限的快乐，并永远保持其青春。但感官之事并非如此；它们一出现就消逝；享受它们的人与没有享受的人相比，处境不会更好，甚至更糟。因为前者如同停泊在港湾里，而后者却将自己暴露在一种激流、疾病围攻的军队中，甚至没有力量承受第一次海浪。\par

因此，但愿这些事不如此，让我们追求节制。这样我们将身体健康，灵魂安妥，并脱离现世和来世的邪恶：愿我们所有人都能脱离这些邪恶，达到天国，藉着我们的主耶稣基督的恩宠和仁慈，愿光荣、权能、尊崇归于祂和圣父及圣神，从现在直到永远，永世无穷。阿们。\par

\SourceRecord{Translated by Talbot W. Chambers.；Nicene and Post-Nicene Fathers, First Series；Vol. 12.；Edited by Philip Schaff.；Buffalo, NY: Christian Literature Publishing Co.,；1889.；<http://www.newadvent.org/fathers/220127.htm>.}

% structure: p0001 source=h1 target=section reason=page-title

\section{格林多前书讲道集 第二十八篇}

% structure: p0002 source=h2 target=subsection reason=relative-heading-rank; base=h2

\subsection{\BibleBook{格林多}{前书} 11:28}

\BibleQuote{人应省察自己，然后才吃这饼，喝这杯。}\par

这些话是什么意思，当我们面前另有题目时？这是\Person{保禄}的习惯，如我先前所说，他不仅处理自己所拟定的主题，而且若出现与他的目的相关联的论点，也会极其勤勉地展开，尤其当它涉及非常必要和紧迫的事时。因此，当他与已婚者谈论时，关于仆人的问题出现了，他就极为有力且详尽地处理了它。又如，当他论及不可在外邦法庭起诉的义务时，又碰到了关于贪心的劝诫，他也就此主题长篇论述。现在，他在这里也做了同样的事：一旦有机会提醒他们关于奥迹的事，他就认为有必要继续那个主题。因为这确实不是一件寻常的事。所以他极其庄重地论及它，为万善之总，即他们以纯洁的良心接近那些奥迹。因此，他不满足于先前所说的话，还加上了这些话，说：\par

\Quote{但人应省察自己：}这也是他在第二封信中所说的：\Quote{你们考验自己，证明自己：}\BibleRef{格后 13:5}不像我们如今这样，因着节期而非出于心灵的真诚而接近。因为我们不考虑如何预备好、清除内在的恶、充满痛悔而接近，只考虑如何在节期及众人皆如此时来到。但\Person{保禄}并未命我们如此前来：他只知一个接近和共融的时节，即人良心的纯洁。因为即使感官所认知的那种盛筵，当我们发烧且充满不良体液时，也不能无死亡危险地享用；更何况我们若以亵渎的情欲（比发烧更严重）触到这一祭台，更是不可容许的。当我说亵渎的情欲时，我指的是身体的、金钱的、愤怒的、恶意的，总之，一切亵渎的情欲。接近者应先使自己脱离这一切，然后才触到那纯洁的祭献。既不应因懒惰和勉强而被节期迫使接近；反之，若真心悔改且预备好了，也不应因不是节期而被阻止。因为节期是善工的展示、灵魂的敬畏、行为的严谨。如果你拥有这些，你随时都可以守节，随时都可以接近。因此他说：\Quote{但各人应省察自己，然后才可接近。}他命令的不是一人审查另一人，而是各人审查自己，使法庭非公开，定罪也无需见证。\par

% structure: p0006 source=h2 target=subsection reason=relative-heading-rank; base=h2

\subsection{\BibleBook{格林多}{前书} 11:29}

2. \Quote{因为那不相称地吃和喝的人，就是吃喝自己的审判。}\par

你说什么呢？请告诉我。这赐予如此多祝福并充满生命的祭台，竟成了审判？不是出于它自身的本性，他说，而是出于接近者的意志。因为正如祂的临在——那带给我们大而不可言喻的祝福——更定那不接受之人的罪；同样，这些奥迹也成为更重惩罚的供应，给那些不相称地领受的人。\par

但他为何吃自己的审判？\Quote{不分辨主的身体：}即不探究、不铭记他面前所陈之物的伟大，如他所应做的；不衡量恩赐的分量。因为如果你准确知道在你面前的是谁，以及那将自己赐予的是谁，又赐给谁，你便不再需要其他辩论，这足以使你全力以赴地警惕；除非你已完全堕落。\par

% structure: p0010 source=h2 target=subsection reason=relative-heading-rank; base=h2

\subsection{\BibleBook{格林多}{前书} 11:30}

\Quote{因此，你们中有许多是软弱的、患病的，不少人也睡了。}\par

这里他不再像论及祭偶像之物时那样从别人引例，叙述古史和旷野中的惩罚，而是从格林多人自己身上引例；这也使得讲论更能刺中他们。因为当他说\Quote{他吃自己的审判}和\Quote{他是有罪的}时，为使他的话不像是空言，他也指向事实，并叫他们自己作见证；这是一种比恐吓更能打动人的方式，表明恐吓已见诸实际。但他并不满足于此，还由此引出并证实关于地狱之火的论点，从两方面恐吓他们；并且解决一个到处被问及的问题。我是说，既然许多人彼此追问：\Quote{为何有不期而遇的死亡？为何有长久不愈的疾病？}他就告诉他们，这些意外事件很多是由于某些罪过。\Quote{那么，那些一直健康，}你们说，\Quote{并安享高寿的人，他们不犯罪吗？}谁敢于这样说？\Quote{那么，}你们说，\Quote{他们为何不受惩罚？}因为他们在那里要遭受更重的惩罚。但我们若愿意，既不用在这里，也不用在那里受罚。\par

% structure: p0013 source=h2 target=subsection reason=relative-heading-rank; base=h2

\subsection{\BibleBook{格林多}{前书} 11:31}

\Quote{因为如果我们省察自己，}他说，\Quote{就不会受审判。}\par

他没有说：\Quote{我们若惩罚自己，若对自己报复，} 而是说只要我们愿意承认我们的过犯，给自己定罪，谴责所行不当的事，我们就可摆脱今世和来世的惩罚。因为那谴责自己的，以两种方式平息天主：既承认自己的罪，又更加小心将来。但既然我们连这轻微的事也不愿意按应当的去做，祂甚至也不容忍我们与世界一同受罚，反而这样宽容我们，只在今世实施惩罚，此处的刑罚是暂时的，安慰却是巨大的；因为结果既是从罪中释放，又对将来有美好的希望，减轻现世的痛苦。他说这些事，既安慰患病者，也使其余的人更加庄重。因此他说：\par

% structure: p0016 source=h2 target=subsection reason=relative-heading-rank; base=h2

\subsection{\BibleRef{格前 11:32}}

\Quote{但我们受审判时，是受主的惩戒。}\par

他没有说我们受惩罚，他也没有说我们被报复，而是说：\Quote{我们受惩戒。} 因为所发生的更属于劝诫而非定罪，属于医治而非报复，属于纠正而非惩罚。不仅如此，他还借着更大邪恶的威胁，使现在的（惩戒）显得轻微，说：\Quote{免得我们和世界一同被定罪。} 你看见他如何也提到了地狱和那可怕的审判台，并表明那试炼和惩罚是必要且无论如何必须的吗？因为如果信友们，就是天主特别关心的人，在一切过犯中都不能免除惩罚（这从现世的事上可以看出），更何况那些不信的人和犯下不可赦免、无法医治之罪的人呢？\par

% structure: p0019 source=h2 target=subsection reason=relative-heading-rank; base=h2

\subsection{\BibleRef{格前 11:33}}

3. \Quote{所以，你们聚集一起吃饭时，要彼此等待。}\par

这样，当他们的恐惧仍高涨，地狱的恐怖仍存留时，他再次选择引入为穷人的劝勉，正为此他说了这一切事；暗示如果他们不这样做，就必须不相称地领受。但如果不分施我们的财物就使我们被排除在那筵席之外，那么强取就更甚了。他没有说：\Quote{所以，当你们聚集时，分施给那些需要的人，} 而是用了更有敬意的说法：\Quote{要彼此等待。} 因为这也为那件事铺路并暗示了它，并以合适的形式引入了劝勉。然后为进一步羞辱他们，\par

% structure: p0022 source=h2 target=subsection reason=relative-heading-rank; base=h2

\subsection{\BibleRef{格前 11:34}}

\Quote{若有人饥饿，在家里吃好了。}\par

通过允许，他阻止了这事，比完全禁止更有力。因为他把他带出教会，送回家中，借此严厉谴责并嘲笑他们，如同肚腹的奴隶，不能自制。因为他没有说：\Quote{若有人藐视穷人，} 而是说：\Quote{若有人饥饿，} 如同对没有耐心的儿童说话；如同沦为食欲奴隶的畜牲。因为如果因他们饥饿就让他们在家中吃，那的确是极其可笑的。\par

然而他不满足于此，还添加了另一件更可怕的事，说：\Quote{你们的聚集不成为审判，} 即你们不得受到惩戒、惩罚，侮辱教会，羞辱你们的弟兄。\Quote{因为你们聚集正是为此，} 他说，\Quote{为能彼此相爱，彼此得益。但如果发生相反的事，你们宁可在家中喂养自己。}\par

但他说这话，是为了更吸引他们归向他。是的，这正是他的目的：既指出由此产生的伤害，又说明定罪非小事，并以各种方式恐吓他们：借着奥迹、患病者、已死的人以及以前提到过的其他事。\par

然后他也以另一种方式再次警告他们，说：\Quote{其余的事，等我来了再安排。} 这或许指其他事，或许就是指这件事。因为很可能他们还有些理由要提出，而不可能通过一封信纠正一切，\Quote{我所吩咐你们的，暂且遵守，} 他说，\Quote{但你们若还有什么事要提，等我来了再说。} 他说这事，或指这件事，如我说过的，或指其他不很紧急的事。他这样行，为从此也能使他们更庄重。因为渴望他的到来，他们就会纠正错误。因为保禄在任何地方的暂居都不是平凡的事：为了表明这一点，他说：\Quote{有些人自高自大，以为我不会到你们那里去。} \BibleRef{格前 4:18} 又在别处说：\Quote{不但我在你们那里时，如今我不在你们那里时更当如此，怀着恐惧战栗，成就你们的救恩。} \BibleRef{斐 2:12} 因此他不仅只是承诺他会来，免得他们不信而更加疏忽；他还陈述了一个他暂居他们那里的必要原因，说：其余的事，等我来了再安排。这暗示着，即使他原本并不渴望，纠正所剩的事也会将他吸引到那里。\par

4. 因此，听了这一切之后，让我们既好好照顾穷人，又节制食欲，摆脱醉酒，小心相称地领受奥迹；无论我们遭受什么，不要为此而痛苦，既不为己也不为人；如遇上夭亡或久病。因为这是摆脱惩罚，这是纠正，这是极好的劝诫。谁说这话？那有基督在他内说话的人。\par

然而即使在这之后，我们许多妇女仍如此愚昧，竟在过度悲伤中甚至超过不信者。有些人因情欲盲目而如此行，但另一些人是为炫耀，并避免外面的人的指责：依我看来，她们最没有借口。因为，\Quote{免得某人指责我，} 她说，\Quote{让天主指责我；免得比畜牲更无知的人定我的罪，让万王之王的律法被践踏。} 啊，这些话岂不应遭雷击多次？\par

再者：若有人在你们遭难后邀请你们赴丧宴，没有人会说什么，因为人的法律规定了这事；但当天主藉祂的法律禁止你们哀悼时，众人却都反对这事。啊，妇人，约伯难道没有进入你的心中吗？你难道不记得他在子女遭祸时所说的话，那些话比万条冠冕更装饰那神圣的头颅，比许多号角更响亮地宣告吗？你难道不重视他灾祸的严重性，那史无前例的沉船，那奇异而可怖的悲剧吗？因为你可能失去一个，或两个，或三个；但他失去了那么多儿女：那个多子的人突然变得无子。他的脏腑不是逐渐被耗尽的；而是一下子，他身体的所有果实都被夺走了。也不是按照自然的共同法则，当他们年老时，而是以既不合时宜又暴力的死亡：而且全部一起，当他不在场，也没有坐在他们旁边，至少让他们能听见临终之言，对他如此悲惨的结局稍得安慰；但出乎一切意料，他对发生的事一无所知，他们全都同时被压倒，他们的房屋成了他们的坟墓和罗网。\par

不仅他们不合时宜的死亡，还有许多其他事令他悲伤：例如他们都在青春年华，全都品德优良、彼此相爱，全都一起，没有留下一个男女，这事不是按自然的共同法则发生，是在如此巨大的损失之后，当他既不觉得自己有罪，也不觉得他们有罪时，他遭受了这些。因为每一种情况本身都足以扰乱心神；但当我们发现它们同时发生时，请想象那浪涛的高度，那风暴的猛烈程度。而尤其比他的丧子更大更坏的是，他甚至不知道这一切为何发生。因此，既然无法为不幸找到原因，他便上升至天主的善意，说：\Quote{主给了，主收回：}正如主所喜悦的，就这样发生了；\Quote{愿主的名永远受赞美。}\BibleRef{约 2:21}当他看见自己，这个追随一切美德的人处于极端困境中，而恶人和骗子却四处享乐、奢华、狂欢时，他说了这些话。他没有说出一些软弱者可能会说的话，如：\Quote{难道我养育儿女，严格训练他们，就是为了这个吗？难道我向所有过路人敞开家门，为了穷困者、赤身者、孤儿奔走多次之后，竟得到这样的回报吗？}相反，他献上了比一切祭献更美的那些话，说：\Quote{我赤身出离母胎，也必赤身归回那里。}然而，若他撕裂衣服、剃了头，不要惊奇。因为他是一位父亲，一位慈爱的父亲：既要显露天性的怜悯，也要显出精神的克制。倘若他没有这样做，或许有人会认为这种克制不过是麻木。因此，他既显出了自然亲情，也显出了虔诚的严谨，而在悲痛中他没有被击倒。\par

5. 是的，甚至当他的试炼进一步进行时，他因回答他的妻子而再次戴上其他冠冕，说：\Quote{我们从主手中得了福，难道不也当忍受祸吗？}\BibleRef{约 2:10}事实上，那时他的妻子是唯一存留下来的，他的一切都已彻底毁灭——子女、财产、甚至他的身体——而她被保留来诱惑和陷害他。这确实是魔鬼不把她与子女一同毁灭、也不求她死亡的原因，因为他预料她会大大有助于陷害那位圣人。因此他把她留作一件工具，一件可怕的工具，供自己使用。\Quote{因为如果我藉着她甚至从乐园里赶出了人类，}他说，\Quote{那么我更能藉着她在粪堆上绊倒他。}\par

请看他的诡计。当牛、驴、骆驼丢失时，他没有使用这个策略，甚至当房子倒塌、子女被埋时也没有，而是如此长久地注视着那位战士，他让她沉默安静。但当虫泉涌出，皮肤开始腐烂脱落，肉体败坏发出最讨厌的分泌物，魔鬼的手以比烤架、火炉和任何火焰更尖锐的痛苦折磨他，四面消耗、吞噬他的身体，比任何野兽更严重，并且在这苦难中过了很长时间；那时他把他带到她面前，他已受尽折磨、疲惫不堪。倘若她在他的不幸一开始就接近他，她既不会发现他如此软弱，也不能用她的话如此夸大那不幸。但现在，当她看到他因时间长久而渴望解脱，渴望结束这紧逼着他的痛苦时，她才来攻击他。为表明他已被完全耗尽，到那时甚至不能呼吸，并且甚至渴望死去，请听他所说的话：\Quote{我但愿我能亲手了结自己，或请求别人替我这样做；}请看，我祈祷，他妻子的邪恶，她从什么话题开始：即从时间的长久，说：\Quote{你要坚持到何时？}\par

如今，即使没有事实，单凭言语就足以使人失去男子气概，那么请想一想，当他听到这些话时，除了言语本身，还有事实本身也在折磨他，他该作何感想？尤其最糟糕的是，说这话的竟是他的妻子，而且是一个完全堕落、自暴自弃的妻子，因此她也要把他拖入绝望。然而，为了更清楚地看到那攻击磐石之墙的器械，让我们听听她的话。那么这些话是什么呢？\Quote{你还要坚持多久？你说：看，我再等一会儿，期待我救恩的望德。} \Quote{不，}她说，\Quote{时间已经暴露了你话的愚蠢，虽然拖延，却未显示任何逃脱之道。}她说这些话，不仅是要把他推入绝望，而且是在嘲笑和戏弄他。\par

因为他总是安慰步步紧逼的她，敷衍她说：\Quote{再等一会儿，这些事很快就要结束了。}因此她责备他说：\Quote{你现在又要说同样的话吗？因为很长时间已经过去了，这些事的终结还没有出现。}请注意她的恶意：她绝口不提牛、羊或骆驼，因为她知道他对这些并不十分烦恼；而是直接指向本性，提醒他他的孩子们。因为当他看见孩子们死去时，他曾撕裂衣服、剃光头发。她不说：\Quote{你的孩子们死了，}而是非常凄惨地说：\Quote{你的纪念从地上消灭了，}\Quote{你的孩子们之所以被渴望的理由。}因为即使在复活已经显明的今天，人们仍然渴望孩子，因为他们保存了逝者的记忆；更何况在当时呢？因此，她的诅咒也因此更加苦涩。因为在那情况下，诅咒者不说：\Quote{愿他的子女被连根拔除，}而是\Quote{愿他的纪念从地上消灭。}\Quote{你的儿女。}因此，虽然她说\Quote{纪念，}但她又准确地提到了两性。\Quote{但如果你，}她说，\Quote{不关心这些，至少想想我的境况。}\Quote{我腹中的疼痛，以及我枉然在忧愁中所忍受的劳苦。}她的意思是：\Quote{我忍受了更多的痛苦，却为了你的缘故而受委屈，承受了辛劳，却被剥夺了果实。}\par

请看，她既未明确提及他财产的损失，也未对此避而不谈、匆匆带过；而是以最凄惨的方式，巧妙地暗示了这一点。因为当她说：\Quote{我也成了流浪者和奴隶，从一个地方到另一个地方，从一家到另一家，}她既暗示了财产损失，也指出了她巨大的痛苦；这些表达甚至加深了那不幸。\Quote{因为我来到别人的门前，}她说；\Quote{我不仅乞讨，而且四处漂泊，服侍一种奇异而不寻常的奴役，到处走动，带着我灾难的标记，向所有人述说我的痛苦；}其中最可怜的是换了一家又一家。她甚至没有停留在这些哀叹上，而是继续说：\Quote{等待太阳落山，那时我将从环绕我的痛苦和患难中得安息，这些痛苦现在正困迫我。}\Quote{因此，}她说，\Quote{对别人来说甜蜜的事——看见光明，对我来说却是痛苦的；而黑夜和黑暗是可渴望的。因为只有这能使我从劳苦中得安息，这成了我痛苦的安慰。但是，你要说些反对主的话，然后死去。}你在这里也看出她的狡猾和邪恶吗？她如何不是立即提出致命的建议，而是先可怜地叙述她的不幸，并拉长悲剧，然后用几句话提出她的建议，甚至没有明确说出来，而是遮掩着，向他展示他极其渴望的解脱，并承诺死亡——这正是他当时最渴望的东西。\par

从这一点也可以看出魔鬼的恶意：因为他知道约伯对天主的渴望，所以不让他的妻子指控天主，免得约伯立即把她当作敌人而避开。因此，她从未提及天主，而是不断地念叨着实际的灾祸。\par

除了以上所说的，你还要加上这一点：提出这建议的是一个女人，一个善于诱惑粗心者的雄辩家。至少有许多人即使没有外在的变故，也仅因女人的建议而跌倒。\par

6. 那么，这位蒙福的圣人，比磐石更坚定的人，做了什么呢？他狠狠地瞪着她，甚至在说话之前，就用表情击退了她的诡计；因为她无疑期望会激起泪泉，但他却变得比狮子更凶猛，满怀着愤怒和义愤——不是因为他的痛苦，而是因为她魔鬼般的建议；他先用眼神表达了他的愤怒，然后用克制的语气给出了责备；因为即使在不幸中，他也保持自制。他说了什么呢？\Quote{你为何说话像愚顽的妇人？}\Quote{我没有这样教导过你，}他说，\Quote{我没有这样培养过你；这就是为什么我现在甚至认不出我自己的配偶。因为这些话是\InnerQuote{愚顽妇人}的建议，是癫狂之人的建议。}你在这里不是看到了适度的伤害，仅仅是足以治愈疾病的打击吗？\par

然后，在灾祸之后，他又提出了足以安慰她的劝告，且非常合理，这样说：\Quote{我们既然从主手里接受了福乐，难道不能忍受灾祸吗？} \Quote{请回想，}他说，\Quote{那些从前的事，并想到那赐予者，你就能高贵地承受这些了。} 你看到这人的谦逊了吗？他完全不将自己的忍耐归功于自己的勇气，而是说这是事情自然结果的一部分。\Quote{因为我们从前得到这些，天主给了我们什么回报？祂偿付了什么报酬？没有，纯粹出于美善。因为那些是恩赐，不是报酬；是恩宠，不是赏报。那么，让我们也高贵地承受这些吧。}\par

这番言论，我们无论男女，都应当记录下来，并将这些话铭刻在我们心中，包括这些和之前那些话：并通过在我们的心中如同绘画般描摹他们的苦难史，我指的是失去财富、丧子之痛、身体疾病、辱骂、嘲笑、他妻子的诡计、魔鬼的陷阱，总而言之，那位义人的一切灾难，并且要精确地描摹，好让我们为自己预备一个最宽广的避难港：好使我们高贵地、感恩地承受一切，从而在此生摆脱一切沮丧，并领受与这种良好态度相应的赏报；借着我们的主耶稣基督的恩宠和怜悯，愿光荣、权能和尊威归于祂，偕同圣父及圣神，自今而永远，永世无尽。阿们。\par

\SourceRecord{Translated by Talbot W. Chambers.；Nicene and Post-Nicene Fathers, First Series；Vol. 12.；Edited by Philip Schaff.；Buffalo, NY: Christian Literature Publishing Co.,；1889.；<http://www.newadvent.org/fathers/220128.htm>.}

% structure: p0001 source=h1 target=section reason=page-title

\section{\WorkTitle{格林多前书第十二章讲道词 第廿九篇}}

% structure: p0002 source=h2 target=subsection reason=relative-heading-rank; base=h2

\subsection{\BibleRef{格前 12:1-2}}

\Quote{弟兄们，论到\OriginalTerm{神恩}{spiritual gifts}，我不愿意你们不明白。你们知道，当你们还是\Person{外邦人}的时候，如何被那些哑巴\OriginalTerm{偶像}{idols}所牵引，无论怎样被牵引。}\par

这段经文非常隐晦，但这隐晦是由于我们对所提及的事实已经生疏，并且这些事实已经停止；它们过去常常发生，但现在不再发生了。为什么它们不再发生了？看，这隐晦又给我们引出另一个问题：就是为什么它们那时发生，而现在不再发生了？\par

不过我们暂且把这个问题搁置，另找时间讨论；现在让我们先说明那时发生了什么事。那时究竟发生了什么呢？凡受洗的人，立刻就说\OriginalTerm{语言之恩}{tongues}；不仅说语言，而且许多人也说预言，还有人行许多其他的奇事。因为他们脱离偶像而来时，对古代的圣经并没有清楚的认识或训练，但他们一受洗就领受了圣神；然而圣神是看不见的，因为祂是无形体；因此天主的\OriginalTerm{恩宠}{grace}赐予了一些可感觉的凭证来证明那能力的运行。有些人立刻讲波斯语，有些人讲拉丁语，有些人讲印度语，还有人讲其他类似的语言；这使外人明白，正是圣神在说话者里面。因此保禄也如此称呼它，说：\BibleQuote{但圣神的显示赐给各人，原是为众人的益处。}\BibleRef{格前 12:7}他把恩赐称为\Quote{圣神的显示}。正如宗徒们首先领受了这个标记，信友们也继续领受它；我指的是语言之恩；不过不仅如此，还有许多其他的恩赐：许多人甚至能复活死人、驱魔，并行其他许多奇事；他们也有大大小小的各种恩赐。但其中最为丰富的是语言之恩；而这在他们中间成了分裂的原因——并非由于它本身的本质，而是由于领受者的乖张：那些拥有较大恩赐的人轻视那些拥有较小恩赐的人；而这些较小恩赐的人又感到悲伤，并嫉妒那些拥有较大恩赐的人。保禄自己在后面的论述中也暗示了这一点。\par

由于这件事使他们爱德破裂，受到致命打击，他非常注意要纠正它。因为在罗马也曾发生这种事，不过方式不同。这就是为什么他在罗马书中提起这事，但隐晦而简短地说：\BibleQuote{就如我们在一个身体上有许多肢体，但肢体并不都有同样的作用；同样，我们众人在基督内，也都是一个身体，彼此之间，每个都是肢体。按我们各人所受的\OriginalTerm{恩宠}{grace}，各有不同的恩赐：如果是预言之恩，就应与信德相称；如果是服务，就当服务；如果是教导的，就当教导。}\BibleRef{罗 12:4}而且，罗马人也因此陷入骄傲，他在那篇讲论的开头就暗示了这一点，这样说：\BibleQuote{我因所赐给我的恩宠，告诉你们每一位：不可把自己估计得太高，而过了分；但应按照天主所分与各人的信德尺度，估计得适中。}\BibleRef{罗 12:3}对于这些人，因为分裂和骄傲的疾病尚未蔓延，他如此谈论；但在这里，他怀着极大的焦虑，因为这恶疾已经大大扩散了。\par

而且，不仅这件事困扰他们，而且当地还有很多占卜者，因为这座城市比平时更沉迷于希腊风俗；加上其他因素，也在他们中间引起了冒犯和混乱。这就是为什么他首先开始陈述\OriginalTerm{占卜}{soothsaying}和\OriginalTerm{预言之恩}{prophecy}之间的区别。为此原因，他们也领受了\OriginalTerm{辨别神类}{discerning of spirits}，以便辨别并知道谁是由纯洁的神说话，谁是由不洁的神说话。\par

因为不可能在说话当时从内部提供所发预言的证据（因为预言之恩的真实性不是在其说出时，而是在事件发生时得到证明）；并且不容易区分真先知和假先知（因为魔鬼，那被诅咒的，也进入了那些说预言的人\BibleRef{列上 22:23}，带来假先知，似乎他们也能预言未来）；此外，人很容易被欺骗，因为所说的话无法在当时受到检验，在预言所涉及的事件尚未发生之前（因为结局才证明假先知和真先知）。——为了使听众在结局之前不受欺骗，他给了他们一个标记，甚至在事件之前就能指示真假。因此，他按照这个顺序开始，并同样继续谈论恩赐，并纠正由此产生的争竞。不过现在，他首先开始谈论占卜者，这样说：\par

\Quote{弟兄们，论到\OriginalTerm{神恩}{spiritual gifts}，我不愿意你们不明白。}他把这些\OriginalTerm{标记}{signs}称为\InnerQuote{属神的}，因为它们是\OriginalTerm{圣神}{Spirit}单独的工作，人的努力对施行这样的奇事毫无贡献。他打算谈论这些，首先，如我所说，他陈述了占卜和预言之恩之间的区别，这样说：\par

\Quote{你们知道，当你们还是外邦人时，是怎样被牵引到那些哑吧偶像那里，无论你们如何被牵引。} 现在他这样说的意思是：\Quote{在偶像庙里，}他说，如果有人曾一度被不洁之神附身并开始行占卜，他就如同被拖走一样，被那神用锁链牵引，对所说出的事一无所知。因为这才是\OriginalTerm{占卜者}{soothsayer}的特点：失去自我、被迫、被推、被拖、像疯子一样被拖曳。但先知并非如此，而是以清醒的头脑、沉静的性情，知道自己在说什么，说出一切。因此，甚至在这事发生之前，你就可以借此辨别占卜者和先知。还要注意他如何使自己的言论免于一切猜疑；他呼吁那些亲身体验过这事的人作证。仿佛他说过：\Quote{我并非撒谎，也非像敌人一样轻率地毁谤外邦人的宗教，你们自己为我作证：你们知道，当你们是外邦人时，是如何被拉扯拖拽。}\par

\Quote{但若有人说这些见证人也被怀疑是信徒，那么，来吧，即使从外邦人中，我也要向你们证明这一点。例如，请听柏拉图这么说：(Apol. Soc. c. 7.) \Quote{正如那些传神谕的人和占卜者说出许多美妙之事，却对所说的一无所知。} 再听另一位诗人，也表达同样的意思。因为藉着某些神秘的仪式和巫术，某人将一个\OriginalTerm{魔鬼}{demon}囚禁在人内，那人行占卜，在占卜中跌倒撕裂，无法忍受魔鬼的暴力，快要在那痉挛中死去；他对那些施行神秘技艺的人说，}\par

\Quote{求你们释放我：\newline{}全能的天主不再能\newline{}拘禁凡人的血肉。}\par

又说：\par

\Quote{解下我的花环，用纯净溪流的水滴洗我的脚，\newline{}擦去这些神秘的线条，\newline{}让我离去。}\par

\Quote{因为这些和类似的事（因为可能还会提到很多），向我们指出了如下两个事实：制服魔鬼并使他们成为奴隶的强制，以及那些一度把自己交给他们的人所遭受的暴力，甚至使他们偏离自然的理性。还有那\OriginalTerm{女巫师}{Pythoness}：因为我现在被迫提出并揭露他们另一个可耻的习俗，本应略过不提，因为提及此类事于我们不合宜；但为使你们更清楚地认识他们的羞辱，有必要提及它，好让至少你们由此认识那些利用占卜者的人的疯狂和极度的嘲弄：据说这位女巫师，身为女性，有时会跨坐在阿波罗的三脚架上，这样那恶灵从下面上升，进入她身体的下部，使女人充满疯狂，她披头散发，开始像酒神狂女一样狂舞，口吐白沫，并这样在狂乱中说出她疯狂的话语。我知道你们听到这些事时会羞愧脸红；但他们却以我所描述的耻辱和疯狂为荣。这些和所有类似的事，正是保禄提醒你们时所说的：\Quote{你们知道，当你们还是外邦人时，是怎样被牵引到那些哑吧偶像那里，无论你们如何被牵引。}}\par

并且因为他与那些熟知的人谈论，他没有精确谨慎地陈述一切，不愿使他们厌烦，而只是提醒他们，唤起他们的记忆，他很快就离开了这一点，急忙转向面前的论题。\par

但\Quote{到那些哑吧偶像那里？}是什么意思？这些占卜者曾被牵引拖曳到它们那里。\par

但如果它们本身是哑的，它们怎么能回答别人呢？为什么魔鬼要把他们领到偶像那里？如同战争中俘虏的人，带着锁链，同时使他的欺骗显得可信。因此，为了使人们不认为它只是一块哑石头，他们竭力把人拴在偶像上，以便将自己的名称和头衔刻在上面。但我们的礼仪并非如此。然而他并没有陈述我们的礼仪，我指的是预言。因为这一切他们都很熟悉，预言在他们中间实行，配得上他们的身份，出于理解和完全的自由。因此，你们看，他们有权说话或不说话。因为他们并非被必然性束缚，而是蒙受殊荣。为此原因，约纳逃跑了；\BibleRef{纳 1:3}为此原因，厄则克耳拖延了；\BibleRef{则 3:15}为此原因，耶肋米亚推辞了。\BibleRef{耶 1:6}天主并不强迫他们，而是劝导、告诫、恐吓；并不使他们的心智昏暗；因为导致分心、疯狂和巨大黑暗，是魔鬼的本份；但天主的本份是光照并深思熟虑地教导必要之事。\par

3. 这是占卜者和先知之间的第一个区别；但第二个、不同的区别是他接下来指出的，说：\par

% structure: p0020 source=h2 target=subsection reason=relative-heading-rank; base=h2

\subsection{格林多前书 12:3}

\Quote{因此我告诉你们，凡说话的人若是在天主之神内，不会说耶稣是可诅咒的：} 然后另一个：\Quote{除非在圣神内，没有人能说耶稣是主。}\par

\Quote{当你们看见}，他说，\Quote{有人不呼号祂的名，或咒骂祂，他就是个占卜者。再者，当你们看见另一个人以祂的名说一切话，要明白他是属神的。} \Quote{那么，}你们说，关于\OriginalTerm{慕道者}{Catechumens}，我们应当怎么说？因为如果，除非在圣神内，没有人能说耶稣是主，我们该怎样说那些确实呼号祂的名，却没有祂之神的人呢？但他的话此时并非关于这些人，因为那时还没有慕道者，而是关于信徒和不信者。那么，难道没有魔鬼呼号天主的名吗？那些\OriginalTerm{附魔者}{demoniacs}不是说过：我们认识你是谁，你是天主的圣者？\BibleRef{谷 1:24} 他们不是对保禄说过：这些人是至高天主的仆人？\BibleRef{宗 16:17} 他们说了，但在鞭打之下，被迫之下；从不自愿且不受鞭打。\par

但在此处，宜探究为何魔鬼说出这些话，又为何保禄斥责他。他效法他的老师；因为基督也曾如此斥责：祂不愿意从他们那里得到见证。魔鬼何以也这样做？意图混淆事物的秩序，篡夺宗徒的尊位，并劝服许多人留意他们：如果这事发生，他们便会轻易由此显得可信，从而引入自己的计谋。因此，为免这些事发生，并使欺骗无从开始，祂即使在他们说真话时也封住他们的口，好使人在他们的谎言中全然不听他们，并且完全掩耳不听他们所说的话。\par

4. 因此，既已藉第一个标记和第二个标记显明了占卜者和先知，他接着讨论奇事；并非无理由地转到这个主题，而是为了消除由此产生的纷争，并劝服那些领受较少部分的人不要忧伤，那些领受较多的人也不要自高。因此他也这样开始。\par

% structure: p0025 source=h2 target=subsection reason=relative-heading-rank; base=h2

\subsection{格林多前书 12:4}

\BibleQuote{恩赐虽有区别，却是同一的圣神。}\par

首先，他关注那领受较小恩赐并为此忧伤的人。他说：\Quote{你为什么沮丧？因为你没有像别人那样领受那么多？然而，你要想到这是恩赐，而不是债务，你就能安抚你的痛苦。}为此，他在一开始就这样说：\BibleQuote{但恩赐虽有区别。}他没有说\Quote{标记}，也没有说\Quote{奇事}，而是说\Quote{恩赐}，以恩赐之名劝服他们不仅不要忧伤，还要感恩。他又说：\Quote{此外，你也要想到这一点：即使你在所得的量上不及他人，但因你从同一源头领受，如同那领受更多的人一样，你享有同等的尊荣。因为你当然不能说：圣神将恩赐给了他，而是一位天使给了你；因为圣神赐给了你，也赐给了他。}因此他加上：\BibleQuote{但同一的圣神。}所以，即使恩赐有区别，但施予者并无区别。因为你和他是从同一泉源汲取。\par

% structure: p0028 source=h2 target=subsection reason=relative-heading-rank; base=h2

\subsection{格林多前书 12:5}

\BibleQuote{职分虽有区别，却是同一的主。}\par

如此，他充实了安慰，又提到了圣子和圣父。并且，他用另一个名称称呼这些恩赐，以此增加安慰。为此他也这样说：\BibleQuote{职分虽有区别，却是同一的主。}因为那听到\Quote{恩赐}并领受了较小份的人，或许会忧伤；但当我们说到\Quote{职分}时，情况就不同了。因为这个词暗示了辛劳和汗水。他说：\Quote{那么，如果祂吩咐另一个人多劳，而让你省力，你为什么忧伤呢？}\par

% structure: p0031 source=h2 target=subsection reason=relative-heading-rank; base=h2

\subsection{格林多前书 12:6}

\BibleQuote{功效虽有区别，却是同一的天主，在一切人身上行一切事。}\par

% structure: p0033 source=h2 target=subsection reason=relative-heading-rank; base=h2

\subsection{格林多前书 12:7}

\BibleQuote{圣神显示在每人身上虽不同，但全是为人的好处。}\par

有人说：\Quote{什么是功效？什么是恩赐？什么是职分？}它们只是名称上的区别，因为事情是相同的。因为所谓的\Quote{恩赐}，就是\Quote{职分}，他也称为\Quote{运作}。如此，你要\BibleQuote{履行你的职分}\BibleRef{弟后 4:5}，以及\BibleQuote{我高举我的职分}\BibleRef{罗 11:13}，并给弟茂德写信说：\BibleQuote{我提醒你把天主藉我的覆手所赋予你的恩赐，再炽燃起来}\BibleRef{弟后 1:6}。又写信给迦拉达人，他说：\BibleQuote{因为那在伯多禄身上运行，使他成为宗徒的，也在我身上运行，使我成为外邦人的宗徒}\BibleRef{迦 2:8}。你看见吗？他暗示圣父、圣子和圣神的恩赐并无区别。不是混淆位格，绝对不可！而是表明本体的同等尊荣。因为圣神所赐的，他也说是天主所做的；这也正是圣子所规定和赏赐的。然而，如果一位比另一位低下，或另一位比它低下，他就不会如此记述，也不会用这种方式安慰那烦恼的人。\par

5. 此后，他又以另一种方式安慰他：考虑到所赐予的尺度是为他有益的，即使不那么宏大。因为说了这是\Quote{同一的圣神}、\Quote{同一的主}、\Quote{同一的天主}，并以此使他恢复平静，他再引入另一种安慰，这样说：\BibleQuote{但圣神显示在每人身上虽不同，全是为人的好处。}以免有人说：\Quote{如果主是同一、圣神是同一、天主是同一，但我领受了较少的份呢？}他说，这样是有益的。\par

他称奇迹为\Quote{圣神的显示}，显然是有道理的。因为对我这信徒来说，那拥有圣神的人因领受圣洗而显明；但对不信者而言，除了奇迹，别无显明之法：所以从这里也再次得到不小的安慰。因为即使恩赐有区别，但证据是同一的：因为无论你领受多或少，你都同样地显明。所以，如果你想要表明这一点，即你有圣神，你就有足够的证明。\par

因此，既然施予者是同一的，所赐予的纯粹是恩惠，且由此有显示，而这对你更为有益；你就不要忧伤，好像被轻视了。因为天主这样做，并非为了羞辱你，也不是要声明你不如别人，而是要饶恕你并着眼于你的益处。领受超过自己能力所及的，反而无益，有害，并是沮丧的适当原因。\par

% structure: p0039 source=h2 target=subsection reason=relative-heading-rank; base=h2

\subsection{格林多前书 12:8}

\BibleQuote{这人从圣神蒙受了智慧的言语，另一人却由同一圣神蒙受了知识的言语；}\par

% structure: p0041 source=h2 target=subsection reason=relative-heading-rank; base=h2

\subsection{格林多前书 12:9}

\BibleQuote{至于另一个人，在同一圣神内，有信德；另一个人，在同一圣神内，有治病的恩惠。}\par

你看见他如何在每一处都加上这个说法，说：\Quote{借着同一圣神，并按照同一圣神？}因为他知道由此而来的安慰是大的。\par

% structure: p0044 source=h2 target=subsection reason=relative-heading-rank; base=h2

\subsection{\BibleRef{格前 12:10}}

\BibleQuote{另一个人能行奇迹；另一个人能说预言；另一个人能辨别神恩；另一个人能说各种语言；另一个人能解释语言。}\par

既然他们以此自夸，因此他把它放在最后，并加上：\par

% structure: p0047 source=h2 target=subsection reason=relative-heading-rank; base=h2

\subsection{\BibleRef{格前 12:11}}

\BibleQuote{但这一切都是同一圣神，他所行的。}\par

他的安慰中所含的普遍良药就是：他们都从同一根源、同一宝藏、同一川流中领受。所以，他一再强调这个说法，以消除表面上的不平等，并安慰他们。确实，他之前指明了圣神、圣子和圣父都分施恩惠，但在这里他满足于只提圣神，好让你即使在此也能理解他们的尊荣是相同的。\par

但什么是\Quote{智慧的言语？}那是\Person{保禄}所有的，是\Person{若望}——雷霆之子所有的。\par

什么是\Quote{知识的言语？}那是大多数信徒所有的，他们确实拥有知识，但并不因此能够教导别人，也不容易把自己所知的传达给别人。\par

\Quote{另一个人有信德：}这里不是指教义的信德，而是指奇迹的信德；关于这信德，基督说：\BibleQuote{如果你们有信德像芥菜子那么大，你们对这座山说：挪开，它就会挪开。}\BibleRef{玛 17:20}宗徒们也为此祈求祂说：\BibleQuote{请增加我们的信德！}\BibleRef{路 17:5}因为这是奇迹之母。但拥有行奇迹的能力和治病的恩惠并不是同一回事：因为拥有治病恩惠的人只施行医治；但拥有行奇迹能力的人也施行惩罚。因为奇迹不仅是医治，也包括惩罚：就如\Person{保禄}使人瞎眼，\Person{伯多禄}使人致死一样。\par

\Quote{另一个人说预言；另一个人辨别神恩。}什么是\Quote{辨别神恩？}即知道谁是属神的，谁不是；谁是先知，谁是欺骗者；正如他对得撒洛尼人说的：\BibleQuote{不要轻视先知之恩，}\BibleRef{得前 5:20}但要辨别一切，持守那美好的。因为那时假先知蜂拥而至，魔鬼暗中以虚假取代真理。\Quote{另一个人能说各种语言；另一个人能解释语言。}因为一个人自己知道自己在说什么，却不能为别人解释；而另一个人却两者兼得，或只得其中之一。这些话的恩赐似乎很大，因为宗徒们首先领受了它，格林多人中的大多数也得到了它。但教导的言语却不是这样。因此他把教导放在前面，把说语言放在最后：因为说语言是为教导的缘故，其余的一切也确实如此：包括预言、行奇迹、说各种语言和解释语言。因为没有一样能与教导相等。因此他也说：\BibleQuote{那些善于管理的长老，尤其是那些在言语和教导上劳苦的，应被视为配得双倍的尊敬。}\BibleRef{弟前 5:17}他又写信给\Person{弟茂德}说：\BibleQuote{你要专注于宣读、劝勉、教导；不要忽略你内在的神恩。}\BibleRef{弟前 4:13}你看，他如何也称它为神恩？\par

6. 接着，他之前所说的安慰——\Quote{同一圣神}——他也放在我们面前，说：\BibleQuote{但这一切都是同一圣神所行的，他随自己的意愿分给各人。}他不仅给予安慰，也堵住了反驳者的口，说：\BibleQuote{他随自己的意愿分给各人。}因为有必要既医治又包扎，正如他在给罗马人的书信中所做的，说：\BibleQuote{你这个人哪，你是谁，竟敢向天主抗议？}\BibleRef{罗 9:20}同样在这里，\BibleQuote{他随自己的意愿分给各人。}\par

而他所指为属于父的，也表明属于圣神。因为论及父，他说：\BibleQuote{但那同在一切内行一切的天主。}同样论及圣神：\BibleQuote{但这一切都是同一圣神所行的。}不过，有人会说：\Quote{祂是在天主的推动下行这事。}不，他没有这样说，而是你捏造的。因为当他说\BibleQuote{那在一切内行一切的}时，这是论人而言；你几乎不能说在那些人中间也包括圣神，即使你疯狂痴迷至极。因为他既然说了\Quote{藉着圣神}，免得你以为\Quote{藉着}一词表示低劣或受推动，他便补充说：圣神\Quote{行事}，而非\Quote{被行事}，并且\BibleQuote{随意}行事，而非受命行事。因为正如论及父，子说\BibleQuote{祂使死者复活，并赋予生命}；同样论及自己，祂说\BibleQuote{祂随自己的意愿赋予生命}。\BibleRef{若 5:21}同样，论及圣神，在另一处也说，祂以权能行一切，没有能阻止祂的（因为\BibleQuote{风随意吹} \BibleRef{若 3:8}，虽说的是风，但有助于确立这一点）；但在此处，祂\BibleQuote{随自己的意愿行一切。}并从另一处得知，祂不属于被推动之物，而属于推动之物。因为他说：\BibleQuote{除了人内里的神，谁能知道人的事？同样，除了天主的圣神，也没有人知道天主的事。}\BibleRef{格前 2:11}那么，\Quote{人的神}即灵魂，无需被推动就能知道自己的事，我想对每个人都是显而易见的。同样，圣神也无需被推动就能\Quote{知道天主的事}。因为祂的意思是这样：\Quote{天主的奥秘}对圣神是知道的，如同人的灵魂知道自己内心的事。但如果灵魂为此目的无需被推动，那么那知道天主的深奥、无需为那知识而被推动的一位，在赐予宗徒恩赐时更无需任何推动之力。\par

但除此以外，我先前说过的那点，我现在也要再次提及。那是什么？就是：如果圣神是低等的，来自另一实体，那么祂的安慰以及我们听到\Quote{同一圣神}的话就没有益处。因为从君王那里领受的人，我承认，发现是祂亲自赐予，这可能是非常慰藉的事；但如果是来自奴隶，那么当人以此责备他时，他反而更加恼怒。所以由此也显而易见，圣神并不来自奴隶的实体，而是来自君王的实体。\par

7. 所以，正如他在说\BibleQuote{服务职分虽有区别，却是同一的主；功效虽有区别，却是同一的天主}时安慰了他们，同样，当他上面说\BibleQuote{神恩虽有区别，却是同一的圣神}，以及此后又说\BibleQuote{这一切都是同一的圣神所行，随自己的心意分别赐给各人}时，也是如此。\par

\Quote{我恳求你们，不要困惑，}他说，\Quote{也不要忧伤说：\InnerQuote{为什么我领受了这个而没有领受那个？}也不要求问圣神的理由。因为如果你知道祂出于眷顾而施赐，就要想到从同一眷顾祂也赐下了尺度，并满足于你所领受的，欢喜快乐；对于你没有领受的，不要埋怨；反倒要承认天主的恩惠，因为你没有领受超过你能力的事。}\par

5. 如果连属神的事也不该过分好奇，何况属世的事？而应安静，不细究为何一人富有，另一人贫穷。因为首先，并非每个富人都从天主成为富人，许多人是出于不义、掠夺和贪婪。因为那禁止人富有的，如何能赐予祂所禁止领受的？\par

但为了远超所需地堵住那些在这些事上反驳我们之人的口，让我们把论述往上推至财富由天主赐予的时代，并回答我。为何亚巴郎富有，而雅各伯甚至缺乏面包？二人岂不都是义人？祂岂不也论及三人说：\BibleQuote{我是亚巴郎的天主，依撒格的天主，雅各伯的天主？}\BibleRef{出 3:6}那么为何一人富有，另一人却做了雇工？或者，为何厄撒乌富有，他是不义之人，谋害兄弟，而雅各伯却长期为奴？为何依撒格一生安逸，而雅各伯却劳苦艰辛？为此他也说：\BibleQuote{我岁月短少，且不美好。}\BibleRef{创 47:9}\par

为何达味，既是先知又是君王，自己一生也劳苦？而他的儿子撒罗满，四十年来比任何人都安稳，享受深切的和平、荣耀、尊贵，遍历各种福乐？又有什么理由，在众先知中，一人受苦更多，另一人更少？因为对每个人这样更相宜。因此我们各人必须说：\BibleQuote{你的判断是多么深奥！}\BibleRef{咏 35:6}因为如果那些伟大奇妙的人并非都同样被天主锻炼，有的藉贫穷，有的藉富有；有的藉安逸，有的藉患难，那么我们现在更当牢记这些事。\par

8. 此外，还应当考虑，许多事情的发生并非按照天主的旨意，而是出于我们的邪恶。因此，不要说：\Quote{为什么一个恶人富有，而一个义人贫穷？} 因为首先，对于这些事情也可以解释：义人从贫穷中不会受到任何损害，反而更增添光荣；恶人在财富中只是为将来储存惩罚，除非他改变；甚至在惩罚之前，财富往往成为他许多邪恶的原因，引他陷入无数陷阱。但天主允许这事，同时表明自由意志的选择，也教导所有人不要为金钱疯狂贪恋。\par

\Quote{那么，当一个恶人富有，且未遭受任何可怕之事时，又怎样呢？} 你说。\Quote{既然如果善人拥有财富，他是正义的；但如果恶人，我们该说什么？} 那就是他更值得怜悯。因为财富加于邪恶，会加重祸害。然而，如果他是善人而贫穷？他并未受损害。那么，他是恶人而贫穷？这是完全公正和应得的，甚至对他有益。\Quote{但某人，} 你说，\Quote{从祖先那里继承了财富，却挥霍在娼妓和寄生虫上，并未遭受任何恶事。} 你说什么？他行淫，而你竟说\Quote{他未遭受恶事}？他醉酒，你以为他在享乐？他挥霍无度，你判断他值得羡慕？有什么比这毁灭灵魂的财富更糟呢？但你，如果身体扭曲残废，你会认为那是值得大哭的；而你看到他整个灵魂被肢解，却还认为他幸福？\Quote{但他没有意识到，} 你说。那么，正因如此，他更值得怜悯，如同所有狂乱的人。因为知道自己生病的人当然会寻求医生并接受治疗；但不知道的人根本不可能得救。告诉我，你称这样的人为幸福吗？\par

但这不足为奇：因为大多数人不认识真正的爱智。因此我们遭受最严厉的惩罚，受责打却仍不远离惩罚。因此有愤怒、沮丧和无休止的骚动；因为当天主向我们展示了没有忧愁的生活，即美德的生活，我们却离开它，标出另一条路，即财富和金钱之路，充满无数邪恶。我们如同这样的人：不懂得辨别身体之美，只归因于衣服和装饰，看见一位拥有天然美貌的女子，便匆匆走过，但看见一位丑陋、畸形、扭曲却穿着华丽衣服的女子，却娶她为妻。如今民众对美德和邪恶的态度也大致如此：他们接纳本性丑陋的，因她外在的装饰；却转离那美丽可爱、因朴素之美而更值得选择的。\par

9. 因此我感到羞愧，因为在愚昧的外邦人中，有人实践这种智慧——即使不在行动上，至少也在判断上；他们知道现世之物的易逝性。而在我们当中，有些人甚至不理解这些，连判断都被败坏；虽然圣经不断在我们耳边回响，说：\Quote{凡作恶的，他眼中看为可憎的；他尊重敬畏主的人。\BibleRef{咏 14:4} 敬畏主超越一切；敬畏天主，遵守祂的诫命，这是人的全部。\BibleRef{训 12:13} 不要嫉妒恶人。\BibleRef{咏 48:16} 凡有血肉的如同草芥，人的一切荣耀如同草上的花。\BibleRef{依 40:7}} 尽管我们每天都听到这些以及类似的话，我们却仍被钉在地上。如同无知的孩子，连续学习字母，但当打乱顺序考问他们时，便会张冠李戴，引人发笑。同样，你们在这里听我们按顺序陈述时，似乎能跟上；但当我们在户外无顺序地问你们，事物中什么应放在首位、什么其次、什么接着什么时，你们不知如何回答，便成为可笑的人。告诉我，这岂非极大的笑柄吗：他们期盼不朽和\BibleQuote{眼所未见，耳所未闻，人心所未曾想到的}美善，却为这世上短暂的东西奋斗，并将之视为可羡慕的？因为如果你还需要学习这些事——财富并非大事，现世之物是影子和梦，如烟消散飞去——那么暂且站在圣所之外，留在门廊里，因为你还不配进入天上宫廷的大门。如果你不能分辨那不稳定、不断消逝之物的本质，何时才能藐视它们呢？\par

但如果你说你知道，就不要再好奇探询并忙碌于推测某人为何富有、某人为何贫穷；因为你发问时，如同四处打听为什么一人俊美而另一人黝黑，或一人鹰鼻而另一人扁平。因为这些事无论是这样还是那样，对我们都毫无区别；贫穷和富有更是如此。一切都取决于我们如何使用它们。无论你是贫穷，可以克己自足地生活；或是富有，如果你逃避美德，你就是所有人中最悲惨的。因为真正与我们相关的是美德之事。如果这些事不被拥有，其余都无用。因此才有那些无休止的问题，因为大多数人认为无关紧要的事对他们重要，而对重要之事却不加理会；因为对我们重要的是德行和智慧。\par

因为那时你们站在我不知道的地方，离她甚远，所以才思想混乱，波涛汹涌，风暴骤起。因为人一旦从天上的荣耀和对天主的爱慕中坠落，就贪求现世的荣耀，成为奴隶和俘虏。你们说：\Quote{我们怎么会贪求这个呢？}这是因为没有强烈地渴慕那个。而这件事本身又是怎么发生的呢？由于疏忽。疏忽又从哪里来？来自轻视。轻视又从哪里来？来自愚昧，来自依恋现世的事物，不愿仔细探究事物的本性。而这后者又从哪里来？来自既不留心阅读圣经，也不与圣善之人交谈，反而追随恶人的集会。\par

因此，为了不让这种事情永远持续下去，也为了避免一波又一波的海浪将我们卷入苦难的深渊，完全淹没并毁灭我们；趁着还有时间，让我们振作起来，站在磐石上——我指的是神圣的教义和话语——俯瞰今生的狂澜。因为这样，我们既能自己逃离同样的命运，也能将其他遭遇海难的人拉上来，从而获得未来的恩赐，藉着恩宠和慈悲等等。\par

\SourceRecord{Translated by Talbot W. Chambers.；Nicene and Post-Nicene Fathers, First Series；Vol. 12.；Edited by Philip Schaff.；Buffalo, NY: Christian Literature Publishing Co.,；1889.；<http://www.newadvent.org/fathers/220129.htm>.}

% structure: p0001 source=h1 target=section reason=page-title

\section{\WorkTitle{论格林多前书第三十篇讲道}}

% structure: p0002 source=h2 target=subsection reason=relative-heading-rank; base=h2

\subsection{\BibleRef{格前 12:12}}

\BibleQuote{就如身体只是一个，却有许多肢体；身体所有的肢体虽多，仍是一个身体；基督也是这样。}\par

在他们因所赐之物是白白恩宠，且都从\Quote{同一圣神}领受，并是为\Quote{有益处}而赐予，甚至藉较小恩赐也得以彰显而受到安抚后，又因当顺从圣神权威的义务而止住他们的口（他说：\Quote{这一切都是这同一圣神所运作，随祂心愿分给每人，}因此不应过于好奇），现在他从另一个常见的比喻着手，同样安慰他们，并照他惯常的做法，诉诸于自然本身。\par

因为当他论及男女发式时，在说完其他一切后，又从中取材来规劝他们，说：\Quote{连本性不也教导你们：男人若留长发，便是他的羞辱；但女人若留长发，便是她的荣耀吗？}\BibleRef{格前 11:14}当论及祭偶像之物，禁止接触时，他也从外邦人的事例中引证，提及奥林匹克运动会，说：\Quote{在赛跑中，众人固然都跑，但只有一个得奖赏。}\BibleRef{格前 9:24}并用牧人、士兵和农夫的事例来证实这些观点。因此，他在这里也引用一个常见的比喻，着力并竭力证明，没有谁真的处于更差的境地：这能显示出令人惊奇的事，并有助于安慰较弱的人，我是指身体的比喻。因为没有什么比确信自己并未因分得较少而有所缺，更能安慰那些心神卑微、恩赐较小的人，并说服他们不必忧伤。因此，保禄为阐明这一点，这样表达：\Quote{就如身体只是一个，却有许多肢体。}\par

你看到他的精确考究吗？他指出了同一事物既是一又是多。因此他又更用力地强调说：\Quote{并且这身体的一切肢体，虽多，仍是一个身体。}他没有说：\Quote{虽多，仍属于一个身体，}而是说：\Quote{这个身体本身就是多。}而这些众多的肢体就是这个一。因此，如果一是多，多又是一，那么差别在哪里？优越在那里？劣势在那里？因为他说：一切都是一个。而且不只是一个，若严格就其本性，即作为身体而言，它们都被发现是一个；但若按其个别本质而言，则出现了差别，而这差别在一切肢体中都一样。因为没有哪个肢体能独自构成身体，每个肢体都同样欠缺构成身体的条件，而需要聚合；只有当多成为一时，才有一个身体。因此，暗示这一点时，他说：\Quote{并且这身体的一切肢体，虽多，仍是一个身体。}他没有说\Quote{优越者与低劣者，}而是说\Quote{虽多，}这是共通的。\par

那么，它们怎么可能是一呢？当你撇开肢体的差别，只考虑身体时。因为眼睛这个，脚也同样，就其作为肢体并构成身体而言。在这方面没有差别。你不能说一个肢体自成身体，而另一个则不然。因为在这点上它们都是平等的，正因为它们同属一个身体。\par

但说了这些，并以此普遍判断清楚地表明后，他加上：\Quote{基督也是这样。}而当他应该说\Quote{教会也是这样}时（这本是自然的结论），他却没有这样说，反而以基督的名号代之，将论述提升到高处，并越来越触动听众的敬畏之心。但他的意思是：\Quote{基督的身体，即教会，也是这样。}因为正如身体和头是一个人，他说教会和基督是一体。因此，他以基督代替教会，将这个名号给予祂的身体。他说：\Quote{正如我们的身体虽由许多部分组成，却是一体；同样，在教会中，我们都是一体。因为虽由许多肢体组成，但这些许多肢体构成一个身体。}\par

2. 你看，他以此常见比喻安抚并振作那个自感轻看的人后，又离开这日常经验的话题，转向另一属神的话题，带来更大的安慰，并表明同等的尊荣。那么，这是什么？\par

% structure: p0010 source=h2 target=subsection reason=relative-heading-rank; base=h2

\subsection{\BibleRef{格前 12:13}}

\Quote{因为我们众人，不论是犹太人，或是希腊人，或是为奴的，或是自主的，都在同一圣神内受洗，成了一个身体。}\par

他的意思是：使我们成为一个身体、并重生我们的，是同一圣神；因为并不是一人受洗于一个圣神，另一人受洗于另一个。不仅那施洗我们的是同一个，而且那使我们受洗归于的，即我们受洗的目的，也是同一个。因为我们受洗，并不是为了形成许多各自的身体，而是为了我们众人彼此保持一个身体的完美本性：即我们众人都是一个身体，我们都受洗归于同一个身体。\par

如此，那形成它的是一位，那形成它的也是一个。他没有说\Quote{为使我们都能成为同一身体}，而是说\Quote{为使我们都能成为一个身体}。因为他总是力求使用更具表现力的词句。他巧妙地说\Quote{我们众人}，也把自己包括在内。他说：\Quote{就连我这位宗徒，在这方面也不比你多什么。因为你同我一样是身体，我同你一样，我们都共有同一个头，经历了同样的产痛。因此，我们也是同一个身体。}他又说：\Quote{我何必提犹太人呢？因为连那些远离我们的外邦人，祂也纳入了同一个身体的整体。}所以他说了\Quote{我们众人}之后，并未止步，又加上\Quote{无论是犹太人或希腊人，无论是为奴的或自主的}。既然从前相距如此遥远，我们都得以联合而成为一体，那么在成为一体之后，我们更无权悲伤沮丧。事实上，差别已不复存在。因为如果祂将同样的恩赐赐给了希腊人和犹太人、为奴的和自主的，那么在这之后，祂怎会因赐予了更完美的合一而使他们分裂呢？\par

\BibleQuote{都饮于同一圣神。}\par

% structure: p0015 source=h2 target=subsection reason=relative-heading-rank; base=h2

\subsection{格林多前书 12:14}

\BibleQuote{原来身体不是只有一个肢体，而是有许多。}\par

即：我们领受了同样的入门礼，共享同一餐桌。他为何不说\Quote{我们被同一身体滋养，饮同一血}？因为他说了\Quote{圣神}，就同时表明了二者，既包括肉身，也包括血。因为我们借着这两者都\Quote{饮于圣神}。\par

但在我看来，他现在指的是圣洗之后、圣事之前在我们内发生的圣神的眷顾。他说\Quote{我们都饮了}，因为这种隐喻的说法极其适合他提出的主题：就好像他论及植物和园子时，说所有树木都被同一泉源浇灌，或同一水灌溉；同样，他说：\Quote{我们都饮了同一圣神，享用了同一恩宠。}\par

既然一位圣神形成我们，并将我们众人聚集在一个身体内——因为\BibleQuote{我们都因一个圣神受洗，成为一个身体}的意义就是如此——，又赐给我们同一餐桌，并给予我们同样的灌溉（因为\BibleQuote{我们都饮于同一圣神}的意义就是如此），将如此分离的人联合起来；而许多事物当它们成为一时，就成为身体：那么，请问，你们为什么总是因着它们的不同而摇摆不定？但如果你们说：\Quote{因为有许多不同的肢体}，要知道这正是身体令人称奇和独有的优越之处，即众多不同的事物合而为一。但如果它们不是许多，那么它们成为一个身体就不会如此奇妙和不可思议；甚至根本就不是一个身体。\par

3. 不过他最后才提出这一点；现在他先论及肢体本身，这样说：\par

% structure: p0021 source=h2 target=subsection reason=relative-heading-rank; base=h2

\subsection{格林多前书 12:15}

\BibleQuote{假使脚说：\InnerQuote{我因为不是手，便不属于身体}；它因此就不属于身体了吗？}\par

% structure: p0023 source=h2 target=subsection reason=relative-heading-rank; base=h2

\subsection{格林多前书 12:16}

\BibleQuote{假使耳说：\InnerQuote{我因为不是眼，便不属于身体}；它因此就不属于身体了吗？}\par

因为如果一个因低贱、另一个因高贵，就不允许它们属于身体，那么整个身体就被废除了。所以不要说：\Quote{因为我卑贱，我就不是身体。}因为脚虽然处于较低的位置，但它仍属于身体：因为是否属于身体，并非由于一个处在这里、另一个在那里（这是造成位置差别的因素），而在于是否结合或分离。因为成为身体或不成为身体，源于是否成为一体。但请你注意他审慎的方式：他如何将他们的言辞应用到我们的肢体上。正如他以上所说的：\BibleQuote{我把我自己和阿颇罗当作一个比喻应用在你们身上}\BibleRef{格前 4:6}，同样，为了使他的论证不带偏见且易于接受，他引入了肢体说话：这样当他们听见本性回答他们时，被经验本身和普遍的声音所驳斥，他们就再没有什么可反对的了。他说：\Quote{因为即使你们愿意，尽管你们抱怨，你们也不能脱离身体。因为就如自然律，恩宠的能力更护守一切，使它们完整。}看看他是如何遵守没有多余之物的规则：他没有在所有肢体上展开论证，而只用了两个而且是最极端的肢体：既指出了最尊贵的眼，也指出了最卑微的脚。他没有让脚与眼对话，而是与稍高于它的手对话；让耳与眼对话。因为我们嫉妒的不是那些远高于我们的人，而是那些略高于我们的人，因此他这样进行对比。\par

% structure: p0026 source=h2 target=subsection reason=relative-heading-rank; base=h2

\subsection{格林多前书 12:17}

\BibleQuote{若全身是眼，听觉在哪里？若全身是听觉，嗅觉在哪里？}\par

这样，因为他落到了肢体的差异上，提到了脚、手、眼和耳，他将他们引向对自己低贱与高贵的思考：请看他又如何安慰他们，暗示这样是合宜的；并且他们众多和不同，正是这一点使他们成为身体。但如果他们全都是同一个，就不成其为身体了。因此他说：\Quote{若全是同一个肢体，身体在哪里？}不过这一点他直到后面才提及；但在此他还指出了更多：除了任何一个人不能成为身体之外，这甚至取消了其余肢体的存在。\par

\BibleQuote{若全身是听觉，嗅觉在哪里？}\par

4. 然而，因为他们仍然感到不安，他之前所做的，现在也照样做。因为正如在那里他首先提出益处来安慰他们，然后严厉地堵住他们的口，说：\BibleQuote{但这一切都是同一圣神所运行，按祂的意愿个别分给各人。}同样，在这里他陈述了理由，表明所有人如此是有益的，然后又将整个事归于天主的旨意，说：\par

% structure: p0031 source=h2 target=subsection reason=relative-heading-rank; base=h2

\subsection{\BibleRef{格前 12:18}}

\BibleQuote{然而天主将各肢体一一安排在身体上，都是按祂的旨意。}\par

正如他论圣神时说\BibleQuote{按祂的意愿}，同样在这里也说\BibleQuote{按祂的旨意}。如今你不要再追问原因：为什么是这样而不是那样。因为即使我们能给出千万个理由，也不如说，正如至高的工匠所愿，事情就这样成就。因为祂的意愿正是按照益处。如果在我们的身体上，我们并不好奇地查问各个肢体，那么在教会中就更当如此。请注意他的细心：他没有提出由本性产生的差别，也没有提出由运作产生的差别，而是提出由位置产生的差别。因为他说：\BibleQuote{如今，天主将各个肢体一一安置在身体上，都是按祂的旨意。}他恰当地说\Quote{各个}，指出功用遍及所有。因为你不能说：\Quote{这个祂亲自安置了，那个却没有；而是每个都按祂的旨意处于各自的位置。}因此，对于脚来说，处于这样的位置也是有益的，不仅对于头如此；如果它颠倒秩序，离开自己的位置而到别处，虽然似乎改善了条件，但会导致整体的崩溃和毁灭。因为它既失去了自己的位置，也达不到另一位置。\par

% structure: p0034 source=h2 target=subsection reason=relative-heading-rank; base=h2

\subsection{\BibleRef{格前 12:19-20}}

\BibleQuote{假如所有肢体都是同一个，那么身体在哪里？}\BibleQuote{然而现在肢体虽然多，身体却是一个。}\par

5. 因此，他们通过天主的安排被充分制止后，\Person{保禄}再次陈述理由。他并非总是这样做，也不是总那样做，而是交替变化他的讲论。因为单方面斥责，会使听者困惑；相反，若使听者习惯于凡事都要求理由，则会在信仰上伤害他们。为此，保禄不断练习两者，使他们既能相信，又不至困惑；在制止他们之后，又同样给出理由。还要留意他在争战中的热忱和得胜的完全。因为他们自以为因巨大的差异而荣誉不均，甚至从这些事中，保禄表明正因为此，他们才荣誉均等。如何？我将告诉你们。\par

\BibleQuote{假如所有都是同一个肢体，}他说，\BibleQuote{那么身体在哪里？}\par

他意思是说，如果你们中间没有巨大的差异，你们就不能成为身体；既然不是身体，就不能成为一；既然不是一，就不能荣誉均等。由此又推出，如果你们都是荣誉均等，你们就不是身体；既然不是身体，就不是一；既然不是一，怎能荣誉均等？然而实际上，因为你们并非都赋有同一种恩赐，所以你们成为身体；既然是身体，你们全都是一个，并在\Quote{是身体}这一点上彼此无别。因此，正是这种差异主要造成你们荣誉的均等。所以他接着说：\BibleQuote{然而现在肢体虽多，身体却是一个。}\par

6. 因此，让我们也思考这些，摒弃一切嫉妒，既不嫉妒那有更大恩赐的人，也不轻视那有较小恩赐的人。因为天主这样愿意了；我们不要反对。但你若仍感不安，就想想你的工作往往是你的兄弟所无法完成的。因此，即使你较为逊色，但在这方面你却占有优势；而他虽然更大，却在这方面较差；如此便达到了均等。因为在身体中，即使是小的肢体似乎也贡献不小，而大的肢体有时也会因它们受损——我是说因它们的缺失而受损。例如，身体中有什么比头发更微不足道呢？但如果你从眉毛和眼睑上将它除去，尽管它微不足道，你却毁坏了面部的全部优雅，眼睛也不再显得同样美丽。然而损失的只是一件小事；但尽管如此，一切美感都被破坏了。不仅美感，而且眼睛的许多功用也丧失了。原因是我们的每个肢体都有其独特的功用和共同的功用；同样，我们有独特的美和共同的美。这些美看似分开，却完全结合在一起，当一个被破坏时，另一个也随之消亡。让我解释：假如有明亮的眼睛、含笑的脸颊、红润的嘴唇、挺直的鼻子、开阔的额头；然而，如果你损害其中最小的一点，你就损害了它们共同的美；一切都显得沮丧；先前如此美丽的一切，看起来都会丑陋；同样，如果你只压碎鼻尖，就给整个面容带来极大的畸形；而这只是对一个肢体的损伤。同样在手部，如果你从一根手指上取走指甲，你会看到同样的结果。如果你在功能上也看到同样的情形，取走一根手指，你会看到其他手指变得不太灵活，不再同样发挥它们的作用。\par

既然一个肢体的损失是共同的畸形，它的安全则是所有人的美丽，让我们不要高傲，也不蹂躏我们的邻人。因为通过那个小肢体，即使是大的肢体也是美丽和悦目的，眼睛也由微不足道的眼睑所装饰。因此，与兄弟争战的人就是与自己争战；因为伤害不仅临到那人，自己也将遭受不小的损失。\par

7. 为使这事不发生，让我们关心我们的邻人如同关心我们自己，也让我们将这身体的图像转移到教会，并关心所有人如同关心我们自己的肢体。因为在教会中有许多不同的肢体：有些更尊贵，有些更欠缺。例如，有贞女歌咏团，有寡妇团体，有在神圣婚姻中发光的兄弟会；简言之，美德有许多等级。同样，在施舍方面也是如此。因为有些人舍弃了自己的一切财物；有些人只关心足够，不求超过必需品；有些人施舍自己多余之物：然而，所有这些人都彼此装饰；如果那更大的轻视那较小的，他就在最大的程度上损害了自己。因此，假设一个贞女轻慢一个已婚妇女，她就切断了自己报酬不少的一部分；那舍弃一切的人若责备没有这样做的人，他就舍弃了自己劳动的许多果实。我为何只提贞女、寡妇和没有财产的人呢？有什么比那些乞讨的人更卑下呢？然而，这些人也在教会中履行着最重要的职务，紧贴着圣所的门，提供其最大的装饰之一；没有他们，教会的圆满就无法完成。这件事，宗徒们似乎也注意到，从起初就制定了法律，如同关于其他一切事物一样，也关于寡妇：他们对此事如此重视，甚至设立了七位执事管理她们。因为当我数算教会成员时，主教、长老、执事、贞女、守贞者都列入我的计数，寡妇也同样。是的，她们所填补的职务并非不重要。因为你愿意时才到这里来；但这些人日夜唱圣咏并专心服务：不仅是为了施舍才这样做；因为如果那是她们的目的，她们可以在市场行走并在小巷中乞讨：但她们也有不小程度的虔敬。至少，请看她们处于何等贫穷的熔炉中；然而你永远不会从她们口中听到一句亵渎或不耐烦的话，不像许多富人之妻那样。有些人时常饿着肚子躺下休息，另一些人经常被寒冷冻僵；然而，她们在感恩和赞颂中度过时光。即使你给一文钱，她们也感谢并愿赠予者获得万福；若你什么也不给，她们也不抱怨，反而仍祝福，并以能享受每日食粮为幸福。\par

有人回答说：\Quote{是的}，\Quote{因为无论愿意与否，她们必须忍受。}为何，告诉我？你为何说出这样苦涩的话？难道没有可耻的技艺给年老的男女带来利益吗？如果她们选择抛弃对正直生活的所有关心，她们难道不能通过这些手段大量养活自己吗？你难道没有看到那个年纪的许多人，通过做淫媒和皮条客以及其他此类服务，不仅活着，而且生活奢侈？这些妇女却不是这样；她们宁愿饿死，也不愿玷污自己的生命并出卖自己的救恩；她们整天坐着，为你预备一份救恩的良药。\par

因为没有一位医师伸出手来使用手术刀，能如此有效地从我们伤口中割除腐败，如同一个穷人伸出右手接受施舍，以除去伤口留下的疤痕一样。真正奇妙的是，他们在施行这种卓越的外科手术时没有痛苦和忧愁：我们这些管理百姓并给你们如此多有益劝告的人，并不比那坐在教堂门前的人通过沉默和面容更真实地讲论。因为我们也每天在你们耳边讲这些话，说：\Quote{人啊，不要高傲；人性是易衰退且容易跌倒的；我们的青春匆匆变为老年，美丽变为丑陋，力量变为软弱，尊荣变为轻视，健康变为疾病，光荣变为卑贱，财富变为贫穷；我们的事如同湍急的水流，永不静止，而是急速向下倾斜。}\par

同样的劝告，它们也通过外表和经验本身给予，而且经验是更明显的劝告。例如，现在坐在外面的那些人中，有多少曾在青春风华时做过大事？这些面目可憎的人中，有多少曾在体力和容貌上超过许多人？不，不要不相信或嘲笑。因为生命确实充满了无数这样的例子。如果许多卑贱低微的人时常成为国王，那么从伟大光荣变为卑贱低微的人有何可奇呢？因为前者更为非凡；而后者却时常发生。所以，不应不相信他们中任何人曾在技艺、武器和财富富足上兴盛过，反而要以极大的怜悯同情他们，并为自己恐惧，免得我们有一天也遭受同样的事。因为我们也是人，也容易经历这种迅速的变化。\par

8. 但也许有些轻率之人，即那些惯于嘲弄之人，会反对我们所说的，并完全嘲笑我们说：\Quote{你们要到几时才能停止不断在讲道中引入穷人和乞丐，向我们预言灾祸，谴责即将来临的贫穷，并想要使我们成为乞丐呢？} 我这样说并非要使你们成为乞丐，人啊，而是急于为你们开启天堂的财富。正如那向健康的人提及病人、述说其痛苦的人，这样说并非要使人生病，而是要藉着对他们灾祸的恐惧，断绝他们的懈怠，以保持健康。贫穷在你看来是可怕的事，甚至连其名字都令人畏惧。是的，因此我们是贫穷的，因为我们害怕贫穷，尽管我们拥有万贯财富。因为不是一无所有的人是穷人，而是那对贫穷战栗的人。因为在人们的灾祸中，我们哀叹并视为不幸的，也不是那些遭受大苦难的人，而是那些不知如何承受的人，即便苦难很小。而那知道如何承受的人，正如众人所知，是配得称赞和冠冕的。为了证明这一点，我们在竞技场上称赞谁呢？是那些被打得很重却不自寻烦恼、反而高昂着头的人，还是那些在第一击之后便逃跑的人？我们不是将前者视为勇敢而高贵的人来加冕，而嘲笑后者是懦弱胆怯的人吗？那么，让我们在生活的事务中也这样做：让我们给那轻易承受一切的人加冕，如同我们给那勇敢的斗士加冕；而哀悼那退缩、害怕危险、在挨打之前就已恐惧而死的人。因为在竞技场上，如果有人还未举起双手，一看到对手伸出右手就逃跑，即使没有受伤，也会因软弱、娇气、不谙此等搏斗而被嘲笑。这正是那些害怕贫穷、甚至无法忍受对贫穷的期待之人的写照。\par

显然，使我们不幸的不是我们，而是你们自己。因为魔鬼怎能不从此嘲笑你们，见你们在挨打之前就已对威胁感到恐惧和战栗？或者更确切地说，当你仅仅视之为一种威胁时，他甚至无需再打击你，而是让你保留财富，却藉着对失去财富的期待，使你变得比蜡还软。因为我们的天性（可以这么说）是在遭受之前，我们所恐惧的对象并不像在遭受之后或未经考验时那样可怕。因此，为了阻止你获得这种德行，他将你囚禁在恐惧的顶峰；藉着对贫穷的恐惧，在你还未经历它之前，就使你如蜡在火中熔化。是的，这样的人比蜡还软，过着比加音本人更悲惨的生活。因为他为所拥有的过多而恐惧，为所没有的而悲伤；此外，他又为所拥有的而战栗，将财富像逃跑的奴隶一样囚禁在内，被我不知道的各种莫名的激情所困扰。因为莫名的欲望、多样的恐惧和焦虑、以及四面八方的战栗，搅扰着他。他们就像一艘船，受到来自各方的逆风，经历着重重巨浪。对这样的人来说，死去比忍受持续的风暴要好得多。因为对加音来说，死去也比永远战栗更可忍受。\par

因此，为了避免我们也遭受这些事，让我们嘲笑魔鬼的诡计，让我们挣断他的绳索，斩断他可怕长矛的尖端，并加固每一个入口。因为如果你嘲笑金钱，他就无处可击，无处可抓。然后你就拔掉了万恶之根；当根不再存在时，也就不会长出任何恶果。\par

9. 好：我们总是说这些事，并且从不停止说；但我们的说教是否有益，那日将会揭晓——即那被火揭示的日子，\BibleRef{格前 3:13}，它检验每个人的工程，显示哪些灯是明亮的，哪些不是。那时，有油的和无油的都将显明。但愿没有人在那时被发现缺乏安慰；相反，愿所有人都带着丰富的仁慈，使灯明亮，与新郎一同进入。\par

因为没有什么比那将要听到新郎说\Quote{我不认识你们}的声音更可怕、更充满痛苦了——这是那些没有大量施舍的人所听见的。\BibleRef{玛 25:12} 但愿我们永远不要听到这声音，而是听到那最愉快、最可羡慕的声音：\Quote{我父所祝福的，你们来吧！承受自创世以来为你们预备的国度吧。} \BibleRef{玛 25:34} 因为这样我们将享受幸福的生活，并拥有所有超过人理解的福乐：愿我们都能达到，藉着恩宠和仁慈等。\par

\SourceRecord{Translated by Talbot W. Chambers.；Nicene and Post-Nicene Fathers, First Series；Vol. 12.；Edited by Philip Schaff.；Buffalo, NY: Christian Literature Publishing Co.,；1889.；<http://www.newadvent.org/fathers/220130.htm>.}

% structure: p0001 source=h1 target=section reason=page-title

\section{格林多前书第三十一篇讲道}

% structure: p0002 source=h2 target=subsection reason=relative-heading-rank; base=h2

\subsection{\BibleRef{格前 12:21}}

\BibleQuote{眼不能对手说：我无需你；头也不能对脚说：我无需你。}\par

他抑制了地位较低者的嫉妒，并消除了他们因他人被赐予更大恩赐而可能感到的沮丧之后，也谦卑了那些领受了更大恩赐者的骄傲。他在先前对前者的讲论中确实也做了同样的事。因为说那是恩赐而非成就，本意就是表明这一点。但如今他更加有力地再次这样做，停留在同一个比喻上。因为从身体及由此产生的合一，他进而对肢体进行实际比较，这正是他们特别渴望得到教导的事。因安慰他们的力量不在于他们都是同一个身体这一事实，而在于确信在所受的恩赐上，他们并未大大落后。他接着说：\BibleQuote{眼不能对手说：我无需你；头也不能对脚说：我无需你。}\par

因为恩赐虽小，却是必须的；正如当某一部分缺失时，许多功能便受阻碍，同样，没有另一部分，教会的圆满也有残缺。他说的不是\Quote{不会说}，而是\Quote{不能说}。因此，即使它愿意说，即使它真那样说了，也是不可能的，且与本性不符。为此，他取了两个极端来论证，首先是手与眼，其次是头与脚，以加强例证的力量。\par

还有什么比脚更卑微的呢？还有什么比头更尊贵、更必需的呢？因为头比任何其他部分都更能代表人。然而，它本身并不足够，也不能独自完成所有事情；因为若是这样，我们的脚就成了多余的附加物。\par

2. 他并未就此止步，还寻求另一种放大，这是他惯常的做法，不仅力求平等，甚至更进一步。因此他又说：\par

% structure: p0008 source=h2 target=subsection reason=relative-heading-rank; base=h2

\subsection{\BibleRef{格前 12:22}}

\BibleQuote{不但如此，身上那些似乎较为软弱的肢体，更是必需的：}\par

% structure: p0010 source=h2 target=subsection reason=relative-heading-rank; base=h2

\subsection{\BibleRef{格前 12:23}}

\BibleQuote{并且我们身上那些我们认为不大体面的，越发增添体面；我们不端正的，越发端正。}\par

在每一句中，他都加上\Quote{身体}这一词，从而既安慰了前者，又遏制了后者。他说：\Quote{我并非只说这一点，大的需要小的，而且大有需要。} 他所说的\Quote{似乎}和\Quote{我们认为也}很恰当，表明这判断并非出于事物的本性，而是出于多数人的意见。因为在我们身上没有什么是可耻的，因为它是天主的工程。既然如此，在我们身上还有什么比生殖器官更不被看重的呢？然而它们却享有更大的尊荣。甚至最贫穷的人，即使身体其余部分裸露，也不能容忍那些部位裸露。然而这决不是可耻之物的状况；相反，它们本应比其余部分更受轻视。因为在家中，一个受羞辱的仆人，不仅得不到更多的关注，甚至不能获得同等的待遇。同理，如果这个肢体是可耻的，它就不该享有特权，甚至不该享有同等权利；但现在它却得到了更多的尊荣：这也是天主的智慧所成就的。因为祂使某些部分天生不需要尊荣；但对于另一些部分，祂虽未赋予它们天生的尊荣，却迫使我们给予它们尊荣。然而它们并不因此可耻。因为动物也天生自足，其中大部分不需要衣服、鞋子和屋顶；但我们的身体并不因此比它们缺少尊荣，因为它需要这一切。\par

诚然，若有人仔细考虑，这些部分按其本性本身既尊贵又必要。事实上，\Person{保禄}本人也效法了这一点，他的判断有利于它们，不是出于我们的关心和它们享有更大尊荣，而是出于事物的本性。\par

因此，当他称它们为\Quote{软弱}和\Quote{不大体面}时，他用了\Quote{似乎}这个说法；但当他称它们为\Quote{必需}时，他不再加上\Quote{似乎}，而是自己作出判断，说\Quote{它们是必需的}，这非常恰当。因为它们是生育儿女和延续人类所必需的。因此，罗马立法者惩罚那些阉割这些肢体、使人成为阉人的，视之为损害我们共同血统、冒犯本性本身的人。\par

但那些放荡、亵渎天主手工的人有祸了。因为许多人惯于因醉酒而诅咒酒，因不洁而诅咒女性；同样，他们也因那些不善用这些肢体的人而视这些肢体为卑贱的。但这是不当的。因为罪并非作为本性的一部分而附着于事物，而是由于冒险行事者的意志产生了过犯。\par

但有些人认为，\Quote{软弱的肢体}、\Quote{不大体面的}、\Quote{必需的}和\Quote{享有更丰盛尊荣的}这些说法，是\Person{保禄}用来指眼睛和脚的，并且他称眼睛为\Quote{更软弱的}和\Quote{必需的}，因为它们虽然力量不足，却在功用上占有优势；但称脚为\Quote{不大体面的}，因为脚也受到我们极大的关注。\par

3. 接着，为了再提出另一个放大，他说：\par

% structure: p0018 source=h2 target=subsection reason=relative-heading-rank; base=h2

\subsection{\BibleRef{格前 12:24}}

\BibleQuote{但我们端正的肢体并不需要。}\par

这就是说，免得有人说：\Quote{为什么说这种话，竟然轻看受尊敬的，而讨好那较少受尊敬的？}\Quote{我们这样做并非出于轻视，}他说，\Quote{而是因为他们\InnerQuote{没有需要}}。并看他如何在此简要地设下多大的赞美，然后迅速推进：这是最适宜且有用的事。他还不满足于此，还加上原因，说：\Quote{但天主将身体调和在一起，把更丰盛的荣誉加给那有欠缺的部分：}\par

% structure: p0021 source=h2 target=subsection reason=relative-heading-rank; base=h2

\subsection{格前 12:25}

\BibleQuote{使身体没有\OriginalTerm{分裂}{schism}}\par

如果天主将之调和在一起，祂就不让那更不雅观的显现出来。因为混合之物成为一体，之前是什么样子就看不出来了：否则我们就不能说它是调和了的。并看他如何不断掠过缺陷，说：\Quote{那有欠缺的}。他说没有\Quote{给那受羞辱的}，\Quote{给那不体面的}，而是说给那有欠缺的（\Quote{那有欠缺的}——怎样？出于本性），赐予更丰盛的荣誉。为什么呢？\BibleQuote{使身体没有\OriginalTerm{分裂}{schism}}。因此，虽然他们享有无穷的安慰，却仍像比别人领受得少那样沉溺于悲伤，他表示他们反倒更受尊荣。因为他的话是：\BibleQuote{把更丰盛的荣誉加给那有欠缺的。}\par

接着他也加上原因，表明为了他们的益处，祂既使那部分有欠缺，又更丰盛地尊荣它。这原因是什么呢？他说：\BibleQuote{使身体没有\OriginalTerm{分裂}{schism}}，（他说\Quote{在肢体中}，而是说\Quote{在身体中}。）因为如果一些肢体既由本性又由我们的预先照顾得到看顾，而别的肢体连一样也得不到，那就会有巨大而不公平的有利条件。那么它们就会因为无法忍受连接而彼此分离。这些被分离后，其余的也会受害。你看到他如何指出，必然有\Quote{更大的荣誉}赐给\Quote{那有欠缺的}？他说：\Quote{否则，伤害就会成为众人的共同伤害。}原因在于，除非这些部分得到我们极大的考虑，它们就会因没有本性帮助而受到粗暴对待；这粗暴对待会导致它们的毁灭；它们的毁灭会分裂身体；身体分裂后，其他远比它们更大的肢体也会灭亡。\par

你看见吗？对这些后者的关心与为那些前者提供照顾是相连的。因为它们的存有并不那么在于自己的本性，而在于借着身体成为一体。所以如果身体灭亡，它们因各自拥有的健康也得不到益处。但如果眼睛或鼻子保留着，保持着它的功能，一旦联合的纽带破裂，此后它们将永远无用；而假若纽带保留，那些肢体受伤，它们既能通过它支持自己，又能迅速恢复健康。\par

但也许有人会说：\Quote{这确实在身体里有道理，即\InnerQuote{那有欠缺的领受了更丰盛的荣誉}，但在人间如何能如此呢？}为什么？在人间你尤其可以看到这事发生。因为那些在第十一时辰来的先领了工资；迷失的羊使牧人撇下九十九只，跑去追它，找到后将它扛在肩上，并没有驱赶；荡子比那循规蹈矩的得了更多的荣宠；盗贼得冠冕，在宗徒们之前被宣布。在塔冷通的比喻中，你也看到这事发生：领了五个塔冷通的和领了两个塔冷通的被赐予同样的赏报；是的，正因他领了两个，他就更蒙受极大的眷顾。因为如果他受托了五个，以他能力的欠缺，他会丧失一切；但他领了二个并尽了自己的责任，就被认为与赚了五个的有同样的价值，从而获得一个优势，即以较少的劳动得到同样的冠冕。而且他也是一个人，与那做五个塔冷通交易的一样。尽管如此，他的主绝不严格追究他，也不迫使他与同伴做同样的事，也不说：\Quote{你为什么不能赚五个？}（虽然祂本可以公正地说），但仍同样赐给他冠冕。\par

4. 因此，你们这些较大的，不要轻看较小的，免得你们反而伤害了自己。因为当它们被割断时，整个身体就毁了。既然，身体若非众多肢体的存在，又是什么呢？正如保禄自己所说：\BibleQuote{身体不是一个肢体，而是许多。}因此，如果这就是身体的本质，我们就当注意使许多肢体继续为许多。因为除非这一点完全保持，打击就会伤及要害；这也是为什么宗徒不只要求这一点，即它们不被分开，而且还要求它们紧密相连。例如，他说了\BibleQuote{使身体没有\OriginalTerm{分裂}{schism}}后，还不满足于此，又加上了\BibleQuote{肢体应彼此同样关心。}加上这另一个原因，也是为了使较小的享有更大的荣誉。因为天主这样安排，不仅是为了避免它们彼此分离，而且也是为了能有丰富的爱和和谐。因为如果每个人的存在都依赖于邻人的安全，就不要跟我提大和小：在这种情况下，没有大和小。只要身体继续存在，你仍可看到差别，但身体一灭亡，就再也看不到了。除非较小的部分也继续存留，否则身体就会灭亡。\par

若现在连较大的肢体，当较小的被折断时，也会毁坏，这些较大的就应当同样照顾较小的，如同照顾自己一样，因为这些较小的安全中，较大的也得以存留。所以，即使你说了万次：\Quote{这样的肢体是受羞辱的和较卑微的，} 但你若不像照顾自己一样照顾它，因它较卑微而忽视它，伤害就会临到你身上。因此，他不仅说\Quote{肢体应当彼此照顾，} 还加上了\Quote{要彼此同样照顾，} 也就是说，小的应当与大的享有同样的上智安排。\par

切不要说，那是一个普通人，但要思考，他是那结合整个身体的肢体之一。正如眼睛，他也使身体成为身体。因为在身体被建立的地方，没有人比邻人拥有更多，因为身体不是由一个部分大而另一个部分小造成的，而是由许多不同的部分造成的。因为正如你，因你较大，有助于构建身体，同样他，因他较小，也有助于构建身体。所以，当身体要被建立时，他的相对不足，在如此高贵的贡献中变得与你有同等的价值；是的，他与你一样有用。这可以从以下事情明显看出。假设没有肢体较大或较小，也没有较多或较少尊贵，而是全部是眼睛或全部是头：身体岂不会毁灭？人人都能看见。同样，如果全部都是较小的，也会发生同样的事。所以，在这方面，较小的也被证明是平等的。是的，如果有人要说更多的话，较小的之所以较小，其目的正是为了使身体得以存留。所以，为了你的缘故，他较小，好使你继续保持大。这就是他要求所有人同样照顾的原因。当他说了\Quote{肢体要彼此同样照顾} 之后，他再次解释\Quote{同样的事}，说道：\par

% structure: p0030 source=h2 target=subsection reason=relative-heading-rank; base=h2

\subsection{\BibleRef{格前 12:26}}

5. \BibleQuote{若一个肢体受苦，众肢体一同受苦；若一个肢体得光荣，众肢体一同欢乐。}\par

\Quote{是的，再无其他目的，} 他说，\Quote{祂使祂所要求的照顾成为共同的事，在如此大的多样性中建立合一，为的是在一切事件中都有完全的共融。因为，如果我们对近人的照顾是共同的保障，那么我们的光荣和悲伤也必须是共同的。} 因此，他在这里要求三件事：不分裂而在完美中合一；彼此有同样的照顾；以及视一切所发生的事为共同的。如上面他所说：\Quote{祂赐予那缺乏的部分更丰厚的尊荣，} 因为它需要它；表明那卑微本身成了更大尊荣的入口；同样，在这里，他也使他们因彼此之间的照顾而平等。因为\Quote{因此祂使他们分享更大的尊荣，} 他说，\Quote{为使他们不会遇到较少的照顾。} 并且不仅从这一点，而且从他们所遭遇的一切好事和痛苦中，肢体彼此相连。例如，当一根刺扎在脚后跟上时，整个身体都感觉到并照顾它：背部弯曲，腹部和大腿收缩，手像卫士和仆人一样伸出，拔出那刺，头低下，眼睛极为关心地注视着它。所以，即使脚因不能抬高而有卑微，但通过使头低下，它获得了平等，并受到同样的尊荣；尤其是每当脚导致头低下时，不是由于恩惠，而是由于脚对头的要求。因此，如果较尊贵的肢体有优势，那么，既如此，它就欠较小的肢体这样的尊荣和照顾，以及同样的同情：这表明了很大的平等。因为有什么比脚后跟更卑贱呢？有什么比头更尊贵呢？然而这个肢体伸到那个肢体，并带动它们与自己一起行动。再者，如果眼睛有了问题，大家都抱怨，大家都无所事事：脚不走，手不工作，胃也不享用惯常的食物；然而病症在于眼睛。你为什么使胃挨饿？为什么让脚静止？为什么束缚手？因为它们与脚相连，整个身体以不可言喻的方式一同受苦。因为如果它没有分担痛苦，它就不会忍受分担照顾。所以，说了\Quote{肢体要彼此同样照顾} 之后，他加上\BibleQuote{若一个肢体受苦，众肢体一同受苦；若一个肢体得光荣，众肢体一同欢乐。} \Quote{它们如何一同欢乐呢？} 你说。头被加冕，整个人受到尊荣。口说话，眼睛欢笑并喜悦。然而功劳不属于眼睛的美丽，而属于舌头。同样，如果眼睛看起来漂亮，整个女人就被装饰：正如这些肢体一样，当笔直的鼻子、挺直的脖颈和其他肢体被赞美时，它们也欢乐并显得愉快；再者，它们为悲伤和不幸大量流泪，虽然它们自己并未受伤。\par

6. 那么，我们众人思考这些事，要效法这些肢体的爱；绝不可做相反的事，践踏邻人的痛苦，嫉妒他的好事。因为这是疯子和神志不清的人的行为。正如挖出自己眼睛的人表现出极大的愚蠢；吞噬自己手的人展示出彻头彻尾疯狂的明显证据。\par

现在，如果肢体间有这样的情况，那么同样地，当它发生在弟兄们中间时，它将愚蠢的名声加于我们，并带来不小的祸害。因为他发光的时候，你的美貌也显出来，整个身子就美丽了。他并不把美丽只限于自己，也容许你荣耀。但如果你熄灭他，你就给整个身子带来共同的黑暗，你造成的灾祸是众肢体所共有的：的确，如果你保持他明亮，你就保持整个身体的青春。因为没有人说：\Quote{眼睛是美丽的：}而是说什么？\Quote{这样的女人是美丽的。}而且如果它也受赞美，那是在公共的颂赞之后。同样地，这事也发生在教会中。我是说，如果有任何杰出的人，团体就分享其美誉。因为敌人不容易分开赞美，而是把它们连在一起。如果有人口才出众，他们不只赞美他，也赞美整个教会。因为他们不只说：\Quote{这样的人是一个奇妙的人，}而是说什么？\Quote{基督徒有一位奇妙的老师：}所以他们使这拥有成为公共的。\par

7. 现在我要问，外邦人彼此联结，而你们却分裂并攻击自己的身子，反对自己的肢体吗？你们不知道这推翻一切吗？因为他说：\BibleQuote{凡一国自相纷争，必将灭亡。}\BibleRef{玛 12:25}\par

但没有什么像嫉妒和妒忌那样分裂和隔绝，那是严重的疾病，完全不可赦免，在某些方面比\BibleQuote{万恶的根源}更坏\BibleRef{弟前 6:12}。因为贪财者在自己得到时快乐；但嫉妒者在别人没有得到时快乐，而不是在自己得到时快乐。因为他认为别人的不幸是对自己的益处，而不是别人的兴盛；他四处走动，是人类的公敌，打击基督的肢体，还有什么比这更接近疯狂呢？魔鬼是嫉妒的，但嫉妒的是人，不是别的魔鬼；但你作为人却嫉妒人，并反对自己同族和同家的人，这连魔鬼也不做。你将得到什么赦免，什么借口？看见弟兄兴盛就战栗变色，而你本应为自己加冕、欢喜和雀跃。\par

如果你确实想效法他，我不禁止；但效法是为了像那被赞许的人一样：不是要压制他，而是为达到同样的崇高境界，显示同样的卓越。这是健康的竞争，无争端的模仿：不为别人的好事悲伤，而为我们自己的坏事烦恼；相反的结果则是嫉妒。因为嫉妒忽视自己的坏事，却在别人的好运中消瘦。因此，穷人不因自己的贫穷而烦恼，却因邻人的富足而烦恼；还有什么比这更令人痛苦呢？是的，在这方面，嫉妒者，如我先前所说，比贪财者更坏；一个因自己的所得而喜乐，另一个则以别人得不到为乐。\par

因此，我恳求你们，离开这条邪恶的路，转向正确的效法（因为这种热忱是猛烈的，比任何火都热），并藉此赢得大福。同样，\Person{保禄}也用来引导那些犹太血统的人进入信仰，他说：\BibleQuote{或者可以激动我骨肉之亲发愤，好救他们一些人。}\BibleRef{罗 11:14}因为那像\Person{保禄}所愿的那样效法的人，看到别人得荣誉时不消瘦，而是看到自己落后时；嫉妒者则不然，而是看到别人兴盛时。他像一只雄蜂，损害别人的劳动；自己从不急于上升，却看到别人上升时哭泣，并尽力把他拉下来。那么，这激情可比作什么呢？在我看来，好像一头迟钝而肥胖的驴子，与一匹有翼的骏马同轭，它自己不愿上升，却试图用自身体重把对方拖下来。因为这人毫不思索也不焦虑摆脱自己的沉睡，却做一切事来绊倒和拉下那飞向天堂的人，成为魔鬼的精确效法者：因为魔鬼看到人在乐园中，并不寻求改变自己的状况，而是把他赶出乐园。又看到他坐在天上，其余的人急忙赶往那里，他就坚持同样的计划，绊倒那些急忙前往的人，并因此为自己更多堆积火炉。因为在每件事上都如此：被嫉妒的人，若警惕，就更加卓越；嫉妒的人，为自己积累更多的恶。同样，\Person{若瑟}变得卓越，\Person{亚郎}司祭也是如此：嫉妒者的阴谋使天主一次又一次地为他投下赞成票，并成为棍杖发芽的契机。这样，\Person{雅各伯}获得了丰富的财富和所有那些祝福。这样，嫉妒者用万恶刺穿自己。我们知道这一切，让我们逃避这样的效法。因为，告诉我，你为什么嫉妒？因为你的弟兄领受了神恩？他从谁领受？回答我。难道不是从天主吗？显然，那恩赐的施予者就是你敌对的对象。你看到这邪恶的趋势了吗？它以什么样的尖端为你罪恶的堆加冠冕？它为你挖掘的报复之坑有多深？\par

那么，亲爱的，让我们逃避嫉妒，既不要嫉妒他人，也不要不为那些嫉妒我们的人祈祷，而要尽力扑灭他们的激情。我们不要像那些愚昧的人一样，他们一心想要惩罚嫉妒者，却竭尽全力点燃他们的火焰。但我们不要这样做；反而要为他们哭泣哀悼。因为他们才是受伤的人，常有虫子啃噬他们的心，积聚的毒液比任何苦胆更苦。来吧，让我们恳求仁慈的天主，既改变他们的心境，也使我们永不陷入这种疾病。因为，对于患有这种溃烂之症的人来说，天堂确实遥不可及；而在天堂之前，连今世的生活也不值得活。因为木头和羊毛被其中的蛀虫啃噬的程度，远不及嫉妒的狂热吞噬嫉妒者的骨髓、摧毁他们灵魂中一切自制力那样彻底。\par

那么，为了使我们都摆脱这无数的祸患，让我们从心中驱除这邪恶的狂热，这比任何坏疽更可怕的病；好使我们重获属神的力量，完成现世的赛程，并获得未来的冠冕；愿我们众人赖我们的主耶稣基督的恩宠和仁慈，得以到达那境界，愿光荣、权能、尊崇归于祂，偕同圣父及圣神，自今而往，及世之世。阿们。\par

\SourceRecord{Translated by Talbot W. Chambers.；Nicene and Post-Nicene Fathers, First Series；Vol. 12.；Edited by Philip Schaff.；Buffalo, NY: Christian Literature Publishing Co.,；1889.；<http://www.newadvent.org/fathers/220131.htm>.}

% structure: p0001 source=h1 target=section reason=page-title

\section{\WorkTitle{论格林多前书第三十二篇讲道}}

% structure: p0002 source=h2 target=subsection reason=relative-heading-rank; base=h2

\subsection{\BibleBook{格林多}{前书} 12:27}

你们便是基督的身体，各自作肢体。\par

因为恐怕有人说：\Quote{身体的比喻与我们何干？身体受本性驱使，而我们的善行出于选择；}他便将此应用于我们自身；为表明我们应当如身体出于本性那样同心合意，他说：\BibleQuote{你们便是基督的身体。}但如果我们的身体不应分裂，基督的身体更不应分裂，且恩宠越强大，越胜于本性。\par

但\Quote{各自\OriginalTerm{各自}{severally}}是什么意思？\Quote{至少就你们而言；且就本性言，应由你们组成一个部分。}因他说了\Quote{身体}，而整个身体并非格林多教会，乃是遍布全世界的教会，所以他用了\Quote{各自}：即你们中间的教会是普世教会及由所有教会组成的身体的一部分；因此，你们不仅当与自身和睦，也当与普世教会和睦，如果你们确是全身的肢体。\par

% structure: p0006 source=h2 target=subsection reason=relative-heading-rank; base=h2

\subsection{\BibleBook{格林多}{前书} 12:28}

2. \BibleQuote{天主在教会内所设立的：第一是宗徒，第二是先知，第三是教师，然后是异能，然后是治病的恩赐、援助、治理、各种语言。}\par

他如今所做的，正是我先前所说的。因为他们以语言自傲，他便将此恩赐置于末位。因他所说的\Quote{第一}与\Quote{第二}并非随意，乃顺序列举，指出较尊贵与较卑微者。故此他将拥有一切恩赐的宗徒放在首位。他并非简单地说\Quote{天主在教会内设立了某些人，宗徒}或\Quote{先知}，乃是用了\Quote{第一、第二、第三}，以表明我所告诉你们的那件事。\par

\BibleQuote{第二是先知。}因为他们曾发预言，如斐理伯的女儿们、阿加波、以及格林多人中的这些人，关于他们他说：\BibleQuote{至于先知，可以两个人或三个人说话。}\BibleRef{格前 14:29}他在致弟茂德的书信中也写道：\BibleQuote{不要疏忽你因预言所得的恩赐。}\BibleRef{弟前 4:14}而且发预言的人更多。若基督说：\BibleQuote{法律和先知直到若翰为止。}\BibleRef{玛 11:13}祂是指那些在祂来临之前先知的先知。\par

\BibleQuote{第三是教师。}因为发预言的人一切都出于圣神；但教导的人有时也发抒己见。故此他也说：\BibleQuote{善于督导的长老，应受加倍的敬奉，尤其是那些在言语和教导上劳苦的人。}\BibleRef{弟前 5:17}然而凡事出于圣神发言的人并不劳苦。因此他将教师置于先知之后：因为先知之恩纯粹是恩赐，而教师之职也包含人的劳苦。教师常凭自己的心思讲论许多事，但与圣经相符。\par

3. \BibleQuote{然后是异能，然后是治病的恩赐。}你看见他如何又将治病与异能分开，如同先前所做？因异能大于治病：有异能的人既能惩罚又能医治，而有治病恩赐的人只能治愈。注意他所作的次序何等卓越，将预言置于异能和治病之前。因为在上面他说\Quote{有人由圣神得了智慧的言语，另有人得了知识的言语}时，并非按次序，而是随意。但此处他标明了第一、第二等阶。为何他将预言放在首位？因为在旧约中也是如此。例如，当依撒意亚与犹太人辩论，显示天主的大能，并引证魔鬼的虚无时，他将预知未来当作神性更大的证据。\BibleRef{依 41:22}基督本人在行了诸多神迹后，也以此作为祂神性不小的标记；并不断补充说：\BibleQuote{这些事我已告诉你们，使你们在事情发生时，相信我就是那一位。}\BibleRef{若 13:19}\par

\Quote{那么，治病的恩赐次于预言是合理的。但为何也次于教导呢？}因宣讲圣言、在听者心中栽种虔诚，与行奇迹并非同一回事；奇迹仅为宣讲之故。因此，凡以言以行教导的人，大于一切。他把那些人特别称为\Quote{教师}，他们既以行为教导，又以言语训诲。例如：这使宗徒们自己成为宗徒。而那些恩赐也给了某些无甚价值的人，起初他们曾说：\BibleQuote{主，我们不是因祢的名发预言，行大能的事吗？}此后却被告诉：\BibleQuote{我从来不认识你们；离开我吧，你们这些作恶的人。}\BibleRef{玛 7:22}但这双重的教导方式——即通过言行——恶人决不会采用。至于他将先知放在首位，不要惊讶。因为他并非指一般先知，而是指那些借预言也教导、并为公众利益而言一切的人；这在后文会越发清楚。\par

\BibleQuote{援助、治理。}\Quote{援助}是什么？支持弱者。请问，这也是恩赐吗？首先，这本身也是天主的恩赐，即适合担任保护人职务的才能；分发属灵之事；此外，他甚至称我们许多善行为\Quote{恩赐}；并非要我们灰心，而是表明在一切事上我们都需要天主的助佑，并使他们感恩，借此激励他们，振奋他们的心神。\par

\BibleQuote{各种语言。}你看见他将此恩赐置于何处？他如何处处将它置于末位？\par

4. 此外，既然又通过这份清单他指出巨大的差异，并激起了之前那些恩赐较小者的不满情绪，他便在接下来的话中极其猛烈地抨击他们，因为他已经给了他们许多证据，表明他们并未被远远落下。我的意思是：因为他们听到这些事很可能会说：\Quote{为什么我们没有都成为宗徒？}——而之前他曾使用更为温和的语调，详细证明这一结果必然性，甚至借用了身体的比喻；因为他说：\BibleQuote{身体不是一个肢体；}又说：\BibleQuote{若全身是同一个肢体，身体在哪里？}以及它们是为使用而赐予的；因为他说：\BibleQuote{圣神的显示}赐给每人，\BibleQuote{为有益处。}并且都从同一位圣神受滋润；以及所赐的是恩惠而非债务；因为他说：\BibleQuote{神恩虽有区别，却是同一圣神。}并且圣神的显示同样通过众人显现；因为他说：\BibleQuote{每人蒙显示，都是藉着圣神。}以及这些事是按照圣神和天主的喜悦而形成的；因为他说：\BibleQuote{这一切都是同一位圣神所行的，他随心所欲分给每人。}又说：\BibleQuote{天主这样安置了肢体，每一个在身体上，按他的意愿。}并且较弱的肢体也是必要的；因为他说：\BibleQuote{那些似乎是较弱的肢体，却是必要的。}因为它们同样是必要的，因着大的也需要小的；因为他说：\BibleQuote{头不能对脚说：我不需要你。}而这些后者甚至享有更大的荣耀；因为他说：\BibleQuote{对于那缺乏的，他赐予更丰富的荣耀。}以及对它们的关怀是共同且平等的；因为他说：\BibleQuote{所有肢体彼此互相关照。}并且全体共享同一荣耀和同一痛苦；因为他说：\BibleQuote{若一个肢体受苦，所有肢体一同受苦；若一个肢体得荣耀，所有肢体一同欢乐。}——我说，他先前曾通过这些话题劝诫他们，而现在从此刻起，他使用语言来压倒并责备他们。因为正如我所说，我们既不应总是劝诫人，也不应总是使他们沉默。因此，\Person{保禄}本人，因为他终于劝诫了他们，从此便猛烈攻击他们，说：\par

% structure: p0016 source=h2 target=subsection reason=relative-heading-rank; base=h2

\subsection{\BibleRef{格前 12:29}}

\BibleQuote{众人都是宗徒吗？众人都是先知吗？众人都有治病的恩赐吗？}\par

他并不停留在第一和第二项神恩上，而是进行到最后一项，或是意指所有人都不能成为一切（正如他所说的：\BibleQuote{若全身是同一个肢体，身体在哪里？}），或是连同这些也确立另一个要点，这要点从另一方面又可以起到安慰作用。那么，这是什么？就是他指明，即使是较小的神恩也与较大的神恩同样被追求，因为连这些也没有绝对赐给所有人？因为他说：\Quote{为什么你因没有治病的恩赐而忧伤呢？想一想，你所有的，即使更为微小，却常常是那有更大恩赐者所没有的。}因此他说，\par

% structure: p0019 source=h2 target=subsection reason=relative-heading-rank; base=h2

\subsection{\BibleRef{格前 12:30}}

\BibleQuote{众人都会说方言吗？众人都会解释吗？}\par

因为正如天主没有把大恩赐赐给所有人，而是给一些人这个，给另一些人那个，同样在较小的恩赐上，他也没有将这些赐给所有人。他这样做，由此获得了深厚的和谐与爱，使每个人都因需要他人而亲近自己的兄弟。这种安排他也设立在艺术中，也设立在元素中，也设立在植物中，也设立在我们的肢体中，并绝对在所有事物中。\par

5. 然后他又添加了最有力的安慰，足以使他们恢复并平息他们烦恼的灵魂。这是什么？\par

% structure: p0023 source=h2 target=subsection reason=relative-heading-rank; base=h2

\subsection{\BibleRef{格前 12:31}}

他说：\BibleQuote{你们要切切渴慕更好的神恩；我还要给你们指示一条更卓越的路。}\par

他说这话，是温和地暗示他们，自己得到较小神恩是他们自己的原因，并且如果他们愿意，他们有能力得到更大的。因为当他说\BibleQuote{你们要切切渴慕}时，他要求他们尽一切努力并渴慕属神之事。他没有说\Quote{更大的神恩}，而是说\BibleQuote{更好的}，即更有用的、更有益处的。他的意思是：\Quote{继续渴慕神恩；我指示给你们神恩的泉源。}因为他没有说\Quote{一个神恩}，而是说\Quote{一条路}，以便更加颂扬他所要提及的。似乎他说：我不是指示给你们一个、两个或三个神恩，而是一条通向所有这些神恩的路；而且不只是一条路，更是\BibleQuote{一条更卓越的路}，一条向所有人普遍开放的路。因为不像神恩那样，赐给一些人这些，赐给另一些人那些，但不是所有人都得到一切；在这事上也是如此：但它是一个普遍性的恩赐。因此他也邀请所有人来获得它。他说：\BibleQuote{你们要切切渴慕更好的神恩，我还要给你们指示一条更卓越的路；}意思是爱德，即对邻人的爱。\par

然后，他打算讲论爱德并赞扬这一德行，便先通过比较贬低其他事物，暗示没有爱德它们便毫无价值；这是非常审慎的。因为若他立刻讲论爱德，并在说了\Quote{我指示给你们一条道路}之后，就加上\Quote{这就是爱德}，而不通过比较来引导论述，那么有些人可能会嘲笑所说的话，因为没有清楚理解所言之事的力道，却仍对这些事物垂涎。因此，他没有立刻展开，而是先用应许激发听者的兴趣，说：\Quote{我指示给你们一条更优越的道路}，并带领他们渴望它之后，也没有立刻就此进行，而是进一步扩大和延展他们的渴望，先讲论这些事物本身，并表明没有爱德它们便毫无价值；迫使他们极度需要彼此相爱；同时也看到，正是由于忽略爱德才产生了导致他们所有邪恶的根源。所以，在这方面，它也可以理当显得伟大，如果这些神恩不仅没有使他们团结，反而使已团结的他们分裂：而爱德却会在许多人如此分裂时，凭自身的力量使他们重新合一，成为一体。然而，他没有立刻这样说，而是将他们主要渴望的事写下来：即爱德是一件礼物，是通向所有神恩的最优越的道路。因此，即使你不愿因友谊而爱你的兄弟，也要为获得更美的标记和丰厚的恩赐，而珍爱爱德。\par

6. 请注意他从何处开始：从他们眼中惊异而伟大的事物，即语言神恩。在提出这一神恩时，他并不是仅按他们拥有的程度提及，而是远超于此。因为他并没有说：\Quote{我若说语言}，而是——\par

% structure: p0028 source=h2 target=subsection reason=relative-heading-rank; base=h2

\subsection{\BibleRef{格前 13:1}}

\Quote{我若说人类的语言——}\par

\Quote{人类的}是什么意思？指全世界所有民族的。他并不满足于这种扩大，还用了另一种更大的扩大，加上这些话：\Quote{并有天使的语言——却没有爱德，我就成了响亮的铜，或叮当作响的钹。}\par

你看到他先是将神恩提升到何种高度，然后又将其降到何种低处并贬低？因为他并非简单地说\Quote{我算不了什么}，而是说\Quote{成了响亮的铜}，一种无感觉无生命的东西。但\Quote{响亮的铜}如何？的确发出声音，却是随意而无益，毫无正当目的。因为除了我毫无益处之外，大多数人都认为我是一个惹人讨厌的人，一种令人厌烦和厌倦的人。你看到没有爱德的人如何与无生命无感觉的东西相似吗？\par

现在他在这里说\Quote{天使的语言}，并不是赋予天使肉体，而是他的意思是：\InnerQuote{即使我能像天使习惯彼此交谈那样说话，若没有爱德，我也算不了什么，反而是一种负担和烦恼。}正如（再举一例）他说：\Quote{一切在天上的、地上的和地下的，因耶稣之名，万膝都要屈伏}\BibleRef{斐 2:10}，他说这些话并非赋予天使膝盖和骨头，绝非如此，而是以我们中间的方式来表达他们朝拜的强烈程度；同样，在这里他称它为\Quote{语言}，并非意指肉体的器官，而是意图以我们中间所知道的方式暗示他们彼此的交谈。\par

7. 然后，为使他的论述令人接受，他不停留在语言神恩上，而是继续论及其他神恩；并贬低所有在缺乏爱德时的神恩，然后描绘爱德的形象。因为他更愿意通过扩大法来论述，便从较小的事物开始，上升到更大的事物。原来，当他指出它们的顺序时，他把语言神恩放在最后，现在却把它排在第一位；正如我说，逐渐上升到更大的神恩。这样，在讲了语言之后，他立刻接着讲先知话；并说：\par

% structure: p0034 source=h2 target=subsection reason=relative-heading-rank; base=h2

\subsection{\BibleRef{格前 13:2}}

\Quote{并且我若有先知话之神恩。}\par

这神恩也带着卓越性。因为正如在那个例子中，他提到的不只是语言，而是全人类的语言，并进一步提到天使的语言，然后表明没有爱德这神恩毫无价值；同样，在这里他提到的不仅是一般的先知话，而是最高的先知话：在说了\Quote{我若有先知话}之后，他加上\Quote{并明白各样的奥秘和各样的知识}；也以强烈的方式表达这神恩。\par

此后，他又继续论及其他神恩。并且，为了避免显得使他们厌烦，他没有逐一列举每一种神恩，而是列出了万母和万源，并以卓越的方式，这样说：\Quote{并且我若有全部的信德。}他还不满足于此，还加上了基督所说的最大的事，说：\Quote{甚至能移山，却没有爱德，我算不了什么。}请考虑，他在这里又如何贬低了语言神恩的尊严。因为就先知话而言，他指出了它所带来的巨大好处，\Quote{明白奥秘，并有全部知识}；就信德而言，则是非平凡的事，即\Quote{移山}；而就语言而言，他只提了神恩本身，便放弃了它。\par

但我求你也要思考这一点：他如何以简短的方式概括了所有神恩，当他提到先知话和信德时：因为奇迹或是言语上的，或是行为上的。基督如何说，信德的最小程度就是能够移山？因为他说：\Quote{假如你们有信德像芥子那么大，你们向这座山说：从这边移到那边去！它必会移去；}\BibleRef{玛 17:20}，祂似乎是在说一件非常小的事；而保禄却说这是\Quote{全部的信德}。那么，该说什么呢？因为这是一件大事，移山，所以他也提到它，并非因为\Quote{全部的信德}只能做这件事，而是因为由于外在物体的巨大，这在较粗俗的人看来是伟大的，因此他也藉此来颂扬他的主题。而他说的是：\par

\Quote{我若有全备的信德，甚至能移山，但我若没有爱德，我算不了什么。}\par

% structure: p0040 source=h2 target=subsection reason=relative-heading-rank; base=h2

\subsection{格林多前书 13:3}

8. \Quote{我若把我所有的财产全施舍给人，并舍身投火被焚，但我若没有爱德，为我毫无益处。}\par

奇妙的增补！因为连这些事他也附加了另一层意思：他没有说\Quote{我若把一半财产施舍给穷人}，或\Quote{两份或三份}，却说\Quote{我若把所有的财产全施舍}。他也没有说\Quote{施舍}，而是说\Quote{分施碎块}，以便在花费之外还加上用心管理的操劳。\par

但即便如此，我仍未指明这卓越性的全部，直到我提出基督关于施舍和死亡的证言。那么祂的证言是什么？祂对富人说：\Quote{你若愿意成为成全的，去变卖你所有的，施舍给穷人，然后来跟随我。}\BibleRef{玛 19:21}同样，论及爱近人时，祂说：\Quote{人若为自己的朋友舍命，再没有比这更大的爱了。}\BibleRef{若 14:13}由此显然可见，即使在天主面前，这也是最大的。但保禄说：\Quote{我告诉你们，即使我们为天主的缘故舍命，且不仅是舍命，更是被焚（因为\InnerQuote{我若舍身被焚}的意思正是如此），若不爱我们的近人，我们也不会有什么大益处。}那么，说神恩若无爱德便无大益，这并不稀奇：因为我们的神恩次于我们的生活方式。无论如何，许多人显示了神恩，却因变得邪恶而受罚：如那些\Quote{因祂的名预言、驱逐许多魔鬼、行了许多大能}的人；又如叛徒犹达斯。而另一些人作为信徒，展现了纯洁的生活，无需别的事物就能得救。因此，神恩如我所说需要爱德，这不稀奇；但严谨的生活若无爱德也毫无益处，这才强烈地突出了这话的紧张感并引起极大的困惑。尤其因为基督似乎将祂伟大的赏报同时归于这两者：即放弃财产和殉道的危险。因为正如我先前所说，祂对富人说：\Quote{你若愿意成为成全的，去变卖你的财产，施舍给穷人，然后来跟随我；}并且对门徒论及殉道说：\Quote{凡为我的缘故丧失自己生命的，必得着生命；}又说：\Quote{凡在人面前承认我的，我在我天上的父面前也必承认他。}因为这成就的劳苦确实伟大，几乎超越本性本身，凡获得这些荣冠的人都深知此事。因为没有任何言语能向我们表达：这行为属于何等尊贵的灵魂，又是何等超凡奇妙。\par

9. 然而，保禄却说，即使这如此奇妙的事，若无爱德，且即使与放弃财产相连，也并无大益。那么他为何这样说？我现在将努力解释，首先询问：一个把自己的全部财产施舍给穷人的人，怎么可能缺少爱德？我承认，那准备被焚且有神恩的人，或许可能没有爱德；但那人不仅施舍财产，而且甚至分施碎块，他怎会没有爱德？那么我们要说什么？或者是说他把假设当作事实来陈述；这是他惯常的做法，当他打算向我们展示某种超绝之事时；正如他写信给迦拉达人时说：\Quote{即使我们或从天上来的天使，给你们宣讲与我们从前所宣讲的不同的福音，当受诅咒。}\BibleRef{迦 1:8}然而，无论是他自己还是天使都不会这样做；但为了表明他要尽可能把问题推至极点，他列举了绝无可能发生的事。又如他写信给罗马人时说：\Quote{无论是天使、掌权者、异能者，都不能使我们与天主的爱隔绝；}因为天使也不会做这事；但这里他也假设了不可能的事；正如在后文中所说：\Quote{任何受造之物，}然而并没有别的受造之物，因为他已包括整个受造界，论及了上下一切。但这里他仍然以假说方式列举了不存在的事，以表明他强烈的渴望。这里他也做了同样的事，说：\Quote{即使人施舍一切，若没有爱德，也毫无益处。}\par

那么，我们或者可以说这点，或者他的意思是：施舍的人还应与那些隐退的人紧密结合，而不仅仅是无同情地施舍，而是怀着怜悯和谦逊，俯身与贫困者一同忧伤。因为天主设立施舍也正是为此：天主本来也可以不靠施舍而养活穷人，但为使我们彼此联结于爱德，并彼此热切相待，祂命令他们由我们供养。因此，另有一处也说：\Quote{良言胜过馈赠；}\BibleRef{德 18:16}又说：\Quote{看，言语超过良好的馈赠。}\BibleRef{德 18:16}祂自己说：\Quote{我喜欢仁爱，而非祭献。}\BibleRef{玛 9:30}因为人们通常爱那些受他们恩惠的人，而受惠者也更友善地对待施惠者，所以祂制定了这条法律，作为友谊的纽带。\par

10. 但上文提出的探究点却是：基督既说两者都属于成全，保禄又怎能断言没有爱德它们就不完全？这不是矛盾，绝对不可：而是与祂和谐一致，且十分精确。因为关于那个富人，祂不仅说\Quote{变卖你的家产，分给穷人}，还加上了\Quote{并且来，跟随我}。然而，跟随祂本身并不能使人完全成为基督的门徒，正如彼此相爱一样。因为祂说：\Quote{众人将因此认出你们是我的门徒：如果你们彼此相爱。}\BibleRef{若 13:35} 又说：\Quote{谁为我的缘故丧失自己的生命，必会获得生命}\BibleRef{玛 10:39}；以及\Quote{凡在人前承认我的，我在天之父也必承认他}；这并不是说不需要爱德，而是宣告这些劳苦的赏报。因为在别处祂强烈暗示殉道同时也要求爱德，这样说：\Quote{你们固然要喝我的杯，也要受我所受的洗}\BibleRef{玛 20:23}，即你们要成为殉道者，为我的缘故被杀；但坐在我的右边和左边（并非真有右边左边之分，而是指最高的优先和尊荣）\Quote{不是我可以赐的，而是给谁预备的，就给谁。}然后指示为谁预备，祂召叫他们说：\Quote{你们当中谁愿作首领，就当作众人的仆人}\BibleRef{玛 20:26}，表明了谦卑和爱。祂要求的是强烈的爱；所以祂不止于此，还加上了：\Quote{正如人子来不是受服事，而是服事人，并交出自己的生命作许多人的赎价}；指出我们应当如此相爱，甚至为所爱的人被杀。因为这首先是爱祂。因此祂也对伯多禄说：\Quote{如果你爱我，就牧放我的羊群。}\BibleRef{若 21:16}\par

11. 为使你明白爱德是何等伟大的德行，让我们用言语描绘它，因为在行为中我们已无处可见；让我们思考，如果它处处充沛，会有多大的好处：那时就不需要法律、法庭、刑罚、报复或任何其他此类事物了；因为如果人人都爱且被爱，就没有人会伤害他人。是的，谋杀、争斗、战争、分裂、掠夺、欺诈和一切邪恶都将被移除，罪恶甚至其名都将不为人知。然而奇迹却不能成就此事；它们反而使那些不警惕的人变得骄傲，趋向虚荣和愚妄。\par

再者：爱德真正奇妙之处在于：所有其他的善事都伴随着它们的恶；例如，舍弃财产的人常常因此自高自大；口才出众的人为虚荣所困扰；谦卑的人，正因如此，内心常自觉不错。但爱德却免于一切这样的弊病。因为没有人能对自己所爱的人感到骄傲。我恳求你，不要只假设一个人爱，而是所有人都同样爱；这样你就能看见它的美德。或者，如果你愿意，先假设一个被爱的人和一个爱他的人，但爱得恰当。他在地上将如在天上生活，处处享受宁静，为自己编织无数冠冕。因为这样的人将从嫉妒、愤怒、猜忌、骄傲、虚荣、邪恶的贪欲、一切不洁的爱、一切紊乱中保守自己的灵魂纯洁。是的，正如没有人会伤害自己，这个人也不会伤害他的邻人。如此之人，即或行走在地上，也将与加俾额尔本身同列。\par

这样的人就是拥有爱德的人。但那行奇迹、有完美知识却没有爱德的人，即使他使万人从死者中复活，也不会获益多少，因为他与众人隔绝，不愿与任何同仆交往。基督正是为此才说，对祂完美之爱的标记就是爱邻人。因为祂说：\Quote{如果你爱我，伯多禄，比这些更深，就牧放我的羊群。}\BibleRef{若 21:15} 你看到祂如何又借此暗中表明，在什么情况下这（爱德）比殉道更伟大吗？因为如果有人有一个爱子，他愿为这孩子舍弃生命，而另一个人爱父亲却丝毫不顾儿子，那么他会大大激怒父亲；父亲也不会感受到那人对他本人的爱，因为忽视了儿子。如果在父子之间如此，那么在神与人之间更是如此，因为天主当然比任何父母更慈爱。\par

因此，祂说了\Quote{第一条且最大的诫命是：你当爱上主你的天主}后，又加上\Quote{第二条——祂没有略过，而是也写下来——也与此相似：你当爱你的邻人如你自己。}且看祂如何以几乎同样的卓越也要求这一条。因为关于天主，祂说\Quote{全心}；关于邻人，祂说\Quote{如你自己}，这等同于\Quote{全心}。\par

是的，如果这条诫命被充分遵守，就不会有奴隶与自由人、统治者与被统治者、富人与穷人、大与小；那时也绝不会知道任何魔鬼：我所说的不仅是撒殚，还有任何其他这样的邪灵，不，即使有一百个或一万个，在爱德存在时，它们都没有力量。因为草能抵挡火，比起魔鬼能抵挡爱德的火焰更容易。爱德比任何墙壁更强固，比任何金刚石更坚牢；或者如果你能说出比这更坚固的材料，爱德的坚固也超越它们。爱德，财富和贫困都不能胜过：不，如果有爱德，就不会有贫穷，也没有无度的财富，而是只有每种境况的好处。因为从前者我们得享丰裕，从后者免于挂虑；我们既不必经历财富的焦虑，也不必经历贫穷的恐惧。\par

12. 为何我要提及由此产生的益处呢？毋宁说，请思考\OriginalTerm{爱德}{love}本身作为祝福是何等伟大：它带来何等喜乐，使灵魂建立在何等大的恩宠中；这正是爱德的卓越特质。因为其他德行各有其伴随的烦恼，如斋戒、节制、警醒，都带有嫉妒、贪欲和轻视。但爱德，除了获益，也拥有极大的快乐，毫无烦恼；它像一只辛勤的蜜蜂，从每朵花中采集甘甜，存入爱者的灵魂。即使有人是奴隶，爱德也使奴隶比自由更甜美。因为爱者并非以命令为乐，而是以服从为乐，尽管命令是甜蜜的；但爱德改变了事物的本性，将一切祝福捧在手中，比任何母亲更温柔，比任何王后更富有，使困难变得轻省容易，使我们的德行变得轻松，而使恶习对我们极其苦涩。例如：花钱似乎令人痛苦，但爱德使之愉快；接受他人之物似乎愉快，但爱德不让它显得愉快，反而塑造我们的心智去避免它如同邪恶。同样，说人坏话似乎对所有人都是愉快的，但爱德使之苦涩，而使人觉得说好话是愉快的；因为没有什么比赞美我们所爱的人更甜蜜。再者，愤怒有一种快感；但在此却不再有，反而它的所有力量都被夺去。即使被爱的人使爱者忧伤，愤怒也绝不显现；只有眼泪、劝告和恳求；爱德如此远离激怒。若她看见有人犯错，她便哀伤痛苦；然而这痛苦本身也带来快乐。因为爱德的眼泪和悲伤比任何欢笑和喜乐更甜蜜。例如：那些笑的人并不像为朋友哭泣的人那样得到更新。若你怀疑这一点，请止住他们的眼泪；他们会像遭受无法忍受的虐待一样抱怨。\Quote{但有一种，}有人说道，\Quote{不当的快乐在爱德中。}走开，住口吧，无论你是谁。因为没有什么像真诚的爱德那样远离这种快乐。\par

请勿对我提论这种寻常的、粗俗和低级的爱，那只是一种病态而非爱德；而是提论保禄所追求的那种爱德，它顾及所受之人的益处；你将会看到，没有哪位父亲比具有这种气质的人更充满慈爱。正如那些爱钱的人不能忍受花钱，宁愿受穷也不愿看到财富减少；同样，心怀善意的人宁愿遭受万般痛苦，也不愿看到他所爱的人受到伤害。\par

13. \Quote{那么，}有人说，\Quote{那爱若瑟的埃及妇人想要害她吗？}因为她以这种魔鬼般的爱去爱。但若瑟却不是这样，而是以保禄所要求的那种爱。那么，请想一想，他的话以及她所说的事迹显示出了多么大的爱。\BibleQuote{凌辱我，使我成为奸妇，亏负我的丈夫，推翻我的全家，并使自己失去对天主的信赖。}这些话表明她不但不爱他，甚至不爱自己。但因为他真心爱她，就设法使她避免这一切。为使你相信他是出于对她的关心，请从他的劝言中了解其性质。因为他不仅推开她，还引入了一段足以扑灭任何火焰的劝诫，即：\BibleQuote{若因我的缘故，我的主人}他说，\BibleQuote{不知道他家里的一切。}他立即提醒她关于她的丈夫，好让她感到羞愧。他没有说\Quote{你的丈夫}，而是说\Quote{我的主人}，这更能约束她，并促使她思考自己是谁，她迷恋的是谁——一个主母迷恋一个奴隶。\Quote{因为若他是主人，那么你就是主母。因此，当羞于与仆人亲近，想一想你是谁的妻子，你想要与谁结合，你对谁忘恩负义、不加考虑，而我要报答他更大的善意。}请看他是如何赞扬他的恩惠。因为那个野蛮而又堕落的妇人无法有高尚的情操，他用人性的考虑使她羞愧，说：\BibleQuote{他因我什么都不知道，}即\BibleQuote{他是我的大恩人，我不能伤害我的保护者的要害。他使我成为他家里的第二个主人，没有人被拦阻不归我管，除了你。}他在这里努力提升她的思想，好使他无论如何能说服她感到羞愧，并表明她所受的荣耀是多么伟大。他不仅止于此，还加上一个足以约束她的称呼，说：\BibleQuote{因为你是他的妻子；我怎能做这恶事？但你说什么？你的丈夫不在场，也不知道他受了亏负？但天主必会看见。}然而她并未因他的劝告而受益，反而仍设法吸引他。因为她想要满足自己的疯狂，并非出于对若瑟的爱，才做了这些事；从她后来的行为就可看出这一点。如她发起审判，提出指控，作假见证，将无辜的人抛给野兽，甚至将他投入监牢；或者更确切地说，她甚至杀了他，因为她如此煽动法官反对他。那么，若瑟也像她一样吗？绝不，完全相反，因为他既没有反驳也没有指控那妇人。\Quote{是的，}也许有人说，\Quote{因为即使他说了也不会被相信。}然而他深受爱戴；这不仅从开始，也从结局可以看出来。因为如果他的野蛮主人不是非常爱他，在他沉默不作任何辩护时，甚至可能杀了他；因为他是埃及人，是统治者，又以为自己在婚姻上受了亏负，而且是被一个仆人，一个曾受他如此大恩的仆人。但这一切都因他对他的尊重以及天主倾注在他身上的恩宠而让步了。除了这恩宠和爱，他还有其他不小的证据可以自证清白，那就是衣服本身。因为如果是她受了强暴，那她的衣服应该被撕破，脸被抓破，而不是她拿着他的衣服。但她说：\BibleQuote{他听见我大声呼喊，就丢下衣服逃出去了。}那么你为什么拿走他的衣服？因为一个受强暴的人，唯一想要的是摆脱闯入者。\par

但不仅从此处，从后来的事件中，我也能指出他的善意和他的爱。是的，即使当他不得不提及他入狱的原因以及他在那里停留的原因时，他也没有将整个故事和盘托出。他说了什么？\BibleQuote{我也没有做什么；但我确实是从希伯来人的地方被偷来的；}他丝毫未提及那奸妇，也没有以此自夸，而一般人即使不为虚荣，也会这样做，以免让人以为他是因为邪恶的原因被关进那牢房。因为如果连正在作恶的人也绝不停止责备同样的事，尽管这样做会带来羞辱，那么他如此纯洁却未提及那妇人的情欲，也未公开她的罪，这难道不值得钦佩吗？当他登上王位，成为全埃及的统治者时，也未曾记恨那妇人所行的恶，也未施加任何惩罚？\par

你看见他是多么关心她吗？但她的不是爱，而是疯狂。因为她爱的不是若瑟，而是想要满足自己的情欲。如果仔细审视她的话，也带有愤怒和极大的嗜血性。因为她说什么？\BibleQuote{你带了一个希伯来仆人来戏弄我们。}她责备丈夫的仁慈，并展示衣服，变得比任何野兽更凶猛；但他却不是这样。我何须提他对她的善意呢？我们知道他对那些要杀他的兄弟们也是如此，无论是在家还是在外面，从未说过他们一句严厉的话。\par

14. 因此，保禄说，我们所说的这种爱德，是万善之母，并将它置于奇迹和一切其他恩赐之上。因为正如那里有金制的衣袍和凉鞋，我们还需要其他衣物来辨别君王；但如果我们看见紫袍和王冠，就不需要再看其他王权的标记：同样在这里，当爱德的冠冕戴在我们头上时，就足以指出基督的真正门徒，不仅向我们自己，也向不信的人。因为主说：\BibleQuote{如果你们彼此相爱，众人因此就可认出你们是我的门徒。}\BibleRef{若 13:35} 所以这个标记当然大于一切标记，因为门徒是由此被认出的。因为即使有人行千万个奇迹，却彼此纷争，他们必为不信者所耻笑。正如他们若不行奇迹，却彼此切实相爱，他们就会受到所有人的尊敬和不可侵犯。因为保禄本人也是因此受到我们钦佩的，并非因为他使死人复活，或洁净了癞病人，而是因为他说：\BibleQuote{谁软弱，我不软弱呢？谁跌倒，我不焦急呢？}\BibleRef{格后 11:29} 因为即使你有千万个奇迹与此相比，你也不会说出任何与此相等的话。因为保禄本人也说，为他预备了极大的赏报，不是因为他行奇迹，而是因为他\Quote{对软弱的人，我就成为软弱的}。他说：\BibleQuote{我的报酬是什么呢？就是传布福音时，使福音不费钱而传布。}\BibleRef{格前 9:18} 当他把自己置于宗徒之前时，他没有说：\Quote{我比他们行了更多的奇迹}，而是说：\BibleQuote{我比他们更劳苦}\BibleRef{格前 15:10} 他甚至愿意为救门徒而死于饥荒。他说：\BibleQuote{我宁愿死，也不愿有人使我这夸耀落了空。}\BibleRef{格前 9:15} 这并不是因为他夸耀，而是免得他似乎责备他们。因为他从不习惯在自己的成就上夸耀，除非时机需要；即使被迫如此，他也称自己是\Quote{愚妄人}。但如果他有时夸耀，那是\Quote{在软弱中}，在受屈中，在深切同情受伤害的人中；正如这里他也说：\BibleQuote{谁软弱，我不软弱呢？} 这些话甚至比危险更伟大。所以他也将它们放在最后，以加强他的论述。\par

我们与祂相比，应当配得什么？我们既不为自己轻看财富，也不放弃我们多余的财物？但保禄却不如此；他反而牺牲自己的灵魂和身体，好让那些用石头砸他和用棍子打他的人，得以获得天国。他说：\Quote{因为基督这样教导我如何去爱。}祂留下了关于爱的新诫命，祂自己也以行动实践了。因为祂是万有的主，具有那有福的本性；对于祂从无中创造、并赐予无数恩惠的人，这些人侮辱祂、向祂吐唾沫，祂却没有转身离开，反而为了他们而降生成人，与娼妓和税吏来往，医治附魔的人，并许下天堂。在这一切之后，他们捉拿祂，用棍子打祂，捆绑祂，鞭打祂，嘲笑祂，最后把祂钉在十字架上。即便如此，祂也没有转身离开，反而当祂高悬在十字架上时，祂说：\BibleQuote{父啊，赦免他们的罪。}但那个先前辱骂祂的强盗，祂却领他进入了乐园；并使迫害者保禄成为宗徒；又为了那些钉死祂的犹太人，把祂自己的门徒，就是那些与祂亲密和完全忠于祂的人，交于死亡。\par

因此，我们在心中追忆这一切，包括天主的事和人的事，让我们效法这些崇高的行为，并拥有那超越一切恩赐的爱德，使我们获得现世和来世的福分。愿我们众人赖我们的主耶稣基督的恩宠和仁慈，获得这一切。愿光荣、权能、尊崇归于祂，同圣父及圣神，从今时直到永远，世世无穷。阿们。\par

\SourceRecord{Translated by Talbot W. Chambers.；Nicene and Post-Nicene Fathers, First Series；Vol. 12.；Edited by Philip Schaff.；Buffalo, NY: Christian Literature Publishing Co.,；1889.；<http://www.newadvent.org/fathers/220132.htm>.}

% structure: p0001 source=h1 target=section reason=page-title

\section{论《格林多前书》第三十三篇讲道}

% structure: p0002 source=h2 target=subsection reason=relative-heading-rank; base=h2

\subsection{格林多前书 13:4-5}

\BibleQuote{爱是含忍的，爱是慈祥的；爱不嫉妒，不夸张，不自大}\par

这样，既然他已表明，信德、知识、先知之恩、语言、恩赐、治愈、完美的生活和殉道，若没有爱，就毫无益处；他接着就必须勾勒出爱的无比美丽，用种种美德的色彩装饰它的形象，并精确地组合其所有部分。但是，亲爱的，你不要匆匆读过这些话，而要仔细考察每一个，好让你知道这事的宝藏和画家的技艺。例如，思考他一开始从何处着手，他首先放置了什么作为一切美善的原因。这又是什么呢？是含忍。这是一切克己的根源。因此，一位智者说：\Quote{含忍的人有真智慧，性急的人只显露愚昧。}\par

并且，他将她比作一座坚城，说它比那座城更坚固。因为她既是无敌的武器，又是一种无法攻克的堡垒，轻易击退一切烦恼。如同火星落入深渊，不能伤害深渊，却很快熄灭：同样，无论什么意外的事落到含忍的灵魂上，那事确实迅速消失，但灵魂却不被搅扰：因为，诚然，没有什么是比含忍更不可穿透的。你可以谈论军队、金钱、马匹、城墙、兵器，或其他任何事物；你找不到任何比含忍更强大的。因为那些被这些包围的人，常常被愤怒击败，像一个无用的孩子一样心烦意乱，使一切充满混乱和风暴；而这人，仿佛安息在港湾里，享受着深沉的安宁。即使你用损失包围他，你也未能移动那磐石；即使你侮辱他，你也未能动摇那高塔；即使你用鞭打伤害他，你也未能划伤那金刚石。\par

是的，因此他被称为含忍，因为他有一种宽广而伟大的灵魂。因为长也被称为大。但这种卓越生于爱，对拥有它的人和享受它的人都有不小的益处。请你不要向我提那些被遗弃的可怜人，他们作恶而不受惩罚，变得更坏；因为这里的坏结果并非来自他的含忍，而是来自那些滥用它的人。所以不要提这些人，而要提那些温良的人，他们从中获得巨大的益处。因为当他们做了坏事而没有受罚时，仰慕受苦者的温良，他们因此收获了非常伟大的自制教训。\par

但是保禄并没有停在这里，而是继续说出爱的其他崇高成就，说：\BibleQuote{是慈祥的。}因为既然有些人实践含忍，不是为了自己的克己，而是为了惩罚那些激怒他们的人，使他们怒火爆发；他说，爱也没有这个缺陷。因此他又加了：\BibleQuote{是慈祥的。}因为他们这样温和地对待他们，绝不是为了点燃那些被愤怒烧灼的人的火，而是为了平息和熄灭它：并且不仅通过高贵地忍受，也通过安抚和安慰，来医治伤痛，愈合情欲的创伤。\par

\BibleQuote{不嫉妒。}因为一个人可能同时含忍又嫉妒，因而那卓越性被破坏。但爱也避免了这一点。\par

\BibleQuote{不夸张}意思是，不鲁莽。因为它使爱的人既谨慎、庄重又坚定。事实上，那些非法爱慕的人的一个标志就是在这方面的缺陷。而那个认识这种爱的人，是所有的人中最完全摆脱这些邪恶的。因为内心没有愤怒时，鲁莽和傲慢都被彻底消除了。爱，像一位优秀的农夫，在内里居于灵魂中，不许任何这些荆棘生长出来。\par

\BibleQuote{不自大。}因为我们看到许多人在这些卓越的方面自视甚高；例如，在不嫉妒、不吝啬、不小气、不鲁莽上：这些邪恶不仅附属于财富和贫穷，甚至附属于本性善良的事物。但爱彻底清除一切。并且思考：含忍的人当然不一定是慈祥的。但如果他不慈祥，那事就变成一种恶，他就有陷入恶意的危险。因此她提供了一剂药，我说是慈祥，并保持这德行的纯洁。再者，慈祥的人常常过于随和；但她也纠正这个。因为他说：\BibleQuote{爱，}他说，\BibleQuote{不夸张，不自大：}慈祥而含忍的人常常喜欢炫耀；但她也除去这个恶。\par

并且看他如何不仅从她所拥有的，也从她所没有的来装饰她。因为他说她既带来德行，又灭绝邪恶，甚至不让它萌发。因此他没有说：\Quote{她确实嫉妒，但克服了嫉妒}；也没有说：\Quote{她傲慢，但惩戒了那情欲}；而是说：\BibleQuote{不嫉妒，不夸张，不自大}；这实在是最令人钦佩的，就是她甚至不需辛劳就成就了她的善事，无需战争和阵势就立起了她的凯旋：她不容许拥有她的人劳苦而后获得冠冕，而是不费力地将奖赏带给他。因为若没有情欲与清醒理智争战，哪里有什么劳苦呢？\par

2. \BibleQuote{不作无礼的事。} \Quote{哪里，}他说，\Quote{我为什么要说\InnerQuote{不自大}，当她远非有那种感觉，以致为了她所爱的人遭受最可耻的事，甚至不认为那事是一种无礼？}再者，他没有说：\Quote{她受无礼却高尚地忍受羞耻}，而是说：\BibleQuote{她甚至根本没有任何羞耻感。}因为如果爱钱之人为他们那污秽的买卖忍受各种辱骂，且不仅不掩面，反倒以此为荣：那么拥有这种可称赞之爱的人，为了他所爱的人的安全，就更不会拒绝任何事了：而且，他所受的任何事都不能使他羞耻。\par

为了我们不至于从任何卑贱的事物中取材，让我们审视这句话应用在\Person{基督}身上的情形，然后我们就会看到所说之话的力量。因为我们的主\Person{耶稣基督}曾被可怜的奴隶吐唾沫和用棍子打；祂不仅不视此为不体面，反而欢欣并称之为光荣；并且将一个强盗和杀人犯与自己一同领进乐园，与一个妓女交谈，在旁人都指责祂时，祂不认为这事可耻，反而允许她亲吻祂的脚，用眼泪浸湿祂的身体，用她的头发擦干，这一切都在一群既是敌人又是仇敌的观众面前；因为\Quote{爱不作不当的事。}\par

因此，父亲们，即使他们是哲学家和演说家中的翘楚，也不耻于与孩子们咿呀学语；目睹此景的人没有一个责备他们，反而此事被视为如此美好和正当，以至于甚至值得祈求。再者，若孩子们变得邪恶，父母继续纠正、照顾他们，减轻他们招惹的责骂，并不感到羞耻。因为爱\Quote{不作不当的事，}反而好像用金色的翅膀遮盖了所爱者的一切过犯。\par

同样，\Person{约纳堂}爱\Person{达味}；当他听见他的父亲说：\BibleRef{撒上 20:30} \BibleQuote{你这逃离家门的少女之子，你这女人养大的，}他并不羞耻，尽管这话充满了极大的侮辱。因为他的意思是这样的：\Quote{你这下贱妓女之子，她们迷恋男人，追逐路人，你这软弱无能的废物，毫无男子气概，活着只为自己和生母的耻辱。}那么怎样呢？他因这些而悲伤吗？他掩面不看他的挚爱吗？绝不，恰恰相反；他展示自己的深情作为装饰。然而，一个当时是国王，而且是国王的儿子，即\Person{约纳堂}；另一个是流亡者和流浪者，我指的是\Person{达味}。但即使如此，他也没有为他的友谊感到羞耻。因为爱不作不当的事。是的，这就是它奇妙的品质：它不仅不让受伤害者悲伤和感到刺痛，反而使他乐于接受。因此，我们所说的这位，在这一切之后，好像有人给他戴上了一顶冠冕，走去扑在\Person{达味}的颈项上。因为爱不知道羞耻为何物。因此，它在别人掩面的事上夸耀。因为羞耻在于不知道如何去爱；不在于当你爱时，为所爱者遭遇危险并忍受一切。\par

但我说\Quote{一切，}时，不要以为我指的也包括有害之事；例如，帮助一个年轻人谈情说爱，或任何人恳求他人为他做任何伤害性的事。因为这样的人并不爱，而且我最近从埃及妇人身上向你们证明了这一点：因为实际上，只有寻求所爱者益处的人才是爱人者；所以如果有人不追求这个，即使追求正当和良善的事，尽管他做出千万次爱的表白，他也比任何敌人都更充满敌意。\par

同样，\Person{黎贝加}从前，因为她极其依恋她的儿子，就实施了一次偷窃，并且不为被察觉而羞耻，也不害怕，尽管风险非同寻常；但即使当她的儿子向她提出顾虑时，她说：\Quote{愿你的诅咒落在我身上，我儿。}你甚至在一个女人身上看到宗徒的心志吗？正如\Person{保禄}选择，若可将小事与大事相比，为犹太人成为诅咒，\BibleRef{罗 9:3}同样，她也为使她的儿子蒙受祝福，选择了无异于受诅咒。她把美好事物让给了他，因为看来她不会与他一同蒙受祝福，但坏事她准备独自承受；然而，她欢喜且急忙，尽管有如此大的危险摆在她面前，她为事情拖延而忧伤：因为她害怕\Person{厄撒乌}可能抢先他们，使她的智慧落空。因此，她也缩短谈话，催促年轻人，仅允许他对所说的话回答，便陈述了一个足以说服他的理由。因为她没有说：\Quote{你说这些毫无道理，你白担心了，你的父亲年迈，视力也不好了；}而是说什么？\Quote{愿你的诅咒落在我身上，我儿。只是你不要破坏计划，也不要失去我们追逐的目标，也不要放弃宝藏。}\par

正是这个\Person{雅各伯}，难道他不是为薪俸在亲属那里服侍了两次七年吗？难道他不是连同奴役一起因那个诡计而遭受嘲笑吗？那么怎样呢？他感受到嘲笑了吗？他认为这是行为不当吗？身为自由人、生来自由、教养良好，却在亲属中忍受奴隶待遇——当一个人受到朋友的侮辱时，这通常是极其恼人的？绝不。原因是他的爱，使时间虽长，却显得短暂。他说：\BibleRef{创 29:20}\BibleQuote{在他眼中好像有几天。}他远未因此奴役而恼怒和脸红。因此，有福的\Person{保禄}说得好：\Quote{爱不作不当的事。}\par

3. 1. \Quote{不求自己的益处，不轻易发怒。}\par

这样说了\BibleQuote{不作无礼的事}之后，他也指出了那使她不作无礼事的心境。那心境是什么？就是她\BibleQuote{不求己益}。因为她把所爱者视为一切，而只有当她不能使他脱离这种无礼时，她才\BibleQuote{作无礼的事}；所以，如果可能通过她自己的无礼来使所爱者受益，她甚至不把那看作无礼；因为对方从此就是你本人，当你爱的时候：因为这就是友谊，即爱者与被爱者不再分开为两个人，而是成为一个单一的人；这事除非因着爱德，否则决不会发生。因此，不要寻求你自己的，好叫你能找到你自己的。因为寻求自己所有的，反找不到自己所有的。因此保禄也说：\BibleQuote{谁也不要求自己的益处，但要求别人的益处。}\BibleRef{格前 10:24}因为你的利益在于你邻人的利益，他的利益也在于你的利益。所以，如同一个人把金子埋在他邻人的家里，若他拒绝去那里寻找挖掘，就永远找不到；同样，如果有人不愿在邻人的益处中寻求自己的利益，他就不会获得为此应得的冠冕：天主自己就这样安排了，为使我们彼此相连；又如同一个人唤醒一个睡着的孩子去跟他的兄弟，当孩子自己不愿意时，就把孩子渴望和期盼的东西放在兄弟手中，好使孩子因渴望得到它而追逐拿着它的人，事情就这样发生了：同样，天主将各人的利益交给了他的邻人，好使我们互相追随，而不致分裂。\par

如果你愿意，也请看在我们这些向你们讲话的人身上。因为我的利益依赖于你们，你们的益处也依赖于我。因此，一方面，你们受教明白中悦天主的事对你们有益，但这事已交托给我，使你们从我这里领受，因而你们被迫奔向我；另一方面，你们变得更好对我有益，因为我为此将得的赏报是大的；但这赏报又在你们身上；所以我也被迫追随你们，使你们更好，好叫我从你们那里得到我的利益。因此保禄也说：\BibleQuote{我的希望是什么呢？不就是你们吗？}又说：\BibleQuote{我的希望、我的喜乐和我夸耀的冠冕。}\BibleRef{得前 2:19}所以保禄的喜乐就是门徒们，他的喜乐他们也有。因此他看见他们丧亡时甚至哭泣。\par

同样，他们的利益也依赖于保禄：因此他说：\BibleQuote{我为以色列的希望才被这条锁链束缚。}\BibleRef{宗 28:20}又说：\BibleQuote{我为选民忍受一切，使他们获得永生。}\BibleRef{弟后 2:10}这一点在世俗之事中也可以看到。他说：\BibleQuote{因为妻子对自己的身体没有主权，而是丈夫有；同样，丈夫对自己的身体也没有主权，而是妻子有。}\BibleRef{格前 7:4}同样，当我们想将任何人捆绑在一起时，我们也这样做。我们不把任何一方留在自己的权下，而是在他们之间伸展一条链条，使一方被另一方所持守，另一方也被这一方所持守。你也在治理者身上看到这一点吗？审判者坐堂审判不是为了自己，而是寻求邻人的利益。被治理者则通过他们的追随、他们的服务和其他一切事物来寻求治理者的利益。士兵为我们拿起武器，因为他们为我们冒险。我们为他们而困苦，因为他们的供给来自我们。\par

但如果你说\Quote{各人这样做是寻求自己的利益}，这我也说，但我补充说，通过别人的益处自己的才被获得。因此，士兵若不为供养他的人作战，就没有人为此目的服侍他；同样，供养者若不供给士兵，就没有人替他武装自己。\par

4. 你看到爱德如何处处延伸并管理一切吗？但不要厌倦，直到你彻底熟悉了这条金链。因为说了\BibleQuote{不求己益}之后，他又提到由此产生的美好事物。这些是什么呢？\par

\BibleQuote{不动怒，不图谋恶事。}请看爱德不仅制伏恶习，甚至不让它兴起。因为他说了，不是\Quote{即使被激怒，她也胜过了}，而是\BibleQuote{不动怒}。他又说了，不是\Quote{不作恶事}，而是\Quote{连想都不想}；也就是说，她不仅不图谋任何恶事，甚至对所爱者也不起疑心。那么她如何能作恶事，或如何被激怒呢？她甚至不容许存有恶意的猜疑；愤怒的泉源又在哪里呢？\par

% structure: p0026 source=h2 target=subsection reason=relative-heading-rank; base=h2

\subsection{格前 13:6-7}

\BibleQuote{不以不义为乐}：即不以受苦的人为乐；不仅于此，还有更重要的，\BibleQuote{与真理同乐}。他说：\Quote{她与那些受称赞的人同乐}，正如保禄说：\BibleQuote{应与喜乐的一同喜乐，与哭泣的一同哭泣。}\BibleRef{罗 12:15}\par

因此，她\BibleQuote{不嫉妒}，因此她\BibleQuote{不自大}：因为她实际上把别人的好事视为自己的。\par

你看到爱德如何逐渐使她所培育的人成为天使吗？因为当他没有了愤怒，纯洁于嫉妒，并免于一切暴虐的情欲时，试想，他从此便脱离了人的本性，而达到了天使的那份宁静。\par

然而，他并不满足于此，还要说出更胜过这些的话：这是按照他先弱后强的论述计划。因此他说，\Quote{含忍一切}。由于她的宽忍，由于她的良善；无论是重担、痛苦、凌辱、鞭打、死亡，或其他任何事。这一点再次可以从真福\Person{达味}的事例中看出。还有什么比看见儿子起来反叛他、图谋篡位、渴求父亲的血更难以忍受的呢？然而这位真福者忍受了这一切，甚至不能对那弑亲者说出一句恶言；反而当他把其余事项都交给将领们时，却严厉地嘱咐要保护他的安全。因为他的爱的基础是坚固的。所以爱也\Quote{含忍一切}。\par

宗徒在此暗示了它的能力，却由下文说明它的良善。因为他说：\Quote{凡事盼望}，他说，\Quote{凡事相信，凡事忍耐}。什么是\Quote{凡事盼望？}它不绝望，他说，\Quote{对被爱的人，即使他一文不值，仍继续纠正、供给、照顾他}。\par

\Quote{凡事相信}。\Quote{因为它不仅盼望，}他说，\Quote{而且因着极大的爱也相信。}即使这些好事没有按照它的盼望实现，反而对方显得更加难以忍受，它也能承受这些。因为他说，它\Quote{凡事忍耐}。\par

% structure: p0033 source=h2 target=subsection reason=relative-heading-rank; base=h2

\subsection{格林多前书 13:8}

5. \Quote{爱永不止息。}\par

你看见他何时把桂冠放在拱顶，以及在这恩赐中什么是一切事物所特有的？因为什么是\Quote{永不止息}？它不被割断，不被忍耐溶解。因为它忍受一切：既然无论发生什么，那爱者永远不会憎恨。这就是它的最大卓越之处。\par

\Person{保禄}就是这样的人。因此他也曾说：\Quote{如果可能，我还能激动我的骨肉之亲发愤；}\BibleRef{罗 11:14} 并且他继续怀着希望。他还给\Person{弟茂德}训令说：\Quote{主的仆人不应当争斗，而应当温良对待众人……以柔和开导那些反对的人，或许天主赐给他们认识真理的心。}\BibleRef{弟后 2:24}\par

有人说：\Quote{那么，如果他们是敌人和外邦人，就必须憎恨他们吗？}必须憎恨，不是他们，而是他们的道理：不是人，而是邪恶的行为、败坏的思想。因为人是天主的作为，而诡诈是魔鬼的作为。所以你不要混淆天主的事物和魔鬼的事物。既然犹太人既是亵渎者、迫害者、伤害者，又说了千万句反对基督的恶言，那么\Person{保禄}，他比任何人都爱基督，会憎恨他们吗？绝不，他反而爱他们，为他们做一切事：一时他说：\Quote{我心里所渴望的和向天主祈求的，就是为使他们得救；}\BibleRef{罗 10:1} 另一时又说：\Quote{我宁愿自己为他们的缘故，被诅咒与基督隔绝。}同样，\Person{厄则克耳}看见他们被杀时说：\Quote{哎，吾主上主，你要灭绝以色列的遗民吗？}\BibleRef{则 9:8} 而\Person{梅瑟}说：\Quote{如果你赦免他们的罪，就赦免吧。}\BibleRef{出 32:32}\par

那么为什么\Person{达味}说：\Quote{上主，我不是憎恨那憎恨你的人吗？不是因你的敌人而憔悴吗？我憎恨他们，以完全的憎恨。}\BibleRef{咏 139:21}\par

首先，\Person{达味}在圣咏集中所说的一切，并不都是他本人的话。因为他自己说：\Quote{我曾在刻达尔的帐幕中寄居；}\BibleRef{咏 119:5} 以及\Quote{当我们坐在巴比伦河畔时，我们哭泣。}\BibleRef{咏 136:1} 然而他既没有见过巴比伦，也没有见过刻达尔的帐幕。\par

但除此之外，我们现在需要更完全的自制。因此，当门徒们求火降下，就像在\Person{厄里亚}的事例中一样时，基督说：\Quote{你们不知道你们是何种心神。}\BibleRef{路 9:55} 因为在那个时代，他们不仅被命令憎恨不虔敬，也要憎恨不虔敬者本人，以免他们的友谊成为他们犯罪的因由。因此他斩断了他们的联系，无论是血缘还是婚姻，并从各方面把他们隔离开。\par

但现在，因为他把我们带到更完全的自制，把我们高置在那祸害之上，他命令我们反而接纳并安抚他们。因为我们不会受到他们的损害，他们反而因我们受益。那么他说什么？我们不可憎恨，而要怜悯。因为如果你憎恨，怎能轻易转化那处于错误中的人？你怎能替不信者祈祷？因为人应当祈祷，请听\Person{保禄}所说：\Quote{所以我首先劝勉众人，要为所有人祈求、祈祷、转求和感谢。}\BibleRef{弟前 2:1} 但那时并非所有人都是信者，我想这是众所皆知的。他又说：\Quote{为君王和所有位高者。}但这些人是不虔敬和违犯者，这也是显而易见的。此外，他提到祈祷的理由时补充说：\Quote{因为这是美好的，在我们的救主天主面前是蒙悦纳的；他愿意所有人得救，并认识真理。}因此，如果他发现一位外邦妻子与信者同居，他并不解除婚姻。还有什么是比男人与妻子更紧密的结合呢？\Quote{二人成为一体，}\BibleRef{创 2:24} 在那实例中有着极大的魅力，和炽热的欲望。但如果我们必须憎恨不虔敬和无法纪的人，我们也会继续憎恨罪人；这样循序渐进，你会与你的多数弟兄隔断，甚至与所有人隔断：因为没有一个人，没有一个人，是没有罪的。因为如果我们的责任是憎恨天主的敌人，那么我们就不仅必须憎恨不虔敬者，也要憎恨罪人：这样我们就比野兽更糟，躲避所有人，并骄傲自满；如同那位法利塞人。但\Person{保禄}没有这样命令我们，而是如何？\Quote{劝戒那散漫的，鼓励那怯懦的，扶持那软弱的，对众人要含忍。}\BibleRef{得前 5:14}\par

6. 那么，当他说：\BibleQuote{若有人不听从我们这封信的话，就记下那人，不要与他交往} \BibleRef{得后 3:14}时，是什么意思呢？首先，他说这话是针对弟兄，但即便如此也不是毫无限制，而是同样带着温柔。因为你不可截断下文，而要加上下一句：他既说了\BibleQuote{不要与他交往}，又加上：\BibleQuote{但不要以他为仇人，而要劝戒他如弟兄} \BibleRef{得后 3:15}。你看见他如何吩咐我们恨恶那邪恶的行为，而不恨恶那人吗？因为将我们彼此分离的实在是魔鬼的作为，它总是竭力消除爱德，以切断纠正之路，使他留在错误中，使你留在敌意中，从而堵塞他得救的道路。因为当医生既憎恶病人又避开他，病人也避开医生时，那患病的人何时才能得痊愈呢？既然一方不会寻求另一方的帮助，另一方也不会去找他？\par

但是，请告诉我，你为什么完全避开他、躲避他？因为他不虔敬？正因如此，你更应欢迎他、照管他，好使你在他患病时扶起他。但即使他病入膏肓，你仍被命令尽你的本分。因为犹达斯也患了不治之症，但天主并没有停止照管他。所以，你也不可厌倦。因为即使你费尽辛劳仍不能将他从不虔敬中解救出来，你仍将得到解救者的赏报，并使他对你的温柔感到惊奇，于是这一切赞美将归于天主。因为即使你行奇迹、复活死人、做任何工作，外教人也不会像看见你展现温和、良善、柔顺的性情时那样惊叹你。这并非小事：因为许多人甚至会完全脱离他们的恶道；没有什么比爱德更能吸引人。因为对于前者，他们反而会嫉妒你，我指的是征兆和奇迹；但对于这爱德，他们既钦佩你又爱慕你；如果他们爱慕，他们也将在适当的时候把握真理。但如果他没有立刻成为信徒，不要惊奇，不要催促，也不要要求一切同时实现，而要容许他目前先赞美、爱慕，然后他将在适当的时候达到这一点。\par

7. 为使你清楚知道这是何等大事，请听连保禄在一位不信的法官面前为自己辩护时所说的话。他说：\BibleQuote{我自己以为是有福的，}他说，\BibleQuote{因为今天要在你面前辩护。} \BibleRef{宗 26:2} 他说这些话，并非要奉承他，绝非如此；而是希望以他的温和赢得他。他确实部分地赢得了那人，那直到当时被视为被定罪的人，竟俘虏了他的法官，而这场胜利由那被俘虏的人自己在众人面前大声承认说：\BibleQuote{你差一点就劝服我作基督徒了。} \BibleRef{宗 26:28} 那么保禄说了什么？他张开了更大的网，说：\Quote{我自己并不避讳它们，也不以它们为耻，但我迁就他们的软弱。因为他们还不能接受我的夸耀；我已从我的主学到：不可将\BibleQuote{一块未漂过的布补在旧衣服上} \BibleRef{玛 9:16}，因此我才这样说。事实上，直到此时他们还听人诋毁我们的教义，厌恶十字架。因此，如果我再加上锁链，他们的憎恨会更深；所以我除去了这些，好使另一条得以被接受。正是如此，在他们看来被捆绑是可耻的，因为他们尚未尝到与我们同在的荣耀。所以必须俯就；当他们学会了真正的生命，那时他们也会知道这铁的美丽，以及这锁链带来的光辉。} 此外，在与其他人谈话时，他甚至称这为恩赐，说：\BibleQuote{因为基督赐予你们的恩宠，不仅是信祂，而且是为祂受苦。} \BibleRef{斐 1:29} 但在当时，对听众来说，不因十字架为耻已是大事；因此他循序渐进。同样，没有人带一个人进入宫殿时，在他还没有看见门廊之前，就强迫还在外面的人细看里面的景象：因为那样的话，它甚至不会显得令人赞叹，除非人进去并亲自认识一切。\par

那么，我们也该这样对待外教人：以俯就，以爱德。因为爱德是一位伟大的导师，它能使人从错误中脱离，能改造品格，能扶着手引导人走向克己，并能从石头中造出人来。\par

8. 你若愿意认识她的力量，请带一个胆小怕事、闻声则惊、见影而颤的人，或一个易怒粗暴、如同野兽而非人的，或一个放荡不羁、完全陷于邪恶的人，把他交在爱的手中，领他进入这所学校；你必会立刻看到那怯懦胆怯之辈变得勇敢豪迈，欣然冒险。而奇妙的是，这些变化并非出于本性的改变，而是在那怯懦的灵魂中，爱彰显了她独特的力量；这正如一个人使一把铅剑，虽未变成钢铁，却仍保持铅的本性，而能发挥钢铁的作用。例如：雅各伯是一个\Quote{恬静的人}，\BibleRef{创 25:27} 居住在屋内，不惯于劳苦和危险，过着一种闲散安逸的生活，如同闺中的少女；他也因此多半被迫坐在家中，看守房屋；远离市场和市场的喧嚣，远离诸如此类的事，甚至安闲宁静。结果如何？爱的火炬点燃他之后，请看它如何使这恬静居家的人变得坚忍耐劳、热爱劳作。这并非我所说，而是这位族长自己的话：他对他的亲戚抱怨说：\Quote{这二十年中，我与你同在。} \BibleRef{创 31:36} 这二十年你是如何度过的？（他也加上这）\Quote{白日受尽炎热，黑夜受尽寒霜，睡眠也离我而去。} 这就是那\Quote{恬静的人，呆在家中}，过着安逸生活的人所说的话。\par

再者，他胆怯是明显的，因为他预期要见厄撒乌，就吓得半死。但再看，这胆怯的人在爱的影响下变得比狮子更勇敢。因为他像冠军一样挺身而出，准备第一个迎接那他所推测的野蛮嗜血的兄弟，并用自己的身体换取他妻子们的安全；他害怕且战栗的那一位，他却渴望自己成为阵中最先看见他的人。因为这种恐惧并不及他对妻子们的爱强烈。你看见了吗？他是如何突然变得冒险，不是改变了他的性格，而是被爱所充满？因为此后他仍然胆怯，这从他不断改变地方可以明显看出。\par

但不要让任何人认为以上所说的有损于那位义人：因为胆怯并非耻辱，这是人的本性；但因胆怯而做出不当之事才是耻辱。因为本性胆怯的人可以藉着虔敬变得勇敢。梅瑟如何？他岂不因惧怕一个埃及人而逃跑，远走他乡吗？然而，这个逃犯，一个连单个人的威胁都受不了，在尝到爱的甘蜜之后，却高贵地、无人强迫，甘愿与他所爱的人一起灭亡。\Quote{现在求你赦免他们的罪，} 他说，请赦免；若不然，\Quote{就求你从你所写的册上涂抹我的名字。} \BibleRef{出 32:32}\par

9. 此外，爱也使暴躁的人温和，使放荡的人贞洁，我们无需再举例子：这对所有人都是明显的。即使一个人比任何野兽更野蛮，爱也能使他变得如绵羊般温顺。例如，有什么比撒乌耳更野蛮、更狂暴呢？但当他的女儿放走他的敌人时，他连一句恶言都没有对她发出。那为了达味毫不留情地用刀杀死所有司祭的人，看到女儿将敌人送出了家门，竟也不向她发怒，连口头上也没有；尽管如此巨大的欺骗是针对他的：因为他被爱的更强缰绳所约束。\par

如今，正如节制是爱的常有效果，贞洁也是。如果一个人以应有的方式爱自己的妻子，即使他再倾向于放荡，也不会因着对她的爱而去看别的女人。\Quote{因为爱情，} \BibleRef{歌 8:5} 有人说得对，\Quote{猛如死亡。} 因此，放荡的行为除了源于缺乏爱，别无其他来源。\par

既然爱是一切美德的设计师，让我们以万全之策将她植入我们自己的灵魂，使她为我们产生许多祝福，并使我们不断结出她的果实，那永远新鲜、永不朽坏的果实。因为这样我们将获得至少永恒的福乐：愿我们众人赖我们的主耶稣基督的恩宠和慈爱，得以获得，祂与圣父及圣神，永生永王，世世无尽。阿们。\par

\SourceRecord{Translated by Talbot W. Chambers.；Nicene and Post-Nicene Fathers, First Series；Vol. 12.；Edited by Philip Schaff.；Buffalo, NY: Christian Literature Publishing Co.,；1889.；<http://www.newadvent.org/fathers/220133.htm>.}

% structure: p0001 source=h1 target=section reason=page-title

\section{\WorkTitle{格林多前书第三十四篇讲道}}

% structure: p0002 source=h2 target=subsection reason=relative-heading-rank; base=h2

\subsection{格林多前书 13:8}

\BibleQuote{可是，先知之恩终必消逝，语言之恩终必停止，知识之恩终必消逝。}\par

讲明了爱德对神恩与生活美德都是必要的，并列举了其一切优点，指出它是严格克己的基础之后，他又从另一个角度，第三个要点，再次指出爱德的宝贵。他这样做，首先是想说服那些被视为次等的人，让他们知道，拥有爱德就能拥有最大的神恩，而且他们不会比拥有神恩的人差，反而更好；其次，对于拥有较大神恩而因此骄傲的人，他也想贬低他们，并表明除非他们有爱德，否则一无所有。这样，他们就会彼此相爱，嫉妒和骄傲也因此被消除；而彼此相爱，又会进一步驱逐这些情感。\BibleQuote{因为爱德不嫉妒，不自夸。}所以，他从各个方面为他们建造了一道不可攻破的城墙，一种多重的合一，这合一首先除去他们的一切混乱，然后自身又变得更加坚固。因此，他也提出无数理由来安慰他们的沮丧。例如：\BibleQuote{同一的圣神，}他说，是赐予者；祂\BibleQuote{赐给每人，用以求利益；随祂的心愿分配，}祂所分赐的是恩赐，不是债务。即使你只领受了一点点，你也同样为身体做出贡献，而且你同样享受很多尊荣。那拥有较大恩赐的人也需要你，那拥有较小恩赐的人。并且，\BibleQuote{爱德是最大的恩赐，是\InnerQuote{更超卓的道路}。}\par

他说这一切，都是为了双倍地使他们彼此连结：一是因为他们拥有爱德就不会觉得自己被轻视；二是因为在追求并获得爱德之后，他们从此就不会再感受到人性的软弱；既拥有一切恩赐的根源，又即使一无所有也不再能争论。因为一旦被爱德所俘虏，就脱离了争论。\par

因此，为了指出他们将从中获得多大的益处，他描绘了爱德的果实：通过赞美爱德来抑制他们的混乱。因为他所提到的每一件事，都足以医治他们的创伤。因此他也说：\BibleQuote{含忍，}对彼此争吵的人；\BibleQuote{慈祥，}对互相疏远、心怀怨恨的人；\BibleQuote{不嫉妒，}对嫉妒比自己强的人；\BibleQuote{不自夸，}对彼此分离的人；\BibleQuote{不傲慢，}对自夸胜过别人的人；\BibleQuote{不作无礼的事，}对不认为谦卑是责任的人；\BibleQuote{不求己益，}对忽略他人的人；\BibleQuote{不动怒，不记仇，}对傲慢无礼的人；\BibleQuote{不以不义为乐，却与真理同乐，}对嫉妒的人；\BibleQuote{凡事包容，}对背信弃义的人；\BibleQuote{凡事盼望，}对灰心失望的人；\BibleQuote{凡事忍耐，永不失败，}对容易分离的人。\par

2. 在他从各方面表明爱德极其伟大之后，又从一个最重要的角度，通过新的比较来提升爱德的尊贵，说道：\BibleQuote{可是，先知之恩终必消逝，语言之恩终必停止。}因为如果这两者都是为了信仰而引入的，当信仰到处传播之后，它们的用处就多余了。但彼此相爱不会停止，反而会进一步增长，无论是在今世还是来世，那时比现在更甚。因为今世有许多削弱我们爱德的事物：财富、事务、肉体的情欲、灵魂的混乱；但在那里，这些都没有。\par

虽然先知之神和语言之恩的消逝不足为奇，但知识之恩的消逝却可能引起一些困惑。因为他也加上了：\BibleQuote{知识之恩终必消逝。}那么，我们将生活在无知中吗？绝非如此。相反，那时我们的知识很可能变得更加深厚。因此他也说：\BibleQuote{那时我就要全知道，如同我全被知道一样。}所以，请注意，为了使你不认为这消逝与先知之恩和语言之恩同等，他在说了\BibleQuote{知识之恩终必消逝}之后，并没有保持沉默，反而加上了它消逝的方式，立即补充说：\par

% structure: p0009 source=h2 target=subsection reason=relative-heading-rank; base=h2

\subsection{格林多前书 13:9-10}

\BibleQuote{因为我们只知道一部分，也只说一部分预言；但当那完全的一来到，那局部的就要消逝。}\par

所以消逝的不是知识，而是我们知识的局部分性。因为我们不仅会知道同样多，甚至更多。但为了举例说明：现在我们知道天主无所不在，但我们不知道如何；我们知道祂从无中生有，但方式我们不知道；我们知道祂由童贞女所生，但如何，我们尚不知道。但那时我们将对这些事知道得更多、更清楚。接着他指出了两者之间的距离有多大，以及我们的欠缺并不不小，说道：\par

% structure: p0012 source=h2 target=subsection reason=relative-heading-rank; base=h2

\subsection{格林多前书 13:11}

\BibleQuote{当我作孩子的时候，说话像孩子，心思像孩子，意念像孩子；可是，当我成了人，就把孩子的事丢弃了。}\par

他又用另一个例子说明了同一件事，说道：\par

% structure: p0015 source=h2 target=subsection reason=relative-heading-rank; base=h2

\subsection{格林多前书 13:12}

\BibleQuote{因为我们现在是对着镜子观看，}接着，因为镜子模糊地呈现实物，他加了\BibleQuote{模糊不清，}以强烈表明现在的知识是极其局部的。\par

\BibleQuote{但那时要面对面。}并不是说天主有面容，而是表达更清晰明白的概念。你看到我们是如何通过逐步增加而学习的吗？\par

\Quote{如今我認知一部分；但那時我將認知，如同我被認知一樣。}你看他如何從兩方面打擊他們的驕傲？既因為他們的認知是部分的，又因為連這部分也不是出於自己。\Quote{因為我不認識祂，但祂使自己被我所認識，}他說。因此，正如如今祂首先認識我，並且親自趨向我，同樣我將在將來比現在更加趨向祂。因為那坐在黑暗中的人，只要他看不見太陽，就不會主動迎向那光芒之美，這光芒其實在開始照耀時就顯現了；但當他感知到它的光輝時，他自己終於也跟隨它的光明。這就是\Quote{如同我被認知一樣}這句話的意思。不是說我們將會認識祂如同祂是那樣，而是說，正如祂如今趨向我們，同樣我們那時也將依附於祂，並且將認識許多如今是隱秘的事，並享受那至為有福的共融與智慧。因為如果知道這麼多的保祿還是小孩子，試想那些事該是怎樣的。如果這些是\Quote{鏡子}和\Quote{謎語}，那麼你由此再次推斷，天主的面對面是何等偉大的事。\par

3. 但為使我為你們開啟這差異的一小部分，並將這思想的一絲微弱光芒傳達給你們的靈魂，我願你們回憶起律法時代的事，如今恩寵已經照耀。因為那些在恩寵之前的事，也有某種偉大而奇妙的外表；然而，請聽保祿在恩寵來臨後如何論及它們：\BibleQuote{那從前有光榮的，在這事上，因那超越的光榮，就算是沒有光榮了。}\BibleRef{格後 3:10}\par

但為使我所說的更加清楚，讓我們將論證應用於當時舉行的某種禮儀，然後你們就會看到何等的差異。如果你們願意，讓我們提出那個逾越節和這個逾越節，然後你們就會明白我們的優越。因為猶太人確實慶祝了它，但他們是\Quote{在鏡子裡，模糊地}慶祝它。但這些隱藏的奧秘，他們從未在心中想過，也不知道它們預表什麼。他們看見一隻羔羊被宰殺，一隻牲畜的血，門框被灑血；但他們從未想過天主聖子降生成人，將被宰殺，並將解放整個世界，並將使希臘人和化外人同樣品嚐這血，並將為所有人敞開天，並將天上的事物賜給全人類，並將帶著祂染血的身體高舉於諸天之上、諸天之天之上，總之，高舉於所有天使、總領天使及其他萬有的高天之上，並使祂發出無法言說的光榮，坐在君王寶座上，在聖父右邊——這些事，我說，無論是他們還是其餘人類，無人預知或曾能想像。\par

4. 但那些無所顧忌的人說什麼呢？他們說\Quote{如今我認知一部分}這話是按分施說的；因為宗徒擁有對天主的完全認識。而今他稱自己是孩子？他如何\Quote{在鏡子裡}看見？如何\Quote{模糊地}，如果他擁有知識的總和？為何他又將此作為聖神的獨有特權，而非受造界的任何其他能力？他說：\BibleQuote{除了人內裡的人靈，誰知道人的事？同樣，除了天主的聖神，誰也不知道天主的事。}\BibleRef{格前 2:11}而基督又說這只屬於祂自己，如此說：\BibleQuote{不是有人看見了聖父，唯獨那從天主來的，祂看見了聖父。}\BibleRef{若 6:46}祂將\Quote{看見}這個詞用於最清楚、最完全的認知。\par

而那認識本體的人，怎能不明白分施呢？因為那認識比這更大。\par

\Quote{那麼，}他說，\Quote{我們不認識天主嗎？}絕非如此。我們知道祂存在，但祂的本體是什麼，我們還不知道。為使你們明白他所說的\Quote{如今我認知一部分}不是關於分施，請聽下文。他接著說：\Quote{但那時我將認知，如同我被認知一樣。}他無疑是被天主所認知，而不是被分施認知。\par

因此誰也不要認為這是一件小的或簡單的過犯，而是雙重的、三重的，甚至是多重的。因為不僅有這種不敬：他們誇口知道那些只屬於聖神和天主的獨生子的事；而且當保祿連這\Quote{部分的}知識若無從上而來的啟示也無法獲得時，這些人卻說他們從自己的推理中得到了全部。因為他們也無法指出聖經何處論及了這些事。\par

5. 然而，讓我們放下他們的狂妄，專心聽從接下來關於愛的話。因為他不僅滿足於這些，又補充說：\par

% structure: p0026 source=h2 target=subsection reason=relative-heading-rank; base=h2

\subsection{格前 13:13}

\BibleQuote{如今常存的有信德、望德、愛德，這三樣；其中最大的是愛德。}\par

因為信德和望德，在所信和所望的美好事物到來時，就會終止。為證明這一點，保祿說：\Quote{因為所見的望德不是望德；誰還盼望他所見的呢？}又說：\BibleQuote{信德是所望之事的擔保，是未見之事的證據。}\BibleRef{羅 8:24}所以，當那些出現時，這些就終止了；但愛德那時最為高漲，變得更加強烈。這是對愛德的另一種讚揚。因為他並不滿足於先前所說的，還努力發現另一種。請注意：他已經說過愛德是極大的恩賜，是更卓越的道路。他說過，若沒有愛德，我們的恩賜就沒有什麼益處；他詳盡地描繪了愛德的形象；他打算再次以另一種方式高舉愛德，並從它的永存性顯示其偉大。因此他也說：\BibleQuote{如今常存的有信德、望德、愛德，這三樣；其中最大的是愛德。}那麼愛德如何更大呢？在於那些會消逝。\par

既然爱德的力量如此伟大，他有充分理由接着说：\Quote{追求爱德。}因为确实需要\Quote{追求}，并且要热切地追赶她：她如此从我们身边飞走，又有那么多事物在这方面绊倒我们。因此，我们始终需要极大的热忱才能追上她。为指出这一点，\Person{保禄}没有说\Quote{跟随爱德}，而是说\Quote{追赶}她；激励我们，点燃我们去抓住她。\par

因为天主从一开始就设计了无数方法将爱德植入我们内。首先，祂赐给全人类同一个头，亚当。为何我们不都是从土中生出？为何不是像他那样成人？为的是使生育和抚养子女，以及从他人而生，能将我们彼此联结。为此，祂也没有用土造女人：因为同一质料的东西不足以使我们因羞愧而同心，除非我们还有同一个祖先，祂也为此做了准备：因为现在，仅仅因地方相隔，我们就认为彼此是外人；那么，若是人类有两个起源，就更会如此。因此，祂如同从同一个头，将整个人类身体联结在一起。因为起初他们看似两个，请看祂如何再次将他们结合，并通过婚姻使它们合一。因为祂说：\Quote{因此，人要离开父亲和母亲，依附自己的妻子，二人成为一体。}\BibleRef{创 2:24}祂没有说\Quote{女人}，而是说\Quote{男人}，因为欲望在男人身上更强烈。是的，为此祂也使欲望更强，好让较强的一方被这情欲的绝对支配所折服，并屈服于较弱的一方。既然也必须引入婚姻，祂将女人的丈夫赐给从她所出的女人。因为在天主眼中，一切都次于爱德。如果事情这样开始后，第一个人立刻变得如此疯狂，魔鬼在他们中间播下了如此大的争战和嫉妒；那么，如果他们不是出自同一个根，魔鬼还会做什么呢？\par

此外，为使一个服从，另一个管理（因为平等常易引起纷争），祂不容许民主制，而是君主制；如同在军队中，这种秩序在每个家庭中都可以看到。例如，丈夫处于君主的地位；妻子处于副官和将军的地位；孩子们被分配第三个指挥位置。然后在他们之后是第四种秩序，即仆人的秩序。因为这些仆人也要管理他们的下属，他们中常有一位被置于全体之上，始终保持着主人的地位，但仍是仆人。同时还有另一种管理，在孩子们中间还有另一种，根据年龄和性别；因为在孩子中，女性没有同等的权威。天主处处设立了相近且密集的小秩序，好使所有人都能保持和谐与良好秩序。因此，甚至在人类增加到众多之前，当只有最初的两个存在时，祂命男人管辖，女人服从。为使男人不会因轻视她为低等而与她分离，请看祂如何尊荣她，甚至在创造她之前就使他们成为一体。因为祂说：\Quote{让我们为男人造一个相称的助手}，暗示她是为他的需要而造的，从而将他引向那为他而造的她：因为凡是为我们而做的，我们更容易善待它们。但另一方面，为使她不会因被赐予他作助手而高傲，也不会破坏这个纽带，祂从男人的肋骨造了她，表示她是整个身体的一部分。为使男人也不会因此高傲，祂不再允许那曾只属于他的东西只属于他，而是通过引入生育子女做了相反的事，并在此也将首尊给了男人，但并未允许全部属于他。\par

你看到天主建立了多少爱德的纽带吗？这些纽带祂通过自然的力量放在我们内，作为和谐的保证。因为我们的同一本质导向此（因为每种动物喜爱同类）；女人从男人而生，孩子又从两人而生。由此也产生许多种情感。因为一个我们爱如父亲，另一个如祖父；一个如母亲，另一个如乳母；一个如儿子或孙或曾孙，另一个如女儿或孙女；一个如兄弟，另一个如侄子；一个如姊妹，另一个如侄女。何必一一列举所有血亲的名称呢？\par

祂还设计了另一种情感的基础。因为禁止同族通婚后，祂引导我们走向外人，并将他们再次拉近我们。因为既然通过这种自然亲属关系，他们不可能与我们相连，祂通过婚姻重新连接我们，通过新娘一个人将整个家庭联合起来，并使整个种族与种族融合。\par

因为主说：\Quote{不可娶，}\BibleRef{肋 18:6}\Quote{你的姊妹，也不能娶你的姑母，也不能娶任何与你有这种血亲的女子，}这些完全阻止婚姻；祂列举了这种亲属关系的等级。你对他们有足够的亲情，因为你也是同一产痛的结果，而其他人与你关系不同。你为什么缩小爱德的广度？你为什么白费心机抛弃一个对她产生情感的基础，好像你可以借此为自己提供一个情感涌流的独特源泉？我是说，从另一个家庭娶妻，通过她得到一连串的亲属：母亲、父亲、兄弟以及他们的关系！\par

7. 你看见祂用多少方式把我们连结在一起了吗？即便如此，祂仍不满足，还使我们彼此需要，以便借此也将我们结合在一起，因为需要尤其能产生友谊。正因如此，祂没有让万物在每个地方都产生，从而也迫使我们彼此交往。祂使我们彼此需要，同时又使交往变得容易。因为若不是这样，事情就会以另一种方式变得痛苦和困难。因为若一个人需要医生、木匠或任何其他工匠，就必须远行到异乡，那么一切都将化为乌有。因此，祂也建立了城市，将众人聚集在一处。并且，为了我们能轻易地与远方保持交往，祂在我们之间铺展了海面，赐予我们风的迅捷，使航行变得容易。起初，祂甚至将所有人聚集在一处，直到那首先接受恩赐的人滥用他们的和谐而犯罪。然而，祂已用各种方式将我们聚拢：借着本性、血亲、语言和地方；正如祂不愿我们从乐园坠落（因为若祂愿意，祂就不会将\Quote{祂所造的人}安置在那里，而是那违背命令的人成了原因），同样，祂也不愿人有多种语言；否则祂一开始就会这样做了。但现在，\Quote{普天之下只有一种语言，大家说同样的话。} \BibleRef{创 11:1}\par

8. 这就是为什么，当大地需要被毁灭时，祂也没有用别的材料造我们，也没有将义人迁走，而是将他留在洪水之中，如同世界的一点火花，从那里重新点燃了我们的种族，借着有福的诺厄。起初祂只设立了一种主权，让男人管理女人。但我们的种族陷入极度混乱之后，祂又设立了其他主权，即主人和统治者的主权，这也是出于爱。也就是说，既然恶习容易瓦解和颠覆我们的种族，祂就在我们的城市中设立了执掌正义的人，如同医生，驱除恶习——这恶习是爱的瘟疫——从而将众人聚集在一起。\par

9. 不仅如此，为了使每个家庭中也充满极大的和谐，祂以统治和优越尊荣了男人；另一方面，祂以欲望武装了女人；生儿育女的恩赐也共同托付给了双方，此外祂还提供了其他促进爱的事物：既没有将一切都托付给男人，也没有将一切都托付给女人，而是\Quote{也将这些事分别给了每一方}；托付给她家，托付给他市场；给他养家的工作，因为他耕种土地；给她纺织的工作，因为织机和纺锤是女人的。正是天主自己赐给了女性纺织的技艺。贪婪有祸了，它不让这种差异显现出来！因为普遍的柔弱已经发展到这种地步，竟将我们的男人引向织机，将梭子、纬线和经线交到他们手中。然而，即便如此，天主眷顾的预见依然闪耀。因为我们仍大大需要女人在其他更必要的事上，并且在我们维系生活的事上需要下属的帮助。\par

10. 这种需要的强制力如此强大，以至于即使一个人比所有人都富有，他也不能摆脱这种紧密的联系，以及对他之下者的需要。因为我们看到，不仅穷人需要富人，富人也需要穷人；而且富人需要穷人更甚于穷人需要富人。为了更清楚地看到这一点，让我们假设两座城，一座只有富人，另一座只有穷人；在富人的城中没有一个穷人，在穷人的城中没有一个富人；让我们彻底排除两者，看看哪一座更能自给自足。因为如果我们发现穷人的城能够自给，那么显然富人会更需要他们。\par

11. 那么，在那座富人的城中，将没有制造商、建筑工、木匠、鞋匠、面包师、农夫、铜匠、制绳匠，也没有任何其他这类手艺。因为富人中的谁会愿意从事这些手艺呢？要知道，连那些从事这些行业的人，一旦富起来，就不再忍受这些工作带来的不便。那么，我们的这座城如何存在呢？有人回答说：\Quote{富人}，\Quote{出钱向穷人购买这些东西。}那么，他们就不能自给自足，因为他们需要别人证明了这一点。但他们如何建造房屋呢？他们也要购买这个吗？但事物的本性不允许这样。因此，他们必须邀请工匠进来，从而破坏我们当初建城时制定的法律。因为你们记得，我们说：\Quote{城中不可有穷人。}然而，需要违背我们的意愿，邀请并将他们带了进来。由此可见，没有穷人就无法存在一座城；因为如果这座城继续拒绝接纳任何这些工匠，它就不再是城，而将灭亡。显然，它不能自给自足，除非它聚集穷人如同某种保护者，留在自己里面。\par

但我们也来看穷人之城，若被剥夺了富人，是否也会陷入同样的匮乏状态。首先，让我们在论述中彻底阐明财富的本质，并将它们清楚地指出来。那么财富是什么呢？金、银、宝石，以及丝绸、紫色和绣金的衣物。既然我们已经看到了财富是什么，就让我们将它们从我们的穷人之城中驱逐出去：如果我们想使之成为一个纯粹的穷人之城，那就不要让任何金子出现在那里，连梦中也不许，也不要有那样的衣物；如果你愿意，连银子、银器也不要。那又怎样呢？因为这样，那座城及其事务就会陷入匮乏吗？告诉我？一点也不。因为，假设首先需要建造房屋：人们不需要金、银、珍珠，而是需要技艺、双手，而且不是任何双手，而是变得粗糙、手指坚硬、有大力气的手，以及木头和石头；假设又要织一件衣服：这里我们也不需要金和银，而如前所述，需要双手、技艺和做工的妇女。又如果需要耕作、挖地？需要的是富人还是穷人？每个人都很清楚，是穷人。当需要打铁或做诸如此类的事情时，我们最需要的就是这一类人。那么，还有什么方面我们需要富人呢？除非需要拆毁这座城。因为，如果那种人进来，这些哲学家（因为我称那些不追求任何多余之物的人为哲学家）开始贪求金银珠宝，沉溺于闲散和奢侈，那么从那天起，他们就会毁掉一切。\par

9. 有人问：\Quote{但除非财富有用，} 有人说，\Quote{否则天主为何赐予它呢？} 又从何证明富有是从天主而来的呢？\Quote{圣经说：\InnerQuote{银子是我的，金子也是我的，并且我愿意给谁，就给谁。}}\BibleRef{盖 2:8} 此刻，若不是有失体统，我真要哈哈大笑，嘲笑那些说这些话的人：因为如同小孩子被允许进入国王的餐桌，他们随着食物将手边的一切都塞进嘴里；同样，这些人也随着神圣的圣经偷偷带进自己的观念。因为我知道，\Quote{银子是我的，金子也是我的} 这话是由先知说的；但 \Quote{我愿意给谁，就给谁} 这话并不是正文，而是这些渣滓加进去的。至于前者为何而说，我将加以解释。哈盖先知因为在犹太人从巴比伦归来后不断应许他们，他将使圣殿恢复从前的样子，有些人对此话存疑，认为殿宇已化为尘土灰烬，要再现原貌几乎不可能——他为了消除他们的不信，以天主的位格说了这些话，仿佛说：\Quote{你们为什么害怕？为什么不肯相信？\InnerQuote{银子是我的，金子也是我的}，我不必向别人借用，就能如此装饰圣殿。} 为了表明这就是意思，他又加了一句：\Quote{这殿后来的荣耀必大过先前的荣耀。} 因此，我们不要将蜘蛛网织在皇袍上。因为若有人被发现紫色外衣上织了假线，就要受最严厉的惩罚，何况在属灵的事上呢？因为由此所犯的，不是寻常的罪。而我说，加添或删减，又何必多言？仅仅一个标点、或朗读时的一种语气，就常常引出许多不敬的思想。\par

有人问：\Quote{那么富人从何而来？} \Quote{因为经上说：\InnerQuote{富贵贫穷来自主。}} 那么，我们要问问那些以此反驳我们的人：所有的财富和所有的贫穷都是来自主吗？不，谁会这么说呢？因为我们也看到，许多人是通过抢劫、邪恶地掘墓、巫术以及其他手段积聚了大量财富，而这些拥有者甚至不配活着。那么，告诉我，我们能说这财富是来自天主吗？决不。那是从哪里来的？来自罪恶。因为，妓女因侮慢自己的身体而致富，英俊的青年常以可耻的方式出卖青春而得金，盗墓者掘开坟墓聚集不义之财，强盗穿墙窃取。因此，并非所有财富都来自天主。有人问：\Quote{那么，我们该如何解释这句话呢？} 你先要了解有一种贫穷不是来自天主，然后我们再来看这句话本身。我的意思是，当一个放荡的青年把财富花在妓女、术士或任何其他邪恶欲望上而变得贫穷，难道不是很清楚这不是来自天主，而是来自他自己的放荡吗？再者，如果有人因懒惰而贫穷，因愚昧而沦为乞丐，因从事危险和非法的勾当而贫穷，难道不是很明显，这些人以及其他类似的人都不是被天主带入贫穷的吗？\par

有人问：\Quote{那么圣经会说谎吗？} 绝对不可！但那些忽略仔细审查所写一切的人，是愚昧的。因为，如果一方面承认圣经不能撒谎，另一方面又证明并非所有财富都来自天主，那么困难的原因就在于不经考虑的读者的软弱。\par

10. 现在我们本应当遣散你们，既已在此为圣经开脱，好让你们因对圣经的疏忽而在我们手中受此惩罚；但因我极其宽恕你们，再也无法忍受看你们困惑不安，让我们也加上解答，先提说说话者，话语何时说，以及对谁说。因为天主不是同样对所有人说话，就像我们对待孩子和成人也不同。那么，这话何时说的，由谁说的，对谁说的？由\Person{撒罗满}在旧约中对犹太人说的，他们只认识感官之事，并借此证明天主的权能。因为这些人说：\Quote{祂也能赐食吗？}以及\Quote{祢显什么神迹给我们看？我们的祖先在旷野吃过玛纳。}——\Quote{他们的天主是自己的肚腹。}\BibleRef{咏 78:24}既然他们借这些事来试探祂，祂便对他们说：\Quote{天主也能使人富足使人贫穷。}并非说当然祂自己使人如此，而是说祂愿意时就能。正如祂说：\Quote{祂斥责海，使它干涸，使一切河流枯竭，}\BibleRef{纳 1:4}然而这事从未发生过。那么先知为何这样说？并非如常发生，而是说祂有能力行这事。\par

那么祂赐下的是何种贫穷，何种富足？请记住圣祖，你们就知晓天主所赐的富足。因为祂使\Person{亚巴郎}富裕，后又使\Person{约伯}富裕，正如\Person{约伯}自己说：\Quote{我们从主领受了福，难道不也领受祸吗？}\BibleRef{约 2:10}\Person{雅各伯}的富足也由此开始。还有一种贫穷来自祂，是受称赞的，就如祂曾想使那富人认识，说：\Quote{你若愿意成全，去变卖你的家产，施舍给穷人，然后来跟随我。}\BibleRef{玛 19:21}又对门徒们立法说：\Quote{不要带金、银，也不要带两件衣服。}\BibleRef{玛 9:10}所以不要说一切富足都是祂的恩赐，因为已指出有些是由凶手、抢劫和无数其他手段聚集的。\par

然而讲论又回到我们先前的问题：即\Quote{若富人丝毫不能帮助我们，他们为何成为富人？}那么我们该说什么？那些自己致富的人没有用处，而那些由天主所造致富的人则极有用处。就从我们刚才提到的人所行的事上学习这一点。\Person{亚巴郎}拥有财富，是为了所有客旅和一切需要者。因为他看见三人走近，以为是凡人，就宰了一头牛犊，揉了三斗细面，并在日头正烈时坐在帐棚门口；请思量他如何慷慨乐意将自己的财物用在众人身上，连同货物也献上自己身体的服事，且年事已高；他成了客旅的港口，给所有有各种需要的人，不将任何东西据为己有，甚至连儿子也不保留：因为按天主的命令，他竟献出了儿子；连同儿子，他也献上自己和全家，当他急忙去救侄子脱离危险时；他这样做并非为利，而是出于纯全的人情。例如，被他所救的人愿将战利品交他处理，他却拒绝一切，甚至\Quote{一根线，一条鞋带。}\BibleRef{创 14:23}\par

有福的\Person{约伯}也是如此。他说：\Quote{我的门对一切来者敞开；}\BibleRef{约 20:15}\Quote{我作盲人的眼，跛者的脚；我是无助者的父亲，客旅不宿在外；无助者无论有何需要，都不落空，也不许一个空怀的人离开我的门。}还有许多类似的事，我们不再一一细数，他继续如此行，把一切财富花在穷人身上。\par

你们也看看那些非由天主而致富的人，好知道他们如何使用财富。请看\Person{拉匝禄}比喻中的那富人，他连一点饼屑也不分给人；看\Person{阿哈布}，连葡萄园也被他强夺；看\Person{革哈齐}；看所有这样的人。因此，那些正当获取财富的人，既从天主领受，就用于天主的诫命；而那些在获取中得罪天主的人，在花费上也同样行：用在娼妓和寄生虫上，或埋藏囤积，而不施舍穷人。\par

\Quote{那么，}有人说，\Quote{天主为何容忍这样的人富有？}因为祂是宽仁的；因为祂要引领我们悔改；因为祂预备了地狱；因为\Quote{祂已定了日子，要按正义审判天下。}\BibleRef{宗 17:31}若祂立即惩罚那不按德性致富的人，则\Person{匝凯}就没有时间悔改，以至四倍偿还一切不义之财，并加上一半家产；\Person{玛窦}也没有机会悔改而成为宗徒，会在未到适当季节就被除去；其他许多人也是如此。因此祂容忍他们，召叫众人悔改。但若他们不愿，仍坚持不变，他们就将听到\Person{保禄}说，他们仗着顽梗不悔的心，为自己积蓄忿怒，直到忿怒的日子，即天主公义审判的日子。\BibleRef{罗 2:5}愿我们逃脱那忿怒，成为天国的富人，追求可称赞的贫穷。这样我们也将获得未来的美善；愿我们众人因我们的主\Person{耶稣}基督的恩宠和慈悯，偕同圣父及圣神，获得光荣、权能和尊崇，自今至永远，无穷世。阿们。\par

\SourceRecord{Translated by Talbot W. Chambers.；Nicene and Post-Nicene Fathers, First Series；Vol. 12.；Edited by Philip Schaff.；Buffalo, NY: Christian Literature Publishing Co.,；1889.；<http://www.newadvent.org/fathers/220134.htm>.}

% structure: p0001 source=h1 target=section reason=page-title

\section{论格林多前书第三十五讲}

% structure: p0002 source=h2 target=subsection reason=relative-heading-rank; base=h2

\subsection{格前 14:1}

\BibleQuote{你们要追求爱德，但也要切慕神恩，尤其是要说先知话。}\par

因此，当他把爱德的一切卓越之处详尽地叙述给他们之后，便在下面的话中劝勉他们欣然把握爱德。所以他也说：\BibleQuote{要追求它。}因为追赶的人只看见所追赶的，并向它努力，直到捉住它才罢休。追赶的人当自己无法追上时，便借助前面的人捉住逃者，恳求那些靠近它的人以极大的热忱将它抓住，并为他保持捉住的状态，直到他赶来。我们也当如此行。当我们自己不能达到爱德时，让我们请那些靠近它的人抓住它，直到我们赶上；当我们捉住它后，不要再让它离开，以免它再次逃脱。因为如果我们不按应有的方式使用它，而将一切置于它之下，它就会不断从我们身边溜走。所以我们必须竭尽全力，以便完全保有它。因为如果做到了这一点，我们今后就不需要太多劳力，甚至几乎不需要；而是安息、度假，在德行的窄路上行进。因此他说：\BibleQuote{要追求它。}\par

然后，为了不让他们以为他引入爱德的谈论只是为了熄灭神恩，他接着说了以下的话：\par

\BibleQuote{但也要切慕神恩，尤其是要说先知话。}\par

% structure: p0007 source=h2 target=subsection reason=relative-heading-rank; base=h2

\subsection{格前 14:2}

\BibleQuote{因为说语言的，不是对人说话，而是对天主说话，因为没有人听得懂；他是在神魂里讲论奥秘的事。}\par

% structure: p0009 source=h2 target=subsection reason=relative-heading-rank; base=h2

\subsection{格前 14:3}

\BibleQuote{但那说先知话的，却是对人说话，建树、劝勉和安慰。}\par

在这里，他对神恩作了比较，降低了说语言的神恩的地位，表明它既非完全无用，也非本身很有益处。因为事实上他们为此很骄傲，认为这神恩是伟大的。它被认为是伟大的，因为宗徒们首先接受了它，并且伴随着如此大的显耀；然而它并不因此就被认为高于其他一切神恩。那么为什么宗徒们比其他人在先接受了它呢？因为他们要到各处去。正如在建塔时期，一种语言分裂为多种；同样，那时多种语言经常汇聚在一个人身上，同一个人用波斯语、罗马语、印度语和多种其他语言说话，圣神在他内发声；这神恩被称为说语言的神恩，因为他能同时说各种语言。请看他是如何既贬低它又提升它。比如，通过说：\BibleQuote{说语言的，不是对人说话，而是对天主说话，因为没有人听得懂，}他贬低了它，暗示它的益处不大；但通过加上\BibleQuote{但他是在神魂里讲论奥秘的事，}他又提升了它，好使它不显得多余、无用和白白赐予。\par

\BibleQuote{但那说先知话的，却是对人说话，建树、劝勉和安慰。}\par

你看见他是通过什么来表明这神恩的优选性质吗？即通过共同的益处。他如何总是把更高的荣誉赋予那有利于众人的东西？难道前者不是也对人说话吗？告诉我。但不是那\Quote{建树、劝勉和安慰}。因此，受圣神感动是两者共有的，即说先知话的和说语言的；但在这一点上，后者（我指的是说先知话的）的优势在于他也对听者有益。因为说语言的人不被那些没有这神恩的人理解。\par

那么，他们不建树任何人吗？他说：\Quote{是的，只建树自己。}因此他也加上了：\par

% structure: p0015 source=h2 target=subsection reason=relative-heading-rank; base=h2

\subsection{格前 14:4}

\BibleQuote{那说语言的，是建树自己。}\par

如果他自己不知道自己在说什么，那又怎么建树呢？目前，他指的是那些理解自己所说的人——他们自己理解，但不知道如何传达给别人。\par

\BibleQuote{但那说先知话的，是建树教会。}正如单个人与教会的差别有多大，这两者之间的间隔也有多大。你看见他的智慧了吗？他并没有排除这神恩使之毫无价值，而是表明它有一些益处，尽管很小，只够拥有者自己吗？\par

2. 其次，免得他们以为他因嫉妒而贬低说语言的神恩（因为多数人有这神恩），为了纠正他们的猜疑，他说：\par

% structure: p0020 source=h2 target=subsection reason=relative-heading-rank; base=h2

\subsection{格前 14:5}

\BibleQuote{我愿意你们都说语言，但更愿意你们说先知话；因为那说先知话的，比那说语言的更大，除非他解释，使教会得到建树。}\par

但\Quote{更愿意}和\Quote{更大}并不表示对立，而是优越。所以由此也清楚，他并非贬低神恩，而是引导他们走向更好的事物，显明了他对他们的关怀和毫无嫉妒的精神。因为他没有说\Quote{我愿意两三个人}，而是\Quote{你们都}说语言，并且不仅于此，还\Quote{愿意你们说先知话}；并且后者胜于前者；\Quote{因为那说先知话的更大。}既然他确立并证明了这一点，接着便也断言；但不是简单地断言，而是有条件地。因此他加上了：\BibleQuote{除非他解释，}因为他若能这样做，我指的是解释，他说：\Quote{他就成了与说先知话的相等的了，}因为那时有许多人受益；这一点尤其值得注意，他如何始终以这为目的，超越一切。\par

% structure: p0023 source=h2 target=subsection reason=relative-heading-rank; base=h2

\subsection{格前 14:6}

\BibleQuote{但如今，弟兄们，如果我去你们那里说语言，对你们有什么益处呢？除非我对你们说话，或借着启示，或借着知识，或借着先知话，或借着教训。}\par

\Quote{我何必说，}他说，\Quote{其余的呢？不，那说舌音的，即使他是\Person{保禄}，对听众也毫无益处。}他说这些话，是为了表明他寻求他们的益处，并非对拥有那恩赐的人怀有任何怨恨；因为即使在他自己身上，他也不避讳指出其无益之处。实际上，他惯常的方式是在自己身上处理令人不悦的话题：正如在书信开头他所说：\Quote{那么\Person{保禄}是谁？\Person{阿颇罗}是谁？\Person{刻法}是谁？}他在这里也这样做，说：\Quote{即使是我，除非我以启示、或先知话、或知识、或教导的方式向你们说话，也对你们无益。}他的意思是：\Quote{若我不说一些你们能理解且清楚明白的事，而只是炫耀我有说舌音的恩赐——你们不懂的舌音，你们将毫无所得。因为你们如何能从不懂的声音中获益呢？}\par

% structure: p0026 source=h2 target=subsection reason=relative-heading-rank; base=h2

\subsection{\BibleRef{格前 14:7}}

\BibleQuote{甚至无生命之物，发出声音，无论是笛或是琴，若声音没有分别，怎知所吹奏的是什么？}\par

\Quote{我为何说，}他说，\Quote{在我们这里这是无益的，唯有清楚易懂的对听众才有用？}因为就是在无生命的乐器上也能看到这一点：无论是笛或是琴，若被杂乱无章、不熟练地弹奏或吹奏，没有恰当的节奏与和声，就不能吸引任何听众。因为就是在这些无声音的乐器中，也需要某种清晰：你若不按艺术规则弹奏或吹气，就一无所成。既然如此，我们从无生命之物尚且要求如此多的清晰、和谐与适当，并努力在这些无声音的乐器中注入如此多的含义，那么在赋有生命和理性的人身上，以及在属神的恩赐中，就更应当以意义为目标了。\par

% structure: p0029 source=h2 target=subsection reason=relative-heading-rank; base=h2

\subsection{\BibleRef{格前 14:8}}

\BibleQuote{因为若号筒发出模糊的声音，谁能准备作战呢？}\par

如此，他从纯装饰性的事物推进到更必要和有用的事物；并且说，不仅在琴上，在号筒上也能看到这种效果。因为在号筒中也有一定的调子；它们时而发出战争的信号，时而发出非战争的信号；又有时引领出战，有时召回。若不明白这一点，就有大危险。这正是他所说的意思，以及其中的危害，当他这样说时：\BibleQuote{谁能准备作战呢？}所以，若不具备这种性质，就是一切毁灭。\Quote{这对我们有何关系？}有人说？这确实特别关系到你们；因此他也加上：\par

% structure: p0032 source=h2 target=subsection reason=relative-heading-rank; base=h2

\subsection{\BibleRef{格前 14:9}}

\BibleQuote{你们也是如此，除非用舌头说出易懂的话，你们将是向空气说话：}即呼无人，说无人。因此处处显出它的无益。\par

4. \Quote{但若它无益，为何赐下？}有人说？是为了使拥有它的人获益。但若要也造福他人，则必须加上解释。他说这话，是为了使他们彼此接近；如果某人自己没有解释的恩赐，他可以找另一个有这恩赐的人，通过他使自己的恩赐有益。因此他处处指出其不完善之处，以便将他们联结在一起。无论如何，认为恩赐自身足够的人，并非称赞它，反而贬低了它，因为不让它在解释中闪耀。因为那恩赐是美好且必要的，但只在有人解释所说的话时才是。手指也是必要的，但若将其与其他肢体分离，就不那么有用了；号角是必要的，但若随意吹响，反而令人烦恼。是的，任何艺术若无适宜的材料，都无法显明；材料若无形式赋予，也不会成型。假设声音如同材料，清晰如同形式，若没有形式，材料就毫无用处。\par

% structure: p0035 source=h2 target=subsection reason=relative-heading-rank; base=h2

\subsection{\BibleRef{格前 14:10}}

\BibleQuote{世上的声音，或许有那么多种类，却没有一种是没有意义的：}即如此多的舌头，如此多的声音：斯基泰人、色雷斯人、罗马人、波斯人、摩尔人、印度人、埃及人，以及无数其他民族。\par

% structure: p0037 source=h2 target=subsection reason=relative-heading-rank; base=h2

\subsection{\BibleRef{格前 14:11}}

\BibleQuote{那么，我若不明白那声音的意义，我在说话者眼中就是化外人。}\Quote{你们不要以为，}他说，\Quote{这事只发生在我们身上；实际上在所有人身上都能看到这一点：我这样说并非贬低那声音，而是表示对我来说，只要它不明了，就是无用的。}接着，为了不使指责令人不快，他对双方提出同样的指责，说：\BibleQuote{他在我眼中是化外人，我在他眼中也是化外人。}你看见他如何一步步把人引向与主题相关的事。这正是他一贯的做法：从远处取来例证，最终以更切合主题的事结束。因为讲了笛与琴（其中有诸多低劣和无益之处）之后，他谈到号筒，一种更有用的东西；接着，从那里他进到声音本身。同样，以前当他论述宗徒接受馈赠并非不允许时，他从农夫、牧人和士兵开始，然后把话题带到更切近主题的事物，即旧约中的司祭。\par

但我求你思量，他如何处处尽力使恩赐免于责难，而将指责转回给接受它的人。因为他没有说：\Quote{我成了化外人，}而是：\Quote{在说话者眼中，我成了化外人。}再者，他没有说：\Quote{说话者成了化外人，}而是：\Quote{说话者在我眼中成了化外人。}\par

5. \BibleQuote{那麼，該怎麼辦呢？}他說。不但不應輕視，反而應該推薦和教導；他自己也正是這樣做的。因為在指責、斥責並指出其無益之後，他繼續勸告他們，說：\par

% structure: p0041 source=h2 target=subsection reason=relative-heading-rank; base=h2

\subsection{\BibleRef{格前 14:12}}

\BibleQuote{你們也這樣，既然你們熱切渴望神恩，就當尋求多多建造教會。}\par

你看到他的目標無處不在嗎？他如何在所有情況下都持續關注一件事：普遍的益處，教會的建立；並以此為某種規則？他沒有說\Quote{好讓你們獲得恩賜}，而是說\Quote{好讓你們多多擁有}，即甚至大大地擁有它們。因此，我絕不願你們不擁有它們，反而願你們多多擁有，只要你們為共同的益處而運用它們。這如何做到？他接著說：\par

% structure: p0044 source=h2 target=subsection reason=relative-heading-rank; base=h2

\subsection{\BibleRef{格前 14:13}}

\BibleQuote{因此，那說語言的，該祈求自己能解釋。}\par

% structure: p0046 source=h2 target=subsection reason=relative-heading-rank; base=h2

\subsection{\BibleRef{格前 14:14}}

\BibleQuote{因為若我用語言祈禱，我的神魂祈禱，但我的理智卻不結果實。}\par

% structure: p0048 source=h2 target=subsection reason=relative-heading-rank; base=h2

\subsection{\BibleRef{格前 14:15}}

\BibleQuote{那麼，該怎麼辦呢？我要用神魂祈禱，也要用理智祈禱；我要用神魂歌唱，也要用理智歌唱。}\par

在此他指出，獲得恩賜在於他們的能力。因為，\Quote{讓他祈禱，}他說，即\Quote{讓他盡自己的本分，}因為你若勤懇祈求，必會獲得。因此祈求，不僅要有語言的恩賜，也要有解釋的恩賜，好使你能對眾人有益，而不將恩賜封閉在自己內。\BibleQuote{因為若我用語言祈禱，我的神魂祈禱，但我的理智卻不結果實。}他說，你看到嗎？他逐步將論證聚焦，表明這樣的人不僅對他人無用，對自己也無用；至少\BibleQuote{他的理智不結果實}？因為若一個人只說波斯語或任何其他外國語言，卻不明白自己所說的，那麼對他本人而言，他也成了野蠻人，不僅對他人，因為他不知道聲音的意義。因為古時有許多人也有祈禱的恩賜，連同語言；他們祈禱，舌頭說話，用波斯語或拉丁語祈禱，但他們的理智不知道所說的是什麼。因此他也說：\BibleQuote{若我用語言祈禱，我的神魂祈禱，但我的理智卻不結果實。}即那賜予我、感動我舌頭的恩賜，\par

那麼，什麼是本身最好、有益的呢？人該如何行動，或向天主求什麼？要用\Quote{神魂}，即恩賜，也用\Quote{理智}來祈禱。因此他也說：\BibleQuote{我要用神魂祈禱，也要用理智祈禱；我要用神魂歌唱，也要用理智歌唱。}\par

6. 他在此也表明同樣的事：既讓舌頭說話，也讓理智不曉得所言。因為若不如此，就會有另一種混亂。\par

% structure: p0053 source=h2 target=subsection reason=relative-heading-rank; base=h2

\subsection{\BibleRef{格前 14:16}}

\BibleQuote{不然，你若用神魂祝福，那佔據無知者席位的人，怎能在你感謝時說\InnerQuote{阿們}？因為他不知道你所說的。}\par

% structure: p0055 source=h2 target=subsection reason=relative-heading-rank; base=h2

\subsection{\BibleRef{格前 14:17}}

\BibleQuote{你固然感謝得好，但別人卻不得建立。}\par

請注意，他再次如何使他的石頭對準準繩，到處尋求教會的建立。他所說的\Quote{無知者}指平信徒，並表明當他們不能說\InnerQuote{阿們}時，也遭受不小的損失。他所說的是：\BibleQuote{你若用野蠻語言祝福，不知道自己所說的，也不能解釋，平信徒就不能回應\InnerQuote{阿們}。因為沒有聽到結尾的\InnerQuote{永遠}一詞，他就不說\InnerQuote{阿們}。}然後，為了安慰他，以免他似乎看輕這恩賜，他如上文所說\Quote{他說奧秘}、\Quote{他向天主說話}、\Quote{他建立自己}、\Quote{他用神魂祈禱}，從這些事中給予不小的安慰，在此也說：\BibleQuote{你固然感謝得好，}因為你受聖神感動而說話；但別人什麼也聽不到，不知道所說的，站在那裡，從中得不到什麼益處。\par

7. 再者，因為他曾貶低擁有此恩賜的人，似乎他們沒有什麼了不得；為了不顯得他輕視他們，好像自己缺乏這恩賜，請看他說什麼：\par

% structure: p0059 source=h2 target=subsection reason=relative-heading-rank; base=h2

\subsection{\BibleRef{格前 14:18}}

\BibleQuote{我感謝天主，我說語言比你們眾人更多。}\par

他在另一處也這樣做：意欲除去猶太主義的優勢，表明它們從今以後無用，他先宣稱自己曾擁有它們，而且極其豐富；然後他稱它們為\Quote{損失}，這樣說：\BibleQuote{若有人自信在肉身方面，我更可以：第八天受割損，是以色列族、本雅明支派，希伯來人所生的希伯來人；就法律說，是法利塞人；就熱忱說，是迫害教會者；就法律的正義說，是無懈可擊的。}\BibleRef{斐 3:4}然後，在表明自己超越眾人之後，他說：\BibleQuote{然而從前對我有益的，我為了基督都看作是損失。}他在此也這樣做，說：\BibleQuote{我說語言比你們眾人更多。}因此你們不要誇耀，好像只有你們有這恩賜。因為我也擁有它，甚至比你們更多。\par

% structure: p0062 source=h2 target=subsection reason=relative-heading-rank; base=h2

\subsection{\BibleRef{格前 14:19}}

\BibleQuote{但在教會中，我寧願用理智說五句話，為能教導別人。}\par

这是什么？\Quote{以我的理智说话，使我也可以教导别人吗？}\Quote{明白我所说的，}和\Quote{既能解释给别人听，又能明智地说话，并教导听者的话语。}\Quote{胜过一万句的语言。}为什么呢？他说：\Quote{为要教导别人。}因为前一种是只有展示而已，后一种则有大益处；这正是他处处寻求的，我指的是公共益处。然而语言之恩是奇怪的，但先知话却是熟悉而古老的，且从前已赐给许多人；这语言之恩反而是那时首次赐下：但他却不太看重它。因此他也没有使用它；不是因为他不拥有它，而是因为他总是寻求更有益处的事：他既完全不受虚荣的束缚，只考虑一件事，即如何使听者变得更好。\par

8. 这就是他具有考虑对自己和他人都有益的事的能力的原因，即因为他摆脱了虚荣。既然那确实被虚荣奴役的人，不但不能明辨什么对别人有益，甚至连自己的益处也不知晓。\par

\Person{西满}就是这样的人，因为他追求虚荣，甚至看不出自己的益处。犹太人也如此，他们因此将自己的救恩献给魔鬼。从此偶像也产生了，外邦哲学家因这种疯狂自我兴奋，在他们虚假的教义中失事沉没。观察这种情欲的乖戾：可见因它，有些人使自己贫穷，另一些人却渴求财富。它的暴政如此强大，甚至在截然相反的事物中也能得势。因此一个人以贞洁自夸，相反另一个人以奸淫自夸；这个以正义，另一个以不义；对于奢华与禁食、谦虚与鲁莽、财富与贫穷也是如此。我说贫穷：因为有些外人，当他们有能力接受时，为了博取赞赏而拒绝接受。但宗徒们并非如此：他们摆脱虚荣，用自己的行为证明了：即当有人称他们为神，并准备向他们献上戴花环的牛时，他们不仅禁止此事，甚至撕裂了自己的衣服。\BibleRef{宗 14:13}后来他们扶起瘸子，当所有人张口凝视他们时，他们说：\Quote{你们为什么这样注视我们，好像凭我们自己的能力使这人行走吗？}那些在赞赏贫穷的人群中，为自己选择了贫穷的状态；但这些却在鄙视贫穷、称赞财富的人群中。并且这些若收到什么，就用来服侍有需要者。因此，并非虚荣，而是仁爱，是他们一切行为的动机。但那些人恰恰相反；他们做这些事如同我们共同本性的敌人和害虫，别无其他。因此一个人为了毫无益处的目的，将全部财产沉入海底，模仿愚人和疯子；另一个人将全部土地转为公地放羊。所以他们做一切事都是为了虚荣。但宗徒们并非如此；相反，他们既接受给他们的东西，又以如此大的慷慨分施给有需要者，以致他们甚至生活在持续的饥饿中。但如果他们迷恋光荣，就不会实行这种接受和分施，因为怕引起对他们的怀疑。因为那为光荣而舍弃自己财产的人，更会拒绝接受别人的东西，以免被认为需要他人，也不招致任何怀疑。但这些宗徒，你们看到他们既服务穷人，又亲自为他们乞讨。他们实在比任何父亲更慈爱。\par

9. 也要观察他们的准则，如何适度且摆脱了一切虚荣。他说：\Quote{只要有食物和遮体之物，我们就当知足。}\BibleRef{弟前 6:8}不像那个\Place{锡诺普}人，他身穿破衣、住在木桶里，毫无益处，使许多人惊奇，却无益于人；而\Person{保禄}并未做这些事（因为他也不着眼于炫耀），反而穿着普通衣服，体面得体，一直住在房子里，并在其他一切德行的实践上完全精确。那位犬儒哲学家轻视这些，生活不洁，公开使自己蒙羞，被他对光荣的疯狂激情拖走。因为如果有人问他住在桶里的原因，会发现唯有虚荣。但\Person{保禄}也为他在罗马所住的房子付了房租。虽然他能做更严厉的事，更能有力量实行这个。但他不看重光荣——那残忍的怪物，那可怕的恶魔，那世界的害虫，那毒蛇。因为，正如那种动物用牙齿撕裂生它的母腹，同样这种情欲也撕裂产生它的人。\par

10. 那么，怎样才能为这多种病症找到药方呢？就要提出那些已将其践踏在脚下的人，并注目于他们的形象，以此规整自己的生活。因为，就如圣祖亚巴郎——但请勿责备我重复，因为我常提到他，且在各种场合提及；这正是最显明他奇妙之处，并使那些拒绝效法他的人毫无借口。因为，如果我们展示某人在这一点上行得好，另一人在那一点上行得好，有人可能会说德行难以达到；因为所有圣人们每人只完成了一部分，而要将它们全部集于一身几乎不可能。但当同一个人被发现拥有一切时，那些在法律和恩宠之后仍不能达到法律和恩宠之前的人所达到的同等程度的人，还有什么借口呢？那么，这位圣祖是如何战胜并征服这个怪物呢？当他与侄儿发生争执时。\BibleRef{创 13:8} 因为事实是，他处于劣势并失去了头份，却未烦恼。但你们知道，在俗人眼中，这类事情中的羞辱比损失更糟糕，尤其当一个人像他当时那样拥有一切权力，并且率先给予尊重，却没有得到回报时。然而，这些事没有一件使他烦恼，他满足于居于次位，并且当被年轻人——年老的他，被侄儿——亏负时，他并未愤慨或介意，反而同样爱他并服侍他。此外，在那场伟大而可怖的战役中得胜，并有力地击退了蛮族\EditorialNote{创 14}之后，他并未在胜利上添加炫耀，也未立起凯旋柱。因为他只愿拯救，而非显扬自己。此外，他款待客旅，却并未在此事上虚荣自大，而是亲自跑去迎接他们，并向他们俯伏，仿佛不是施予，而是领受恩惠；他称他们为主，却不知来者是谁，并将自己的妻子当作婢女献上。\EditorialNote{创 18}而在此前在埃及时，当他显得如此非凡，并重新得回这女人——他的妻子，且享有了如此大的尊荣\EditorialNote{创 12}，他并未向任何人炫耀。尽管当地人称他为王子，他却亲自付了墓地的价钱。\BibleRef{创 23:6}当他派人去为儿子聘娶妻子时，他没有吩咐要高高在上、威严地谈论他，\EditorialNote{创 24}而只是要把新娘带来。\par

11. 你也要查考那些在恩宠之下的人的行为吗？那时从各方面都有教义的荣耀环绕他们，你也要看到这情欲被驱逐吗？请你想想，这位讲论这些事的宗徒自己，是如何总是将一切归于天主，如何不断提及自己的罪过，而从不提及自己的善行——除非为了纠正门徒或许需要；即使被迫这样做，他也称之为愚妄，并将首位让给伯多禄，不以与普黎史拉和阿桂拉一同劳苦为耻，处处热切地显示自己谦卑，不在市场上招摇，也不带领人群，而是置身于无名之辈中。故此他也说：\BibleQuote{他身在的却是软弱。}\BibleRef{格后 10:10} 即，易于被轻视，毫无炫耀之意。又说：\BibleQuote{我求你们不要作恶，并非要我们显得中悦。}再者，他若轻视这荣耀又何足为奇？因为他轻视天堂的荣耀、天国和地狱，为了那中悦基督的事：他愿自己为基督的荣耀而被弃绝于基督。因为他说他愿意为犹太人受这事，乃是为此，免得那些无知的人将应许归于自己。因此，他若准备放弃那些事，那么他轻视属人之事又何足为奇呢？\par

12. 然而我们这个时代的人却被万事压倒，不仅被好名的欲望，也被侮辱和怕受羞辱所压倒。因此，若有人称赞，它就使你自大；若有人责备，它就使你沮丧。如同软弱的身体易受常见伤害，匍匐于地上的灵魂也是如此。因为这样的人不仅贫穷能毁掉，财富也能；不仅忧愁，喜乐亦然；顺境比逆境更甚。因为贫穷迫使你明智，而财富往往引你陷入某种大恶。正如发烧的人在任何事上都难以取悦，同样，心灵败坏的人从各方面都会受伤。\par

因此，既知道这些事，我们不要逃避贫穷，也不要羡慕财富；而要预备我们的灵魂足以应付一切处境。因为，建造房屋的人并不考虑如何既不让雨水下降，也不让阳光照射（因为这是不可能的），而是考虑如何使它能够承受一切。同样，造船的人并不设计或想出任何办法来阻止波浪冲击船身，或在海中兴起风暴（因为这也是不可能的），而是使船的两侧准备好迎接一切。再者，照料身体的人并不追求温度没有不均衡，而是使身体易于承受这一切。所以，我们也要这样对待灵魂：既不焦虑地逃避贫穷，也不渴望成为富有，而是调节二者以获得我们自身的平安。\par

因此，让我们放下这些事，使我们的灵魂既能面对财富也能面对贫穷。因为即使没有人类所遭遇的灾祸（这几乎是不可能的），那不求财富但知道如何轻松承受一切的人，也总是比那永远富有的人更好。为什么呢？首先，这样的人安全来自内在，而另一个人安全来自外在。正如一个依靠自己身体力量和战斗技巧的士兵，比一个只靠盔甲力量的人更好；同样，依靠财富的人，与以德行为保障的人相比，是低劣的。其次，即使他不陷于贫穷，他也免不了烦恼。因为财富有许多风暴和烦恼；但德行并非如此，只有喜乐与平安。是的，它使人远离那些为他设下陷阱的人，但财富恰恰相反，使人容易遭受攻击和捕获。如同动物中，鹿和兔因其天生的胆怯最容易被捕获，而野猪、公牛和狮子则不易落入埋伏者的手中；同样，在富人和自愿过贫穷生活的人身上，也可以看到这一点。前者如同狮子和公牛，后者如同鹿和兔。富人惧怕谁呢？岂非有强盗、权贵、嫉妒者和告密者？何必提强盗和告密者，当一个人甚至怀疑自己的家人时？\par

13. 我为何说\Quote{当他活着的时候？}即使死后，他也未能脱离强盗的恶行，死亡也无法使他安全，恶人甚至在他死后仍然劫掠他，财富是如此危险的东西。因为他们不仅挖进房屋，甚至撬开坟墓和棺材。还有什么比这人更可悲的呢？因为连死亡也不能给他这种安全，那悲惨的身体，即使失去了生命，也不能脱离生活的邪恶——那些行此恶事的人甚至与尘土和灰烬争战，且比活着时更猛烈？因为那时，他们或许进入他的库房，搬走他的箱子，但不会触碰他的身体，也不会剥去他身上的衣物；但现在盗墓者那可咒的手甚至不放过这些，而是移动它、翻转它，并以极大的残酷侮辱它。因为当它被埋入地下后，他们剥去它上面的泥土和裹尸布，就这样把它丢弃。\par

有什么敌人像财富那样致命呢？它既毁灭活人的灵魂，又侮辱死人的身体，甚至不许它安葬在地里——连判了罪的人和犯了最卑鄙罪行的人也能享有入土为安的待遇。因为立法者对他们判处死刑后，不再追问；但财富甚至在死后还对其主人施加最痛苦的惩罚，使他们赤身露体、无人埋葬，成为可怕而凄惨的景象：因为那些因判决和法官的愤怒而遭受此难的人，所受的痛苦也不及这些。因为前者只是一两天无人埋葬，然后就被下葬；而这些人在被埋葬后，却被剥光衣服并受侮辱。如果盗贼没有连棺材一起拿走，那也不是因为财富的功劳，而是因为贫穷的保护。因为正是贫穷守护了它。而十分肯定的是，如果我们连棺材也用财富来装饰，不用石头而用金子来打造，那么我们连这个也会失去。\par

财富是如此不可靠的东西；它与其说属于拥有它的人，不如说属于那些企图夺取它的人。因此，试图证明财富不可抗拒的论点实属多余，因为连在死亡之日，拥有者也不能得到安全。然而，有谁不与逝者和解呢？无论是野兽、魔鬼还是别的什么？仅凭景象就足以使一个铁石心肠、麻木不仁的人软化。因此，你知道，当人看见一具尸体时，即使是他公开或私下的仇敌，他也会像最亲爱的人一样为他哭泣；他的怒气随着生命而消散，怜悯随之而来。在哀悼和送葬的时候，不可能区分敌人和非敌人。所有人都如此尊重共同的本性和由此引入的习俗。但财富即使在死后也不放松对主人的愤怒；相反，它使那些从未受其伤害的人成为死者的敌人——至少剥去尸体是心怀怨恨和敌意的行为。本性使甚至仇敌在那时也与他和解；但财富却使那些无可指责的人成为敌人，并残忍地对待那孤苦无依的尸体。而在那种情况下，有许多事会引人怜悯：它是尸体，不能动弹，趋向尘土和腐朽，无人援助；但这些都不能软化那些可咒的恶人，因为他们被卑鄙的贪欲所控制。因为贪婪之情，如同无情的暴君，就在眼前，命令那些不人道的命令，使他们变成野兽，然后带他们到坟墓前。是的，如同野兽攻击死者，如果死者的肢体对他们有用，他们甚至不会不食其肉。这就是我们享受财富的结果：甚至在死后受侮辱，被剥夺连最凶恶的罪犯也能得到的埋葬。\par

14. 我们还爱这样可怕的敌人吗？我恳求你们，不要，我的弟兄们！让我们不回头地逃离它；如果它落入我们手中，不要把它藏在里面，而要用穷人的手把它捆绑起来。因为这是更能束缚它的锁链，从这些宝库中它再也不会逃脱；这样，这个不可靠的东西在将来就变得可靠、驯服、温顺，因为慈悲的右手对它有这种效果。\par

正如我先前所说的，如果它曾来到我们这里，就让我们把它交给她；但若它不来，就不要追求它，也不要烦恼，更不要认为它的拥有者是幸福的；因为这是怎样的幸福观念呢？除非你也会说，那些与野兽搏斗的人是值得羡慕的，因为那些价格高昂的野兽被这些竞赛的提议者关闭并保留给自己：然而他们自己却不敢接近或触摸它们，反而因它们而恐惧战栗，便将别人交给它们。我说，富人的情形与此类似，当他们把财富关闭在宝库中，如同一些凶猛的野兽，并且日复一日地从它那里受到无数的创伤：后者与野兽不同：因为野兽，当你牵它们出来时，那时才会伤害遇到它们的人；但这财富，当它被关闭和保存时，反而摧毁它的拥有者和囤积者。\par

但让我们使这野兽变得温顺。它将变得温顺，如果我们不把它关闭起来，而是把它交到所有需要的人手中。这样，我们将从这一方面获得最大的祝福，既在今世安全地生活并怀有美好的希望，又在将来的日子里坦然无惧地站立。愿我们众人借着恩宠和慈悲等，都能达到此境。\par

\SourceRecord{Translated by Talbot W. Chambers.；Nicene and Post-Nicene Fathers, First Series；Vol. 12.；Edited by Philip Schaff.；Buffalo, NY: Christian Literature Publishing Co.,；1889.；<http://www.newadvent.org/fathers/220135.htm>.}

% structure: p0001 source=h1 target=section reason=page-title

\section{\WorkTitle{格林多前书讲道集 第三十六篇}}

% structure: p0002 source=h2 target=subsection reason=relative-heading-rank; base=h2

\subsection{\BibleRef{格前 14:20}}

\BibleQuote{弟兄们，在心思上不要作小孩子，但在恶事上要作婴孩，在心思上却要作成人。}\par

正如所料，在他长篇论证之后，他采用了更加激烈的语气和大量的责备；并提出了一个适合主题的例子。因为孩子们也常常对琐事垂涎并惊慌，但对非常伟大的事物却没有如此钦佩。既然这些人也拥有说\OriginalTerm{舌音}{tongues}的恩赐，这是所有恩赐中最卑微的，却以为他们拥有了全部；因此他说：\BibleQuote{不要作小孩子}，即：在你们应当深思熟虑的地方不要没有理解，但在不义、虚荣、骄傲的地方，你们要作孩子般单纯。因为在邪恶上作婴孩的人也应当有智慧。因为智慧与邪恶并存就不是智慧，同样，单纯与愚昧并存也不是单纯，必须在单纯中避免愚昧，在智慧中避免邪恶。因为正如苦药和甜药过量都无益，单纯本身和智慧本身也是如此：这就是为什么基督命令我们将两者混合，说：\BibleQuote{你们要机警如同蛇，纯朴如同鸽子。}\BibleRef{玛 10:16}\par

但在邪恶上作婴孩是什么意思？甚至不知道邪恶是什么：是的，他愿意他们如此。因此他也说：\BibleQuote{听说你们中间有\OriginalTerm{奸淫}{fornication}的事。}\BibleRef{格前 5:1}他没有说\BibleQuote{行了}，而是说\BibleQuote{听说}：好像他说：\Quote{你们不是不知道这件事；你们已经听说过一段时间了。}我说，他愿意他们既是成人又是孩子；一个在邪恶上，另一个在智慧上。因为如此，男人甚至可能成为男人，如果他也是一个孩子：但只要他不在邪恶上作孩子，他就不会是男人。因为邪恶的人，不是成熟，而是愚昧。\par

% structure: p0006 source=h2 target=subsection reason=relative-heading-rank; base=h2

\subsection{\BibleRef{格前 14:21}}

\BibleQuote{在法律上记载：我要用异语之舌、外邦之唇，向这人民说话；即便如此，他们仍不听从我，主说。}\par

然而，法律中确实没有记载，但正如我先前所说，他总称全部旧约为\Quote{法律}：包括先知书和历史书。他引用了\Person{依撒意亚}先知的见证，又一次为了他们的益处而暗地贬低这恩赐的光荣；尽管如此，他仍用赞美的方式陈述。因为\Quote{即便如此}这一措辞有力地指出这奇迹足以使他们惊异；如果他们不信，错在他们。天主为何行奇事，如果他们不信？为的是他在每种情况下都显示出尽了自己的本分。\par

2. 然后，他从预言中甚至已经表明，所讨论的标记不是很有用，他补充说：\par

% structure: p0010 source=h2 target=subsection reason=relative-heading-rank; base=h2

\subsection{\BibleRef{格前 14:22}}

\BibleQuote{因此，舌音是为作\OriginalTerm{标记}{sign}，不是为相信的人，而是为不信的人；但说预言是为作标记，不是为不信的人，而是为相信的人。}\par

% structure: p0012 source=h2 target=subsection reason=relative-heading-rank; base=h2

\subsection{\BibleRef{格前 14:23}}

\BibleQuote{因此，如果整个教会聚集在一起，众人都说舌音，有不懂的或不信的人进来，他们岂不说你们疯了？}\par

% structure: p0014 source=h2 target=subsection reason=relative-heading-rank; base=h2

\subsection{\BibleRef{格前 14:24}}

\BibleQuote{但如果众人都说预言，有一个不信的或不懂的人进来，他就会被众人责备，被众人审断：}\par

% structure: p0016 source=h2 target=subsection reason=relative-heading-rank; base=h2

\subsection{\BibleRef{格前 14:25}}

\BibleQuote{这样，他心里的秘密就显露出来；于是他就会俯伏在地，朝拜天主，宣告说：天主真在你们中间。}\par

在这里，人们似乎从所说的内容中发现了巨大的困难。因为如果舌音是为不信的人作标记，他怎么说，如果不信的人看见你们说舌音，他们会说\Quote{你们疯了}？如果预言\Quote{不是为不信的人，而是为相信的人}，不信的人又怎能受益？\par

\Quote{因为如果有不信的人进来，}他说，\Quote{当你们说预言时，他就会被众人责备，被众人审断。}\par

而且不仅如此，在此之后又产生了另一个问题：因为舌音反而显得比预言更大。因为如果舌音是为不信的人作标记，而预言是为相信的人，那么吸引陌生人并使之成为家人，比管理家人更大。那么，这句话的意思是什么？没有什么困难或晦涩，也不是与前面矛盾的，而是非常一致，如果我们留意的话：即预言适合两者，而舌音则不然。因此，当他提到舌音时说\Quote{它是作标记}，他补充说\Quote{不是为相信的人，而是为不信的人}，并为他们作\Quote{标记}，即为了惊奇，而非为了教导。\par

\Quote{但在预言方面，}有人说，\Quote{他也做了同样的事，说：\InnerQuote{但预言不是为不信的人，而是为相信的人。}因为相信的人不需要看见标记，只需要教导和讲解。那么，你怎么说预言对两者都有用，当\Person{保禄}说\InnerQuote{不是为不信的人，而是为相信的人}？}如果你们仔细察验，就会明白所说的。因为他不是说\Quote{预言对不信的人没有用}，而是\Quote{不是为作标记}，如同舌音，即一个没有益处的单纯标记：舌音对相信的人也毫无用处；因为它唯一的工作是使人惊奇和混乱；\Quote{标记}这个词是可能有两种含义的词之一：正如当他说\BibleQuote{求你给我一个标记}\BibleRef{咏 85:17}并加上\Quote{为善}；又说\BibleQuote{我成了许多人惊异的事}\BibleRef{咏 70:7}，即一个标记。\par

为了向你说明，他在这里引入\Quote{标记}一词，并非指它必定产生某种益处，他还加上了由此而来的结果。结果是什么呢？他说：\Quote{他们会说，你疯了。}这并非出于标记的本性，而是出于他们的愚昧。但当你听到不信者时，不要以为每次所指的都是同一类人；有时指那些顽疾不治、固执不改的人，有时指那些可能改变的人，如宗徒时代那些听闻天主的奇事而惊叹的人，如科尔乃略的例子。他的意思因此是这样：先知之恩无论对于不信者还是信者都有益处；至于说舌音之恩，当不信者和轻率者听到时，非但不能从中获益，反而讥讽说话的人，好像他们是疯子。事实上，舌音之恩对他们来说只是一个标记，即仅仅为了使他们惊奇；而明白的人也能从中获益，标记正是为此而赐下的。正如当时不仅有人指责他们喝醉了，也有许多人惊叹他们讲述天主的奇事。可见，嘲笑者是那些没有见识的人。因此，保禄并非简单地说：\Quote{他们会说，你疯了}，而是加上了：\Quote{无学问的和不信者}。\par

但先知之恩不仅仅是标记，它对于信德和裨益，对两类人都适用且有用。这一点，即使没有直接说明，保禄在后面也更清楚地解释了，他说：\Quote{他被众人所责。因为如果大家都讲先知之恩，有一个不信的或无学问的人进来，他就要被众人责，被众人审；他心中的隐秘就这样显露出来，于是他必将面伏于地，朝拜天主，声称天主真在你们中间。}\par

因此，先知之恩的伟大不仅在于它对每一类人都有用，还在于它能吸引不信者中那些更无耻的人。因为伯多禄揭穿撒斐拉时，那是先知之恩的作为，与他讲舌音之恩时的情形并不相同；在前一种情形中，众人无不退缩；而当他讲舌音之恩时，他甚至被视为疯狂。\par

3. 说了舌音之恩没有益处，并又通过将责难转向犹太人而缓和了这一说法之后，他接着指出它甚至会造成伤害。\Quote{那么它为什么被赐下呢？}为的是要与翻译一同发出；因为没有翻译，它在无知识的人中间甚至会产生相反的效果。他说：\Quote{如果大家都讲舌音之恩，有一个不信者或无知者进来，他们会说，你疯了。}正如宗徒们也曾被怀疑是喝醉了：因为经上说：\BibleQuote{这些人喝醉了新酒。}\BibleRef{宗 2:13}但这并非标记的错，而是他们不熟练的缘故；因此他加了\Quote{无学问的和不信者}，以表明这想法源于他们的无知和缺乏信德。因为正如我以前所说，他的目的是将这个恩赐归入并非被贬低、而是没有太大益处的恩赐之列，为要抑制他们，使他们不得不寻求翻译者。因为既然大部分人不是着眼于此，而是利用它来炫耀和争竞，他特别就是要把他们从这种倾向中拉回来，暗示他们的名声受损，他们为自己招致了疯狂的嫌疑。这正是保禄一直试图确立的原则，当他想要引导人远离某事物时，他就指出那人会在他所渴望的那些事物方面遭受损失。\par

你也要如此：如果你想引导人远离快乐，就指出那事物是苦涩的；如果你想引他们远离虚荣，就指出那事物充满耻辱；保禄常常这样做。当他想使富人摆脱贪财之心时，他不仅说财富是有害的事物，还指出它会引人陷入诱惑。他说：\BibleQuote{那些想要发财的人，就陷入诱惑。}\BibleRef{弟前 6:9}这样，因为财富似乎能使人脱离诱惑，他却赋予它富人所以为的相反特质。另一些人固守外邦的智慧，似乎用它来支持基督的教义；他指出它不但不能助益十字架，反而使它落空。他们坚持在外人面前诉讼，认为受自己人审判不合适，好像外人更聪明；他指明在外人面前诉讼是可耻的。他们坚持吃祭偶像之物，好像显示完全的知识；他暗示这是不完全知识的标志，即不知道如何处理与邻人有关的事。同样，在这里，因为他们出于爱虚荣而对舌音之恩狂热，他指明这个恩赐反而比什么都使他们蒙羞，不仅剥夺了他们的荣耀，还使他们陷入疯狂的嫌疑。但他没有立即说出这一点，而是在说了许多话之后，当他的言论变得可接受时，才引出这个与他们意见如此相悖的话题。实际上，这不过是普通的法则：一个想要彻底动摇根深蒂固的见解、并使人们转向相反观点的人，不能立即陈述相反的观点；否则，他在那些受相反信念先入为主的人眼中就显得可笑。因为过于出乎意料的事物一开始并不容易被人接受，你必须先用其他论点加以充分铺垫，然后才让它转向相反的方向。\par

例如，当他谈论婚姻时就是这样做的：我的意思是，既然许多人认为婚姻带来安逸，而他想暗示不结婚才是安逸；如果他立刻这样说，就不容易让人接受；但现在，他在讲了许多其他内容之后，恰当地安排了这个话题，就深深地打动了听众。他在守贞一事上也是如此。因为在这之前他说了很多，在这之后又说了，最后他说：\Quote{我宽容你们}和\Quote{我愿意你们无挂虑}。见：\BibleRef{格前7:28}\par

他在论及语言之恩时也是如此，表明语言之恩不仅剥夺荣耀，还会使拥有者在不信者眼中蒙羞。相反，先知之恩在不信者中既不受责备，又具有极大的信誉和用处。因为对于说预言，没有人会说\Quote{他们疯了；}也没有人会嘲笑那些说预言的人；相反，他们会感到惊奇并钦佩他们。因为\Quote{他被众人所谴责，}即他心中的事被揭露并显示给众人：现在，一个人进来看到一个人说波斯语，另一个人说叙利亚语，与进来听到自己心中的秘密，是不同的；无论是作为试探者心怀恶意而来，还是真诚而来；或者他做了这样那样的事，设计了那样的事。因为这事比另一件事更可畏、更有益。因此，关于语言之恩，他说\Quote{你们疯了；}然而他并非自己肯定这一点，而是依据他们的判断：即\Quote{他们会说，}他说，\Quote{你们疯了；}在这里，相反地，他既使用事实的判断，也使用那些受益者的判断。\Quote{因为他是被众人所谴责的，}他说，他被众人所判断；这样他心中的秘密就被揭示出来；于是他将俯伏在地敬拜天主，宣称天主确实在你们中间。你看见这事不容两种解释了吗？在前一种情况下，所做的事可能被怀疑，而不信者可能会将其归为疯狂；而在这里却没有这样的事，他反而会既惊奇又敬拜，首先通过行为来宣认，然后也通过言语。同样，\Person{拿步高}也敬拜了天主，说：\BibleQuote{的确，你们的天主是那启示奥秘的天主，因为你能揭示这个奥秘。}见：\BibleRef{达2:47}你看到先知之恩的能力了吗？它如何改变了那个野蛮人，使他受教并引入信德？\par

% structure: p0029 source=h2 target=subsection reason=relative-heading-rank; base=h2

\subsection{格林多前书14:26}

\BibleQuote{那么，弟兄们，该怎么办呢？当你们聚集时，每人有诗歌，有教导，有语言，有启示，有解释。凡事都当为建立德行而行。}\par

你看见基督徒的根基和规矩了吗？正如工匠的工作是建造，基督徒的工作就是在一切事上使邻人受益。\par

但是，因为他曾激烈地贬低这个恩赐，免得它显得多余；为了打击他们的骄傲，他这样做：——他又将它与其他恩赐并列，说：\Quote{有诗歌，有教导，有语言。}因为从前他们也凭恩赐创作诗歌，凭恩赐教导。尽管如此，\Quote{让这一切都朝向同一件事，}他说，就是纠正邻人：不要随意行事。因为如果你不来建立你的弟兄，你究竟为什么来这里呢？事实上，我并不太在意恩赐的差别。有一件事令我关注，一件事是我的愿望，就是凡事都\Quote{为建立德行}而行。因此，拥有较小恩赐的人，如果这没有缺失，将会超越拥有更大恩赐的人。是的，恩赐正是为此赋予的，使每个人都能被建立；因为除非这发生，恩赐反而会变成领受者的定罪。请告诉我，说预言有什么用处？复活死人有什么用处，如果没有人从中受益？但如果这是恩赐的目的，并且有可能通过没有恩赐的其他方式实现它，那么不要因那些征兆而自夸，也不要因恩赐被剥夺而自怜。\par

% structure: p0033 source=h2 target=subsection reason=relative-heading-rank; base=h2

\subsection{格林多前书14:27}

\BibleQuote{若有人说语言之恩，只可由两人，或最多三人，且要轮流；并要有一人解释。}\par

% structure: p0035 source=h2 target=subsection reason=relative-heading-rank; base=h2

\subsection{格林多前书14:28}

\BibleQuote{但若没有解释者，就应在教会中缄默，只对自己和天主说话。}\par

你说什么，请告诉我？既然讲了这么多关于舌音的事，说那恩赐若没有翻译者，便是无益的、多余的事，你还命令再说舌音吗？他说：我不命令，也不禁止；就如他说：\BibleQuote{若有不信者请你们赴席，你们愿意去，}他说这话，并不是为他们去立一条法律，而是不拦阻他们：这里也是如此。\BibleQuote{让他只对自己和天主说话。}如果他不能保持沉默，他说，而是这样好胜虚荣，\BibleQuote{让他独自说话。}如此，仅因许可这件事，他就大大地制止了他们，使他们羞愧。他在别处也这样做，论及与妻子的交往时说道：\BibleQuote{但我说这话是因为你们的无节制。}但他论及预言时，却不是这样说的。那是怎样说的呢？是以命令和立法的口吻：\BibleQuote{让先知们两个人或三个人说话。}他在这里毫不寻求翻译者，也不阻止那说预言的人的口，如在前一条中那样说：\BibleQuote{若没有翻译者，叫他缄默；}因为事实上，说舌音的人自己并不够。因此，若有人兼有两种恩赐，就让他说话。但若他没有，却仍想说话，就让他借助翻译者来说。因为先知是翻译者，却是天主的；而你是人的翻译者。\BibleQuote{但若没有翻译者，叫他缄默：}因为不应做任何多余的事，任何出于虚荣的事。只有\BibleQuote{让他只对自己和天主说话；}即，在心里，或安静无声地：至少，如果他要说的话。因为这肯定不是立法的口吻，而是可能通过许可更加羞辱他们；就如他说：\BibleQuote{但若有人饥饿，可在家里先吃，}似乎给予许可，却因此更尖锐地触动他们。\BibleQuote{因为你们聚会不是为此，}他说，\BibleQuote{要显示你们有恩赐，而是为建立听者；}这也是他之前所说的：\BibleQuote{一切事都当为建立而行。}\par

% structure: p0038 source=h2 target=subsection reason=relative-heading-rank; base=h2

\subsection{格前 14:29}

6. \BibleQuote{让先知们两个人或三个人说话，并让其余的人辨别。}\par

他从未加上\Quote{最多}一词，如同在舌音的例子中那样。有人问：这是怎么回事？因为他认为预言本身也不够，如果至少他将判断权交给别人。不，肯定足够；这就是为什么他没有阻止先知的口，如对另一个人那样，当没有翻译者时；也没有像在他（说舌音者）的情况下他说：\BibleQuote{若没有翻译者，叫他缄默，}同样在先知的情况下，他说：\Quote{若没有辨别者，就不要让他说预言；}但他只保护了听者；因为为了听者的满意，他说了这话，免得以占卜者身份混入他们中间。因为这事他在开头也已叫他们提防，当他引入占卜与预言的区别时，现在他叫他们辨别并侦察这事，使任何撒殚的教师不能秘密潜入。\par

% structure: p0041 source=h2 target=subsection reason=relative-heading-rank; base=h2

\subsection{格前 14:30}

\BibleQuote{但若启示临到另一个坐着的人，先说话的人就当缄默。}\par

% structure: p0043 source=h2 target=subsection reason=relative-heading-rank; base=h2

\subsection{格前 14:31}

\BibleQuote{因为你们众人都能一个一个地讲预言，使众人学习，使众人得到安慰。}\par

这是所说的什么？他说：\BibleQuote{若你讲预言时，正在说话，另一个人受感动，你从此就缄默。}因为他在舌音的情况下所说的，这里也要求，即应\Quote{轮流}，只是在这里方式更属预言性质。因为他没有用同样的措辞\Quote{轮流吗？}，而是\BibleQuote{若启示临到另一个人。}因为还有什么必要，当第二个人受感动要讲预言时，第一个人还在说话呢？难道他们两人都该说？不，这是渎神的，会引起混乱。难道该让第一个人说？这也是不合宜的。因为当一个人在说话时，圣神感动另一个人，正是为了让他也能说些什么。\par

所以，为了安慰那个被制止的人，他说：\BibleQuote{因为你们众人都能一个一个地讲预言，使众人学习，使众人得到安慰。}你看见他如何再次陈述他做一切事的原因吗？因为如果说舌音的人若没有翻译者，他完全禁止他说话，因为无益；那么，如果预言没有这种品质，反而造成混乱、扰乱和不合时宜的喧闹，他也有理由命令约束预言。\par

% structure: p0047 source=h2 target=subsection reason=relative-heading-rank; base=h2

\subsection{格前 14:32}

\BibleQuote{先知的神是服从先知的。}\par

你看见他如何认真而可畏地使他羞愧吗？为使那人既不争斗也不结党，他表明恩赐本身是受管辖的。因为在这里，他用\Quote{神}意指其实际运作。但若神是服从的，那么它的拥有者更不能有理由争辩。\par

7. 然后他表示这事也为天主所悦纳，接着说道：\par

% structure: p0051 source=h2 target=subsection reason=relative-heading-rank; base=h2

\subsection{格前 14:33}

\BibleQuote{因为天主不是混乱的天主，而是和平的天主，正如我在众圣徒的教会中一样。}\par

你看见他用了多少理由来使他缄默并安慰他，同时又在向另一个人让步？第一点也是主要的一点，就是他不会因这样的做法而被阻止；他说：\BibleQuote{因为你们众人都能讲预言，} \BibleQuote{一个一个地。}第二点，这件事在圣神看来是好的；\BibleQuote{因为先知的神是服从先知的。}除此以外，这事符合天主的心意；他说：\BibleQuote{因为天主不是混乱的天主，而是和平的天主；}第四点，这习惯在全世界都通行，没有奇怪的事加给他们。因为他说：\BibleQuote{我在众圣徒的教会中这样教导。}\par

现在还有什么比这些更可怕的呢？因为事实上，那时的教会确实是天堂，圣神治理一切，推动每位首领，使他受到感动。但现在我们只保留了那些恩赐的标记。因为现在我们也是两三个人轮流发言，当一人沉默时，另一人开始。但这只是那些事物的迹象和纪念。因此，当我们开始发言时，百姓回应说：\Quote{与你的圣神，}表明他们从前就是这样说话的，不是出于自己的智慧，而是受圣神感动。但现在却不是这样了（我至今只谈我自己的情况）。然而，现在的教会就像一个从昔日繁荣中堕落的妇女，在许多方面只保留了那古老繁荣的标记；她展示着她金饰的匣子和箱子，却失去了财富。现今的教会正像这样。我这样说不是指恩赐：因为如果只是这样，那并不算稀奇；而是指生活和德行。因此，当时她的寡妇名单和童贞女的歌队给教会带来了极大的装饰；但现在她变得荒凉空虚，只剩下标记。现在确实有寡妇，也有童贞女，但她们没有保持那些准备参加这种竞赛的妇女应有的装饰。因为童贞女的特殊标记是只关心天主的事，毫无分心地服侍祂；而寡妇的标记也不仅仅是不要再婚，而是其他方面：对穷人的爱德、款待、恒心祈祷，以及保禄在写给弟茂德的书中严格要求的其他一切。我们也可以看到已婚妇女在我们中间表现出极大的端庄。但这不是唯一的要求，而是那种对穷人的勤恳关注，通过这个，那些古代的妇女曾最光辉地闪耀。不像现今一般情况。因为那时，她们不用金子，而是用施舍的美丽衣裳来装饰自己；但现在，她们放弃了这些，反而被自己罪链编织的金线从四面八方装扮起来。\par

我还要再提一个也耗尽了继承光彩的库房吗？从前他们全都聚在一起，共同咏唱圣咏。现在我们也是这样做的；但那时，所有人都是同心合意；而现在，连在一个灵魂里也看不到那种同心，反而处处都是极大的争战。\par

如今，在教会中主持的人，仿佛进入圣父的房屋，也为众人祈求说：\Quote{平安，给众人}；但这平安的名字常被提及，而事实却无处可见。\par

8. 那时，房屋本身就是教堂；但现在，教堂本身成了一所房屋，甚至比任何房屋更糟。因为在一所房屋里，可以看到很好的秩序：主母端庄地坐在她的椅子上，少女们安静地纺织，每个仆人都手中有指定的工作。但在这里，却是极大的喧哗、极大的混乱，我们的聚会与酒店毫无区别，笑声如此响亮，骚动如此之大；如同在浴室、在市场上，喊声和嘈杂声无处不在。而这些东西只有在这里：因为在其他地方，即使在教堂里也不允许与邻人交谈，即使你接回一位久别的朋友也不可，但这些事是在外面做的，而且非常恰当。因为教堂不是理发店或香水店，也不是市场里任何商人的仓库，而是天使的地方、总领天使的地方、天主的宫殿、天堂本身。因此，如果有人劈开了天堂，带你进去，即使你看见你的父亲或兄弟，你也不敢说话；同样，在这里，除了这些属神的声音，也不应发出任何其他声音。因为，事实上，这个地方的事物也是一个天堂。\par

如果你不信，就看看这祭台，想一想它是为谁而设，以及为何而设：想一想谁将要在这里显现；甚至在那一刻之前就要战栗敬畏。因为，当一个人只看见君王的宝座时，他心里就起身，期待君王的显现。因此，你也要在那令人战栗的时刻之前就感到战栗：提升你自己，甚至在看见帷幕拉开、天使歌队前行之前，你就要升到天堂本身。\par

但未入门者不明白这些事。那么，为了他，也有必要引入其他话题。因为对于他，我们也不缺少能够彻底激励他、使他翱翔的理由。\par

那么，你这位不明白这些事的人，当你听到先知说：\Quote{主这样说}时，就要离开大地，也升到天堂，想一想是谁藉着他与你说话。\par

但事实上，为了一个引人发笑的小丑，或为了一个淫荡堕落的女人，竟有如此庞大的一群观众聚集起来，完全安静地聆听所说的内容，而且没有人命令安静；既没有喧哗，也没有喊叫，更没有一丝噪音；但当天主从天上就如此可畏的主题说话时，我们却比狗还要无礼，甚至对娼妓的尊敬超过对天主的尊敬。\par

听到这些事令你毛骨悚然吗？不，那么，当你们做这些事时，更应毛骨悚然。\par

9. 保禄论到那些轻看穷人、独自宴乐的人所说的话：\BibleQuote{怎么，你们没有房屋可以吃喝吗？还是你们轻视天主的教会，羞辱那些没有的人呢？}\BibleRef{格前 11:22}——请允许我也同样对那些在此处制造骚动和交谈的人说：\Quote{什么？你们没有房屋可以闲谈吗？还是你们轻视天主的教会，甚至败坏那些本来端庄安静的人呢？}\Quote{但是与朋友交谈是甜蜜愉快的。}我不禁止这个，但应在家中、市场、浴场去做。因为教会不是交谈的地方，而是施教的地方。但现在它与市场无异；不，若说得不过分，恐怕也与剧场无异；来此聚会的妇女们打扮得比那里的淫荡者更加放荡。因此我们看到许多放荡之徒被她们引诱到这里；若有人试图或有意败坏一个妇女，我想没有比教会更合适的地方了。若有什么买卖之事，教会似乎比市场更方便。因为在这里谈论这类话题比在店铺本身还多。若有人想说或听什么闲话，你会发现这里也比外面的广场更多。若你想听什么政治的事、私人家务或军营的事，不要去审判厅，也不要坐在药铺里；因为这里，我说这里，有人更准确地报告这一切；我们的集会什么都是，就不是教会。\par

我是否刺痛了你们？我本人认为没有。因为你们继续同样的行为，我怎能知道你们被所说的话触动了呢？因此我必须再次谈论同样的话题。\par

这些事难道还能容忍吗？这些事还能忍受吗？我们每日使自己疲倦分心，好使你们不至于没有学到一些有益的事就离开：但你们中没有人变得更好，反而更受伤害。是的，并且\Quote{你们聚在一起成为审判，}不再有遮盖你们罪的借口，并且你们赶出更谦逊的人，用你们的愚行从四面扰乱他们。\par

但众人都说什么？有人\Quote{我听不见读的是什么，}又说\Quote{也不知道所说的话是什么。}因为你们制造骚动和混乱，因为你们没有带着虔诚的灵魂来。你说什么？\Quote{我不知道所说的事情。}那么，正是为此你应当留心。但若连晦涩都不能激起你的灵魂，那么若事情清楚，你更会匆匆放过。是的，这就是为什么不是所有事情都清楚，免得你纵容怠惰；也不是所有都晦涩，免得你绝望。\par

然而，那个太监和外国人\BibleRef{宗 8:20}没有任何此类抱怨，反而在如此众多重要事务的包围中，且在旅途中，手不释卷地阅读：而你们，既有丰盛的教师，又有别人私下给你们阅读，却向我提出你们的借口和托词？你们不知道所说的话吗？那么为何不求知？但确实不可能一无所知。因为许多事情本身是明显清楚的。再者，即使你们一无所知，也应当安静，不要妨碍那些专心的人；那样，天主因着你们的安静和敬畏，或许会使晦涩的事也明白。但你们不能保持沉默吗？那么，出去吧，不要也成为别人的祸害。\par

因为事实上，在教会中总只应有一个声音，正如只有一个身体一样。因此，读经者独自发出声音，主教也甘愿安静坐着；咏唱者独自咏唱；虽然众人回应，但声音如同从一张嘴发出。讲道者独自讲道。但当有许多人在谈论许多不同的话题时，我们为何要徒劳地扰乱你们？因为除非你们认为我们只是无益地扰乱你们，否则你们就不会在我们谈论如此崇高的事的时候，谈论无关紧要的事情。\par

10. 因此，不仅在你们的品行中，而且在你们对事物的评估中，都有极大的偏差。你们渴望多余之物，舍弃真理，追求各种影子和梦境。现今的一切不都是影子和梦境，甚至比影子更糟吗？因为它们在出现之前就飞逝；在飞逝之前，它们带来的烦恼很多，超过快乐。让一个人在今世获得大量财富并埋藏在地下，然而当夜晚过去，他必赤身离去，这并不奇怪。因为那些只在梦中富有的人，从床榻上起来时，一无所有，好像睡时所有的。贪得无厌的人也是如此；或者不如说，不是如此，而是更糟的处境。因为那梦见自己富有的人，既没有他幻想的金钱，醒来时也没有发现这幻想给他带来任何别的损害；但这个人既被剥夺了财富，又必须离去，满身是从财富中产生的罪恶；他在财富中只享受了幻想，但财富带来的恶果他不再在幻想中看见，而是在事情的真相中；他的快乐是在梦中，但快乐随之而来的惩罚不再是梦，而是实际的经历。是的，甚至在惩罚之前，就在此地，他受到最重的惩罚，在聚敛财富的过程中，使自己充满了无数的悲伤、焦虑、控告、诽谤、骚动、烦乱。\par

因此，为使我们能从梦境和并非梦境的邪恶中得以解脱，让我们选择施舍代替贪婪，选择怜悯人代替掠夺。因为这样，我们将获得今生和来世的美物，借着我们的主耶稣基督的恩宠和仁慈，愿光荣、权能、尊崇归于祂，偕同圣父及圣神，自今而始，及于永远，世世无尽。阿们。\par

\SourceRecord{Translated by Talbot W. Chambers.；Nicene and Post-Nicene Fathers, First Series；Vol. 12.；Edited by Philip Schaff.；Buffalo, NY: Christian Literature Publishing Co.,；1889.；<http://www.newadvent.org/fathers/220136.htm>.}

% structure: p0001 source=h1 target=section reason=page-title

\section{讲道集第37篇——论格林多前书}

% structure: p0002 source=h2 target=subsection reason=relative-heading-rank; base=h2

\subsection{格林多前书 14:34}

\Quote{妇女在集会中应当缄默；她们不许发言，但要服从，正如法律所说的。}\par

在平息了因说语言和说先知话引起的混乱，并制定了避免混乱的法则，即说语言的人应轮流进行，而说先知话的人应在另一个人开始时静默之后，他接着处理由妇女引起的无序，制止她们不合时宜的放肆言论；而且时机非常恰当。因为即使对那些有神恩的人，也不允许他们轻率发言，也不允许他们随心所欲地发言，尽管他们是受圣神感动；更何况那些闲谈无益的妇女呢？因此，他带着很大的权威压制她们的喋喋不休，并援引法律，以此封住她们的口；不仅在这里劝诫或提出建议，甚至通过援引古代关于该主题的法律，强烈地命令她们。因为他说了\Quote{妇女在集会中应当缄默；}和\Quote{不许她们发言，但要服从}之后，他加了\Quote{正如法律所说的。}这法律在哪里说呢？\Quote{妳的愿望应归于你的丈夫，他要管辖妳。}\BibleRef{创 3:16}你看到保禄的智慧了吗？他引用了什么样的见证，不仅命令她们沉默，而且是带着畏惧的沉默；这种畏惧之大，如同一个使女应当保持安静一样。因此，他自己说了\Quote{不许她们发言}之后，没有加\Quote{但要缄默}，而是用更重的\Quote{服从}代替了\Quote{缄默}。如果对丈夫尚且如此，更何况对教师、父亲和整个教会的集会呢。\Quote{但如果她们连发言也不允许，}有人说，\Quote{甚至不能提问，那她们出席集会有什么用？}她们应当听该听的；至于疑问点，让她们在家里从丈夫那里学习。因此他又加了：\par

% structure: p0005 source=h2 target=subsection reason=relative-heading-rank; base=h2

\subsection{格林多前书 14:35}

\Quote{如果她们要学习什么，可以在家里问自己的丈夫。}\par

这样，\Quote{不仅，看来，不允许她们随意发言，}他说，\Quote{甚至也不允许在集会中提问。}如果她们不应该提问，那么随意发言就更违背规定了。他为什么把她们置于如此大的服从之下呢？因为女人在某种程度上是较软弱的受造物，容易动摇和轻浮。你看，他为何设置她们的丈夫为她们的教师，这对双方都有益处。因为这样他既使妇女有秩序，又使丈夫们认真，因为他们必须非常精确地把所听到的传授给妻子。\par

此外，因为他们以为公开说话是她们的装饰，他又把话题转到相反的一面，说：\Quote{因为妇女在集会中发言是可耻的。}也就是说，他首先从天主的法律中证明了这一点，然后从普遍的理性和我们接受的习惯来证明；甚至当他在与妇女谈论长发时，他说：\Quote{不是本性也教导你们吗？}\BibleRef{格前 11:14}处处你都可以发现这是他的风格，不仅从圣经中，也从普遍习惯中，来使他们羞愧。\par

2. 除此之外，他也通过大家一致同意和到处规定的考虑来使他们羞愧；这个论题他也在这里提出，说：\par

% structure: p0010 source=h2 target=subsection reason=relative-heading-rank; base=h2

\subsection{格林多前书 14:36}

\Quote{什么？难道天主的话是从你们出来的？或是只有你们才领受了？}\par

这样，他引进了其他也持守这法律的教会，既通过事物的新颖性来平息混乱，又通过普遍的声音使他的话易于接受。因此，他在别处也说：\Quote{谁要使你们记念我在基督内的行径，就如我在各处各教会中所教导的。}\BibleRef{格前 4:17}又说：\Quote{因为天主不是混乱的天主，而是平安的天主，如同在各圣人们的教会一样。}\BibleRef{格前 14:33}而在这里：\Quote{什么？难道天主的话是从你们出来的？或是只有你们才领受了？}即\Quote{你们既不是首先，也不是唯一相信的人，而是全世界。}这正如他写信给哥罗森人时所说的：\Quote{这福音在全世界结果实并增长，}\BibleRef{哥 1:6}谈论福音。\par

但另有些时候，他也把它转向鼓励听众；例如当他说他们的初果并显于众人时。这样，写信给得撒洛尼人时他说：\Quote{因为天主的话从你们那里传扬出来，}\Quote{你们对天主的信德传遍各处。}\BibleRef{得前 1:8}又写信给罗马人时：\Quote{你们的信德传遍了全世界。}因为两者都能使人羞愧并激励人，既因为受到他人称赞，又因为有别人参与他们的判断。因此，在这里他也说：\Quote{什么？难道天主的话是从你们出来的？或是只有你们才领受了？}\Quote{因为你们不能说，\InnerQuote{我们成了其余人的教师，我们不需要向别人学习；}}也不能说，\Quote{信德只留在这地方，不应接受来自其他地区的先例。}你看，他用多少论据使他们羞愧？他援引法律，指出事情的可耻，提出其他教会。\par

3. 接着，他把最有力的放在最后，说：\Quote{天主甚至现在通过我命令这些事。}\par

% structure: p0015 source=h2 target=subsection reason=relative-heading-rank; base=h2

\subsection{格林多前书 14:37}

这样：\Quote{若有人自以为是一位先知或属神的人，就该承认我所写的给你们的是主的诫命。}\par

% structure: p0017 source=h2 target=subsection reason=relative-heading-rank; base=h2

\subsection{格林多前书 14:38}

\Quote{但若有人无知，就让他无知吧。}\par

他为什么加上这句话？暗示他并非用暴力或争辩，这正是那些不愿建立自己之事、而求他人得益之人的标志。因此，他在另一处也说：\BibleQuote{但如果有人似乎好争辩，我们却没有这样的惯例。}\BibleRef{格前 11:16} 但他并非处处如此，只在过失不太严重时才这样做，且更多是出于使其羞愧。因为当他论及其他罪时，他并不这样说。而是怎样？\BibleQuote{你们不要自欺：行淫者、作娈童的，都不能承受天主的国。}\BibleRef{格前 6:9} 又说：\BibleQuote{看，我保禄告诉你们：如果你们接受割损，基督对你们就毫无益处。}\BibleRef{迦 5:2} 但在这里，由于他的话题是关于沉默，他并没有强烈责备他们，反而以此更吸引他们。然后，按他惯常的作法，他回到先前的话题，即他离开去说这些话的源头，如下所述：\par

% structure: p0020 source=h2 target=subsection reason=relative-heading-rank; base=h2

\subsection{格前 14:39}

\BibleQuote{所以，弟兄们，你们要切慕说预言，也不要禁止说方言。}\par

这也是他的习惯：不仅解决眼前的问题，而且由此出发，纠正任何在他看来与之相关的事物，然后再回到前者，以免显得偏离主题。例如，当他论及他们在宴席上的和谐时，他岔开话题谈到他们在奥迹中的共融，并以此使他们羞愧，然后又回到前题，说：\BibleQuote{所以，你们聚在一起吃饭时，要彼此等待。}\BibleRef{格前 11:33}\par

同样在这里，他论及在恩赐中要有秩序，以及既不可因较小的恩赐而气馁，也不可因较大的而自高；然后从那里岔开去谈论妇女的庄重，并确立这点之后，他回到主题，说：\BibleQuote{所以，弟兄们，你们要切慕说预言，也不要禁止说方言。}你看他如何一直保持这两者的区别？又如何表明一个非常必要，另一个并非如此？因此，关于前者他说\Quote{要切慕，}关于后者他说\Quote{不要禁止。}\par

4. 然后，作为简要总结，他纠正一切，加上这些话：\par

% structure: p0025 source=h2 target=subsection reason=relative-heading-rank; base=h2

\subsection{格前 14:40}

\BibleQuote{凡事都要规规矩矩地按着次序行。}\BibleRef{格前 14:40}\par

他再次打击那些无缘无故行为不当、招致疯狂名声的人，以及那些不守自己的本位的人。因为没有什么比好的秩序、和平、爱更能建立人，正如它们的反面倾向于拆毁。不仅在属灵的事上，在一切其他事上我们也可以观察到这一点。无论是舞蹈、船只、战车、军营，如果你混乱了秩序，把较大的从他们的位置上除去，把较小的带入他们的行列，你就毁坏了一切，事情就颠倒了。因此，我们也不要破坏我们的秩序，不要把头放在下面，把脚放在上面：当我们推倒理性，把我们的情欲、激情和快乐置于理性之上时，就发生了这种情况；从此波涛汹涌，混乱巨大，风暴难忍，一切都陷入黑暗。\par

如果你愿意，让我们首先检查由此产生的不雅，然后检查损失。那么，如何能清楚并完全了解呢？让我们提出一个处于那种状态的人：迷恋一个妓女，被不光彩的情欲征服；然后我们就会看到由此产生的嘲笑。因为有什么比一个男子在妓女的房门前守候、被娼妓殴打、哭泣、哀叹、使他的荣耀变为羞辱更卑贱呢？如果你也愿意看到损失，我恳求你，回忆金钱的花费、极度的风险、与情敌的争斗、在那种争斗中受伤、被鞭打。\par

那些被贪财之欲所捆绑的人也是如此；或者他们行为更不雅。因为前者完全专注于一个人，而贪婪的人却同样忙碌于所有人的财产，无论穷富，渴望不存在的东西；这尤其表明他们情欲的狂野。因为他们说不仅仅是\Quote{我渴望拥有这个或那个人的财产，}而且他们希望山变成金子，房子和所有看见的一切；他们进入另一个世界，这种情欲达到无限的程度，从不停止他们的贪欲。有什么言语能描绘出那些思想的暴风雨、波涛和黑暗？当波涛和风暴如此巨大时，能有什么快乐？一点也没有；有的只是喧闹、痛苦和乌云，不降雨却带来极大的心碎：这种情况通常发生在那些迷恋不属于自己的美色的人身上。因此，那些根本没有迷恋之情的人，比任何情人都更快乐。\par

5. 这一点无人会否认。但依我看来，即使那爱慕却克制自己激情的人，也比他不断享受情妇的人活得更愉快。尽管这证明颇为困难，但即便如此，我们也必须冒险提出这个论点：困难增加的原因不在于事情的本性，而在于缺乏适合听取这种高尚道德的听众。那么，请问：对情人来说，是被所爱者鄙视更愉快，还是被尊敬并俯视她更愉快？显然是后者。那么，告诉我，妓女会更看重谁？是那做她奴隶、已被她意志俘虏的人，还是那高于她的罗网、飞得比她的箭更高的人？每个人都看得出，是后者。她又会更关心谁？是那已堕落的人，还是尚未如此的人？当然是尚未如此的人。谁会更令人渴望？是被征服者，还是尚未被俘者？是至今尚未被俘者。你若不信，我就从你们内心所发生的事来证明。比如：一个男人会更迷恋哪种女人？是轻易顺从、委身于他的，还是拒绝并给他带来麻烦的？显然是后者；因为这样渴望会更强烈地燃烧。当然，在女人身上也会发生完全相同的事。她们会更尊重和钦佩那些俯视她们的人。若果真如此，那么另一方面，更受尊敬和爱慕的人享受更大的快乐，这也是真的。因为将军也放过已被攻取的城市，却全力围困那坚持抵抗的城；猎人在捕获猎物后，把它关在黑暗中，如同妓女对待她的情人，而去追赶那逃避他的猎物。\par

但有人会对我说：\Quote{前者享受了他的欲望，后者却没有。}然而，免于耻辱、免于做她暴虐命令的奴隶、不被她当作苦力牵着走、不被殴打、被吐口水、被头朝下扔——你认为这是微小的快乐吗？请告诉我。不，若有人准确考察这些事，并能将他们的侮辱、抱怨、永无止境的争吵——有些出于脾气，有些出于放纵——他们的敌意，以及所有其他只有亲身感受者才知道的事——集中起来，他会发现没有哪场战争比这种可怜的生活有更多休战。那么，你说的是什么快乐？是交合时短暂而短暂的享受？但这种享受很快就会被争吵、风暴、愤怒以及同样的疯狂所覆盖。\par

6. 我们说这些事，如同与放荡的青年谈话，他们很不耐心地听我们谈论天国和地狱的道理。\par

现在我们也提出这些主题时，甚至无法说出节制者有多大的快乐；若有人在心中描绘他的冠冕、他的赏报、他与天使的交往、他在世界面前的宣扬、他的勇气、那些有福而不朽的希望。\par

\Quote{但交合有一种快乐：}他们不断重复这点；\Quote{而节制者在与本性的暴政搏斗时不断受苦。}不，你会发现恰恰相反的结果。因为这种暴力和骚动更多地存在于不洁者身上：他的身体里有猛烈的风暴，暴风雨中的海也不像他那样痛苦；他从未抵抗过自己的激情，而是不断被它打击；如同附魔者和那些不断被恶灵撕裂的人。而节制者，像一位高贵的斗士，不断打击它，收获最好的快乐，比千万次那种快乐更甜蜜；这胜利和他的良心，以及那些光辉的凯旋，是他不断装饰自己的饰物。\par

至于另一个人，若他在交合后得到片刻喘息，这算不了什么。因为风暴又来了，又起波浪。但自制的人根本不让这种骚动抓住他，不让海兴起，不让野兽咆哮。即使他在克制这种冲动时忍受了一些暴力，但另一个也不断受到打击和刺伤，无法忍受那刺痛；这好比一匹狂野的马奋力挣扎，一人用缰绳勒住它，巧妙地把持；而另一人则松开缰绳以求免于麻烦，却被它拖着到处跑。\par

若我说这些事比适当更直白，愿无人责备我。因为我并不想以言辞的庄重来炫耀，而是要使我的听众庄重。\par

因此，先知们也不回避这样的言辞，希望铲除犹太人的放荡，甚至比我们现在所说的更赤裸裸地斥责他们。因为医生要除去溃疡，并不考虑如何保持双手干净，而是如何使病人摆脱溃疡；那想要高举卑微者的人，首先自己卑微；那想要杀死谋反者的人，也染上和他一样的血，这使他更光彩。因为若有人看见一个士兵从战场归来，满身血污和脑浆，他不会因此厌恶他或避开他，反而会更钦佩他。所以我们也当这样做，当我们看见任何人杀死自己的邪欲归来，满身是血，我们当更钦佩他，参与他的战斗和胜利，并向那些放纵这种疯狂爱情的人说：\Quote{请向我们展示你们从情欲中获得的快乐；因为节制者从他胜利中得着快乐，而你却从任何方面都得不到。但若你提到与犯罪相连的快乐，另一方的快乐更明显更满足。因你从享受中得到短暂而不明显的快乐，而他因自己的良心，享受到更大、持久、更甜的喜乐。难道女人的陪伴能与自制相比？自制能保全灵魂不受搅扰，并给它安上翅膀。}\par

那么：正如我所说，节欲的人就这样明显地让我们看到他的快乐；但在你身上，我却看到因失败而产生的沮丧，但快乐呢，我渴望看到，却找不到。因为你认为快乐的时刻是哪一个？是犯罪行动之前的时刻吗？不，并非如此，因为那是疯狂、谵妄和狂乱的时候：咬牙切齿、失去自己，这并非快乐；即使它是快乐，它也不会在你身上产生与受苦者相同的效果。因为用拳头打人和被打的人都咬牙切齿，而分娩的妇女因痛苦而神志错乱，也是如此。所以，这并非快乐，而是狂乱、混乱和骚动。那么，我们是否可以说行动之后的时刻？不，这也不是。因为我们不能说一个刚分娩的妇女是快乐的，而只是从某些痛苦中解脱。但这绝不是快乐，而是软弱和衰退：这两者之间有巨大的差别。那么，快乐的时刻是什么，请告诉我？没有。但若有，也短暂到不明显。至少，我们多次努力寻找并捕捉它，却未能做到。但贞洁者的快乐不是这样，它更宽广且对所有人显而易见。更确切地说，他全部的生活都在快乐中，他的良心得着冠冕，波浪平息，内心没有来自任何方面的搅动。\par

既然这个人的生活更在于快乐，而沉溺于快乐的生活则充满沮丧和不安；让我们逃离放荡，持守节欲，好使我们也能获得将来的美好事物，赖恩宠和仁慈等等，等等。\par

\SourceRecord{Translated by Talbot W. Chambers.；Nicene and Post-Nicene Fathers, First Series；Vol. 12.；Edited by Philip Schaff.；Buffalo, NY: Christian Literature Publishing Co.,；1889.；<http://www.newadvent.org/fathers/220137.htm>.}

% structure: p0001 source=h1 target=section reason=page-title

\section{论格林多前书第三十八篇讲道}

% structure: p0002 source=h2 target=subsection reason=relative-heading-rank; base=h2

\subsection{格林多前书 15:1-2}

\BibleQuote{弟兄们，我愿意你们认清，我从前所传给你们的天主福音，你们也接受了，在其中站稳了，藉着这福音也得了救——我用什么话语传给你们。}\par

讲完神恩的论题后，他转向那最必需的事，即复活的主题。因为在这事上，他们也大为动摇。正如人的身体，当热病实际侵入其坚硬的部分——我指神经、静脉和主要元素——时，除非悉心照料，否则灾祸便不可收拾；同样，那时也几乎发生同样的事。因为那灾祸正蔓延至虔敬的基本要素。因此，保禄也极其恳切。因为此后他的论述不再关乎道德，或某人行淫，某人贪财，某人蒙头；而是关乎一切美善的总纲。因为关于复活本身，他们意见分歧。因为这是我们的整个希望，魔鬼便对此猛烈攻击：一时完全颠覆它，一时又说它\Quote{已经过去}；保禄在写给弟茂德的信中将此称为坏疽，我指这邪恶的教义，并给那些引人陷于此的人烙上印记，说：\BibleQuote{其中有依默纳约和非肋托，他们在真理上偏离了，说复活已过，因而颠覆了一些人的信德。}\BibleRef{弟后 2:17} 他们一时这样说，另一时又说身体不再复活，但灵魂的净化就是复活。\par

但那邪恶的恶魔怂恿他们这样说，不仅想推翻复活，还想表明为我们所做的一切都是虚构。因为若他们相信身体不复活，他便逐渐使他们相信基督也没有复活。随之他也会适时引入这一点：祂没有来，也没有行祂所行的事。因为魔鬼的诡计就是如此。因此保禄也称之为\Quote{\OriginalTerm{诡计}{cunning craftiness}}，因为他不立即表明要达成的意图，以免被识破，而是戴上一副面具，制造另一类伎俩：如同狡猾的敌人攻击一座有城墙的城市，从地下秘密挖掘，从而难以防备，并达成目的。因此，既然他的这些陷阱不断被揭露，这些狡猾的伏击被这位可敬的勇士追捕，他说：\BibleQuote{因为我们不是不知道他的阴谋。}\BibleRef{格后 2:11} 同样在此，他揭露了其全部诡计，指出所有策略，并详述了保禄极力想成就的一切。是的，因此他在这之后才提出这个要旨，既因为它极为必要，也因为它涉及我们的全部状况。\par

请看他的体谅：他先稳固自己的立场，然后才继续论述，并大大地使外人在无言可辩。他稳固自己的立场，不是靠推理，而是靠已经发生且他们自己已接受并相信的事：这是最能使他们羞愧、并抓住他们的事。因为此后若他们不愿相信，那不再是不信保禄，而是不信他们自己：这正责备那些一旦接受后又改变心意的人。因此他也从这里开始，暗示他无需其他证人来证明他说真话，正是那些被欺骗的人。\par

2. 但为了使我说的话更清楚，我们需要在接下来的经文中关注原话。那么这些是什么？他说：\BibleQuote{弟兄们，我愿意你们认清，我从前所传给你们的天主福音。}你看见他以何等谦逊开始吗？你看见他如何从一开始就指出他并未引入任何新奇或奇怪的事吗？因为那\Quote{认清}已认识但后来被遗忘的事，就是藉回忆使之显明。\par

当他称他们为\Quote{弟兄们}时，即使从这里，他也为其论述的证明奠定了重要基础。因为成爲\Quote{弟兄们}别无他因，只因基督按肉身的奥迹。而这正是他如此称呼他们的原因，同时安抚、劝勉他们，并提醒他们无数的恩惠。\par

接下来的内容也同样证明了这一点。那么这是什么？\Quote{福音}。因为福音的总纲源于此：天主成为人，被钉十字架，并复活。这个福音也是加俾额尔向童贞女宣讲的，也是先知向世界宣讲的，也是所有宗徒宣讲的。\par

\BibleQuote{我从前所传给你们的天主福音，你们也接受了，在其中站稳了，藉着这福音也得了救——我用什么话语传给你们；你们若持守这福音，除非你们徒然相信。}\par

你看见他如何叫他们自己成为所言之事的见证人吗？他并不说：\Quote{你们听见了}，而是说：\Quote{你们接受了}，要求他们视之为某种寄托，表明他们不仅用言语，也用行动、神迹和奇事接受了它，并应妥善保存。\par

接着，因他论及久远之事，他也指向当下，说：\Quote{在其中站稳了}，占据有利地位，使他们即使极想否认也无能为力。这也是为何他在开始时不说\Quote{我教导你们}，而是\Quote{我愿意你们认清}已显明的事。\par

他又怎能说那些被风浪如此摇荡的人\Quote{站着}呢？他假装不知道，为使他们受益；在对迦拉达人的事上他也这样做，但方式不同。因为在那个场合，他无法假装不知，便用另一种方式组织他的言辞，说：\Quote{我在主内信任你们，你们不会有别的想法。}\BibleRef{迦 5:10}他没有说\Quote{你们没有别的想法}，因为他们的过错是公认而明显的，但他为将来担保；然而这也不确定；但这是为了更有效地吸引他们到自己这里来。然而在这里，他确实假装不知道，说：\Quote{你们也站在这福音上。}\par

接着是益处：\Quote{借着这福音你们也得以得救，用我所传给你们的话。}\Quote{那么，现在的这番解释是为了教义的清晰和阐明。因为教义本身你们不需要，}他说，\Quote{去学习，而是要受提醒和改正。}他说这些话，是不给他们留下一次陷入轻率的余地。\par

但\Quote{我用什么话传给你们}是什么意思？\Quote{我说}，他说，\Quote{复活是怎样发生的？因为我相信你们并不怀疑有复活：但你们或许寻求更清楚地了解那句话。那么我将为你们提供这一点：因为我确实确信你们持守这教义。}接着，因为他直接断言\Quote{你们也站在这福音上}，为了不因此使他们更加懈怠，他又警告他们说：\Quote{你们持守这话，除非你们是徒然相信。}暗示打击是在主要头上，争战不是为了寻常之事，而是为了整个信仰。目前他有所保留地说，但随着他继续并变得热烈，他揭开面纱并呼喊说：\Quote{但如果基督没有复活，那么我们的宣讲就是空的，你们的信德也是空的：你们还在你们的罪中。}但开头并非如此：因为这样进行，温和而逐渐地，是适宜的。\par

% structure: p0016 source=h2 target=subsection reason=relative-heading-rank; base=h2

\subsection{格前 15:3}

\Quote{因为我当日所传给你们的，首先就是我所领受的。}\par

这里他也没有说\Quote{我对你们说}，也没有说\Quote{我教导你们}，而是再次使用同样的措辞，说：\Quote{我传给你们的，也是我所领受的。}他在这里也没有说\Quote{我受教}，而是说\Quote{我领受}：确立了这两件事：第一，一个人不应从自己引入任何东西；第二，通过他的行为论证，他们完全信服，而不是通过空言：并且逐渐地，当他使他的论证可信时，他将一切归功于基督，并表明在这些教义中没有任何人的成分。\par

但这是什么意思：\Quote{因为我当日所传给你们的，首先就是}？因为那是他的话。\Quote{起初，不是现在。}这样说，他把时间当作见证，并表明那些如此长时间被说服的人如今改变主意，是最大的耻辱：不仅如此，而且也表明这教义是必要的。因此它也是\Quote{传}在\Quote{首先}，并且从起初就立即传了。那么你传了什么？告诉我。但他没有立即这样说，而是先说：\Quote{我领受了。}那么你领受了什么？\Quote{基督为我们的罪死了。}他没有立即说我们身体的复活，然而实际上他确实建立了这一点，但通过其他话题从远处进行，说\Quote{基督死了}，并奠定了关于复活教义的坚固基础和不可反驳的根基。因为他不是简单地说\Quote{基督死了}；尽管这本身就足以宣告复活，但他加上一句：\Quote{基督为我们的罪死了。}\par

3. 但首先值得听听那些受\OriginalTerm{摩尼教}{Manichaean}教义感染的人在这里说些什么，他们既是真理的敌人，也反对自己的救恩。那么这些人说什么呢？他们说，这里的死亡，保禄无非指我们在罪中的状态；而复活则指我们从罪中获得释放。你看到没有比错误更软弱的东西吗？它如何被自己的翅膀所捕获，不需要外来的争战，而是自行刺穿自己？例如，考虑这些人，他们如何也被自己的陈述刺穿了自己。因为如果这是死亡，并且基督没有取身体，如你们所假设，然而他死了，那么按照你们，他就在罪中。因为我确实说，他取了身体，并且他的死亡，我说，是肉身的死亡；但你否认这一点，将被迫肯定另一件事。但如果他在罪中，他怎么说：\Quote{你们中间谁能指证我有罪？}以及\Quote{这世界的统治者来了，他在我内毫无所有？}\BibleRef{若 8:46}并且又说：\Quote{我们理当这样履行完所有的义德？}\BibleRef{玛 3:15}而且，如果他本身在罪中，他怎能为了罪人而死？因为那为罪人而死的人，自己应当无罪。因为如果他自己也犯罪，他怎能为了其他罪人而死？但如果他为别人的罪而死，他死了，他是无罪的：如果他是无罪的而死了，他死了——不是罪的死亡；因为他既是无罪的，怎能呢？——而是肉身的死亡。因此保禄并不是简单地说\Quote{他死了}，而是加上\Quote{为我们的罪}：既强迫这些异端者违背他们的意志承认他身体的死亡，也借此表明在死亡之前他是无罪的：因为那为别人的罪而死的人，按必然之理，自己应当是无罪的。\par

他也不满足于此，反而加上\Quote{照经上所记载的：}他借此既再次使论证可靠，又暗示他所说的是何种死亡：因为圣经处处宣扬的是身体的死亡。因为，\BibleQuote{他们刺透了我的手和我的脚，}祂说，以及\BibleQuote{他们要瞻望他们所刺透的那一位。}\BibleRef{咏 20:18}\BibleRef{若 19:37}还有许多其他例证，不必一一列举，部分以言词、部分以预象，可见于经卷中，彰显祂在肉身中被宰杀，且祂为我们的罪而被杀。因为，\Quote{因我百姓的罪恶，}一人说，\Quote{祂走向死亡：}以及\Quote{主为我们的罪恶交付了祂：}以及\Quote{祂因我们的悖逆而受创。}\EditorialNote{依撒意亚 53}但若你不忍受旧约，请听若望呼喊并宣告二者，即祂在肉身中的宰杀及其原因：他说，\BibleQuote{看，天主的羔羊，除免世罪者！}\BibleRef{若 1:29}以及保禄说：\BibleQuote{祂使那不知罪的，替我们成了罪，好叫我们在祂内成为天主的正义。}\BibleRef{格后 5:21}以及再次：\BibleQuote{基督由法律的咒骂中赎出了我们，为我们成了可咒骂的。}\BibleRef{迦 3:13}以及再次：\BibleQuote{解除了率领者和掌权者，公开羞辱了他们，战胜了他们；}\BibleRef{哥 2:15}以及成千上万的其他说法，表明祂在肉身死亡时所发生的事，以及因我们的罪所发生的事。是的，基督自己说：\BibleQuote{我为他们祝圣我自己}以及\BibleQuote{今世的首领已被定罪；}表明祂虽无罪却被杀。\par

% structure: p0022 source=h2 target=subsection reason=relative-heading-rank; base=h2

\subsection{格林多前书 15:4}

4. \Quote{并且祂被埋葬。}\par

这也证实了前文，因为被埋葬的无疑是一个身体。这里他不再加上\Quote{照经上所记载的。}他原本可以，却没有加。原因何在？或因埋葬一事那时和现在都是众人皆知的，或因\Quote{照经上所记载的}这句话是两者共用的。那么，他为什么在此处加上\Quote{照经上所记载的}，即\Quote{并且祂第三天复活，照经上所记载的，}而不满足于前面那样笼统的说法？因为这对大多数人来说也是隐晦的：因此他再次凭着灵感引入\Quote{圣经}，构思出如此智慧而神圣的思想。\par

那么，他在论及祂的死亡时为何也这样做呢？因为在那种情况下，虽然十字架对众人是显明的，祂当着所有人的面被悬在上面，但原因却不再那么清楚。众人确实都知道祂死了，但祂为世人之罪而受此苦，却不再同样为大众所知。因此他引入圣经的见证。\par

然而，我们所说的已足以证明这一点。但圣经在何处说祂被埋葬，并在第三天复活呢？藉着约纳的预象，主自己也引用说：\BibleQuote{正如约纳在大鱼腹中三天三夜，同样，人子也要在地心三天三夜。}\BibleRef{玛 12:40}藉着旷野中的荆棘。因为正如荆棘燃烧却没有烧毁，\BibleRef{出 3:2}同样，那身体也确实死了，但没有持续被死亡拘禁。还有达尼尔中的龙也预示了这一点。因为正如龙吞下先知给的食物，便从中裂开；同样，阴府吞下那身体，也被撕裂，那身体自行撕裂其子宮而复活。\par

如今，若你也愿意以言语聆听那些你在预象中所见之事，请听依撒意亚说：\BibleQuote{祂的生命从地上被夺去，}\BibleRef{依 53:8}并且主乐意洁净祂的创伤……把光显示给祂；以及他之前的达味说：\BibleQuote{祢不会将我的灵魂遗弃在阴府，也不让祢的圣者见到腐朽。}\BibleRef{咏 15:10}\par

因此，保禄也引导你们到圣经，使你们知道这些事不是无缘无故或随意发生的。因为当这么多先知预先描述和宣告它们时，它们怎能是偶然呢？圣经在提及我们主的死亡时，绝无意指罪的死亡，而是身体的死亡，以及同一性质的埋葬和复活。\par

% structure: p0029 source=h2 target=subsection reason=relative-heading-rank; base=h2

\subsection{格林多前书 15:5}

5. \Quote{并且祂显现给\Person{刻法}：}他立即提名最可信之人。\Quote{然后显现给那十二位。}\par

% structure: p0031 source=h2 target=subsection reason=relative-heading-rank; base=h2

\subsection{格林多前书 15:6}

\Quote{随后祂同时显现给五百多位弟兄；其中大多数至今仍在，但有些人已经睡了。}\par

% structure: p0033 source=h2 target=subsection reason=relative-heading-rank; base=h2

\subsection{格林多前书 15:7}

\Quote{然后显现给雅各伯；随后显现给众宗徒。}\par

% structure: p0035 source=h2 target=subsection reason=relative-heading-rank; base=h2

\subsection{格林多前书 15:8}

\Quote{末了，也显现给我这个像流产儿的人。}\par

因此，既然他提到了圣经的证明，他也加上事件的证明，列出先和众先知之后，宗徒和其他信徒作为复活的证人。相反，如果他指的是另一种复活，即从罪中得救，那么他说\Quote{祂显现给某人}就是徒劳的；因为这是建立身体复活之人的论证，而非模糊教导从罪中得救的人。为此，他也没有笼统地说一次\Quote{祂显现了，}尽管他本可以这样做，只用一个普遍的表达：但他现在说了两次、三次，几乎在每一个见过祂的人身上都使用了它。因为他说：\Quote{祂显现给\Person{刻法}，祂显现给五百多位弟兄，也显现给我。}然而福音书却说相反的事，即祂先显现给玛利亚。\BibleRef{谷 16:9}\par

但祂这里所说的是哪十二位宗徒呢？因为祂被接升天之后，玛弟亚才被选入数中，并非在复活之后立刻。但很可能祂甚至在升天之后也显现了。无论如何，我们的这位宗徒自己在祂升天之后也被召叫，并看见了祂。因此，他也没有定下时间，而是简单而不加限定地叙述了显现。因为很可能有许多次显现；故此若望也说：\Quote{这是第三次祂显现。} \BibleRef{若 21:14}\par

\Quote{以后同时显现给五百多位弟兄。}有些人说\Quote{以上}是指从天以上，即\Quote{不是在地上行走，而是在上面、在头顶上显现给他们：}并补充说，保禄的目的不仅是证实复活，也是证实升天。另有些人说，\Quote{五百多位}这句话的意思是\Quote{超过五百人。}\par

\Quote{其中大半至今还在。}因此，他说：\Quote{虽然我讲述的是古老的事，但我仍有活着的见证人。}\Quote{但也有已经长眠的。}他没有说\Quote{死了}，而是说\Quote{长眠了}，以此再次证实复活。\Quote{以后显现给雅各伯。}我想，这是祂的兄弟。因为据说主亲自祝圣了他，并首先立他为耶路撒冷的主教。\Quote{然后显现给众宗徒。}因为还有其他的宗徒，如那七十二人。\par

\Quote{最后，也显现给我这个像流产儿的人。}这与其说是别的，不如说是谦卑的表达。因为并非因我是最小的，所以祂在其余人之后才显现给我。即使祂最后召叫了我，我仍比许多在我之前的人更显赫，甚至比所有人都更显赫。那五百多位弟兄也不一定比雅各伯更好，因为祂先向他们显现。\par

为什么祂不向所有人同时显现呢？为了先播下信德的种子。因为最先看见祂并确切深信的人，将这事告诉其余的人；他们的报告先来到，使听者期待这一伟大奇迹，并为信德之眼见铺平道路。因此，祂既没有向所有人一起显现，也没有在一开始向许多人显现，而是先向一人显现，且是全体中的领袖和最信实的人；因为确实需要极其信实的灵魂才能首先领受这一显现。那些在别人看见之后才看见祂，并从他们那里听到的人，在他们的见证中得到了对自己信德不小的贡献，并有助于预先准备他们的心思；但第一个被算为配得看见祂的人，正如我所说，需要伟大的信德，才不会被如此出乎意料的景象所困惑。因此，祂首先显现给伯多禄。因为那首先承认祂是基督的人，理应首先配得看见祂的复活。不仅如此，也因为他否认了祂，为了更充分地安慰他，并表明他并没有被绝望，在其余人之前，祂赐予他这一显现，并将祂的羊群首先托付给他。因此，祂也首先显现给妇女们。因为这个性别被视为低等，所以祂的诞生和祂的复活中，这个性别首先尝到了祂的恩宠。\par

但在伯多禄之后，祂也间隔地显现给每一个人，有时向较少的人，有时向较多的人，以此使他们彼此成为见证人和教师，并使祂的宗徒们在所说的一切事上值得信赖。\par

6. \Quote{最后，也显现给我这个像流产儿的人。}这里的谦卑言辞是什么意思？在什么情况下是合宜的？因为如果他希望表明自己值得信任，并将自己列在复活的见证人中，他就在做与他愿望相反的事：因为那时他应当赞扬自己，表明自己是伟大的，而他在许多地方正是这样做的，当时机要求时。然而，他在这里也谦卑地说话的原因正是他即将这样做。但不是立即，而是以他自己特有的良好判断力：首先谦卑地说话，并堆积许多对自己的指责，然后才颂扬关于自己的事。原因何在？为了当他即将说出关于自己的那伟大而崇高的言辞\Quote{我比众位更劳苦}时，他的话语能更被接纳，既因前文，也因这是作为前文的结果而非主题。因此，当致书弟茂德时，他打算说关于自己的大事，却先列出对自己的指控。因为所有的人在赞扬别人时，都自由而大胆地说话；但被迫赞扬自己的人，特别是当他也要召唤自己作证时，会感到不安和脸红。因此，这位蒙福者也先陈述自己的可怜，然后说出那崇高的言辞。他这样做，部分是为了减轻谈论自己的冒犯，部分是为了借此使之后要说的内容更容易被相信。因为那真实陈述自己可耻之事、毫不隐瞒的人——例如他迫害教会、摧毁信德——他由此也使自己的荣耀之事无可怀疑。\par

请注意他极大的谦卑。在说了\Quote{最后，也显现给我}之后，他还不满足：因为\BibleQuote{有许多在先的，将要在后；在后的，将要在先}，\BibleRef{玛 20:16}他说。因此他加上了\Quote{像流产儿一样。}他也没有停在这里，还加上自己的判断并附有理由，说道：\par

% structure: p0046 source=h2 target=subsection reason=relative-heading-rank; base=h2

\subsection{格林多前书 15:9}

\BibleQuote{因为我原是宗徒中最小的，不配称为宗徒，因为我迫害过天主的教会。}\par

他没有说\Quote{只有十二宗徒}，而是说\Quote{所有其他宗徒}。他说这些话，一方面是因为他谦卑地说话，另一方面也是因为他确实如我所说那样心怀谦逊，预先为将要说的话作好安排，使之更易被接受。因为如果他站出来说：\Quote{你们应当相信我，基督从死者中复活了；因为我看见了他，而且我是所有人中最值得信赖的，因为我劳苦更多}，这话可能会冒犯听众。但现在，他先谈论那些令人羞愧和带有指责的话题，从而消除了这种叙述中可能令人不快之处，并为听众相信他的证言铺平了道路。\par

因此，正如我所说，他不仅简单地宣称自己是最末的、不配称为宗徒，而且还说明了原因，说：\Quote{因为我迫害了教会。}然而，所有这些都已得到宽恕，但他自己从未忘记它们，想要彰显天主恩宠的伟大。因此，他继续说：\par

% structure: p0050 source=h2 target=subsection reason=relative-heading-rank; base=h2

\subsection{\BibleRef{格前 15:10}}

7. \Quote{然而，因天主的恩宠，我成了今日的我。}\par

你看到又一种极度的谦卑吗？他将缺点归咎于自己，但对善行却只字不提；反而将一切归于天主。接着，为了不使听者因之而懈怠，他说：\Quote{他赐给我的恩宠没有落空。}这又是有所保留的：他没有说\Quote{我显示了与他恩宠相称的勤勉}，而是说\Quote{没有落空}。\par

\Quote{我比他们众人更劳苦。}他没有说\Quote{我受尊敬}，而是说\Quote{我劳苦}；当他提到危险和死亡时，用\Quote{劳苦}一词再次缓和了语气。\par

然后，他又实践了惯常的谦卑，迅速将这一点也归诸天主，说：\Quote{不是我，而是与我同在的天主的恩宠。}有什么比这样的灵魂更令人敬佩？他在多方贬低自己之后，只说了一句崇高的话，却连这句话也不认为是自己的；从前面和后面的事上，他处处设法收敛这句高调的话，因为那是被迫说出的。\par

但请看他是如何充满谦卑的表达。他说：\Quote{他最后显现给我这个像流产儿的人。}因此，他在自己前面没有提任何人，并说自己是\Quote{宗徒中最小的}，甚至不配这个称号。他还不满足于此，为了不让人以为他只是嘴上谦卑，他还陈述了理由和证据：说自己是\Quote{流产儿}，是因为他最后见到了耶稣；说不配称为宗徒，是因为他\Quote{迫害了教会}。因为一个单纯谦卑的人不会这样做；但一个陈述理由的人，是从痛悔的心说出这一切。因此，他在别处也提到这些事，说：\Quote{我感谢那赐给我能力的，我们的主基督耶稣，因他认为我忠信，派我服事他，尽管我先前是亵渎者、迫害者和施暴者。} \BibleRef{弟前 1:12}\par

但他为什么说出那句高调的话：\Quote{我比他们众人更劳苦}？他看出是情势所迫。因为如果他不说这句话，只贬低自己，他又怎能大胆地为自己作证，将自己与其他人并列，并说：\par

% structure: p0057 source=h2 target=subsection reason=relative-heading-rank; base=h2

\subsection{\BibleRef{格前 15:11}}

\Quote{总之，不论是我，或是他们，我们都这样宣讲。}\par

因为证人应当是可信的，且是一个伟人。至于他如何\Quote{比他们众人更劳苦}，他在前面已经指出，说：\Quote{难道我们没有权利吃喝，像其他宗徒一样吗？}又说：\Quote{对没有法律的人，我就如同没有法律的人。}这样，在需要严格的地方，他超越众人；在需要迁就的地方，他同样显示出极大的优越。\par

但有些人引用他被派往外邦人那里，以及他跑遍了世界大部分地方。由此可见，他得到了更多的恩宠。因为如果他劳苦更多，恩宠也更多；但他得到更多恩宠，是因为他也显示了更多的勤勉。你看，通过那些他用来试图掩没自己事迹的细节，他反而显示出自己是第一位的。\par

8. 当我们听到这些时，让我们也公开显扬自己的缺点，而对我们的优点则闭口不言。如果机会迫使我们谈及，也要有所保留，并将一切归诸天主的恩宠。宗徒也正是这样做的，不断给自己的以前生活打上坏标记，而将以后的状态归诸恩宠，为要表明天主从各方面而来的慈悲：从祂拯救了这样的我，并在拯救后又使我成为今日的我。因此，犯罪者不要绝望，有德者不要自负；前者要极其恐惧，后者要进取。因为懒惰的人不能坚守德行，勤勉的人也不能轻易摆脱邪恶。有福的达味就是这两方面的例子：他稍一打盹，就大大跌倒；当他心中刺痛时，又迅速恢复到从前的高度。事实上，两者都是邪恶：绝望和懒惰；前者能使人从高天迅速堕落；后者则不让跌倒的人重新站起。因此，保禄针对前者说：\Quote{自以为站得稳的，要小心，免得跌倒。}\BibleRef{格前 10:12}；但对于后者，他说：\Quote{今天你们如果听从他的声音，不要再心硬。}\BibleRef{希 4:7}；又说：\Quote{你们要举起下垂的手和麻木的膝盖。}\BibleRef{希 12:12}。对于那些犯了奸淫但悔改的人，他很快安慰他，\Quote{免得这样的人被过度的忧愁所吞没。}\BibleRef{格后 2:7}\par

既然如此，人啊，你在其他悲伤中为何沮丧？因为如果为了罪——唯在罪中悲伤才有益处——过度悲伤会造成很大损害，那么其他一切事情更当如此。你为何悲伤呢？因为你损失了钱财？请看那些连面包都吃不饱的人，你很快就会得到安慰。在每一件令你痛苦的事中，不要为所发生的事哀叹，而要感谢那些未曾发生的灾难。你曾有钱财，现在失去了？不要为损失流泪，而要感谢你享有它的时光。如同约伯所说：\BibleQuote{我们由天主手里接受了好，难道就不能接受坏吗？}\BibleRef{约 2:10}同时还要运用这个论点：即使你失去了钱财，但你的身体仍然健康，你并未因贫穷而哀叹它也已残缺。但你的身体也遭受了伤害吗？然而这还不是人类苦难的底线，你仍在木桶的中段被承载。因为许多人除了贫穷和残疾，还与魔鬼搏斗，在荒野漂泊；还有人遭受比这些更痛苦的事。愿我们永远不必承受一个人所能忍受的一切。\par

因此，要常常思量这些事，记念那些处境更糟的人，不为任何事烦恼；但当你犯罪时，才叹息，才哭泣；我不禁止你，反倒嘱咐你这样做；不过即使那时也要节制，记着还有回头之路，还有和好。但你是否看见别人奢侈，自己贫穷；别人穿华服，居高位？然而不要只看这些，也要看从中产生的痛苦。在你的贫穷中，不要只看乞讨，也要把从那里产生的快乐算进去。因为财富确有一副欢快的面具，但内里充满阴郁；贫穷则相反。如果你展开每个人的良心，在穷人的灵魂里你会看到极大的安全和自由；但在富人的灵魂里，却是混乱、无序和风暴。如果你因看见他富裕而悲伤，他也因看见比自己更富裕的人而更加烦恼。正如你畏惧他，他也畏惧另一个人，在这方面他并不比你强。但你看见他在职位上而烦恼，因为你是个平民，是被管辖者之一。可是请回想他卸任的那一天。就在那天之前，还有骚动、危险、劳累、谄媚、不眠之夜，以及一切痛苦。\par

9. 我们这些话是对那些无意于高尚道德的人说的；因为如果你知道这一点，还有别的更重大的事可以用来安慰你；但眼下我们必须用较粗浅的论点来与你辩论。因此，当你看见一个富人时，要想想比他更富的人，你就会看见他和你处于同样的境况。然后，再看看比你更穷的人，要想到有多少人饿着肚子上床，失去了家产，住在地牢里，每日求死。因为贫穷不会产生悲伤，财富也不会产生快乐，两者都是由我们自己的思想产生的。让我们从最低处开始想：清道夫悲伤烦恼，因为他无法摆脱这既悲惨又令人鄙视的营生；但如果你使他摆脱这个，让他安稳地拥有生活必需品，他会因为没有超出所需而再次悲伤；如果你给他更多，他又想加倍，因此和以前一样烦恼；如果你给他两倍或三倍，他又会因为没有参与政事而沮丧；如果你也给他这个，他又会因自己不是最高官员而自认不幸；当他得到这个荣誉时，他会因不是统治者而哀伤；当他成了统治者，又因不是整个民族的统治者而哀伤；统治了整个民族，又因不是许多民族的统治者；统治了许多民族，又因不是所有民族的统治者；当他成为副王时，又会因自己不是国王而烦恼；如果是国王，又因不是唯一的国王而烦恼；如果是唯一的，又因不是蛮族的国王而烦恼；如果是蛮族的，又因不是整个世界的国王而烦恼；如果是全世界的，为何不也是另一个世界的呢？如此，他的思想无止境地循环，从来不允许他快乐。你看见了吗？即使你让一个人从卑贱贫穷变成国王，你若没有先矫正他贪求更多的思想倾向，你就不能消除他的沮丧。\par

来，让我也向你展示相反的情形：即使你将一个具有思虑的人从高位降到低位，你也不会使他陷入沮丧和忧愁。若你愿意，让我们沿着同一阶梯下行，请将总督从他的宝座上降下，并假设剥夺他那尊位。我说，他并不会因此烦恼，若他选择牢记我所说的那些事。因为他不会计算他所被剥夺的事物，而是计算他仍拥有的，即从他职位而来的荣耀。但若你也剥夺了这点，他会计算那些身居平民位置、从未升至如此权位的人，而他的财富足以安慰他。若你再将他逐出财富，他会看向那些拥有中等产业的人。若你甚至夺去中等财富，只允许他享用必要的食物，他可能会想到那些连这点也没有、却与不断的饥饿搏斗并生活在监狱中的人。即使你将他投入那牢狱，当他想到那些患有不治之症和不可医治的痛苦的人时，他会看到自己处于更好的境况。正如先前所说的清道夫，即使被立为王，也不会收获任何快乐；同样，这个人若成为囚犯，也绝不会使自己烦恼。所以，财富不是快乐的根基，贫穷也不是悲伤的根基，而是我们自己的判断，以及这个事实：我们心灵的眼睛并不纯净，也不固定在任何地方而安住，而是无限制地四处游荡。正如健康的身体，若只以面包为养料，便处于良好而有力的状态；但那些患病的，即使享受丰盛多样的饮食，却变得更加虚弱；灵魂的情况也是如此。心胸狭窄者，即使戴上王冠和享有不可言说的尊荣也不能快乐；但克己者，即使身带锁链和脚镣并处于贫穷，也会享受纯粹的快乐。赖我们的\Person{主耶稣基督}的恩宠和慈悯，愿光荣、权能、尊威归于\Person{圣父}及\Person{圣神}，同祂一起，从现今直到永远，万世无疆。阿们。\par

10. 这些事既铭记在心，让我们始终看向那些比我们低的人。我承认，确实有另一种安慰，但那是道德上的高深造诣，超脱于多数人的粗俗。这是什么？那就是：财富是虚无，贫穷是虚无，耻辱是虚无，尊荣是虚无，它们只在短暂的时间和言辞上彼此有别。与此相伴还有另一个更伟大的抚慰话题：对将来之事的思虑，无论是恶是善，那些真正为恶和真正为善的事，以及因它们而得的安慰。但正如我所说，既然许多人远离这些教理，我们不得已而讨论其他话题，以便循序渐进地将之前所讲的内容之接受者引向它们。\par

那么，让我们考虑这一切，以各种方式正确塑造自己，我们就决不会为这些意外之事悲伤。因为即使我们看见画中人物富有，我们也不会说他们是可羡慕的，正如看见画中穷人，我们不会称他们为可怜可悲的；尽管那些画中人物比我们所认的富人更为持久。因为一个人在画中保持富有比在事物本性中更为长久：画中人物常常持续百年仍是富有的样子，而现实中的富人有时连安逸地享受其财产一年都没有，就突然被剥夺了一切。那么，默想这一切，让我们从四面八方建立快乐，作为对抗我们非理性忧愁的堡垒，好使我们既能愉快地度过今生，又能获得将来的福乐，赖我们的\Person{主耶稣基督}的恩宠和慈悯，愿光荣、权能、尊威归于\Person{圣父}及\Person{圣神}，同祂一起，从现今直到永远，万世无疆。阿们。\par

\SourceRecord{Translated by Talbot W. Chambers.；Nicene and Post-Nicene Fathers, First Series；Vol. 12.；Edited by Philip Schaff.；Buffalo, NY: Christian Literature Publishing Co.,；1889.；<http://www.newadvent.org/fathers/220138.htm>.}

% structure: p0001 source=h1 target=section reason=page-title

\section{\WorkTitle{格林多前书第三十九篇讲道}}

% structure: p0002 source=h2 target=subsection reason=relative-heading-rank; base=h2

\subsection{格前 15:11}

\BibleQuote{不拘是我，或是他们，我们就这样宣讲，你们也就这样相信了。}\par

在赞扬了宗徒并抑制了自己之后，他又高举自己超越他们，以便建立平等（因为他指出自己与他们各有优势，从而实现了平等），并由此证明自己值得信任；即便如此，他也没有遣散他们，而是再次将自己与他们并列，指出他们在基督内的和谐一致。然而，他这样做并非要让自己显得是附属于他们，而是要使自己也出现在同等的行列中。因为这样对福音有益。因此，他同样热切地努力：一方面，不让自己显得忽视他们；另一方面，也不因他们所受的尊荣而被自己管辖的人轻视。因此，他现在又使自己与他们平等，说道：\BibleQuote{不拘是我，或是他们，我们就这样宣讲。} \Quote{你们愿意从谁学习，就从谁学习，}他说，\Quote{我们之间没有区别。}他没有说：\Quote{如果你们不信我，就当信他们；}而是使自己值得信任，并说自己本身是足够的，同时也为他们本身同样肯定。因为人的差别没有效果，他们的权威是相等的。在 \WorkTitle{迦拉达书} 中他也这样做，将他们与自己同列，并非因为需要他们，而是表明他自己已经足够：\BibleQuote{那些有声誉的人并没有增加我什么：} \BibleRef{迦 2:6}尽管如此，我还是追求与他们的合一。\BibleQuote{因为他们将右手给了我，}他说，\BibleRef{迦 2:9}因为如果保禄的声誉总是依赖他人并依靠他人的证词来确认，门徒们就会因此受到无限的伤害。因此，他这样做并非要抬高自己，而是为福音担忧。因此，他在这里也使自己平等，说道：\BibleQuote{不拘是我，或是他们，我们就这样宣讲。}\par

\Quote{不拘是我，或是他们，我们就这样宣讲。} \Quote{你们愿意从谁学习，就从谁学习，}他说，\Quote{我们之间没有区别。}他没有说：\Quote{如果你们不信我，就当信他们；}而是使自己值得信任，并说自己本身是足够的，同时也为他们本身同样肯定。因为人的差别没有效果，他们的权威是相等的。在 \WorkTitle{迦拉达书} 中他也这样做，将他们与自己同列，并非因为需要他们，而是表明他自己已经足够：\BibleQuote{那些有声誉的人并没有增加我什么：} \BibleRef{迦 2:6}尽管如此，我还是追求与他们的合一。\BibleQuote{因为他们将右手给了我，}他说，\BibleRef{迦 2:9}因为如果保禄的声誉总是依赖他人并依靠他人的证词来确认，门徒们就会因此受到无限的伤害。因此，他这样做并非要抬高自己，而是为福音担忧。因此，他在这里也使自己平等，说道：\BibleQuote{不拘是我，或是他们，我们就这样宣讲。}\par

他说得好：\BibleQuote{我们宣讲，}表明他极大的胆量。因为我们不是秘密地说话，也不是在角落里，而是发出比号角更清晰的声音。他没有说：\BibleQuote{我们宣讲过，}而是\BibleQuote{如今我们仍然这样宣讲。} \BibleQuote{你们也就这样相信了。}这里他没有说：\BibleQuote{你们相信，}而是\BibleQuote{你们相信了。}因为他们心意动摇，所以他回到从前的时候，并接着引用他们自己的见证。\par

% structure: p0007 source=h2 target=subsection reason=relative-heading-rank; base=h2

\subsection{格前 15:12}

2. \BibleQuote{如果基督被宣讲了，说祂从死者中复活了，怎么你们中间有人说没有死者的复活呢？}\par

你看到他是何等巧妙地推理，并从基督复活的事实证明复活，而且首先以多种方式确立了前者吗？他说：\Quote{因为先知们也预言了这事，主自己也借着显现表明了，我们宣讲，你们也相信了；}这样编织了他的四重见证：先知的见证、事态结果的见证、宗徒的见证、门徒的见证；或者更确切地说是五重见证。因为连这个原因本身也暗示了复活，即祂为别人的罪而死。因此，如果这事已得到证明，那么显然另一件事也随之而来，即其他死去的人也同样复活。这就是为什么他像对待一个公认的事实一样，质问他们，说：\BibleQuote{如今基督若已经复活了，怎么你们中间有人说没有死者的复活呢？}\par

由此他又一次减轻了反对者的狂妄；因为他说：\Quote{怎么你们说，}而是\Quote{怎么你们中间有人说。}他既不指控所有人，也不公开指名他所控告的人，以免使他们更加鲁莽；另一方面，他也没有完全隐瞒，以便纠正他们。为此，他将他们与群众分开，准备与他们争辩，借此既削弱又迷惑他们，并让其余的人在与这些人的冲突中对真理更坚定，也不让他们投向那些企图毁灭他们的人：他其实已经准备好采取激烈的言辞。\par

再者，免得他们说：\Quote{基督复活这事对所有人都是清楚明白的，没有人怀疑；但这并不必然意味着另一件事，即人类的复活：}——因为前者曾受预告并确实发生，并由祂的显现所证实，即基督复活的事实；但后者还在希望中，即我们自己的分——请看他是如何做的；他从另一方面又加以证明，这是大能的明证。他说：\BibleQuote{为什么有人说没有死者的复活呢？}当然，如果是这样，那么前者反过来也会被这推翻，即基督复活的事实。因此他又加上说：\par

% structure: p0012 source=h2 target=subsection reason=relative-heading-rank; base=h2

\subsection{格前 15:13}

\BibleQuote{若没有死者的复活，基督也就没有复活。}\par

你看到保禄的精力和他战斗的精神是多么不可战胜吗？他不仅从明显的事证明可疑的事，而且还从可疑的事试图向反对者证明先前明显的事？不是因为已经发生的事需要证明，而是为了表明这事与那事同样值得相信。\par

3. \Quote{这是什么推论呢？}有人说。\Quote{因为若基督没有复活，那么别人也不会复活，这是必然的；但若别人没有复活，基督也不会复活，这又有什么理由呢？}既然这看来并不十分合理，请看他如何明智地推理，他从起初，甚至从福音的根基开始，预先撒下种子，例如：祂\BibleQuote{为我们的罪死了，}并且复活了，而且祂是\BibleQuote{长眠者的初果。}因为初果——若不是为那些复活的人，祂能是谁的初果呢？如果那些祂作为初果的人不复活，祂又怎能是初果呢？那么他们怎么不复活呢？\par

再者，如果死者不被复活，基督为何复活？祂为何来临？祂为何取了肉身，若祂不打算再复活肉身的呢？因为祂自己并非需要它，而是为了我们。但后来他继续下去时会写下这些事；目前他说：\Quote{如果死者不被复活，基督也未被复活}，好像那与此相连。因为若祂不打算复活自己，祂就不会成就那另一项工作。你看见那整个\OriginalTerm{救恩计划}{œconomy}如何逐渐被他们的话语和他们对复活的不信所推翻了吗？但至今他还没有谈到降生成人，而是谈到复活。因为并非祂成为肉身，而是祂的死亡，除去了死亡；因为当祂有肉身时，死亡的暴政仍统治着。\par

% structure: p0017 source=h2 target=subsection reason=relative-heading-rank; base=h2

\subsection{格林多前书 15:14}

\BibleQuote{如果基督未被复活，那么我们的宣讲就是徒然的，你们的信德也是徒然的。}\par

虽然按顺序本应接着说的是：\Quote{但如果基督没有复活，你们就是与明显的事实作对，与那么多先知和事实的真理作对}；然而，他说了更令他们恐惧的事：\Quote{那么，我们的宣讲就是徒然的，你们的信德也是徒然的。}因为他想要彻底震动他们的心思：\Quote{我们失去了一切，}他说，\Quote{一切都完了，若祂没有复活。}你看见那\OriginalTerm{救恩计划}{œconomy}的奥秘何其伟大吗？正如这样：如果死后祂不能再复活，那么罪就没有被解除，死亡没有被除去，咒诅也没有被移除。是的，不仅我们宣讲徒然，你们也相信徒然。\par

4. 不仅借此他显明了这些邪恶教义的不敬，他还进一步竭力反驳它们，说：\par

% structure: p0021 source=h2 target=subsection reason=relative-heading-rank; base=h2

\subsection{格林多前书 15:15}

\BibleQuote{是的，并且我们被发现是天主的假证人，因为我们为祂作证说祂复活了基督；而祂并没有复活祂，如果死者不被复活的话。}\par

但如果这是荒谬的（因为这是对天主的控诉和诽谤），并且祂没有复活祂，如你们所说，那么不仅这一点，其他荒谬之事也会随之而来。\par

他再次确立这一切，并重新提起，说：\par

% structure: p0025 source=h2 target=subsection reason=relative-heading-rank; base=h2

\subsection{格林多前书 15:16}

\BibleQuote{因为如果死者不被复活，基督也未被复活。}\par

因为若祂不打算这样做，祂就不会来到世界。但他没有提到这一点，而是提到了结局，也就是祂的复活；通过它吸引万有。\par

% structure: p0028 source=h2 target=subsection reason=relative-heading-rank; base=h2

\subsection{格林多前书 15:17}

\BibleQuote{如果基督没有被复活，你们的信德就是徒然的。}\par

他用一切清晰和公认的事，不断围绕基督的复活，借着更强的论点，使那似乎软弱和可疑的，也变得坚强和清晰。\par

\BibleQuote{你们仍在你们的罪中。}因为如果祂没有被复活，祂也没有死；如果祂没有死，祂也没有除去罪：祂的死就是除罪。\BibleQuote{看，}有人说，\BibleQuote{天主的羔羊，除免世罪者。}\BibleRef{若 1:29} 但\Quote{除免}如何？借着祂的死。因此，祂也称祂为羔羊，如被宰杀的一样。但如果祂没有复活，祂也没有被宰杀；如果祂没有被宰杀，罪也没有被除去；如果它没有被除去，你们就在其中；如果你们在其中，我们宣讲徒然；如果我们宣讲徒然，你们相信你们已被和好也是徒然。此外，死亡仍然是不死的，如果祂没有复活。因为如果连祂也被死亡抓住，没有解除它的痛苦，那么祂自己仍被它抓住，如何释放所有其他人呢？因此他又加上：\par

% structure: p0032 source=h2 target=subsection reason=relative-heading-rank; base=h2

\subsection{格林多前书 15:18}

\BibleQuote{那么那些在基督里安眠的人也灭亡了。}\par

\Quote{我为什么提到你们呢？}他说，\Quote{当所有那些也灭亡了的人，他们做了一切，不再受未来不确定性的支配？}但借着\Quote{在基督里}这个表达，他指的是\Quote{在信德中}，或者\Quote{那些为祂而死的人，他们忍受了许多危险，许多痛苦，走了窄路。}\par

那些口出秽言的\OriginalTerm{摩尼教徒}{Manichees}在哪里呢？他们说这里的复活指的是从罪中得解放。因为这些紧凑而连续的推理，反之亦然，并不表明他们所说的任何东西，而是表明我们所肯定的。确实，\Quote{再次起来}是指那跌倒的人说的：这就是为什么他不断解释，不仅说祂被复活，还加上\Quote{从死者中}这话。格林多人也并非怀疑罪的赦免，而是怀疑身体的复活。\par

但究竟有什么必要，除非人类不是没有罪的，基督自己也该如此？然而，如果祂不复活人类，很自然会说：\Quote{祂为何来临，取了我们的肉身并复活了呢？}但按照我们的假设并非如此。是的，无论人类犯罪或不犯罪，在天主那里总是不能犯罪的；发生在我们身上的事触及不到祂，一个情况与另一个情况也不以转换的方式相呼应，如同身体的复活这事一样。\par

% structure: p0037 source=h2 target=subsection reason=relative-heading-rank; base=h2

\subsection{格林多前书 15:19}

4. \BibleQuote{如果我们只在今生寄望于基督，我们就是众人中最可怜的了。}\par

你说什么，\Person{保禄}？\Quote{只在今生我们怀有希望}是什么意思，如果我们的身体不被复活，而灵魂存留且不死？因为即使灵魂存留，即使它无限地不死（它确实如此），没有肉体，它将不会接受那些隐藏的美事，也同样不会受惩罚。因为众人都要在基督的审判台前显露出来，\newline{}\BibleQuote{好使每人借着身体，按其所行的，或善或恶，领取报应。}\newline{}\BibleRef{格后 5:10} 因此他说：\BibleQuote{如果我们只在今生寄望于基督，我们就是众人中最可怜的了。}因为如果身体不再复活，灵魂就无冠冕，没有那在天上的福乐。如果这样，那时我们将一无所享；如果那时一无所享，赏报就在今生。\Quote{那么在这方面，有谁能比我们更悲惨呢？}他说。\par

但他说这些话，既是为坚定他们关于身体复活的道理，也是为劝服他们关于那不朽的生命，使他们不至于以为我们的一切都止于今世。因为他已藉前面的论证充分确立了他的意图，并说了：\BibleQuote{如果死人不复活，基督也就没有复活；但如果基督没有复活，我们就灭亡了，我们仍然在罪中。}他又加上这话，彻底摧毁了他们的傲慢。因为每当他要引入任何必要的道理时，他先藉恐惧彻底动摇他们铁石的心肠。他在这里正是如此，既在上面播下了那些种子，使他们焦虑，如同失去一切的人；如今又换一种方式，用他们最深刻感受的方式，做着同样的事，说：\Quote{\BibleQuote{我们是最可怜的人}，如果在经过如此多的斗争、死亡和无数灾祸之后，我们还要失去如此大的福分，而我们的福乐仅限于今世。}因为事实上一切都依赖于复活。由此也明显可见，他的讲话不是关于从罪中复活，而是关于身体，以及今世和来世的生命。\par

% structure: p0041 source=h2 target=subsection reason=relative-heading-rank; base=h2

\subsection{\BibleRef{格前 15:20}}

5. \BibleQuote{但现在基督已从死人中复活，成为睡了之人初果。}\par

在表明不信复活会产生多大的祸害之后，他再次接过话题说：\BibleQuote{但现在基督已从死人中复活}，不断加上\BibleQuote{从死人中}，以堵住异端者的口。\BibleQuote{成为睡了的人初果。}如果他们是初果，那么他们自己也必会复活。如果他是讲从罪中复活，而没有一个人是无罪的——因为保禄也说：\BibleQuote{我虽然不觉得自己有错，却也不能因此成义}——那么按你们的说法，有谁会复活呢？你们看见他的讲话是关于身体的吗？而且为了使其可信，他不断提出那在肉身中复活的基督。\par

接着他也提出一个理由。因为，如我所说，当一个人宣称却不陈述理由时，他的话不易被大众接受。那么理由是什么呢？\par

% structure: p0045 source=h2 target=subsection reason=relative-heading-rank; base=h2

\subsection{\BibleRef{格前 15:21}}

\BibleQuote{因为死亡既因一人而来，死人的复活也因一人而来。}\par

但如果藉一个人，无疑他是具有身体的。请观察他的用心，如何在另一基础上也使他的论证无可辩驳。正如他所说：\Quote{那被打败的，必须亲自再发起战斗，那被击败的本性也必须亲自取得胜利。因为这样才洗刷了羞辱。}\par

但让我们看看他所说的是哪一种死亡。\par

% structure: p0049 source=h2 target=subsection reason=relative-heading-rank; base=h2

\subsection{\BibleRef{格前 15:22}}

\BibleQuote{因为正如在亚当内众人都死了，照样在基督内众人都要复活。}\par

那么，请告诉我：众人在亚当内都死于罪吗？那么诺厄怎么会在他那世代是义人？亚巴郎怎么是？约伯怎么是？其余所有人怎么是？再者，请问：众人都要在基督内复活吗？那么那些被引入地狱之火的人又在哪里？可见，如果这是指身体而言，道理就成立；但如果是指义德和罪，就不再成立了。\par

此外，免得你听到复活是众人的普遍事，就以为罪人也能得救，他又加上：\par

% structure: p0053 source=h2 target=subsection reason=relative-heading-rank; base=h2

\subsection{\BibleRef{格前 15:23}}

\BibleQuote{但各人要按自己的次序。}\par

因为不要因为你听见复活，就以为所有人都享受同样的福分。既然在惩罚中，并非所有人受同样的苦，差别很大；那么在有罪人和义人的地方，分别就更大了。\par

\BibleQuote{首先是初果基督，然后是属于基督的人}，即信徒和蒙悦纳的人。\par

% structure: p0057 source=h2 target=subsection reason=relative-heading-rank; base=h2

\subsection{\BibleRef{格前 15:24}}

\BibleQuote{再后才是结局。}\par

因为当这些人复活时，一切都要终结，不像现在基督复活后事物仍悬而未决。因此他加上\BibleQuote{在祂来临时}，叫你明白他是指那时而言，\BibleQuote{当祂把国度交付于天主父，当祂废除一切率领者、一切掌权者和大能者的时候。}\par

6. 这里，请留心听我，看我所说的没有一句从你面前溜走。因为我们的争斗是与敌人；因此我们首先必须实践\OriginalTerm{归谬法}{reductio ad absurdum}，保禄也常这样做。因为这样我们就最容易发现他们所说的。那么让我们先问他们：这句话是什么意思：\BibleQuote{当祂把国度交付于圣父的时候？}如果我们照字面理解，而不以合乎天主性的意义理解，那么此后祂就不再保有它。因为交付给别人的，就不再拥有。这不仅荒谬，而且接收的那一方也显得在接收之前不是拥有者。因此按他们的说法，圣父以前就不是君王，掌管我们的事；圣子在这些事之后也不再是君王。那么，首先关于圣父，圣子自己怎么说：\BibleQuote{我父一直工作到现在，我也工作}\BibleRef{若 5:17}；而达尼尔论及祂说：\BibleQuote{祂的王国是永远的王国，必不消灭}\BibleRef{达 7:14}。你看见当人以人的方式理解这话时，产生了多少与圣经相悖的荒谬吗？\par

但那么，他这里说基督\Quote{废除}的是什么\Quote{率领}？是天使的率领吗？绝不是。是信徒的率领吗？也不是。那么是什么率领？是魔鬼的率领，关于这个他说：\BibleQuote{我们的战斗不是对抗血和肉，而是对抗率领者，对抗掌权者，对抗这黑暗世界的统治者}\BibleRef{弗 6:12}。因为现在它还没有完全\Quote{废除}，他们在许多地方工作，但那时他们将停止。\par

% structure: p0062 source=h2 target=subsection reason=relative-heading-rank; base=h2

\subsection{\BibleRef{格前 15:25}}

\BibleQuote{因为祂必须为王，直到把一切仇敌放在祂脚下。}\par

此外，由此又产生另一荒诞之说，除非我们也以与\OriginalTerm{天主性}{Deity}相称的方式理解这一点。因为\Quote{直到}这一表述具有终极和限制的含义；但就天主而言，这并不存在。\par

% structure: p0065 source=h2 target=subsection reason=relative-heading-rank; base=h2

\subsection{\BibleRef{格前15:26}}

\BibleQuote{最后被废除的仇敌是死亡。}\par

何以是最后？毕竟，在魔鬼之后，在其它一切之后。因为在起初，死亡也是最后出现的；先是魔鬼的计谋，然后是我们的不顺服，接着是死亡。因此，实际上它现在已然被废除，但确切地实现是在那时。\par

% structure: p0068 source=h2 target=subsection reason=relative-heading-rank; base=h2

\subsection{\BibleRef{格前15:27}}

7. \BibleQuote{因为祂使万有屈服在祂的脚下。但当祂说万有都已屈服时，显然，那使万有屈服于祂的除外。}\par

% structure: p0070 source=h2 target=subsection reason=relative-heading-rank; base=h2

\subsection{\BibleRef{格前15:28}}

\BibleQuote{当万有都屈服于祂以后，那时，子自己也要屈服于那使万有屈服于祂的。}\par

然而在此之前他并未说是圣父\BibleQuote{使万有屈服于祂}，而是祂自己\BibleQuote{废除}。因为他说：\BibleQuote{当祂废除一切率领者、掌权者时}，又说：\BibleQuote{因为祂必须为王，直到把一切仇敌置于祂的脚下。}那么他在这里为什么说\Quote{圣父}呢？\par

不仅存在这一明显的困惑，而且他还怀着一种难以理解的恐惧，用修正的说法说：\Quote{那使万有屈服于祂的除外}，好像有人会怀疑圣父自己可能不屈服于圣子；还有什么比这更不合理呢？尽管如此，他还是害怕这一点。\par

那么情况如何呢？因为实际上有许多问题接连而来。那么，请仔细听我讲；事实上，我们有必要先谈谈保禄的宗旨和他的思想——这处处闪耀可见——然后再附上我们的解答；这解答本身也是我们解答的组成部分。\par

那么保禄的思想和他的习惯是什么呢？当他单独论及\OriginalTerm{天主性}{Godhead}时，他说话是一种方式；而当他进入\OriginalTerm{救恩计划}{economy}的论述时，则是另一种方式。因此，一旦抓住我们主的肉身，他此后便自由地使用一切降低祂的说法；毫无畏惧，仿佛那肉身能承受所有这类表述。因此，我们也要看看这里，他的论述是关于纯粹的天主性，还是由于降生成人而说出的那些话；或者，让我们首先指出他是在何处做了我所说的这事。他在哪里做了呢？在写给斐理伯人的信中，他说：\BibleQuote{祂虽具有天主的形体，却没有将与天主同等视为应紧持不放的珍宝，反而使自己空虚，取了奴仆的形体，成为人的样式；且以人的形态出现后，祂谦抑自己，听命至死，且死在十字架上。因此，天主极其高举祂。}\BibleRef{斐2:6}\par

你可看到，当他单独论及天主性时，他说出那些伟大的话——祂\Quote{具有天主的形体}，祂与生祂的那位\Quote{同等}，并将一切归于祂？但当他向你显示祂成为肉身时，他又降低了言谈。因为除非你区分这些，否则所说的话便有很大出入。因为，如果祂\Quote{与天主同等}，祂如何能高举与自己同等者呢？如果祂\Quote{具有天主的形体}，祂又怎能\Quote{赐给}祂\Quote{一个名}呢？因为赐予者是赐予没有的人，高举者是高举先前卑微的人。那么，在祂接受\Quote{高举}和\Quote{那名}之前，祂便是不完美和缺乏的，其它许多荒谬结论也会随之而来。但如果你加上降生成人，说这些事便不会错。因此，在这里也考虑这些事，并以这种心态接受这些表达。\par

8. 现在，除了这些，我们还要陈述为什么这段圣经如此写的其它理由。但目前有必要提到：第一，保禄的论述是关于复活——一件被认为不可能并备受怀疑的事；其次，他写信给格林多人，其中有许多哲学家，总是嘲笑这样的事。因为尽管在其他事情上互相争吵，但在这一点上，他们却异口同声地合谋，教条式地宣称没有复活。因此，为这样一个如此被怀疑和嘲笑的主题而争论，既因已形成的偏见，也因事情的困难；为了证明其可能性，他首先从基督的复活来论证。并且，既已从先知、从目击者、从信者那里证明了这一点，当他获得了一个公认的\OriginalTerm{归谬法}{reductio ad absurdum}后，他在后文中也证明了人类的复活。他说：\BibleQuote{因为如果死人不复活，基督也就没有复活。}\par

接着，在紧密地推进了前几节中的这些反诘论证后，他又尝试另一种方式，称祂为\BibleQuote{初果，}并指向祂\BibleQuote{废止一切宰制、权威和异能，以及最后的死亡。}\Quote{死亡怎能被废除，}他说，\Quote{除非它先释放它所持有的身体？}既然他曾论及独生子说了伟大的事，即祂\BibleQuote{交出王国，}也就是说，祂亲自成就这些事，祂自己是战争中的胜利者，并\BibleQuote{将万物置于祂脚下，}为了纠正群众的不信，他又补充说：\BibleQuote{因为祂必须为王，直到祂把一切仇敌都放在祂脚下。}他用\Quote{直到，}不是要终结王国，而是使所说的话值得相信，并促使他们有信心。因为\Quote{不要，}他说，\BibleQuote{因为你们听说祂将废止一切宰制、权威和异能，}即魔鬼、众恶魔的集团（尽管有许多）、不相信者的群众、死亡的暴政，以及一切邪恶：不要恐惧，好像祂的力量已耗尽。因为直到祂完成这一切，祂\BibleQuote{必须为王；}不是说，在祂实现之后祂就不再为王，而是确立另一方面：即使现在不是，将来无疑会是。因为祂的王国不被断绝：是的，祂统御、得胜并持续，直到祂调整好一切事物。\par

但这种说话方式在旧约中也能找到；例如，当它说：\BibleQuote{但主的话永远长存；}\BibleRef{咏 118:89} 以及，\BibleQuote{祢是同样的，祢的岁月不会穷尽。}\BibleRef{咏 101:27} 现在先知在讲述那些需要长时间才能成就且必定发生的事时，就说这些以及类似的话；目的是驱除迟钝听者的恐惧。\par

但为了说明\Quote{直到，}和\Quote{乃至，}论及天主时并不表示终结，请听某人所说：\BibleQuote{从永远到永远，祢是天主：}\BibleRef{咏 89:2} 以及，\BibleQuote{我是，我是，}和\BibleQuote{直到你们年老，我仍是那位。}\BibleRef{依 46:4}\par

正是为此缘故，他把死亡放在最后，以便从战胜其余之物中，不信者也能轻易接受这一点。因为当祂消灭了那引入死亡的魔鬼时，祂更会终结那魔鬼的作为。\par

九、既然他将一切都归于祂，即\BibleQuote{废止宰制和权威，}祂王国的成全（我指信友的得救、世界的和平、邪恶的消除，因为这就是成全祂的王国），以及结束死亡；并且他没有说\Quote{圣父藉着祂，}而是说\Quote{祂自己要废除，祂自己要放在自己脚下，}而且他丝毫没有提及那位生了祂者；他后来担心，在较无理性的人中，这可能使圣子显得比圣父更大，或被视为某个不同的、非受生的本原。因此，他温和地护卫自己，限定自己话语的份量，说：\BibleQuote{因为祂将万物置于祂脚下，}再次将这些伟大成就归给圣父；这不是说圣子没有能力——因为若是这样，他先前为何为祂见证如此伟大的事，并将所说的一切都归于祂呢？——而是出于我所提到的原因，并且为了表明凡为我们所成就的一切，都是圣父和圣子共通的。因为单独祂自己就足以\Quote{将万物置于祂脚下，}再听保禄说\BibleRef{斐 3:21}\BibleQuote{祂要以那能使万有归顺自己的能力，改变我们卑贱的身体，使之同化于祂光荣的身体。}\par

随后他又用了一个修正，说：\BibleQuote{但当祂说，一切都已服从时，显然，那使万物服从于祂的，是被排除在外的，}借此为独生子见证不小的光荣。因为如果祂是较小且远为低下的，他绝不会产生这种担忧。他对此还不满足，又加上了如下的另一件事。我说，以免有人怀疑地问：\Quote{但如果圣父没有被\InnerQuote{置于祂之下？}这毫不妨碍圣子更为强大；}担心这种不虔的假设，因为那个表达不足够指明这点，他又添加了远远超过它的内容：\BibleQuote{但当万物都已服从于祂时，圣子自己也要服从于那使万物服从于祂的，为使天主成为万有中的万有。}显示祂与圣父极大的和谐，并且祂是一切其他美好事物的本原和第一因，祂生了如此伟大、在能力和成就上的那一位。\par

10. 若他所言超过了主题所需，不必惊奇。因为他效法他的师傅如此行：当祂也意在显示与生祂者的和谐，并表明祂并非不按祂的旨意而来时，祂下降至此，我所说的，并非和谐证明所要求，而是在场之人的软弱所需。因为祂向圣父祈祷，别无其他原因；并陈述理由说：\BibleQuote{为叫他们相信祢派遣了我。}\BibleRef{若 11:42}因此，保禄效法祂，在此的说话方式超出了所需；并非要你怀疑一种被迫的奴役，绝非如此；而是为了更彻底地铲除那些不敬的教义。因当他决意连根拔除某物时，他惯常如此做，并且做得更多。同样，例如，当他谈论相信的妻子与不信的丈夫按婚姻法律同居时，为使妻子不认为自己因与不信者的交合和拥抱而被玷污，他并没有说\Quote{妻子是不洁的}，也没有说\Quote{她决不会受不信者伤害}，反而说了更甚的话：\Quote{不信者甚至因她而\InnerQuote{圣化}了}，意思并非指外教人因她而成圣，而是藉着这极其有力的表达来消除她的恐惧。同样在此，他以极强烈的言辞除去那不敬教义的热忱，乃是他如此表达的原因。因为怀疑圣子软弱是极端不敬：（因此他纠正道：\Quote{祂必将一切仇敌放在祂脚下}；）所以另一方面，认为圣父低于祂更是不敬。因此他也以极大的力量除去此念。且注意他如何表达。他并非简单地说\Quote{那将万有置于祂之下的，是被排除在外的}，而是说\Quote{这是显然的}，\Quote{即使被承认，}他说，\Quote{我仍要确定它。}\par

为使你明白这就是所言之事的原因，我要问你一个问题：到那时，圣子是否又经历一次\Quote{服从}？这如何能不算是亵渎且不合天主身份？因为最大的服从与顺命乃是：祂身为天主却取了奴仆的形象。那么祂怎能被\Quote{服从}？你看到吗，为除去这亵渎的观念，他使用了这种表达？且是在恰当而有所保留的意义上？因为祂身为圣子和一个神圣位格，故祂服从；不是按人的方式，而是作为自由行动、拥有一切权柄者。否则，祂如何同登宝座？如何说\BibleQuote{就如父复活死人，同样祂也使祂愿意的人复活}？\BibleRef{若 5:21}又如何说\BibleQuote{凡父所有的都是祂的，而祂所有的都是父的}？\BibleRef{若 16:15}因为这些说法向我们表明，祂的权柄与生祂者的权柄完全相等。\par

11. 但\Quote{当祂将王国交出}是什么意思？圣经承认天主的两个王国：一是藉特别拥有，一是藉创造。如此，就创造而言，祂是万有之王，包括希腊人、犹太人、魔鬼和祂的仇敌；但就使信者、自愿者与服从者成为祂自己的子民而言，祂是他们的王。这个王国也被说成有一个开始。因为关于这事，祂也在第二篇圣咏中说：\BibleQuote{你向我请求，我必将外邦人赐你作产业。}\BibleRef{咏 2:8}同样，祂亲自对门徒说：\BibleQuote{天上地下的一切权柄都已被授予了我。}\BibleRef{玛 28:18}将一切归因于生祂者，并非表示祂自己不足，而是为表明祂是圣子，而非非受生者。然后，这个王国祂确实\Quote{交出}，即\Quote{导向圆满的终结}。\par

\Quote{那么，}有人说，\Quote{祂为何对圣神只字不提？}因为现在他并非在论述圣神，也不把所有事物混为一谈。既然当他说\Quote{只有一个天主圣父，和一位主耶稣}时，无疑并非容许圣神是低等的，他之所以沉默，乃因当时并非紧要，故如此表达。因为他惯于只提圣父，但我们不可因此排除圣子；他也惯于只提圣子和圣神，但我们也不可因此否认圣父。\par

但\Quote{好叫天主成为万有中的万有}是什么意思？意为万有皆依赖祂，无人以为有两个无始的权威，或另有分开的王国；没有任何事物独立于祂。因为当仇敌伏在圣子脚下，而圣子在将仇敌置于脚下后与父毫无差异，且与祂完全和谐，那时祂自己将成为\Quote{万有中的万有}。\par

但有人说，他说此话是为宣示罪恶的消除，仿佛从此万有都会屈服，无人抵抗或行恶。因为当没有罪恶时，显然\Quote{天主将成为万有中的万有}。\par

12.但如果身体不会复活，这些事如何为真？因为最凶恶的敌人——死亡，仍存留着，任意而为。有人说：\Quote{不，因为他们不再犯罪。}但那又怎样？因为他这里讨论的不是灵魂的死，而是身体的死？那么他怎么被\Quote{制服？}因为胜利就是夺取被夺走和被扣留的。但如果人的身体要被扣在土里，那么死亡的暴政仍然存在，这些身体被抓住，而没有别的身体可以让他被战胜。但如果保禄所说的事发生了，正如它无疑将要发生，天主的胜利就会显现，而且是一个荣耀的胜利，因为祂能复活被死亡抓住的身体。因为敌人也是在一个人夺取战利品时才被征服，而不是当他允许他们留在对方手中时；但除非有人敢于夺取他的东西，我们怎能说他是被征服的呢？基督自己在福音中就是以这种胜利的方式说祂已经得胜，这样说：\BibleQuote{他捆绑了壮士后，才能夺取他的家财。}\BibleRef{玛 12:29}。因为如果不是这样，这根本就不是明显的胜利。因为就如灵魂的死，\BibleQuote{已死的人是脱离了罪}\BibleRef{罗 6:7}。（然而我们不能说这是一种胜利，因为胜利者不是不再作恶的人，而是消除了先前情欲束缚的人。）同样，在这种情况下，我也不该称死亡停止吞噬人的身体为光荣的胜利，而应称死亡前此所抓住的身体被从他手中夺走。\par

但若他们仍要争辩，并说这些事是关于灵魂的死，那么\Quote{最后被毁灭}这说法如何成立？因为在每个人受圣洗时，它已完全被毁灭。然而，如果你是指身体，这话是可以接受的；我是说，这样的说法即它\Quote{最后被毁灭}。\par

但若有人怀疑，为何在讨论复活时，他没有提出在主时代复活的那些身体，我们的回答也许是这样：这不能用来支持复活。因为指出那些复活后又死去的人，不适合一个力图证明死亡被完全毁灭的人。是的，这正是他说\Quote{最后被毁灭}的原因，以免你再怀疑他的复活。因为当罪被除去时，死亡更会终止：源头干涸，从它流出的溪流仍存在，这是完全不合理的；根被灭，果实仍存，也是如此。\par

13. 既然在末日，天主的仇敌——死亡、魔鬼和恶神——都要被毁灭，我们不要因天主的仇敌的亨通而沮丧。因为主的仇敌在他们荣耀和升高的时刻就衰败；\BibleQuote{是的，他们如烟消散。}\BibleRef{咏 36:20}。当你看到任何天主的仇敌富有，有武装随从和许多谄媚者时，不要沮丧，而要哀叹、哭泣、呼求天主，使祂将他列在祂的朋友中；他既然是仇敌，越享通，你就越要为他哀悼。因为对罪人我们总该哀哭，尤其当他们享受财富和丰裕的日子时；如同一个人该为生病却大吃大喝的人哀哭一样。\par

但有些人，当他们听到这些话时，心性如此不幸，以致于为此悲叹说：\Quote{我什么都没有，因此眼泪是当得的。}你说得好：\Quote{什么都没有}，不是因为你没有别人所有的，而是因为你把这事看得如同能称为有福；是的，为此你值得无尽的哀哭：就像一个健康的人把躺在软榻上的病人看作有福，这病人并不像他那么悲惨可怜，因为病人不意识到自己的福气。在这些人身上也可观察到这样的结果：不，甚至我们的整个生活都被搅乱和颠倒了。因为这些话毁了许多人，把他们出卖给魔鬼，使他们比那些受饥荒折磨的人更可怜。是的，那些贪求更多的人，比乞丐更悲惨，因为他们被更大更苦的悲伤所占据，这从下面的话可以明显看出。\par

有一次旱灾降临我们的城，众人都为最后的灾祸而战栗，恳求天主救他们脱离这恐惧。人们可以看到那时梅瑟所说的话：\BibleQuote{天变为铜}\BibleRef{申 28:23}，以及最可怕的死亡日日等待。但后来，当仁慈的天主认为合宜时，出乎一切意料，天上降下大量雨水，从此众人都在过节欢宴，如同从死亡的门槛上回来。但在如此大的祝福和众人的共同喜乐中，有一个极其富有的人，面带愁容，垂头丧气，因悲伤而几乎死去；当许多人询问原因，为何在众人的共同喜乐中他独自忧伤时，他竟不能压抑自己的野蛮情欲，被疾病的暴虐所驱使，当众宣布了原因。他说：\Quote{为什么？我拥有上万石麦子，却没有办法处置它们了。}那么，告诉我们，对于这些话，他该被用石头砸死，我们还认为他有福吗？那比任何野兽更残忍的人，共同的敌人？人哪，你说什么？你之所以悲伤，难道是因为众人没有灭亡，好让你敛财吗？你没有听过撒罗满所说的话吗？\BibleQuote{屯积粮米，人民必咒骂他。}\BibleRef{箴 11:26}却成了世界祝福的共同敌人，世界之主的慷慨的敌对者，玛门的友人，毋宁说是奴隶？不，那吐出这些话的舌头不该被割掉，那颗心不该被熄灭吗？\par

14. 你看见金子如何不让人做人，而是做野兽和魔鬼吗？有什么比这个富人更可怜的呢？他每日祈祷要有饥荒，好让他得一点金子。是的，他的情欲到如今已经转向了它的反面：他甚至在丰收时也不欢乐，反而因此忧愁（他已到了这种地步），因为他的财产无限。虽然一个拥有很多的人应当喜乐，但这个人却因此沮丧。你看见吗，如我所说，富人从现有的东西中所得的快乐，不及他们为尚未增添的东西所受的忧愁。那个拥有无数小麦的人比那挨饿的人更忧愁哀叹。而那个仅仅得到必需食物的人，却自夸、跳跃、感谢天主；另一个拥有如此之多的人却烦恼，以为自己完蛋了。那么，给我们带来快乐的不是富余，而是自制的心。因为没有它，即使一个人获得并拥有一切，他也会感到如同被剥夺了一切而哀伤。正如我们现在谈论的这个人，即使他按自己希望的高价卖掉了所有，他也会因为没卖得更高而忧愁；如果他能得到更高价，他又会寻求再提价；如果他用一磅银子的价格卖掉了一斗，他还会因半斗卖不到同样的价钱而烦恼。如果起初价格没有定得那么高，不要惊讶。因为醉汉也不是一开始就发烧，但当他们喝了很多酒后，就会把火焰煽得更旺；同样，这些人抓得越多，就越感到贫穷；比别人赚得多的人，正是最缺乏的人。\par

15. 我说这些话不仅是对这个人，也是对那些有这样疾病的人说的——那些抬高货物价格、利用邻居的贫穷做交易的人。因为没有人关心人性，而是到处都有贪婪的欲望在出售时带出许多。油和酒有的卖得快，有的卖得慢，但都不是出于对他人的关心；而是一个寻求利润，另一个避免因产品变质而损失。因此，因为大多数人不重视天主的法律，把所有东西都关在门内，天主用其他方式引导他们走向人性——即使出于必要，他们也会做点好事——在他们心中注入了对更大损失的恐惧，不让地里出产的东西长久保存，好让他们仅仅出于害怕腐败的损失，甚至违背自己的意愿，将那些被邪恶地埋藏在家的东西卖给穷人。然而，尽管如此，有些人是如此贪得无厌，甚至因此也不改正。例如，许多人甚至倒空了整桶的酒，不给穷人一杯，也不给需要的人一文钱，但酒变成醋后，他们把它倒在地上，连同果实一起毁掉了桶。另一些人，不肯给饥饿的人一小块饼，却把整仓的粮食扔进河里；因为他们不听天主吩咐他们施舍给穷人的话，却在蛾子的命令下，即使不情愿，也倒空了屋里所有的东西，彻底毁坏和浪费；同时招致了损失和许多嘲笑咒骂。\par

他们在世上的事就是这样；但关于来世，还有什么言语能向我们描述呢？因为正如这些人在世上把被虫蛀了的、没用的谷粒扔进河里，同样，做这些事的人，也因此变得无用，天主把他们扔进火河里。因为正如谷粒被蛾子和虫子吞吃，他们的灵魂也被残忍和不人道吞吃。这些事情的原因是他们被钉在现世的事物上，只渴望今生。因此，这样的人充满了无限的悲伤；因为无论你提到什么快乐，对结局的恐惧足以消灭一切，这样的人\BibleQuote{活着的时候已经死了}\BibleRef{弟前 5:6}。\par

如今，不信者有这种感受不足为奇；但那些分享了如此伟大的奥迹、学习了如此高超的关于来世克己道理的人，却喜欢停留在现世的事物上，他们该得到什么宽恕呢？\par

16. 那么，他们喜欢停留在现世事物上的原因是什么？来自于他们沉溺于奢侈，养肥肉体，使灵魂娇嫩，加重他们的负担，使他们的黑暗变大，使他们的帷幕变厚。因为在奢侈中，较好的部分被奴役，而较差的部分占上风；前者被蒙蔽，处处被拖在残缺的状态中；而后者到处引导和带领人，尽管它本应处于被带领的等级。\par

由於靈魂與身體之間的紐帶實在強大，造物主如此安排，免得有人誘使我們因身體陌生而憎惡它。因為天主確實吩咐我們愛仇敵，但魔鬼竟得逞到使一些人甚至憎恨自己的身體。當有人說身體屬於魔鬼時，他除了證明這一點外別無他物；這是昏聵之極。因為若它屬於魔鬼，那這如此完美的和諧何以解釋，這和諧使身體完全適於自制靈魂的活動呢？\Quote{不然，}有人說，\Quote{若它適宜，身體如何能蒙蔽它呢？}蒙蔽靈魂的不是身體，遠非如此，人啊，而是奢華。但我們為何渴望奢華？絕非因為我們有身體，而是出於邪惡的選擇。因為身體需要飲食，而非暴食；身體需要滋養，而非瓦解崩潰。你看，奢華不僅對靈魂有害，也對接收滋養的身體本身有害。因為身體變得軟弱而非強壯，柔軟而非堅實，多病而非健康，沉重而非輕快，羸弱而非結實，醜陋而非俊美，難聞而非芬芳，不潔而非潔淨，痛苦而非安舒，無用而非有益，衰老而非年輕，腐朽而非強壯，遲鈍呆滯而非敏捷，殘缺而非健全。倘若身體屬於魔鬼，它就不該受魔鬼之物的傷害，即罪惡的傷害。\par

17. 然而，身體或食物並非屬於魔鬼，唯獨奢華才是。因為藉此，那惡者行出了其無數的邪惡。他由此使整個民族成為犧牲品。有人說：\Quote{蒙愛者肥胖了，}\Quote{長厚了，擴大了，踢跳了。}\BibleRef{申 32:15}從此也開始了索多瑪的那些雷擊。為宣告此事，厄則克耳說：\Quote{但這是索多瑪的罪孽：在驕傲、麵包飽足和安逸中，他們縱慾放蕩。}\BibleRef{則 16:4}因此保祿也說：\BibleRef{弟前 5:6}\Quote{那沉湎於享樂的，雖生猶死。}這是為何？因為她像一座墳墓，懷著自己的身體，被綁在無數的邪惡上。如果身體如此腐朽，靈魂將如何受影響？她將充滿何種混亂、何種波濤、何種風暴？由此，你看，她變得不合適一切職責，無力輕易說話、傾聽、商議或做任何必要之事。但如同一個舵手，當風暴勝過其技藝時，連同船隻和水手一起沉入深淵：同樣，靈魂與身體一同沉淪於麻木的痛苦深淵中。\par

因為事實上，天主在我們的身體中設置了胃作為一種磨，賦予它相應的力量，並規定了它每日應磨的固定量。因此，若有人投入過多，未被消化之物便傷害全身。由此產生疾病、虛弱和畸形：因為確實，奢華不僅使美麗女子多病，也使她面貌可憎。因為她不斷發出難聞的呼氣，呼出陳酒的氣味，面色過於紅潤，破壞了女子應有的勻稱，失去了一切端莊，身體變得鬆弛，眼瞼充血腫脹，體積過大，肌肉成為無用的負擔；想想這一切引起多大的厭惡。\par

此外，我曾聽一位醫生說，許多人未能達到應有的身高，最大的障礙莫過於奢華的生活。因為氣息被投入的大量食物所阻塞，並忙於消化這些東西，本應用於生長的力氣消耗在這多餘之物的消化上。何必談論痛風、遍體的風濕、由此產生的其他疾病，以及整個可憎之事？因為沒有什麼比一個用大量食物嬌養自己的女人更令人厭惡的了。因此，在較貧窮的女子中，人們可以看到更多的美麗：多餘之物被消耗，而不附著於她們，如同一些多餘的無用黏土，毫無益處。因為她們的日常鍛煉、勞作、艱辛、節儉的餐桌和清淡的飲食，為她們帶來大量的身體健康，並由此帶來大量容光。\par

18. 但若你談論奢華的快樂，你會發現它只到喉嚨為止：因為一旦它經過舌頭，便飛走了，在身體中留下許多令人厭惡的東西。因為，我祈求你，不要只看宴席上的縱慾者，而是當你看到他們起身時，然後跟隨他們，你將看到他們更像是野獸和沒有理性的受造物，而非人類。你將看到他們頭痛、膨脹、拘束、需要床榻和長椅以及大量休息，如同在巨大風暴中顛簸的人，需要他人拯救，並渴望他們尚未膨脹到破裂之前的狀態：他們挺著如同孕婦般沉重的肚子，幾乎無法邁步，幾乎看不見，幾乎說不出話，幾乎做不了任何事。但若他們偶然小睡片刻，又看到奇怪的夢境，充滿各種幻想。\par

對於他們那一種瘋狂，該說什麼呢？我指的是情慾的瘋狂，因為這也有其根源於此。是的，如同發情的公馬追逐母馬，他們被醉酒的刺驅策，撲向一切，比牠們更無理性，跳躍更加瘋狂；並犯下更多不體面之事，甚至提及都是不合法的。因為他們實際上不再知道自己受什麼苦，也不知道自己在做什麼。\par

但是，那戒绝奢侈的人却不如此：相反，他安坐于港湾，目睹他人的船难，享有纯洁而持久的快乐，追随那适合自由人的生活。因此，知道这些事，让我们逃避奢侈的邪恶宴席，紧守俭朴的餐桌；这样，我们身心都处于良好的状态，就能实践一切德行，并获得未来的美善，借着我们的主\Person{耶稣基督}的恩宠和慈悲，愿光荣、权能、尊崇归于祂、及圣父、及圣神，自今至永远，及世之世。阿们。\par

\SourceRecord{Translated by Talbot W. Chambers.；Nicene and Post-Nicene Fathers, First Series；Vol. 12.；Edited by Philip Schaff.；Buffalo, NY: Christian Literature Publishing Co.,；1889.；<http://www.newadvent.org/fathers/220139.htm>.}

% structure: p0001 source=h1 target=section reason=page-title

\section{\WorkTitle{格林多前书讲道词 第40篇}}

% structure: p0002 source=h2 target=subsection reason=relative-heading-rank; base=h2

\subsection{格前 15:29}

\BibleQuote{否则，那些为死人受洗的，将作什么呢？如果死人总不复活，为什么还为他们受洗呢？}\par

他再一次着手另一个主题，有时从天主所作的来证实他所说的话，有时从他们所行的事来证实。对于任何论点的辩护来说，这也不是一个小理由，如果人把反对者本身带上来，叫他们用自己的行动来证明他所肯定的事。那么，他所说的是什么意思呢？或者，你们愿意我先提一下，那些受\OriginalTerm{马尔西翁派异端}{Marcionite heresy}影响的人是如何曲解这句话的？我确实知道，我会引起很多笑声；然而，正是为此，我最要提出来，使你们更完全地避免这种弊病：即当他们中间有慕道者去世时，他们把活人藏在死者的床下，然后走近尸体与他谈话，问他是否愿意受洗；当他没有回答时，藏在下面的人代替他说，他当然愿意受洗；于是他们代替已亡者给他施洗，如同在舞台上开玩笑的人一样。魔鬼对粗心罪人的灵魂有如此大的力量。然后受质询时，他们援引这句话，说连宗徒也说过：\BibleQuote{那些为死人受洗的。}你们看见他们极端可笑吗？对这些事还值得回答吗？我想不必；除非必须对疯人谈论他们疯狂中所说的话。但为了使更单纯的人不被迷惑，必须屈尊回答这些人。例如，如果这是保禄的意思，为什么天主威胁那不受洗的人呢？因为从此以后不可能有人不受洗了，既然这已被编造出来；而且，过错不再在于死人，而在于活人。但是，他对谁说了：\BibleQuote{你们若不吃我的肉，并喝我的血，在你们内便没有生命。}\BibleRef{若 6:53}是对活人还是对死人，请告诉我？又说：\BibleQuote{人除非由水和圣神而生，不能进天主的国。}\BibleRef{若 3:5}如果这是允许的，而且不需要领受者的心意或他活着时的赞同，那么，有什么阻止希腊人和犹太人都这样成为信徒呢？其他人可以在他们死后代替他们做这些事吗？\par

但为了避免无益的辛劳，就像割断他们那些微不足道的蜘蛛网一样，来，让我们为你们解释这句话的含义。那么，保禄在说什么呢？\par

2. 但首先，我愿意提醒你们这些已受洗的人，关于那晚引导你们进入奥迹的人要你们所作的回答；然后，我也要解释保禄的话：这样，对你们来说也会更清楚。我们在其他一切之后，加上保禄现在所说的这句话。我确实愿意明确地说出来，但为了未入门者的缘故，我不敢；因为这些人给我们的讲解增加了困难，迫使我们要么不清晰地讲，要么向他们宣布那不可言喻的奥迹。然而，我将尽我所能，如同透过一层帷幕来说话。\par

例如：在宣读了那些神秘而可畏的话语，以及那从天上降下的可怕教义规则之后，我们在即将施洗时又加上这一句，叫他们说：\Quote{我信死人的复活，}并在这信仰上我们受洗。因为在我们与其余的一起宣认了之后，最后我们才降入那神圣水泉之中。因此，保禄提醒他们，说：\BibleQuote{如果没有复活，你们为什么还为死人受洗呢？}即为了死去的身体。因为实际上，你们受洗正是为了这个，为了你们死去的身体的复活，相信它不再留在死亡中。你确实在言语中提到了死人的复活；但司铎，如同在一种图像中，以真实的行为向你表明你在言语中所相信和宣认的事。当你没有征兆而相信时，他就也给你征兆；当你尽了你自己的部分，天主也就完全保证。如何并以何种方式？通过水。因为受洗、浸入然后升起，是下降到\OriginalTerm{阴府}{Hades}并从那里返回的象征。因此，保禄也称洗礼为埋葬，说：\BibleQuote{我们藉着洗礼已归于死亡与他同葬了。}\BibleRef{罗 6:4}藉此，他也使将来之事变得可信，我指的是我们身体的复活。因为赦免罪过比复活身体大得多。基督在宣布这一点时说：\BibleQuote{说：你的罪赦了，或说：起来行走，那一样更容易呢？}\BibleRef{玛 9:5}\Quote{前者是更难的，}祂说，\Quote{但因为你们不信它，以为它是隐藏的，而用较容易的事代替较难的事来证明我的能力，我也不拒绝给你们这个证据。}然后祂对瘫子说：\BibleQuote{起来，拿起你的床，回家去罢。}\par

\Quote{这有什么难的呢？}有人说，\Quote{当这也可能发生在君王和统治者身上时？因为他们也赦免奸淫者和杀人者。}你是在开玩笑，人啊，说这些话。因为赦罪唯独天主才有可能。但统治者和君王，无论是他们赦免奸淫者还是杀人者，确实免除了他们现世的惩罚；但并未涤除他们的罪。即使他们提拔被赦免者到官职，即使给他们穿上紫红袍，即使将冠冕戴在他们头上，他们也只会使他们成为君王，却不能使他们从罪中获得释放。唯独天主能成就此事；因此，在\OriginalTerm{重生之洗}{Laver of Regeneration}中，祂将使之实现。因为祂的恩宠触及灵魂的深处，从那里将罪连根拔起。这就是为何被君王赦免的人，其灵魂仍可看出不洁；但受洗者的灵魂不再如此，反而比阳光更纯洁，并且如同起初被造时的状态，甚至比那更好。因为它由一位圣神所祝福，从各方面点燃它，使其圣德炽烈。如同当你重铸铁或金时，使它们变得纯净而崭新；同样，圣神在圣洗中如同在熔炉中重铸灵魂，烧毁其罪过，使其比一切最纯的金子更加闪耀。\par

此外，我们身体复活的可靠性，他又从下文中向你表明：既然罪带来了死亡，如今既然根已枯干，此后就不应再怀疑果实会被毁灭。因此，他先提到\Quote{赦罪}，你随后也宣认\Quote{死者复活}；一个引导你，如同用手牵着你，走向另一个。\par

然而，因为\Quote{复活}一词不足以表明全部：因为许多人复活后又离去了，如旧约中的那些人，如拉匝禄，如基督受难时复活的那些人；因此，他被命令说：\Quote{并永生}，以便没有人再以为那复活之后还会有死亡。\par

因此，保禄将这番话提醒他们，说：\Quote{那些为死者受洗的，将做什么呢？}\Quote{因为如果没有复活，}他说，\Quote{这些话不过是空谈。如果没有复活，我们怎能说服他们相信我们并不赐予的东西呢？}如同一个人吩咐另一人交付一份文件，表明他已收到某笔款项，却从未支付文件中所写金额，然而在签字后却向对方追索指定金额。那么，签字者既然已使自己承担责任，却没有收到他所承认已收到的款项，他还能做什么呢？这里，保禄对受洗者也是这样说。\Quote{那些受洗的，}他说，\Quote{既然签署了承认死人复活，却没有得到它，反而受骗，他们将做什么呢？如果事实没有随之发生，这宣认又有什么必要呢？}\par

% structure: p0012 source=h2 target=subsection reason=relative-heading-rank; base=h2

\subsection{\BibleRef{格前 15:30-31}}

3. \Quote{为什么我们也每时每刻身处危险？}\par

\Quote{我指着我在我们主耶稣基督内以你们为荣的那夸耀，我天天冒死。}\par

请看，他又是如何竭力从自己的见证来确立这教义的；或者不如说，不仅从他自己的，也来自其他宗徒的见证。而这也不是小事：你们所提出的导师们充满了强烈的确信，不仅用言语，也用行为表明这确信。因此，你看，他没有简单地说\Quote{我们确信}，因为仅此不足以说服他们，他还用事实提供证据；好像他应该说：\Quote{用言语承认这些事，对你们来说或许并不奇妙；但如果我们也向你们呈现行为所发出的声音，你们对此又能说什么呢？请听，我们的危险如何日复一日地宣认了这些事。}他没有说\Quote{我}，而是说\Quote{我们}，把所有宗徒都与他一起，从而既说话谦虚，又增加了话语的可信度。\par

因为你们能说什么呢？可能说我们在传讲这些事时欺骗你们，我们的教义来自虚荣？不，我们的危险不允许你们如此判断。因为谁会选择无目的地、无效果地持续身处危险呢？因此他也说：\Quote{为什么我们也每时每刻身处危险？}因为即使有人可能因虚荣而选择这样做，那也只是一次两次，不会像我们一样终生如此。因为我们已将全部生命都为此目的而奉献。\par

\Quote{我指着我在我们主耶稣基督内以你们为荣的那夸耀，我天天冒死：}这里的\Quote{夸耀}，意指他们的进步。因为既然他暗示他的危险很多，唯恐他这样说像是在哀叹，他说：\Quote{我远非悲伤，反而因为了你们的缘故受苦而夸耀。}而且，他说，他为此双重喜悦：既因为了他们而身处危险，也因看到他们的进步。然后，他做了他惯常的事：因为他说了伟大的事，便将一切都归于基督。\par

但他如何\Quote{天天冒死}呢？通过对那事件的准备和预备。他为何说这些话？也是为了通过这些事确立复活的教义。\Quote{因为若不是在此生之后有复活和生命，谁会选择经历如此多的死亡呢？是的，即使相信复活的人也很少会为此置身危险，除非他们心灵极其高尚；更不用说非信徒（他这样说）会选择经历如此多、如此可怕的死亡了。}请看，他如何一步步攀登到更高处。他曾说\Quote{我们身处危险}，又加了\Quote{每时每刻}，然后\Quote{天天}，然后他说：\Quote{我不但\InnerQuote{身处危险}，而且\InnerQuote{死亡}。}最后还指出了他们是怎么样的死亡；这样说：\par

% structure: p0019 source=h2 target=subsection reason=relative-heading-rank; base=h2

\subsection{\BibleRef{格前 15:32}}

\Quote{如果我按人的方式在厄弗所与野兽搏斗，那对我有什么益处呢？}\par

\BibleQuote{若按人的方式？} \BibleQuote{就人而论，我与野兽搏斗；因为天主若把我从那些危险中救出，又如何？所以我正是最应关心这些事的人；我，忍受如此大的危险，尚未得到任何回报。因为如果补偿的时刻并不临近，而我们的赏报被封闭在今世，那么我们的损失更大。因为你们没有冒险就相信了，但我们是天天被杀。}\par

但他说这一切，并非因为他即使在受苦中也没有益处，而是因为许多人的软弱，并为了坚定他们在复活道理上：并非他自己为雇价奔跑；因为做天主喜欢的事，对他来说就是足够的报酬。所以，当他加上说：\BibleQuote{如果我们只在今生寄望于基督，我们就是众人中最可怜的了。}他再次为他们这样说，为的是借这悲惨的恐惧，推翻他们对复活的不信。并且俯就他们的软弱，他如此说。因为事实上，巨大的赏报是时时中悦基督；而且除报酬之外，为祂的缘故冒险本身就是一个极伟大的报答。\par

4. \BibleQuote{如果死人不复活，我们就吃喝吧，因为明天我们就要死了。}\par

这句话，无疑，是嘲讽说的：因此他也不是自己提出来的，而是召来了声音最高昂的先知依撒意亚，他在谈论某些麻木不仁和堕落的人时用了这些话：\BibleQuote{他们宰牛杀羊，吃肉喝酒，说：\InnerQuote{我们吃喝吧，因为明天我们就要死了。}这些事已启示给万军的上主耳中，这罪孽必不得赦免，直到你们死亡。}\BibleRef{依 22:13}如果当时那些这样说的人被剥夺了宽恕，那么在恩宠时期就更甚了。\par

然后，为使他的讲论不过于严厉，他没有长时间停留在他那\Quote{\Emph{reduc} \Emph{tio ad absurdum}，}而是再次将他的讲论转向劝勉，说：\par

% structure: p0026 source=h2 target=subsection reason=relative-heading-rank; base=h2

\subsection{格林多前书 15:33}

\BibleQuote{你们不要被欺骗：邪恶的交往会败坏良好的品行。}\par

他说这话，既是为责备他们好像没有理智（因为在这里，他用了仁慈的说法，把容易受骗的称为\Quote{良善的}），也是尽可能为以往的事对他们稍加原谅，好使他们回头，并把他大部分的责备从他们身上转移到别人身上，并借此吸引他们悔改。他在\WorkTitle{迦拉达书}中也同样做，说：\BibleQuote{那搅扰你们的，无论他是谁，必承当他的审判。}\BibleRef{迦 5:10}\par

% structure: p0029 source=h2 target=subsection reason=relative-heading-rank; base=h2

\subsection{格林多前书 15:34}

\BibleQuote{你们要公义地醒来，不要犯罪。}\par

好像他在对醉汉和疯子说话。因为突然把一切从他们手中抛开，是醉汉和疯子的事，他们不再看见他们所看见的，也不相信他们先前所承认的。但什么是\Quote{公义地？}是着眼于有益有用的事。因为可能不公义地醒来，当一个人完全被唤醒去伤害自己的灵魂。他加上\Quote{不要犯罪}也很好，暗示不信仰的罪由此而来。在许多地方他也暗示这一点，即腐败的生活是邪恶教义的根源；当他这样说：\BibleQuote{贪爱钱财是万恶的根源，有些人追求它，就偏离了信德。}\BibleRef{弟前 6:10}是的，许多意识到自己的邪恶且不愿承担惩罚的人，因这种恐惧也在关于复活的信德上受到损害；而非常行善的人甚至渴望天天看到它。\par

\BibleQuote{因为有些人没有对天主的认识；我说这话是叫你们羞愧。}\par

看，他如何再次把责备转移到别人身上。因为他不是说：\Quote{你们没有知识}，而是说：\Quote{有些人没有知识}。因为不相信复活，是出于一个人没有完全意识到天主的能力是不可抗拒的，并足以成就一切。因为如果祂从无中造出存在之物，那么祂更能使已消散的复活。\par

因为他触到了他们的痛处并大大地嘲弄了他们，指责他们贪食、愚蠢、疯狂；他缓和那些说法，说：\Quote{我说这话是叫你们羞愧}，也就是说，借你们的羞愧，使你们站直、回头、变得更好。因为他怕如果割得太深，会使他们逃避。\par

5. 但让我们不要以为这些话只是对他们说的，也是针对如今所有罹患同样疾病、过着败坏生活的众人。因为事实上，不仅持有错误教义的人酩酊大醉、狂乱，那些被重大罪恶束缚的人也是如此。因此，对他们也大可以正当地说：\Quote{醒寤}，尤其是对那些被贪饕的昏睡所压制、邪恶地掠夺的人。因为有一种掠夺是善的，即掠夺天堂，这并不伤害人。虽然在金钱方面，一人若不为富，除非另一人先变穷；但在属灵之事上却并非如此，恰恰相反：没有人能在不使他人丰裕的情况下致富。因为若你不帮助任何人，你便无法变得富足。因此，在暂时之事上，给予导致减少；在属灵之事上，相反地，给予带来增加，而不给予则带来极大的贫穷和极重的惩罚。那埋藏塔冷通的人所象征的就是这个。是的，那拥有智慧言语的人，通过分施给别人增加自己的富足，使多人成为智者；但将其埋藏在家中的人，因疏忽赚取众人的利益而剥夺了自己的富足。同样，拥有其他恩赐的人，通过治愈多人增进了自己的恩赐：他并未因分施而空虚，反而用自己的属灵恩赐充满了许多人。在一切属灵之事上，这个法则坚定不移。同样，在天国中，那使多人同享天国的人，将因此而更完全地收获其果实；但那不愿有任何人分享的人，自己将从那许多福乐中被驱逐。因为如果这个感官世界的智慧，即使万人强取也不会耗尽；工匠也不会因培养许多工匠而失去自己的技艺；那么，那夺取天国的人更不会使之减少，反而只有当我们为此召叫多人时，我们的财富才会增加。\par

让我们夺取那些不会耗尽、反而在我们夺取时增加的事物；让我们夺取那些无人能通过诬告欺骗我们、无人能因之嫉妒我们的事物。因为倘若有一个地方，有金泉不断涌流，且越汲取越涌流；另有一个地方，埋有宝藏于土中，你愿从哪个致富？难道不是从第一个吗？显然是的。但为使这不只是言辞中的概念，请以空气和太阳为例思考。这些被众人汲取，并满足众人。然而，无论人享用与否，它们依然不变、不减；但我所说的事物则更为伟大：因为属灵智慧在分施与否时并非保持不变，反而在分施中增长。\par

但若有人不能忍受我所说的，仍固守世俗事物的贫穷，攫取那些会减少的事物，那么即使就这些事物而言，也让他想起玛纳食物的例子\BibleRef{出 16:20}，并因那惩罚的榜样而战栗。因为在那个事例中所发生的，如今在贪婪之人身上也能看到。但那时发生了什么？从他们的贪婪中生了虫子。如今在他们身上也是如此。因为食物的度量对众人相同；我们只有一个胃要填满；只是你享乐饮食，有更多要排出的。正如那时，那些在家中收取超过合法数量的人，收取的不是玛纳，而是更多的虫子和腐烂；同样，在奢侈和贪婪中，贪食和醉酒的人收取的不是更多的美味，而是更多的败坏。\par

6. 尽管如此，今人比他们更糟糕，因为他们经历了一次便受了纠正；而这些人每天将远比那更可怕的虫子带进自己家中，却既不察觉，也不满足。因为那些事在无用的辛劳上与此相似（至于惩罚，这些更严重）：这里有证据供你思考。\par

请问，富人与穷人有什么不同？他没有一个身体要穿衣？一个肚腹要喂养？那么他在哪方面有优势？在忧虑中，在自我消耗中，在违背天主中，在败坏肉体中，在浪费灵魂中。是的，这些就是他比穷人优越之处：因为若他有许多胃要填满，或许他还有话可说，比如他的需要更大，花费的必要性也更大。但甚至有人说：\Quote{如今他们可以}，\Quote{回答说，他们填满了许多肚腹，即仆婢的肚腹。}但这并非出于需要，也非出于人道，而是出于纯粹的骄傲；因此无法容忍他们的借口。\par

因为你为什么有许多仆人呢？既然在我们的衣着上，我们只当遵从需要，在我们的餐桌上也是如此，那么在仆人的事上也当这样。那么有什么需要呢？完全没有。因为事实上，一个主人只需要一个仆人；或者更确切地说，两三个主人共用一个仆人。但如果这令你难以忍受，就想想那些没有仆人却得到更及时服侍的人吧。因为天主造人，足以服侍自己，甚至也足以服侍邻人。如果你不信，请听保禄说：\BibleQuote{这双手供应了我的需要，以及和我在一起的人的需要。}\BibleRef{宗 20:34}之后，这位世界的导师、堪当天堂的人，都不耻于服侍无数的他人；而你却认为，除非你带着成群的奴隶，否则就是耻辱，岂不知这恰恰是使你蒙羞的事吗？因为天主赐给我们双手和双脚，正是为了使我们不需要仆人。因为奴隶这一阶层绝不是因需要而被引入的，否则早在亚当之时就造了一个奴隶了；但它是罪的惩罚和不服命的刑罚。然而当基督来时，祂也终结了这事。\BibleQuote{因为在基督耶稣内，没有奴隶或自由人。}\BibleRef{迦 3:28}所以没有必要拥有奴隶；如果实在有必要，也只需一个，至多两个。那些成群的仆人是什么意思呢？因为如同卖羊的和贩奴者，我们的富人也在澡堂和广场上四处寻觅。\par

然而，我不愿过于苛求。我们允许你保留第二个仆人。但如果你收集许多仆人，你这样做并非出于人道，而是出于放纵。因为如果是出于照顾他们，我吩咐你不要让他们任何人服侍你，而是当你买下他们，教他们手艺以维持生计时，就让他们自由。但当你鞭打他们、给他们戴上锁链时，那就再也不是人道之举了。\par

我知道我使听众感到厌恶。但我必须做什么呢？因为我被安置于此，我将不会停止说这些话，无论它们是否产生效果。你在市场里为自己清道是什么意思呢？难道你是在野兽中行走，以至于要驱赶那些遇见你的人吗？不要害怕，那些靠近你、走在你身边的人没有一个会咬你。但你认为与其他人并肩行走是一种侮辱吗？这是何等的疯狂，何等的惊人愚蠢！当一匹马紧跟着你时，你并不认为它给你带来任何侮辱；但如果是一个人，除非他被赶出一百英里以外，你就认为他羞辱了你。而且你为什么还有仆人手持束棒，把自由人当作奴隶使用，或者更确切地说，你自己活得比任何奴隶都更不体面呢？因为事实上，比任何仆人都卑贱的，正是那带着如此多骄傲的人。\par

因此，那些把自己奴役于这种痛苦激情的人，将看不到真正的自由。不，如果你必须驱赶和清除，那么不要驱赶那些靠近你的人，而应当驱除你自己的骄傲；不是通过你的仆人，而是通过你自己；不是用这鞭子，而是用属神的鞭子。因为现在你的仆人驱赶你旁边行走的人，而你自己却被自己的任性更耻辱地从高位上驱赶下来，比任何仆人驱赶你的邻人更甚。但如果你从马上下来，借着谦卑驱除骄傲，你就坐得更高，将自己置于更大的尊荣中，不需要仆人来作此事。我是说，当你变得谦逊，行走在地上时，你将坐在那辆载你上达高天的谦卑之车上，那辆有翼之车；但如果你从它上面坠落，进入骄傲之车，你将与那些被拖在地上行乞的人无异，甚至比他们更可怜、更可悲：因为是他们身体的缺陷迫使他们被拖行，而你是你自身骄傲的疾病。\BibleQuote{凡高举自己的，}他说，\BibleQuote{必被贬抑。}\BibleRef{玛 23:12}为了使我们不被贬抑，反而被高举，让我们趋向那高举吧。因为这样我们也必按照天主的圣言\BibleQuote{找到灵魂的安息}，并获得真实且最崇高的尊荣；愿我们众人，藉着恩宠和慈悲等等，获得这尊荣。\par

\SourceRecord{Translated by Talbot W. Chambers.；Nicene and Post-Nicene Fathers, First Series；Vol. 12.；Edited by Philip Schaff.；Buffalo, NY: Christian Literature Publishing Co.,；1889.；<http://www.newadvent.org/fathers/220140.htm>.}

% structure: p0001 source=h1 target=section reason=page-title

\section{格林多前书第41篇讲道}

% structure: p0002 source=h2 target=subsection reason=relative-heading-rank; base=h2

\subsection{格林多前书 15:35-36}

但有人会说：\BibleQuote{死人怎样复活？他们带着什么样的身体来呢？愚昧的人哪，你所种的若不死，就不会活过来。}\par

宗徒虽然处处极其温和谦卑，但在这里因反对者的不敬，采用了颇为尖刻的语气。但他并不满足于此，还运用了理由和例证，以此制服甚至最争辩的人。而且上面他说：\BibleQuote{既然死是因一人而来，死者的复活也是因一人而来}；但在这里他解决了一个外邦人提出的异议。再看他是如何再次缓和其责备的强烈程度；他没有说：\Quote{但也许你要说}，而是不确定地列出反驳者，以便尽管完全自由地使用他激烈的风格，也不至于过分伤害听众。他提出了两个难点：一个是关于复活的方式，另一个是关于复活身体的性质。因为他们两方面都提出了疑问，说：\BibleQuote{那已被分解的怎样复活呢？}以及\BibleQuote{他们带着什么样的身体来呢？}但\BibleQuote{带着什么样的身体}是什么意思？好像他们说：\Quote{带着这个已经消耗、已经朽坏的身体，还是带着别的身体呢？}\par

然后，为了指出他们所探讨的对象并非有争议而是公认之点，他立刻更尖锐地反驳他们，说：\BibleQuote{愚昧的人哪，你所种的若不死，就不会活过来。}我们在面对那些否认公认之事的人时也惯常这样做。\par

2. 为什么他没有立刻诉诸天主的能力？因为他是在与不信者辩论。因为当他向信友发言时，他不太需要理由。因此，在别处说过：\BibleQuote{祂要改变我们卑贱的身体，使之成为与祂光荣的身体相似的。}\BibleRef{斐 3:2}而且指示了比复活更多的事，他没有陈述类比，而是代替任何证明，提出了天主的能力，接着又说：\BibleQuote{依照祂能使万有归顺自己的德能。}但在这里他也提出了理由。也就是说，从圣经确立之后，他还加上后面这些额外的话，顾及那些不顺从圣经的人；他说：\BibleQuote{愚昧的人哪，你所种的：}意思是：\Quote{你自己从这些事得到证明，凭你每天所做的事，你还怀疑吗？因此我称你为愚昧，因为你对自己每天所做的事无知，而且你本身是复活的制造者，却怀疑天主。}因此他非常强调地说：\BibleQuote{你所种的，}你这属死、属朽坏的人。\par

并且看他如何使用适合他所怀目的的表达：因此他说：\BibleQuote{不会活过来，}又说：\BibleQuote{除非死了。}你看，他舍弃了适合种子的术语，如\Quote{发芽}、\Quote{生长}、\Quote{腐烂}，而采用了与我们的肉身相应的术语，即\BibleQuote{活过来}和\BibleQuote{除非死了}；这些术语本来不属于种子，而是属于身体。\par

他没有说：\Quote{死后它活着，}而是说了一件更伟大的事：\Quote{所以它活着，因为它死了。}你看到吗？我常注意到，他不断地将他们的论证颠倒过来。因此，他们以为我们不再复活的确实标记，他反而用来说明我们确实复活。因为他们说：\Quote{身体不再复活，因为它死了。}那么他怎样反驳他们的论证呢？他说：\Quote{不，除非它死了，它就不能复活；因此它复活了，因为它死了。}正如基督更清楚地表明这件事，说：\BibleQuote{一粒麦子若不落在地里死了，无论有多少单独存在；但若是死了，就结出许多子粒来。}\BibleRef{若 12:24}保禄也从这里吸取了这个例子，他没有说：\Quote{它不能活，}而是说：\BibleQuote{不能活过来；}再次假定天主的能力，并表明不是土地的本性，而是天主亲自成就了一切。\par

那么，他为何没有提出那个更贴近主题的例证呢？我是说，人类的种子？（因为我们的生育也是从某种腐朽开始，如同谷粒一样。）因为它不具同等说服力，而后者是个更完全的例证：因为他需要一个完全朽坏的事物为例，而前者只是部分；因此他宁可举出另一个。此外，前者是从有生命的身体发出，落入有生命的子宫；但这里，不是肉体，而是土地，种子被投入其中，并在其中分解，如同已死的身体。因此，就这一点而言，那个例证也更合适。\par

% structure: p0010 source=h2 target=subsection reason=relative-heading-rank; base=h2

\subsection{格林多前书 15:37}

3. \BibleQuote{并且播种的人所播种的不是将来要出现的身体。}\par

因为前面说的话回应了\Quote{他们怎样复活}的问题；而这里，回应了\Quote{他们带着什么样的身体来}的疑惑。但\BibleQuote{你所播种的不是将来要出现的身体}是什么意思？不是完整的谷穗，也不是新的谷粒。因为这里他的话语不再涉及复活事实本身，而是复活的方式，即复活的身体是什么性质；是与原先同类，还是更好、更光荣。他从同一个类比中引出两者，暗示复活的身体要更好得多。\par

但异端者毫不考虑这些事，就冲进我们当中说：\Quote{一个身体倒下，另一个身体又复活。那么怎能说是复活？因为复活的是那个倒下的。但如果一个身体倒下而另一个复活，那战胜死亡的奇妙惊人战利品在哪里？因为他似乎不再归还他所掳获的。而且，所谓的类比如何与前述之事相符？}为什么，被撒下的不是一种本体，被举起的另一种，而是同一本体得到了提升。否则，基督在成为复活初果时也不会重取同一身体；但依你们，祂抛弃了先前的身体（虽然它没有犯罪）而取了另一个。那么另一个从何而来？因为这一身体来自童贞女，而那一身体来自哪里？你看见这论证已陷入何等荒谬吗？因为祂展示钉痕本身是为了什么？岂不是为了证明被钉的是同一身体，而同一身体又从死者中复活？而且祂的约纳预表又是什么意思？因为被吞下的不是一个约纳，而被抛在旱地上的另一个。祂又说\Quote{拆毁这座圣殿，三天内我要把它重建起来}是什么意思？因为那被拆毁的，显然祂又把它重建了。因此，福音作者也补充说\Quote{祂是指祂身体的圣殿。}\BibleRef{若 2:19}\par

那么他说\Quote{你所撒下的不是那将来要有的形体}是什么意思？即不是谷穗：因为它是同一个又不是同一个；同一个，因为本体相同；不是同一个，因为这一形体更卓越，本体保持不变而美丽变得更大，且同一身体崭新地复活。因为若非如此，就不需要复活了，我的意思是如果它不是为了提升而复活。因为祂为什么全然拆毁祂的房屋，除非祂要把它建造得更荣耀？\par

你看，这是他对认为它是完全朽坏的人说的。接着，免得有人从这里猜疑是另一个身体在说话，他修饰了这隐晦的说法，并亲自解释他所说的话，不允许听众由此转向任何其他方向。那么还需要我们的推理做什么？听他亲自说话并解释\Quote{你所撒下的不是那将来要有的形体}这话。因为他立刻加上\Quote{乃是赤裸的籽粒，或许是麦子，或其他种子}；即它不是那将来要有的形体；例如，不是这样穿着，没有秆和穗，而是\Quote{赤裸的籽粒，或许是麦子，或其他种子}。\par

% structure: p0016 source=h2 target=subsection reason=relative-heading-rank; base=h2

\subsection{格林多前书 15:38}

\Quote{但天主随祂意愿给籽粒一个形体。}\par

有人说：\Quote{是的，}但依你看来，\Quote{但在那种情况下是自然的运作。}告诉我，是哪种自然的运作？因为在那情况下也同样是天主做了全部；不是自然，不是土地，不是雨水。因此，使这些显明的事，他省略了土地、雨水、大气、太阳和农人的手，而加上\Quote{天主随祂意愿给籽粒一个形体。}所以当你听到天主的权能和旨意时，不要好奇追问，也不要忙着查究如何和以什么方式。\par

\Quote{各给每种籽粒自己的形体。}那么他们所说的外来质料在哪里？因为他给每种\Quote{它自己的。}所以当他说\Quote{你所撒下的不是那将来要有的形体}时，他不是说一个本体被举起代替另一个，而是说它被提升、更荣耀。因为他说：\Quote{各给每种籽粒自己的形体。}\par

4. 从接下来他引入那时将要有的复活的差别。因为不要以为籽粒被撒下而都长成谷穗，所以在复活中也有同等的荣耀。因为首先，即使在籽粒中也不是只有一个等级，而是有些更有价值，有些低劣。因此他也加上\Quote{各给每种籽粒自己的形体。}\par

然而，他不满足于此，而是寻求另一个更大更明显的差别。为使你听见（如我所说）都复活时，不以为众人都享同样的赏报；他在前面的经文中就撒下了这一思想的种子，说\Quote{但各人按着自己的次序。}但他在这里也更清楚地表达出来，说：\par

% structure: p0022 source=h2 target=subsection reason=relative-heading-rank; base=h2

\subsection{格林多前书 15:39}

\Quote{凡有血肉的，不是同样的血肉。}因为他说：何必论籽粒呢？关于身体，让我们讨论这一点，就是我们正在谈论的。因此他又加上说：\par

\Quote{人的血肉是一种，牲畜的血肉是另一种，鸟类的血肉是另一种，鱼类的血肉又是另一种。}\par

% structure: p0025 source=h2 target=subsection reason=relative-heading-rank; base=h2

\subsection{格林多前书 15:40}

\Quote{也有天上的形体，也有地上的形体；但天上的荣耀是一种，地上的荣耀是另一种。}\par

% structure: p0027 source=h2 target=subsection reason=relative-heading-rank; base=h2

\subsection{格林多前书 15:41}

\Quote{日有日的荣耀，月有月的荣耀，星有星的荣耀；因为此星与彼星在荣耀上也有分别。}\par

祂用这些言辞表示什么呢？祂为何从身体复活的话题转到谈论星体和太阳？祂并非转移话题，也非偏离目的；绝非如此：祂仍紧扣主题。因为祂既已确立关于复活的道理，便在接下来的话中暗示：虽然只有一次复活，但将来的光荣会有极大的差别。目前祂将整体分为两类：\Quote{属天的身体}和\Quote{属地的身体}。祂以种子表明身体得以复活；但祂以此表明并非所有人都享有同样的光荣。因为对复活的不信使人懈怠，而以为所有人都获得同样赏报的想法同样使人懒惰。因此祂纠正了这两种错误。祂在前几节中已经完成了前者，但现在开始处理后者。祂将义人和罪人分为两等，又将这两等再细分为许多部分，表明义人和罪人不会得到同样的赏报；义人之间也并非全然相同，罪人之间也是如此。\par

你看，祂首先在义人和罪人之间作出区分，说：\Quote{属天的身体}和\Quote{属地的身体}：以\Quote{属地的}暗示后者，以\Quote{属天的}暗示前者。然后祂进一步区分罪人与罪人，说：\Quote{凡肉体不都是同样的肉体；人的肉体是一种，鱼类的肉体是另一种，鸟类的肉体是另一种，兽类的肉体是另一种。}虽然都是身体，但有些更为卑贱，有些则不那么卑贱。在生活方式乃至构造上也是如此。\par

说了这些之后，祂又升到天上，说：\Quote{太阳的光荣是一种，月亮的光荣是另一种。}因为在属地的身体中有差别，在属天的身体中也有差别；并且这差别并非寻常，而是极其巨大：不仅太阳和月亮、星辰之间有差别，就是星辰与星辰之间也有差别。它们虽然都在天上，但有些享有更大份的光荣，有些则享有较小的。我们由此学到什么？尽管所有人都将在天主的国里，但并非所有人都享有同样的赏报；尽管所有罪人都将在地狱里，但并非所有人都承受同样的刑罚。因此祂接着说：\par

% structure: p0032 source=h2 target=subsection reason=relative-heading-rank; base=h2

\subsection{格前 15:42}

\Quote{死人的复活也是这样。}\par

\Quote{这样，}怎样？带着相当大的差别。然后祂离开这个道理，认为它已充分证明，又回到复活本身的证明及其方式，说：\par

5. \Quote{播种的是可朽坏的，复活起来的是不可朽坏的。}请注意祂的谨慎。如同种子的情形，祂用了属于身体的术语，说：\Quote{若不是死了，就不能获得生命；}同样在身体的情形，祂用了属于种子的说法，说：\Quote{播种的是可朽坏的，复活起来的是不可朽坏的。}祂没有说\Quote{生长出来}，免得你以为那是大地的工作，而是说\Quote{复活起来}。祂在这里所说的播种，不是指我们在母胎中的受生，而是指我们死去的身体被埋葬在地里、它们的朽坏和灰烬。因此，在说了\Quote{播种的是可朽坏的，复活起来的是不可朽坏的}之后，祂又加上：\par

% structure: p0036 source=h2 target=subsection reason=relative-heading-rank; base=h2

\subsection{格前 15:43}

\Quote{播种的是可羞辱的。}因为有什么比朽坏的尸体更难看呢？\Quote{复活起来的是光荣的。}\par

\Quote{播种的是软弱的。}因为不到三十天，整个身体就消失了，血肉不能维系自身，连一天也维持不住。\Quote{复活起来的是强而有力的。}因为将来没有任何事物能胜过它。\par

这就是为什么祂需要先前的那些类比，免得许多人听到这些事，即身体\Quote{以不可朽坏、光荣和强而有力的方式复活起来}，就以为所有复活的人之间没有差别。因为所有人确实都复活起来，既有能力又不可朽坏；但在这种不可朽坏的光荣中，并非所有人都处于同样的尊荣和安全状态。\par

% structure: p0040 source=h2 target=subsection reason=relative-heading-rank; base=h2

\subsection{格前 15:44}

\Quote{播种的是属灵魂的身体，复活起来的是属神的身体。既有属灵魂的身体，也有属神的身体。}\par

你说什么？难道\Quote{这个}身体不是属神的吗？它确实是属神的，但那时将更是如此。因为现在，当人犯了大罪时，圣神的丰沛恩宠常常离去；再者，当圣神继续临在时，血肉的生命依赖于灵魂：结果便是一个空洞，没有圣神。但在那一天并非如此；相反，祂永恒地住在义人的血肉中，胜利将属于祂，属灵魂的灵魂也存在。\par

因为祂借\Quote{属神的身体}所暗示的，或许是这类事，或是那身体将更轻、更精细，甚至能在空中飘浮；或者祂两者都意指。如果你不相信这个道理，请看那些天上的身体，它们如此光荣，（暂时）如此耐久，在永不朽坏的安宁中存留；你要从这里相信，天主也能使这些可朽坏的身体成为不可朽坏的，并且比那些可见的身体更卓越。\par

% structure: p0044 source=h2 target=subsection reason=relative-heading-rank; base=h2

\subsection{格前 15:45}

6. \Quote{也这样记载着：\BibleRef{创 2:7}\InnerQuote{第一个人亚当成了有灵的生魂，}最后的亚当成了赐生命的圣神。}\par

然而，前者是记载的，后者却没有记载。那么祂怎么说\Quote{记载着}呢？祂按照事件的结果修改了表达方式：这是祂惯常的做法，也是每一位先知的做法。例如，先知说耶路撒冷应被称为\Quote{正义之城}\BibleRef{依 1:26}，但它从未被称为这样。那么，先知说谎了吗？绝没有。因为他是在说事件的结果。同样，基督应被称为\Quote{厄玛奴耳}\BibleRef{依 7:14}，但祂从未被称为这样。事实却发出了这个声音；所以这里也一样，\Quote{最后的亚当成了赐生命的圣神。}\par

他说这些话，为叫你明白，今生和来生的标记和抵押都已经临到我们身上；即今生的\Person{亚当}，和来生的\Person{基督}。因为他把更美的事作为\OriginalTerm{望德}{hope}之物，就表明它们的开端已经发生，它们的根和泉源已经显明。但如果根和泉源对众人都明显，就不必怀疑果实。因此他说：\Quote{最后的\Person{亚当}成了赋予生命的\OriginalTerm{圣神}{Spirit}。}在别处他也说：\Quote{祂必借着住在你们内的\OriginalTerm{圣神}{Spirit}，使你们有死的身体活起来。}\BibleRef{罗 7:11}可见，使人活是\OriginalTerm{圣神}{Spirit}的工作。\par

此外，免得有人说：\Quote{为什么较坏的事是更早的？为什么一种，即属自然的，不仅作为初果，而是完全发生；另一种却只作为初果？}——他表明每样事的原则也是如此安排的。\par

% structure: p0049 source=h2 target=subsection reason=relative-heading-rank; base=h2

\subsection{\BibleRef{格前 15:46}}

他说：\Quote{那不是在先的，是属\OriginalTerm{圣神}{Spirit}的；而是属\OriginalTerm{自然}{natural}的在先，然后才是属\OriginalTerm{圣神}{Spirit}的。}\par

他并不问为什么，而是满足于\Person{天主}的安排，以事实的见证证明\Person{天主}绝妙的经纶，并暗示我们的状况总是在进步之中；同时以此也增加他论点的可信度。因为如果较小的已经发生，我们更应当期待那更好的。\par

7. 既然我们要享受如此大的福分，让我们站在这阵列中，不要哀悼去世的人，而要哀悼那些生命终结不善的人。因为农夫看到种子溶解时，并不哀伤；相反，只要他看到种子在地里保持坚实，他就恐惧战栗，但当看到它溶解时，他就喜乐。因为未来庄稼的开始就是它的溶解。所以我们也应当欢喜当可朽坏的房屋倒下，当人被播种时。不要惊奇他把埋葬称为\Quote{播种；}因为，事实上，这是更好的播种：因为那次播种之后是死亡、劳苦、危险和忧虑；但这一次，如果我们生活得好，就是冠冕和赏报；那一次是败坏和死亡，这一次是不朽和永生，以及那些无限的福分。那种播种伴随拥抱、快乐和睡眠；但这播种，只有从天降下的一声声音，一切立刻臻于完美。那复活的人不再被引向充满辛劳的生活，而是到了痛苦、忧愁和叹息都已逃离的地方。\par

如果你需要保护，因此哀悼你的丈夫，就投奔到\Person{天主}那里，他是众人共同的\OriginalTerm{保护者}{Protector}、\OriginalTerm{救主}{Saviour}和\OriginalTerm{恩主}{Benefactor}，投向那不可抵抗的联盟，那随时可得的援助，那永存的避难所，它无所不在，且从四方做我们的城墙。\par

\Quote{但你的交往是可羡慕和可爱的。}我也知道。但如果你用正确的理智来对待这悲伤，并思考是谁带走了他，而且通过高尚的忍受，你将你的心灵作为祭献献给我们的\Person{天主}，那么这波浪也不会强到使你无法抵挡。时间所带来的，你也要通过自制来实现。但如果你屈服于软弱，你的情感固然会随时间消失，但不会给你带来赏报。\par

同时，连同这些理由，也收集一些例子，有些在今世，有些在\OriginalTerm{圣经}{Holy Scriptures}中。思量\Person{亚巴郎}杀了自己的儿子，既没有流一滴眼泪，也没有说一句苦毒的话。\Quote{但他，}你说，\Quote{是\Person{亚巴郎}。}不，你确实被召叫到更崇高的事业中。\Person{约伯}固然悲伤，但只是作为一位爱孩子、非常关心逝者的父亲所应有的程度；而我们现在所做的，倒像是仇人和敌人的行为。因为如果一个人被带到法庭并接受加冕，你却捶胸痛哭，我不会说你是那加冕者的朋友，而是大仇人和敌对者。\Quote{不，}你说，\Quote{我哀悼他，并非为他，而是为我自己。}那么，这不是一个慈爱的人所应为，为了自己的缘故，宁愿他仍在冲突中，面对未来的不确定性，而本来他可以加冕并靠岸；或者宁愿他在汪洋中颠簸，而本来他可以在港口中。\par

\Quote{但我不知道他去了哪里，}你说。为什么你不知道，告诉我？因为他生活得好或坏，就显明他要去哪里。\Quote{不，正因此我哀悼，}你说，\Quote{因为他死时是个罪人。}这不过是一个借口和托词。因为如果这就是你为死者哀哭的理由，你本应在他活着的时候塑造和纠正他。事实是你处处只关心自己，而不是他。\par

但假设他带着罪死去，即便为此，我们也应当欢喜，因为他在罪中被制止，没有增加他的过犯；并要尽可能地帮助他，不是用眼泪，而是用祈祷、恳求、施舍和奉献。因为这些事情并非毫无意义地被设计，我们在神圣奥迹中纪念已亡的人，并为他们祈求天主，恳求在我们面前、除免世罪的羔羊——并非徒然，而是为使一些安息因此而临到他们。不是徒然的，当那站在祭台旁的人在举行可敬的奥迹时高呼：\Quote{为所有在\OriginalTerm{基督}{Christ}内安息的人，以及那些为他们举行纪念的人。}因为如果没有为他们举行的纪念，这些话就不会被说；因为我们的礼仪不是纯粹的戏剧表演，绝无此事！是的，这些事是因\OriginalTerm{圣神}{Spirit}的诫命而行的。\par

让我们因此援助他们，并为他们举行纪念。因为如果约伯的子女因父亲的祭献而得洁净，为何你怀疑当我们也为亡者奉献时，他们也会得到一些安慰？因为天主惯于应允那些为他人祈求者的恳求。这正是保禄所暗示的，他说：\Quote{因着多人的代求，使我们为你们所领受的恩赐，因许多人而藉着感谢，归于我们。}\BibleRef{格后1:11} 因此，让我们不要厌倦于以奉献和为他们祈祷来援助亡者：因为世界的共同赎罪祭就在我们面前。因此，我们大胆地为整个世界恳求，并将他们的名字与殉道者、宣信者、司铎的名字一同提及。因为事实上我们都是一个身体，尽管有些肢体比别的更荣耀；并且有可能从各种源头为他们获得赦免：从我们的祈祷，从我们为他们奉献的礼品，从那些与他们一同被提及名字的人。那么，你为何悲伤？为何哀恸，当你有能力为亡者积聚如此多的赦免时？\par

9. 难道你变得孤苦伶仃，失去了保护者吗？不，千万不要这样说。因为你并没有失去你的天主。所以，只要你拥有祂，祂比丈夫、父亲、子女和亲属对你更好：因为即使他们活着时，也是祂成就了一切。\par

这些事你当思想，并同达味说：\Quote{上主是我的光明，我的救援，我还怕谁？}\BibleRef{咏26:1} 你说：\Quote{你是孤儿的慈父，寡妇的审判者，}\BibleRef{咏67:5} 并呼求祂的援助，你必会得到祂的照顾，比从前更多，因你处于更大的困境中。\par

或者你失去了一个孩子？你没有失去他，不要这样说。这件事是睡眠，不是死亡；是迁移，不是毁灭；是从较差到更好的旅途。不要因此激怒天主，而要平息祂。因为如果你高贵地忍受它，那么对亡者和你自己都会带来一些慰藉；但若相反，你会更加点燃天主的怒火。因为如果一个仆人被主人责打时，你站在一旁抱怨，你会更加激怒主人来反对你。所以不要这样，而要感谢，好使这悲伤的云雾也从你身上消散。你要同那位有福者说：\Quote{上主赐的，上主收回。}\BibleRef{约1:21} 要思想，在祂眼中蒙喜悦的人中，有多少从未得到过子女，也从未被称为父亲。你说：\Quote{我但愿从未如此，因为未曾经历比为享乐后又失去更好。}不，我恳求你，不要这样说，也不要这样激怒主：但要为你所领受的感谢祂，为你未得至终的归光荣于祂。约伯并没有像你这样忘恩负义地说\Quote{未曾领受更好}，而是为前者感谢说\Quote{上主赐了}，为后者祝福天主说\Quote{上主收回，愿上主的名永远受赞美。}并且他这样制止了他的妻子，为自己辩白，说出那些令人钦佩的话：\Quote{我们从天主手中得福，难道不也受祸吗？}然而此后，更猛烈的考验临到他；但他甚至没有为此而气馁，反而同样高贵地忍受并光荣了天主。\par

你也要这样做，并思考：不是人取走了他，而是天主，造他的那一位，祂比你更关心他，知道什么对他有益；祂不是敌人，也不是埋伏者。看，多少活着的人使父母的生活难以忍受。你说：\Quote{难道你没有看见那些正直的人吗？}我也看见这些人，但他们还不如你的孩子安全。因为尽管他们目前被认可，但他们的结局如何不可知；至于他，你不再有任何恐惧，也不再担心他会发生什么事或经历任何变化。\par

你也要这样思想有关一位贤妻和你的家室之看守者，并为一切感谢天主。即使你失去一位妻子，也要感谢。也许天主的旨意是引导你走向节欲，祂召唤你进入更高贵的战场，祂乐意让你从这捆缚中解脱。如果我们这样管束自己，我们不仅会获得今生的喜乐，还会获得将来的荣冠，等等。\par

\SourceRecord{Translated by Talbot W. Chambers.；Nicene and Post-Nicene Fathers, First Series；Vol. 12.；Edited by Philip Schaff.；Buffalo, NY: Christian Literature Publishing Co.,；1889.；<http://www.newadvent.org/fathers/220141.htm>.}

% structure: p0001 source=h1 target=section reason=page-title

\section{格林多前书讲道稿第四十二篇}

% structure: p0002 source=h2 target=subsection reason=relative-heading-rank; base=h2

\subsection{\BibleRef{格前 15:47}}

\BibleQuote{第一个人出于地，属于土；第二个人是主，出自天上。}\par

说了\Quote{那属血肉的在先}和\Quote{那属神的在后}之后，他又指出另一种差别，论及\Quote{属土的}和\Quote{属天的}。因为第一种差别是在现世生命与来世生命之间；而这种差别是在恩宠之前与恩宠之后。他这样说，是为了卓越的生活方式——（如我所说，为阻止人们因如此信赖复活而忽略操行与成全，又从这个话题使他们焦虑并劝勉德行，）——说：\Quote{第一个人出于地，属于土；第二个人是主，出自天上。}以\Quote{人}的名称称呼整体，而称呼这一个是出于更好的部分，另一个是出于最坏的部分。\par

% structure: p0005 source=h2 target=subsection reason=relative-heading-rank; base=h2

\subsection{\BibleRef{格前 15:48}}

\Quote{那属土的怎样，凡属土的也怎样；}他们必朽坏，归于终结。\Quote{那属天的怎样，凡属天的也怎样；}他们必存留不朽而荣耀。\par

那么怎样呢？这位不也死了吗？他确实死了，但未受损伤，反而借此消灭了死亡。你是否看见，他在这部分主题上也用死亡来确立复活的道理？因为，如我之前所说，\Quote{有了开端和元首，}他说，\Quote{就不必怀疑整个身体。}\par

此外，他由此也构成了关于最佳生活方式的劝言，提出了高尚严格的生活与并非如此的生活的标准，并引出了两者的原则：一个是基督，另一个是亚当。因此他并非简单地说\Quote{出于地}，而是\Quote{属土的}，即\Quote{粗俗的，被钉在现世事物上；}又关于基督，相反地，\Quote{主出自天上。}\par

但如果有人说，因此主没有身体，因为据说他是\Quote{出自天上}——尽管前面所说已足以堵住他们的口——但没有什么阻止我们由此考虑来使他们沉默：即所谓\Quote{主出自天上}是什么意思？他说的是祂的本性，还是祂最完善的生活？我想每个人都清楚他指的是祂的生活。因此他也加上：\par

% structure: p0010 source=h2 target=subsection reason=relative-heading-rank; base=h2

\subsection{\BibleRef{格前 15:49}}

\Quote{我们既然带了那属土的肖像，}即我们行了恶，\Quote{我们也要带那属天的肖像，}即我们实践一切善。\par

此外，我还想问你们：说\Quote{那出于地的，属于土}和\Quote{主出自天上}，是论及本性吗？有人说：\Quote{是。}那么怎样呢？亚当只是\Quote{属土的}，还是也有另一种与属天和非物质实体类同的质料？圣经称之为\Quote{灵魂}和\Quote{神}？人人都看见他也有这个。因此主也不只是来自高处，尽管祂被说成是\Quote{出自天上}，但祂也取了我们的肉身。但保禄的意思是这样：\Quote{我们既然带了那属土的肖像，}即恶行，\Quote{我们也要带那属天的肖像，}即在天上的生活方式。而如果他在谈论本性，这就不需要劝勉或忠告。所以由此也明显，这话关系到我们的生活方式。\par

因此他也以劝告的方式引入这话，并称之为\Quote{肖像}，在这里再次表明他是在谈论行为，而非本性。因为我们成为属土的，是因为行了恶；并非因我们最初被塑造为\Quote{属土的}，而是因为我们犯了罪。因为罪先来了，然后死亡，然后判决：\Quote{你原是尘土，仍要归于尘土。}\BibleRef{创 3:19}于是各种情欲也成群进入。因为一个人\Quote{出于地}并不简单地使他成为\Quote{属土的}（因为主也来自这同一团泥土），而是由于做属土的事；同样，他也成为\Quote{属天的}，是因为行那适于天国的事。\par

但够了：因为我何必过分劳苦来证明这一点呢？宗徒自己接着向我们阐明这思想，这样说：\par

% structure: p0015 source=h2 target=subsection reason=relative-heading-rank; base=h2

\subsection{\BibleRef{格前 15:50}}

\BibleQuote{弟兄们，我告诉你们，血肉之躯不能继承天主的国。}\par

你是否看见他又解释自己，为我们省去麻烦？这是他常做的事：因为他此处用\Quote{血肉}指人的恶行，正如他在别处也这样做的；例如当他说：\Quote{但是你们并不属于血肉，}以及：\Quote{凡属于血肉的人，都不能中悦天主。}\BibleRef{罗 8:8}所以他说\Quote{现在我告诉你们}时，意思无非就是：\Quote{因此我说了这些话，好使你们知道恶行不引人进入天国。}这样他从复活也立即引入了天国的道理；因此他也加上：\Quote{朽坏的也不能继承不朽坏的，}即邪恶不能继承那荣耀和不朽坏之物的享受。因为他在许多其他地方称邪恶为这名字，说：\Quote{那随从肉体播种的，必由肉体收获败坏。}\BibleRef{迦 6:8}现在如果他是在谈论身体而不是恶行，他就不会说\Quote{朽坏}。因为他绝不以\Quote{朽坏}称呼身体，因为身体不是朽坏，而是可朽坏的：因此当他接下来也论及身体时，他不称它为\Quote{朽坏}，而是\Quote{可朽坏的}，说：\Quote{这必朽坏的必须穿上不朽坏的。}\par

接着，在完成了关于我们生活方式的劝告后，他根据自己惯常的做法，将主题紧密相连，又转到身体复活的道理，如下：\par

% structure: p0019 source=h2 target=subsection reason=relative-heading-rank; base=h2

\subsection{\BibleRef{格前 15:51}}

\BibleQuote{看，我告诉你们一个奥秘。}\par

他即将要说的事是可怕、不可言喻、且众人所不知的；这也表明了他赐予他们的尊荣是何等伟大：我是说，他向他们讲述奥秘。但这奥秘是什么呢？\par

\BibleQuote{我们不全都要长眠，但我们全都要被改变。} 他的意思是：\Quote{我们不全都要死亡，\InnerQuote{但我们全都要被改变}}， 就连那些不死的人也要改变，因为他们也是可朽的。\Quote{因此，你不可因为死亡而害怕，} 他说，\Quote{好像你不会复活似的：因为有一些人甚至能逃脱死亡，然而这并不足以使他们获得复活，就连那些不死去的身体也必须被改变，转化为不朽。}\par

% structure: p0023 source=h2 target=subsection reason=relative-heading-rank; base=h2

\subsection{\BibleBook{格林多}{前书} 15:52}

\BibleQuote{在片刻之间，在一眨眼之间，在最后的号角声中。}\par

在他详细论述了复活之后，他适时地指出了复活极其惊人的特性。正如他所说：\Quote{不仅这一点} 是奇妙的：我们的身体先化为朽坏，然后才复活；也不仅复活后的身体比现今的更好；也不仅它们进入一个更美好的状态；也不仅各人收回自己的，而非别人的身体；而是那么多、那么大、超越人类一切理智和想象的事，竟在\BibleQuote{片刻之间} 成就，即在瞬间之内。为了更清楚地说明这一点，他说：\BibleQuote{在一眨眼之间，} \Quote{就在人眨眼的工夫。} 再者，因为他所说的事极其重大且令人惊叹——如此多、如此大的效果竟如此迅速地发生——他援引行这事者的可信性来证明这一点，如下所述：\BibleQuote{因为号角将吹响，死者将复活成为不朽的，我们也要被改变。} 他使用\Quote{我们} 一词，不是指他自己，而是指那时仍活着的人。\par

% structure: p0026 source=h2 target=subsection reason=relative-heading-rank; base=h2

\subsection{\BibleBook{格林多}{前书} 15:53}

\BibleQuote{因为这可朽的必须穿上不朽。}\par

因此，免得有人听见\BibleQuote{血肉之躯不能承受天主的国} 就以为我们的身体不会复活，他补充说：\BibleQuote{这可朽的必须穿上不朽，这可死的必须穿上不死。} 如今身体是\Quote{可朽的}，身体是\Quote{可死的}：所以身体确实存留，因为穿上它的正是身体；但它的可死性和朽坏性消失了，当不死和不朽临到它身上时。因此，既然你已听到它成为了不朽，就不要再问它怎能活到永远。\par

% structure: p0029 source=h2 target=subsection reason=relative-heading-rank; base=h2

\subsection{\BibleBook{格林多}{前书} 15:54}

4. \BibleQuote{但当这可朽的穿上了不朽，这可死的穿上了不死时，那时就应验了所记载的话：\InnerQuote{死亡被胜利吞灭了。}}\par

因此，既然他在谈论伟大而隐秘的事，他再次引用先知的话\BibleRef{欧 13:14}来证实他的话语。\BibleQuote{死亡被胜利吞灭}，即彻底地吞灭；连一点碎片、一丝复返的希望都不留，因为不朽吞灭了朽坏。\par

% structure: p0032 source=h2 target=subsection reason=relative-heading-rank; base=h2

\subsection{\BibleBook{格林多}{前书} 15:55}

\BibleQuote{死亡啊，你的刺在哪里？坟墓啊，你的胜利在哪里？}\par

你看见他高贵的灵魂了吗？他如同为胜利献祭的人，受到灵感激发，已视未来的事如过去的事，跳起来践踏倒在脚下的死亡，对着躺在下面的死神的头欢呼胜利，大声呐喊说：\BibleQuote{死亡啊，你的刺在哪里？坟墓啊，你的胜利在哪里？} 它已完全消失，已灭亡，已彻底消逝，你从前所做的一切都徒然了。因为祂不仅剥夺了死亡的武装并战胜了它，甚至毁灭了它，使它完全不复存在。\par

% structure: p0035 source=h2 target=subsection reason=relative-heading-rank; base=h2

\subsection{\BibleBook{格林多}{前书} 15:56}

\BibleQuote{死亡的刺就是罪，而罪的能力就是法律。}\par

你看见这段讲论是关于身体的死亡吗？因此也是关于身体的复活。因为如果这些身体不复活，死亡怎会被\Quote{吞灭}？不仅如此，\Quote{法律} 又怎会是\Quote{罪的能力}？因为\Quote{罪} 确实是\Quote{死亡的刺}，比死亡更苦，并且罪借死亡而有能力，这是明显的；但\Quote{法律} 又怎会是罪的能力呢？因为没有法律，罪是软弱的，虽然罪恶发生，却不能如此完全地定人的罪：因为虽然恶事发生，但它没有被如此清楚地指出来。所以，法律所带来的改变不小：首先使我们更清楚地认识罪，继而加重了刑罚。如果法律意欲制止罪恶，却反而更可怕地发展了它，这并非医生的过错，而是滥用药物之过。正如基督的临在加重了犹太人的负担，但我们不可因此责怪基督，反而更应钦佩祂，同时更恨恶那些人，因为他们本应受益的东西反而伤害了他们。是的，为了表明法律本身并没有给罪加添力量，基督自己完全成全了法律，并且没有罪。\par

但我希望你注意，他也是从这一主题证实了复活。因为如果死亡的原因是我们犯罪，而基督来了，除去罪，借着圣洗救我们脱离罪，并且连同罪也一并结束了法律（罪在于违反法律），那么你为什么还要怀疑复活呢？因为此后，死亡从何处得胜？借着法律吗？不，法律已被废除。借着罪吗？不，罪已被彻底消灭。\par

% structure: p0039 source=h2 target=subsection reason=relative-heading-rank; base=h2

\subsection{\BibleBook{格林多}{前书} 15:57}

\BibleQuote{但感谢天主，祂借着我们的主耶稣基督赐予我们胜利。}\par

因为祂亲自立了凯旋柱，却使我们也有分于荣冠。这并非出于债务，而纯粹出于怜悯。\par

% structure: p0042 source=h2 target=subsection reason=relative-heading-rank; base=h2

\subsection{\BibleBook{格林多}{前书} 15:58}

5. \BibleQuote{因此，我所亲爱的弟兄们，你们要坚定不移，不可动摇。}\par

在这一切之后，这劝勉是公正而合时的。因为没有什么比认为我们无端或无益地遭受打击更使人不安的了。\par

\BibleQuote{常常在主的工作上富足有余：} 即在于纯洁的生活。他没有说\Quote{行善}，而是说\Quote{富足有余}，好使我们丰丰富富地行善，并且超越界限。\par

\Quote{知道你们的劳苦在主内不是徒然的。}\par

你说什么？又是劳苦？但随之而来的是冠冕，而且是超乎诸天之上的。因为从前那因人被逐出乐园的劳苦，是他过犯的惩罚；但这劳苦却是将来赏报的缘由。所以，它实在不能算是劳苦，既因这个缘故，也因它所领受的从上而来的巨大帮助：这也是为何他加上\Quote{在主内}的原因。因为前者的目的是要我们受罚，而后者则是要我们获得将来的美善。\par

所以，我亲爱的，我们不要沉睡。因为不可能，绝不可能有人凭着懒惰就进入天国，那些生活奢侈安逸的人也不能。是的，若我们克制自己、\Quote{攻克己身}、忍受无数的劳苦，才能达到那些福乐，这已是了不起的事。你们岂不见天地之间的距离何等遥远？眼前的争战何其巨大？人性多么易于向恶？罪是多么容易\Quote{缠绕我们？}路上有多少陷阱？\par

那么，为什么我们还要在自然的忧虑之外，为自己招来那么多烦恼，给自己增添更多麻烦，使担子更重呢？难道我们为食物、衣服和住房操心还不够吗？难道为必需品思虑还不够吗？虽然基督甚至从这些中引导我们出来，说：\Quote{不要为你们的生命忧虑吃什么，也不要为你们的身体忧虑穿什么。}\BibleRef{玛 6:25} 但如果连必要的食物和衣服，以及明天都不该忧虑，那么那些堆积如此多的废物并将自己埋葬其中的人，何时才能有力量浮出呢？你们岂没听见保禄说：\Quote{没有当兵的会让自己的生活纠缠于世务？}\BibleRef{弟后 2:4} 我们却生活奢侈，饮食无度，为外物忍受打击，而在属天的事上却懦弱不堪。你们不知道这应许对人而言太高了吗？一个行走在地上的人不可能登上穹苍之拱。我们甚至不愿像人一样生活，反倒比禽兽更甚。\par

你们不知道我们将在怎样的审判台前站立吗？你们不考虑我们连言语和思想都要交账，却连自己的行为都不在意。\Quote{凡注视妇女而贪恋她的，}祂说，\Quote{已经在心里犯奸淫了。}\BibleRef{玛 5:28} 然而那些连一个漫不经心的眼神都要负责的人，竟不拒绝在罪中自身腐烂。\Quote{凡向弟兄说\InnerQuote{愚妄}的，应被投入地狱之火。}\BibleRef{玛 5:22} 但我们甚至用千般辱骂羞辱他们，并狡猾地谋害他们。\Quote{爱那爱自己的人，并不比外邦人强。}\BibleRef{玛 5:46} 我们却嫉妒他们。那么，当我们被命令要超越旧有的界限时，我们却用比他们更苛刻的尺度编织自己的人生，我们还有什么宽恕可言？什么借口能救我们？当我们受罚时，谁能站起来帮助我们？没有人；而是必要在哀哭切齿中，痛苦地被拖入那无光的黑暗，那祈祷也无法避免的剧痛，那无法缓解的惩罚。\par

因此，我恳求、哀求，并抱住你们的双膝，趁着我们还拥有这狭小的生命行粮，愿你们被所说的这些话刺透内心，愿你们悔改，愿你们变得更善；免得我们像那富人一样，在离世后在那个世界徒然哀叹，从此持续不治的哀号。因为即使你有父亲或儿子或朋友或任何对天主有信心的人，没有一人能拯救你，因为你的行为已经毁灭了你。因为那个审判台只凭我们的行为来判断，没有其他方式能在那里得救。\par

我说这些话，并非要使你们忧伤或令你们绝望，而是免得你们被虚妄而冷淡的盼望滋养，将信心寄托在某个人或另一个人身上，而忽略了自己应有的善行。因为如果我们懒惰，就没有义人、先知、宗徒或任何人能站在我们身边；但如果我们勤奋，有了每个人自己行为的充分辩护，我们将满怀信心地离去，并得到那为爱天主之人所预备的美善；愿我们都能达到，等等。\par

\SourceRecord{Translated by Talbot W. Chambers.；Nicene and Post-Nicene Fathers, First Series；Vol. 12.；Edited by Philip Schaff.；Buffalo, NY: Christian Literature Publishing Co.,；1889.；<http://www.newadvent.org/fathers/220142.htm>.}

% structure: p0001 source=h1 target=section reason=page-title

\section{\WorkTitle{格林多前书}第四十三篇讲道}

% structure: p0002 source=h2 target=subsection reason=relative-heading-rank; base=h2

\subsection{\BibleRef{格前 16:1}}

\BibleQuote{关于为圣人捐款的事，我怎样吩咐了迦拉达的众教会，你们也该照样做。}\par

在完成了关于教义的讲论，并即将进入更属于道德的领域时，他放下了其他一切，转而讨论诸善之首，即关于施舍的事。他并非一般地讨论道德，而只是处理了这件事后就停下来了。这显然与他处处所做的不同；因为在其他书信中，他在结尾处才讨论施舍、节制、温良、忍耐以及其他一切。那么，他为什么只在这里处理实践道德的这一部分呢？因为之前所说的话中，大部分也属于伦理性质：我是说，当他惩罚行淫者时，当他纠正那些在外邦人面前诉讼的人时，当他恐吓醉酒者和贪食者时，当他谴责叛乱者、好争辩者以及爱居首的人时，当他将那些不相称地接近奥迹的人交付给那不可容忍的判决时，当他谈论爱德时。因此，我说，最紧迫的主题，即对圣人所需的援助，他唯独在这里提及。\par

请注意他的考虑。当他已说服他们关于复活的事，并使他们更加恳切时，然后，且只在此时，他才讨论这一点。\par

的确，他以前曾对他们论及这些事，当时他说：\Quote{若我们给你们播种了属神的事物，而收获你们的属世之物，难道是大事吗？} 又说：\Quote{谁栽种葡萄园，而不吃它的果子呢？} 但因他知道这一道德成就的伟大，他并不拒绝在书信的结尾再作一次提及。\par

他称呼这捐献为\OriginalTerm{捐项}{\GreekText{λογίαν}}（即\Quote{捐献}），一开始就使事情变得容易。因为当众人一同捐献时，分派给每人的负担就变得轻省。\par

但在论及捐献之后，他并没有立即说：\Quote{你们每人要自己储存；} 虽然这当然是自然的结果；而是先说了：\Quote{我怎样吩咐了迦拉达的众教会，} 然后加上这一点，借着叙述他人的善行来激发他们的好胜心，并以叙述的形式提出。他给罗马人写信时也这样做；因为对他们，他在看似叙述他为何要去耶路撒冷时，便随之引入了关于施舍的讲论：\Quote{但现在我往耶路撒冷去，为圣人服务；因为马其顿和阿哈雅喜欢凑集一笔捐款，为接济圣人中的穷人。} \BibleRef{罗 15:25} 他只是借提及马其顿人和格林多人来激励他们；这里则借迦拉达人。因为他说：\Quote{我怎样吩咐了迦拉达的众教会，你们也照样做；} 因为他们此后若被发现不如格林多人，必会感到羞愧。他并没说：\Quote{我劝告，} 或：\Quote{我建议，} 而是说：\Quote{我吩咐，} 这更具权威。他也不是举出一个城市、两个或三个，而是整个地区；他在教义训导中也这样做，说：\Quote{如同在圣人的众教会中一样。} 因为如果这对于教义的劝服是有效的，那么对于行为的效法就更为有效。\par

% structure: p0009 source=h2 target=subsection reason=relative-heading-rank; base=h2

\subsection{\BibleRef{格前 16:2}}

2. \Quote{那么我问，你吩咐了什么？}\par

\Quote{每周的第一日，就是主日，你们每人要按自己的进项储存。} 请注意他如何甚至从时间上劝勉他们：因为确实这一日足以引导他们施舍。因此他说：\Quote{你们要回想，} 在这日你们获得了什么：所有不可言喻的祝福，以及我们生命的根源和开始，都发生在这日。但不仅在这一方面，这季节便于激发热心的仁爱，还因为它有休息和免除劳作的自由：灵魂在摆脱劳苦后，变得更愿意、更适于显示怜悯。此外，在那日参与如此伟大不朽的奥迹，也注入了极大的热忱。因此，\Quote{你们每人，} 不仅这个或那个人，而是\Quote{你们每人，} 无论贫穷或富有，妇女或男子，奴隶或自由人，\Quote{都要自己储存。} 他没有说：\Quote{让他把钱带到教会，} 恐怕他们因数目微小而羞耻；而是说：\Quote{通过逐渐增加而增多他的捐献，然后当我来到时，就把它拿出来。但目前先储存起来，} 他说，\Quote{在家里，使你的家成为一座教堂，你的小盒子成为一个宝库。成为圣洁财富的守护者，穷人自命的管家。你的仁爱之心将这司祭职赋予你。}\par

如今我们的宝库正是这一点的标记：但标记仍在，而事物本身却无处可寻。\par

3. 我知道许多会众又会在这类主题上指责我，说：\Quote{求你们，不要严厉，不要使听众不悦。体谅他们的性情；迁就听者的心意。因为这样你实在使我们羞愧；你使我们脸红。}但我不能容忍这样的话：因为连\Person{保禄}也不羞于在这类问题上不断麻烦人，说些乞丐才说的话。我确实承认，如果我说\Quote{给我}和\Quote{存放在我家}，或许我所说的话中有可羞耻之处；然而即使如此也几乎谈不上羞耻；因为我们读到：\Quote{侍奉祭台的人，}我们读到，\Quote{与祭台有份。}\BibleRef{格前 9:13}然而，或许有人会指责，仿佛我在为自己的利益立论。但现在我是为穷人恳求；不，不单是为穷人，更是为你们施予者的缘故。因此我也敢直言。因为说\Quote{给你饥饿的主吃；给他赤身裸体的穿上衣服；接待他这个陌生人}有什么可羞耻的呢？你的主在全世界面前这样说并不羞耻：\Quote{我饿了，你们没有给我吃，}他是一无所缺、一无所求的。而我却要羞耻、犹豫吗？不要这样。这种羞耻来自魔鬼的陷阱。那么我不再羞耻，反而要大胆地说：\Quote{给穷人施舍；}我要比穷人自己更大声地说。诚然，如果有人能指出并证明，说这些话是为了吸引你们归向我们，并假借穷人之名自己获利，那么这种行为，我说不但是可耻，甚至该受千万雷霆；活着对他们都是过于宽恕。相反，如果靠着天主的恩宠，我们丝毫不为自己麻烦，反而把福音不花钱地传给你们；虽然决不像\Person{保禄}那样劳苦，但满足于自己的——我将以全部坦率说：\Quote{给穷人施舍；}是的，我不会停止这样说，对那些不行施舍的人，我将是严厉的控告者。因为如果我是将军，拥有士兵，我不会羞于为我的士兵要求食物：因为我热切关心你们的救恩。\par

4. 为使我的论证更有力、更有效，我将以\Person{保禄}为同伴，像他那样讲论说：\Quote{你们各人要照自己所得兴旺的程度，自己储蓄。}现在也要注意他如何避免成为负担。他没有说\Quote{这么多}或\Quote{那么多}，而是说\Quote{照自己所得兴旺的程度}，无论多少。他也没有说\Quote{任何人赚了多少}，而是说\Quote{照自己所得兴旺的程度；}表明供应来自天主。不仅这样，而且他不要求他们一次全部存起来，使他的建议变得容易：因为一点点积攒起来，使人感觉不到负担和花费。这里你也看到他不要求他们立即拿出，而是给他们充分时间的原因；他加上原因说：\Quote{免得我来的时候才现凑}；意思是，免得到了该捐献的季节你们才被迫收集。这也在不寻常的程度上鼓励了他们：对\Person{保禄}的期待必然使他们更加热切。\par

% structure: p0015 source=h2 target=subsection reason=relative-heading-rank; base=h2

\subsection{格林多前书 16:3}

\Quote{及至我来到，你们所认可的人，我就打发他们带着信，把你们的慷慨送到耶路撒冷去。}\par

他没有说\Quote{这个人}和\Quote{那个人}，而是说\Quote{你们所认可的人}，即你们所选的人，这样就使他的服务免于嫌疑。因此他把选择运送者的投票权留给他们。他远不会说\Quote{捐款是你们的，但选择运送者的特权不是你们的。}然后，为了不让他们以为他完全不在，他加上自己的书信，说：\Quote{你们所认可的人，我就打发他们带着信去。}好像他在说：我也要和他们一起，通过我的书信分享这一服务。他没有说\Quote{我要打发这些人去运送你们的施舍，}而是说\Quote{你们的慷慨；}以表示他们在行大事；表明他们自己也是受益者。并且在别处，他称之为\Quote{祝福}和\Quote{分享。}\BibleRef{格后 9:5} 一个是为了不使他们懈怠，另一个是为了不使他们骄傲。但他从未称它为\Quote{施舍}。\par

% structure: p0018 source=h2 target=subsection reason=relative-heading-rank; base=h2

\subsection{格林多前书 16:4}

\Quote{如果我也应当去，他们就与我同去。}\par

这里他又劝勉他们慷慨。如此说：\Quote{如果这么多，}他说，\Quote{以至于需要我亲自去，我也不推辞。}但他起初并没有承诺这一点，也没有说\Quote{等我来了，我亲自送去。}因为他如果一开始就这样定好，就不会如此重视它。然而后来他适时地加上了这一点。这里你就看到了他为什么不立即承诺，也不完全缄默的原因：但说了\Quote{我要打发人}之后，才加上了他自己。这里他又一次将决定权留给他们；在说\Quote{如果我也应当去}时，这取决于他们，也就是要他们的捐款足够大；大到甚至影响他的计划，使他亲自成行。\par

% structure: p0021 source=h2 target=subsection reason=relative-heading-rank; base=h2

\subsection{格林多前书 16:5}

\Quote{但我必要到你们那里去，}他说，\Quote{当我经过了马其顿以后。}这话他以前也说过；但那时带着怒气：至少他加上了\BibleRef{格前 4:19}\Quote{并且我要知道的，不是那些自高自大之人的言语，乃是他们的能力：}但这里，更温和；好叫他们甚至渴望他的来临。然后，免得他们说：\Quote{你为什么尊崇马其顿人超过我们呢？}他没有说：\Quote{当我离开时，}而是说：\Quote{当我经过了马其顿以后；因为我正经过马其顿。}\par

% structure: p0023 source=h2 target=subsection reason=relative-heading-rank; base=h2

\subsection{格前 16:6}

\Quote{但或许我会与你们同住，甚至过冬。}因为我根本不希望只是顺路看望你们，而是要在你们中间继续停留一段时间。因为他写这封信时，正在厄弗所，而且是冬天；这从他说的话可以知道：\Quote{直到五旬节，我要留在厄弗所；但此后我要往马其顿去，经过那里之后，夏天我将与你们同在；或许我甚至会与你们一起过冬。}为什么他说\Quote{或许}，而没有肯定地断言呢？因为保禄并非预知一切；这是有良好目的的。因此他没有绝对肯定地说，为的是如果事情没有成就，他可以有所凭借：首先，他先前提到这事时是不确定的；其次，圣神的能力引领他到哪里，不是他自己所愿的。这一点他在第二封书信中也表明了，为自己延迟而辩解时说：\Quote{或者我所打算的，是照着肉体打算的，以致在我有\InnerQuote{是、是}和\InnerQuote{否、否}吗？}\BibleRef{格后 1:17}\par

\Quote{好叫你们送我上路，无论我往哪里去。}这也是爱的标志，和强烈情感的体现。\par

% structure: p0026 source=h2 target=subsection reason=relative-heading-rank; base=h2

\subsection{格前 16:7}

\Quote{因为我不愿现在顺路见你们；我盼望与你们同住一些时候，如果主许可。}\par

他说这些话，既是为表明他的爱，也是为了吓唬罪人，但不是公开地，而是带着友谊的表象。\par

% structure: p0029 source=h2 target=subsection reason=relative-heading-rank; base=h2

\subsection{格前 16:8}

\Quote{但我要留在厄弗所直到五旬节。}\par

正如我们所料，他详细地告诉他们一切，像对待朋友一样告知。因为这也正是友谊的标志：说出自己为何不与他们同在、为何延迟、以及他现在所在的地方。\par

% structure: p0032 source=h2 target=subsection reason=relative-heading-rank; base=h2

\subsection{格前 16:9}

\Quote{因为有一扇又大又有功效的门为我敞开了，并且反对的人很多。}\par

如果门是\Quote{大}的，怎么会有\Quote{反对的人}呢？正因如此，反对的人才多，因为人的信德大；因为入口又大又宽。\Quote{大}门是什么意思？有许多人准备好接受信德，许多人预备好就近并归依。那里有一个宽阔的入口为我敞开，因为如今情况已到了这种地步：那些前来之人的心思已处于信德服从的成熟阶段。为此之故，魔鬼吹起的狂风猛烈，因为他看见许多人离他而去。\par

由此可见，他留下是出于两方面的需要：一方面是因为收获丰富，另一方面是因为争战巨大。\par

他同时也以此鼓舞他们，即说，从此圣言处处运行，并且迅速发芽生长。即便有许多人阴谋反对它，这也是福音进展的标志。因为那邪恶的魔鬼除了看到自己的财物被大量掠夺时，是不会那样狂暴的（\EditorialNote{玛窦福音12}）。\par

6. 那么，当我们想要成就任何伟大而崇高的事工时，不要看这带来的劳苦有多大，而要看向那收获。请看保禄，他没有因此停留，没有因此退缩，尽管\Quote{反对的人很多}；而是因为\Quote{有一扇又大又有功效的门}，他奋勇前进，坚持到底。是的，正如我所说，这正是魔鬼被剥夺的标志；因为，请相信，它不是因为微小而平庸的事工才惹怒那邪恶的怪兽。所以，当你看见一个义人做着伟大而卓越的事工，却遭受无数患难时，不要惊奇；相反，如果魔鬼受了那么多打击却保持安静、温顺地承受伤口，那才更值得惊奇。正如一条蛇，如果不断地被刺戳，就会变得凶猛并扑向刺它的人，你对此不会感到惊讶。没有一条蛇像魔鬼那样向你猛扑，它跳起来攻击所有的人；它像蝎子一样竖起毒刺，挺身而立。因此，不要让这事扰乱你；因为一个从战争、胜利和杀戮中归来的人，必然浑身是血，而且常常也受了伤。那么，当你看见一个人行哀矜，并做其它无数善工，从而削减魔鬼的权势，然后陷入诱惑和危险之中时，不要因此烦恼。这正是他陷入诱惑的原因，因为他狠狠地打击了魔鬼。\par

\Quote{天主怎么会允许这事呢？}你会说。为的是让他更显著地得到冠冕；为的是让那一个受到更严重的打击。因为当一个人行了善事之后受苦，而且是严重的苦，却仍然不断感恩，这对魔鬼是一个打击。因为，即使在我们的境遇顺遂时显出慈悲并坚守德行，这是一件伟事；但在严重的灾难中仍不放弃这崇高的行为，这更是伟大；这样的人可真正说是为天主而做的。所以，亲爱的诸位，即使我们处于危险中，即使我们遭受极大的痛苦，让我们更热切地致力于为德行的劳作。因为现在根本不是报应的时候。\par

那麼，我們在此切莫要求我們的冠冕，免得當冠冕按時來臨時，我們減少了賞報。因為就如工匠的情形：那些自食其力而工作的人得到更高的工資；而那些靠僱主供養的人，工資則被削減不少；同樣，對於聖人也是如此：那廣行善事並忍受極大苦難的人，他的賞報絲毫無損，並且得到更豐厚的報酬，不僅為他所做的善事，也為他所受的苦難。但那在此享受安逸和奢華的人，在那裏卻沒有如此光輝的冠冕。所以，我們切莫在此尋求我們的賞報。反而在我們行善卻受苦的時候，尤其要喜樂。因為天主在來世為我們存留的，不只是我們善行的賞報，也包括我們試探的賞報。\par

為了更清楚地說明我的意思：假設有兩位富有的慈善家，他們都施捨給窮人；然後一位繼續富有，享受一切順遂；另一位卻陷入貧窮、疾病和災禍，但仍感謝天主。當這兩位離世到另一個世界時，誰會得到更大的賞報？豈不是顯然是那患病和處於逆境的人，因為他雖行善受苦，卻沒有按人性的軟弱而感覺？我想每個人都明白這一點。事實上，這才是金剛石般的雕像，這才是體貼的僕人。\BibleRef{玛 25:21}但如果我們不應當為了天國的希望而行善，而只是因為這是天主所喜悅的（這超越任何天國），那麼，那因為在此沒有得到賞報而對德行變得更加懈怠的人，該當如何呢？\par

所以，當我們看到某個曾經款待寡婦、經常設宴的人，他的房子被火燒毀，或遭受其他類似的災禍時，我們不要煩惱。是的，正因為此事，他將得到他的賞報。因為約伯之所以受讚賞，與其說是由於他的施捨，不如說是由於他後來的苦難。為此，他的朋友們也不被看重，被視為無足輕重，因為他們尋求今世的賞報，並以此為根據來判斷這位義人。所以，我們不要在此尋求我們的回報；不要變得貧窮匱乏；因為當天堂以及超越天堂的事物被應許時，卻還注視著地上之物，實在是極其卑劣的行徑。我們絕不可這樣做；無論遭遇何事，我們都要緊緊持守天主的誡命，並聽從真福保祿的教導。\par

7. 讓我們在家中為窮人做一個小錢箱；放在你祈禱的地方附近；每當你進去祈禱時，先放下你的施捨，然後再獻上你的祈禱；就如你不願以未洗的手祈禱一樣，也不要沒有施捨就祈禱；因為即使掛在床頭的福音書，也不及為你積攢的施捨重要；如果你掛著福音書卻什麼也不做，它不會帶給你這麼大的益處。但如果你有這個小錢箱，你就有了對抗魔鬼的防禦，你就為你的祈禱增添了翅膀，你就使你的家成為聖潔，因為你在那裏為君王\BibleRef{玛 25:34}積存了食物。為此，也讓這個小錢箱放在床邊，這樣夜晚就不會被幻象困擾。只是不可投入任何不義之財。因為這是愛德；而愛德絕不可能出自鐵石心腸。\par

你是否也願我提及你應從哪些資源中進行捐獻，以便在這方面也使這種貢獻變得容易？例如，手工匠人，如鞋匠、皮革匠、銅匠或其他工匠——當他賣出他的任何一件製品時，讓他將售價的初果獻給天主；讓他從自己的份額中投入一小部分在這裏，即使份額相當微薄，也要將一部分歸於天主。因為我並不要求什麼大事；而是如同猶太人中那些充滿無數罪惡的幼稚之人一樣，我們這些嚮往天堂的人，也當投入同樣的份額。我說這話，並非制定法律，也不禁止更多，而是建議捐獻不少於十分之一。這件事不僅在售賣時要實行，在購買和收取報酬時也要如此。那些擁有土地的人，也要對他們的租金遵守這條法律；是的，讓所有以正當方式獲得收入的人都遵守這條法律。因為對那些放高利貸的人，以及那些以暴力對待他人、利用鄰居的災難謀取私利的士兵，我與他們無關。因為天主不會從那裏接受任何東西。但這些話我是對那些以正義勞動積累財產的人說的。\par

是的，如果我們養成了這種習慣，日後若忽略了這條規則，我們的良心就會刺痛；過了一段時間，我們甚至不會覺得這是難事；逐漸地，我們會達到更大的德行，通過練習輕視財富，並拔除萬惡的根源，我們將在平安中度過今生，並獲得來生；願我們眾人都能有分於那來生，等等，等等。\par

\SourceRecord{Translated by Talbot W. Chambers.；Nicene and Post-Nicene Fathers, First Series；Vol. 12.；Edited by Philip Schaff.；Buffalo, NY: Christian Literature Publishing Co.,；1889.；<http://www.newadvent.org/fathers/220143.htm>.}

% structure: p0001 source=h1 target=section reason=page-title

\section{格林多前书讲道 第44篇}

% structure: p0002 source=h2 target=subsection reason=relative-heading-rank; base=h2

\subsection{格林多前书 16:10}

\Emph{如今，若弟茂德到你们那里去，} \Emph{要叫他在你们那里无所畏惧。}\par

或许有人认为，在这劝言中有些与\Person{弟茂德}的勇气不相称之处。但说这话并非为了\Person{弟茂德}，而是为了听者的缘故：免得他们因谋害他而伤害自己，因为他方面，他总是置身于危险的道路上。\par

\Quote{因为儿子怎样服事父亲，}他说，\Quote{他也照样与我一同为福音劳碌。} \BibleRef{斐 2:22} 但恐怕他们因为对门徒的鲁莽，也进而攻击老师，变得更坏，他便从远处警告他们，说：\Quote{使他可以在你们那里无所畏惧；} 意思是，使那些绝望的人没有一个起来反对他。因为他或许有意要责备他们关于保禄所写的事：保禄确实自称为此而打发他。\Quote{因为我打发弟茂德到你们那里去，}他说，\BibleRef{格前 4:17} \Quote{他必提醒你们我在基督内的行径，正如我在各处各教会中所教导的。} 这样，免得他们倚仗自己的高贵出身和财富、民众的支持以及外来的智慧，而攻击他、唾弃他、谋害他，因他的责备而恼怒；或者免得他们因老师的责备而向他寻求满足，如此惩罚别人。因此他说：\Quote{使他在你们那里无所畏惧。} 好像他说：\Quote{不要对我提外面的人，外邦人和不信者。我要求你的，正是你们，为你们我也写了整封书信，} 就是那些在开头时他曾吓唬的人。所以他说：\Quote{在你们那里。}\par

然后，他基于弟茂德的职分，将他陈明为一个完全可信赖的人，说：\Quote{因为他作主的工作。} 这就是说，\Quote{不要看，}他说，\Quote{他既不富有，也未受高等教育，也不年老：而要看他负有什么命令，他在做什么工作。\InnerQuote{因为他作主的工作。}} 这代替了一切高贵、财富、年岁和智慧。\par

他不仅满足于此，还加上：\Quote{我也一样。} 又在前面某处说：\Quote{他是我的爱子，在主内忠信；他必提醒你们我在基督内的行径。} 既然他既年轻，又被单独托付改善如此众多的人民，这二者都使他易被轻视，他便加上，正如我们所料，\par

% structure: p0008 source=h2 target=subsection reason=relative-heading-rank; base=h2

\subsection{格林多前书 16:11}

\Quote{所以，谁也不要轻视他。} 他不仅要求他们做到这一点，还要求更大的尊敬；因此他也说：\Quote{要平安地送他前行；} 意思是，毫无恐惧，不惹起争斗或争端，不结仇或怨恨，而要向他表示一切顺服，如同对待师长。\par

\Quote{使他可以到我这里来，因为我正期待他和弟兄们一同来。} 这也是一个要恐吓他们的人所说的话。也就是说，为使他们的行为更谨慎，知道无论\Person{弟茂德}受怎样的待遇，一切都会告诉他，所以他加上了\Quote{因为我期待他。} 此外，借此他也表明\Person{弟茂德}值得他们信赖；因为他正要出发时还等候他；也表明他对他们的爱，看来他为了他们的缘故，打发走了对他如此有用的人。\par

% structure: p0011 source=h2 target=subsection reason=relative-heading-rank; base=h2

\subsection{格林多前书 16:12}

\Quote{至于阿颇罗兄弟，我曾再三劝他同弟兄们一起到你们那里去。}\par

此人似乎既受过良好教育，也比\Person{弟茂德}年长。免得他们因此说：\Quote{他为什么不打发那个成年人，却打发这个年轻人代替他？} 请注意他如何也缓和了这点，既称他为兄弟，又说曾再三劝他。因为免得他看来似乎给\Person{弟茂德}的尊荣比给那人的更多，而更加抬举他，因此没有打发那人去，以致激起他们的嫉妒更加猛烈，他便加上：\Quote{我曾再三劝他来。} 那么，难道他没有让步，没有同意吗？难道他抗拒，表现出好争吗？他没有这样说，而是为了不引起对他的偏见，也为自己辩护，他说：\Quote{他现在完全不愿意来。} 然后，免得他们说这一切都是借口和托辞，他接着说：\Quote{但他一有机会，就会到你们那里去。} 这既为那人辩护，又因给予盼望，使那些渴望见他的人得到安慰。\par

2. 随后，他指出他们不应在师长身上，而应在自己身上寄予救恩的希望，他说：\par

% structure: p0015 source=h2 target=subsection reason=relative-heading-rank; base=h2

\subsection{格林多前书 16:13-14}

\Quote{你们要警醒，在信德上站稳。}\par

不是在于外来的智慧，因为在那里不能站立，只会被冲走；正如同你们可以在\Quote{在信德上} \Quote{站稳。} \Quote{你们要振作，要刚强。} \Quote{你们所做的一切，都要以爱德行之。} 他说这些话时，似乎是在劝告；但他是在责备他们懒惰。因此他说\Quote{警醒，}好像他们在睡觉；\Quote{站稳，}好像他们摇摆不定；\Quote{振作，}好像他们胆小如鼠；\Quote{你们所做的一切，都要以爱德行之，}好像他们分裂不和。第一个警告是针对欺骗者，即\Quote{警醒，} \Quote{站稳；}其次，是对付谋害我们的人，\Quote{振作；}第三，是对付制造派别、企图分裂的人，\Quote{你们所做的一切，都要以爱德行之；}这爱德是\Quote{圆满的联系，}是一切美善的根源和泉源。\par

但是\Quote{凡事都应在爱德中？}是什么意思？他说：\Quote{无论谁责备，}他说，\Quote{或治理或被治理，或学习或教导，一切都当在爱德中：}因为事实上，上面所提的一切都源自忽略爱德。如果爱德没有被忽略，他们就不会自高自大，就不会说：\Quote{我是属\Person{保禄}的，我是属\Person{阿颇罗}的。}如果爱德存在，他们就不会到外邦人面前告状，甚至根本不会打官司。如果爱德存在，那个臭名昭著的人就不会娶他父亲的妻子；他们就不会轻视软弱的弟兄；他们中间就不会有纷争；他们就不会因神恩而自夸。因此他说：\Quote{一切都当在爱德中行。}\par

% structure: p0019 source=h2 target=subsection reason=relative-heading-rank; base=h2

\subsection{格林多前书 16:15-16}

3. \Quote{如今我劝勉你们，弟兄们；——你们知道\Person{斯特法纳}的家，是\Place{阿哈雅}的初结果实，并且他们献身于服务圣人。}\par

在开头他也提到这个人，说：\Quote{我也给\Person{斯特法纳}的家施了洗。}如今他称他为\Quote{初结果实}，不仅是\Place{格林多}的，也是全\Place{希腊}的。这也不是小小的赞词，因为他是最先来到基督面前的人。因此，在\WorkTitle{罗马书}中，因这个缘故称赞某些人时，他说：\BibleQuote{他们也比我先在基督内。}\BibleRef{罗 16:7}他并没有说他们是首先相信的，而是说他们是\Quote{初结果实}；暗示他们连同信仰也显出了极其卓越的生活，各方面证明自己堪当，如同果实一样。因为初结果实应当比其他那些它作为初果的东西更好：\Person{保禄}用这一表达也把这种称赞归给了这些人：就是说，他们不仅有真诚的信德，而且展现了极大的虔诚、德行的顶峰，以及在施舍上的慷慨。\par

不仅从此处，而且从另一话题他也表明他们的虔诚，即他们也使全家充满了虔敬。\par

而且他们也在善工上繁盛，他通过下文表明，说：\Quote{他们献身于服务圣人。}你们听，对他们的好客之赞词何其多？因为他没有说\Quote{他们服务}，而是说\Quote{献身于}：他们完全选择了这种生活，这是他们一直在忙碌的事业。\par

你们也要服从这样的人，也就是说，\Quote{你们要在金钱支出和亲自服务上与他们分担：与他们一同分享。}因为当有同伴时，他们的劳苦将变得轻省，他们积极行善的成果将扩展到更多。\par

他不仅说\Quote{作同工}，而且补充说\Quote{他们指示什么，你们就当服从}，暗示最严格的服从。而且为了不显得偏袒他们，他加上\Quote{也帮助所有在工作和劳苦中帮助的人。}\Quote{让这个，}他说，\Quote{成为一般规则：因为我不是专门说他们，但如果有任何人像他们一样，让他也享有同样的益处。}因此当他开始称赞时，他呼吁他们自己作证，说：\Quote{我劝勉你们，你们知道\Person{斯特法纳}的家。}\Quote{因为你们自己也意识到，}他说，\Quote{他们如何劳苦，不需要从我们这里学习。}\par

% structure: p0026 source=h2 target=subsection reason=relative-heading-rank; base=h2

\subsection{格林多前书 16:17}

\Quote{但我因\Person{斯特法纳}、\Person{福图纳图斯}和\Person{阿哈依科}的到来而欢喜，因为你们所欠缺的，他们补足了。}\par

% structure: p0028 source=h2 target=subsection reason=relative-heading-rank; base=h2

\subsection{格林多前书 16:18}

\Quote{因为他们使我的灵和你们的灵得了安息。}\par

因此，既然他们很自然会对这些人极为恼怒，因为正是他们前来向他报告了关于分裂的一切，并且也是通过他们写信询问关于童贞女和已婚者的问题——请注意他如何缓和他们的情绪；一方面在书信开头说：\Quote{因为\Person{克罗厄}家的人向我报告了}；这样既隐藏了这些人又提出了另一些人（因为似乎后者是通过前者提供信息的）；另一方面在这里又说：\Quote{他们补足了你们的欠缺，并使我的灵和你们的灵得了安息}，表明他们是代替所有人来的，并选择为你们承担如此长途的旅行。那么，他们这独特的赞词如何成为共同的？\Quote{如果你们藉着对他们的善意，来弥补我因你们所缺欠的；如果你们尊重他们，接待他们，与他们一同行善。}因此他说：\Quote{你们要承认这样的人。}在称赞那些来的人的同时，他也把另一些人包含在称赞之中，即派遣者与受遣者一起：他说：\Quote{\InnerQuote{他们使我的灵和你们的灵得了安息，因此你们要承认这样的人}，因为为了你们的缘故，他们离开了家乡。}你们看出他的体贴吗？他暗示他们不仅使\Person{保禄}受惠，也使\Place{格林多}人受惠，因为他们自身就承载着整个城市。这件事既增加了他们的信誉，也不允许其他人与他们分离，因为在他们身上，全城的人都呈现在\Person{保禄}面前。\par

% structure: p0031 source=h2 target=subsection reason=relative-heading-rank; base=h2

\subsection{格林多前书 16:19}

\Quote{\Place{亚细亚}的众教会问候你们。}他不断地通过问候使肢体联合并合而为一。\par

\Quote{\Person{阿桂拉}和\Person{普黎史拉}在主内多多问候你们；}——因为他与他们同住，是个制帐篷的——\Quote{并问候在他们家里的教会。}这件事也不是小小的优点，即他们将自己的家变成了教会。\par

% structure: p0034 source=h2 target=subsection reason=relative-heading-rank; base=h2

\subsection{格林多前书 16:20}

4. \Quote{诸位弟兄都问候你们。你们要以圣吻彼此问候。} 他仅仅在此处添加了\Quote{圣吻}。原因是什么？他们因彼此说\InnerQuote{我是属保禄的，我是属阿颇罗的，我是属刻法的，我是属基督的}、因\InnerQuote{一个饥饿，另一个醉酒}、因他们之间有争端、嫉妒和诉讼，而彼此大为不和。此外，因神恩而有许多嫉妒和极大的骄傲。既然他藉劝勉将他们联结起来，自然吩咐他们也使用圣吻作为合一的方式：因为这合而为一，产生同一身体。当没有欺骗和虚伪时，这就是圣的。\par

% structure: p0036 source=h2 target=subsection reason=relative-heading-rank; base=h2

\subsection{格林多前书 16:21}

\Quote{我保禄亲笔问候。} 这暗示这书信是以极大的认真写成的；因此他加上，\par

% structure: p0038 source=h2 target=subsection reason=relative-heading-rank; base=h2

\subsection{格林多前书 16:22}

\Quote{若有人不爱我们的主耶稣基督，就让他受诅咒。}\par

藉此一语，他震慑了所有人：那些使自己的肢体成为娼妓肢体的人；那些因祭偶像之物在弟兄们的路上设置绊脚石的人；那些以人的名字自称的人；那些不信复活的人。他不仅震慑，还指出了德行之路与万恶之源，即：当我们对祂的爱变得强烈时，没有任何罪不被它熄灭和驱逐；而当它太弱时，则使同样的罪滋长起来。\par

\Quote{ \Emph{玛兰阿塔}。} 为什么使用这个词？又为什么用希伯来语？既然傲慢是万恶之因，而这傲慢是由外来的智慧产生的，并且这是万恶的总和与实质，尤其使格林多人分裂；在压制他们的傲慢时，他甚至连希腊语都不用，而用希伯来语：表明他非但不以那种纯朴为耻，反而非常热切地拥抱它。\par

但是\Quote{ \Emph{玛兰阿塔}}是什么意思？\Quote{我们的主已经来了。} 那么他为什么特别使用这个短语？为确认降生奥迹的教理：在众多主题中，他特别从这类主题中汇集了那些作为复活种子的论证。不仅如此，也是为责备他们；仿佛他说：\Quote{众人的共同主屈尊降下至此，而你们仍处于原状，仍停留在你们的罪中吗？你们难道不因祂的爱之超越、祂恩赐的冠冕而激动吗？是的，只考虑这一点，} 他说，\Quote{就足使你们在一切德行中进步，并且你们将能熄灭一切罪。}\par

% structure: p0043 source=h2 target=subsection reason=relative-heading-rank; base=h2

\subsection{格林多前书 16:23}

\Quote{我们的主耶稣基督的恩宠与你们同在。}\par

这正如一位老师，不仅以劝言帮助，也以祈祷帮助。\par

% structure: p0046 source=h2 target=subsection reason=relative-heading-rank; base=h2

\subsection{格林多前书 16:24}

\Quote{我的爱在基督耶稣内与你们众人同在。阿们。}\par

为阻止他们以为他这样结尾是奉承他们，他说：\Quote{在基督耶稣内。} 这没有任何人性或肉性的成分，而是属于一种属灵的性质。因此它是完全真诚的。因为这句话实在是一位深爱者的话。正如这样：因为他与他们在地方上分离，仿佛伸出右手，以他爱的双臂环绕他们，说：\Quote{我的爱与你们众人同在；} 正如他说：\Quote{我与你们众人同在。} 由此暗示所写之事并非出于恼怒或愤怒，而是出于眷顾，因为如此严厉的指责之后，他并没有转身离开，反而爱他们，并在他们远离时拥抱他们，藉着这些书信和著作投入他们的怀抱。\par

5. 因为纠正者应当这样做：因为那仅出于愤怒行事的人，只是满足自己的情绪；但他在纠正罪人之后也付出爱的行动，表明他所说的一切责备的话都是出于深情的爱。同样，让我们也彼此管教；让纠正者不要发怒（因为这不属于纠正，而属于情绪），也让被纠正者不要心怀怨恨。因为所做的事是治疗，而非侮辱。如果医师使用烧灼术而不被责备，而且常常如此，尽管他们完全偏离了目标；但即使在痛苦中，接受烧灼和截肢的人仍将引起这种痛苦的人视为恩人；那么，接受责备的人更应当如此，如同对待医师一样听从纠正者，而不是像对待仇敌。让我们这些责备的人也以极大的温和、极大的谨慎靠近。如果你们看见一个弟兄犯罪，如同基督所命令的，不要公开责备，而\BibleQuote{只在你和他之间}：\BibleRef{玛 18:15} 不当他跌倒时责备或侮辱他，而要带着痛苦和柔情的心。并让自己也准备好接受责备，如果你在任何事上犯错。\par

为使我说的话更清楚，让我们设想一个事例来检验我们的原则。因为天主禁止我们真的需要这样一个例子。假设一个弟兄与一位贞女同住一屋，保持着荣誉和贞洁，却仍不能完全逃避恶评。如果你听到他们同住的事，不要轻蔑，也不要说：\Quote{怎么，他没有理智吗？他自己不知道什么对他有益吗？无端获得爱，但不要无端招来仇恨。为什么我要承担无端的敌意？} 这些是野兽的痴言，甚至说是魔鬼的话：因为这样做是为了纠正弟兄而被人憎恨，并非无缘无故，而是为了极大的祝福和说不尽的冠冕。\par

但若你说：\Quote{甚么？他难道没有理智吗？}你将听我说：他没有；他被自己的情欲灌醉了。因为在异教法庭上，受冤屈的人不得在激怒时为自己辩护（尽管那种同情并无过失），何况被恶习控制的人呢？因此我说，尽管他多有智慧，他的心却不清醒。因为有什么比达味更有智慧呢？他曾说：\Quote{祢的智慧中黑暗和隐密的事，祢已让我知晓。}\EditorialNote{第267页第6行，参照七十子译本第6行}但当他用不义的目光注视那个士兵的妻子时，按他自己论及在怒海上航行的人所说的\BibleRef{咏 106:27}：\Quote{他所有的智慧都被吞没了；}他需要别人纠正他，甚至没有察觉自己身处何等邪恶之中。因此，他悲叹自己的过犯说：\Quote{如同沉重的重担，他们沉重地压在我身上：我的伤口发臭流脓，因我的愚昧。}\BibleRef{咏 37:5}所以，犯罪的人没有理智。因为他醉了，在黑暗中。不要说这些话，也不要加上另一句：\Quote{我毫不在意。\InnerQuote{因为各人要担自己的担子。}}\BibleRef{迦 6:5}不，这也成为对你的严厉指控，因为你看见一个人陷于错误而不去挽回他。因为若不按犹太人的法律\BibleRef{出 23:4}，连仇人的牲畜也不可忽视；那么，谁若不轻视仇人的牲畜、也不轻视敌人灭亡的灵魂，却轻视朋友的灵魂，他将得到什么宽恕呢？\par

是的，他有理智也不足以成为我们的借口：因为我们自己，经过许多和多种劝告，也没有做到良善，也没有对自己有用。那么，对于那陷于错误的人，也要记住这一点：他更有可能从你而不是从自己那里得到最好的劝告。\par

不要说：\Quote{但这些事与我何干？}你要畏惧那第一个说这话的人；因为那句话：\Quote{难道我是我兄弟的看守者吗？}\BibleRef{创 4:9}与此相同。把我们自己身体的事视为不相干，这是一切邪恶之母。你说什么？你不管你的兄弟吗？那么谁管他呢？那不信者，为他欢喜、责备他、侮辱他吗？还是那魔鬼，催促他、陷害他？\par

这从何而来？他说：\Quote{我怎知道我能成就什么呢？尽管我说出和劝告正确的事。}但是你怎么清楚你行不出善呢？唉，这又是极端的愚昧，因为结局尚在暗昧之中，却招致承认冷漠的诸多责备。然而，预先看见未来的天主，常常说话却行不出善；但祂也并未因此放弃；而且，祂知道祂连人都不能说服。那么，如果预先知道不会得到什么益处的那一位，仍然不停止纠正的工作，你完全不知道未来，却灰心丧气、麻木不仁，你还有什么借口呢？是的，许多人都因屡次尝试而成功：他们越是绝望，越能达到目的。即使你得不到什么益处，你也已尽了自己的本分。\par

那么，不要不近人情，不要冷酷无情，也不要漠不关心：因为这些话出于残忍和冷漠，从以下事实可以显明：即，当你身体有一个肢体疼痛时，你为什么不说\Quote{我何在乎？}然而又怎能肯定，若适当照顾，它就能复原呢？但你并没有留下任何未做之事，即使你无益，也不得因遗漏任何应做之事而自责。因此我问，我们如此照顾我们身体的肢体，却忽略基督的肢体吗？不，这样的事怎能得到宽恕？因为若我对你说\Quote{要照顾你自己的肢体}不能感动你，为了使你哪怕因恐惧而变好，我提醒你想起基督的身体。但看见祂的肉体腐烂却忽视它，怎能不是一件可怕的事呢？若你有一个奴隶或一头驴长了坏疽的疮，你也不忍心忽视它：但你看见基督的身体长满了疥癣，却匆匆走过？难道你不认为这样的事应受无数雷霆之击吗？因为这一切都颠倒了，因了我们的不近人情，因了我们的漠不关心。因此现在，我恳求你，让这种残忍从我们中间除去。\par

6. 你要亲近我所指的那人，如同与一位贞女同住，并说一些对你兄弟的小小赞美，从他拥有的其他优点中提取。用你的称赞如同温水般敷他，从而缓和伤口的肿胀。也要谈到你自己是可怜可悲的；指责整个人类种族；指出我们都在罪恶中；请求宽恕，说你正承担对你而言过于重大的事，但爱催逼你冒险做一切。然后，在给予劝告时，不要以专横的方式，而要以兄弟的方式。当你通过这些方法消了肿胀，平息了因即将来临的锐利责备而产生的痛苦，当你反复恳求他不要生气：当你用这些事束缚了他，然后才用刀；既不要太紧迫，也不要放松；使他既不因一方而逃脱，也不因另一方而轻视。因为如果你不触及要害，就一无所成，如果你用力过猛，就会使他惊退。\par

因此，即使在此之后，当指责即将发出时，仍要在你的责备中再次混入赞许。既然这种行为本身不能成为称赞的对象（因为与一位童贞女同居一室并非可称赞之事），那么就让行此事之人的动机成为你达成此目的的论题，并说：\Quote{我确实知道你是为了天主这样做，而且那可怜妇人的孤苦无依触动了你的眼目，使你向她伸出援手。} 即使他可能并非出于此意，你仍要如此说；之后再添加以下的话，再次自我开脱并说：\Quote{我说这些话不是要指导你，而是要提醒你。你是为了天主这样做；我也知道。但让我们看看是否因此产生了另一种恶。如果没有，就让她留在你家中，坚守这美好的目的。没有人会阻碍你。但如果由此产生的祸患超过了益处，我恳求你谨慎，免得我们急于安慰一个灵魂，却给一万人设置了绊脚石。} 不要立即加上那些使人跌倒的人应受的惩罚，而要借用他自己的见证，说：你不需要从我这里学习这些；你自己知道，\InnerQuote{若有人使这些小子中的一个跌倒}，会受到何等大的惩罚。这样，你的言语变得甘甜，平息了他的怒气，然后再施以纠正的良药。如果他再次强调她的孤苦，你甚至不要这样揭穿他的伪装，而要对他说：\Quote{不要让这类事使你害怕：你将有一个充分的理由，就是给人留下的绊脚石：因为你放弃这个目的，并非出于冷漠，而是出于对他们的关心。}\par

让你的建议内容简短，因为无需过多的教导；但另一方面，要多多表达宽容，且接连不断。要不断诉诸爱的话题，掩盖你话语中的痛苦，并让他拥有全部的权力，说：\Quote{这是我的忠告和建议；但关于是否采纳，你只有判断权：因为我不强迫你，而是将整个事情交由你自己斟酌。}\par

如果我们这样处理我们的责备，我们就能轻易纠正那些犯错的人；而我们现在的所作所为，确实更像是野兽或无理性之物的行为，而非人的行为。因为如果有人现在发现某人犯了这类错误，他们根本不与本人交谈，而是像喝醉了的愚蠢老妇人一样，彼此交头接耳。那句谚语\Quote{无缘无故得爱，却不要无缘无故惹恨}在他们看来毫无意义。但是，当他们想要说坏话时，他们不介意\Quote{无缘无故被恨}，更确切地说，我应该说\Quote{受惩罚}；因为由此产生的不仅是恨，还有惩罚。但是当需要纠正时，他们却拿这个和无数其他借口来推脱。然而，当你讲坏话、诽谤时，正是考虑这些事情的时候。我指的是那句\Quote{不要无缘无故被恨}、\Quote{我无能为力}和\Quote{这与我无关}。但事实上，在前一种情况下，你充满激烈而无益的好奇心，毫不介意恨和无数恶果；但当你要关心你弟兄的救恩时，你却乐于做那种不多管闲事、不冒犯他人的人。然而，从恶言中产生的仇恨，既有来自天主的，也有来自人的；但这对你来说并不重要；但是通过私下给予忠告和那样的责备，他反而和天主都会成为你的朋友。即使他恨你，天主反而更加爱你的。事实上，他甚至不会像因你的恶言而导致的恨那样恨你；而在那种情况下，他会像躲避仇敌一样躲避你，而现在他会认为你比任何父亲都更可敬。即使表面上他对此感到恼火，内心和私下里他也会非常感激你。\par

7. 因此，要牢记这些事，让我们关心自己的肢体，不要彼此磨砺舌头，也不要说\Quote{伤人的话}，抹黑邻人的名誉，如同在战争和战斗中互相攻击。因为当舌头沉醉，在比狗肉更不洁的餐桌上大快朵颐；当它贪嗜鲜血，倾泻污秽，使口成为阴沟的通道，甚至比那更可憎时，那么禁食或守夜又有什么益处呢？因为从那里出来的东西污染了身体，但从舌头出来的东西常常窒息灵魂。\par

我说这些话，并非为那些被冤枉遭受恶名的人担忧；因为当他们高尚地承受他人所说的话时，他们甚至配得冠冕；而是为你们这些如此说话的人担忧。因为对于被冤枉遭受恶名的人，圣经宣告他是\Quote{有福的}；但恶言者，圣经将他们从神圣奥迹中驱逐，甚至从圣殿前院中逐出。因为经上记载：\BibleRef{咏 100:5} \BibleQuote{暗中诋毁邻舍的人，我必驱逐出去。} 他又说，这样的人不配诵读圣经。因为，祂说：\BibleRef{咏 49:16} \BibleQuote{你为何宣扬我的律法，将我的盟约挂在嘴上？} 然后，祂附加原因说：\BibleRef{咏 49:20} \BibleQuote{你坐着毁谤你的弟兄。} 在这里，祂并没有明确指出所说的是真实还是虚假。但他在别处也将这点纳入禁令中：祂暗示，即使你说的是实话，你也不该说出这类事情。因为，祂说：\Quote{你们不要判断人}，免得你们受判断：\BibleRef{玛 7:1} 因为那毁谤税吏的人也受到了谴责，尽管他控告邻人的话是真实的。\par

\Quote{那么，}你将会说，\Quote{若有一个人胆大而污秽，难道我们不该纠正他吗？难道不该揭发他吗？}我们既应揭发，也应纠正：但要按我先前提到的方式。但如果你以斥责的方式去做，小心不要效法那位法利塞人而导致你落入他的状态。因为这样做毫无益处：对你说话者无益，对听你的人无益，对被指控的人无益。后者反而变得更无所顾忌：因为当他未被发现时，他尚有羞耻感；但一旦他显明而众所周知，他也抛弃了那感觉加给他的约束。\par

而听者也会受到更大的损害。因为无论他是否自觉有善行，他都会因对别人的指控而变得骄傲自满；或者如果他自觉有错，他就会更热衷于作恶。\par

第三，说话者自己既会招致听者的恶评，也会激起天主对他更大的愤怒。\par

因此，我恳求你们，让我们弃绝一切不入耳的话语。若有什么话有益于建立德性，就让我们说这话。\par

但你想要对那人报复吗？为何不惩罚他反而惩罚你自己呢？不，你，如此热切地向那些惹恼你的人寻求补偿，当像保禄所推荐的那样去报复。\Quote{如果你的仇人饿了，就给他吃；如果他渴了，就给他喝。}\BibleRef{罗 12:20}但你若不这样做，反而图谋害他，你便是把剑指向自己。\par

因此，若那人说坏话，当以赞美和嘉许回应他。因为这样你既能向他报复，又能使你自己免于恶意的猜疑。因为那听见说自己不好而感到痛苦的人，被认为是出于某种罪恶的自觉；但那嘲笑所说的话的人，则显示出他最无可争辩的证据，表明他自觉没有任何恶事。\par

既然你既无益于听者，也无益于自己，也无益于被指控的人，反而把剑指向自己，即使从这些考量中，你也当学习更节制。因为人本当受天国以及天主所喜悦的事所感动；但既然你性情粗俗，像野兽一样咬人，那就藉此受教吧；好让这些论证纠正你之后，你能够单单从考量天主所喜悦的事来调整自己；并且，在超越一切情欲之后，获得天上的福乐——愿天主使我们众人藉我们的主耶稣基督的恩宠和祂对人类的仁慈，得以获得这些福乐；愿光荣、权能、尊崇归于祂，及圣父和圣神，从今时直到永远。阿们。\par

\SourceRecord{Translated by Talbot W. Chambers.；Nicene and Post-Nicene Fathers, First Series；Vol. 12.；Edited by Philip Schaff.；Buffalo, NY: Christian Literature Publishing Co.,；1889.；<http://www.newadvent.org/fathers/220144.htm>.}
